\section{The full profinite group of the CM curve}
\label{part:A}

Throughout this part $K=\Q(i)$, $E/K$ is $y^2=x^3-x$,
$U=E\setminus\{O\}$, and $\pi=\pi_1(U_{\overline K},\vec v)$ is
the full geometric profinite group. We write $\Pi=\pi_1(U_K,\vec v)$,
$\Delta=\Gal(\overline K/K^{\rm ab})$, and $I=I_O\subset\pi$ for
the actual cusp inertia selected by the tangential basepoint. A packet section
is an actual continuous splitting $s_b:G_K\to\Pi$ with image in the
decomposition group of this lift of $O$. Its action on $I\simeq\wh\Z(1)$
is cyclotomic. No rationality of its completed Kummer parameter is assumed.

\begin{statement}{A.1}{Proved}{The full fixed group is inertia}
For every such packet section,
\[
                 C_\pi(s_b(\Delta))=\pi^{s_b(\Delta)}=I.
\]
The same equality holds after replacing $\Delta$ by any open subgroup of
$\Delta$. In fact, for any hyperbolic curve over any number field, the full
geometric centralizer of a point section restricted to the closed commutator
subgroup of the absolute Galois group is trivial; for a cuspidal section it is
the specified cusp inertia. These assertions use the valuation uniqueness
theorem of Mochizuki--Tsujimura, and do not use the unrefereed local section
conjecture preprint of October 2026.
\end{statement}
\begin{proof}
We give the comparison with the local theorem explicitly. The input is
\citet[Theorem A, Definition 2.2(i), Proposition 2.3(v), and
Proposition 2.4(vii)]{MT2024}. With all primes retained, every hyperbolic curve
over a finite extension $k/\Q_p$ satisfies the required resolution of
nonsingularities condition. If $J\subset G_k$ is closed and
$J\cap I_k$ has a continuous quotient isomorphic to $\Z_r$ for a prime
$r\ne p$, an actual section over $J$ fixes at most one primitive
$p$-valuation of the universal etale function field. The extra condition that
the $r$-adic cyclotomic image be open is needed for existence, not uniqueness.
The proof on printed page 54 separates these two assertions, using
Proposition 2.3(ix) for uniqueness. We will supply the valuation ourselves.

First consider the local group. Put
$J=\overline{[G_k,G_k]}$. Wild inertia $P_k$ is pro-$p$, and the tame
quotient has the form
\[
 G_k/P_k=\left(\prod_{r\ne p}\Z_r(1)\right)\rtimes\wh\Z,
\]
where an arithmetic Frobenius acts by the residue cardinality $q$.
The unramified quotient is abelian, so $J\subset I_k$.
The image of $J$ in tame inertia contains
$(q-1)\prod_{r\ne p}\Z_r(1)$: take the commutator of a Frobenius lift
with a tame element. Its image in each $\Z_r(1)$ is therefore a nonzero
closed subgroup, hence isomorphic to $\Z_r$. Projection onto that image
is the quotient required by the uniqueness theorem. This remains true for
every open subgroup of $J$, since its image has finite index in that of $J$.

Let $X/k$ be our local curve, let $L=k(\widetilde X)$ be the union of
the function fields of its connected finite etale covers, and choose a
compatible lift $\widetilde b$ of the point or cusp defining the section.
The horizontal valuation ring $A\subset L$ at this lift has residue field
$\overline k$. It is the union of the local rings at the compatible points
on the smooth compactifications of the finite covers. At cusps those covers
may ramify; this does not change the residue field. Write
$\operatorname{ev}:A\twoheadrightarrow\overline k$ for the residue map and
define
\[
 R=\operatorname{ev}^{-1}(\mathcal O_{\overline k}).
\]
This is the composite of the horizontal valuation with the unique extension
of the $p$-adic valuation to $\overline k$. It is a $p$-valuation in the
precise sense of Definition 2.2(i). It is primitive: if $\mathfrak m_A$
is the horizontal maximal ideal, then
$\mathfrak m_A=p\mathfrak m_A\subset pR$, while
$R/\mathfrak m_A=\mathcal O_{\overline k}$ and
$\sqrt{p\mathcal O_{\overline k}}=\mathfrak m_{\overline k}$.
Consequently the only prime of $R$ containing $p$ is its maximal ideal.
Furthermore
\[
                         R[1/p]=A.                       \tag{A.1}
\]
Indeed, multiply the residue of any horizontal unit by a sufficiently large
power of $p$ to make it integral; elements of $\mathfrak m_A$ already lie
in $R$. The reverse inclusion follows because $p$ is a unit in $A$.

The actual point section, or any section into the chosen cuspidal
decomposition group, preserves $A$ and its residue valuation, hence $R$.
Let $g\in\pi_1(X_{\overline k})$ centralize its restriction to $J$.
Then both $R$ and $gR$ are primitive valuations fixed by that restriction.
Uniqueness gives $gR=R$. Since $g$ fixes $p$, (A.1) implies $gA=A$.
The geometric stabilizer of this horizontal valuation is trivial at an
ordinary point: finite geometric Galois covers are etale there and their
residue fields are algebraically closed. At a cusp the stabilizer is exactly
its full geometric inertia, since decomposition and inertia agree over an
algebraically closed residue field. This proves the desired containment.
The reverse containment at a cusp holds because all cyclotomic characters
are trivial on $J$.

Now let $k_0$ be a number field and choose a finite place $v$. A chosen
prolongation identifies its decomposition group $D_v\subset G_{k_0}$ with
$G_{(k_0)_v}$. In characteristic zero the base-change map identifies the
geometric etale groups of $X_{\overline{k_0}}$ and
$X_{\overline{(k_0)_v}}$; the local arithmetic sequence is the pullback
of the global sequence along $D_v\hookrightarrow G_{k_0}$. These are
comparisons of full profinite groups, with the paths fixed. Since
\[
 \overline{[D_v,D_v]}
       \subset \overline{[G_{k_0},G_{k_0}]},
\]
an element fixed by the global commutator subgroup is fixed by the local
subgroup just treated. The local result puts it in the stated horizontal
stabilizer. In the cuspidal case the global commutator subgroup acts trivially
on all of inertia, proving equality.

Finally, if $\Delta_0$ is open in $\Delta$, choose an open subgroup
$G_{k_1}\subset G_K$ with $G_{k_1}\cap\Delta\subset\Delta_0$.
Then $\overline{[G_{k_1},G_{k_1}]}\subset\Delta_0$.
Apply the same argument over the finite extension $k_1/K$ and use that
$\Delta_0$ still acts trivially on inertia. This proves the strengthened
assertion. The argument at a cusp also verifies the step used in the
attached ANABELIAN\_149\_PUSH C.7h, while replacing its local cyclotomic
kernel by the smaller local commutator subgroup.
\end{proof}

% The residual audit is included after the basic fixed-group argument.

\begin{statement}{A.3}{Proved}{The homology fixed by an open subgroup of the abelian-field kernel}
Let $Y=Y_H$, $\overline Y$ its compactification, $S$ its geometric cusp
set, $r=|S|$, and $J=\operatorname{Jac}(\overline Y)$. Let $J_E\subset
J_{\overline K}$ be the sum of all abelian subvarieties isogenous to $E$.
Use the principal polarization of $J$ to project its rational isogeny
decomposition onto $J_E$. Let $D_E\subset J_E$ be the connected
Zariski closure of the group generated by the projected classes
$[s-s_0]$, $s\in S$, for one $s_0\in S$. Rational projection is
interpreted after multiplication by a common denominator; this does not
alter $D_E$. Put
\[
 B_H=(J_E\cap D_E^\perp)^0,
 \qquad
 W_{S,\ell}=\Q_\ell(1)^S/\Q_\ell(1)\cdot(1,\ldots,1).
\]
Then the subspace of $H_1(Y_{\overline K},\Q_\ell)$ whose vectors
are fixed by some open subgroup of $\Delta_H$ is, after a finite field
extension defining the displayed data, canonically
\[
                    W_{S,\ell}\oplus V_\ell B_H.         \tag{A.6}
\]
Here ``canonically'' means the unique Galois-equivariant splitting on
this subspace, once the natural boundary inclusion and proper-homology
projection are specified. This proves Lead 2 in its rational,
eventually-fixed sense, with this precise meaning of the sum.

Let $b=\dim B_H$. After a further finite field extension, the lines in
(A.6) stable under an open subgroup of $G_K$ are exactly as follows.
Every line in $W_{S,\ell}$ is such a line, with character $\chi_\ell$.
If $\ell\equiv1\pmod4$, there are in addition all lines in each of
the two $b$-dimensional CM eigenspaces of $V_\ell B_H$, with the two
CM characters of $E$. At $\ell=2$ and at $\ell\equiv3\pmod4$ there
are no additional $\Q_\ell$-lines. There are no lines mixing these
different character spaces.
\end{statement}
\begin{proof}
We separate the proper and the boundary questions. We first prove that
for any abelian variety $A$ over a finite extension of $K$, its
potential $\Delta$-fixed rational Tate subspace is exactly $V_\ell A_E$.
One direction follows by defining an isogeny $E^d\to A_E$ over a
finite extension: $\Delta$ fixes every torsion point of $E$, so the
corresponding finite-index subgroup of $\Delta$ fixes $V_\ell A_E$.

For the reverse direction, make the field of definition $L$ large enough
that a proposed fixed vector is fixed by $\Delta\cap G_L$, and take
$R=\operatorname{Res}_{L/K}A$. Its Tate representation is
$\operatorname{Ind}_{G_L}^{G_K}V_\ell A$. In the restriction of this
induction to $\Delta$, the orbit of the identity coset gives
$\operatorname{Ind}_{\Delta\cap G_L}^{\Delta}V_\ell A$; its invariant
vectors identify with $(V_\ell A)^{\Delta\cap G_L}$. Thus the proposed
vector occurs in $(V_\ell R)^\Delta$. Decompose $R$, up to $K$-isogeny,
into $K$-simple factors. A factor with nonzero $\Delta$-invariant
Tate module has torsion of unbounded $\ell$-power order over $K^{\rm ab}$:
scale an invariant vector into its Tate lattice and reduce it modulo
$\ell^n$ for every $n$.
The theorem of \citet[\S1, Theorem; \S3, Case II]{Ruppert1998} says that
its endomorphism algebra over $K$ is a field $F$ of degree twice its
dimension.

We explain why such a factor is geometrically $E$-isotypic, a point
specific to the imaginary quadratic field $K$. If its dimension is $d$,
the differential representation is faithful on endomorphisms in
characteristic zero and makes $\operatorname{Lie}(A)$ a nonzero
$F\otimes_\Q K$-module, faithful on $F$, of $K$-dimension $d$. If $F$ and $K$ were
linearly disjoint, this algebra would be a field of degree $2d$ over $K$,
which cannot act faithfully on a $d$-dimensional vector space. Hence
$K\subset F$. Now $F\otimes_\Q K\simeq F\times F$, each factor
having $K$-dimension $d$. The Lie representation uses exactly one factor.
Its CM type is therefore induced from one of the two embeddings of $K$.
A $K$-basis of its rational first homology identifies its rational Hodge
structure with $d$ copies of the Hodge structure of a CM elliptic curve
with multiplication by $K$. Such a curve is isogenous to $E$.
The rational Hodge isomorphism, after clearing denominators in the
lattices, is an isogeny of complex abelian varieties and hence of the
abelian varieties over $\overline K$. Projection back from the Weil
restriction proves that our fixed vector lies in $V_\ell A_E$.
This argument applied to all factors also proves equality of the full
subspaces, not merely existence of one CM constituent.

We next supply the particular Kummer input needed here. Over a finite
extension $L/K$ with all relevant endomorphisms defined, let
$P_1,\ldots,P_m\in E(L)$ be linearly independent modulo torsion over
$K=\operatorname{End}^0(E)$. The joint Kummer map on
$G_{L K^{\rm ab}}$ has image spanning $(V_\ell E)^m$ over $\Q_\ell$.
Here $K^{\rm ab}=K(E_{\rm tors})$, so this is exactly the full torsion
field, rather than an uncontrolled further extension of it.
To prove the spanning assertion, let $W$ be the span of the image.
It is stable under $G_L$ and under the commutative algebra generated by
its action. This algebra on $V_\ell E$ is $K\otimes\Q_\ell$.
For example, a Frobenius above $5$ when $\ell\ne5$, or above $13$
when $\ell=5$, generates it: its eigenvalue quotient is a non-root of
unity in $K$, and this remains so after every finite residue-degree
power. The two eigenvalues may be taken as $-1\pm2i$ and $3\pm2i$,
respectively, from the elementary point counts $8$ and $8$.
Thus, if $W$ were proper, a nonzero row
$(a_1,\ldots,a_m)\in(K\otimes\Q_\ell)^m$ would annihilate it.

The global cocycle $\sum a_j\kappa(P_j)$ would then restrict to zero
over $LK^{\rm ab}$. Its class is zero already over $L$. Indeed its
possible inflation kernel is $H^1(\Gamma_L,V_\ell E)$, where
$\Gamma_L=\Gal(LK^{\rm ab}/L)$ is abelian; a good Frobenius $g$
has $\rho(g)-1$ invertible, and the commutative cocycle identity gives
\[
 b(h)=(\rho(h)-1)(\rho(g)-1)^{-1}b(g).
\]
So this inflation group is zero. The rational Kummer map
$E(L)\otimes\Q_\ell\hookrightarrow H^1(G_L,V_\ell E)$ is injective,
by the finite-coefficient Kummer exact sequences and the finite
generation of $E(L)$. Independence of the $P_j$ over $K$ persists
after tensoring with $\Q_\ell$ and forces all $a_j=0$, a contradiction.
This proves the assertion for every $\ell$, including $2$.
More generally, for a point of $E^d$ its Kummer image spans the Tate
module of the connected algebraic subgroup it generates: apply the
argument to a basis of its coordinates over $K$, and then transport
through an isogeny. Compactness upgrades full rational span to openness
in the appropriate integral lattice. This proves the CM Kummer fact
needed below directly; it does not rely on applying a large-prime
mod-$\ell$ statement at the exceptional primes.

Over a finite $L$ splitting $S$, the generalized Jacobian gives the exact
sequence of rational Galois representations
\[
 0\longrightarrow W_{S,\ell}\longrightarrow H_1(Y_{\overline K},\Q_\ell)
       \longrightarrow V_\ell J\longrightarrow0.         \tag{A.7}
\]
The extension by the coordinate of $s-s_0$ is the extension of $J$
by $\mathbf G_m$ classified by $[s-s_0]\in J^\vee(L)$ under the
principal polarization. One can see its Tate extension class by
choosing compatible division points of $[s-s_0]$: its cocycle, paired
with $V_\ell J$ through the Weil pairing, is the corresponding Kummer
cocycle. This description also fixes the comparison map and the sign
ambiguity, which has no effect on its kernel.

We have proved that the proper image of a potentially $\Delta$-fixed
vector lies in $V_\ell J_E$. On that space $\Delta\cap G_L$ acts
trivially after increasing $L$. The connecting map of (A.7) is therefore
the family of homomorphisms obtained by pairing with the restricted
Kummer classes of the projected cusp divisors. The preceding Kummer
argument identifies their common annihilator with $V_\ell B_H$.
A nonzero homomorphism from a profinite group to a rational Tate vector
space cannot vanish on an open subgroup: otherwise its image is finite,
and the additive target has no finite subgroup. Thus taking the union
over open subgroups does not enlarge this annihilator.

The pullback of the generalized-Jacobian extension to $B_H$ is classified
by torsion points of $B_H^\vee$: this is precisely the definition of the
connected orthogonal complement of $D_E$. Its rational Tate extension
splits over a finite number field. The splitting is unique, since
$\operatorname{Hom}_{G_L}(V_\ell B_H,W_{S,\ell})=0$: at a good
Frobenius the former has weight one and the latter weight two. This
proves (A.6) with the asserted meaning and independence of a further
finite extension.

Finally define an isogeny $E^b\to B_H$ over a finite extension. If $\ell$
splits in $K$, $K\otimes\Q_\ell=\Q_\ell\times\Q_\ell$, and the
two distinct CM characters act scalarly on the two $b$-dimensional
summands. Their ratio is not finite, as the Frobenius calculation above
shows. They are also distinct from $\chi_\ell$ on every open subgroup,
by weight. Hence precisely the stated lines are stable. At an inert or
ramified prime, $K\otimes\Q_\ell$ is a quadratic field, and a stable
$\Q_\ell$-line would yield a unital homomorphism of that field into
$\Q_\ell$, since the same Frobenius generates the algebra. This is
impossible. The weight distinction rules out mixed lines and completes
the classification.
\end{proof}

\begin{statement}{A.4}{Proved}{All first-cover models, Jacobians, and cusp divisors}
Put
\[
 (a,b)\in\mathcal T=\{(0,-1),(6,1),(-6,1)\},\qquad
 E_{a,b}:y^2=x^3+ax^2+bx.
\]
Each $E_{a,b}$ has a degree-two isogeny to $E$, with kernel
$\{O,(0,0)\}$. The three geometric $S_3$-covers and three geometric
$D_8$-covers of $U$ have smooth compactifications
\[
 C_{3,a,b}:v^2=w^7+aw^4+bw,\qquad
 C_{4,a,b}:v^2=w^8+aw^4+b,                         \tag{A.8}
\]
respectively. Their maps to $E_{a,b}$ are $(w^3,wv)$ and
$(w^4,w^2v)$; compose with the indicated isogeny to $E$.
All six curves have genus three. The $C_3$ cusp set is
$\{P_0,P_\infty\}$, and the $C_4$ cusp set consists of the two
points over $w=0$ and the two over $w=\infty$.

For the $C_3$ models take $k=-1$ if $(a,b)=(0,-1)$ and $k=1$
otherwise, so $k^3=b$. For the $C_4$ models choose $c^4=b$.
The following are geometric isogeny decompositions and quotient models:
\[
\begin{array}{c|c|c}
\text{curve}&\text{elliptic complementary factor}&\text{Jacobian}\\ \hline
C_{3,a,b}&F_{3,a,k}:u^2=t^3-3kt+a&E_{a,b}\times F_{3,a,k}^{\,2}\rule{0pt}{13pt}\\
C_{4,a,b}&F_{4,a,c}:u^2=t^4-4ct^2+a+2c^2&E_{a,b}\times F_{4,a,c}^{\,2}
\end{array}                                                     \tag{A.9}
\]
The genus-two quotients by a reflection have equations
\[
\begin{split}
 D_{3,a,k}:&\quad z^2=(t^2-4k)(t^3-3kt+a),\\
 D_{4,a,c}:&\quad z^2=(t^2-4c)(t^4-4ct^2+a+2c^2).
\end{split}                                                     \tag{A.10}
\]
Their Jacobians are isogenous to $E_{a,b}\times F_{3,a,k}$ and
$E_{a,b}\times F_{4,a,c}$, respectively. The $D_3$ quotient has
one cusp at infinity; the $D_4$ quotient has two. For $D_8$ the choices
$c$ and $ic$ give the two reflection conjugacy classes; $c$ and $-c$
give isomorphic quotients.

The $E$-isotypic dimensions in (A.9) are three for $C_{3,0,-1}$
and one for each of the other five curves. In (A.10) they are two
for $D_{3,0,-1}$ and one otherwise. In more explicit terms,
\[
\begin{array}{c|c}
\text{complementary factor}&j\\ \hline
F_{3,0,-1}&1728\rule{0pt}{13pt}\\
F_{3,\pm6,1}&-216\\
F_{4,0,c}&8000\\
F_{4,\pm6,c}&10976\text{ or }128
\end{array}                                                     \tag{A.11}
\]
The last two values are isogenous complementary factors, not additional
$E$ factors.

The unique geometric $Q_8$-cover has compactification
\[
 C_Q:X^4+Y^4=2Z^4,
\quad x_E=(xy)^2,\quad y_E=\frac{i}{2}xy(x^4-y^4),
\quad x=X/Z,\ y=Y/Z.                                            \tag{A.12}
\]
Its four cusps are $Z=0$, and $\operatorname{Jac}(C_Q)\sim E^3$.
It has no genus-two quotient, even by an automorphism outside its
$Q_8$ deck group. Thus an expectation of a genus-two quotient for
this particular cover must be corrected.

The groups generated by degree-zero cusp divisor classes are exactly
\[
\begin{array}{c|c}
 C_{3,a,b}&\Z/2\rule{0pt}{13pt}\\
 D_{3,a,k}&0\\
 C_{4,a,b}&(\Z/4)^2\\
 D_{4,a,c}&\Z/4\\
 C_Q&(\Z/4)^2.
\end{array}                                                     \tag{A.13}
\]
In particular all projected cusp divisors in A.3 are torsion for
these covers. The dimensions of their potentially $\Delta$-fixed
rational homology are, in the same order, $7$ for $C_{3,0,-1}$ and
$3$ for the other $C_3$ curves; $4$ for $D_{3,0,-1}$ and $2$ for
the other $D_3$ curves; $5$ for every $C_4$; $3$ for every $D_4$;
and $9$ for $C_Q$. Their potentially Galois-stable rational lines
are then exactly the boundary and CM lines described in A.3.
\end{statement}
\begin{proof}
We give constructions and the divisor calculations, so no computer-algebra
decomposition of a Jacobian is assumed. The standard degree-two isogeny
with kernel $(0,0)$ is
\[
 (x,y)\longmapsto
 \left(x+a+\frac b x,\ y\left(1-\frac b{x^2}\right)\right)
\]
from $E_{a,b}$ to $y'^2=x'^3-2ax'^2+(a^2-4b)x'$. For
$a=\pm6,b=1$, the substitution $x' =4x_E+2a/3$, $y'=8y_E$
identifies the target with $E$. For $a=0,b=-1$, use
$x_E=-ix'/2$, $y_E=-(1+i)y'/4$; this is the isogeny $[1+i]$
with the indicated coordinates. Direct substitution proves both
covering maps in (A.8).

Let $\omega^3=1$, $\omega\ne1$. The $C_3$ deck generators are
\[
 \tau(w,v)=(\omega w,\omega^{-1}v),\qquad
 j(w,v)=(k/w,-k^2v/w^4).
\]
They satisfy $\tau^3=j^2=1$, $j\tau j=\tau^{-1}$.
The second map induces translation by $(0,0)$ on $E_{a,b}$.
For $C_4$, the corresponding generators are
\[
 \tau(w,v)=(iw,-v),\qquad j(w,v)=(c/w,-c^2v/w^4),
\]
and satisfy $\tau^4=j^2=1$, $j\tau j=\tau^{-1}$.
The polynomials in (A.8) have simple roots for our three parameter
pairs, so their indicated smooth projective models have genus three.
The only points above the removed kernel of the isogeny are exactly
the displayed cusps. The maps have total degrees six and eight, and
the displayed deck groups have those orders, hence these are the
required Galois covers.

For completeness they exhaust the covers, rather than merely giving
examples. Surjections from the free rank-two geometric group to the
three finite groups are generating pairs modulo group automorphisms.
The executed independent program \texttt{a\_first\_covers.py}
enumerates $18,24,24$ generating pairs in $S_3,D_8,Q_8$.
Their automorphism groups have orders $6,8,24$, and the actions on
generating pairs are free. Thus the numbers of covers are $3,3,1$.
In each of the first two cases the three displayed covers have the
three different degree-two intermediate covers of $E\setminus O$,
corresponding to the three nonzero two-torsion subgroups, so they are
distinct. Although the two sources with $a=6$ and $a=-6$ have equal
$j$, their dual kernels on $E$ are the two different nonzero points
other than $(0,0)$; the source isomorphism covers $[i]$ on $E$, not
the identity. This also identifies the unique $[i]$-stable $S_3$ cover.

Let $h(w,v)=(w,-v)$ be the hyperelliptic involution. On $C_3$ set
\[
 t=w+k/w,\qquad u=v/w^2,\qquad z=(w-k/w)u.
\]
The map $(t,u)$ is the quotient by $hj$, and $(t,z)$ the quotient
by $j$. The identities
\[
 w^3+k^3/w^3=t^3-3kt,\qquad (w-k/w)^2=t^2-4k
\]
give the first equations in (A.9)--(A.10). Similarly, on $C_4$ use
$t=w+c/w$, $u=v/w^2$, and $z=(w-c/w)u$. The identity
\[
 w^4+c^4/w^4=t^4-4ct^2+2c^2
\]
gives the second equations. Each $j$ is fixed-point-free: a fixed
point would have $w^2=k$ or $w^2=c$ and $v=0$, which our parameter
values exclude. Riemann--Hurwitz therefore gives genus two for its
quotient. The quotient by $hj$ has genus one.

Here is also a direct explanation of the isogeny factors. The quotient
by the rotation group has Jacobian $E_{a,b}$. Its differential is the
one-dimensional rotation-invariant summand of
$\langle dw/v,w\,dw/v,w^2dw/v\rangle$. The remaining two-dimensional
summand is the standard rational representation of $S_3$, respectively
$D_8$. Its simple rational group-algebra component is $M_2(\Q)$.
The two reflection projectors give isogenous one-dimensional abelian
factors; the quotient by $hj$ supplies one of them. Clearing the
denominators of the rational idempotents gives the isogenies in (A.9).
The $j$-invariant part similarly gives (A.10). Equivalently the images
of the elliptic quotient maps span the differential spaces, which
makes the induced map of abelian varieties an isogeny.

The $j$ values follow either from the cubic equation or from
\[
 j(u^2=t^4+Bt^2+C)=
       \frac{16(B^2+12C)^3}{C(B^2-4C)^2}.
\]
After scaling $t$ by a square root of $c$, the $C_4$ complements
have $B=-4$ and $C=2+a/c^2$, hence $C=2,8,-4$.
It remains to justify that none of the complements with $j\ne1728$
is geometrically isogenous to $E$; inequality of $j$ alone would
not justify that assertion. For $F_{3,6,1}$ the trace at $5$ is
$-1$, with Frobenius field $\Q(\sqrt{-19})$, whereas $E$ has
trace $-2$ and field $\Q(i)$. The eigenvalue quotients are not
roots of unity, so all positive powers still generate those
different quadratic fields. An isogeny over any finite extension
would identify some powers of these Frobenius polynomials, which
is impossible. The $a=-6$ factor is a quadratic twist and gives
the same conclusion. For $F_{4,0,c}$, use its normalized quartic
at $3$: its trace is $2$, hence its reduction is ordinary, while
$E$ has supersingular reduction. For the other two quartics use
$7$, where their traces are $-4$ and $4$, while $E$ has trace zero
and is supersingular. All these primes are primes of good reduction
for the indicated equations. Geometric isogeny preserves the
ordinary/supersingular distinction. The independent point counts
are recorded in \texttt{a\_first\_covers.json}.

For (A.12) the automorphisms
\[
 \alpha(x,y)=(ix,-iy),\qquad \beta(x,y)=(y,-x)
\]
satisfy $\alpha^2=\beta^2$, $\alpha^4=1$ and
$\beta\alpha\beta^{-1}=\alpha^{-1}$, and generate $Q_8$.
Their common central involution has exactly the four points $Z=0$
as fixed points; the other nonidentity elements outside the center
have none. The invariant functions in (A.12) satisfy $y_E^2=x_E^3-x_E$
and have degree eight, so give the quotient. The three sign
involutions on $X,Y,Z$ yield elliptic quotients with $j=1728$.
Their differential projectors span the three-dimensional space,
giving $\operatorname{Jac}(C_Q)\sim E^3$.

The smooth plane quartic $C_Q$ is nonhyperelliptic, and its canonical
embedding is this plane model. Every involution extends to a
projective linear involution of the plane and fixes a line. That
line meets the curve in four distinct points: tangency would give
derivative one for the involution at a fixed point of the curve,
forcing the involution to be the identity in characteristic zero.
Thus no involution is free. A map from genus three to genus two
would, by Riemann--Hurwitz, have degree two and be unramified.
Such a map cannot exist. This proves the assertion about genus-two
quotients without relying merely on the list of subgroups of $Q_8$.

We finish with the exact cusp-divisor groups. On $C_3$,
$\operatorname{div}(w)=2(P_0-P_\infty)$. The class is nonzero:
a function with divisor $P_0-P_\infty$ would give a degree-one map
from a positive-genus curve to $\mathbf P^1$. The $D_3$ quotient
has only one cusp.

For $C_4$ write $s=c^2$, so $s^2=b$, and label the cusps
$A=(0,s)$, $B=(0,-s)$, $C=(\infty,+)$, $D=(\infty,-)$.
The order-four automorphism $(w,v)\mapsto(iw,v)$ fixes all four.
Its quotient is the rational conic $v^2=q^2+aq+b$, $q=w^4$,
with parameter
\[
 T=(v-s)/w^4,\qquad w^4=\frac{a-2sT}{T^2-1}.
\]
The four cusp images are $a/(2s),\infty,1,-1$; they are distinct.
Pullbacks of rational functions on this line kill four times each
cusp difference, and
$\operatorname{div}(w)=A+B-C-D$ is an additional relation.
These are all relations. Indeed a function whose divisor is
supported on the four cusps is an eigenfunction for the order-four
automorphism, since that automorphism fixes its divisor; its
quotient with its transform is constant. It therefore has the form
$w^e R(T)$, $0\le e<4$, and its valuations are exactly those
generated by the stated relations. The resulting Smith form is
$(1,4,4)$ on the degree-zero divisor lattice, so the group is
$(\Z/4)^2$.

For $D_4$ write $Q_+,Q_-$ for the two points at infinity and
$\delta=[Q_+-Q_-]$. Pushforward from $C_4$ shows $4\delta=0$.
Let $W_+,W_-$ be its Weierstrass points $t^2=4c$.
The connected etale double cover $C_4\to D_4$ corresponds to
the nonzero two-torsion class
\[
 \eta=[W_++W_--Q_+-Q_-],
\]
as is seen by adjoining $\sqrt{t^2-4c}$. The identity
\[
 z^2-(t^3-4ct)^2=(a+2c^2)(t^2-4c)
\]
shows that
$\operatorname{div}(z-t^3+4ct)=W_++W_-+Q_+-3Q_-$,
with the labels chosen by $z/t^3=+1$ at $Q_+$.
Consequently $\eta=-2\delta\ne0$, so $\delta$ has exact order four.

Finally the four points $Z=0$ on $C_Q$ are the totally ramified
points of the cyclic fourth-root cover
\[
 (Z/Y)^4=((X/Y)^4+1)/2.
\]
The same eigenfunction argument gives all relations in their
degree-zero divisor group: four times each difference, and the
relation whose coefficient vector in the basis $P_j-P_4$ is
$(1,1,1)$. The latter is also visible from the line $Z=0$ and a
hyperflex tangent line, whose divisor is $4P_4$. Its Smith form
is again $(1,4,4)$, proving (A.13).

All divisor assertions used in A.3 are therefore torsion assertions,
not inferred just from the CM decompositions. Applying A.3 now
gives each listed dimension as $2\dim J_E+|S|-1$.
\end{proof}

\begin{statement}{A.10}{Proved}{The complete first-cover calculation}
For every first cover and genus-two quotient specified in Task A.4,
at every prime and for every stated packet, the subgroup fixed by
the permitted open subgroup of $\Delta$ is exactly the selected
cusp inertia. Every nontrivial closed procyclic subgroup with
open arithmetic normalizer is contained in a cusp inertia group,
up to conjugacy. Together with A.4 this completes Task A.4.

The explicit models, isogenies, cusp classes, and rational homology
of Task A.4 are determined by A.4. A.22 determines the full fixed subgroups for covers with
$\ell$-group Galois closure, including every $D_8/Q_8$ case at
two. A.27 determines the stable genus-two symmetric-cover fixed groups
for every original packet and every odd prime. A.30 determines the
stable genus-three symmetric-cover fixed groups for every original
packet at every odd prime, including three. A.31 supplies the
two-adic comparison and settles both stable symmetric curves at
two as well. A.32 determines the two nonstable symmetric
genus-two quotient families at every prime.
A.33 settles the two nonstable symmetric genus-three covers at two
and reduces their odd-prime leading fixed classes to
$\mathbb L(t,\Omega)$ in every degree.
A.35 removes the entire proper image in the quaternion case at every odd
prime, and A.36 determines its full punctured fixed group at every prime.
A.37 determines all three dihedral genus-three fixed groups at every prime.
A.38 determines the two nonstable symmetric genus-three fixed groups
at every prime, using the weighted cyclotomic criterion.
A.39 determines every dihedral genus-two reflection quotient at
every prime, including both reflection classes of each cover.
A.14 and A.25 settle the procyclic normalizer classification for every
displayed genus-three cover and every displayed genus-two quotient,
at all primes. The corresponding universal assertion in Task A.5
remains Open. Potentially fixed homology is not asserted to lift to
fixed elements of $H^{(\ell)}$.
\end{statement}
\begin{proof}[Proof and comparison with the earlier reductions]
The displayed cases exhaust the covers and quotients in A.4.
The cited theorems give genuine actions with the chosen cusp lift,
and include passage to every permitted open subgroup of $\Delta$.
Their fixed-subgroup equalities, together with A.14 and A.25 for
the normalizer assertion, prove the statement. We record why the
abelian calculations alone had not sufficed.

For the $[i]$-stable $S_3$ cover the first geometric claims in Lead 3
are proved: it is $v^2=w^7-w$, its cusp class has order two, and
its Jacobian is isogenous to $E^3$. Thus its seven-dimensional
rational open-curve homology becomes fixed by an open subgroup of
$\Delta$. This is a statement about the associated abelian quotient.
The hyperelliptic involution acts as $-1$ on compact-curve homology
but as $+1$ on the boundary generator, so a parity argument treating
all seven generators alike would be invalid.

A.28 computes the actual stable marked graph at two: three elliptic
tails attached to a rational component containing both cusps.
Its separating monodromy and a CM Frobenius conjugate reduce every
possible fixed leading class to $\mathbb L(t,\Omega)$. A.29 proves
that any nonzero first Johnson matrix kills every class in that
algebra except $t$, without an invertibility assumption. A.30 computes the actual primitive cusp-path extension using the
universal Collino cycle and its elliptic norm. Its nonzero regulator
is detected by the unit $2+\sqrt3$, so the possible zero-Johnson
packet does not exist. A.31 replaces tame monodromy at two by
Fontaine's monodromy on the crystalline realization of the global
orbit span of a fixed class. Its explicit logarithmic tangential
comparison handles arbitrary packets over ramified local fields.
The stable symmetric-cover fixed groups are therefore determined
at every prime.

For the other first covers, the non-$E$ factors in (A.11) already
exclude their proper homology from the potentially fixed subspace,
but do not identify the nonabelian kernel. A.1 and A.9 completely
settle which fixed quotient elements admit a fixed full-profinite
lift: exactly those in the image of the specified cusp inertia.
They do not decide whether additional quotient-fixed elements exist.
The last symmetric case was a nonstable $S_3$
cover at an odd prime. A.33 gives its three-tail graph and
uses the derived special linear action on the non-CM factors
to reduce the kernel to $\mathbb L(t,\Omega)$. The primitive
first Johnson class takes values in a symmetric-square module;
the CM elliptic norm in A.30 does not evaluate it.
Instead, A.38 uses the full Galois action on the actual fixed
subgroup: the Tate-graded weighted kernel forces the cyclotomic
kernel to act trivially on that subgroup, after which proper-cover
homology and Ribet's theorem exclude every nonperipheral element.
A.39 applies the same criterion to the remaining dihedral
genus-two quotients. Thus no unevaluated first Johnson class
remains in Task A.4. The assertion on arbitrary finite covers,
where these explicit weighted kernels are unavailable, remains
the universal open problem of Task A.5.
\end{proof}

\subsection{Proper-homology detection and the first-cover normalizers}

\begin{statement}{A.11}{Proved}{A noncuspidal pro-prime normalizer is visible in proper homology}
Let $Y$ be an affine hyperbolic curve over a number field, with smooth
compactification $\overline Y$. Let $P=\pi_1(Y_{\overline K})^{(\ell)}$.
If a nontrivial closed procyclic subgroup $A\subset P$ has arithmetic
normalizer with open image and is not contained in cusp inertia, then
\[
                  A\longrightarrow V_\ell\operatorname{Jac}(\overline Y)
\]
has nonzero image. Its rational span is a line stable under an open
arithmetic subgroup. Consequently, if finitely many morphisms
$Y\to E_j\setminus\{O_j\}$ detect every potentially stable line in
proper homology, and each target satisfies the pro-$\ell$ cuspidal
normalizer criterion, then every arithmetically normalized procyclic
subgroup of $P$ is cuspidal.
\end{statement}
\begin{proof}
Choose free pro-$\ell$ generators consisting of $2g$ handle generators
and $r-1$ cusp generators, with weights one and two respectively.
The associated weighted Magnus filtration is separated: it is the
restriction of weighted degree in the noncommutative power-series ring,
and every generator has positive weight. Its rational graded pieces
are generated by the proper first homology in weight one and the
boundary Tate classes in weight two. At a good prime with residue
cardinality $q$, every eigenvalue in degree $d$ therefore has complex
absolute value $q^{d/2}$. This remains true after finite extension and
after splitting the cusp set. The construction is Galois invariant:
its degree-two subgroup is the kernel of the map to proper homology,
and commutators add weights.

Let $u$ generate $A\simeq\Z_\ell$ and let $d$ be its first nonzero
weighted degree. The leading Magnus term is nonzero in a torsion-free
module. If $\theta$ is the conjugation character of the arithmetic
normalizer, geometric inner automorphisms act trivially on that leading
term; cancellation shows that $\theta$ factors through the arithmetic
image. The resulting line in degree $d$ has character $\theta$.

By \citet[Lemma 2.1.4]{Nakamura1994}, a noncuspidal element is detected
outside the boundary in the mod-$\ell$ homology of a finite cover.
If necessary replace $u$ by a nonzero $\ell$-power and replace $P$ by
an open subgroup, as in the one-cusp case of that lemma. A nonzero
power of a noncuspidal element is still noncuspidal: otherwise its
centralizer would put the original element in the same inertia group.
The detector is an open subgroup of a pro-$\ell$ group, so its
finite cover has an $\ell$-group Galois closure and its pro-$\ell$
topology is induced. Its proper integral homology is torsion-free, so its
nonzero mod-$\ell$ image gives a nonzero rational image. Descend that
finite cover to a finite number field and intersect with the original
normalizer. Its arithmetic image is still open. The same character
$\theta$ now occurs in proper homology and has Frobenius weight one.
At one prime good for both curves, comparison gives $d=1$. By the
definition of the weighted filtration, the original image in proper
homology is nonzero, as asserted.

For the last assertion, a noncuspidal $A$ would give such a line.
One of the displayed morphisms has nonzero proper-homology image on
it, so the image of $A$ in the target pro-$\ell$ group is nontrivial.
The image of its arithmetic normalizer still has open arithmetic
projection. The target criterion puts this image in cusp inertia,
whose proper homology is zero, a contradiction. These are actual
morphisms of punctured curves; a homomorphism of Jacobians alone
would not justify this argument.
\end{proof}

\begin{statement}{A.12}{Proved}{The two CM elliptic targets}
For an elliptic curve with geometric full CM
by $\Z[i]$ or $\Z[\sqrt{-2}]$, every nontrivial closed procyclic
subgroup of the geometric pro-$\ell$ group of its once-punctured curve
with open arithmetic normalizer is contained in cusp inertia.
This holds for every prime $\ell$. In the second CM case the split
prime $3$ uses the additional rational nonvanishing proof A.21 below.
\end{statement}
\begin{proof}
All assertions are unchanged by finite extension of the field of
definition. We may work with the CM eigenspaces after an extension
defining all endomorphisms.
At a nonsplit prime, including the ramified prime $2$, proper Tate
homology has no rational line stable under an open Galois subgroup.
Indeed an ordinary good Frobenius generates the quadratic CM algebra
over $\Q_\ell$, and so does every positive power: its two eigenvalues
have quotient of infinite order, since their ideals have different
support at the two primes over its residue characteristic. A stable
line would give a unital map from that quadratic field to $\Q_\ell$.
Statement A.11 now proves the criterion at these primes.

At a split prime a stable proper line must be one of the two CM
eigenlines. We recall the rational obstruction to either line, giving
the character separation explicitly because CM by $\Z[\sqrt{-2}]$
has no automorphism of order four. Put $L=\mathbb L_{\Q_\ell}(x,y)$,
$\omega=[x,y]$, and let $b_w$ denote Lyndon bracketing with $x<y$.
In the special derivation algebra set
\[
\begin{aligned}
 \epsilon_4(x)&=\ad_x^4(y),\\
 \epsilon_4(y)&=[y,\ad_x^3(y)]-[[x,y],\ad_x^2(y)],\\
 T&=[F,\epsilon_4],\qquad F(x)=y,\quad F(y)=0.
\end{aligned}
\]
The direct tensor calculation in \texttt{a\_normalizer\_certificates.py}
gives
\[
 T(\omega)=0,\quad
 T(x)=-2b_{xxxyy}-b_{xxyxy},\quad
 T(y)=2b_{xxyyy}+3b_{xyxyy}.                    \tag{A.15}
\]
In particular $T(x)\notin[x,L_4]$ and $T(y)\notin[y,L_4]$.
For example the relevant bidegree of $[x,L_4]$ is the one-dimensional
span of $b_{xxxyy}$, while the coefficient of $b_{xxyxy}$ in $T(x)$
is $-1$. This is a rational obstruction, so it also applies to a
nonprimitive integral eigenvector.

The degree-four special outer module has the three distinct CM
characters $\psi^3\bar\psi$, $\psi^2\bar\psi^2$, and
$\psi\bar\psi^3$; its middle line is $\Q_\ell T$.
The special modules in degrees one and three vanish, and the
degree-two special module is inner, generated by $\ad\omega$.
These assertions follow either by the boundary equation
$[D(x),y]+[x,D(y)]=0$ or by its Lyndon calculation.
Nakamura's metabelian formula
\cite[Theorem A, Corollary 4.12]{Nakamura1995} identifies the middle
degree-four coordinate with a nonzero rational scalar times the
elliptic Soul\'e character $\kappa_{(2,2)}$. For split $\ell\ge5$
this character is nonzero by \cite[Theorem 4.5 and Corollary 4.10]{Ishii2025};
the fields have class number one, the multiplication is by the full
integer ring, and split primes are ordinary after finite extension.
For $\Z[\sqrt{-2}]$ at $3$, A.21 supplies the same rational assertion.
For $\Z[i]$ there is no split prime below five.

Here is why the nonzero middle coordinate gives the needed individual
derivation even if the two outside coordinates are also present.
The rational span of the degree-four Galois image is stable under the
CM torus. The ratio $\psi/\bar\psi$ has infinite image on every open
Galois subgroup, so the three displayed characters remain distinct.
Choose an element separating their three eigenvalues; polynomial
projectors in its action show that a nonzero middle projection puts
the entire line $\Q_\ell T$ in this span. One can compute the image
on the kernel of the homological representation. All lower outer
degrees are zero; their possible degree-two inner term can be
removed by a rational power of the actual boundary element, whose
logarithm starts with $\omega$. This keeps the boundary fixed, so
the special degree-three vanishing still applies. Thus this is a statement in the
actual first non-inner outer quotient, not merely in a universal
space of special derivations. Nonvanishing persists on every open
subgroup, since an additive $\ell$-adic character with nonzero image
cannot have finite image.

Suppose a procyclic subgroup with first term $x$ were normalized
arithmetically. An element in the kernel of the homological action
must act on this subgroup with exponent one. After the inner
normalization just described, its first non-inner action is in
degree four. Writing a proposed generator as $x+$ higher terms in
the rational nilpotent Lie algebra through degree five, invariance
up to conjugation forces $D(x)\in[x,L_4]$. To justify the degree of
the conjugator, cancel its successive components proportional to
the proposed generator; the free-Lie centralizer of $x$ is
$\Q_\ell x$, so its first effective component has degree four.
Contributions of the
higher terms start in degree six. This is a linear condition on
$D$, so it holds on the rational span of the Galois image. It
contradicts (A.15) for $D=T$. The other eigenline is excluded in
the same way. A.11 now finishes the proof in the original pro-$\ell$
group. We have used the Malcev quotient only for a necessary
condition, not for a claim about integral conjugacy.
\end{proof}

\begin{statement}{A.13}{Proved}{Punctured elliptic maps for the quaternion and stable dihedral covers}
Work over an algebraic closure of $\mathbf Q(i)$; every curve, map,
origin, and isogeny below is defined over a finite extension.
Let $S$ be the cusp set specified in A.4.

For the quaternion cover
\[
 C_Q:X^4+Y^4=2Z^4,\qquad S=\{Z=0\},
\]
each of the three sign-involution elliptic quotients can be followed
by an isogeny so that the inverse image of its target origin is
\emph{exactly} $S$. For the sign changes of $X$ and $Y$ these
composites have degree four and have target
$E':B^2=A^3-32A$, with affine formulas
\[
 \phi_X(x,y)=(4y^2,-8i\,yx^2),\qquad
 \phi_Y(x,y)=(4x^2,-8i\,xy^2),\qquad x=X/Z,\quad y=Y/Z.
                                                        \tag{A.13.1}
\]
Their fibres over $O$ are reduced. For the sign change of $Z$,
choose one cusp image as origin on the elliptic quotient and
follow its degree-two quotient map by multiplication by two.
The resulting degree-eight map has fibre $2\sum_{P\in S}P$
over the origin. All three elliptic targets have CM by
$\mathbf Z[i]$, and their pulled-back differential lines span
$H^0(C_Q,\Omega^1)$.

For the stable dihedral genus-two curve
\[
 D_c:z^2=(t^2-4c)(t^4-4ct^2+2c^2),\qquad c^4=-1,
\]
write $Q_+,Q_-$ for its two points at infinity. There is a
degree-two map
\[
 D_c\longrightarrow E_c^{-2}:\quad (t,z)\longmapsto(X,Y)=(t^2,z),
 \qquad
 Y^2=X^3-8cX^2+18c^2X-8c^3.                         \tag{A.13.2}
\]
This elliptic curve has $j=8000$ and CM by
$\mathbf Z[\sqrt{-2}]$. The origin fibre is precisely
$Q_++Q_-$. It is the complementary elliptic factor to the
$\mathbf Q(i)$ factor of A.4 and defines a punctured map with
no extra points removed from the source.

For the stable symmetric-group genus-two curve
\[
 D_3:z^2=(t^2+4)(t^3+3t),\qquad P=\infty,
\]
every nonconstant morphism to an elliptic curve whose origin fibre
is supported only at $P$ has degree three and ramification divisor
$2P$. All such maps pull back elliptic differentials to the same
one-dimensional subspace of $H^0(D_3,\Omega^1)$; they cannot detect
the second elliptic factor independently. In fact they coincide
up to an isomorphism of their pointed elliptic targets.
Thus the proposed second independent punctured elliptic map
for this one-cusp genus-two curve does not exist.
\end{statement}
\begin{proof}
For the sign change of $X$, the quotient has coordinates
$t=Y/Z$, $u=X^2/Z^2$ and equation
\[
                         u^2=2-t^4.
\]
This is a smooth genus-one quartic. Its two points at infinity,
$P_+$ and $P_-$, have $u/t^2=+i$ and $-i$, respectively.
The sign involution has no fixed cusp, so the four cusps map
in pairs to these two points and are unramified there.
The identity
\[
                   (u-it^2)(u+it^2)=2
\]
and the expansions at infinity give
\[
             \operatorname{div}(u-it^2)=2P_+-2P_-.
\]
The difference $P_--P_+$ is nonzero: otherwise a function with
one simple pole would give a degree-one map from this elliptic
curve to a projective line. Choose $P_+$ as origin. The other
point is therefore an exact two-torsion point, and the isogeny
with kernel these two points has precisely their union as its
origin fibre. Pulling back gives exactly the four original cusps.

There are explicit coordinates for this isogeny. Put
$v=-iu$, so $v^2=t^4-2$, and set
\[
       a=2(v+t^2),\qquad b=4t(v+t^2).
\]
Then $b^2=a^3+8a$, the chosen origin is infinity, and the other
infinity is $(0,0)$. The standard degree-two isogeny with kernel
$(0,0)$ is
\[
              (a,b)\longmapsto
                   \left(a+\frac8a,b\left(1-\frac8{a^2}\right)\right).
\]
Its two coordinates simplify to $4t^2$ and $8tv$.
This proves the first formula in (A.13.1); interchange $X,Y$ to
obtain the second. The target has the automorphism
$(A,B)\mapsto(-A,iB)$, and is isomorphic to $y^2=x^3-x$
after adjoining a fourth root of $32$.

For the sign change of $Z$, use $t=X/Y$, $u=Z^2/Y^2$ on
the quotient. Its equation is $2u^2=t^4+1$. The four cusp
images are exactly the four points $u=0$. Choose one of these
as origin. The involution $u\mapsto-u$ fixes this origin and
is the elliptic involution $[-1]$; hence those four points are
exactly the full two-torsion subgroup. Multiplication by two
has precisely that subgroup as its origin fibre. Each original
cusp is ramified with index two in the sign-$Z$ quotient, so the
composite fibre has multiplicity two at each of the four cusps.
The quartic has $j=1728$, either by its order-four automorphism
or by the elementary quartic invariant formula. Its target under
$[2]$ is the same elliptic curve.

Finally the canonical forms on the plane quartic are the
adjunction forms corresponding to $X,Y,Z$. Under the three
sign involutions their respective invariant lines are distinct
and together span the three-dimensional differential space.
These are the pullbacks from the three genus-one quotients.
An isogeny in characteristic zero has nonzero differential,
so composing with the displayed isogenies leaves those three
lines unchanged. This proves the claimed simultaneous detection
of all three proper elliptic factors.

For (A.13.2), direct substitution $X=t^2$, $Y=z$ gives the cubic.
Its discriminant and invariant are
\[
                       \Delta=512c^6,\qquad c_4=160c^2,
\]
so it is smooth and $j=c_4^3/\Delta=8000$. At each infinity
of $D_c$, the functions $X$ and $Y$ have pole orders two and
three, respectively. Thus each point maps unramified to the
cubic's origin, and their total degree two exhausts the fibre.

The CM assertion can also be proved without consulting a table
of singular moduli. Translate $x=X-4c$ to write the cubic as
\[
                         y^2=x^3+4cx^2+2c^2x.
\]
The two-isogeny with kernel $(0,0)$ has target
$y'^2=x'^3-8cx'^2+8c^2x'$. This target is isomorphic to the
source by $(x',y')=(-2x,2\sqrt{-2}\,y)$.
Compose the two-isogeny with the inverse isomorphism. The resulting
degree-two endomorphism has differential $-\sqrt{-2}$, so its
square is multiplication by $-2$: endomorphisms of an elliptic
curve in characteristic zero are determined by their differentials.
This embeds the maximal quadratic order $\mathbf Z[\sqrt{-2}]$
in its endomorphism ring. It is therefore the CM ring.
It differs from $\mathbf Q(i)$, so this elliptic quotient is
the complementary factor in A.4.

For the last assertion, the point at infinity on the odd-degree
hyperelliptic model $D_3$ is a Weierstrass point. If
$f:D_3\to E_0$ is nonconstant and $f^{-1}(O)$ is supported
only there, then $f^*O=dP$, where $d=\deg f$. The ramification
divisor $R_f$ has degree two by Riemann--Hurwitz and contains
$(d-1)P$. Thus $2\le d\le3$. The case $d=2$ is impossible:
$R_f$ is the divisor of the pullback of a nonzero elliptic
differential, hence is canonical. The unique canonical divisor
containing the Weierstrass point $P$ is $2P$, whereas a
degree-two map has ramification multiplicity at most one at
each point. Therefore $d=3$ and $R_f=2P$.
Consequently
\[
                 f^*H^0(E_0,\Omega^1)=H^0(D_3,\Omega^1(-2P)),
\]
a one-dimensional space. The degree-three map already supplied
by A.4 gives precisely this line.

For completeness, equality of these lines gives uniqueness of
the pointed target quotient. Embed the field of definition in
$\mathbf C$. The induced homomorphism
$J(D_3)\to E_0$ has connected kernel. Otherwise $f$ factors
through a nontrivial finite étale isogeny of elliptic curves;
its degree divides three, and the remaining map would have
degree one, which a genus-two curve cannot have. The connected
kernel is the elliptic subvariety whose tangent space is the
annihilator of the displayed differential line. A connected
complex subtorus is determined by its tangent space. Hence all
these kernels are equal, and the quotient elliptic curves are
isomorphic. Since $P$ maps to the chosen origins, no translation
ambiguity remains.

All displayed functions extend to morphisms of smooth proper
curves, and the exact fibre computations show that they restrict
to morphisms of the indicated punctured curves. After compatible
geometric base points and paths are chosen, they induce the
usual continuous homomorphisms of arithmetic and geometric
fundamental groups. A change of path conjugates these maps by
a geometric element; the proper-homology detection assertions
are unaffected. The executable certificate
\texttt{code/a\_cover\_maps.py} verifies the map identities,
the two-isogeny identities, and the elliptic invariants.
\end{proof}

\IfFileExists{sections/A_prime_three.tex}{\begin{statement}{A.21}{Proved}{The diagonal elliptic Soul\'e character at the split prime three}
Let \(k=\mathbb Q(\sqrt{-2})\), let
\(3=\mathfrak p\overline{\mathfrak p}\), and let \(E/k\)
be a CM elliptic curve with endomorphism ring
\(\mathbb Z[\sqrt{-2}]\), good reduction at three, and
its standard puncture at the origin. In the homological
normalization \(\psi\overline\psi=\chi_3\), the diagonal
elliptic Soul\'e character
\(\kappa_{(2,2)}\) is nonzero after tensoring with
\(\mathbb Q_3\). Its restriction to any finite-index
subgroup of its defining Galois group remains nonzero.
This assertion at three is proved below; it is not an
application of a theorem whose statement assumes
\(p\ge5\).
\end{statement}
\begin{proof}
Write \(S\) for the two primes above three and the
infinite places. We first prove the precise cohomological
input, including the residual character:
\[
 H^2(G_{k,S},\mathbb Z_3(2))=0.                         \tag{A.21.1}
\]
Put \(L=k(\zeta_3)=\mathbb Q(\sqrt{-2},\sqrt{-3})\).
The discriminants of its three quadratic subfields are
\(-8,-3,24\); hence \(\operatorname{disc}(L)=576\).
Its Minkowski bound is \(36/\pi^2<4\). There is no
prime ideal of norm two: at two the extension supplied
by \(\sqrt{-3}\) is the unramified quadratic extension,
so the residue degree is two. The only prime ideals of
norm below four are the two primes above three. Each
has square principal, generated by
\(1+\sqrt{-2}\) or \(1-\sqrt{-2}\). Thus the class
group is generated by elements of order dividing two,
and \(\operatorname{Cl}(L)[3]=0\).

There are exactly two primes in \(S_L\), both fixed by
\(A=\operatorname{Gal}(L/k)\simeq C_2\). Kummer theory
and the Brauer invariant-sum sequence give
\[
 H^2(G_{L,S_L},\mu_3)
  =\ker(\mathbb F_3^2\xrightarrow{\sum}\mathbb F_3)
  \simeq\mathbb F_3,
\]
with trivial \(A\)-action. Passing from \(\mu_3\) to
\(\mathbb F_3(2)\) tensors this action by the nontrivial
cyclotomic character of \(A\). Averaging over \(A\)
is exact, and therefore
\(H^2(G_{k,S},\mathbb F_3(2))=0\). The coefficient
sequences prove vanishing for every
\(\mathbb Z/3^n(2)\). Finite-ramification cohomology
with finite coefficients is finite, so the
Mittag--Leffler property and continuous comparison
prove (A.21.1). This uses the same precise finiteness
and comparison theorems as D.1.

We next explain why the cyclotomic elliptic-unit map
has nonzero rational image, without any integral
claim that twelve is a three-adic unit. Put
\(L_\infty=k(\mu_{3^\infty})\),
\(\Delta=\operatorname{Gal}(k(\mu_3)/k)\simeq C_2\),
and \(\Gamma=\operatorname{Gal}(L_\infty/k(\mu_3))\simeq\mathbb Z_3\).
Let \(\mathcal E_L\) be the inverse limit under norms
of the three-adic completions of global units in this
tower, and let \(\mathcal A_L\) be the inverse limit
of the three-primary ideal class groups. A superscript
\((1)\) in the next formulas means the cyclotomic
Tate twist. The norm-and-Kummer construction gives
\[
 e_L:\mathcal E_L(1)_{\Delta\times\Gamma}
              \longrightarrow H^1(G_{k,S},\mathbb Z_3(2)).
                                                               \tag{A.21.2}
\]
Both its kernel and the twisted class-group coinvariants
are finite. Here is the cohomological verification of
the rational assertion needed, so that no restriction
\(p\ne3\) from a regulator theorem is imported.

Form the inverse corestriction complex
\(C_{\rm Iw}=R\varprojlim_nR\Gamma(G_{L_n,S},\mathbb Z_3(1))\).
Shapiro's cochain isomorphism, followed by the two-term
resolution for \(\Gamma\), identifies its derived
specialization at the character \(\chi_3^{-1}\), with
the corresponding \(\Delta\)-component, with
\(R\Gamma(G_{k,S},\mathbb Z_3(2))\).
The finite-level Kummer sequences identify its first
cohomology with norm-compatible \(S\)-units and its
second cohomology with an extension of
\(\mathcal A_L\) by the local Brauer-invariant modules.
The valuation and Brauer modules at the two primes
above three have trivial \(\Delta\)-action, whereas
\(\chi_3^{-1}|_\Delta\) is nontrivial. They therefore
vanish in this component; ordinary units and
\(S\)-units give the same component. The resulting
one-variable specialization sequences are
\[
\begin{split}
0&\longrightarrow\mathcal E_L(1)_{\Delta\times\Gamma}
 \longrightarrow H^1(G_{k,S},\mathbb Z_3(2))
 \longrightarrow\mathcal A_L(1)^{\Delta\times\Gamma}
 \longrightarrow0,\\
H^2(G_{k,S},\mathbb Z_3(2))
 &\simeq\mathcal A_L(1)_{\Delta\times\Gamma},
\end{split}                                                   \tag{A.21.3}
\]
up to finite groups, which is sufficient here. More
explicitly, any finite error from passing between
finite-level Kummer sequences and completed units
vanishes after tensoring with \(\mathbb Q_3\); the
completed cohomology groups are compact; inverse limits
of their exact sequences are exact by compactness. At
each finite coefficient level the groups are finite,
so the corresponding diagonal towers are Mittag--Leffler. The maps in (A.21.3) are the
norm-and-Kummer maps defining (A.21.2), not different
abstract comparison maps. The class-group Iwasawa
module is finitely generated torsion. For such a
one-variable module, finite twisted coinvariants imply
finite twisted invariants, by its elementary-divisor
decomposition. Equation (A.21.1) proves this finiteness
and makes (A.21.2) an isomorphism after tensoring with
\(\mathbb Q_3\). The global Euler characteristic formula
then gives
\(\dim_{\mathbb Q_3}H^1(G_{k,S},\mathbb Q_3(2))=1\).
Thus the rational image of (A.21.2) is nonzero.

Let \(\mathcal C_L^{\rm R}\subset\mathcal E_L\) be
Rubin's elliptic-unit module. Rubin's original
hypotheses hold at three: the Hilbert class field of
\(k\) is \(k\), its roots of unity have order two,
\([k(\mu_3):k]=2\), and three splits in \(k\).
In particular \cite[\S4, Theorem 4.1(i)]{DE_Rubin1991}
allows this prime; its restriction is
\(p\nmid\#\mu(H)[k_0:k]\), not \(p\nmid6\).
The first characteristic-ideal equality in that
theorem identifies
\(\operatorname{char}(\mathcal E_L/\mathcal C_L^{\rm R})\)
with \(\operatorname{char}(\mathcal A_L)\).
The finite specialization in (A.21.3), the
one-variable torsion structure theorem, and
\cite[Corollary 7.8]{DE_Rubin1991} show that
\[
 (\mathcal C_L^{\rm R}(1)_{\Delta\times\Gamma})\otimes\mathbb Q_3
 \longrightarrow
 (\mathcal E_L(1)_{\Delta\times\Gamma})\otimes\mathbb Q_3
                                                               \tag{A.21.4}
\]
is an isomorphism. The same conclusion follows by
taking the invariants and coinvariants of
\(0\to\mathcal C_L^{\rm R}\to\mathcal E_L
\to\mathcal E_L/\mathcal C_L^{\rm R}\to0\): both
end terms from the quotient are finite. Consequently
Rubin units have nonzero image under (A.21.2).

We transfer this image from the cyclotomic tower to
\(k(3^\infty)\). Write \(\mathcal C^{\rm R}\) for
its two-variable elliptic-unit module. The kernel and
cokernel of the norm map
\[
 (\mathcal C^{\rm R})_{\operatorname{Gal}(k(3^\infty)/L_\infty)}
                   \longrightarrow\mathcal C_L^{\rm R}
\]
are annihilated by the product of the annihilator of
\(\mathbb Z_3(1)\) and the augmentation ideals of the
two local decomposition groups, by
\cite[Theorem 7.7(ii)]{DE_Rubin1991}. After twist \((1)\)
and coinvariants, a displayed annihilator evaluates as
\[
 (\chi_3(\sigma)^{-1}-\chi_3(\sigma))
 (\chi_3(\tau_1)^{-1}-1)
 (\chi_3(\tau_2)^{-1}-1).                              \tag{A.21.5}
\]
The cyclotomic character maps each local decomposition
group onto an open subgroup of \(\mathbb Z_3^\times\).
Choose all three character values equal to a common
\(u=4^a\), \(a\ge1\), lying in these images. Then
(A.21.5) is \((u^{-1}-u)(u^{-1}-1)^2\ne0\).
For \(u=4\) its value is \(-135/64\).
Localization at the evaluation character therefore
annihilates the kernel and cokernel before taking
coinvariants, so no nonexact tensor step is being
silently used. The rational norm map on the required
specializations is an isomorphism. Norm compatibility
of Kummer corestriction proves that the two-variable
Rubin units also have nonzero rational Soul\'e image.

Now compare the units actually defining the geometric
character. Let \(\mathcal C^{\theta}\) be formed from
the theta values of \cite[\S2.2]{DE_Ishii2024}.
At every finite level, the definitions give
\[
 12 C^{\rm R}\ \subset\ C^{\theta}\ \subset\ C^{\rm R}
                                                               \tag{A.21.6}
\]
in additive notation, after adjoining the same roots
of unity. A Rubin generator is an appropriate twelfth
root of the corresponding theta-unit generator;
raising it to twelve gives that generator, and theta
units themselves belong to Rubin's group. These
inclusions commute with norms and use the same
integer twelve at every level. Hence (A.21.6) also
holds for the inverse-limit modules and makes their
rationalizations equal. Roots of unity cause no
additional image: after twist \((1)\) they have
character \(\chi_3^2\), whose specialization at the
trivial character is rationally zero. Thus theta
units have nonzero rational Soul\'e image at three.
This is the step for which the published argument
uses \(p\ge5\); replacing an integral equality by the
uniform inclusions (A.21.6) proves the rational
assertion at three.

Finally the distribution and Kummer identities in
\cite[\S3.1 and \S4.1]{DE_Ishii2024} identify that
image with multiples of \(\kappa_{(2,2)}\). We record
the specialization explicitly. Take
\(\pi=1+\sqrt{-2}\), so \(\pi\overline\pi=3\),
and \(w_k=2\). For any auxiliary ideal \(\mathfrak a\)
prime to six, the primitive theta-unit character
\(\kappa_{(2,2),\mathfrak a}\) satisfies
\[
 \underbrace{2\left(\frac1{1-3}+\frac1{1-3}-1\right)}_{-4}
       \kappa_{(2,2),\mathfrak a}
   =\left(N\mathfrak a-(N\mathfrak a)^{-1}\right)
       \kappa_{(2,2)}.                                  \tag{A.21.7}
\]
The identity follows by separating the two coordinates
of the primitive torsion sum, applying the theta
distribution relation, and summing the geometric
series of ratio three. Its denominators \(1-3\)
are three-adic units, so this calculation is valid at
three. The factor on the right is nonzero for every
nontrivial auxiliary ideal. The theta units in the
two-variable inverse limit are generated by the
primitive theta values and their Galois translates:
the one-coordinate units vanish in this inverse
limit, because norms along the other tower raise
those units to arbitrarily high powers of three.
Taking \(\sigma-1\) simply multiplies their specialized
Soul\'e images by \(\chi_3(\sigma)^{-1}-1\).
Thus a zero \(\kappa_{(2,2)}\) would make the entire
nonzero theta-unit image vanish, contradicting the
preceding paragraph. This proves nonvanishing.
For example \(\mathfrak a=(5)\) has norm twenty-five
and the right factor in (A.21.7) is \(624/25\ne0\).

All arguments just given concern rational nonvanishing;
no assertion about integral surjectivity is required.
The character is a continuous homomorphism to a
characteristic-zero three-adic line on its defining
kernel. If its restriction to a finite-index subgroup
were zero, its image would be finite and hence zero.
This proves the final assertion.
\end{proof}
}{}

\begin{statement}{A.14}{Proved}{The normalized procyclic subgroups in the first covers}
For every prime $\ell$, the nontrivial closed procyclic subgroups
with open arithmetic normalizer in the following pro-$\ell$ groups
are exactly the nontrivial closed subgroups of conjugates of cusp
inertia: all three $S_3$ covers $C_3$, all three dihedral covers
$C_4$, the quaternion cover $C_Q$, all dihedral genus-two quotients
$D_4$, and the two nonstable $S_3$ genus-two quotients $D_3$.
The same assertion holds for the stable $S_3$ genus-two quotient
$D_3$ at $\ell=2$. Statement A.25 below settles its odd-prime
normalizer problem by a further monodromy detector. The requested
nonabelian fixed subgroups are computed in A.22, A.24, and A.27 in the
specified cases; the remaining fixed-group problems are recorded
separately in A.10.
\end{statement}
\begin{proof}
First note that every nontrivial closed subgroup of cusp inertia
has open arithmetic normalizer, after making that cusp rational.
We prove the converse using A.11 and the explicit models of A.4.

On the stable $C_3$, its map to $E_{0,-1}$ followed by the
two-isogeny to $E$ has both cusps as its entire origin fibre.
The elliptic quotient by $hj$ also maps exactly these two cusps
to its origin. This quotient has $j=1728$. Composing it with a
cubic rotation gives a second independent elliptic differential
on the complementary two-dimensional summand. Together the three
maps detect all of proper homology. Statements A.11--A.12 apply.

For a nonstable $C_3$, the complementary elliptic factor
$u^2=t^3-3t+6$, up to twist, is non-CM. Its good ordinary
Frobenius discriminants at $5$ and $11$ are $-19$ and $-35$,
respectively, as independently counted in
\texttt{a\_normalizer\_certificates.json}.
A characteristic-zero CM field would inject into each ordinary
reduction's geometric endomorphism field, giving two different
quadratic fields, which is impossible. Similarly the nonstable
$C_4$ complement $u^2=t^4-4t^2+8$ has discriminants $-16$ and
$-12$ at $5$ and $7$, and is non-CM. Its other reflection quotient
is isogenous to it, so has the same conclusion. The sign changes
of the parameter give twists and do not change geometric CM.

An elliptic curve with geometric endomorphism ring $\Z$ has no
Tate line stable under any open arithmetic subgroup. Indeed Tate
semisimplicity and the endomorphism theorem, over that finite
extension, would give endomorphism dimension at least two if its
two-dimensional representation had an invariant line. The actual
endomorphism algebra is $\Q_\ell$. A line in a direct sum of
copies would project to such a line in one copy. Thus in each
nonstable $C_3$ or $C_4$, all potentially stable proper lines are
in the $E$ factor, which the original covering map to $U$ detects.
This proves these cases, and the same reasoning proves their
genus-two quotient cases, since the descended map to $U$ has
exactly the indicated cusps above $O$.

For stable $C_4$, the proper factors are $E$ and two copies of
the elliptic curve with $j=8000$. Its elliptic reflection quotient
is $u^2=t^4-4ct^2+2c^2$. The degree-two map
$(t,u)\mapsto(X,Y)=(t^2,tu)$ has target
$Y^2=X^3-4cX^2+2c^2X$, again with $j=8000$.
Its origin fibre is exactly the two infinity points.
Consequently the composite map from $C_4$ has the four cusps as
its entire origin fibre. Rotate the reflection to obtain the
second independent complementary differential. By A.13 these
targets have CM by $\Z[\sqrt{-2}]$, and A.12 applies at every
prime. The stable $D_4$ has the additional explicit degree-two
map of A.13 to a once-punctured elliptic curve with this same CM,
which detects its complementary factor. The original map to $E$
detects the remaining factor.

The three punctured elliptic maps of A.13 for $C_Q$ have $j=1728$
targets and span its three-dimensional differential space.
They therefore detect proper homology, proving the quaternion case.

Finally the stable $C_3\to D_3$ is an etale double covering,
also at the compactification. At $\ell=2$ its pro-$2$ group is
an open subgroup of index two of that of $D_3$; the pro-$2$
topology on an index-two subgroup is induced. A nontrivial
arithmetically normalized procyclic subgroup intersects it in
a nontrivial subgroup, which is cuspidal by the already proved
$C_3$ case. The centralizer of this nonzero inertia subgroup is
the containing inertia in $D_3$, so the original subgroup is
cuspidal as well. At odd primes this inclusion argument is not
valid. A.13 proves that no second independent elliptic map with
only the single cusp above its origin can supply the missing
detector, so the same method cannot silently be extended there.
\end{proof}

\begin{statement}{A.22}{Proved}{Fixed subgroups on covers with pro-prime Galois closure}
Let $K=\mathbf Q(i)$, $U=E\setminus\{O\}$, and use the chosen
geometric base point, cusp inertia $I$, and packet section $s_b$
of Part A. Write
\[
 \pi=\pi_1(U_{\overline K}),\qquad
 P_\ell=\pi^{(\ell)},\qquad
 I_\ell=\operatorname{im}(I\to P_\ell).
\]
Let $H\subset\pi$ be open, and suppose the finite quotient by
its normal core is an $\ell$-group. Then the natural continuous
map $H^{(\ell)}\to P_\ell$ is injective with open image
$\overline H=\operatorname{im}(H\to P_\ell)$.
If $\Delta_0\subset\Delta$ is any open subgroup such that
$s_b(\Delta_0)$ preserves $H$, then, in this embedding,
\[
       (H^{(\ell)})^{s_b(\Delta_0)}
           =\overline H\cap I_\ell
           =\operatorname{im}(H\cap I\to H^{(\ell)}).
                                                        \tag{A.22.1}
\]
The last group is the inertia of the cusp selected by the coset
$H\cdot1$ and the chosen base-point path.

In particular this computes the full fixed subgroup at $\ell=2$
for all three $D_8$ genus-three covers, the $Q_8$ genus-three cover,
and every genus-two reflection quotient of a $D_8$ cover in A.4.
If $c$ generates the original $I_2$, their fixed subgroup maps
to $\overline{\langle c^2\rangle}$ in $P_2$; intrinsically it is
the full pro-$2$ inertia of the selected cusp. These statements
hold for every packet section, including those with vanishing
degree-two Kummer character, and for every stated $\Delta_0$.
\end{statement}
\begin{proof}
We first verify the finite-index restriction of the fixed-group
input R1 (LM8 G.10). The following elementary observation applies
to any finite-rank free pro-$\ell$ group $P$. If an automorphism
$\alpha$ acts trivially on $P^{\mathrm{ab}}$ and
$\alpha^N(u)=u$ for a positive integer $N$, then $\alpha(u)=u$.
Indeed $\alpha$ acts trivially on every lower-central graded
piece, because these are generated from degree one by brackets.
They are torsion-free $\mathbf Z_\ell$ modules, by the integral
Magnus realization of the free Lie algebra. If
$v=\alpha(u)u^{-1}\ne1$, choose its first nonzero lower-central
class $v_d$. Modulo the next term, the product giving
$\alpha^N(u)u^{-1}$ has class $Nv_d$, which cannot vanish.
The lower-central filtration is separated: every finite
$\ell$-group quotient is nilpotent. This proves the observation,
also when $\ell$ divides $N$.

CM over $K$ makes the action of $\Delta$ on
$P_\ell^{\mathrm{ab}}=T_\ell E$ trivial. If $u$ is fixed by an
open subgroup $\Delta_0$, every $g\in\Delta$ has a positive
power in $\Delta_0$ (one may first take its finite-index normal
core). Apply the observation to the action of $s_b(g)$.
It follows that
\[
           P_\ell^{s_b(\Delta_0)}
             =P_\ell^{s_b(\Delta)}=I_\ell.
                                                        \tag{A.22.2}
\]
The final equality is the all-prime packet statement R1,
whose attached proof LM8 G.10 separates the genuine degree-two
inner case from the nonzero higher metabelian-detected
derivation case. The finite-index step just proved needs
neither a nonzero degree-two coefficient nor a normal
$\Delta_0$; it does not assume that fixed points commute with
pro-$\ell$ completion.

We next prove the claimed completion comparison. Put
$C=\bigcap_{g\in\pi}gHg^{-1}$, so $\pi/C$ is a finite
$\ell$-group. Given an open normal subgroup $N\triangleleft H$
with $H/N$ a finite $\ell$-group, $N\cap C$ is normal in $C$.
Intersect its finitely many conjugates under $\pi/C$ to obtain
$M\triangleleft\pi$, with $M\subset N$.
The quotient $C/M$ embeds in a finite product of finite
$\ell$-groups $C/g(N\cap C)g^{-1}$. Hence both $C/M$ and
$\pi/M$ are finite $\ell$-groups. This proves that the
pro-$\ell$ topology of $H$ is induced from that of $\pi$.
Equivalently, writing $R_\ell(\pi)$ for the kernel of
$\pi\to P_\ell$, one has
\[
            R_\ell(H)=R_\ell(\pi)\subset C\subset H.
\]
Thus $H^{(\ell)}=H/R_\ell(\pi)$ identifies with
$\overline H\subset P_\ell$, and its index is $[\pi:H]$.
The identification is equivariant for every automorphism
preserving $H$, in particular for $s_b(\Delta_0)$.
Intersecting (A.22.2) with $\overline H$ now gives the first
equality of (A.22.1).

Since $H$ is the inverse image of $\overline H$ under
$\pi\to P_\ell$, every element of $\overline H\cap I_\ell$
has a lift in $H\cap I$. The converse inclusion is immediate.
The geometric covering interpretation of the latter subgroup
is precisely the inertia of the selected cusp. This proves
the second equality and identifies the actual subgroup, not
only its conjugacy class after a further quotient.

For the listed applications the Galois closure is $D_8$ or
$Q_8$, so it is a $2$-group. Their puncture monodromy, the
commutator of a generating pair, is the nontrivial central
element of order two. It is absent from a reflection subgroup
of $D_8$, so it has orbits of length two also on the four
cosets defining each genus-two quotient. Thus the ramification
index of every displayed cusp over $O$ is exactly two.
The original inertia generator therefore enters each $H$
first as $c^2$, proving the final formula. The independent
finite certificate \texttt{code/a\_fixed\_subcovers.py}
checks the commutators of all generating pairs in these two
groups and the reflection-coset orbit lengths.
\end{proof}

\begin{statement}{A.24}{Proved}{Generic packet twists on the stable symmetric covers}
Let
\[
 C:\ v^2=w^7-w,\qquad
 D:\ z^2=(t^2+4)(t^3+3t),
\]
with the two cusps $P_0,P_\infty$ removed from $C$ and the single
cusp $Q_\infty$ removed from $D$. Use the maps
$C\to D\to E$ and the selected cusp paths of A.4. Fix a prime
$\ell$, any original packet section $s_b$ on $E\setminus O$,
and a rational prime $a\ne\ell$. Let
$\kappa_{a,\ell}:G_K\to\Z_\ell(1)$ be its Kummer cocycle, and
regard this as a cocycle in the $\ell$-component of
$I=\prod_p\Z_p(1)$. For $n\in\Z_\ell$ put
\[
 s_n(g)=c^{\,n\kappa_{a,\ell}(g)}s_b(g),
                                                        \tag{A.24.1}
\]
where all components of the displayed exponent other than the
$\ell$-component are zero.

There is one finite extension $L/K$, independent of $n$, such
that all these sections restrict to the selected cusp packets
on both covers over $G_L$. Put $\Delta_0=\Delta\cap G_L$.
Write $P_{C,\ell}$ and $P_{D,\ell}$ for the full pro-$\ell$
geometric groups of these two open curves, and $I_{C,\ell}$,
$I_{D,\ell}$ for their selected cusp inertia groups.

There is a set $B_{C,\ell}\subset\Z_\ell$ of cardinality at most
three such that, for every $n\notin B_{C,\ell}$ and every open
subgroup $\Delta_1\subset\Delta_0$,
\[
             P_{C,\ell}^{\,s_n(\Delta_1)}=I_{C,\ell}.
                                                        \tag{A.24.2}
\]
For $D$ the same assertion holds outside an at most countable
set $B_{D,\ell}$:
\[
             P_{D,\ell}^{\,s_n(\Delta_1)}=I_{D,\ell}.
                                                        \tag{A.24.3}
\]
At $\ell=2$ one may take $B_{D,2}=B_{C,2}$, so the bound is
again three. At every prime the complement of the indicated
exceptional set is dense in $\Z_\ell$ and has Haar measure one.
The sets are explicit in terms of the actual first nontrivial
Johnson matrices described in the proof. No value of those
arithmetic matrices, no membership decision for a specified
packet, and no assertion for all exceptional twists is made.
\end{statement}
\begin{proof}
\paragraph{A common arithmetic group and the factor $1/3$.}
The two maps to the base have ramification index three at the
selected cusps. For $C\to E_{0,-1}$, the coordinate map is
$(w^3,wv)$; its local parameter $x/y=w^2/v$ has order three
at $P_\infty$, and the map to $E$ is an unramified isogeny.
The involution $j$ giving $C\to D$ is fixed-point-free, so this
map is etale, also at the cusps, and the index for $D\to E$
is also three. With compatible generators the actual full
inertia inclusions are therefore
\[
                 c_C\longmapsto c^3,\qquad
                 c_D\longmapsto c^3.                    \tag{A.24.4}
\]
The image of $s_b$ lies in the arithmetic subgroup of the
chosen finite cover after restriction to a finite extension:
take the inverse image of that open subgroup under $s_b$.
The additional twist in (A.24.1) lies in the cover subgroup
automatically when $\ell\ne3$. When $\ell=3$, impose the one
further finite condition $\kappa_{a,3}=0\bmod3$.
Choose $L$ also to define the cusps, the tangent vectors, and
the isogenies $J_C\sim E^3$, $J_D\sim E^2$ of A.4.
These finitely many choices give a common $L$ for every $n$.
On $\Delta_0$ the induced twist in either cover is exactly
\[
                  c_C^{\,n\kappa_{a,\ell}/3}
             \quad\hbox{or}\quad
                  c_D^{\,n\kappa_{a,\ell}/3}.             \tag{A.24.5}
\]
At $\ell=3$ the numerator has been made divisible by three
in $\Z_3$; at all other primes three is a unit. Thus (A.24.5)
is an equality of integral pro-$\ell$ group actions, not
only a rational formula.

\paragraph{Tangential equivariance and trivial homology action.}
Let $h$ be the hyperelliptic involution of either curve.
It fixes the selected cusp, and its derivative there is $-1$.
There is a path in the punctured tangent line from a selected
tangent vector to its negative whose successive root
coordinates are roots of unity. All these coordinates are
fixed by $\Delta$, since $\mu_\infty\subset K^{\rm ab}$.
Use this path to make $h$ a based geometric automorphism.
Functoriality of the tangential fibre functor now gives an
actual commutation relation with the tangential section on
$\Delta_0$. This based automorphism fixes the chosen inertia
generator, so it also commutes with every twist of that
section by an inertia cocycle. It consequently commutes with
each packet action used here.

No assertion that this based lift has order two is needed.
Its square can be an inner automorphism by a cusp power.
Such an inner automorphism acts trivially on the
lower-central graded Lie algebra, and the induced action on
homology is the usual hyperelliptic one.

On $C$, that action is $+1$ on the one-dimensional boundary
space and $-1$ on proper homology. Over $\Q_\ell$ the two
eigenspaces split the homology extension equivariantly.
The isogeny $J_C\sim E^3$ makes the proper action of $\Delta_0$
trivial; the boundary action is the cyclotomic character,
also trivial on $\Delta_0$. Thus $\Delta_0$ acts trivially on
all of $H_1(C\setminus\{P_0,P_\infty\},\Z_\ell)$, since this
lattice is torsion-free. This argument does not divide an
integral lattice by two when $\ell=2$; the splitting is used
over $\Q_\ell$ only. For $D$, its one cusp contributes no
boundary homology, and $J_D\sim E^2$ gives the same
triviality directly. Every action in question is therefore
an IA automorphism of a free pro-$\ell$ group.

\paragraph{A nonzero Kummer value on $\Delta_0$.}
The restriction $\kappa_{a,\ell}|_\Delta$ is nonzero. To see
this directly, take a place $v$ of $K$ over $a$. Tame
$\ell$-inertia at $v$ has a generator $\tau$, and an
arithmetic Frobenius $F$ acts on it by the residue
cardinality $q_v$. The $\ell$-adic Kummer value of $a$ on
$\tau$ is its nonzero valuation $e_v$ times a chosen tame
generator. On the local commutator $[F,\tau]$, its value is
\[
                        (q_v-1)e_v\ne0
                       \quad\hbox{in }\Z_\ell(1).
                                                        \tag{A.24.6}
\]
Possible wild-inertia components do not affect this
calculation, because the residue characteristic is
$a\ne\ell$. The local commutator is in the global closed
commutator subgroup $\Delta$. On $\Delta$ the Kummer
cocycle is an additive character. A positive power of this
element belongs to the finite-index subgroup $\Delta_0$
and still has nonzero Kummer value. Fix such a
$\delta\in\Delta_0$, and put
\[
                         \beta=\kappa_{a,\ell}(\delta)/3
                              \in\Z_\ell\setminus\{0\}.
                                                        \tag{A.24.7}
\]
A compatible identification of $\Z_\ell(1)$ has been used
on $\Delta_0$ throughout.

\paragraph{The ordinary first Johnson map on $C$.}
Use the ordinary lower-central filtration $\gamma_d$ on
$P_{C,\ell}$. Its rational graded Lie algebra is
\[
                  \mathbb L(\Q_\ell t\oplus V),\qquad
                         \dim V=6,
\]
where $t$ is the class of $c_C$ and $V$ is the negative
hyperelliptic eigenspace. The class $t$ is primitive in
integral homology: a punctured genus-three surface with
two punctures has a free presentation in which either
selected cusp is a free generator.
The ordinary degree of $t$ is one. Its natural geometric
weight is two; these are different filtrations, and no
identification of them is used.

For the action of $s_b(\delta)$, the first Johnson map is
the degree-one derivation induced by
$g\mapsto s_b(\delta)g s_b(\delta)^{-1}g^{-1}$
modulo $\gamma_3$. It kills $t$, because the actual cusp
generator is fixed. Equivariance for $h$ forces it to send
$V$ into the negative eigenspace of
$\Lambda^2(\Q_\ell t\oplus V)$, namely $[t,V]$.
It therefore has the form
\[
             D_0(t)=0,\qquad D_0(v)=[t,A_\delta v],
                \qquad A_\delta\in\operatorname{End}_{\Q_\ell}(V).
\]
The integral twist (A.24.5) changes this to
\[
       D_n(t)=0,\qquad
       D_n(v)=[t,(A_\delta+n\beta\,1_V)v].               \tag{A.24.8}
\]
In geometric weights this part raises weight by two, which
explains why it appears in ordinary Johnson degree one.
There is also a useful constraint from the other cusp.
Let $\omega\in\Lambda^2V$ be its nondegenerate symplectic
class, written in the tensor algebra as
$\omega=\sum J_{ij}v_iv_j$ for a nonsingular skew matrix
$J$. We claim
\[
                         A_\delta J=J A_\delta^T.
                                                        \tag{A.24.13}
\]
For this calculation only, use the rational completed
free Lie algebra with $t=\log c_C$ as one generator.
The logarithm of a representative $c'$ of the other
cusp has the form
\[
               \log c'=-t-\omega+[t,u]
                         +\text{terms of degree at least three}.
\]
The term $[t,u]$ allows for the choice of the proper
homology complement and of the path to that cusp.
The packet action preserves the conjugacy class of
$c'$ and its exponent, since $\chi_\ell(\delta)=1$.
If $f_1$ is the linear term of a conjugator, comparison
in degree two gives $[f_1,t]=0$, hence $f_1=a t$.
Comparison in degree three, using that the action
fixes the actual $t$, now gives
$D_0(\omega)\in[t,\mathbb L_2(\Q_\ell t\oplus V)]$.
The extra term $D_0([t,u])$ is itself in that subspace.
Project to words containing one $t$ and two letters of
$V$. The words $v_i t v_j$ therefore have coefficient
zero. Their coefficient matrix in $D_0(\omega)$ is
$-A_\delta J+J A_\delta^T$, which proves (A.24.13).
This is the peripheral constraint on the actual action,
not a condition imposed on a freely chosen derivation.

It follows that $(A_\delta+n\beta\,1_V)J$ is skew.
The polynomial
\[
 p_\delta(n)=
 \frac{\operatorname{Pf}((A_\delta+n\beta\,1_V)J)}
      {\operatorname{Pf}(J)}
 \quad\hbox{satisfies}\quad
 \det(A_\delta+n\beta\,1_V)=p_\delta(n)^2.
                                                        \tag{A.24.14}
\]
Here $\operatorname{Pf}$ is the Pfaffian of a six-by-six
skew matrix. It has degree three in its entries, so
$p_\delta$ has degree three in $n$ and leading
coefficient $\beta^3\ne0$.
Let $B_{C,\ell}$ be its roots in $\Z_\ell$.

We prove the Lie-kernel assertion needed here in all
degrees. For a vector space $V$ over a characteristic-zero
field and an invertible $A$, let
$D(t)=0$, $D(v)=[t,Av]$. Then
\[
                    \ker D=\Q_\ell t.                  \tag{A.24.9}
\]
Indeed let $V_j=\operatorname{ad}_t^jV$ for $j\ge0$.
The ideal generated by $V$ is freely generated by
$W=\bigoplus_{j\ge0}V_j$. One can verify this elimination
identity without an additional relation theorem:
on $\mathbb L(W)$ the shift $V_j\to V_{j+1}$ is a
derivation; adjoining $t$ acting by that shift gives the
universal Lie algebra on $t$ and $V_0$.
The derivation under consideration satisfies
$D(v_j)=(Av)_{j+1}$.
Embed $\mathbb L(W)$ in the free associative algebra.
For a nonzero polynomial choose its largest occurring jet
index $j$. In its derivative, the terms containing a new
letter from $V_{j+1}$ come only from replacing a letter in
$V_j$. If those terms vanished, substitute
$V_{j+1}\to A^{-1}V_j$. The result is the Euler operator
counting the occurrences of $V_j$ in each word.
It cannot annihilate the part involving $V_j$ in
characteristic zero. Thus the derivative is nonzero on the
whole ideal. This proves (A.24.9).

If $u\ne1$ is fixed by $s_n(\delta)$, its first nonzero
lower-central class is killed by $D_n$. When
$n\notin B_{C,\ell}$, (A.24.9) forces that class to be
$a t$ in degree one. Its coefficient is in $\Z_\ell$,
because $t$ is primitive in the integral lattice.
Subtract the actual fixed cusp power $c_C^a$.
Any nonidentity remainder would have first degree at least
two, contradicting (A.24.9). The filtration is separated,
so $u=c_C^a$. This is an argument about actual pro-$\ell$
elements, with no lifting of a rational Malcev element.
It proves (A.24.2) for $\Delta_0$.
For any open $\Delta_1$, take a positive power
$\delta^N\in\Delta_1$; its first Johnson map is $ND_n$,
which has the same kernel. This proves the assertion for
every stated open subgroup.

\paragraph{The ordinary degree-two pencil on $D$.}
Now $P_{D,\ell}$ is free pro-$\ell$ of rank four. Write
$W=H_1(D\setminus Q_\infty,\Q_\ell)$ and choose a symplectic
integral basis so that the initial cusp class is
\[
                 \omega=[x_1,y_1]+[x_2,y_2]
                      \in\operatorname{gr}_2P_{D,\ell}.
\]
This class is primitive over $\Z_\ell$, and the cusp has
no degree-one term. The hyperelliptic involution acts as
$-1$ on $W$ and as $+1$ on $\mathbb L_2(W)$.
Its commutation with the packet action therefore forces
the degree-one Johnson map to vanish. The degree-two map
for $s_n(\delta)$ is consequently a well-defined
homogeneous derivation of the form
\[
          T_n=T_0+n\beta\,\operatorname{ad}(\omega),
                         \qquad T_n(\omega)=0.           \tag{A.24.10}
\]
This formula uses the actual first degree-two term of the
action and asserts no formula for $T_0$.

For clarity, the required centralizer fact has a short
direct proof:
\[
                C_{\mathbb L(W)}(\omega)=\Q_\ell\omega.
                                                        \tag{A.24.11}
\]
In the tensor algebra order the letters so that the
leading word of $\omega$ is $x_1y_1$.
This word is not a proper power. If a homogeneous
associative polynomial $f$ commutes with $\omega$, their
leading words commute, so the leading word of $f$ is
$(x_1y_1)^m$. Subtracting the matching multiple of
$\omega^m$ would leave a commuting polynomial with a
strictly smaller leading word, which is impossible by
the same word argument. Thus $f$ is a scalar multiple of
$\omega^m$. A power $\omega^m$ with $m>1$ is not primitive
for the tensor Hopf coproduct in characteristic zero,
whereas every Lie element is primitive. This proves
(A.24.11).

In each degree $d\ne2$,
$\operatorname{ad}(\omega):\mathbb L_d(W)\to
\mathbb L_{d+2}(W)$ is therefore injective. In degree two
use the source $\mathbb L_2(W)/\Q_\ell\omega$.
For each of these finite-dimensional source spaces
choose a full-rank minor which is nonzero for
$\operatorname{ad}(\omega)$.
The corresponding minor for (A.24.10) is a polynomial
in $n$ with nonzero leading coefficient. Let $B_{D,\ell}$
be the union, over all $d\ge1$, of its roots in $\Z_\ell$.
This set is at most countable. Outside it,
\[
                         \ker T_n=\Q_\ell\omega.
                                                        \tag{A.24.12}
\]
The first nonzero class of a fixed pro-$\ell$ element is
then $a\omega$ in degree two. Primitivity makes
$a\in\Z_\ell$; subtract the actual cusp power $c_D^a$.
A nonidentity remainder has first degree at least three
and contradicts (A.24.12). This proves (A.24.3).
Replacing $\delta$ by a positive power proves it for
all open $\Delta_1$ just as above.

\paragraph{The sharper bound at two and the scope.}
For $\ell=2$ the etale double cover $C\to D$ induces an
embedding of $P_{C,2}$ as an open subgroup of $P_{D,2}$,
by the pro-prime topology argument in A.22.
If $u\in P_{D,2}$ is fixed, then $u^2\in P_{C,2}$ is
fixed. Outside $B_{C,2}$ it belongs to the selected
$I_{C,2}$, whose image is $I_{D,2}$ since this double
cover is unramified at the cusp.
The latter inertia is root-closed, as follows here
directly from (A.24.11). If $v$ commutes with
$c_D^b\ne1$, its first lower-central class commutes with
$b\omega$. It is therefore $a\omega$ in degree two,
with $a\in\Z_2$ by primitivity. Subtracting $c_D^a$
leaves an element still commuting with $c_D^b$ and
having first degree at least three, which (A.24.11)
excludes. Hence $C_{P_{D,2}}(c_D^b)=I_{D,2}$.
If $u^2=c_D^b\ne1$, apply this to $u$ itself.
If $u^2=1$, any first nonzero lower-central class of
$u$ would be killed by two in a torsion-free
$\Z_2$-module, so $u=1$ as well.
This proves (A.24.3) outside $B_{C,2}$.

Finite and countable subsets of $\Z_\ell$ have Haar
measure zero, and their complements are dense because
every nonempty open ball is uncountable.
The independent certificate
\texttt{code/a\_generic\_packets.py} checks the tangent
signs, the ramification orders, and finite-degree Lie
ranks for deliberately resonant invertible matrices.
It does not compute $A_\delta$ or $T_0$ and is not used
to substitute for the unbounded Lie proofs above.
\end{proof}

\begin{statement}{A.25}{Proved}{The stable genus-two normalizer at every odd prime}
Let
\[
 D:\ z^2=(t^2+4)(t^3+3t),\qquad
 Y=D\setminus\{Q_\infty\},
\]
and let $\ell$ be odd. Every nontrivial closed procyclic
subgroup of $\pi_1(Y_{\overline K})^{(\ell)}$ whose
arithmetic normalizer has open image in $G_K$ is contained
in a conjugate of cusp inertia.
Together with A.14 at $\ell=2$, this completes the
pro-$\ell$ normalizer classification for every first cover
and every displayed genus-two quotient in Task A.4.

The extra detector is the degree-two separating monodromy
at a place above $2$, acting on the proper genus-two group.
It is an actual arithmetic derivation and does not require
a second cusp-compatible elliptic map or a lifting through
the odd-prime quotient $P_{C,\ell}\twoheadrightarrow P_{D,\ell}$.
\end{statement}
\begin{proof}
\paragraph{The degree-three map.}
The map to the original elliptic curve is
\[
 f:D\longrightarrow E,\qquad
 x_E=-\frac{i}{2}(t^3+3t),\qquad
 y_E=-\frac{1+i}{4}z(t^2+1).
                                                        \tag{A.25.1}
\]
Its coordinates satisfy $y_E^2=x_E^3-x_E$.
The pole order of $x_E$ at the unique point at infinity
is six, so $\deg f=3$ and
\[
                         f^*(O)=3Q_\infty.              \tag{A.25.2}
\]
Riemann--Hurwitz then shows that the indicated totally
ramified cusp is all the ramification. In particular
$f:Y\to E\setminus O$ is finite etale. Formula (A.25.1)
also follows by descending the $C_3$ map in A.4:
$w^3-w^{-3}=t^3+3t$ and
$w^3+w^{-3}=(w+w^{-1})(t^2+1)$.

\paragraph{The stable reduction above two.}
Work over a finite extension of $K_2$ on which $E$ has
good reduction, $D$ has stable reduction, and all
components are defined. Such an extension exists.
The potential good reduction of $E$ follows from its
integral $j$-invariant; the exact criterion is
\cite[Proposition VII.5.5]{fghSilvermanAEC}.
Its geometric special fibre $E_0$ has $j=0$ in
characteristic two and is supersingular. For example,
over an algebraic closure its equation is
$y^2+y=x^3$, whose only geometric two-torsion point is
the origin. Hence there is no etale degree-two isogeny
onto $E_0$.

By A.4, $J_D\sim E^2$ after a finite extension.
Good reduction of an abelian variety is invariant
under isogeny, as follows from the
Neron--Ogg--Shafarevich criterion. Thus $J_D$ has good
reduction here. The dual graph of the stable fibre of
$D$ is a tree: for a nodal curve the torus in its
generalized Jacobian has rank the first Betti number
of the dual graph. Consequently the stable fibre is
either smooth of genus two or two smooth elliptic
curves meeting in one separating node. In genus two
there are no further stable compact-type possibilities:
a rational vertex of valence at least three would
require at least three positive-genus leaves.

The map $f$ extends to this stable model with values in
the good elliptic model $\mathcal E$. Here is the
extension argument, including the nodal case.
Resolve the rational map on a regular semistable model
by blowing up closed points. Every exceptional
geometric curve is rational and maps constantly to
the elliptic special fibre. The same is true of every
rational component contracted to obtain the stable
model. The map is therefore constant on the connected
fibres of the contraction. Its graph image over the
stable model is proper and quasi-finite, hence finite,
and birational. The stable total space is normal
(locally at a node it is $xy=\varpi^e$), so this finite
birational graph is the stable model itself. This
proves descent of the morphism.

Suppose first that the stable fibre $D_0$ were smooth.
The pullback of $\mathcal O_{\mathcal E}(O)$ has degree
three on every fibre, so $f_0:D_0\to E_0$ has degree
three and is nonconstant. There is no vertical
component in the Cartier divisor $f^*(O)$.
Equation (A.25.2) therefore extends as the relative
divisor identity $f^*(O)=3\overline{Q_\infty}$, and
\[
                         f_0^*(O)=3(Q_\infty)_0.
\]
The degree-three map is separable in characteristic
two. Its ramification index three at this point is
tame and contributes two to the different.
Riemann--Hurwitz gives total different degree two,
so there is no other ramification.

The Galois closure of this separable cubic map has
group $C_3$ or $S_3$. In the cyclic case a Kummer
description would require a rational function on
$E_0$ whose divisor has exactly one valuation
nonzero modulo three, contradicting degree zero.
In the $S_3$ case inertia at the sole branch point
is $C_3$, and the quotient by $A_3$ is a connected
etale double cover of $E_0$. It is a degree-two
isogeny after a choice of origin, contradicting
supersingularity. Thus $D_0$ is not smooth.

We have proved that
\[
                  D_0=E_1\cup_N E_2                    \tag{A.25.3}
\]
with both components elliptic. The sum of the degrees
of the two restrictions $f_i:E_i\to E_0$ is three,
by the same fibrewise degree calculation. Neither
restriction is constant. Indeed, suppose the degrees
were $(3,0)$. The nonconstant restriction is a
separable degree-three isogeny of elliptic curves,
up to translation, hence is etale and has three
distinct geometric preimages of $O$.
On the open part $E_1\setminus\{N\}$, no vertical
component of the Cartier divisor $f^*(O)$ is present:
the nonconstant component itself is not a component
of this divisor, and the only other component meets
it at $N$. Its horizontal part is exactly
$3\overline{Q_\infty}$. Thus the support of
$f_1^{-1}(O)$ outside $N$ has at most one point.
An etale degree-three isogeny has at least two of
its three preimages outside $N$, a contradiction.
This argument also covers specialization of the
cusp to the node or to the contracted component.
Interchanging the indices excludes $(0,3)$.
After relabelling, therefore,
\[
                         \deg f_1=1,\qquad \deg f_2=2.
                                                        \tag{A.25.4}
\]
Both restrictions are isogenies after translation.
Their maps on rational $\ell$-adic homology are
isomorphisms for every odd $\ell$, including when
the degree-two map in characteristic two is
inseparable.

\paragraph{The two nonzero components of a candidate line.}
Suppose $A$ is a noncuspidal arithmetically normalized
procyclic subgroup of the pro-$\ell$ group of $Y$.
By A.11 a generator has a nonzero proper homology
class $v$. The map (A.25.1), together with the
elliptic normalizer criterion A.12, gives
$f_*v=0$: a nonzero image would produce a
noncuspidal normalized procyclic subgroup on
$E\setminus O$.

Good reduction of the Jacobian and compact-type
specialization identify
\[
 H_1(D_{\overline K},\Q_\ell)=V_1\oplus V_2,\qquad
 V_i=H_1(E_i,\Q_\ell).
                                                        \tag{A.25.5}
\]
This identification is compatible with $f_*$, whose
restrictions are $(f_i)_*$. By (A.25.4) both are
isomorphisms. Hence
\[
                  v=v_1+v_2,\qquad v_1\ne0,\quad v_2\ne0.
                                                        \tag{A.25.6}
\]
For example, if $v_1=0$, the equality
$(f_1)_*v_1+(f_2)_*v_2=0$ and injectivity of
$(f_2)_*$ force $v_2=0$.

\paragraph{The actual monodromy and its Lie obstruction.}
Pass now to the proper pro-$\ell$ group of $D$.
The image of $A$ is still a nontrivial procyclic
group, since its homology class is $v\ne0$, and
its arithmetic normalizer still has open image.
This passage does not assert an injective comparison
from a prime-to-$\ell$ covering group.

The exact monodromy input is the non-abelian
Picard--Lefschetz theorem
\cite[Definition 3.1.5 and Theorem 3.1.6]{fghBettsDogra}.
For a regular semistable model it identifies the
action of tame inertia on the rational unipotent
fundamental group with edge twists of its reduction
graph. Resolving the node (A.25.3) merely subdivides
its one edge; its total length is a positive number
$m$. Choose a base point specializing to the second
elliptic component, after another finite extension
if needed. The twist fixes its vertex group and
conjugates the first vertex group by the $m$-th
power of the separating boundary. Its first
Johnson degree is two. With
\[
 \omega_1=[a,b],\qquad \omega_2=[c,d],\qquad
 \operatorname{gr}\mathfrak p_D
       =\mathbb L(a,b,c,d)/(\omega_1+\omega_2),
                                                        \tag{A.25.7}
\]
the degree-two derivation, up to the harmless common
nonzero tame scalar and an inner derivation, is
\[
        T(x)=m[\omega_1,x]\ (x\in V_1),\qquad
        T(x)=0\ (x\in V_2).                              \tag{A.25.8}
\]
The graph presentation in Definition 3.1.5 gives
(A.25.8) directly; it does not infer a Galois
derivation from an ambient formal possibility.

We record the necessary normalizer condition
carefully. An open arithmetic image contains a
nonzero power of a tame inertia element after all
the finite extensions made above. Inertia acts
trivially on the proper homology, since the
Jacobian has good reduction. Thus its conjugation
character on the proposed procyclic subgroup is
one. In the rational unipotent group, its standard
action must consequently carry the image of a
generator $u$ to a conjugate of $u$.
The action in (A.25.8) has zero Johnson degree one.
Comparing degree two of this conjugacy equation
gives $[q_1,v]=0$ in
$\Lambda^2(V_1\oplus V_2)/\Q_\ell(\omega_1+\omega_2)$,
where $q_1$ is the linear term of a conjugator.
A decomposable bivector cannot be a nonzero
multiple of the nondegenerate rank-four symplectic
form, so $q_1$ is a scalar multiple of $v$.
Multiplying the rational conjugator by the inverse
of the corresponding rational power of $u$ removes
that term. This is only a necessary equation in a
rational quotient, not a claimed integral lift.
Comparison in degree three now yields
\[
            T(v)\in[v,\operatorname{gr}_2\mathfrak p_D].
                                                        \tag{A.25.9}
\]
Changing the initial representative of the outer
monodromy adds an inner derivation and does not
change this necessary condition.

Choose symplectic bases with $a=v_1$ and $c=v_2$;
this is possible by (A.25.6). There is a Lie map
from the algebra in (A.25.7) to the free Lie
algebra on $x,y$ defined by
\[
                 a\mapsto x,\quad c\mapsto-x,\quad
                 b\mapsto y,\quad d\mapsto y.
                                                        \tag{A.25.10}
\]
It kills both $v=a+c$ and the relation
$[a,b]+[c,d]$. But it carries $T(v)$ to
\[
                          m[[x,y],x]\ne0.
                                                        \tag{A.25.11}
\]
It therefore kills the right side of (A.25.9)
and not the left, a contradiction. This excludes
the proposed subgroup at every odd prime.

All nontrivial closed subgroups of cusp inertia
do have open arithmetic normalizer after the
finite extension making that cusp rational.
This proves the stated classification. The
independent exact certificate
\texttt{code/a\_D3\_monodromy.py} verifies (A.25.1),
the origin-fibre pole orders, and the nonzero
free-Lie witness (A.25.10)--(A.25.11).
\end{proof}

\begin{statement}{A.26}{Proved}{All odd-prime stable packets lose the proper first layer}
Let $\ell$ be odd and let $s_b$ be any original cusp
packet section in Part A. On the stable $S_3$ cover
$C:v^2=w^7-w$ and its three genus-two reflection
quotients, restrict to any open subgroup $\Delta_0$
over which the selected cusp lifts and packet
actions are defined. Then
\[
 \operatorname{im}\!\left(
    P_{D,\ell}^{s_b(\Delta_0)}
       \longrightarrow H_1(D_{\overline K},\Z_\ell)
                    \right)=0,                         \tag{A.26.1}
\]
and
\[
 \operatorname{im}\!\left(
    P_{C,\ell}^{s_b(\Delta_0)}
       \longrightarrow H_1(C\setminus\{P_0,P_\infty\},
                           \Z_\ell)
                    \right)
              =\Z_\ell[c_C].                           \tag{A.26.2}
\]
The first formula refers to each of the three
reflection quotients. Thus a fixed element on $C$
can be multiplied by an actual fixed cusp power
to put it in $\gamma_2P_{C,\ell}$; a fixed element
on $D$ already lies in $\gamma_2P_{D,\ell}$.
These conclusions include every exceptional packet
left by the determinant tests in A.24.
They do not identify the full nonabelian fixed
subgroup in those exceptional cases.
\end{statement}
\begin{proof}
It suffices first to treat the model $D$ of A.25.
Let $u$ be fixed and let $v$ be its rational proper
homology class. Under the degree-three punctured
map $f:Y\to E\setminus O$, its image is fixed by
the restriction of the original $s_b$.
The input R1, with its finite-index extension
proved in A.22, puts that image in the selected
elliptic cusp inertia. Consequently $f_*v=0$.
The integral conclusion will follow from the
rational one because proper homology is
torsion-free over $\Z_\ell$.

Use the stable reduction above two proved in A.25.
If $v\ne0$, the component degrees $(1,2)$ force
its projections to both elliptic components to
be nonzero. Unlike the normalizer proof, we must
now exhibit the needed monodromy inside
$\Delta_0$, not merely inside an open arithmetic
group. Choose a finite number field extension
whose intersection with $\Delta$ is contained
in $\Delta_0$, and enlarge its completion at two
until the stable model and its components are
defined. This does not change the two positive
component degrees.

Let $\tau$ be a tame $\ell$-inertia element with
nonzero monodromy and $F$ a Frobenius element in
this local Galois group. On the semistable
prime-to-two outer fundamental group, the
commutator $[F,\tau]$ has monodromy equal to
$(q-1)$ times that of $\tau$, where $q>1$ is the
residue cardinality. This follows either from
the tame relation $F\tau F^{-1}=\tau^q$ modulo
wild inertia, or from the edge-twist description
in \cite[Definition 3.1.5 and Theorem 3.1.6]{fghBettsDogra}.
Wild inertia has trivial semistable outer
prime-to-two action. The scalar $q-1$ is nonzero
in $\Q_\ell$, whether or not it is an $\ell$-adic
unit.

The local commutator lies in the global closed
commutator subgroup and in the chosen finite
extension subgroup; hence it is in $\Delta_0$.
The fixedness of $u$ makes its image in the
proper rational unipotent group invariant up
to conjugacy for this monodromy. Changing from
the original cusp base point to the smooth
component base point in A.25 changes the
representative by an inner automorphism only.
The degree-two and degree-three comparison
in that proof therefore gives the same
necessary condition
\[
                 T(v)\in[v,\operatorname{gr}_2
                                      \mathfrak p_D].
\]
The explicit map (A.25.10) kills the right side
and maps the left to the nonzero scalar multiple
$(q-1)m[[x,y],x]$. This contradiction proves
$v=0$, and then (A.26.1).

The other two reflection quotients are obtained
by conjugating the reflection by the cubic
rotation of $C$. After one common finite field
extension the same argument applies to their
induced packet actions and the original map
to $E\setminus O$.
The three pushforward maps from proper
$H_1(C,\Q_\ell)$ to these quotient homologies
are jointly injective. Indeed A.4 gives the
$S_3$ representation as a two-dimensional
trivial part and a standard two-dimensional
representation tensored with a two-dimensional
elliptic homology space. The quotient by a
reflection detects its positive eigenspace.
The positive reflection eigenspaces span the
standard representation, and the trivial part
is detected by every quotient. Equivalently
the sum of the three operators $1+j$ is
invertible on this representation.

The proper homology image of a fixed element
of $P_{C,\ell}$ therefore vanishes. Its open
homology class lies in the boundary lattice,
which is precisely the primitive rank-one
lattice $\Z_\ell[c_C]$. Every cusp power is
fixed by $\Delta_0$, so equality holds in
(A.26.2), and subtraction of that actual
integral cusp power has the claimed effect
on the lower-central filtration.
\end{proof}

\begin{statement}{A.27}{Proved}{Every odd-prime packet on the stable genus-two cover}
Let
\[
 D:\ z^2=(t^2+4)(t^3+3t),\qquad
 Y=D\setminus\{Q_\infty\},
\]
and let $\ell$ be odd. For every original packet section $s_b$
in Part A and every open subgroup $\Delta_0$ over which its
selected lift to $Y$ is defined,
\[
                 P_{D,\ell}^{\,s_b(\Delta_0)}=I_{D,\ell}.
                                                        \tag{A.27.1}
\]
The same assertion holds for all three reflection quotients
of the stable $S_3$ cover, with their selected cusps. In
particular, at odd primes none of the exceptional packet
parameters left for $D$ in A.24 is exceptional for the actual
fixed-subgroup conclusion.

The algebraic input is an unbounded kernel calculation.
Put $L=\mathbb L_k(a,b,c,d)$ over any characteristic-zero
field, and write
\[
 \omega_1=[a,b],\quad \omega_2=[c,d],\quad
 \Omega=\omega_1+\omega_2.
\]
Let $T_i$ act by $\operatorname{ad}(\omega_i)$ on the two
generators of its own component and by zero on the other
component. For $m_1m_2\ne0$ and any $\lambda\in k$, set
\[
             D_\lambda=m_1T_1+m_2T_2
                                +\lambda\operatorname{ad}(\Omega).
                                                        \tag{A.27.2}
\]
Then
\[
 \ker(D_\lambda:L\longrightarrow L)=
 \begin{cases}
   \mathbb L_k(\omega_1,\omega_2),&\lambda=0,\\
   k\Omega,&\lambda\ne0.
 \end{cases}                                           \tag{A.27.3}
\]
The arithmetic proof uses the full separating monodromy
and one of its actual Galois conjugates.
\end{statement}
\begin{proof}
\paragraph{A tensor-algebra normal form.}
First prove (A.27.3). Work in
\[
 A=k\langle a,b,c,d\rangle=A_1*A_2,
 \qquad A_1=k\langle a,b\rangle,
 \quad A_2=k\langle c,d\rangle,
\]
and put $C=k\langle\omega_1,\omega_2\rangle$.
These two quadratic generators freely generate their
associative subalgebra: their leading words $ab$ and $cd$
give distinct leading words for distinct products.

Order words in $a,b$ lexicographically with $a>b$, within
each fixed length. The leading word of $\omega_1$ is $ab$,
which has no proper self-overlap. Let $\mathcal B_1$ be the
set of nonempty words that neither begin nor end with $ab$.
Every word that is not a power of $ab$ has a unique expression
\[
                        (ab)^r w(ab)^s,
 \qquad w\in\mathcal B_1,
\]
obtained by removing its maximal initial and terminal
powers of $ab$. The absence of a self-overlap guarantees
that these removals do not create an additional ambiguity.
Consequently the elements
\[
 \{\omega_1^n:n\ge0\}\ \sqcup\
 \{\omega_1^r w\omega_1^s:
              w\in\mathcal B_1,\ r,s\ge0\}
                                                        \tag{A.27.4}
\]
are a basis of $A_1$: in each finite-dimensional homogeneous
piece their leading words give a triangular matrix with
diagonal entries one. Thus, as a bimodule,
\[
 A_1=k[\omega_1]\ \oplus\!!\bigoplus_{w\in\mathcal B_1}
                          k[\omega_1]w k[\omega_1].
                                                        \tag{A.27.5}
\]
The analogous statement holds for $A_2$, using $cd$ and
$\mathcal B_2$. The displayed summands with $w$ are free
bimodules; no finiteness assertion for $\mathcal B_i$ is
being made.

Apply these two bases to the usual alternating-block
basis of the free product $A_1*A_2$. A basis word either
belongs to $C$, or has a first nonboundary core
$w\in\mathcal B_i$. All the pure $\omega$-blocks preceding
that core, together with the left power of $\omega_i$
in its own block, form an arbitrary word $u$ in $C$.
What follows is a tail beginning with that fixed core.
This is a unique normal form. In particular there are
tail spaces $M_i$ and a vector-space direct sum
\[
                    A=C\ \oplus\ (C\otimes M_1)
                           \ \oplus\ (C\otimes M_2).
                                                        \tag{A.27.6}
\]
Here a tail begins with a nonboundary core of type $i$,
and the tensor notation records its uniquely determined
prefix, not an assertion that $A$ is a free algebra over $C$.

Give a prefix its length as a word in the two letters
$\omega_1,\omega_2$. The derivation
$D_0=m_1T_1+m_2T_2$ kills every pure boundary block.
On the first core of type $i$ it is
\[
                      D_0(w)=m_i(\omega_iw-w\omega_i).
\]
The first term increases the prefix length by one, by
right multiplication of the prefix by $\omega_i$.
The second term and the action on all later cores preserve
that prefix length. These operations cannot remove or
change the first core: multiplying by $\omega_i$ changes
only the corresponding left or right power in (A.27.4).
Likewise the left part of $\lambda\operatorname{ad}(\Omega)$
increases prefix length by one, whereas its right part
preserves it. Appending a pure boundary word does not
remove a nonboundary core.

It follows that the part of $D_\lambda$ increasing the
prefix length on $C\otimes M_i$ is
\[
           (\lambda L_\Omega+m_i R_{\omega_i})\otimes1.
                                                        \tag{A.27.7}
\]
This operator is injective on $C$. If $\lambda=0$, this
is just nonzero right multiplication in a free algebra.
If $\lambda\ne0$, let $j\ne i$ and take the largest
number of occurrences of $\omega_j$ in a nonzero finite
linear combination of words. In its image under (A.27.7),
the terms with one more such occurrence come only from
$\lambda L_{\omega_j}$, and cannot cancel. This also
proves injectivity after tensoring with any tail space.

Now choose the largest prefix length in a nonzero element
of $C\otimes M_i$. Its image of prefix length one greater
is nonzero by (A.27.7); terms of smaller prefix length
cannot contribute there. Thus no nonzero element involving
a nonboundary core is killed by $D_\lambda$, and
\[
                         \ker(D_\lambda|_A)\subset C.
                                                        \tag{A.27.8}
\]
On $C$ the derivation is simply
$\lambda\operatorname{ad}(\Omega)$. For $\lambda=0$ its
kernel is all of $C$. For $\lambda\ne0$ its kernel is
$k[\Omega]$: replace the two free generators of $C$ by
$\Omega,\omega_1$, and use the elementary centralizer
calculation for a free generator. Finally $\omega_1$ and
$\omega_2$ are primitive in the tensor Hopf algebra.
The primitive elements of $C$ form
$\mathbb L_k(\omega_1,\omega_2)$, and those of
$k[\Omega]$ form $k\Omega$. Intersecting (A.27.8) with
$L$ proves (A.27.3) in every degree.

\paragraph{The cusp lies on the rational bridge.}
Use the stable reduction of A.25 at a place above $2$.
Its two elliptic components $E_1,E_2$ have degrees $1,2$
under the extension of
\[
 f:D\longrightarrow E,\qquad
 x_E=-i(t^3+3t)/2,\quad
 y_E=-(1+i)z(t^2+1)/4,
\]
and $f^*(O)=3Q_\infty$ on the generic fibre. Neither
component is contracted, so the Cartier divisor $f^*(O)$
on the normal stable total space has no vertical component.
Its Weil divisor is therefore $3\overline{Q_\infty}$.
If the cusp section met a smooth point of just one
elliptic component, the support of this divisor would
miss the other component entirely. The restriction of
$f$ to that other elliptic component is nonconstant and
surjective, so it must have a preimage of $O$, a
contradiction. Hence the cusp specializes to the node.

The stable pointed model consequently inserts a rational
component carrying the cusp between $E_1$ and $E_2$.
After a further finite extension and regular subdivision,
the distances from this rational component to the two
elliptic components are both strictly positive. Denote
their total edge lengths by $e_1,e_2>0$. Choose an ordinary
base point reducing to a smooth point of the rational
component away from its two nodes and the cusp. Such a
point exists after a finite unramified extension and lifts
by smoothness.

For this affine curve and ordinary base point,
\cite[Definition 3.1.5 and Theorem 3.1.6]{fghBettsDogra}
gives the based edge-twist description. The rational root
component has three boundary loops: the cusp and the two
nodes. Its vertex group is fixed by monodromy. Transport
the two elliptic handle pairs to that root. The first
nonzero ordinary Johnson derivation of tame inertia is
\[
                            e_1T_1+e_2T_2,
                                                        \tag{A.27.9}
\]
with a possible simultaneous sign depending on the
chosen inertia convention. Each $T_i$ fixes both
$\omega_1,\omega_2$, and hence the cusp class $\Omega$.
There is no first-degree Johnson term: every conjugating
boundary word starts in lower-central degree two.

\paragraph{Actual packet monodromy and its only inner term.}
Choose a finite number field $L$ defining all these maps,
cusps and CM isogenies, such that $\Delta\cap G_L$ is
contained in the prescribed $\Delta_0$. Enlarge $L$ to
realize the local semistable situation above, and put
\[
                      \Gamma=\overline{[G_L,G_L]}.
\]
This is a normal subgroup of $G_L$ contained in
$\Delta_0$. Its actual packet action on $P_{D,\ell}$ is
trivial on homology, since $J_D\sim E^2$ and the CM
characters are abelian. By the tangential hyperelliptic
equivariance established in A.24, the first Johnson
degree of this action is zero: the involution acts as
$-1$ on homology and $+1$ on its second Lie degree.
The square of its based lift may be inner; only the
induced graded signs are used here.

Let $F$ be a local Frobenius above $2$ and $\tau$ a tame
$\ell$-inertia element with nonzero monodromy. Then
$\delta=[F,\tau]\in\Gamma$. Its affine outer monodromy
is $(q-1)$ times that of $\tau$, by the tame relation,
where $q>1$ is the residue cardinality. In the semistable
prime-to-two unipotent outer action wild inertia is
trivial. Thus an outer representative for the first
nonzero derivation of $\delta$ is
\[
             m_1T_1+m_2T_2,
                \qquad m_i=(q-1)e_i\ne0,
                                                        \tag{A.27.10}
\]
up to the same nonzero normalization for both coefficients.

Compare this ordinary-base-point representative with
the actual selected packet action. They have the same
\emph{affine} outer action, so their difference is inner.
Their first Johnson degrees are both zero. Consequently
the linear term of an inner conjugator is zero: a linear
free-Lie element commuting with all four generators is
zero. In degree two the difference therefore has the form
$\operatorname{ad}(w_2)$, with $w_2\in L_2$.
The actual packet fixes its actual cusp element on
$\Gamma\subset\Delta$; since its first Johnson degree
vanishes, its degree-two derivation kills the leading
cusp class $\Omega$. The representative (A.27.10) kills
$\Omega$ as well. Hence
\[
                              [w_2,\Omega]=0.
\]
The free-Lie centralizer calculation proved in A.24
implies $w_2=\lambda\Omega$. We have obtained the
\emph{actual based} degree-two derivation (A.27.2), with
both $m_i\ne0$ and an arbitrary $\lambda\in\Q_\ell$.
This comparison requires no tangential extension of the
ordinary-base-point theorem of Betts--Dogra.

\paragraph{A Galois conjugate removes the case $\lambda=0$.}
If $\lambda\ne0$, (A.27.3) already gives the required
graded kernel. Suppose $\lambda=0$. Since $\Gamma$ is
normal in $G_L$, the degree-two image of its actual
action is stable under conjugation by $s_b(G_L)$.
Conjugation on this first nonzero degree depends only
on the action on homology. We exhibit one such
conjugate whose kernel meets that of $D_0$ only in
$\Q_\ell\Omega$.

Extend coefficients to split the CM field. Choose
symplectic local component bases $a,b$ and $c,d$ with
\[
              f_*a=x,\quad f_*b=y,\quad
              f_*c=x,\quad f_*d=2y,
                                                        \tag{A.27.11}
\]
where $x,y$ is a CM eigenbasis on $E$. These
normalizations respect the polarization because the
two component maps have degrees one and two.
Adjointness identifies the original elliptic subspace
with $\langle a+2c,b+d\rangle$, and the complementary
subspace with $\langle a-c,2b-d\rangle$.
The automorphism $\rho(t,z)=(-t,-iz)$ satisfies
$f\rho=[-i]f$. It acts by $[-i]$ on the original
elliptic factor and by $[i]$ on the complementary one,
as seen on $dt/z$ and $t\,dt/z$. It preserves the two
special components because their map degrees differ.
Thus a symplectic CM eigenbasis is
\[
 p_1=a+2c,\quad p_2=(2b-d)/3,\qquad
 q_1=(b+d)/3,\quad q_2=c-a,
                                                        \tag{A.27.12}
\]
where the $p$'s have one CM character $\psi$ and the
$q$'s the other $\bar\psi$. These are rational
coordinates, also when $\ell=3$; no integral division
by three is asserted.

Put $P=[p_1,p_2]$, $Q=[q_1,q_2]$,
$A=[p_1,q_1]$, and $B=[p_2,q_2]$. Direct substitution
gives
\[
 \Omega=A+B,\qquad
 \omega_1=\tfrac13(P+A+2B+2Q),\qquad
 \omega_2=\tfrac13(-P+2A+B-2Q).
                                                        \tag{A.27.13}
\]
Take an actual Frobenius $g\in G_L$ at an ordinary
good prime of $E$ split in its CM field and different
from $\ell$. Its two CM eigenvalues $\alpha,\bar\alpha$
have ratio of infinite order: as algebraic numbers,
their prime-ideal supports above that rational prime
are disjoint. Passing to a positive residue power
does not change this fact. If $t$ is the action of
$g$ on homology, (A.27.13) shows
\[
 U\cap tU=k\Omega,\qquad
                         U=\langle\omega_1,\omega_2\rangle.
                                                        \tag{A.27.14}
\]
Indeed both planes contain $\Omega$ up to its
cyclotomic scalar, whereas the ratio of their $P$ and
$Q$ coefficients is changed by
$(\alpha/\bar\alpha)^2\ne1$.

Inside $V^{\otimes 2r}$, group the tensor factors in
pairs. For subspaces $U,U'\subset V^{\otimes2}$,
\[
                  U^{\otimes r}\cap(U')^{\otimes r}
                            =(U\cap U')^{\otimes r}.
                                                        \tag{A.27.15}
\]
This follows by choosing a vector-space basis adapted
to $U\cap U'$, then to $U$ and $U'$, and finally to
their sum. Hence the associative algebras freely
generated by $U$ and $tU$ intersect in $k[\Omega]$.
Their primitive parts intersect in $k\Omega$.
By (A.27.3), the two actual degree-two derivations
from $\delta$ and $g\delta g^{-1}$ therefore have
common free-Lie kernel exactly $k\Omega$.
Descent from the coefficient extension gives the
same assertion over $\Q_\ell$.

\paragraph{From the graded obstruction to actual fixed elements.}
Let $u\in P_{D,\ell}$ be fixed by $s_b(\Delta_0)$.
The closed lower-central filtration of this free
pro-$\ell$ group is separated and has torsion-free
free-Lie graded pieces over $\Z_\ell$. If $u\ne1$,
its first nonzero rational graded class is killed
by the actual degree-two derivations above. It is
therefore $a\Omega$ in degree two. The integral
leading class of the selected cusp element $c_D$
is primitive: in an integral symplectic basis it
is $[a,b]+[c,d]$, up to orientation. Thus $a\in\Z_\ell$.

The actual cusp power $c_D^a$ is fixed by
$s_b(\Delta_0)$. Multiplying $u$ by its inverse
removes that leading term and leaves a fixed
element in $\gamma_3P_{D,\ell}$. A nontrivial
remainder would have a leading class of degree
at least three in the same common kernel, which
is impossible. Hence $u=c_D^a$. The reverse
inclusion in (A.27.1) follows from the trivial
cyclotomic action on $\Delta_0$.

The three reflection quotients and their chosen
cusp lifts are conjugate after a finite extension,
so the argument applies to each. The proof uses
the place above two and therefore asserts
(A.27.1) only for odd $\ell$.
\end{proof}

\begin{statement}{A.28}{Proved}{The three-tail reduction and the remaining two-boundary algebra}
Let $C:v^2=w^7-w$, with its two cusps $P_0,P_\infty$.
After a finite extension at the place above two, its stable pointed
special fibre has three elliptic components attached to one rational
component carrying both cusps. The three elliptic components map
isomorphically to the good special fibre of $E_{0,-1}$ under the
cyclic cubic covering
\[
                C\longrightarrow E_{0,-1},\qquad
                      (w,v)\longmapsto(w^3,wv).
                                                        \tag{A.28.1}
\]
The cubic rotation permutes the three tails, whose edge lengths are
equal and strictly positive.

For every odd $\ell$ and every original packet, the monodromy of a
local Frobenius--tame commutator has the following alternatives on
the free rational graded Lie algebra
\[
           \mathbb L(kt\oplus V_1\oplus V_2\oplus V_3),
           \qquad \dim V_i=2,\quad k=\mathbb Q_\ell.
\]
Here $t$ is the primitive class of the selected cusp. If its first
Johnson derivation is nonzero, the actual fixed subgroup is precisely
that cusp inertia. Otherwise its second Johnson derivation, together
with one actual CM Frobenius conjugate, has common kernel
\[
                         \mathbb L_k(t,\Omega),
              \qquad \Omega=\omega_1+\omega_2+\omega_3,
                                                        \tag{A.28.2}
\]
where $\omega_i$ is the symplectic form of $V_i$.
Consequently every leading graded class of an actual fixed element
lies in this displayed algebra. The algebra contains higher
commutators, so this statement leaves the corresponding exceptional
full fixed subgroup undecided.
\end{statement}
\begin{proof}
We first determine the geometry, keeping the distinction between the
cyclic cubic quotient and the degree-six map to the original $E$.
On $E'=E_{0,-1}$, the function $x$ has divisor $2T-2O$, with
$T=(0,0)$ of order two. The cubic cover is the Kummer cover $w^3=x$.
Both branch points specialize to the origin after potential good
reduction above two: the special fibre is supersingular, and its
geometric two-torsion consists only of the origin. The stable
two-pointed model of this good elliptic curve therefore consists of
the elliptic component and a rational component bearing $T$ and $O$,
joined by an edge of positive length. One obtains this model by
successively blowing up the coincidence of the two sections, then
contracting unmarked valence-two rational components.

Normalize this marked model in $w^3=x$, after a finite extension
which makes the vertical Kummer valuations divisible by three.
This is a tame operation, since the residue characteristic is two.
The local description by a cube root of a monomial in the parameters
shows that its only possible branch points on each component are the
marked points and the node. At the elliptic component the node is the
only such point. A cyclic cubic cover cannot branch at just one
point: the divisor of a Kummer function has degree zero, so its sole
possibly nonzero order modulo three must be zero. Thus the cover is
unramified at this node.

Over the rational component the two marked orders are $2$ and $-2$
modulo three. Their sum is zero, again giving no ramification at
the node. Its normalization is a connected rational curve, by
Riemann--Hurwitz for a cyclic cubic cover of a projective line with
exactly two totally ramified points. There are three points over its
node. Over the elliptic component there are exactly two possibilities:
a connected etale degree-three isogeny, or three disjoint copies of
that elliptic curve. These statements also verify semistability
directly: the ramified marked-point coordinates have the form
$u=z^3$, and each unramified node has three ordinary nodal preimages.
After the preceding finite extension there are no further positive-
genus components introduced by normalization or regular subdivision.

The first possibility would give one elliptic component and one
rational component joined at three nodes, hence a dual graph of
first Betti number two. It is impossible. Indeed A.4 gives
$J_C\sim E^3$ after a finite extension; its good reduction then
follows from that of $E$ and Neron--Ogg--Shafarevich. The torus rank
of the generalized Jacobian of a nodal fibre is the first Betti
number of its dual graph, so that graph must be a tree. We obtain
the asserted three-tail fibre. The rotation fixes the central
component, permutes its three nodes, and permutes the elliptic
components. It preserves the node thicknesses, proving equality
and positivity of the three edge lengths. Both cusps lie on the
central rational component, at the totally ramified marked points.

Choose an ordinary base point on that central component, after a
finite extension. The affine edge-twist theorem used in A.27 gives
zero first Johnson degree and second-degree derivation
\[
 D_0(t)=0,\qquad D_0|_{V_i}=m_i\operatorname{ad}(\omega_i),
                     \qquad m_i\ne0.                   \tag{A.28.3}
\]
For a Frobenius--tame commutator, its coefficients are the positive
edge lengths multiplied by the common nonzero scalar $q-1$.
The argument needs only their nonvanishing.

Compare this affine representative with the actual packet action.
Their difference is inner. Both actions fix the leading selected
cusp class $t$. The first-degree inner term is therefore
$\lambda\operatorname{ad}(t)$: the centralizer of a free generator
in the free Lie algebra is its own line. If $\lambda\ne0$, its
kernel is $kt$. Primitivity of the actual cusp class, subtraction
of its actual power, and separation of the lower-central filtration
give the full fixed subgroup assertion, just as in A.27.
If $\lambda=0$, the first-degree action vanishes. The second-degree
inner correction is $\operatorname{ad}(z_2)$ and must kill $t$,
so $z_2=0$. Thus (A.28.3) is the actual second-degree derivation,
with no unspecified inner term.

The first-core normal form of A.27 applies to any number of elliptic
blocks and an additional free generator $t$. Treat $t$ and the
$\omega_i$ as the boundary-prefix letters. On a first nonboundary
core of type $i$, the only term increasing its prefix is nonzero
right multiplication by $m_i\omega_i$. Hence the associative
kernel is exactly $k\langle t,\omega_1,\omega_2,\omega_3\rangle$;
its primitive part is the corresponding free Lie algebra.

We compute its intersection with an actual arithmetic conjugate.
Extend coefficients to split the CM field, and choose symplectic
bases $a_i,b_i$ on the tails with $(f_i)_*a_i=x$ and
$(f_i)_*b_i=y$ on $E'$, where $x,y$ is a CM eigenbasis. This is
possible because each component map has degree one. Put
$A=\sum a_i$ and $B=\sum b_i$.
The original elliptic subspace is $\langle A,B\rangle$ and its
orthogonal complement is the sum-zero subspace in the two triples.
The automorphism $(w,v)\mapsto(-w,iv)$ acts as $[-i]$ on the
original elliptic factor and as $[i]$ on the complementary two
factors: on the three holomorphic differentials its eigenvalues
are $i,-i,i$. It preserves each special elliptic tail, since it
has order four and commutes with their three-cycle. On every tail
its action is identified by (A.28.1) with $[-i]$ on $E'$.
It follows that the global CM eigenspaces are
\[
          \langle A\rangle\oplus\{\textstyle\sum u_i b_i:
                        \sum u_i=0\},\qquad
          \langle B\rangle\oplus\{\textstyle\sum u_i a_i:
                        \sum u_i=0\}.                  \tag{A.28.4}
\]

An actual good ordinary Frobenius with the two CM eigenvalues
$s,r$ acts, in these coordinates, by
\[
 a_i\longmapsto r a_i+(s-r)A/3,\qquad
 b_i\longmapsto s b_i+(r-s)B/3.                         \tag{A.28.5}
\]
Take its place away from $\ell$ and pass to any required positive
residue power. Its eigenvalue ratio is not a root of unity and has
complex absolute value one. Hence $s-r\ne0$ and $s+2r\ne0$.
Put $P=\frac13(1)_{3\times3}$ and
$M=rI+(s-r)P$. The other matrix is $srM^{-1}$.
The plane $U=\langle\omega_1,\omega_2,\omega_3\rangle$
identifies with the diagonal matrices in the $a_i\wedge b_j$
coordinates. The conjugate plane is $srMUM^{-1}$.
Every entry of $M$ is nonzero. If a diagonal matrix $D$ satisfies
$DM=MD'$ for another diagonal matrix $D'$, the $(i,j)$ entry gives
$D_{ii}=D'_{jj}$ for every $i,j$. Both are therefore scalar. Thus
\[
                             U\cap gU=k\Omega.         \tag{A.28.6}
\]
The exact symbolic matrix calculation is also executed in
\texttt{a\_C3\_stable\_graph.py}.

The boundary generator $t$ is carried to a nonzero scalar multiple
of itself. Decompose tensor words according to the occurrences of
$t$ and then group each intervening even-length $V$ block into
pairs. The tensor-intersection argument of A.27 on each block
shows that the two associative kernel algebras intersect in
$k\langle t,\Omega\rangle$. Their primitive intersection is
(A.28.2). The actual commutator group is normal in a sufficiently
small arithmetic group, so the Frobenius conjugate used here is
another actual derivation fixing every element under consideration.
Finally descend the kernel statement to $\mathbb Q_\ell$.

Unlike the one-cusp case, the free Lie algebra in (A.28.2) is not
one-dimensional. In particular $[t,\Omega]$ survives this argument.
The three genus-two quotient maps do not remove it at leading
weight, because each sends the two generators to proportional
symplectic forms. A further actual cusp-path or higher-monodromy
calculation is required; no existence of a fixed element with
this leading term is being asserted.
\end{proof}

\begin{statement}{A.29}{Proved}{At most one odd-prime stable genus-three packet remains}
Let $\ell$ be odd and let
$C:v^2=w^7-w$ with its two cusps removed. Use an induced
original packet action and a permitted open subgroup
$\Delta_0$ as in A.24 and A.28. If its ordinary first
Johnson homomorphism is nonzero, then
\[
                    P_{C,\ell}^{\,s_b(\Delta_0)}=I_{C,\ell}.
                                                        \tag{A.29.1}
\]
No invertibility hypothesis on any first-Johnson matrix
is needed.

For each odd $\ell$, there is at most one original
pro-$\ell$ packet class, up to conjugation by original
cusp inertia, for which the induced first Johnson
homomorphism vanishes after finite descent. Every other
original packet class satisfies (A.29.1) for every
permitted open subgroup. In particular, along every
nontrivial Kummer direction in A.24, at most one
parameter can remain undecided on $C$.

This does not assert that the exceptional class exists,
identify its tangent parameter, or determine its full
fixed subgroup. Vanishing here means vanishing on an
open subgroup of $\Delta$, rather than just on the
possibly smaller group $\overline{[G_L,G_L]}$.
\end{statement}
\begin{proof}
\paragraph{A nonzero matrix is enough after the monodromy reduction.}
Write the ordinary graded Lie algebra as
$\mathbb L(kt\oplus V)$, where $k=\mathbb Q_\ell$,
$t$ is the selected cusp class, and $V$ is the
six-dimensional negative hyperelliptic eigenspace.
For every $\delta\in\Delta_0$, A.24 gives
\[
 D_\delta(t)=0,\qquad
 D_\delta(v)=[t,A_\delta v],\qquad
                  A_\delta J=JA_\delta^T,
                                                        \tag{A.29.2}
\]
where $\Omega=\sum J_{ij}v_iv_j$ in the tensor algebra
and $J$ is nonsingular and skew. Define the quadratic
Lie element
\[
                     \Omega_{A_\delta}
                              =\sum(A_\delta J)_{ij}v_iv_j.
\]
It is a Lie element because $A_\delta J$ is skew, and
it is nonzero precisely when $A_\delta\ne0$.
The same tensor expansion that gives the peripheral
constraint in A.24 gives the exact identity
\[
                   D_\delta(\Omega)
                             =[t,\Omega_{A_\delta}].
                                                        \tag{A.29.3}
\]
The terms with a $t$ between two $V$ letters cancel
by (A.29.2).

We claim that, whenever $A_\delta\ne0$,
\[
 \ker D_\delta\ \cap\ \mathbb L_k(t,\Omega)=kt.
                                                        \tag{A.29.4}
\]
If $\Omega_{A_\delta}=a\Omega$ for $a\ne0$, the
restriction is $a\operatorname{ad}(t)$, so this is
the free-generator centralizer calculation. Otherwise
the two quadratic tensors $\Omega,\Omega_{A_\delta}$
are linearly independent. Their tensor algebra
embeds into $T(V)$ by grouping factors in pairs,
and adjoining the independent letter $t$ still
gives a free associative algebra. There is therefore
a homomorphism on this algebra that fixes $t,\Omega$
and sends $\Omega_{A_\delta}$ to $\Omega$.
For every $u\in\mathbb L(t,\Omega)$ it sends
$D_\delta(u)$ to $[t,u]$, by (A.29.3) and the
derivation rule. Thus $D_\delta(u)=0$ implies
$[t,u]=0$, proving (A.29.4).

Let $u$ be an actual fixed pro-$\ell$ element.
If the first local monodromy alternative in A.28
occurs, that statement already proves (A.29.1).
Otherwise its two actual second-degree derivations
put the first nonzero class of $u$ in
$\mathbb L(t,\Omega)$. These derivations come from
the normal commutator group used there, contained
in $\Delta_0$. We may also apply the first
derivation of any $\delta\in\Delta_0$ with
$A_\delta\ne0$. By (A.29.4), the leading class
is a multiple of $t$. Its coefficient lies in
$\mathbb Z_\ell$ because $t$ is primitive in
integral homology. Subtract the corresponding
actual fixed cusp power. A nontrivial remainder
would have a leading class of degree at least
two in the same common kernel $kt$, which is
impossible. This proves (A.29.1).

\paragraph{The exceptional original packet class is unique if it exists.}
The first Johnson map is an additive continuous
homomorphism on any permitted $\Delta_0$, since
this group acts trivially on integral homology.
Its target is a torsion-free finite-dimensional
$\mathbb Q_\ell$-space. If it vanishes on an open
subgroup, it vanishes on the whole group:
each element has a positive power in that
subgroup. Hence the zero-Johnson condition is
independent of further finite descent.

Suppose two original packet pro-$\ell$ classes
have this condition. Choose a common finite
descent field and a common open subgroup
$\Delta_0$ for the two induced actions.
Let $\kappa:G_K\to\mathbb Z_\ell(1)$ be their
original cusp Kummer difference. The actual
cusp ramification index for the stable $S_3$
cover is three. Exactly as in A.24, the two
first-Johnson matrices on the common subgroup
differ by
\[
                        \frac{\kappa(\delta)}3\,1_V.
                                                        \tag{A.29.5}
\]
At $\ell=3$ the common finite descent ensures
that the numerator is divisible by three in
the integral cusp action. At other primes
three is a unit. In all cases zero on both
sides forces $\kappa|_{\Delta_0}=0$.

The restriction of $\kappa$ to $\Delta$ is an
additive character into a torsion-free group.
Vanishing on its open subgroup $\Delta_0$
therefore implies $\kappa|_\Delta=0$.
It descends to a cyclotomic cocycle on
$G_K/\Delta$. The odd-prime calculation in
A.6 gives
\[
              H^1_{\mathrm{cts}}
                 (G_K/\Delta,\mathbb Z_\ell(1))=0.
                                                        \tag{A.29.6}
\]
Thus $\kappa$ is a cusp coboundary and the
two original pro-$\ell$ packet classes coincide.
This proves the asserted uniqueness, without
assuming any value for the actual matrices.

Finally a nontrivial Kummer direction used
in A.24 has nonzero restriction to every
permitted open subgroup of $\Delta$, by
the local Frobenius--tame commutator there.
Two distinct parameters on that ray cannot
both have zero first Johnson map, by
(A.29.5). This gives the one-parameter
conclusion as well.
\end{proof}

\begin{statement}{A.30}{Proved}{The explicit Collino unit removes the stable genus-three exception}
Let $\ell$ be odd. For every original packet whose action is
induced on the stable $S_3$ cover
\[
                  C^\circ:\quad v^2=w^7-w,\qquad w\ne0,\infty,
\]
and every permitted open subgroup $\Delta_0$, one has
\[
                    P_{C,\ell}^{\,s_b(\Delta_0)}=I_{C,\ell}.
                                                        \tag{A.30.1}
\]
In particular the possible zero-first-Johnson packet in A.29
does not exist. This conclusion includes $\ell=3$.

The arithmetic input is an explicit higher-Chow pushforward.
Put $L=\mathbb Q(i,\sqrt3)$, $\zeta=(-1+i\sqrt3)/2$,
$\epsilon=2+\sqrt3$, and $F:Y^2=X^3+3X$. There is a
degree-four map $q:C\to F$ for which the Collino cycle of
the two removed Weierstrass points pushes forward to
\[
                  \zeta^2\epsilon^{-4}\in
                    \operatorname{CH}^1(F,1)=L^\times.
                                                        \tag{A.30.2}
\]
After subtracting twice the degree-two reflection
pushforward, this is a functional on the primitive
exterior-square representation. Its rational
$\ell$-adic regulator is $-4\kappa_\epsilon$, which
remains nonzero on every open subgroup of $\Delta$.
The comparison with the actual Johnson class is proved
below on universal moduli and then specialized; it does
not assume faithfulness of a Hodge realization.
\end{statement}
\begin{proof}
\paragraph{1. A canonical universal extension for the twice-punctured curve.}
We first work on the moduli stack $\mathcal H$ of smooth
hyperelliptic curves of genus $g\ge2$ with all
Weierstrass points labelled. Choose three of its
sections, denoted by $P,Q,R$. Write $\mathcal C^\circ
=\mathcal C\setminus(Q\cup R)$, and use the ordinary
section $P$ as base point. The quotient by the
hyperelliptic involution is a projective-line bundle
with three disjoint sections, so it has the unique
coordinate $f$ with values $0,\infty,1$ at $Q,R,P$.
On each fibre,
\[
                         \operatorname{div}(f)=2Q-2R.
\]
All these constructions are defined over $\mathbb Q$.
They may equally be performed on a finite level cover
and descended; the stack notation keeps the geometric
mapping class group unchanged.

Let $\mathcal V=H_1(\mathcal C/\mathcal H,\mathbb Q_\ell)$,
let $t$ be the positive local class about $Q$, and
let $\mathfrak p$ be the relative rational unipotent
Lie algebra of $\mathcal C^\circ$ based at $P$.
Use the ordinary lower central series $\Gamma_\bullet$.
The hyperelliptic involution $h$ fixes $P$ and is an
actual involution commuting with all monodromy.
Its eigenspaces on $H_1(\mathcal C^\circ)$ are
$\mathcal V$ with eigenvalue $-1$ and
$\mathbb Q_\ell(1)t$ with eigenvalue $+1$.
The affine fundamental group is free, so the negative
eigenspace in its class-two quotient gives the
canonical exact sequence of local systems
\[
 0\longrightarrow \mathbb Q_\ell(1)\otimes\mathcal V
 \ \xrightarrow{\ t\otimes v\mapsto[t,v]\ }\
 (\mathfrak p/\Gamma_3\mathfrak p)^-
 \longrightarrow\mathcal V\longrightarrow0.
                                                        \tag{A.30.3}
\]
In particular no choice of a global Johnson splitting
is being assumed.

Its extension class has coefficients
$\operatorname{End}(\mathcal V)(1)$. Polarization
identifies this with $\mathcal V\otimes\mathcal V$.
Apply the alternating projection and then divide by
the polarization line to obtain a canonical class
\[
 a_\ell\in H^1_{\mathrm{cts}}
       (\mathcal H,\mathcal W_\ell),\qquad
 \mathcal W_\ell=
       \Lambda^2\mathcal V/\mathbb Q_\ell\Omega .
                                                        \tag{A.30.4}
\]
Here $\Omega$ denotes the polarization tensor, so the
displayed line is intrinsically a copy of
$\mathbb Q_\ell(1)$. We normalize the alternating
identification so that a self-adjoint symplectic
projection $A$ is sent to $\Omega_A=AJ$, in the tensor
notation of A.29.

On any subgroup acting trivially on homology, the
restriction of (A.30.3) is represented by the literal
first-Johnson homomorphism
\[
                  v\longmapsto[t,A_\gamma v].
                                                        \tag{A.30.5}
\]
Thus, when the peripheral relation makes $AJ$ skew,
the restriction of $a_\ell$ is exactly the class of
$\Omega_{A_\gamma}$ modulo $\Omega$.

\paragraph{2. The universal regulator and its precise comparison target.}
Let $\mathcal J$ be the relative Jacobian. The
Collino cycle is represented by
\[
 Z_{Q,R}=(\iota_Q(\mathcal C),f)
                  +(\iota^R(\mathcal C),f),\qquad
 \iota_Q(x)=[x-Q],\quad \iota^R(x)=[R-x].
                                                        \tag{A.30.6}
\]
The divisors cancel: they are $2[0]-2[R-Q]$ and its
negative. This defines a relative, and hence a total,
higher-Chow cycle of codimension $g$ and degree one.
The normalization $f(P)=1$ fixes its decomposable
ambiguity. The alternative second embedding $x-R$
gives the same cycle here, because the hyperelliptic
involution sends $x-R$ to $R-x$ and fixes $f$.
Rescaling $f$ by a function on the base adds a
decomposable class with coefficient twice the
polarization line, so it does not change the
primitive class below. Thus this normalization
also agrees with the primitive normal function
used in the quoted monodromy calculation.

Take the rational continuous \'etale cycle class in
$H^{2g-1}(\mathcal J,\mathbb Q_\ell(g))$. There is a
canonical way to extract its Leray degree-one part:
fix an integer $n>1$ and take the $n^{2g-2}$
eigenspace of $[n]^*$. On
$R^j\pi_*\mathbb Q_\ell(g)$ the eigenvalue is $n^j$,
so in total degree $2g-1$ this eigenspace identifies
canonically with
\[
 H^1\bigl(\mathcal H,
          R^{2g-2}\pi_*\mathbb Q_\ell(g)\bigr)
       =H^1(\mathcal H,\Lambda^2\mathcal V).
\]
Distinct eigenvalues split the finite Leray
filtration and kill its possible differentials.
Thus neither a noncanonical edge map nor a
choice of a lift through that filtration is involved.
Projecting away from $\Omega$ gives
\[
                    r_\ell(Z_{Q,R})
                      \in H^1(\mathcal H,\mathcal W_\ell).
                                                        \tag{A.30.7}
\]

We use the usual comparison compatibility of higher
cycle classes, rather than an injectivity conjecture
for realizations. In degree one this compatibility
can also be read directly from the construction:
a symbol $(Y,u)$ is the Gysin image of the Kummer
class of $u$ on the complement of its divisor;
the boundary is the divisor of $u$. The cancellation
in (A.30.6) closes the corresponding total cycle.
Over $\mathbb C$, the comparison sends the Kummer
class to the winding-number class of $u$, and
commutes with the Gysin and localization maps.
It therefore sends the \'etale class of (A.30.6)
to its Betti cycle class. The latter is the
underlying topological class of the regulator
normal function; its Leray degree-one class is
the monodromy of that normal function.
The $[n]^*$ projector commutes with all these
comparisons and with specialization.

For a published formulation of the same comparison
by universal Chern classes, see
\cite[Section 2, especially Lemma 2.2]{fghAsakuraSaito}.
That lemma is displayed there in codimension two;
its proof uses the universal Chern-class compatibility
in every degree. The degree-one symbol construction
just given supplies the required codimension-$g$
case. The passage from a universal higher-Chow
class to arithmetic by geometric comparison and
vanishing primitive invariants is also carried out
in \cite[Claim 3.3 and the following comparison on
pp.~171--172]{fghAsakuraSaito}. The analogous
fundamental-group comparison for the Ceresa class
is \cite[Corollary 6.7 and Theorem 8.1]{fghHainMatsumoto}.

\paragraph{3. The geometric comparison is a direct Dehn-twist calculation.}
Embed the base field in $\mathbb C$. The geometric
mapping class group of $\mathcal H$ is the
hyperelliptic group with all Weierstrass points
fixed; write $T\mathcal H$ for its Torelli subgroup.
Both the Betti version of (A.30.4) and the
monodromy class of (A.30.7) restrict to ordinary
homomorphisms on $T\mathcal H$.

Let $d$ be a symmetric separating simple closed
curve, with positive Dehn twist $T_d$. If $d$
does not separate $Q$ from $R$, its class in
$H_1(C\setminus\{Q,R\})$ is zero. The twist
conjugates the generators on one side by a
loop with that homology class and fixes those
on the other side, so its ordinary first
Johnson map is zero.

Suppose instead that $d$ separates $Q$ from $R$.
Write $V_Q,V_R$ for the proper homology of the
two sides and $\Omega_Q,\Omega_R$ for their
symplectic tensors. The first homology class
of the boundary of either side is $t$ or $-t$.
The same based twist formula now gives
$A_{T_d}=\pm\operatorname{pr}_{V_Q}$ or
$\mp\operatorname{pr}_{V_R}$, according to
which side contains $P$ and the fixed orientation
convention. These matrices differ, after using
one fixed convention, by a scalar identity.
Consequently their primitive tensor is one
fixed sign times $\Omega_R\bmod\Omega$.
This also verifies independence from the
position of the ordinary base point.

The actual Collino-regulator monodromy is
$4\Omega_R\bmod\Omega$ in this case, and zero
in the preceding case, by
\cite[Theorem 8.7]{fghLuoWatanabe}.
This is a statement about the regulator
normal function itself. It does not use an
identification of our open class-two extension
with a compact-curve Johnson projection.
Theorem A of
\cite{fghBrendleMargalitPutman}
says that symmetric separating twists generate
the hyperelliptic Torelli group. Hence the two
geometric classes restrict to homomorphisms
satisfying
\[
                    r_{\mathrm B}|_{T\mathcal H}
                         =\eta\,4a_{\mathrm B}|_{T\mathcal H},
                  \qquad \eta\in\{1,-1\},
                                                        \tag{A.30.8}
\]
where $\eta$ merely records a common choice of
the twist and boundary orientations.

This equality holds in $H^1$ of the full geometric
group as well. Indeed its restriction to Torelli
is injective for the coefficient representation
$\Lambda^2V/\mathbb Q\Omega$:
\cite[Theorem 5.1 and Section 5.3]{fghLuoWatanabe}
identifies this $H^1$ with the equivariant
homomorphisms out of the abelianization of
the relative unipotent radical, and proves
that the Torelli image is Zariski dense in
that radical. An equivariant homomorphism
zero on Torelli is therefore zero.
After tensoring with $\mathbb Q_\ell$, the
Betti--\'etale comparison above proves
(A.30.8) for the geometric \'etale classes.
More explicitly, the topological geometric
fundamental group has dense image in its
\'etale profinite completion. Restriction
of continuous $H^1$ to this dense subgroup
is injective: if a continuous cocycle is
the coboundary of a vector on the dense
subgroup, subtract that coboundary and
use continuity to make it zero everywhere.
Thus no separate comparison assertion for
arbitrary $\mathbb Q_\ell$ local systems
is needed at this step.

Finally the geometric homology monodromy contains
$\operatorname{Sp}_{2g}(\mathbb Z)[2]$
and is Zariski dense in $\operatorname{Sp}_{2g}$;
see \cite[Theorem 3.6 and Section 4.1]{fghLuoWatanabe}.
The only invariant in $\Lambda^2V$ is the
polarization line, so
$H^0(\mathcal H_{\overline{\mathbb Q}},
\mathcal W_\ell)=0$. The first
Hochschild--Serre terms now make
\[
 H^1(\mathcal H,\mathcal W_\ell)
       \longrightarrow
 H^1(\mathcal H_{\overline{\mathbb Q}},\mathcal W_\ell)
                                                        \tag{A.30.9}
\]
injective. We obtain the unconditional arithmetic
identity
\[
                       r_\ell(Z_{Q,R})=\eta\,4a_\ell.
                                                        \tag{A.30.10}
\]
No assertion that two arbitrary Hodge extensions
have equal \'etale realizations has been used.
In particular the corrected augmentation-space
choices in Colombo's comparison, discussed in
\cite[Introduction and Theorem 4.14]{fghSarkarSreekantan},
are not premises of this proof.

\paragraph{4. Specialization and the tangential base point.}
Specialize to $C:v^2=w^7-w$, $Q=P_0$, $R=P_\infty$,
and $P=(1,0)$. Then the normalized function is
exactly $f=w$. Over $L$ all Weierstrass points
are rational, so this is a point of $\mathcal H$.

For completeness, the primitive extension class
is unchanged when it is compared with the
tangential packet action. Rationally the
hyperelliptic involution can be made into an
actual commuting involution at a tangential
point at $Q$: its derivative is $-1$, and the
Kummer class of $-1$ vanishes in
$H^1(G_L,\mathbb Q_\ell(1))$. The resulting
invariant rational path from a tangent to its
negative has square zero as a loop, because
$\mathbb Q_\ell(1)^{G_L}=0$. Thus the corrected
involution has order two. Twisting a packet
by cusp inertia commutes with this involution.

Choose an involution-equivariant rational
unipotent path from this tangential point to
$P$. Such a path exists by successive
averaging in the central quotients of its
unipotent path torsor: the order-two group
has vanishing positive cohomology on
$\mathbb Q_\ell$-vector spaces. The two
Galois actions then differ by an inner
cocycle invariant under the involution.
Its linear part lies in the positive
eigenspace $\mathbb Q_\ell t$. Its effect
on (A.30.5) is therefore a scalar
$\operatorname{ad}(t)$, which disappears
in $\mathcal W_\ell$.
The same argument applies to every induced
original packet after the permitted common
finite descent. At $\ell=3$ the integral
packet change still has the divisibility
described in A.24; the present comparison
is rational and introduces no integral
division assertion.

It follows that (A.30.10), restricted to
any permitted subgroup of $\Delta$, is
the primitive tensor of the actual matrices
$A_\delta$ in A.29, multiplied by the
fixed nonzero rational factor $\eta4$.

\paragraph{5. A degree-four elliptic projection and its exact norm.}
Put
\[
 x_0=w-w^{-1},\quad
 x_1=\zeta w-\zeta^2w^{-1},\quad
 x_2=\zeta^2w-\zeta w^{-1},\qquad
 u=v/w^2 .
\]
Then $x_0+x_1+x_2=0$ and
$u^2=x_j^3+3x_j$ for $j=0,1,2$.
Hence $q_0=(x_0,u)$ and $q_1=(x_1,u)$
are degree-two maps to $F$.
With $[i](X,Y)=(-X,iY)$, define
$q=q_0+[i]q_1$. The addition formula gives
\[
 \begin{split}
 X_q&=\frac{-2i(w^3-w^{-3})}{x_2^2}-x_0+x_1,\\
 Y_q&=u\left(\frac{(1-i)(x_0-X_q)}{x_0+x_1}-1\right).
 \end{split}                                             \tag{A.30.11}
\]
After cancellation, $X_q=N(w)/D(w)$ with
$\deg N=4$, $\deg D=3$, $D(0)=0$, and
$D$ a nonzero constant times
$w(w^2-\zeta^2)$. The numerator and denominator
are relatively prime. Thus $q$ has degree
four and
\[
                       q^{-1}(O)
                       =\{P_0,P_\infty,P_\zeta,P_{-\zeta}\}
                                                        \tag{A.30.12}
\]
with ramification index one at all four points.
The poles of $X_q$ have order two on the
hyperelliptic curve, including at the two
finite Weierstrass points $\pm\zeta$.

For a generic target point $(X,Y)$, $Y_q=vR(w)$
recovers $v$ uniquely from $w$. The four conjugates
of $w$ are therefore exactly the roots of
$N(w)-XD(w)$. Their product is
$N(0)/\operatorname{lc}(N)$, independent of $X$.
Equivalently, the coefficients of $w$ at infinity
and $w^{-1}$ at zero in $X_q$ are
\[
          a_\infty=\zeta-1-2i\zeta^2,\qquad
          a_0=1-\zeta^2+2i\zeta .
\]
Since $D$ is proportional to $w(w^2-\zeta^2)$,
\[
 \operatorname{Norm}_{L(C)/L(F)}(w)
       =-\zeta^2\,\frac{a_0}{a_\infty}
       =\zeta(2+\sqrt3)^{-2}.
                                                        \tag{A.30.13}
\]
These are exact algebraic identities.
The retained program \texttt{code/a\_C3\_collino\_unit.py}
verifies (A.30.11), the target equation, cancellation,
degrees, all four poles, and (A.30.13) over
$\mathbb Q(i,\sqrt3)(w)$.

Both $Q$ and $R$ map to $O$ under $q$.
The two symbols in (A.30.6) consequently push
forward to the same constant norm, so their
product is $\zeta^2\epsilon^{-4}$, as claimed
in (A.30.2). For comparison, the equation
$w^2-Xw-1=0$ for $q_0$ gives
$\operatorname{Norm}_{q_0}(w)=-1$ and
$(q_0)_*Z_{Q,R}=1$.

\paragraph{6. Removing the scalar part and proving nonvanishing on $\Delta$.}
A map $C\to F$ sending $Q$ to $O$ induces a
homomorphism $J(C)\to F$, and therefore a
functional
\[
       \Lambda^2H_1(C,\mathbb Q_\ell)
                  \longrightarrow\mathbb Q_\ell(1).
\]
It sends $\Omega_C$ to
$\deg(q)\Omega_F$. This follows either from
$q_*q^*=\deg(q)$ on $H^1(F)$ and the
intersection pairing, or from the adjoint
description of the induced Jacobian
homomorphism. Consequently
\[
                        \Phi=q_*-2(q_0)_*
                                                        \tag{A.30.14}
\]
kills $\Omega_C$, because $4-2\cdot2=0$.
It is an algebraic functional on the
primitive quotient, defined over $L$.
Functoriality of higher cycle classes and
(A.30.13) give
\[
             \Phi\bigl(r_\ell(Z_{Q,R})\bigr)
                    =-4\kappa_\epsilon
                    \quad\hbox{in }H^1(G_L,\mathbb Q_\ell(1)).
                                                        \tag{A.30.15}
\]
The root of unity $\zeta^2$ has zero rational
Kummer class. This calculation is why an
off-diagonal elliptic map is needed: the
individual reflection projection has zero
regulator, while the combination (A.30.14)
also removes the scalar polarization
component.

The unit $\epsilon=2+\sqrt3>1$ has infinite
order, so $\kappa_\epsilon\ne0$ over $L$:
the unit group modulo torsion is a finitely
generated free abelian group, whose
$\ell$-adic Kummer map is injective after
tensoring with $\mathbb Q_\ell$.
Moreover $L\subset K^{\mathrm{ab}}$ for
$K=\mathbb Q(i)$. Thus $\Delta\subset G_L$
and $G_L/\Delta$ is abelian. Some element
of this quotient has cyclotomic character
different from one. Sah's cocycle identity
therefore gives
\[
                   H^1(G_L/\Delta,\mathbb Q_\ell(1))=0,
\]
since multiplication by that nonzero
$\ell$-adic scalar is invertible.
Inflation--restriction makes
\[
 H^1(G_L,\mathbb Q_\ell(1))
           \longrightarrow H^1(\Delta,\mathbb Q_\ell(1))
\]
injective. The restriction of $\kappa_\epsilon$
is hence nonzero. On $\Delta$ the coefficients
are trivial, so it is an additive character;
such a character cannot vanish on an open
subgroup unless it vanishes everywhere.
This proves nonvanishing on every permitted
$\Delta_0$.

By (A.30.10), (A.30.15), and the base-point
comparison, the actual primitive first
Johnson homomorphism is nonzero on every
such subgroup, for every packet. In particular
some actual $A_\delta$ is nonzero.
Statement A.29 now applies and gives
(A.30.1). Its integral cusp-power subtraction
is already part of that statement, so no
fixed element of a rational completion is
silently lifted to the pro-$\ell$ group.
\end{proof}


\begin{statement}{A.31}{Proved}{The stable \(S_3\) fixed subgroups at the prime two}
Let
\[
 C:\ v^2=w^7-w,\qquad
 D:\ z^2=(t^2+4)(t^3+3t),
\]
with the two indicated cusps removed from \(C\), and the unique
cusp removed from \(D\). For every original packet and every
permitted open subgroup \(\Delta_0\) on which the selected lift
is defined,
\[
 P_{C,2}^{\,s_b(\Delta_0)}=I_{C,2},
 \qquad
 P_{D,2}^{\,s_b(\Delta_0)}=I_{D,2}.                 \tag{A.31.1}
\]
The second assertion applies to each reflection quotient.
Together with A.27 and A.30, these statements give the full
fixed-subgroup equality for the stable \(S_3\) genus-three cover
and all its genus-two reflection quotients at every prime.

The replacement for prime-to-two tame monodromy is Fontaine's
monodromy on the semistable realization. It is used on a
crystalline realization of a global orbit span, never as an
operator on an individual etale vector.
\end{statement}
\begin{proof}
\paragraph{The semistable comparison used at a cusp.}
We first specify the comparison input, including its tangential
form. Let \(F/\mathbb Q_p\) be finite. Suppose that a marked
curve has a regular semistable model over \(\mathcal O_F\),
with disjoint marked sections in its smooth locus. Give the
model the log structure of the special fibre and the marked
sections. At an integral nonzero tangent to a marked section,
every finite-dimensional quotient of the rational unipotent
fundamental Lie algebra is semistable. Its \(D_{\rm st}\)
realization is the corresponding logarithmic crystalline
Lie algebra, compatibly with brackets, the lower-central
filtration, and the cusp map \(\mathbb Q_p(1)\to\mathfrak u\).
Its monodromy is the logarithmic Gauss--Manin residue.

Here is the passage from the ordinary to the tangential fibre
in this assertion. The ramified semistable comparison is that
of \cite[Corollary 5.1 and Theorem 5.2]{fghAIK}: it is an
association of the universal unipotent local systems with
logarithmic crystalline objects on the whole marked model,
before taking the fibre. The open-curve scope is explicit in
\cite[Remark 1.10 and Section 2]{fghAIK}.
Here is the literal local period map in this pullback.
Write \(W=W(k)\), choose a uniformizer \(\varpi\) of \(F\),
and set
\[
 \mathcal O=W[[Z]],\qquad Z\longmapsto\varpi,\qquad
 \mathcal O_{\rm cris}
   =\widehat{\operatorname{DP}}_{\,(p,P_\varpi(Z))}
                                      W[[Z]],
\]
as in Section 2.1 of \cite{fghAIK}.
On the lifted smooth chart use \(W[[Z,z]]\), with
horizontal marked divisor \(z=0\), and choose the
tangent with normal coordinate \(1\).
Let \(\overline R\) be the integral Kummer tower of
this chart used in the definition of its period sheaf.
A tangential specialization of this tower gives a
map \(\overline R\to\mathcal O_{\mathbb C_p}\).
It sends the underlying elements \(z^{1/p^n}\) to
zero, but sends their logarithmic generators to
compatible formal generators \(e_h/p^n\), whose
structure map is zero. It consequently induces
both a map of tilt rings and a map of their log
monoids. This is the specialization construction
in \cite[Section 9.25, especially (9.25.4)--(9.25.7)]
{fghOlssonTangential}.

Combine its Witt-vector map with \(z\mapsto0\)
and \(Z\mapsto Z\). By the universal property of
the log divided-power envelope, the result is
a morphism from the AIK local log-period
envelope to
\[
       \bigl(A_{\log}(F),\,
                         M_{\rm vertical}\oplus\mathbb N e_h\bigr).
                                                        \tag{A.31.12}
\]
The horizontal copy from the lifted chart and the
horizontal copy from the tilt both map to \(e_h\).
Thus their \emph{log exactification unit} \(u_h\)
maps to \(1\), and every positive divided power
of \(u_h-1\) maps to zero. This is a map of log
envelopes; no numerical quotient \(0/0\) is taken.
In contrast, the vertical exactification unit
\([\varpi^\flat]/Z\) is retained. The underlying
target ring is exactly the vertical \(A_{\log}(F)\)
of AIK, with its additional horizontal log
generator of structure value zero. The
Eisenstein relation \(P_\varpi(Z)\) and its
divided powers remain in the vertical base.
This explains explicitly why ramification of
\(F\) causes no new horizontal relation.

Frobenius sends both log generators to their
\(p\)-fold multiples and sends every
exactification unit to its \(p\)-th power.
It therefore commutes with (A.31.12).
The tangential splitting makes the formal
\(e_h/p^n\) fixed by the arithmetic \(G_F\)
section, so the map is \(G_F\)-equivariant.
It also commutes with the vertical logarithmic
connection, since \(u_h\) has vertical derivative
zero and the vertical unit is unchanged.
The divided-power filtration is respected.
Inverting the Tate period gives the same map
with target \(B_{\log}(F)\).
Pullback of a crystal along this morphism of
logarithmic enlargements is its evaluation
on the target; hence the AIK coefficient
association pulls back to an association on
the horizontal log point. This is the
base-change identity (9.26.2)--(9.26.4) of
\cite{fghOlssonTangential}, with the vertical
log base retained.

Relative to that vertical base, a unipotent
object at this log point is a finite module
with a nilpotent horizontal residue \(R\).
Nilpotence follows from its geometric
unipotent filtration: \(R\) is zero on each
constant graded quotient. The tangent
retraction evaluates this object on the
chosen \(e_h\), giving its usual tangential
fibre. Equivalently, with
\(\ell_h=\log u_h\), the logarithmic Taylor
identification is the finite exponential
\(\exp(-\ell_hR)\), followed by
\(\ell_h=0\). Integrability says that \(R\)
is horizontal for the vertical connection;
thus this identification commutes with
vertical monodromy. Frobenius compatibility
follows from \(\varphi(\ell_h)=p\ell_h\)
and the corresponding residue relation
obtained from \(d\log(\varphi z)=p\,d\log z\).
On the etale side horizontal inertia is
\(\mathbb Q_p(1)\); its comparison is exactly
the Tate-period calculation of
\cite[Proposition 9.5, Corollary 9.16, and
Section 9.24]{fghOlssonTangential}.
This verifies the arithmetic fibre
association and all needed compatibilities.

The intermediate coefficient ring here is
\(B_{\log}\), not a quotient identified
directly with \(B_{\rm st}\).
The fibre is \(B_{\log}\)-admissible.
\cite[Proposition 2.1]{fghAIK}, namely
Breuil's equivalence, converts this into
\(B_{\rm st}\)-admissibility and identifies
the monodromy with the vertical residue.
This construction uses Olsson's local
log-point calculation and AIK's ramified
coefficient comparison. It does not apply
Olsson's global unramified good-reduction
theorem to a ramified semistable model.

One can take the AIK universal objects based at any ordinary
point through the smooth locus, evaluate them at this log
point, and reconstruct its unipotent group by their tensor
matrix coefficients. These finite Lie quotients are subquotients
of finite sums and tensor products of the associated fibres.
Semistable representations are closed under these operations,
so all such Lie quotients are semistable. The same construction gives the
path torsors and the compatibility with the cusp.
Finally \cite[Lemma 6.2(ii) and the deformation argument in
Section 6.2]{fghAIK} identifies the vertical residue with
Fontaine's \(N\). Deformation to a complex marked semistable
curve identifies it with the logarithm of the usual edge
twists. If a node has equation \(xy=\varpi^e\), its
coefficient is \(e\), up to the single common nonzero
normalization of \(N\). Rational subdivisions add these
coefficients along an arm. This statement is about the
residue operator, and hence gives the edge-twist formulas
in every finite unipotent quotient.

\paragraph{Every cusp packet remains semistable.}
The preceding comparison applies first to an integral
algebraic tangent. It also applies to every packet used
here, after rationalization. Locally the difference from
such a tangent is
\[
                  \kappa\in H^1(G_F,\mathbb Q_p(1)).
\]
Every such class is semistable. Indeed Kummer theory identifies
this finite-dimensional vector space with the rationalized
pro-\(p\) completion of \(F^\times\). A uniformizer and a
basis of principal units generate it over \(\mathbb Q_p\).
Their Kummer extensions split over \(B_{\rm st}\):
the logarithm of the compatible unit roots belongs to
\(B_{\rm cris}\), and the uniformizer logarithm belongs to
\(B_{\rm st}\). The explicit construction is also given in
\cite[Sections 11.7--11.10]{fghOlssonSemistable}; this
Kummer-period calculation works over every finite \(F\).
Thus there is
\[
 b_\kappa\in B_{\rm st}\otimes_{\mathbb Q_p}\mathbb Q_p(1),
 \qquad
                 \kappa(g)=g(b_\kappa)-b_\kappa .
                                                        \tag{A.31.2}
\]
The sign can of course be reversed with the convention for
the Kummer cocycle.

Let \(c:\mathbb Q_p(1)\to\mathfrak u\) be the actual cusp
map. The twisted action is
\(\rho_\kappa(g)=\exp(\operatorname{ad}c(\kappa(g)))\rho(g)\).
Over \(B_{\rm st}\), conjugation by
\(\exp(\operatorname{ad}c(b_\kappa))\), with the sign matching
(A.31.2), identifies this with the untwisted action.
The exponentials are finite in each nilpotent quotient.
This proves semistability for every packet, not just for
packets represented by an algebraic element of \(F^\times\).
It also shows directly why a packet can add a cusp-inner
term to \(N\). We shall allow that term rather than choose
a special packet.

The selected lift to either cover is itself cuspidal after
finite descent. Its ramification index over the original
cusp is three; at \(p=2\) division by three is integral.
Consequently the preceding local argument applies to all
the induced original packets in the statement.

\paragraph{A semistable representation with abelian image is crystalline.}
We give the linear-algebra proof needed below. Let \(W\)
be semistable and suppose that its Galois image is abelian.
The finite-dimensional commutative algebra generated by
that image splits after a finite extension of coefficient
fields. Decompose \(W\) into its generalized-character
blocks. Each character \(\eta\) is a semistable rank-one
subquotient, and hence crystalline: a nilpotent endomorphism
of its rank-one \(D_{\rm st}\) module is zero.
Twist its block by \(\eta^{-1}\). All composition factors
are now the trivial representation. Exactness of
\(D_{\rm st}\) implies that \(\varphi^f\) on this block
is unipotent, where \(f\) is the residue degree of \(F\).
The relation
\[
                       N\varphi^f=p^f\varphi^fN
                                                        \tag{A.31.3}
\]
then forces \(N=0\): it would carry the generalized
eigenspace with eigenvalue \(1\) into the generalized
eigenspace with eigenvalue \(p^{-f}\), which is zero.
Twisting back by the crystalline character preserves
this conclusion. It descends through the coefficient
extension, proving the assertion. This also covers
unipotent abelian representations; no semisimplicity
of the original image is assumed.

\paragraph{A fixed element supplies crystalline leading classes.}
Enlarge a finite number field \(L\) so that the cover,
packet lift, cusps, CM isogenies and marked semistable
models at a place \(v\mid2\) are defined, and
\[
       \Gamma=\overline{[G_L,G_L]}\subset\Delta_0 .
\]
This is possible because \(\Delta_0\) is open in \(\Delta\).
Put \(F=L_v\). If \(u\) is an actual fixed pro-\(2\)
element, its logarithm in a finite rational nilpotent
quotient \(\mathfrak u_n\) is fixed by \(\Gamma\).
The span
\[
 M_n=\left\langle \rho(g)\log(u)_n:g\in G_L
                                            \right\rangle_{\mathbb Q_2}
                                                        \tag{A.31.4}
\]
is also fixed by \(\Gamma\), since \(\Gamma\) is normal
in \(G_L\). Its Galois image is therefore abelian.
The first two paragraphs make \(\mathfrak u_n\), and
hence \(M_n\), semistable locally at \(v\).
The preceding paragraph gives
\[
                         N D_{\rm st}(M_n)=0.           \tag{A.31.5}
\]

Both proper Jacobians have good reduction after this
extension. Thus their \(H_1\) representations are
crystalline. For \(C\), ordinary open homology splits
rationally into the boundary line \(\mathbb Q_2(1)\)
and proper homology: the hyperelliptic involution acts
by \(+1\) on the former and by \(-1\) on the latter.
The projectors \((1\pm h)/2\) are legitimate over
\(\mathbb Q_2\). It follows that all ordinary graded
pieces \(\operatorname{gr}_d\mathfrak u\), free-Lie
constructions from \(H_1\), are crystalline.
After the chosen finite extension the global action
on \(H_1\) is abelian, given by the two CM characters
and, for \(C\), the cyclotomic boundary character.

Suppose that \(u\) has first nonzero degree \(d\), and
let \(W_d\) be the global \(G_L\)-span of its leading
class. For \(n\) sufficiently large, \(M_n\) is in
lower-central degree at least \(d\) and maps onto
\(W_d\). If \(N_r\) is the first nonzero graded
derivation of \(N\), exactness of \(D_{\rm st}\)
and (A.31.5) give
\[
                 D_{\rm cris}(W_d)\subset\ker N_r .
                                                        \tag{A.31.6}
\]
Indeed lift a vector of \(D_{\rm cris}(W_d)\) to
\(D_{\rm st}(M_n)\) and take the component of \(N\)
in degree \(d+r\). All smaller degree components
vanish, so that component is \(N_r\) of the leading
vector.

There is also a genuine CM-stability statement on
the left side of (A.31.6). Every global Galois operator
on \(H_1\) commutes with the local \(G_F\)-action,
because the global homology action is abelian.
It is therefore a morphism of crystalline
representations. The same is true on each free-Lie
graded piece and on \(W_d\). Applying \(D_{\rm cris}\)
shows that \(D_{\rm cris}(W_d)\) is stable under the
induced global CM operators. In particular it lies
in the intersection of \(\ker N_r\) with every one
of its CM translates.
All component maps, polarizations and CM
endomorphisms used in A.27--A.28 are compatible
with crystalline realization. Their displayed
linear formulas consequently hold here after a
coefficient extension splitting the CM field.
This proves the required transport statement
entirely inside the crystalline realization.
It does not identify \(N\) with an operator on
the original \(\mathbb Q_2\)-vector space.

\paragraph{The genus-two curve.}
The pointed stable model is the two-tail model of
A.27, with the cusp on its rational bridge.
Its ordinary-based residue has no degree-one term
and has degree-two term
\[
                           m_1T_1+m_2T_2,
                     \qquad m_1m_2\ne0.                \tag{A.31.7}
\]
Compare it with the actual packet's \(N\).
They have the same affine outer residue. Over a
characteristic-zero coefficient field a unipotent
path torsor has a point, so the two based residues
differ by an inner derivation. Homology is
crystalline, so both have degree zero equal to
zero. Their possible degree-one difference is
\(\operatorname{ad}(w_1)\).
The actual cusp map is crystalline, and thus
\(N(c)=0\). Its leading class is
\(\Omega=\omega_1+\omega_2\), so
\([w_1,\Omega]=0\). The free-Lie centralizer of
\(\Omega\) is its own line, in degree two;
therefore \(w_1=0\).
The degree-two difference is then
\(\operatorname{ad}(w_2)\), with
\([w_2,\Omega]=0\), and hence
\[
                  N_2=m_1T_1+m_2T_2
                                +\lambda\operatorname{ad}(\Omega)
                                                        \tag{A.31.8}
\]
for some scalar \(\lambda\). This allows every
packet-inner ambiguity.

The unbounded kernel calculation (A.27.3) applies
over this characteristic-zero field. If
\(\lambda\ne0\), its kernel is the line \(\Omega\).
If \(\lambda=0\), its kernel is
\(\mathbb L(\omega_1,\omega_2)\).
Choose a global ordinary CM Frobenius away from
two with eigenvalue ratio of infinite order.
The crystalline versions of (A.27.11)--(A.27.15)
give the intersection of this kernel with its
translate as the line \(\Omega\).
Equation (A.31.6) and CM stability therefore imply
that every nonzero leading fixed class is a
multiple of \(\Omega\) in degree two.
Faithfulness of crystalline comparison descends
this containment to etale graded homology.

\paragraph{The genus-three curve.}
Use the three-tail pointed model of A.28.
Its ordinary-based residue has zero first degree
and degree-two part
\[
 N_2(t)=0,\qquad
 N_2|_{V_i}=m_i\operatorname{ad}(\omega_i),
                         \quad m_i\ne0.                \tag{A.31.9}
\]
The same affine outer-residue comparison applies.
Because the actual selected cusp \(c\) is crystalline
and has leading degree-one class \(t\), the possible
inner degree-one term must be
\(\lambda\operatorname{ad}(t)\).
If \(\lambda\ne0\), its kernel is the line \(t\),
and (A.31.6) already gives the desired conclusion.
If \(\lambda=0\), the degree-two inner correction
is \(\operatorname{ad}(w_2)\) with \([w_2,t]=0\);
the centralizer of \(t\) has no degree-two part,
so \(w_2=0\). Thus (A.31.9) is the actual \(N_2\).

The first-core proof of A.28 gives its kernel
\(\mathbb L(t,\omega_1,\omega_2,\omega_3)\).
The component maps have degree one, so the
crystalline CM matrices are exactly (A.28.5).
Their diagonal-plane intersection and its
tensor-algebra consequence give
\[
         D_{\rm cris}(W_d)\subset
              \mathbb L\bigl(t,\Omega\bigr)_d,
                \qquad \Omega=\omega_1+\omega_2+\omega_3 .
                                                        \tag{A.31.10}
\]
The right side is itself the crystalline realization
of the etale subspace generated by the boundary
line and the proper polarization tensor. Taking
the crystalline image of the map to its quotient,
and using faithfulness, descends (A.31.10) to
\(W_d\subset\mathbb L(t,\Omega)_d\).

The first-Johnson nonvanishing established inside
the proof of A.30 is valid over \(\mathbb Q_2\)
as well. Its universal regulator comparison is
rational; the factors \(2\) and \(4\) are invertible
in this field. The explicit primitive functional
is still \(q_*-2(q_0)_*\), and its value is
\[
                      -4\,\kappa_{\,2+\sqrt3}.
                                                        \tag{A.31.11}
\]
The unit \(2+\sqrt3\) has infinite order, so its
rational \(2\)-adic Kummer class is nonzero after
every finite number-field extension. The restriction
to \(G_L\cap\Delta\) is injective on this class:
the abelian quotient has a nontrivial cyclotomic
scalar, so Sah's lemma kills its rational
\(H^1(-,\mathbb Q_2(1))\).
The restricted class stays nonzero on every open
subgroup, since it is then an additive map to a
characteristic-zero vector space.
Packet changes add only a scalar first-Johnson
matrix, killed by this primitive functional.
Hence the actual first Johnson matrix
\(A_\delta\) is nonzero for some
\(\delta\in\Delta_0\), for every packet.

The algebraic implication (A.29.4) uses only
characteristic zero and \(A_\delta\ne0\);
it has no odd-prime hypothesis. Combining it
with (A.31.10) puts every nonzero leading fixed
class on the line \(t\).
This use of A.29 does not invoke its separate
odd-prime integral uniqueness assertion.

\paragraph{Return to the actual pro-\(2\) groups.}
The lower-central filtration of a finitely generated
free pro-\(2\) group is separated, and its graded
pieces are the torsion-free integral free Lie
pieces. On \(D\) the selected cusp has primitive
leading class \(\Omega\) in degree two.
On \(C\) it has primitive class \(t\) in degree one.
Thus the leading coefficient of an actual fixed
element on the indicated line belongs to
\(\mathbb Z_2\). Subtract the corresponding
\emph{actual} cusp power. The remainder is fixed
and has strictly larger first degree; the
graded restrictions above force it to be the
identity. All cusp powers are fixed on
\(\Delta_0\), since its cyclotomic character is
trivial. This proves both equalities (A.31.1)
and completes the assertion for every reflection
quotient.
\end{proof}

\begin{statement}{A.32}{Proved}{The other symmetric genus-two quotients at every prime}
For \(a=6\) or \(-6\), put
\[
 D_a:\quad z^2=(t^2-4)(t^3-3t+a),\qquad
 Y_a=D_a\setminus\{Q_\infty\}.
\]
Let \(P_{a,\ell}\) be the geometric pro-\(\ell\) group of
\(Y_a\), based at the selected lift of the original cusp
packet of Part A. For every prime \(\ell\), every original
packet, and every permitted open \(\Delta_0\subset\Delta\),
\[
             P_{a,\ell}^{\,s_b(\Delta_0)}=I_{a,\ell}.
                                                        \tag{A.32.1}
\]
Here \(I_{a,\ell}\) is the full inertia of that selected cusp.
The assertion applies to every reflection quotient of the two
nonstable symmetric covers. Together with A.27 and A.31 it
determines all the symmetric genus-two quotient fixed groups
at every prime.
\end{statement}
\begin{proof}
\paragraph{The arithmetic and the pointed reduction.}
The actual map to the original elliptic curve is
\[
 f_a:D_a\longrightarrow E,\qquad
 x_E=\frac{t^3-3t+a/3}{4},\qquad
 y_E=\frac{z(t^2-1)}8.                                  \tag{A.32.2}
\]
Substitution, using \(a^2=36\), gives \(y_E^2=x_E^3-x_E\).
The pole order of \(x_E\) is six, so the map has degree
three and \(f_a^*O=3Q_\infty\). It is precisely the
descent of the cover in A.4, and is etale on the displayed
open curves. The complementary elliptic curve is
\[
                  F_a:\quad u^2=t^3-3t+a,
                  \qquad j(F_a)=-216.
\]
A.4 gives \(J_{D_a}\sim E\times F_a\) over a finite
number field. The non-CM assertion for \(F_a\) was proved
in A.14: for \(a=6\) its ordinary Frobenius discriminants
at five and eleven are \(-19\) and \(-35\); the other
sign is a geometric twist. An imaginary quadratic
endomorphism field could not be both of these fields.

Choose a finite number field \(L\) defining these data
and with \(\Gamma=\overline{[G_L,G_L]}\subset\Delta_0\).
To see this last inclusion, choose an open neighborhood
\(V\) of the identity in \(G_K\) with
\(V\cap\Delta\subset\Delta_0\), then take \(G_L\subset V\).
Enlarge \(L\) further to realize semistable reduction at a
place above two. Both elliptic factors have integral
\(j\), so both acquire good reduction there. Thus the
Jacobian acquires good reduction and the stable dual
graph of \(D_a\) is a tree.

We spell out why the two-tail model of A.25 still applies.
Its proof did not require the two elliptic factors to be
isogenous. The good special fibre of the original \(E\)
is supersingular in characteristic two. If \(D_a\) had
smooth stable reduction, (A.32.2) would specialize to a
separable cubic map with its only ramification a tame
point of index three. A cyclic Galois closure is
impossible by the degree-zero divisor condition modulo
three. An \(S_3\) closure would give a connected etale
double cover of the supersingular elliptic target,
also impossible. The compact-type genus-two fibre is
therefore two elliptic curves meeting in one node.

The map extends to the stable model by resolving its
rational map and contracting the rational exceptional
components, exactly as in A.25. Its two component degrees
are positive: a degree-three map on one elliptic component
would be etale and have three distinct points over the
origin, whereas the horizontal divisor \(3Q_\infty\)
and the node supply at most two. Their sum is three,
so the degrees are one and two. The divisor
\(f_a^*O=3Q_\infty\) has no vertical component, since
both component maps are nonconstant. If the cusp met
the smooth part of one component, the other component
would have no preimage of the origin, contradicting
surjectivity. Thus the cusp specializes to the node.
The pointed model inserts a rational bridge containing
the cusp, with two strictly positive arm lengths.

\paragraph{The separating derivation.}
Write \(V=V_1\oplus V_2\) for component homology and
choose symplectic bases \(a_1,b_1\) and \(a_2,b_2\).
Put \(\omega_i=[a_i,b_i]\) and
\(\Omega=\omega_1+\omega_2\).
For odd \(\ell\), apply the affine edge-twist theorem
to an ordinary point on the rational bridge. A local
Frobenius--tame commutator belongs to \(\Gamma\), acts
trivially on homology, and has first possible outer
Johnson degree two, with representative
\[
                        m_1T_1+m_2T_2,
                         \qquad m_1m_2\ne0.             \tag{A.32.3}
\]
Here the \(T_i\) are those of A.27, and \(m_i\) is
the corresponding positive arm length times a common
nonzero tame scalar. The actual packet action differs
only by an inner action. Its possible degree-one term
is \(\operatorname{ad}(w_1)\). Since it fixes the actual
cusp, \([w_1,\Omega]=0\), and the free-Lie centralizer
of \(\Omega\) gives \(w_1=0\). In degree two the same
argument makes the inner term a scalar multiple of
\(\operatorname{ad}(\Omega)\). Hence the actual
derivation is \(D_\lambda\) of A.27. This reasoning
uses only the local commutator's trivial homology
action; the whole \(\Gamma\)-action on homology need
not be trivial.

At \(\ell=2\), A.31's tangential semistable comparison
applies to this same pointed model and to every
induced packet. In its crystalline graded realization,
the identical boundary-centralizer argument gives
\[
                       N_1=0,\qquad N_2=D_\lambda.
                                                        \tag{A.32.4}
\]
The two coefficients in (A.32.3) are nonzero node
thicknesses, with their common monodromy normalization.
No comparison of a pro-two group with its prime-to-two
covering quotient is involved.

\paragraph{The non-CM factor supplies the missing invariants.}
Let \(V_E=f_a^*H_1(E)\) and let \(V_F\) be its
polarization-orthogonal complement. They are actual
\(G_L\)-subrepresentations, with
\(V=V_E\oplus V_F\), and \(V_F\simeq H_1(F_a)\).
The action of \(\Gamma\) on \(V_E\) is trivial,
because \(E\) has CM defined over \(L\).
Serre's open-image theorem for the non-CM curve
\(F_a\) gives an open image in
\(\mathrm{GL}_2(\mathbb Z_\ell)\), for every \(\ell\)
\cite[Introduction (1)--(2), Theorem 3]{A_Serre1972}.
Its commutator image is Zariski dense in
\(\mathrm{SL}(V_F)=\mathrm{Sp}(V_F)\): the Lie
algebra of that image contains
\([\mathfrak{gl}_2,\mathfrak{gl}_2]=\mathfrak{sl}_2\),
and its determinant is one. Consequently every
\(\Gamma\)-fixed vector of \(\mathbb L_d(V)\) is
fixed by the algebraic group \(\mathrm{Sp}(V_F)\),
acting as the identity on \(V_E\).

We compute the relevant intersection over any
characteristic-zero coefficient field. The component
degrees one and two permit normalizations
\[
 f_*a_1=x,\quad f_*b_1=y,\quad
 f_*a_2=x,\quad f_*b_2=2y.
\]
Adjointness gives
\[
 \begin{split}
 V_E&=\langle A,B\rangle,
       \qquad A=a_1+2a_2,\quad B=b_1+b_2,\\
 V_F&=\langle C,D\rangle,
       \qquad C=a_1-a_2,\quad D=2b_1-b_2.
 \end{split}                                             \tag{A.32.5}
\]
The two pairings are \(\langle A,B\rangle
=\langle C,D\rangle=3\); all cross pairings vanish.
Choose \(g\in\mathrm{Sp}(V_F)\) given by
\(C\mapsto2C\), \(D\mapsto D/2\), and the identity
on \(V_E\). In the \(a\)-coordinates its matrix is
\[
 M=\frac13\begin{pmatrix}5&-1\\-2&4\end{pmatrix},
\]
and its \(b\)-matrix is \((M^{-1})^T\).
Every entry of \(M\) is nonzero. Thus, for the
quadratic plane \(U=\langle\omega_1,\omega_2\rangle\),
\[
                              U\cap gU=k\Omega.         \tag{A.32.6}
\]
Indeed this is the intersection of diagonal matrices
with \(M\operatorname{Diag}M^{-1}\). If
\(DM=MD'\) with \(D,D'\) diagonal, its four entries
force all diagonal entries to be the same.
The tensor-intersection proof of A.27 then gives
\[
              \mathbb L(U)\cap\mathbb L(gU)=k\Omega.
                                                        \tag{A.32.7}
\]
The explicit matrix \(g\) is an algebraic symmetry
of the invariant space; we do not claim it is itself
a specified Galois element.

\paragraph{The prime-two invariant comparison.}
For odd \(\ell\), a leading class of a fixed element
is in \(\ker D_\lambda\) and is fixed by
\(\mathrm{Sp}(V_F)\). A.27 immediately gives the
line \(k\Omega\) when \(\lambda\ne0\); when
\(\lambda=0\), (A.32.7) gives the same conclusion.

For \(\ell=2\), let \(W_d\) be the global orbit
span of the leading class. The corresponding orbit
span in a sufficiently large finite Lie quotient
is fixed by \(\Gamma\) and has abelian global
Galois image. A.31 proves it crystalline, so
\[
                    D_{\rm cris}(W_d)\subset\ker N_2.
                                                        \tag{A.32.8}
\]
We must also transport the algebraic invariants;
non-CM global operators need not commute with the
local Galois action. Define the invariant subspace
instead as the kernel of the canonical infinitesimal
action morphism
\[
 \partial:\mathbb L_d(V)\longrightarrow
       \mathbb L_d(V)\otimes\mathfrak{sl}(V_F)^\vee,
              \qquad \partial(w)(T)=T(w).                \tag{A.32.9}
\]
This is a morphism of crystalline representations:
both summands of \(V\) are crystalline, and the
map is formed from their projections, tensors and
brackets. Exact tensoriality of \(D_{\rm cris}\)
therefore identifies its kernel with the analogous
kernel on crystalline realizations. Since
\(\mathrm{Sp}(V_F)\) is connected in characteristic
zero, its group invariants equal these Lie-algebra
invariants. Thus \(D_{\rm cris}(W_d)\) is fixed
by every element of
\(\mathrm{Sp}(D_{\rm cris}V_F)\).

Component maps and their polarization adjoints
give (A.32.5) also in crystalline homology; the
degree-two component map is invertible there
over the coefficient field. Apply the abstract
symplectic matrix in (A.32.6) within this
realization. Equations (A.32.8), (A.32.7) and
A.27 put \(D_{\rm cris}(W_d)\) on the line
\(\Omega\) in degree two. That line is the
crystalline image of the etale cusp line.
Faithfulness of the crystalline functor applied
to the quotient by that line gives
\(W_d\subset\mathbb Q_2\Omega\).
This argument never applies Fontaine monodromy
to an individual etale vector.

\paragraph{The full pro-prime conclusion.}
For every prime we have proved that a nonidentity
fixed element has first nonzero degree two and
leading class on the actual cusp line. In the
integral free Lie algebra that cusp class is
primitive: it is the sum of the two standard
handle commutators. Its coefficient is therefore
in \(\mathbb Z_\ell\), including at two and three.
Subtract the corresponding actual cusp power.
The remainder is fixed and has strictly larger
first degree, which the proved restriction
excludes. Separation of the lower-central
filtration gives the identity. The reverse
inclusion follows from the trivial cyclotomic
character on \(\Delta_0\). This proves (A.32.1).
\end{proof}

\begin{statement}{A.33}{Proved}{The nonstable symmetric genus-three covers at two and their odd-prime reduction}
For $a=6$ or $-6$, let
\[
 C_a:\quad v^2=w^7+aw^4+w,\qquad
 X_a=C_a\setminus\{P_0,P_\infty\}.
\]
Use the lift of any original packet of Part A at a selected
cusp, and any permitted open subgroup $\Delta_0\subset\Delta$
preserving this cover and that lift. Then the full pro-$2$
fixed subgroup is exactly the selected cusp inertia:
\[
 \pi_1(X_{a,\overline K})^{(2),\,s_b(\Delta_0)}
                       =I_{a,2}.                       \tag{A.33.1}
\]
For every odd prime $\ell$, every leading ordinary
lower-central class of a fixed element belongs, after the
geometric identifications in the proof, to
\[
 \mathbb L_{\mathbf Q_\ell}(t,\Omega)
 \ \subset\
 \mathbb L_{\mathbf Q_\ell}
       (\,\mathbf Q_\ell t\oplus V_1\oplus V_2\oplus V_3\,),
 \qquad \Omega=\sum_{i=1}^3[a_i,b_i].                  \tag{A.33.2}
\]
Here $t$ is the primitive homology class of the chosen
cusp, and $a_i,b_i$ are symplectic generators of the three
elliptic components at two. Their ordinary degrees are
$\deg t=1$ and $\deg\Omega=2$.
If the actual local commutator has nonzero first Johnson
degree, the odd-prime fixed subgroup is also exactly
the chosen inertia. Otherwise (A.33.2) is a reduction
of the remaining problem; it does not assert the
existence of an extra fixed element.
Statement A.38 below completes this remaining case
by upgrading to the weight filtration and applying
the cyclotomic fixed-subgroup criterion.
\end{statement}
\begin{proof}
\paragraph{The unramified double cover settles the prime two.}
The fixed-point-free reflection is
\[
 j(w,v)=(1/w,-v/w^4).
\]
Its quotient map is
\[
 t_D=w+1/w,\qquad
 z_D=(w-1/w)v/w^2,
\]
and gives the actual curve of A.32:
\[
 z_D^2=(t_D^2-4)(t_D^3-3t_D+a).
\]
The identities follow by substitution. A fixed point
of $j$ away from the cusps would have $w=\pm1$ and
$v=0$. The two values of the polynomial are $a+2$
and $a-2$, both nonzero. The two cusps are exchanged,
so there are no remaining fixed points. Thus
$C_a\to D_a$ is an unramified double cover of proper
curves, and $X_a\to Y_a$ is its restriction.
The two cusps map unramifiedly to the unique cusp
of $Y_a$. The chosen tangential base point therefore
maps to the chosen tangential base point there.

Write $H_C\triangleleft H_D$ for the corresponding
geometric profinite groups. Their quotient is $C_2$.
The elementary core argument in A.22, applied with
ambient group $H_D$, proves that $H_C^{(2)}$ embeds
as an open subgroup of index two in $H_D^{(2)}$.
This assertion does not use any CM hypothesis:
an open normal subgroup with $2$-group quotient has
its pro-$2$ topology induced from the ambient group.
The embedding is equivariant after a permitted
finite extension. A.32 consequently gives
\[
 (H_C^{(2)})^{s_b(\Delta_0)}
       =H_C^{(2)}\cap (H_D^{(2)})^{s_b(\Delta_0)}
       =H_C^{(2)}\cap I_{D,2}=I_{a,2}.
\]
The last equality uses ramification index one at
the selected cusp. If an auxiliary extension was
needed, prove this first for its smaller open
subgroup of $\Delta_0$, and then use the inclusions
of fixed subgroups and the fact that $\Delta_0$
fixes the selected inertia. This proves (A.33.1)
for the originally stated $\Delta_0$.
For odd $\ell$, the double cover instead induces
a surjection on maximal pro-$\ell$ quotients;
the preceding embedding argument is unavailable.

\paragraph{The pointed three-tail reduction at two.}
The cyclic cubic quotient is
\[
 f:C_a\longrightarrow E_{a,1},\qquad
          (w,v)\longmapsto(w^3,wv),
 \qquad E_{a,1}:\ y^2=x^3+ax^2+x.
\]
It is the Kummer cover $w^3=x$, and
$\operatorname{div}(x)=2T-2O$ with $T=(0,0)$
of order two. The curve $E_{a,1}$ is $2$-isogenous
to the original $E$. At the place above two,
pass to a finite extension over which it has good
reduction. The special elliptic fibre is
supersingular. Its geometric two-torsion is
trivial, so $T$ and $O$ meet in the special fibre.
The stable two-pointed model is an elliptic
component joined to a rational component bearing
the two marks. A regular model may contain a
chain of intervening unmarked rational components;
contracting them replaces the chain by a node
of strictly positive thickness.

Normalize this marked model in $w^3=x$.
After a finite extension, all vertical Kummer
valuations are divisible by three. The remaining
normalization is tame, and its local equations
are cube roots of monomials in the marked and
nodal parameters. On the elliptic component,
the node is the only possible branch point.
The degree-zero divisor condition modulo three
excludes branching at just one point. On the
rational component, the two marked orders are
$2,-2$ modulo three, and its node is unramified.
Its connected cyclic cubic cover has genus zero
and three points above the node. At an unramified
node the strict local normalization has three
ordinary nodal branches. Any regular subdivisions
add only rational vertices.

There are two possible covers above the elliptic
component: one connected unramified cubic isogeny,
or three disjoint copies. The connected case has
three edges joining two components, hence graph
first Betti number two. It is excluded by the
Jacobian. A.4 gives
$J_{C_a}\sim E\times F_a^2$, where
$F_a:u^2=t^3-3t+a$ has $j=-216$.
Both factors have potentially good reduction at
two, so the Jacobian has good reduction after a
finite extension. Its torus rank, equal to the
first Betti number of the stable dual graph,
is zero. The split case is therefore forced:
three elliptic tails attached to one rational
center bearing both cusps. The rotation permutes
the tails and preserves their positive arm
lengths. The degree-one maps of the tails to
the good elliptic quotient identify their
symplectic homology spaces.

\paragraph{The actual separating derivation.}
The hyperelliptic involution acts by $+1$ on the
boundary class $t$ and by $-1$ on proper homology.
It therefore gives the equivariant rational
splitting $H_1(X_a)=\mathbf Q_\ell t\oplus H_1(C_a)$
used here. On the three tails it is the elliptic
involution, so their homology classes lie in this
same negative eigenspace. This identifies the
component generators with actual open-curve
homology without adding unspecified boundary
components.
Choose a finite number field $L$ defining the
cover, its deck action and the required data,
with $\Gamma=\overline{[G_L,G_L]}\subset\Delta_0$,
as in A.32. For odd $\ell$, a local tame inertia
element acts trivially on homology after the
preceding semistable extension. Its commutator
with a local Frobenius belongs to $\Gamma$.
The affine edge-twist theorem, applied at an
ordinary point on the rational center, gives
zero first Johnson degree and second-degree
derivation
\[
 D_0(t)=0,\qquad
 D_0|_{V_i}=m_i\operatorname{ad}(\omega_i),\qquad
 \omega_i=[a_i,b_i],\qquad m_i\ne0.
\]
The $m_i$ are the positive arm lengths times a
common nonzero tame commutator scalar.
This uses the same based affine monodromy
statement as A.27 and A.28.

The actual packet representative has the same
outer action. Its difference is inner. Since
both representatives fix the selected cusp,
the possible degree-one inner term is
$\lambda\operatorname{ad}(t)$. If $\lambda\ne0$,
its kernel on the free rational Lie algebra is
$\mathbf Q_\ell t$. Integral primitivity of
the cusp class allows subtraction of the
corresponding actual cusp power. Any nonidentity
remainder would have a higher leading fixed
class, which this kernel excludes. The separated
lower-central filtration then gives precisely
the inertia fixed subgroup.

If $\lambda=0$, both first-degree actions vanish.
Their degree-two difference is
$\operatorname{ad}(z_2)$ and fixes $t$, forcing
$z_2=0$ by the free-Lie centralizer of $t$.
Thus $D_0$ is the actual second Johnson
derivation. The all-degree first-core argument
of A.27, with three elliptic blocks and the
extra letter $t$, gives
\[
 \ker D_0=\mathbb L_{\mathbf Q_\ell}
                      (t,\omega_1,\omega_2,\omega_3).
\]
The proof applies because each $m_i$ is nonzero;
no equality of their values is required.

\paragraph{The non-CM invariant intersection.}
Put $A=\sum a_i$ and $B=\sum b_i$.
The original elliptic subrepresentation is
$V_E=\langle A,B\rangle$, and its symplectic
orthogonal complement is
\[
 V_F^{\oplus2}
   =\left\langle \delta a_i,\delta b_i:1\le i\le3\right\rangle,
 \qquad
 \delta a_i=a_i-A/3,\quad \delta b_i=b_i-B/3.
\]
The deck group acts by permutations on the
three tail indices. Indeed its action on the
elliptic quotient is either the identity or
translation by $T$, both trivial on homology.
The sum-zero part is its standard rational
representation with multiplicity space
$H_1(F_a)$. Thus the arithmetic action on
that part is the same two-dimensional action
on every pair $\delta a_i,\delta b_i$.

The $\Gamma$-action is trivial on $V_E$ and
Zariski dense in $\mathrm{SL}(H_1(F_a))$ on
the multiplicity space, by the non-CM assertion
and Serre theorem checked in A.32. Every
leading class of a $\Gamma$-fixed group
element is therefore fixed by this algebraic
special linear group. In particular it is
fixed by its symplectic unipotent element
\[
 g(a_i)=a_i,\qquad
 g(b_i)=b_i+\delta a_i.
\]
This element fixes $A,B,t$ and commutes with
the deck permutations. It satisfies
\[
 g(\omega_i)=\omega_i+\frac13[A,\delta a_i].
\]
Let $U=\langle\omega_1,\omega_2,\omega_3\rangle$.
If $g(\sum d_i\omega_i)\in U$, its component
in $\bigwedge^2\langle a_1,a_2,a_3\rangle$
is zero. Hence $[A,\sum d_i\delta a_i]=0$.
The vector in the bracket belongs to the
sum-zero plane, whereas the degree-one
centralizer of $A$ is its line. It is
therefore zero, which means $d_1=d_2=d_3$.
We have proved
\[
                         U\cap gU=\mathbf Q_\ell\Omega.
\]

Finally, decompose associative words by the
occurrences of $t$, and pair the letters
in every intervening even-length block.
On each paired tensor block the intersection
of the two subalgebras is the tensor algebra
on $U\cap gU$. This is the same exact
tensor-intersection argument as in A.27.
It gives
\[
 \mathbb L(t,\omega_1,\omega_2,\omega_3)
       \cap g\mathbb L(t,\omega_1,\omega_2,\omega_3)
                      =\mathbb L(t,\Omega).
\]
An actual leading fixed class belongs to the
first algebra by monodromy, and to its
conjugate because it is fixed by $g$.
This proves (A.33.2), including $\ell=3$:
all divisions by three here occur in the
rational graded argument, not in an asserted
integral group operation. The full pro-$2$
result was obtained independently by the
actual induced-topology embedding.

The exact program
\texttt{code/a\_nonstable\_C3\_fixed.py} verifies
the curve and quotient identities, the
symplectic matrix, its deck commutation,
and the rank-five sum of the two
three-dimensional boundary spaces.
Its JSON and stdout are retained.
\end{proof}

\subsection{An all-prime fixed-subgroup comparison for the tripod}

\begin{statement}{A.34}{Proved}{The tripod over an abelian closure, for every cusp packet}
Let $k$ be a number field in $\overline{\Q}$, let
$V=\mathbf P^1_k\setminus\{0,1,\infty\}$, and fix a prime $\ell$.
Write $F_\ell=\pi_1(V_{\bar k},\overrightarrow{01})^{(\ell)}$
and $I_0=\langle x\rangle\simeq\Z_\ell(1)$ for its actual inertia
at zero. Let $s_b:G_k\longrightarrow F_\ell\rtimes G_k$ be any
continuous section in this cusp packet, with its representative
chosen in $I_0\rtimes G_k$. Put
$\Delta_k=\operatorname{Gal}(\bar k/k^{\rm ab})$.
For every open subgroup $\Delta_0\subset\Delta_k$ we have
\[
                   F_\ell^{s_b(\Delta_0)}=I_0.        \tag{A.34.1}
\]
The same assertion holds at either of the other two cusps after
transporting the base point and inertia by the corresponding
automorphism of $V$. It includes $\ell=2,3$, arbitrary completed
tangent parameters, and packet representatives whose Kummer
character vanishes on $\Delta_0$.

The elementary algebra behind this assertion is the following
unbounded calculation. On the free Lie algebra
$L=\mathbb L_{\Q_\ell}(X,Y)$ set
\[
 F_3=[X,[X,Y]]+[Y,[X,Y]],\qquad
 dX=0,\qquad dY=[Y,F_3].                               \tag{A.34.2}
\]
Then $\ker(d:L\to L)=\Q_\ell X$.
The comparison (A.34.1) is a statement about the actual free
pro-$\ell$ group, not just its rational Lie algebra. It supplies
a target fixed-subgroup calculation for Task A.4; it does not
assert that a cover's maximal pro-$\ell$ quotient embeds in
that of this target.
\end{statement}
\begin{proof}
We first prove the algebraic claim over any field of
characteristic zero. Embed $L$ in the free associative algebra
$T=\Q_\ell\langle X,Y\rangle$. On words of a fixed length use
lexicographic order with $X>Y$. Expanding (A.34.2) gives
\[
\begin{split}
 dY={}&-X^2Y^2+2XYXY-2YXYX+Y^2X^2\\
      &+3Y^2XY-3YXY^2+XY^3-Y^3X .
\end{split}
                                                        \tag{A.34.3}
\]
Its leading word is $-X^2Y^2$. Every word containing $Y$ has a
unique expression $w=X^aYv$ where the displayed $Y$ is its first.
In the Leibniz expansion of $dw$, differentiation at that first
$Y$ has leading word
\[
                         -X^{a+2}Y^2v.               \tag{A.34.4}
\]
Differentiating any later $Y$ leaves a $Y$ in position $a+1$,
whereas (A.34.4) has an $X$ there. All such terms are smaller.
The leading words (A.34.4), as $w$ varies in one degree, are
distinct and preserve the lexicographic order: different values
of $a$ are compared at the first $Y$, and equal values are
compared in $v$. Thus a nonzero linear combination of words
containing $Y$ cannot have derivative zero. Consequently
\[
                    \ker(d:T\to T)=\Q_\ell[X].
\]
Its intersection with $L$ is $\Q_\ell X$. For example this last
fact follows from the primitive coproduct: $X^n$ is not
primitive for $n>1$ in characteristic zero. This proves the
kernel assertion in every degree. The same word argument, or
comparison of the first and last letters, gives
\[
       \ker(\operatorname{ad}X:L\to L)=\Q_\ell X.    \tag{A.34.5}
\]

We next obtain an actual Galois element with first derivation a
nonzero multiple of $d$. The arithmetic input is the nonzero
degree-three Soul\'e class in the actual outer Johnson filtration.
Precisely, with
\[
 I_m^\ell G_\Q=\ker\{G_\Q\longrightarrow
             \operatorname{Out}(F_\ell/\gamma_{m+1}F_\ell)\},
\]
the rational vector space
$(I_3^\ell G_\Q/I_4^\ell G_\Q)\otimes\Q_\ell$ is nonzero.
This is the $m=1$ case of
\cite[Remark 8.6, together with Theorem 8.4 and
the filtration comparison in Proposition 8.2]{A_HainMatsumoto2003}.
That source states the nonvanishing for every prime; it is not
an assertion only for regular primes or only for $\ell>3$.
The comparison with the based tangential action is specified
in its Section 8.5. Only this degree-three nonvanishing is
used here, not the freeness of the full arithmetic Lie algebra.

Choose $h\in I_3^\ell G_\Q$ with nonzero rational degree-three
class. For the unit tangential section its genuine action is
\[
 \rho_0(g)x=x^{\chi_\ell(g)},\qquad
 \rho_0(g)y=f_g^{-1}y^{\chi_\ell(g)}f_g,
                  \qquad f_g\in F_\ell'.             \tag{A.34.6}
\]
Here the interval path from zero to one fixes the normalization;
in particular $f_g$ has zero abelianization and satisfies the
two-cycle identity $f_g(x,y)f_g(y,x)=1$.
We check that passing from the outer filtration to (A.34.6)
introduces no lower inner term. Since $\chi_\ell(h)=1$ and
$f_h\in\gamma_2F_\ell$, this action is already trivial modulo
$\gamma_3F_\ell$. If its possible degree-two derivation were
inner, its inner element's degree-one part would have zero
bracket with both $X,Y$, and hence be zero. Its degree-two
part would commute with $X$, and hence be zero by (A.34.5).
Therefore the degree-two derivation vanishes. The degree-two
part of $f_h$ is a multiple of $[X,Y]$; its bracket with $Y$
is zero only for the zero multiple. Thus
$f_h\in\gamma_3F_\ell$.

Write its degree-three part as
$a[X,[X,Y]]+b[Y,[X,Y]]$. Interchanging $X,Y$ sends these two
basis elements to the negatives of one another. The two-cycle
identity consequently gives $a=b$. Nonzero outer degree three
forces $a\ne0$, since a zero degree-three part would make the
actual action trivial modulo $\gamma_5F_\ell$. The first
derivation of $\rho_0(h)$ is therefore $a d$, up to the
irrelevant simultaneous sign convention for group commutators.

Take a positive integer power of $h$ to put it in the open
subgroup $G_k$; its first derivation is multiplied by that
integer and stays nonzero. Choose $g\in G_k$ with
$\chi_\ell(g)^3\ne1$, which is possible because the cyclotomic
image of a number field is open in $\Z_\ell^\times$. Conjugation
by $\rho_0(g)$ multiplies a degree-three derivation by
$\chi_\ell(g)^3$: on homogeneous degree $n$ the underlying
associated graded action is multiplication by
$\chi_\ell(g)^n$. Hence the actual commutator
\[
                         e=[g,h]\in\Delta_k
\]
has first derivation
$(\chi_\ell(g)^3-1)a d\ne0$. It still belongs to
$I_3^\ell G_\Q$. Taking a further positive integer power puts
$e$ in the specified open subgroup $\Delta_0$, again without
killing this rational leading derivation. This establishes
the required arithmetic element in every $\Delta_0$ by actual
group operations.

Every representative of the packet has the form
$s_b(g)=x^{\kappa(g)}s_0(g)$ for an integral continuous
cocycle $\kappa:G_k\to\Z_\ell(1)$. Thus
\[
                 \rho_b(g)=\operatorname{Inn}(x^{\kappa(g)})
                                  \rho_0(g).          \tag{A.34.7}
\]
For the element $e$ just constructed, there are two cases.
If $\kappa(e)\ne0$, the first derivation of $\rho_b(e)$ is
the degree-one derivation $\kappa(e)\operatorname{ad}X$.
If $\kappa(e)=0$, its first derivation is the nonzero multiple
of $d$ already obtained. Both kernels on $L$ are exactly
$\Q_\ell X$. This treats the zero-Kummer packet as well as all
other packets; no generic tangent condition is imposed.

Finally, the closed lower central filtration of the free
pro-$\ell$ group is separated, with
\[
 \operatorname{gr}_\gamma F_\ell
                  \simeq\mathbb L_{\Z_\ell}(X,Y).
\]
One can use Hall commutators to verify this integral comparison:
in each nilpotent quotient their ordered exponent coordinates
are in $\Z_\ell$, and their successive commutator coordinates
are the free Lie basis. Separation also follows directly by
detecting a nonidentity element in a finite $\ell$-group,
which is nilpotent. In particular every nonzero graded class
remains nonzero after tensoring with $\Q_\ell$.

If $u\in F_\ell$ is fixed by $\rho_b(e)$, its first nonzero
lower central class must be killed by the first derivation.
For a class of degree $n$, this follows by taking the
degree $n+j$ component of $\rho_b(e)u\,u^{-1}$, where $j$ is
the first derivation's degree. The two kernel calculations
show that $n=1$ and the abelianization of $u$ is $cX$.
Its coefficient $c$ lies in $\Z_\ell$ because the integral
abelianization is $\Z_\ell X\oplus\Z_\ell Y$. The element
$ux^{-c}$ is still fixed and belongs to $\gamma_2F_\ell$;
the same argument forces it to be one. We conclude
$F_\ell^{\rho_b(e)}=\langle x\rangle$. Since all of
$\Delta_0$ has cyclotomic character one, it fixes $x$;
intersecting fixed groups proves (A.34.1).

The program \texttt{code/a\_tripod\_fixed.py} independently
expands (A.34.3) over the integers and verifies the unique,
order-preserving leading word on every word through degree
twelve. These finite checks support the displayed identities;
the unbounded proof is the word argument above.
\end{proof}

\subsection{The quaternion cover: pointed reduction and the proper fixed group}

\begin{statement}{A.35}{Proved}{The quaternion cover has no proper fixed image at odd primes}
Let \(C=C_Q\) be the smooth quartic \(X^4+Y^4=2Z^4\), and
let \(S=\{Z=0\}\). After a finite extension the pair \((C,S)\)
has the cyclic description
\[
             U^8=-64B(B-1)^2,\qquad
             C\setminus S\longrightarrow
                    \mathbf P^1\setminus\{0,1,\infty\}.
                                                        \tag{A.35.1}
\]
This is an \'etale cyclic cover of degree eight. There are respectively
one, two, and one points above \(0,1,\infty\), with ramification
indices \(8,4,8\).

In the unpointed stable reduction at \(2\), the three elliptic
components meet a central rational component. All four points of
\(S\) specialize to the same node: the node adjacent to the unique
elliptic component on which the central involution of the original
\(Q_8\)-cover acts trivially. In particular, the four marked points
are not distinct smooth points of that central component.

Let \(P\in S\), let \(k\) be a number field over which \(P\) and
the curve are defined, and let \(\bar s_P\) be the section of the
proper arithmetic fundamental group defined by \(P\). For every
odd prime \(\ell\) and every open subgroup
\(\Delta_0\subset\operatorname{Gal}(\bar k/k^{\mathrm{ab}})\),
\[
 \left(\pi_1(C_{\bar k},P)^{(\ell)}\right)^{
                                  \bar s_P(\Delta_0)}=1.
                                                        \tag{A.35.2}
\]
Consequently every fixed element in the corresponding punctured
pro-\(\ell\) group, for every tangent packet at \(P\), belongs to
the closed normal subgroup generated by the four cusp inertias.
This consequence does not identify that punctured fixed subgroup
with the selected inertia; the normal subgroup in question is
much larger than a single inertia group.
\end{statement}

\begin{proof}
\paragraph{The cyclic model and the specialization of the cusps.}
Over a finite extension, put the quartic in the form
\(W^4=1-T^4\), with \(S=\{W=0\}\). Set
\[
 z=\frac{T-1}{i(T+1)},\qquad B=z^2,\qquad
 U=W(1-iz).
\]
Then \(T=(1+iz)/(1-iz)\) and
\[
 (1-T^4)(1-iz)^4=8iz(z^2-1).
\]
Squaring gives (A.35.1). Conversely \(z=U^4/(8i(B-1))\),
so these formulas identify the function fields. The Kummer
exponents at \(0,1,\infty\) are \(1,2,-3\), giving the stated
indices and numbers of points. The four ramification points are
exactly \(W=0\). The program
\texttt{code/fgh\_q8\_models.py} checks the displayed identities
over \(\mathbf Q(i)\).

The unpointed stable fibre with a central rational component
and three elliptic tails is the explicit computation of
\cite[Example 10.4.39, pp.\ 548--549]{fghLiu};
it is also recorded in
\cite[Table 7.1, p.\ 50]{fghQuarticReduction}.
Only that established Fermat computation is used here, not
the latter paper's conjectural general reduction criteria.
The elliptic tails are supersingular: the good special fibre
of the Jacobian is their product and is isogenous to a power
of the supersingular reduction of the CM elliptic curve
with \(j=1728\).

Let \(\tau\) generate the cyclic group in (A.35.1), and put
\(\sigma=\tau^4\). In the quartic coordinates
\(\tau^2\) acts by \(W\mapsto iW\), while
\(\sigma\) acts by \(W\mapsto-W\).
Thus \(C/\langle\tau^2\rangle\) has genus zero and
\(C/\langle\sigma\rangle\) has genus one. On rational
prime-to-two Tate homology, the \(\tau^2\)-invariants are
therefore zero and the \(\sigma\)-invariants have dimension two.

The permutation of the three tails induced by \(\tau\) is a
transposition. Indeed its order divides eight, so the only
alternatives are a transposition and the identity. In the latter
case it would act on each elliptic tail as an automorphism
fixing its attachment point. Such an automorphism cannot have
order eight: a finite-order elliptic endomorphism has an
integral quadratic characteristic polynomial, whereas a
primitive eighth root has degree four over \(\mathbf Q\).
Its fourth power would therefore be the identity on each
tail. This contradicts the dimension of the
\(\sigma\)-invariant homology just calculated.

Write \(E_0\) for the fixed tail and \(E_1,E_2\) for the
exchanged tails. On \(E_0\), the action of \(\tau\) has order
four: its fourth power is the identity, and its square cannot
have invariant homology. Its only geometric fixed point is
the attachment point. To see the uniqueness, take that point
as origin. The endomorphism \(\tau-1\) has degree two, and a
supersingular elliptic curve in characteristic two has no
nonzero geometric two-primary torsion. On the central
rational component, \(\tau\) exchanges two of the three
attachment points and fixes the third. A nonidentity
involution of \(\mathbf P^1\) in characteristic two has a
unique fixed point, namely this third attachment point.

The cusps above \(B=0,\infty\) are fixed by \(\tau\). Their
specializations cannot lie on the two exchanged tails, and
the preceding fixed-point calculations force them to the
node adjacent to \(E_0\). Finally the original \(Q_8\) deck
group acts transitively on \(S\) and commutes with its central
involution \(\sigma\). This involution acts trivially on
\(E_0\) and as minus one on each of the other two tails:
the fixed-tail contribution already exhausts its
two-dimensional invariant homology. The \(Q_8\) action
therefore preserves \(E_0\) and its attachment node.
Transitivity puts all four cusp specializations at that node.
This proves the claimed pointed correction without requiring
any assertion about the full graph permutation of a coordinate
sign involution.

\paragraph{The actual based monodromy.}
Retain just the one marked point \(P\). In the stable pointed
model, its specialization lies on a rational bridge inserted
between \(E_0\) and the original central component. Let the
two new edge lengths be \(a,b>0\), with \(a\) on the \(E_0\)
side, and let \(d_1,d_2>0\) be the two remaining elliptic-arm
lengths. These inequalities are strict: locally the original
node has equation \(uv=\varpi^e\), and a section through
the node has \(v(u(P)),v(v(P))>0\); these are precisely the
two lengths after a finite extension. Rational components
of valence two may be contracted, adding their lengths.

Put \(V_i=H_1(E_i,\mathbf Q_\ell)\), choose symplectic bases
\(x_i,y_i\), and write \(\omega_i=[x_i,y_i]\).
The rational graded Lie algebra of the proper curve is
\[
 \mathfrak s=
 \mathbb L(V_0\oplus V_1\oplus V_2)/
                  (\omega_0+\omega_1+\omega_2).
                                                        \tag{A.35.3}
\]
The based nonabelian Picard--Lefschetz formula
\cite[Definition 3.1.5 and Theorem 3.1.6]{fghBettsDogra},
applied with the base point on the displayed bridge, gives
the first monodromy derivation
\[
 \begin{split}
 N(v_0)&=a[\omega_0,v_0],\\
 N(v_i)&=[b(\omega_1+\omega_2)+d_i\omega_i,v_i],
                                      \qquad i=1,2.
 \end{split}                                             \tag{A.35.4}
\]
The term with coefficient \(b\) is the twist around the side
containing both \(E_1,E_2\). Thus (A.35.4) is a based formula;
an unspecified inner term is not being suppressed.
Equivalently, set \(m_0=a+b\), \(m_i=d_i\) for \(i=1,2\),
and \(w=-b\omega_0\). Then
\[
       N(v_i)=[w+m_i\omega_i,v_i].
                                                        \tag{A.35.5}
\]
All nodes are separating, so the degree-one monodromy
derivation vanishes and \(N\) has degree two. Filling the
chosen cusp kills every change of tangential packet, leaving
the proper point section used in this formula.

\paragraph{An unbounded kernel calculation in the surface quotient.}
We prove, in the actual quotient (A.35.3), that
\[
                    \ker N=\mathbf Q_\ell\omega_0.
                                                        \tag{A.35.6}
\]
Here is a normal form that makes the quotient issue explicit.
Present its enveloping algebra using \(u_1,u_2,x_i,y_i\),
put \(u_0=-u_1-u_2\), and impose the three relations
\[
                    x_i y_i=y_i x_i+u_i,\quad i=0,1,2.
                                                        \tag{A.35.7}
\]
The leading words \(x_i y_i\) have no self-overlaps or mutual
overlaps. Ordering the \(x_i\)'s above the \(y_i\)'s and the
\(u_j\)'s gives terminating reduction (give \(u_j\) degree two).
Consequently the words avoiding all three leading words
form a basis. In particular
\(B=\mathbf Q_\ell\langle u_1,u_2\rangle\) embeds.

Group these normal words by the ordered sequence of their
tail letters \(x_i,y_i\). Between consecutive tail letters
are slots containing words in \(B\). If consecutive tail
letters form \(x_i y_i\), their separating slot must be
nonempty; the other slots are unrestricted. The derivation
inserts only \(u\)-letters, so it preserves every such
normal-word component. On \(B\) it is \(\operatorname{ad}w\).
In a component with at least one tail letter, let its first
tail letter have index \(i\). On the first \(B\)-slot the
operator that increases that slot's degree is exactly
\[
                      L_w+m_iR_{u_i}.
                                                        \tag{A.35.8}
\]
For instance, \(N(f)=[w,f]\) and the left part of
\(N(x_i)=[w+m_i u_i,x_i]\) together give
\(wf+fm_i u_i\). Contributions from later slots do not
increase the first-slot degree.

The operator (A.35.8) is injective. For \(i=0\) it is
\(-bL_{u_0}+(a+b)R_{u_0}\). Choose \(u_0\) and a transverse
linear letter as free generators of \(B\). On pure powers
of \(u_0\), the coefficient is \(a\ne0\). Among words
containing the transverse letter, choose the greatest
initial \(u_0\)-prefix length: left multiplication increases
it, right multiplication does not, and \(-b\ne0\).
For \(i=1,2\), the linear letters \(w,u_i\) are independent.
Count occurrences of \(w\); the left summand increases
their number, while the right summand does not.
Taking the greatest first-slot degree in any finite
linear combination now proves that a component with a
tail letter has zero kernel.

Thus the associative kernel lies in \(B\), where it is
the centralizer of the linear letter \(w\), namely
\(\mathbf Q_\ell[w]\). The elements \(u_1,u_2\) are
primitive, so \(B\) is an embedded Hopf algebra. In
characteristic zero the only primitive polynomials in
one primitive letter are its scalar multiples. Intersecting
with \(\mathfrak s\) proves (A.35.6). This argument does not
pass a free-algebra kernel through the surface relation.
The certificate \texttt{code/a\_Q8\_proper\_fixed.py}
checks the absence of overlap ambiguities and the slot
matrices through degree nine for positive sample lengths;
the degree-independent proof is (A.35.8).

\paragraph{A separated integral filtration.}
For completeness, we use an integral augmentation filtration,
so no rational Lie element is silently lifted to a group
element. Write \(P_\ell=\pi_1(C_{\bar k},P)^{(\ell)}\).
Its pro-\(\ell\) presentation has six generators and the
single relation \(\prod_i[\alpha_i,\beta_i]=1\).
In \(\mathbf Z_\ell\langle\!\langle X_i,Y_i\rangle\!\rangle\),
impose the exact relation
\[
        \prod_i[1+X_i,1+Y_i]-1=0.
                                                        \tag{A.35.9}
\]
Its initial quadratic part is \(\sum_i[X_i,Y_i]\), up to
one harmless sign convention. Orient the coefficient-one
term \(X_0Y_0\) as a reduction rule. This word has no
self-overlap. At each finite degree cutoff, reduction
terminates: higher-degree terms decrease the truncated
order, and at a fixed degree the remaining quadratic
terms are lexicographically smaller. Disjoint reductions
commute; there is no overlap ambiguity. The compatible
normal-word bases therefore give a complete quotient
algebra \(\mathcal A\), free over \(\mathbf Z_\ell\) in
each graded degree, with
\[
 \operatorname{gr}\mathcal A=
 \mathbf Z_\ell\langle X_i,Y_i\rangle/
                       \left(\sum_i[X_i,Y_i]\right).
                                                        \tag{A.35.10}
\]
This is also the primitive one-relator situation of
\cite[Theorem 2 and its corollary, pp.\ 142--143]{fghLabute1967};
the direct monic reduction establishes what is needed here
without restricting the prime.

The map \(P_\ell\to\mathcal A^\times\) is injective.
For this continuous group map, give the power-series algebra
its compact \((\ell,X_i,Y_i)\)-adic topology. Its relevant
unit subgroup is an inverse limit of finite \(\ell\)-groups.
This compact topology is distinct from the degree filtration
used to extract first terms.
Indeed every finite \(\ell\)-group quotient \(Q\) permits
evaluation in \(\mathbf Z_\ell[Q]\). Its augmentation ideal
is topologically nilpotent for the \(\ell\)-adic topology,
since the augmentation ideal in \(\mathbf F_\ell[Q]\) is
nilpotent: some power is contained in
\(\ell\mathbf Z_\ell[Q]\), and iterating gives convergence.
Hence all the displayed series converge under
evaluation, and a nonidentity image in \(Q\) is detected
in its group algebra. Finite \(\ell\)-quotients separate
the points of \(P_\ell\).
The induced augmentation filtration on \(P_\ell\) is
therefore separated. Every first nonzero class remains
nonzero after tensoring with \(\mathbf Q_\ell\), and is
primitive for the associated graded Hopf coproduct.
Indeed the closed relation ideal is a Hopf ideal, since
the word set equal to one in (A.35.9) is group-like.
The primitive subspace of
\(\operatorname{gr}\mathcal A\otimes\mathbf Q_\ell
 =U(\mathfrak s)\) is exactly \(\mathfrak s\).
Thus (A.35.6) applies to the first class of every actual
nonidentity pro-\(\ell\) element.

\paragraph{From the local kernel to the global fixed group.}
Enlarge the number field to \(L\), with
\(G_L\cap\Delta_k\subset\Delta_0\), so that the stable model,
the components, and the point are defined at a place over
two. Put \(\Gamma=[G_L,G_L]\), with closure understood.
Then \(\Gamma\subset\Delta_0\), and
\[
                       H=P_\ell^{\bar s_P(\Gamma)}
\]
is a closed subgroup preserved by \(G_L\).
For \(\ell\ne2\), choose a local tame inertia element with
nonzero \(\ell\)-adic tame parameter. The actual commutator
of that element with a local Frobenius belongs to \(\Gamma\).
Its first derivation is a nonzero multiple of \(N\):
the multiplier is the tame parameter times \(q-1\), where
\(q>1\) is the residue cardinality. This is the tame
Frobenius relation applied to the based monodromy action.

By (A.35.6), every nonidentity element of \(H\) has first
augmentation degree two, on the line
\(\mathbf Q_\ell\omega_0\). Hence \(H\) is contained in
degree two and meets degree three trivially. Its map to
the degree-two augmentation quotient is an injective
homomorphism with a closed image on a one-dimensional
rational line in a finite free \(\mathbf Z_\ell\)-module.
If nonzero, that image is isomorphic to \(\mathbf Z_\ell\).
Consequently \(H\) is either trivial or procyclic.

Suppose it is nontrivial, with topological generator \(h\).
Its arithmetic conjugation character \(\theta\) occurs
in degree two of the proper curve's homology-generated
filtration. At every good Frobenius it therefore has
weight two. We spell out the incompatible weight-one
detector. Choose a finite \(\ell\)-group quotient \(Q\)
in which \(h\) is nonidentity, and let \(P'\) be the
inverse image of the cyclic subgroup generated by its
image. The homomorphism
\[
              P'\longrightarrow\langle\bar h\rangle
                      \longrightarrow\mathbf Z/\ell
\]
is nonzero on \(h\). Thus \(h\) has nonzero image in
the integral homology of \(P'\), which is torsion-free:
\(P'\) is the pro-\(\ell\) group of a finite proper
\'etale cover, by the induced topology for covers with
\(\ell\)-group Galois closure.

After replacing \(L\) by a finite extension, its arithmetic
section preserves \(P'\). There are only finitely many
subgroups of its fixed index in the finitely generated
pro-\(\ell\) group, so this stabilizer is open. The cover
and this action are then defined over that number field.
The nonzero rational line generated by \(h\) in its
proper homology carries the same character \(\theta\).
It has weight one. At a prime good for both curves, and
different from \(\ell\), the complex absolute values
would have to be both \(q\) and \(q^{1/2}\), a contradiction.
Therefore \(H=1\), proving (A.35.2).

The punctured-to-proper map is equivariant, and its kernel
is the closed normal subgroup generated by the filled
cusp inertias. This proves the final assertion and also
identifies why this proper-image argument alone does not settle
the full punctured fixed-subgroup problem; A.36 supplies that step.
\end{proof}

\subsection{The quaternion cover: the complete fixed subgroup}

\begin{statement}{A.36}{Proved}{All quaternion-cover fixed elements are in the selected cusp inertia}
For the quaternion first cover \(C_Q\setminus S\) of A.4,
every prime \(\ell\), every packet section induced as in
Part A, and every open subgroup \(\Delta_0\) preserving the
chosen cover and cusp lift, the fixed subgroup of the actual
geometric pro-\(\ell\) group is exactly the chosen cusp inertia:
\[
       \left(\pi_1((C_Q\setminus S)_{\bar K})^{(\ell)}
                          \right)^{s_b(\Delta_0)}
                              =I_{P,\ell}.              \tag{A.36.1}
\]
At odd primes the proof applies more generally to every
cusp packet on this punctured quartic over a finite number
field of definition, with its own derived Galois subgroup.
No nonzero Kummer condition on the packet is required.

The pointed stable reduction at \(2\) has two rational
components each carrying exactly two of the four cusps
and meeting the rest of the fibre at one node. For the
cyclic-eight map of A.35, the cusps above \(B=0,\infty\)
form one such pair; the two cusps above \(B=1\) form the
other. This pairing and the actual first ordinary
monodromy derivation, rather than injectivity of the
whole covering homomorphism on pro-\(\ell\) groups,
give the required comparison.
\end{statement}

\begin{proof}
\paragraph{The four-marked rational tree.}
We use the common-node specialization proved in A.35.
The original \(Q_8\) action preserves both components
adjacent to that node. On the metric edge of the unpointed
stable model, it fixes both ends, hence acts identically.
The four cusps have the same retraction point on this
edge, since they form a single \(Q_8\)-orbit. This point
is interior: the cusps specialize to the node, not to a
smooth point of either adjacent component.
After a finite extension, subdivide at that point.
The resulting rational bridge \(R\) has two attachment
points fixed by \(Q_8\).

Every element of the induced \(Q_8\) action on \(R\)
has order a power of two. A nonidentity automorphism
of \(\mathbf P^1\) of two-power order in characteristic
two is unipotent and has exactly one geometric fixed
point. It cannot fix both attachments. Thus the action
on \(R\) is trivial. Transitivity now implies that all
four marked sections meet the same smooth point of
\(R\) before their marked separation.
Their separation therefore attaches a rooted rational
tree to \(R\), with four marked points and no positive
genus component in this branch.

We record the short stability argument determining this
tree. Its root is fixed by \(Q_8\), so transitivity of
the four marks forces their marked vertices to have
the same distance from the root. A marked vertex
cannot have a descendant: such a descendant subtree
has a marked leaf, at strictly greater distance.
Hence all marked vertices are leaves. Every leaf has
at least two marks, by stability, while every unmarked
internal vertex has at least two outgoing edges.
With four marks, there are just two possibilities:
one component with four marks, or an unmarked root
component with two children, each carrying two marks.
There are no intervening unmarked vertices of valence
two in the stable tree.

The first possibility is excluded by the cyclic-eight
automorphism \(\tau\) from A.35. It preserves the root,
fixes the two cusps over \(B=0,\infty\), and exchanges
the other two. On a single four-mark component it would
fix the attaching point and two distinct marked points.
An automorphism of \(\mathbf P^1\) fixing three points
is the identity, contradicting that exchange. We
therefore obtain the two two-mark leaves.
The two \(\tau\)-fixed cusps lie on the same leaf:
if they lay on different leaves, \(\tau\) would preserve
both leaves and fix each leaf's other marked point as
well. This again contradicts the exchange.
Each marked leaf is consequently a three-punctured
rational component.

The metric statements used here have their elementary
local interpretation in \(uv=\varpi^e\): retraction
records \(v(u(P))\) on the interval \([0,e]\).
An automorphism preserving the two branches preserves
this coordinate, and equal coordinates produce the
same inserted bridge. Thus no choice of a complex
picture or of a conjectural reduction criterion
enters the marked-tree argument.

\paragraph{The first ordinary derivation and its packet ambiguity.}
Fix a cusp \(P\), and call the other cusp on its leaf
\(Q\); call the remaining pair \(R_1,R_2\).
Let \(r_1>0\) be the edge length from the \(P,Q\) leaf
to the unmarked root of this rational branch, and
\(r_2>0\) the length from that root to the other marked
leaf. All other edges lead toward the genus components.

The geometric pro-\(\ell\) group is free of rank nine.
Use six handle generators and three independent cusp
generators, with first homology classes
\[
         v_1,\ldots,v_6,\quad p,q,r.
\]
Here \(p,q,r\) correspond to \(P,Q,R_1\); the remaining
cusp has first class \(-p-q-r\). Put \(s=p+q\).
We use the \emph{ordinary} lower central degree in this
paragraph, so all nine displayed generators have degree
one. The first derivation of the edge twists, with a
base point on the \(P,Q\) leaf, is
\[
 \begin{split}
 D_0p&=D_0q=0,\\
 D_0v_j&=-r_1[s,v_j]\quad(1\le j\le6),\\
 D_0r&=-(r_1+r_2)[s,r].
 \end{split}                                             \tag{A.36.2}
\]
Indeed the first edge encloses all generators outside
the base leaf and its boundary has first class
\(-p-q=-s\). The edge toward the other marked leaf
has first boundary class \(r+(-p-q-r)=-s\).
Edges enclosing only genus components have boundary
in the commutator subgroup and contribute first in
ordinary degree two, not in (A.36.2). This verifies both
coefficients and their signs from the based edge-twist
formula of \cite[Definition 3.1.5 and Theorem 3.1.6]{fghBettsDogra}.

There is a small base-point issue worth making explicit:
that theorem is stated for ordinary points. Choose an
ordinary local point in the punctured residue disc of
\(P\), and blow up its cusp half-edge until it extends
as a smooth point. The theorem applies directly there.
The added segment along that half-edge contributes
only an inner multiple of \(\operatorname{ad}p\).
Comparison with a tangential fibre takes place inside
the punctured formal disc at \(P\). Its geometric
prime-to-two covers are the tame Kummer covers
\(z\mapsto z^{\ell^n}\); the comparison path's cocycle
therefore lies in the actual cusp inertia \(I_{P,\ell}\).
Changing to any completed tangent packet also changes
the action by an inner cusp cocycle. Thus the genuine
degree-one derivation has the form
\[
                   D_\lambda=D_0+\lambda\operatorname{ad}p,
                  \qquad \lambda\in\mathbf Q_\ell.       \tag{A.36.3}
\]
No value of \(\lambda\), and no nonvanishing assertion
about that value, is assumed.

Enlarge the global field to \(L\), so that the curve,
marked model, selected packet, and subsequent covering
map are defined, and so that
\(\Gamma=[G_L,G_L]\subset\Delta_0\).
The commutator of local Frobenius at two with a tame
inertia element having nonzero \(\ell\)-adic tame
parameter belongs to \(\Gamma\). Its first derivation
is a nonzero common multiple of (A.36.3).
Here \(\ell\ne2\); semistability makes the prime-to-two
inertia action tame, and the Frobenius--tame relation
multiplies the parameter by \(q_{\mathrm{res}}-1\ne0\).
After rescaling the derivation, (A.36.3) applies
with some \(\lambda\).

\paragraph{An unbounded free-Lie kernel.}
Inside the free associative algebra on \(p,s,v_1,\ldots,v_6,r\),
let \(B=\mathbf Q_\ell\langle p,s\rangle\).
The derivation on \(B\) is \(\lambda\operatorname{ad}p\).
On any of the seven other letters \(z\), it has the form
\[
                  D_\lambda z=[\lambda p+c_zs,z],
                 \qquad c_z\ne0.                        \tag{A.36.4}
\]
Group words by their ordered sequence of these seven
letters, leaving words in \(B\) in the intervening slots.
The component with an empty such sequence is \(B\).
For any other component, take its first outside letter
\(z\). On its first \(B\)-slot the operator increasing
that slot's degree is
\[
                      \lambda L_p+c_zR_s.                \tag{A.36.5}
\]
All later contributions preserve that first-slot degree.
If \(\lambda=0\), (A.36.5) is injective by right
multiplication by the free letter \(s\).
If \(\lambda\ne0\), count occurrences of \(p\):
the left summand raises their number while the right
summand does not. It is again injective.
The greatest first-slot degree in a finite expression
therefore shows that there is no kernel outside \(B\).
Within \(B\) the kernel is all of \(B\) for
\(\lambda=0\), and is \(\mathbf Q_\ell[p]\) otherwise.
Intersecting with the primitive free Lie algebra gives
\[
 \ker D_\lambda=
 \begin{cases}
    \mathbb L(p,q),&\lambda=0,\\
    \mathbf Q_\ell p,&\lambda\ne0 .
 \end{cases}                                             \tag{A.36.6}
\]
This proof applies in every degree and to arbitrary
nonzero values of the seven coefficients in (A.36.4).

\paragraph{The actual target map removes the remaining two-letter kernel.}
The \(Q_8\) deck action is transitive on the marked pairs.
Compose the cyclic-eight map of A.35 with a deck
automorphism so that \(P,Q\) map respectively to
\(0,\infty\). After enlarging \(L\) if needed this is
an actual morphism of punctured curves
\[
          f:C\setminus S\longrightarrow
                    \mathbf P^1\setminus\{0,1,\infty\}.
                                                        \tag{A.36.7}
\]
Both selected cusp ramification indices are eight.
On first homology, \(p,q\) map to \(8X,8Z\), where
\(X,Z\) are the classes at \(0,\infty\). They are a
basis of the target rational first homology.
Consequently the induced map on the free Lie algebra
\(\mathbb L(p,q)\) is injective in every degree.
This asserts injectivity only on that explicitly
identified leading algebra.

Let \(g\) be an actual element fixed by the specified
\(\Delta_0\). Its image under (A.36.7) is fixed by
\(\Gamma=\Delta_L\) in the induced cusp packet.
Statement A.34 puts this image in the actual inertia
at zero. Since eight is a unit in \(\mathbf Z_\ell\),
there is an actual integral power \(c_P^a\) of the
selected cusp generator with the same image.
Replace \(g\) by \(h=g c_P^{-a}\).
This element is still fixed by \(\Delta_0\) and maps
to one under \(f\).

The ordinary lower central filtration of a finite-rank
free pro-\(\ell\) group is separated and has torsion-free
free-Lie graded pieces over \(\mathbf Z_\ell\).
If \(h\ne1\), its first nonzero class survives rationally
and is killed by the genuine first derivation
\(D_\lambda\). By (A.36.6) that class lies in
\(\mathbb L(p,q)\), including when \(\lambda\ne0\).
But its image in the target graded Lie algebra must
vanish because \(f(h)=1\). This contradicts the
injectivity just proved. Thus \(h=1\) and
\(g\in I_{P,\ell}\). The opposite inclusion holds
because \(\Delta_0\) has cyclotomic character one.

This proves (A.36.1) for every odd prime, with no
integral Malcev lifting. At \(\ell=2\), A.22 applies
to the original \(Q_8\) Galois closure and gives the
same conclusion. The finite certificate
\texttt{code/a\_Q8\_cusp\_fixed.py} checks the degree-one
edge coefficients, the rank-two inertia map, and the
kernel dimensions in small degrees; the unbounded
argument is (A.36.5).
\end{proof}

\subsection{The three dihedral genus-three covers: the complete fixed subgroup}

\begin{statement}{A.37}{Proved}{All first dihedral genus-three covers have exactly cusp fixed subgroups}
For each of the three curves
\[
 C=C_{4,a,b}:v^2=w^8+aw^4+b,
 \qquad (a,b)\in\{(0,-1),(6,1),(-6,1)\},
\]
let \(S\) be the four cusps of A.4. For every prime \(\ell\),
every packet section induced as in Part A, and every open subgroup
\(\Delta_0\) preserving the selected cover and cusp lift \(P\),
\[
 \left(\pi_1((C\setminus S)_{\bar K})^{(\ell)}\right)^{s_b(\Delta_0)}
       =I_{P,\ell}.                                      \tag{A.37.1}
\]
At odd primes the proof works for every cusp packet over a finite
number field of definition, with that field's own derived Galois
subgroup. No nonzero Kummer condition on the packet is imposed.

The pointed stable reduction at two has exactly two rational
components carrying the cusps. Each carries two cusps and meets
the rest of the fibre at a single node. The other components can
have different configurations for the three curves. The assertion
uses only these two marked leaves and does not identify the whole
pro-\(\ell\) group with its image under a covering map.
\end{statement}

\begin{proof}
\paragraph{A harmonic cyclic-four map.}
Choose \(s^2=b\) and put \(T=(v-s)/w^4\). Direct substitution gives
\[
 w^4=\frac{a-2sT}{T^2-1},\qquad v=s+Tw^4.               \tag{A.37.2}
\]
Thus \(\mu(w,v)=(iw,v)\) generates a cyclic-four quotient
to \(\mathbf P^1_T\), totally ramified at precisely the four
cusps, whose images are
\[
                 a/(2s),\quad\infty,\quad 1,\quad-1.
                                                               \tag{A.37.3}
\]
Every nonidentity power of \(\mu\), including
\(\sigma=\mu^2\), fixes precisely these four points:
on the affine chart its fixed points must have \(w=0\),
and at infinity the two branches \(v/w^4=\pm1\)
are individually fixed. There are no additional generic
ramification points to be marked.
For \(a=0\) this is the harmonic set \(\{0,\infty,1,-1\}\).
For \((a,s)=(6,1)\), the substitution \(z=(T-1)/2\)
identifies it with the same set; for \((a,s)=(-6,1)\), use
\(z=(T+1)/2\). In every case its cross-ratios are
\(-1,2,1/2\), so this four-pointed projective line has bad
marked reduction at two, after every finite extension.
Indeed four distinct reductions, after normalizing three
of them to \(0,1,\infty\), would make both the fourth
coordinate and its difference from one units. None of
these three cross-ratios has that property at two.

The original \(D_8\) deck group from A.4 is generated by
\[
 \alpha(w,v)=(iw,-v),\qquad
 j(w,v)=(c/w,-c^2v/w^4),\qquad c^4=b.
\]
It acts transitively on the four cusps and commutes with
\(\sigma=\alpha^2=\mu^2\). The quotient by \(\mu\) has
genus zero, while the quotient by \(\sigma\) has genus one.
Consequently their invariant subspaces on proper rational
\(\ell\)-adic homology have dimensions zero and two,
respectively. These dimensions persist in a compact-type
stable fibre.

\paragraph{The possible unpointed special fibres.}
All proper factors of the Jacobians in A.4 have integral
\(j\)-invariants reducing to zero at two: the relevant values
are \(1728,8000,10976,128\). After a finite extension they
therefore have good supersingular reduction. The Jacobian of
\(C\) has good reduction with \(2\)-rank zero. Its stable
fibre is of compact type, and each positive-genus component
has \(2\)-rank zero. An unmarked stable genus-three tree
can consequently have one genus-three component, one
genus-two and one elliptic component, or three elliptic
components. In the last case its only shapes are a chain
of three elliptic components or three elliptic tails
attached to a rational centre. There are no additional
rational vertices: stability gives each such vertex
valence at least three, and the tree has at most three
positive-genus vertices.

We need a short characteristic-two restriction. Every
nontrivial two-power automorphism of a smooth curve of
\(2\)-rank zero has exactly one geometric fixed point;
all nontrivial elements of a two-group have the same
fixed point. This is the Deuring--Shafarevich argument in
\cite[Lemma 3.2 and Remark 3.1]{fghGiuliettiKorchmaros}.
For a faithful cyclic-four action there is therefore
one ramification point. If \(u_1,u_2\) are its upper
jumps, then
\[
 u_1>0\text{ is odd},\qquad
 u_2=2u_1\quad\text{or}\quad u_2>2u_1\text{ is odd},
                                                               \tag{A.37.4}
\]
by \cite[Theorem 1.1(c)]{fghObusPries}.
The lower jumps are \(u_1,2u_2-u_1\), and the different
is \(u_1+2u_2+3\). Riemann--Hurwitz gives
\[
      g=4g_{\rm quot}+\frac{u_1+2u_2-3}{2}.             \tag{A.37.5}
\]
In particular genus three is impossible: a positive quotient
genus gives \(g\ge5\), while genus zero would require
\(u_1+2u_2=9\), contrary to (A.37.4).
For genus two the only possibility is
\(g_{\rm quot}=0,(u_1,u_2)=(1,3)\); the quotient by
the unique involution also has genus zero, since its
different is six.

In the smooth genus-three case the specialization of
\(\sigma\) is nonidentity: its homology fixed space
has dimension two, rather than six. Thus \(\mu\)
specializes to a faithful order-four automorphism, and
(A.37.5) rules this case out.

If the fibre has genera two and one, \(\mu\) preserves
both components. It must have order four on the genus-two
component: otherwise \(\sigma\) would act identically
there and have at least four homology invariants.
The preceding genus-two calculation shows that
\(\sigma\) has no invariants on that component.
It is therefore the identity on the elliptic component,
where \(\mu\) is nontrivial of order two. Each component's
unique \(\mu\)-fixed point is their joining node.
All four \(\mu\)-fixed cusp sections specialize to that node.

Suppose the fibre is an elliptic chain. If \(\mu\)
preserved each component, its action on the middle
component would be nontrivial, because \(\mu\) has no
homology invariants. It would fix two distinct attaching
points, contradicting the fixed-point restriction.
Hence \(\mu\) exchanges the two outer components.
On those two components \(\mu^2\) is nontrivial;
otherwise their exchanged homology spaces would supply
\(\mu\)-invariants. All \(\sigma\)-invariants therefore
come from the middle elliptic component, on which
\(\mu\) is an involution. Its unique fixed point is
smooth, since its two attaching points are exchanged.
All four cusps specialize to this one smooth point.

It remains to consider the three-tail star. If \(\mu\)
exchanges two tails, it acts on the rational centre as
a nonidentity involution. Its unique fixed point there
is the attaching node of the third tail. On that third
tail \(\mu\) is a nontrivial involution, again with
unique fixed point that node. All four cusps specialize
to the same node. If instead \(\mu\) preserves all
tails, it is the identity on the centre, since it fixes
all three attaching points. It acts with order two
on exactly one tail \(E_0\), and order four on the
other two: this is forced by the two-dimensional
\(\sigma\)-invariant space. Its only fixed point on
each tail is the attaching node. All cusp specializations
therefore lie on the rational centre, including possibly
at its nodes. The \(D_8\) action preserves \(E_0\), as
it commutes with \(\sigma\). Its action on the centre
embeds in the permutations of the other two attaching
points, so has order at most two. Transitivity on the
four cusps shows that their specialization set consists
of either one point or two points, with clusters of
four or two cusps respectively.

\paragraph{Separating the marked clusters.}
In every case above there are one or two cusp clusters,
each of even size, and each can be separated in a
rational branch containing no positive-genus component.
For a cluster at a smooth point this is the ordinary
marked resolution. For a cluster at a node, its
stabilizer in \(D_8\) preserves both adjacent components
and hence fixes the two ends of the metric edge.
It acts identically on the edge. Since it is transitive
on the cluster, all its cusp sections have the same
interior retraction point. After subdividing there,
the stabilizer fixes both attachments of the inserted
rational bridge. A nontrivial two-power automorphism
of a projective line in characteristic two has only
one geometric fixed point, so this bridge action is
trivial. The cluster thus meets a common smooth point
on that bridge, and its marked separation attaches a
rational branch containing no genus.

The same rooted-tree stability argument as in A.36
now applies within each cluster. Transitivity forces
all marked vertices to be leaves, each with at least
two marks. A two-point cluster gives a single two-mark
leaf. A four-point cluster gives either one four-mark
leaf or two two-mark leaves.

We exclude the four-mark leaf using the actual cyclic
quotient (A.37.2). The four cusps are precisely all its
ramification points. The stable marked model is therefore
\(\langle\mu\rangle\)-semistable in the sense of
\cite[Definition 3.1(c)]{fghBouwWewers}.
Its quotient is a semistable model of
\(\mathbf P^1_T\), and the four branch sections specialize
to distinct smooth points, by
\cite[Lemma 3.2(i)--(ii)]{fghBouwWewers}.
If a source component carried all four cusps, \(\mu\)
would fix these four points and hence act identically
on that component. Its image under the finite quotient
would be one component containing all four distinct
smooth branch marks. Marked stabilization preserves
that component and contracts every other component:
the target has genus zero and no marks elsewhere.
It would give good reduction of the four-pointed
projective line, contradicting (A.37.3).
This argument uses the semistable quotient theorem
with its marked-ramification hypotheses; it does not
identify the wild scheme quotient with the geometric
special-fibre quotient without justification.

There are consequently precisely two two-mark rational
leaves and no cusps on the remaining components.

\paragraph{The genuine first derivation and its kernel.}
Fix the selected leaf with cusps \(P,Q\), and denote
the other leaf's cusps by \(R_1,R_2\). Let
\(p,q,r\) be the classes of \(P,Q,R_1\) in ordinary
homology, put \(s_0=p+q\), and choose six handle
generators \(v_1,\ldots,v_6\). The last cusp class
is \(-s_0-r\). For a base point on the selected leaf,
the based edge-twist formula gives
\[
 \begin{split}
 D_0p&=D_0q=0,\\
 D_0v_j&=-a_j[s_0,v_j],\qquad a_j>0,\\
 D_0r&=-L[s_0,r],\qquad L>0.
 \end{split}
                                                        \tag{A.37.6}
\]
Here \(L\) is the length of the path between the two
marked leaves. For a handle on a positive-genus
component, \(a_j\) is the distance from the selected
leaf to the attachment of that component's path to
this leaf-to-leaf path. Each \(a_j\) is positive
because the selected component is a genus-zero leaf.
Every edge on this path has first boundary class
\(-s_0\), while edges off the path enclose only
genus and have zero first boundary class. This proves
(A.37.6) directly from
\cite[Definition 3.1.5 and Theorem 3.1.6]{fghBettsDogra}.

The ordinary-point to tangential-point comparison in
A.36 takes place inside the selected cusp's punctured
formal disc. It changes this genuine first derivation
only by \(\lambda\operatorname{ad}p\), for some
\(\lambda\in\mathbf Q_\ell\). This includes every
completed tangent packet. Choose a finite global field
\(L_0\) with \(\Gamma=\overline{[G_{L_0},G_{L_0}]}\subset\Delta_0\).
An actual commutator of local Frobenius at two with
tame inertia has first derivation a nonzero multiple
of
\[
                 D_\lambda=D_0+\lambda\operatorname{ad}p.
                                                               \tag{A.37.7}
\]
As in A.36, no value of \(\lambda\) is assumed.

The all-degree first-slot proof in A.36 applies
without change: on the free subalgebra generated by
\(p,s_0\), the derivation is \(\lambda\operatorname{ad}p\);
on every outside letter \(z\), it is
\([\lambda p+c_zs_0,z]\), with \(c_z\ne0\).
The first-slot operator \(\lambda L_p+c_zR_{s_0}\)
is injective. Hence, inside the ordinary free Lie algebra,
\[
                  \ker D_\lambda\subset\mathbb L(p,q).
                                                               \tag{A.37.8}
\]
This is an unbounded statement, including \(\lambda=0\).

\paragraph{The harmonic target detects the remaining leading algebra.}
After a projective change of \(T\), the four branch
points are \(\{0,\infty,1,-1\}\). The two maps
\[
                   z\longmapsto z^2,
          \qquad z\longmapsto\left(\frac{z-1}{z+1}\right)^2
                                                               \tag{A.37.9}
\]
are degree-two maps to the tripod, each having exactly
these four preimages of \(\{0,1,\infty\}\).
The first folds \(\{1,-1\}\); the second folds
\(\{0,\infty\}\). Thus at least one sends \(P,Q\)
to distinct target cusps. Compose it with (A.37.2).
The result is an actual degree-eight etale morphism
of punctured curves, with selected cusp indices
\(e_P,e_Q\in\{4,8\}\).
On first homology it sends \(p,q\) to nonzero
multiples of two distinct tripod cusp classes.
It is therefore injective on \(\mathbb L(p,q)\)
in every degree over \(\mathbf Q_\ell\).

If \(g\) is an actual \(\Delta_0\)-fixed element,
A.34 puts its image in the target's selected cusp
inertia. At odd \(\ell\), \(e_P\) is a unit in
\(\mathbf Z_\ell\), so subtract an actual integral
power of the selected cusp generator to obtain an
element \(h\) which is still fixed and maps to one.
If \(h\ne1\), its first nonzero ordinary lower-central
class survives rationally, is killed by (A.37.7),
and hence lies in (A.37.8). Its image cannot vanish
under the target map, a contradiction. Thus \(h=1\).
The opposite inclusion follows from trivial cyclotomic
action on \(\Delta_0\). Finally, \(\ell=2\) is A.22,
since each original \(D_8\) cover has two-group
Galois closure. This proves (A.37.1).

The executed exact certificate \texttt{code/a\_C4\_fixed.py}
checks (A.37.2), all three harmonic normalizations,
the four branch indices, the cyclic-four jump calculation,
and the homology determinants for all six possible
choices of a two-cusp leaf. Its output is retained
in the adjacent \texttt{.json} and \texttt{.out} files.
\end{proof}

\subsection{A cyclotomic criterion and the remaining symmetric covers}

\begin{statement}{A.38}{Proved}{The nonstable symmetric genus-three fixed groups at every prime}
For \(a=6,-6\), put
\[
 C_a:v^2=w^7+aw^4+w,\qquad
 X_a=C_a\setminus\{P_0,P_\infty\}.
\]
For every prime \(\ell\), every original packet of Part A,
and every permitted open subgroup \(\Delta_0\) preserving
the cover and the selected cusp lift \(P\), one has
\[
 \left(\pi_1(X_{a,\overline K})^{(\ell)}\right)^{s_b(\Delta_0)}
                =I_{P,\ell}.                         \tag{A.38.1}
\]
At odd primes the argument also applies to every cusp
packet over a finite number field of definition, with
that field's own derived Galois subgroup. It includes
the packet with zero first Johnson coefficient.
No nonzero primitive Collino regulator is required.

More generally, let \(X/k\) be a hyperbolic affine curve
over a number field with at least two geometric cusps,
let \(s\) be a packet at a rational cusp, and write
\(F=\pi_1(X_{\bar k})^{(\ell)}\),
\(H=F^{s(\overline{[G_k,G_k]})}\).
Give its rational unipotent completion the weight
filtration in which proper homology has weight one
and boundary generators have weight two. Suppose the
Lie algebra of the unipotent closure of \(H\) has
zero odd graded pieces, and \(G_k\) acts on its
weight \(2m\) piece by the scalar \(\chi_\ell^m\).
Then \(H\) is exactly the selected cusp inertia.
This criterion concerns the closure of an actual
subgroup; it does not assert that a rational fixed
Lie element lifts to \(F\).
\end{statement}

\begin{proof}
\paragraph{The arithmetic criterion.}
We first prove the general assertion, keeping the
given field and the given actual subgroup \(H\).
The actual packet makes \(G_k\) act on \(F\), and
\(H\) is stable under this action. On \(H\) the
action factors through \(G_k^{\rm ab}\), since the
closed commutator subgroup fixes every element of
\(H\). This remains true on every algebraic unipotent
closure of its image.

Let \(\mathfrak h\) denote the inverse limit of the
Lie algebras of those closures in the finite weight
quotients of the rational unipotent completion of
\(F\). Images of homomorphisms of unipotent algebraic
groups in characteristic zero are closed. Thus the
transition maps between these closures are surjective,
and their Lie algebras are the finite quotients of
\(\mathfrak h\) with its induced weight filtration.
All these quotients are finite dimensional.

Choose a finite place of residue cardinality \(q\ne1\)
and residue characteristic different from \(\ell\),
and choose any arithmetic Frobenius lift \(\varphi\).
We use only \(\chi_\ell(\varphi)=q\). On a finite
filtered quotient of \(\mathfrak h\), the action of
\(\varphi\) on weight \(2m\) is the scalar \(q^m\).
These scalars are distinct for distinct \(m\).
A filtered linear operator which is scalar on every
graded piece with distinct such scalars is semisimple:
its minimal polynomial divides the product of the
distinct factors \(T-q^m\), as follows successively
through the filtration. Its eigenspace for \(q^m\)
maps isomorphically to the weight \(2m\) graded piece.

Now let \(\gamma\in\ker\chi_\ell\). It commutes
with \(\varphi\) on \(\mathfrak h\), because the
action there is abelian. It induces the identity on
every associated graded piece. It therefore fixes
each \(\varphi\)-eigenspace pointwise, and acts as
the identity on every finite quotient of
\(\mathfrak h\). Exponentiation gives the same
conclusion on its unipotent groups.

The map from the finite-rank free pro-\(\ell\)
group \(F\) to its rational unipotent completion is
injective. Indeed its integral Magnus expansion is
injective, its ordinary lower-central graded pieces
are torsion free, and the weighted filtration is
separated: assigning only weights one and two makes
it comparable with the ordinary degree filtration.
We conclude for the actual group that
\[
                     H\subset F^{s(\ker\chi_\ell)}.
                                                            \tag{A.38.2}
\]
This uses rational completions only to establish an
identity on existing integral elements.

We next prove that every element of the right-hand
side of (A.38.2) is peripheral. Suppose \(h\) were
not peripheral. The purely group-theoretic detection
lemma of \cite[Lemma 2.1.4, printed pp.96--97]{Nakamura1994},
applied to the full class of finite \(\ell\)-groups,
gives an open subgroup \(U\subset F\) containing
\(h\) in which its image in proper homology modulo
\(\ell\) is nonzero. More explicitly,
\[
 h\notin\overline{[U,U]U^\ell
                  \langle Z\cap U\rangle},          \tag{A.38.3}
\]
where \(Z\) is the union of all conjugates of all
cusp inertia groups. Killing \(Z\cap U\) is exactly
proper compactification of the corresponding cover.
The proper integral homology is a free
\(\mathbf Z_\ell\)-module, so (A.38.3) also gives
a nonzero rational homology vector.

There are finitely many subgroups of the given index
in the finitely generated group \(F\). After a finite
extension \(k'/k\), its actual packet action preserves
\(U\). The group \(U\rtimes s(G_{k'})\) is then
an actual arithmetic finite-cover subgroup. The
corresponding geometric cover has \(\ell\)-group
Galois closure; the core argument of A.22 identifies
its full geometric pro-\(\ell\) group with \(U\).
Let \(J\) be the Jacobian of its smooth proper
compactification. The vector just found belongs to
\(V_\ell J\) and, by (A.38.2), is fixed by
\(G_{k'(\mu_{\ell^\infty})}\).

This contradicts \cite[Theorem 3, printed p.315,
proof pp.317--318]{A_RibetCyclotomic1981}.
That theorem gives
\[
       (V_\ell J)^{G_{k'(\mu_\infty)}}=0             \tag{A.38.4}
\]
for every abelian variety over every number field
and every prime. The subgroup fixing all roots of
unity is contained in the subgroup fixing the
\(\ell\)-power roots, so (A.38.4) gives the required
vanishing as well. All hypotheses concern the
Jacobian of this actual finite cover; no claim is
made about homology of an infinite cover.

Thus every nonidentity element of \(H\) is peripheral.
Since \(X\) has at least two geometric cusps, every
cusp has a nonzero primitive class in ordinary
geometric pro-\(\ell\) abelianization. Every nonzero
power of an inertia generator therefore has nonzero
ordinary homology. Every commutator of two elements
of \(H\), on the other hand, has zero ordinary
homology. It must be the identity. Hence \(H\) is
abelian. It contains the selected inertia \(I\),
because the cyclotomic character is trivial on
the closed commutator subgroup. The centralizer
of a nontrivial element of this primitive free
factor is \(I\) itself, by
\cite[Lemma 2.1.2]{Nakamura1994}. Consequently
\(H=I\), proving the criterion.

\paragraph{The weighted monodromy for the two symmetric curves.}
We apply this criterion to the curves in the statement
and odd \(\ell\). Choose a finite field of definition
\(k\) so that
\(\Gamma=\overline{[G_k,G_k]}\subset\Delta_0\).
All marked and deck data of A.33 are defined there.
It is enough to compute the \(\Gamma\)-fixed group:
it contains the desired \(\Delta_0\)-fixed group,
and \(\Delta_0\) already fixes the selected inertia.

A.33 establishes the actual pointed reduction at
two: a rational centre bearing the two cusps, with
three elliptic tails. Let \(t\) denote the logarithm
of the selected primitive cusp generator, and let
\(a_i,b_i\) be the six handle generators attached
to these tails. In the weight-graded rational Lie
algebra put
\[
 \operatorname{wt}(a_i)=\operatorname{wt}(b_i)=1,
 \quad \operatorname{wt}(t)=2,
 \quad \omega_i=[a_i,b_i],\quad
 \Omega=\omega_1+\omega_2+\omega_3.                  \tag{A.38.5}
\]
It is the free weighted Lie algebra on the handle
space and \(t\). This filtration is distinct from
the ordinary lower-central filtration used in A.33.

For an ordinary base point on the rational centre,
the first nonzero weighted monodromy derivation is
\(D_0v_i=m_i[\omega_i,v_i]\), \(D_0t=0\), with
all \(m_i\ne0\). This is the actual based edge
formula of A.33 and Betts--Dogra. The edges enclose
only their indicated genus-one components. Their
boundary logarithms have first weight two term
\(\omega_i\); higher terms cannot change this
degree-two derivation.

Compare the ordinary point with the selected tangent
through its punctured formal residue disc, as in A.36.
The comparison is inner by the selected cusp group.
Equivalently, at weight two an inner discrepancy
fixing \(t\) has conjugator in the degree-two
centralizer of \(t\), which is exactly
\(\mathbf Q_\ell t\). There is no weight-one
inner discrepancy fixing \(t\). Thus for every
packet the genuine derivation of an actual local
Frobenius--tame-inertia commutator in \(\Gamma\)
is a nonzero common multiple of
\[
 \begin{split}
 D_\lambda(t)&=0,\\
 D_\lambda(v_i)&=[m_i\omega_i+\lambda t,v_i]
       \quad(v_i=a_i,b_i),
 \qquad\lambda\in\mathbf Q_\ell.
 \end{split}
                                                        \tag{A.38.6}
\]
Both summands now have weight shift two. In
particular the zero first ordinary Johnson case
does not require a separate nonvanishing argument.

\paragraph{The unbounded weighted kernel.}
Inside the free associative enveloping algebra set
\(B=\mathbf Q_\ell\langle t,\omega_1,\omega_2,
\omega_3\rangle\). These are freely independent
associative generators: introduce symbols \(u_i\)
and orient
\[
                  a_i b_i\longmapsto b_i a_i+u_i.
\]
The leading words have no overlaps, and the normal
words give an injective copy of the free algebra
on \(t,u_1,u_2,u_3\). This is the same normal-form
construction as A.27, with the additional letter
\(t\).

Group normal words by their ordered sequence of
handle letters, and place the words in \(B\)
in the intervening slots. On \(B\), (A.38.6)
is \(\lambda\operatorname{ad}t\). For a nonempty
handle sequence beginning in the \(i\)-th block,
the operator increasing its first slot's length is
\[
                       \lambda L_t+m_iR_{u_i}.       \tag{A.38.7}
\]
It is injective. For \(\lambda=0\) this is right
multiplication by a free letter; for
\(\lambda\ne0\), the greatest number of occurrences
of \(t\) in a finite input is increased only by the
left summand. All other parts of the derivation
preserve that first-slot length. A component of
greatest first-slot length in a finite polynomial
therefore cannot lie in the kernel. There is no
kernel outside \(B\).

Within \(B\) the kernel is all of \(B\) when
\(\lambda=0\), and \(\mathbf Q_\ell[t]\) otherwise.
Intersecting with primitive elements gives, in
particular,
\[
 \ker D_\lambda\subset
               \mathbb L(t,\omega_1,\omega_2,\omega_3).
                                                            \tag{A.38.8}
\]
This is a proof in every weight. It applies to
the associated graded of the unipotent closure
of actual \(\Gamma\)-fixed elements, because
the commutator defining (A.38.6) fixes that closure
pointwise.

The non-CM invariant intersection of A.33 is unchanged
in the weighted grading. We recall the exact linear
map to make this point explicit. Put
\(A=\sum a_i\), \(B_1=\sum b_i\), and
\(\delta a_i=a_i-A/3\). The Zariski closure of
\(\Gamma\) on the two copies of the non-CM elliptic
factor contains the symplectic transformation
\[
 g(a_i)=a_i,\qquad g(b_i)=b_i+\delta a_i,\qquad g(t)=t.
\]
For \(U=\langle\omega_1,\omega_2,\omega_3\rangle\)
one computes
\[
 g(\omega_i)=\omega_i+\tfrac13[A,\delta a_i],
 \qquad U\cap gU=\mathbf Q_\ell\Omega.               \tag{A.38.9}
\]
The second identity follows by comparing the
\(a_i\wedge a_j\) components: the sum-zero vector
\(\sum d_i\delta a_i\) must commute with \(A\),
and hence vanish, so all \(d_i\) are equal.
In associative words, split at each occurrence
of \(t\), and pair consecutive handle letters
in each intervening block. Tensor-space intersection
in each pair replaces \(U\) and \(gU\) by their
intersection. Taking primitives yields
\[
 \mathbb L(t,U)\cap\mathbb L(t,gU)
                       =\mathbb L(t,\Omega).         \tag{A.38.10}
\]
Both \(t\) and \(\Omega\) have weight two.

The full arithmetic action on these two generators
is \(\chi_\ell\): for \(t\) this is actual cusp
inertia, and for \(\Omega\) it is the symplectic
similitude character on proper homology. It follows
from (A.38.8)--(A.38.10) that every weight \(2m\)
graded piece of the fixed subgroup's unipotent
closure has scalar character \(\chi_\ell^m\),
and every odd piece vanishes. The criterion proved
above applies and gives exactly \(I_{P,\ell}\).

At \(\ell=2\), A.33 already gives the assertion
by the actual unramified double-cover comparison.
All steps at odd primes are over \(\mathbf Q_\ell\)
only for detecting identities on existing elements;
the final cusp group and equality are integral.
This includes \(\ell=3\), despite the harmless
rational denominators in (A.38.9).

The executed certificate \path{code/a_cyclotomic_fixed.py}
checks the first-slot operators over \(\mathbf Q\)
for the four boundary letters through slot length
four, and checks a filtered Frobenius matrix with
four two-dimensional scalar graded pieces and
nonzero extension blocks. Its commutation equations
have no nonzero strictly filtration-increasing
solution. These finite checks support the preceding
unbounded proof; they are not used to realize
rational Lie classes as actual group elements.
\end{proof}

\subsection{The dihedral genus-two quotients: the complete fixed subgroup}

\begin{statement}{A.39}{Proved}{Every dihedral genus-two first quotient has exactly cusp fixed subgroups}
Let
\[
 (a,b)\in\{(0,-1),(6,1),(-6,1)\},\qquad c^4=b,
\]
and let
\[
 D_{a,c}:z^2=(x^2-4c)(x^4-4cx^2+a+2c^2),
 \qquad Y_{a,c}=D_{a,c}\setminus\{Q_+,Q_-\},
\]
where \(Q_+,Q_-\) are its two points at infinity.
These include both reflection conjugacy classes for each
of the three \(D_8\) covers in A.4. For every prime
\(\ell\), every original packet of Part A, and every
permitted open subgroup \(\Delta_0\) preserving the
curve and its selected cusp lift \(P\),
\[
 \left(\pi_1(Y_{a,c,\bar K})^{(\ell)}\right)^{s_b(\Delta_0)}
                      =I_{P,\ell}.                    \tag{A.39.1}
\]
At odd primes the same argument works for every cusp
packet over a finite number field of definition, with
that field's own derived Galois subgroup.

The proof uses the actual pointed reduction at two,
an unbounded weighted monodromy kernel, and the
cyclotomic criterion of A.38. It does not transfer an
embedding through the unramified double cover
\(C_4\to D_{a,c}\) at an odd prime, nor require a
primitive Collino regulator for the conjugate cusps.
\end{statement}

\begin{proof}
\paragraph{The two cusps meet in the unpointed stable model.}
The proper Jacobian is isogenous to \(E\times F\),
where the complementary \(j\)-invariant is \(8000\),
\(10976\), or \(128\), by A.4. Together with
\(j(E)=1728\), all have positive valuation at two.
After a finite extension the Jacobian has good
supersingular reduction. The unpointed stable fibre
is of compact type and has \(2\)-rank zero.
A stable unmarked compact-type genus-two curve is
either a smooth genus-two curve or two smooth
elliptic curves joined at a node. It has no rational
components: a rational vertex in a stable tree needs
at least three branches, each eventually containing
positive genus, which would give total genus at
least three.

Write \(\delta=[Q_+-Q_-]\). A.4 proves
\(4\delta=0\), in fact exact order four, for every
\(a,c\) under consideration. Only the displayed
annihilation is needed here. A useful uniform
identity behind it is
\[
 \begin{aligned}
 A&=x^3-4cx,\qquad d=a+2c^2,\qquad
 z^2-A^2=d(x^2-4c),\\
 f&=\frac{(z-A)^2}{x^2-4c}=d\frac{z-A}{z+A},\qquad
 \operatorname{div}(f)=4(Q_+-Q_-).
 \end{aligned}                                      \tag{A.39.2}
\]
Here \(d\ne0\) in all the indicated cases.
Since the geometric special Jacobian has \(2\)-rank
zero, it has no nonzero geometric points of order
a power of two. The specialization of \(\delta\)
is therefore zero.

Take the Abel map defined by either cusp as its
base point. By
\cite[Theorem 1.1]{fghDupuyRabinoff}, it extends to
the compact-type model and is a closed immersion
because the unpointed special fibre has no rational
components. Thus \(Q_+\) and \(Q_-\) specialize to
the same point of that fibre. For the two-component
case this can also be read directly from
\cite[Proposition 3.2]{fghDupuyRabinoff}: the special
Abel map embeds the two elliptic components as the
two Jacobian axes, meeting only at the origin.
In particular the assertion includes specialization
to a node, not just to the smooth locus.

If the unpointed fibre is smooth, separating the two
sections attaches a rational component carrying both
cusps and a single node. The edge to the genus-two
component has strictly positive length. Contracting
unmarked rational subdivisions merely adds their
lengths.

If the unpointed fibre is \(E_1\cup E_2\), consider
the hyperelliptic involution \(h(x,z)=(x,-z)\).
It exchanges the two cusps and acts as \(-1\) on
proper homology. It cannot exchange the elliptic
components, since an involution exchanging two
two-dimensional homology spaces has a two-dimensional
invariant space. Thus it preserves both components
and acts as their elliptic involution with origin
the common node. Each supersingular elliptic
component has only this one geometric fixed point.
The common specialization of \(Q_+,Q_-\), being
fixed by \(h\), must consequently be the node.

The involution preserves the two branches of that
node, hence fixes the metric edge pointwise. Its
exchange of the two cusp sections shows that their
interior retraction points agree. Subdivide at that
point after a finite extension. On the inserted
rational bridge, \(h\) fixes the two attachments.
In characteristic two a nonidentity involution of
\(\mathbf P^1\) has only one geometric fixed point,
so this bridge action is the identity. The two
sections therefore still meet at the same smooth
point of the bridge. Separating them attaches a
rational two-mark leaf, with no genus in that
branch. In the stable pointed model the other end
of its attaching edge is the rational bridge
joining the two elliptic components. Denote the
leaf edge length by \(r>0\), and the two arm lengths
by \(m_1,m_2>0\). Their positivity follows because
the common retraction is interior to the original
node edge and the cusp sections collide on its
inserted bridge.

\paragraph{Actual weighted monodromy.}
For odd \(\ell\), give the free rational Lie algebra
of the punctured genus-two curve its weight filtration:
the four proper handle generators have weight one,
and the logarithm \(t\) of the selected cusp has
weight two. Write
\[
 \omega_1=[a_1,b_1],\quad\omega_2=[a_2,b_2],
 \qquad\Omega=\omega_1+\omega_2.
\]
Use an ordinary base point reducing to the rational
two-mark leaf. Its other cusp and the node form the
other two boundaries of that rational tripod.
The node encloses all proper genus and has first
weight two boundary class \(\Omega\), up to the
common orientation convention. The based affine
edge-twist formula in
\cite[Definition 3.1.5 and Theorem 3.1.6]{fghBettsDogra}
therefore gives the following first nonzero weighted
derivation. In the smooth genus-two case it is
\[
                 Nt=0,\qquad Nv=r[\Omega,v]
                  \quad(v=a_1,b_1,a_2,b_2),\quad r>0.
                                                               \tag{A.39.3}
\]
In the two-elliptic case it is
\[
             Nt=0,\qquad
             Nv_i=[r\Omega+m_i\omega_i,v_i]
                  \quad(v_i=a_i,b_i),\quad r,m_i>0.
                                                               \tag{A.39.4}
\]
The common leaf edge contributes \(r\Omega\) to
both handle blocks, while each arm contributes its
own \(m_i\omega_i\). All these terms have weight
shift two. There is no weight-one edge term.

The ordinary-point comparison with the chosen
tangent takes place through the selected cusp's
punctured formal residue disc, exactly as in A.36
and A.38. Moving the ordinary point out along the
cusp half-edge contributes only a cusp inner term.
The comparison cocycle and every completed tangent
packet change lie in the actual selected inertia.
Thus the genuine derivation is
\[
                         D_\lambda=N+\lambda\operatorname{ad}t,
                   \qquad\lambda\in\mathbf Q_\ell.   \tag{A.39.5}
\]
Equivalently, a weight-two inner discrepancy fixing
\(t\) has conjugator in its centralizer
\(\mathbf Q_\ell t\); no weight-one discrepancy
can fix \(t\).

Choose a finite global field \(k\) defining the data
and realizing the indicated local semistable model,
with \(\Gamma=\overline{[G_k,G_k]}\subset\Delta_0\).
An actual local Frobenius--tame-inertia commutator
at two belongs to \(\Gamma\), and its first weighted
derivation is a nonzero common multiple of (A.39.5).
Indeed its tame parameter is multiplied by
\(q_{\rm res}-1\ne0\). The common factor can be
divided out in the rational kernel calculation.
No nonvanishing condition on \(\lambda\) is assumed.

\paragraph{The smooth-component all-degree kernel.}
First use (A.39.3). The leading associative word of
\(\Omega\), in a degree-lexicographic order with
\(a_1>b_1>a_2>b_2\), is \(a_1b_1\), with coefficient
one and no proper self-overlap. The triangular
bimodule construction of A.27 therefore applies to
\(A_V=\mathbf Q_\ell\langle a_1,b_1,a_2,b_2\rangle\)
and \(\Omega\). Specifically, let \(\mathcal C\)
be all nonempty handle words which neither begin
nor end with \(a_1b_1\). Stripping initial and
terminal powers of that word gives the basis
\[
 \{\Omega^n:n\ge0\}\ \sqcup\
 \{\Omega^u w\Omega^v:w\in\mathcal C,\ u,v\ge0\}.
                                                               \tag{A.39.6}
\]
Each element has its distinct indicated leading
word, so this is a basis in every finite homogeneous
degree. Adjoin the free letter \(t\), and set
\(B=\mathbf Q_\ell\langle t,\Omega\rangle\).
A word outside \(B\) has a first nonboundary handle
core in (A.39.6), with an arbitrary \(B\)-prefix.
On that prefix the part of \(D_\lambda\) increasing
its length is
\[
                          \lambda L_t+rR_\Omega.
                                                               \tag{A.39.7}
\]
It is injective: if \(\lambda=0\), use right
multiplication by \(\Omega\); otherwise count
occurrences of \(t\). Subsequent terms cannot
remove the first core and preserve its prefix
length. The kernel is therefore contained in
\(B\). There the derivation is
\(\lambda\operatorname{ad}t\), so its associative
kernel is \(B\) for \(\lambda=0\) and
\(\mathbf Q_\ell[t]\) otherwise.

\paragraph{The two-component all-degree kernel.}
For (A.39.4), use the two separate handle-core
bases of A.27 and put
\(B=\mathbf Q_\ell\langle t,\omega_1,\omega_2\rangle\).
On this freely generated algebra,
\[
 N_Bt=0,\qquad N_B\omega_i=r[\Omega,\omega_i],
 \qquad D_\lambda|_B=N_B+\lambda\operatorname{ad}t.
                                                               \tag{A.39.8}
\]
For a first core of type \(i\), the part increasing
the length of its \(B\)-prefix is the operator
\[
 K_i=N_B+\lambda L_t+R_{r\Omega+m_i\omega_i}.
                                                               \tag{A.39.9}
\]
This follows by combining \(D_\lambda\) on the
prefix with the left half of
\([r\Omega+m_i\omega_i+\lambda t,w]\) on the
core; the right \(\lambda t\) terms cancel.
All other terms preserve the prefix length and
the first core.

If \(\lambda\ne0\), \(N_B\) and both right
multiplications preserve the number of \(t\)'s,
while \(\lambda L_t\) raises it. The largest
\(t\)-count proves injectivity of \(K_i\).
If \(\lambda=0\), use the free generating system
\(\Omega,\omega_i,t\) of \(B\).
In this system \(N_B\Omega=N_Bt=0\) and
\(N_B\omega_i=r[\Omega,\omega_i]\), so \(N_B\)
preserves \(\omega_i\)-count. Right multiplication
by \(r\Omega\) also preserves it, whereas
\(m_iR_{\omega_i}\) raises it. Since \(m_i\ne0\),
\(K_i\) is again injective. Taking a greatest
prefix length proves that no element outside
\(B\) belongs to the kernel.

It remains to compute the kernel in \(B\), where
\(\Omega,t,\omega_1\) are free generators.
On \(B_0=\mathbf Q_\ell\langle t,\Omega\rangle\)
the derivation is \(\lambda\operatorname{ad}t\),
and
\[
                 D_\lambda\omega_1
                           =[r\Omega+\lambda t,\omega_1].
\]
The first-slot argument of A.36 now has operator
\(\lambda L_t+rR_\Omega\), which is injective.
Thus the associative kernel is \(B_0\) if
\(\lambda=0\) and \(\mathbf Q_\ell[t]\) otherwise.

In both geometric cases, all displayed boundary
generators are primitive in the tensor Hopf
algebra. Intersecting their associative kernels
with primitives gives the exact free-Lie formula
\[
 \ker D_\lambda=
 \begin{cases}
     \mathbb L_{\mathbf Q_\ell}(t,\Omega),&\lambda=0,\\
     \mathbf Q_\ell t,&\lambda\ne0.
 \end{cases}                                           \tag{A.39.10}
\]
This holds in all weights; both \(t\) and
\(\Omega\) have weight two.

\paragraph{The actual fixed subgroup.}
Let \(H\) be the actual \(\Gamma\)-fixed subgroup
of the geometric pro-\(\ell\) group. Its unipotent
closure is pointwise \(\Gamma\)-fixed. Every
leading weighted Lie class in that closure is
therefore killed by the actual derivation (A.39.5),
and hence belongs to (A.39.10). The full \(G_k\)
action on \(t\) is \(\chi_\ell\), by the selected
cuspidal packet; its action on \(\Omega\) is the
same character, by the symplectic similitude on
proper homology. The closure has no odd weight
pieces, and its weight \(2m\) piece has scalar
character \(\chi_\ell^m\).

The criterion of A.38 applies, giving \(H=I_{P,\ell}\).
Since \(\Gamma\subset\Delta_0\) and the original
\(\Delta_0\) fixes the selected inertia, this proves
(A.39.1) for every odd prime and every stated packet.
At \(\ell=2\), A.22 applies to each reflection
subgroup, whose Galois closure is the original
\(D_8\) cover. This gives the remaining prime.

The exact certificate \texttt{code/a\_D4\_fixed.py}
checks (A.39.2), the two-component derivation
identities on \(B\), the triangular basis used for
a smooth genus-two component, and finite weighted
consistency checks for both geometric cases and
several packet scalars. The proof of (A.39.10)
is the unbounded argument above, not an inference
from those finite checks.
\end{proof}

% A.40: independently audited and registered by the parent agent.
\subsection{All genus-one covers}

\begin{statement}{A.40}{Proved}{Every genus-one cover has exactly cusp fixed groups and cusp normalizers}
Let \(Y_H\to U\) be any finite connected geometric cover
whose smooth compactification has genus one. After a finite
field extension and a choice of one cusp as origin, its
compactification map is an isogeny
\[
                \varphi:E'\longrightarrow E,
       \qquad S=\ker\varphi,
       \qquad Y_H=E'\setminus S.
\]
For every prime \(\ell\), every original packet, and every
permitted open subgroup \(\Delta_0\) preserving the cover
and its selected cusp lift \(P\),
\[
       \left(\pi_1(Y_{H,\bar K})^{(\ell)}\right)^{s_b(\Delta_0)}
                              =I_{P,\ell}.             \tag{A.40.1}
\]
Every nontrivial closed procyclic subgroup of this pro-\(\ell\)
group whose arithmetic normalizer has open Galois image is
contained in a conjugate of a cusp inertia group.
Thus Task A.5 holds for every genus-one compactification,
with no restriction on the isogeny degree.

The fixed-group proof includes the following more general
geometric criterion. Suppose a smooth proper curve of positive
genus has potentially good reduction at a finite place of
residue characteristic \(p\), and two distinct cusps, one of
them the selected cusp, have the same specialization in its
good model. For \(\ell\ne p\), the leading weighted Lie
classes of its actual derived-fixed subgroup lie in the free
Lie algebra generated by the boundary classes and the total
proper symplectic class. Consequently that fixed subgroup is
the selected inertia, by A.38. Other cusps may specialize at
different points; their locations do not affect the assertion.
\end{statement}

\begin{proof}
\paragraph{A single positive-genus component and its marked trees.}
First prove the geometric criterion. Enlarge the finite number
field \(k\) so all cusps and a pointed semistable model at the
chosen place are defined, the proper curve has good reduction,
and \(\Gamma=\overline{[G_k,G_k]}\) lies in the specified
\(\Delta_0\), when one is prescribed. Let \(P\ne Q\) be
the two cusps with equal specialization.

Separating the marked sections in the smooth proper model,
and then stabilizing, produces one positive-genus component
\(C\), with rational marked trees attached at its smooth
points. Every finite edge has a strictly positive rational
length. The cusps \(P,Q\) lie in the same attached tree.
This follows directly by blowing up the finitely many common
specializations of sections: a blowup introduces a rational
exceptional component and preserves the tree structure;
contracting an unstable rational component adds the lengths
of its adjacent edges. In particular the paths from \(P\)
to \(C\) and to \(Q\) separate strictly before reaching
\(C\).

Root this metric tree at the component carrying \(P\).
Write \(e\) for its finite edges, \(m_e>0\) for their
lengths, and orient each edge away from the root. A cusp
already on \(C\) is simply assigned to that vertex. Let
\(E_e\), \(Q_e\), and \(j_e\) be the indicators that
\(C\), \(Q\), and cusp \(j\), respectively, lie on the
side of \(e\) away from the root.

Use handle generators \(a_i,b_i\), \(1\le i\le g\),
and cusp generators \(t_j\), of weights one and two
respectively. Set
\[
                  \Omega=\sum_{i=1}^g[a_i,b_i].
\]
The first peripheral relation is
\(\Omega+\sum_jt_j=0\). Eliminate \(t_Q\), keeping
the selected \(t_P\) as an independent generator. The
weighted graded Lie algebra is the free Lie algebra on
the handles and these remaining cusp generators.

For an ordinary base point in the selected cusp disc,
the actual based edge-twist formula of
\cite[Definition 3.1.5 and Theorem 3.1.6]{fghBettsDogra}
gives the first weighted monodromy derivation. Put
\[
 \delta_e=E_e\Omega+\sum_{j:\,j_e=1}t_j,
 \qquad
 W=\sum_e m_eE_e\delta_e,
 \qquad W_j=\sum_e m_ej_e\delta_e.                    \tag{A.40.2}
\]
With a common orientation convention it has
\[
         Nv=[W,v]\quad\text{for every handle }v,
             \qquad Nt_j=[W_j,t_j].                 \tag{A.40.3}
\]
Changing the common orientation negates the entire
derivation and has no effect on its kernel. All
edge boundaries have weight two. The comparison with
the chosen tangent takes place in its actual cusp
disc. As in A.36 and A.39, an arbitrary packet changes
the genuine first derivation only by
\(\lambda\operatorname{ad}(t_P)\),
\(\lambda\in\mathbf Q_\ell\). Thus replace
\(W,W_j\) by \(W+\lambda t_P,W_j+\lambda t_P\).
This change leaves every coefficient of \(\Omega\)
below unchanged. There is no weight-one discrepancy:
such an inner discrepancy would centralize \(t_P\)
in the free Lie algebra, whose centralizer there
contains no weight-one element.

After elimination of \(Q\), the coefficient of
\(\Omega\) in \(\delta_e\) is \(E_e-Q_e\).
Consequently the coefficient in \(W\), denoted \(r\),
and that in \(W_j\), denoted \(c_j\), are
\[
 r=\sum_e m_eE_e(1-Q_e),\qquad
 c_j=\sum_e m_ej_e(E_e-Q_e).                         \tag{A.40.4}
\]
The path from the branching point of \(P,Q,C\) to
\(C\) has positive length, and its edges are precisely
those counted by \(r\). Hence \(r>0\). Some \(c_j\)
can be negative or equal to \(r\); no sign or strict
inequality for them will be used. In the special case
where all cusps lie in one attached tree, one also has
\[
 r-c_j=\sum_e m_e\bigl(E_e(1-Q_e)(1-j_e)
                   +(1-E_e)Q_ej_e\bigr)>0,          \tag{A.40.5}
\]
since the last edge incident to \(C\) contributes.

An actual local Frobenius--tame-inertia commutator in
\(\Gamma\) has first derivation a nonzero common
multiple of (A.40.3) with this packet correction.
The multiplier is \(q_{\rm res}-1\ne0\); after the
semistable extension, the prime-to-\(p\) monodromy is
the actual edge action just used. Thus it suffices to
bound the kernel of this genuine rational derivation.

\paragraph{The unbounded prefix calculation.}
Let \(A\) be the free associative algebra on the handles
and retained cusp letters, and put
\[
               B=\mathbf Q_\ell\langle\Omega,t_j:
                                                j\ne Q\rangle.
\]
It is a free algebra on the displayed generators.
For completeness, the leading word of \(\Omega\)
is \(a_1b_1\), with coefficient one and no proper
self-overlap, after choosing an appropriate order.
The handle algebra has the triangular basis
\[
 \{\Omega^n:n\ge0\}\ \sqcup
 \{\Omega^u w\Omega^v:u,v\ge0,\quad
         w\ne\varnothing\text{ begins and ends with no }a_1b_1\}.
                                                               \tag{A.40.6}
\]
Indeed, strip initial and terminal powers of
\(a_1b_1\) from each handle word; the indicated
leading words are distinct and exhaust all words.
Adjoining the cusp letters by free product gives a
normal basis whose words outside \(B\) have a first
nonboundary handle core with an arbitrary \(B\)-prefix.

Let \(D_B\) denote the restriction of the corrected
derivation to \(B\). On every handle word the
derivation is the commutator with the same \(W\),
including its packet correction. On the prefix of
the first core, its portion increasing boundary-word
length is therefore
\[
                            K=D_B+R_W.              \tag{A.40.7}
\]
All terms acting after that core preserve its prefix
length. They cannot remove the core: left or right
multiplication by \(\Omega\) remains in its displayed
bimodule, and a cusp letter merely adds a free-product
separator. Thus injectivity of \(K\) excludes all
elements outside \(B\) from \(\ker D\).

Filter \(B\) by the number of occurrences of
\(\Omega\). Since
\(D_B\Omega=[W,\Omega]\), this term does not raise
that count. The part of \(K\) which raises it by one
is
\[
  K^+=D_c+rR_\Omega,
  \qquad D_c\Omega=0,
  \qquad D_ct_j=c_j[\Omega,t_j].                    \tag{A.40.8}
\]
All remaining terms, including every packet term,
preserve the count.

Fix the ordered cusp word \(i_1\cdots i_n\). Its
span of insertions
\[
          \Omega^{a_0}t_{i_1}\Omega^{a_1}
                         \cdots t_{i_n}\Omega^{a_n}
\]
identifies linearly with the polynomial ring
\(\mathbf Q_\ell[Z_0,\ldots,Z_n]\), by sending
the insertion to \(Z_0^{a_0}\cdots Z_n^{a_n}\).
Under this identification \(K^+\) is multiplication
by
\[
 c_{i_1}Z_0+
       \sum_{h=1}^{n-1}(c_{i_{h+1}}-c_{i_h})Z_h
                       +(r-c_{i_n})Z_n.             \tag{A.40.9}
\]
Its coefficients sum to \(r\ne0\), so this
polynomial is nonzero, regardless of the signs of
the \(c_j\). For \(n=0\), the multiplier is simply
\(rZ_0\). A polynomial ring over \(\mathbf Q_\ell\)
is a domain, so \(K^+\) is injective on every one
of these mutually independent cusp-word summands.
It follows that \(K\) is injective by taking the
largest \(\Omega\)-count in a nonzero polynomial.
Taking a largest first-core prefix length now proves
\[
                \ker D\subset B\cap\mathbb L
                    =\mathbb L(t_j:j\ne Q,\Omega).  \tag{A.40.10}
\]
The last equality is the primitive-element identity
for the tensor Hopf algebra on these freely generated
primitive elements. The argument works in every
weight; it is not a finite-degree dimension count.

Every leading weighted class of the unipotent closure
of the actual \(\Gamma\)-fixed subgroup is killed
by \(D\), hence lies in (A.40.10). Each retained
cusp class has character \(\chi_\ell\), and so does
\(\Omega\), by the proper symplectic similitude.
All its generators have weight two. Thus the closure
is Tate-graded in the precise sense of A.38, which
gives the selected inertia as the actual fixed group.
This proves the geometric criterion.

\paragraph{A collision for every prime dividing an isogeny degree.}
Return to \(\varphi:E'\to E\), and let
\(N=\deg\varphi>1\). We claim that for every
prime \(p\mid N\), some place above \(p\), after
a finite extension, has a nonzero element of
\(S[p^\infty]\) specializing to zero in the good
model of \(E'\).

Both curves acquire good reduction
by \cite[Corollary VII.7.2]{fghSilvermanAEC}, and have the
same ordinary or supersingular reduction type after
such an extension: their isogeny extends between
the good models. The degree of this special isogeny
equals its generic degree, since the degree of the
pullback of an ample degree-one line bundle is
constant in a smooth proper family. At \(p=2\),
\(E\) is potentially supersingular, as already used
in A.39. At odd \(p\equiv3\pmod4\), it is
supersingular by
\cite[Theorem V.4.1(a), Example V.4.5]{fghSilvermanAEC}:
the expansion of \((x^3-x)^{(p-1)/2}\) has no
term \(x^{p-1}\). In both cases the geometric
special fibre has no nonzero \(p\)-power torsion,
so every point of \(S[p^\infty]\) specializes to
zero, proving the claim.

Suppose now \(p\equiv1\pmod4\). Write
\(p=\pi\bar\pi\) in \(\mathbf Z[i]\), and let
\(p^n\) be the exact \(p\)-part of \(N\).
The reductions are ordinary. At a place of \(K\)
dividing \(\pi\), the isogeny \([\pi]\) of \(E\)
has degree \(p\) and zero special differential, so
its reduction is purely inseparable. The isogeny
\([\bar\pi]\) has invertible differential and
separable reduction. Therefore, among
\(E[p^n]\), the points specializing to zero are
exactly \(E[\pi^n]\). At the conjugate prime they
are exactly \(E[\bar\pi^n]\). This can also be
seen by decomposing
\(E[p^n]=E[\pi^n]\oplus E[\bar\pi^n]\): the
first summand has purely inseparable reduction and
the second is finite etale at the selected place.

Assume for contradiction that \(S[p^\infty]\)
specializes injectively at every place above \(p\).
Prime-to-\(p\) torsion always specializes injectively
in good reduction. Thus all \(N\) kernel points
remain distinct in the special kernel, and
\(\widetilde\varphi\) is separable of degree \(N\).
The dual isogeny satisfies
\(\widehat\varphi\varphi=[N]\). In ordinary
characteristic \(p\), \([N]\) has inseparable
degree \(p^n\); hence \(\widehat\varphi\) has
inseparable degree \(p^n\) and separable degree
\(N/p^n\). In particular the special kernel of
\(\widehat\varphi\) has no nonzero geometric
\(p\)-primary point. Its generic \(p\)-primary
kernel \(T\subset E[p^n]\), which has order
\(p^n\), therefore specializes to zero at both
CM primes. It follows that
\[
                  T\subset E[\pi^n]\cap E[\bar\pi^n]=0,
\]
because the two ideals are coprime. This contradicts
\(|T|=p^n>1\) and proves the collision claim.
The dual-isogeny and inseparable-degree facts here
are \cite[Theorem III.6.1, Corollary III.6.4,
Theorem V.3.1]{fghSilvermanAEC}; prime-to-\(p\)
specialization is \cite[Proposition VII.3.1(b)]{fghSilvermanAEC}.

Translation by any selected cusp \(P\in S\)
extends to the good model. Thus a nonzero kernel
point specializing to zero gives a distinct cusp
\(Q\in S\) specializing to the same point as
\(P\). The collision claim applies to every
selected cusp, not just the chosen origin.

\paragraph{All coefficient primes and the normalizers.}
Fix \(\ell\). If \(N\) has a prime divisor
\(p\ne\ell\), choose its collision place and
apply the geometric criterion above. If \(N\)
is an \(\ell\)-power, the geometric isogeny
cover has Galois group \(S\), an \(\ell\)-group,
so A.22 gives (A.40.1) directly. Degree one is
the original curve and is covered by R1, with
the same open-subgroup strengthening used in A.22.
This proves (A.40.1) for every isogeny degree,
every coefficient prime, and every stated packet.

For the normalizer assertion, first enlarge the
finite field of definition so the isogeny, the
cusp origin, and all required paths are defined.
The map \(Y_H\to U\) induces a continuous map
of geometrically pro-\(\ell\) arithmetic groups
over the identity of this open Galois group. A
closed procyclic subgroup with open normalizer
therefore has normalized image with open arithmetic
image whenever that image is nontrivial.

If the source subgroup were noncuspidal, A.11
would give it a nonzero class in the proper
\(\mathbf Q_\ell\)-homology of \(E'\). The
isogeny induces an isomorphism
\(V_\ell E'\to V_\ell E\), even when
\(\ell\mid N\). Its image subgroup is thus
nontrivial and has nonzero proper homology.
By R2, independently verified in A.11--A.12,
the target subgroup is cuspidal, which gives
zero proper homology, a contradiction. Restriction
to a finite field leaves the original normalizer
image open, so this proves the stated classification.

Finally every finite map between the genus-one
compactifications is unramified by Riemann--Hurwitz.
Choosing a point above \(O\) as origin makes it
an isogeny, and its removed set is exactly its
kernel. These cover all genus-one cases in the
statement.

The program \texttt{code/a\_isogeny\_fixed.py}
checks the rooted-tree coefficient identity on
all labelled trees of small orders, the exact
slot polynomial formula, and the prefix injectivity
on an explicit branched marked tree, including
negative coefficients and arbitrary sample packet
scalars. The all-degree and all-tree assertions
are proved above, independently of those checks.
\end{proof}

% A.41: full independent parent audit passed; registered after checkpoint six.
% The tensor-intersection and diagonal-conjugacy lemmas also passed DE audit.
\subsection{A compact-type criterion for higher-genus covers}

\begin{statement}{A.41}{Proved}{A rational marked path in compact-type reduction forces cusp fixed groups}
Let \(X/k\) be a smooth affine hyperbolic curve over a number
field, whose smooth compactification has positive genus and
whose complement contains at least two geometric points.
Let \(P\) be a rational cusp and let \(s\) be any packet
at \(P\). Suppose that, after a finite extension at a place
of residue characteristic \(p\), the pointed stable model
has a tree as its dual graph. Suppose also that a distinct
cusp \(Q\) has this property: the unique path between the
vertices carrying \(P\) and \(Q\), including its endpoints,
contains only genus-zero components. Then, for every prime
\(\ell\ne p\) and every open subgroup
\(\Gamma_0\subset\overline{[G_k,G_k]}\),
\[
 \left(\pi_1(X_{\bar k})^{(\ell)}\right)^{s(\Gamma_0)}
                         =I_{P,\ell}.                 \tag{A.41.1}
\]
The same conclusion holds for the permitted open
\(\Delta_0\) in Part A whenever the cover and selected
cusp are preserved and a finite field of definition has
its closed derived subgroup contained in \(\Delta_0\).

A sufficient condition for the geometric hypothesis is
compact-type reduction with a special Jacobian of
\(p\)-rank zero and a \(p\)-primary torsion cusp divisor
\([P-Q]\). More generally it suffices that \(P,Q\) have
equal specialization under the extended Abel--Jacobi map.
This gives a criterion for higher-genus instances of
Task A.5; it makes no assertion for a dual graph with a
cycle, or for a path which contains a positive-genus vertex.
\end{statement}

\begin{proof}
\paragraph{Field restriction and the local graph.}
Enlarge the number field to \(L\) so that all cusps and
components are split, the stated pointed stable model is
defined at a place above \(p\), and the Zariski closure of
the proper \(\ell\)-adic homology representation is
connected. We may also arrange
\(\Gamma=\overline{[G_L,G_L]}\subset\Gamma_0\), or
\(\Gamma\subset\Delta_0\) in Part A. Indeed an open
subgroup of \(G_k\) can be chosen whose intersection
with the closed derived group lies in the prescribed open
subgroup, and its own closed derived group lies in that
intersection. Further finite extensions preserve the graph
hypothesis: edge lengths are rescaled and any inserted
valence-two rational vertices may be contracted. It is
enough to prove that the actual \(\Gamma\)-fixed group
is \(I_{P,\ell}\). The specified larger group already
fixes this inertia, since its cyclotomic character is one.

Let \(C_1,\ldots,C_m\) be the positive-genus components
of the compact-type special fibre. The good reduction of
its Jacobian and its special product decomposition give
an orthogonal symplectic decomposition
\[
 V=H_1(\overline X_{\bar L},\mathbf Q_\ell)
          =\bigoplus_{i=1}^m V_i,
 \qquad \Omega=\sum_{i=1}^m\Omega_i,
 \qquad \Omega_i\in\Lambda^2V_i.                    \tag{A.41.2}
\]
Here \(\Omega_i\) is the nondegenerate symplectic
bivector on the indicated component. This is a local
symplectic decomposition; the individual summands are not
asserted to be globally Galois stable. The compact-type
Jacobian product is recalled in
\cite[Section 1, p.1]{fghDupuyRabinoff}. Equivalently,
(A.41.2) is the homology decomposition obtained by cutting
the surface along the separating nodes in the based
comparison of \cite[Theorem 3.1.6]{fghBettsDogra}.

Root the metric tree at the vertex carrying \(P\), and
orient its finite edges away from that root. Write
\(m_e>0\) for their lengths. For each edge let
\(i_e,Q_e,j_e\) indicate whether component \(C_i\),
the cusp \(Q\), and the cusp \(j\), respectively,
lie on its side away from the root. Use handle generators
of weight one for \(V\), and cusp generators \(t_j\)
of weight two. The first peripheral relation is
\(\Omega+\sum_j t_j=0\). Eliminate \(t_Q\) and
retain \(t_P\) among the independent cusp generators.
The weighted graded geometric Lie algebra is then
\[
                         \mathbb L(V\oplus T),
 \qquad T=\langle t_j:j\ne Q\rangle,
 \quad \operatorname{wt}V=1,\quad\operatorname{wt}T=2.
                                                               \tag{A.41.3}
\]

\paragraph{The genuine weighted derivation.}
Before eliminating \(Q\), put
\[
 \delta_e=\sum_i i_e\Omega_i+\sum_{j:\,j_e=1}t_j,
 \qquad W_i=\sum_e m_ei_e\delta_e,
 \qquad W_j=\sum_e m_ej_e\delta_e.                    \tag{A.41.4}
\]
The actual based edge-twist formula of
\cite[Definition 3.1.5 and Theorem 3.1.6]{fghBettsDogra}
gives the first weighted monodromy derivation as
\[
                  Dv=[W_i,v]\ (v\in V_i),
                  \qquad Dt_j=[W_j,t_j].             \tag{A.41.5}
\]
As in A.36 and A.40, choose the ordinary comparison point
in the actual selected cusp disc. Comparison with the
specified tangential packet changes the first derivation
only by \(\lambda\operatorname{ad}(t_P)\), with
\(\lambda\in\mathbf Q_\ell\). Thus in (A.41.4)--(A.41.5)
replace every \(W_i,W_j\) by \(W_i+\lambda t_P,W_j+
\lambda t_P\). There is no weight-one inner discrepancy:
it would centralize the independent cusp generator
\(t_P\), whose Lie centralizer has no weight-one part.
A common change of edge orientation merely negates \(D\).

After semistable extension, the prime-to-\(p\) monodromy
has this actual edge action. A local Frobenius--tame-inertia
commutator, regarded inside \(\Gamma\), has first
weighted derivation a nonzero common multiple of \(D\),
with factor \(q_{\mathrm{res}}-1\ne0\). Thus the
leading classes of the unipotent closure of actual
\(\Gamma\)-fixed elements are killed by \(D\).
All terms in (A.41.5) shift weight by two, and the cusp
packet term has exactly the same shift.

The coefficient of \(\Omega_i\) in the eliminated
\(\delta_e\) is \(i_e-Q_e\). Its coefficient in
its own \(W_i\) is therefore
\[
               r_i=\sum_e m_ei_e(1-Q_e)>0.           \tag{A.41.6}
\]
The equality describes the distance from the branching
point of \(P,Q,C_i\) to \(C_i\). Since the entire
\(P\)--\(Q\) path has genus zero, \(C_i\) is off
that path and the indicated distance is strictly positive.
This is the only positivity assertion needed. Coefficients
of \(\Omega_i\) in other \(W\)'s may be negative,
zero, or equal to \(r_i\), and \(\lambda\) is arbitrary.

\paragraph{An unbounded calculation with several handle blocks.}
Work over \(F=\mathbf Q_\ell\) in the free associative
enveloping algebra \(A=T(V\oplus T)\), and put
\[
             U=\langle\Omega_1,\ldots,\Omega_m\rangle,
             \qquad B=F\langle T,U\rangle\subset A.
                                                               \tag{A.41.7}
\]
This is a free associative algebra on the displayed
primitive generators. To make its complementary basis
explicit, choose a symplectic basis in each \(V_i\),
with leading word \(a_{i1}b_{i1}\) for \(\Omega_i\).
Its coefficient is one and it has no proper self-overlap.
Each handle algebra has the triangular basis
\[
 \{\Omega_i^n:n\ge0\}\ \sqcup
 \{\Omega_i^u w\Omega_i^v:u,v\ge0,
   \ w\ne\varnothing\text{ has neither initial nor terminal }
                                      a_{i1}b_{i1}\}.           \tag{A.41.8}
\]
Indeed strip the maximal initial and terminal powers of
\(a_{i1}b_{i1}\) from every word; the resulting distinct
leading words exhaust the usual word basis. Taking the
free product of these component algebras and the free
cusp algebra gives a normal basis in which every word
outside \(B\) has a first nonboundary handle core,
of some type \(i\), preceded by an arbitrary \(B\)-word.
Pure boundary powers adjoining that core are retained
in the adjoining slots. A separator from a different
component or a cusp letter cannot disappear.

By (A.41.5), the derivation on every handle word of
component type \(i\) is exactly its commutator with
\(W_i\). Also
\(D\Omega_i=[W_i,\Omega_i]\), so \(D\) preserves
\(B\); write this restriction as \(D_B\). For a
fixed first core of type \(i\), the part increasing
the length of its first boundary slot is
\[
                             K_i=D_B+R_{W_i}.        \tag{A.41.9}
\]
All terms acting after that core preserve the first-slot
length. They cannot remove its core: multiplication by
\(\Omega_i\) at either end stays in its displayed
bimodule, whereas a different boundary letter adds a
separator. The derivative of an existing boundary letter
retains that letter in both commutator terms. Thus the
top first-slot operator acts independently for each fixed
first-core tail, with no contribution from later slots.

To prove \(K_i\) injective, count occurrences of the
single letter \(\Omega_i\) in \(B\). Its part which
raises that count by one is
\[
                K_i^+=D_{c_i}+r_iR_{\Omega_i},
 \quad D_{c_i}\Omega_i=0,
 \quad D_{c_i}u=c_u[\Omega_i,u]\quad(u\ne\Omega_i), 
                                                               \tag{A.41.10}
\]
where \(c_u\) is the coefficient of \(\Omega_i\)
in the \(W\) belonging to the boundary letter \(u\).
For \(u=\Omega_j\) use \(W_j\) for component \(j\);
for a cusp use its cusp \(W\). The self coefficient
in \(D_B\Omega_i=[W_i,\Omega_i]\) vanishes, so
that expression preserves this count. Every omitted
term in \(K_i\) also preserves it.

Fix a word \(u_1\cdots u_n\) in all the boundary
letters other than \(\Omega_i\). Insert powers of
\(\Omega_i\) in its \(n+1\) slots and identify the
resulting vector space with \(F[Z_0,\ldots,Z_n]\).
The operator (A.41.10) becomes multiplication by
\[
 c_{u_1}Z_0+\sum_{h=1}^{n-1}(c_{u_{h+1}}-c_{u_h})Z_h
                           +(r_i-c_{u_n})Z_n.        \tag{A.41.11}
\]
Its coefficients telescope to \(r_i\ne0\). When
\(n=0\), its multiplier is \(r_iZ_0\). Therefore
it is injective on every mutually independent word
summand. Taking the largest \(\Omega_i\)-count of
a finite nonzero input proves that \(K_i\) is
injective. Now take a largest first-slot length among
the normal words of a nonzero polynomial outside
\(B\). Formula (A.41.9) cannot vanish there. Hence,
in every weight,
\[
                   \ker D\subset B\cap\mathbb L(V\oplus T)
                                      =\mathbb L(T\oplus U).
                                                               \tag{A.41.12}
\]
The final equality holds because \(B\) is the tensor
Hopf algebra on the indicated freely generated primitive
elements. This argument proves the assertion in all
weights, and does not infer it from finitely many ranks.

\paragraph{The exact intersection of the global translates.}
Let \(\mathfrak h\) be the Lie algebra of the
unipotent closure of the actual \(\Gamma\)-fixed
subgroup, with its induced weight filtration, as in
A.38. Its leading graded pieces are stable under
\(G_L\) and, by (A.41.12), lie in
\(\mathbb L(T\oplus U)\). On the weighted ambient
grading, \(V\) is the canonical proper homology,
\(T\) is the independent boundary subspace with
character \(\chi_\ell\), and \(\Lambda^2V\) is
the canonical space of brackets of weight-one classes.
In particular global Galois action does not mix \(T\)
with \(\Lambda^2V\) in weight two.

Let \(G\subset\operatorname{GSp}(V)\) be the
connected Zariski closure of that action. Each finite
weighted piece of \(\mathbb L(V\oplus T)\) is an
algebraic representation of \(G\), with the action on
\(T\) given by its similitude character. The same
leading piece of \(\mathfrak h\) therefore lies in
all translates \(\mathbb L(T\oplus gU)\),
\(g\in G\). Set
\[
                           U_* =\bigcap_{g\in G}gU.
                                                               \tag{A.41.13}
\]
The following intersection is an equality of actual
subspaces of the free weighted Lie algebra:
\[
        \bigcap_{g\in G}\mathbb L(T\oplus gU)
                                    =\mathbb L(T\oplus U_*).
                                                               \tag{A.41.14}
\]
Here is a direct verification. The block map
\[
          T\bigl(T\oplus(V\otimes V)\bigr)
                               \longrightarrow T(V\oplus T)
\]
is injective: cusp letters are individually recognizable,
and each maximal string of handle letters in its image
has even length, with a unique consecutive pairing from
its left end. Thus, in weight \(2n\), intersections
can be taken inside the \(n\)-fold tensor powers of
\(T\oplus(V\otimes V)\). For subspaces \(A_1,A_2\)
of a vector space, a basis adapted to their intersection
gives
\(A_1^{\otimes n}\cap A_2^{\otimes n}
=(A_1\cap A_2)^{\otimes n}\). Iterate this identity.
The intersection of any family of finite-dimensional
subspaces is already a finite subintersection, so the
same assertion holds for the family in (A.41.13).
Taking primitive elements gives (A.41.14). No passage
through a surface-group quotient is involved.

\paragraph{Why the remaining component classes are Tate.}
Use the total nondegenerate bivector \(\Omega\) to
identify a bivector \(\beta:V^*\to V\) with
\(A_\beta=\beta\Omega^{-1}\in\operatorname{End}(V)\).
Then \(A_{\Omega_i}\) is the projector onto the
orthogonal block \(V_i\). Consequently \(U\)
identifies with the commutative diagonal algebra
\[
                 \mathcal D=\bigoplus_{i=1}^mF P_i,
          \qquad P_i|_{V_i}=1,\quad P_i|_{V_j}=0\ (j\ne i).
                                                               \tag{A.41.15}
\]
For a symplectic similitude \(g\),
\[
               A_{g\beta}=\chi_\ell(g)\,gA_\beta g^{-1}.
                                                               \tag{A.41.16}
\]
The subspace corresponding to \(U_*\) is conjugation
stable and lies in \(\mathcal D\). For any element
\(A\) in that subspace, its entire \(G\)-conjugacy
orbit is contained in \(\mathcal D\) and has the
same multiset of eigenvalues as \(A\). There are
only finitely many elements of \(\mathcal D\) with
that multiset: each fixed nonzero block must receive
one of the finitely many eigenvalues. This remains
true with repeated eigenvalues or unequal block sizes.
After extension to an algebraic closure the orbit of
the connected group \(G\) is connected, and so this
finite orbit is a singleton. Hence \(G\) centralizes
\(A\). Formula (A.41.16) proves
\[
                         g\beta=\chi_\ell(g)\beta
                      \quad\text{for all }\beta\in U_*.
                                                               \tag{A.41.17}
\]
No reductivity assumption is used, and none of the
individual local projectors \(P_i\) is claimed to
come from a global endomorphism.

Every generator in \(T\oplus U_*\) now has weight
two and character \(\chi_\ell\). Equations
(A.41.12)--(A.41.17) show that \(\mathfrak h\) has
zero odd graded pieces and scalar character
\(\chi_\ell^n\) in weight \(2n\). Apply the
actual-subgroup criterion in A.38. It gives that the
\(\Gamma\)-fixed group is exactly \(I_{P,\ell}\),
which proves (A.41.1) and its stated Part A variant.

\paragraph{A torsion-divisor criterion for the rational path.}
The extended Abel--Jacobi map exists on the pointed
compact-type model by
\cite[Theorem 1.1]{fghDupuyRabinoff}. On any component,
its projection to that component's Jacobian is its
Abel map, up to translation, and all other component
coordinates are constant; this is
\cite[Proposition 3.2, pp.5--6]{fghDupuyRabinoff}.
It follows that two distinct smooth marked points with
the same Abel specialization cannot have a positive-genus
component on their joining path. Indeed, projection to
the Jacobian of such a component gives the difference
of the two points at which the path enters and leaves
it, using the marked endpoint if necessary. Those
points are distinct: two different incident edges have
different nodes, and a marked point is disjoint from
all nodes. The Abel map of a smooth positive-genus
component is injective on points, so this projected
difference is nonzero, a contradiction. Thus their
path contains only rational components.

Finally suppose the special Jacobian has \(p\)-rank
zero and \([P-Q]\) is \(p\)-primary torsion. Its
specialization is a geometric \(p\)-power torsion
point of that abelian variety, hence is zero. The
Abel specializations are therefore equal, giving the
sufficient condition in the statement.

The exact program \texttt{code/a\_compact\_type\_fixed.py}
checks the rooted-tree distance identity, multicore
prefix operators with negative and repeated coefficients,
and tensor intersections for a symplectic mixing matrix.
The all-degree and connected-group arguments are proved
above, independently of those finite certificates.
\end{proof}

% A.42: registered after independent full root and DE audits, both PASS.
% This adds the coefficient-prime-equals-residue-characteristic case to A.41.
\subsection{The residue-characteristic coefficient prime}

\begin{statement}{A.42}{Proved}{The compact-type rational-path criterion holds at every prime}
Under the geometric hypotheses of A.41, its fixed-subgroup
conclusion also holds for \(\ell=p\). Consequently for
every prime \(\ell\), every cusp packet at \(P\), and
every indicated open derived subgroup, the full fixed group
of the actual geometric pro-\(\ell\) fundamental group
is exactly \(I_{P,\ell}\).

The crystalline argument uses the connected global algebraic
monodromy group only after a change of fibre. It does not
regard individual global Galois matrices as morphisms of
local crystalline representations, and it does not apply
Fontaine's monodromy to an individual etale vector.
\end{statement}

\begin{proof}
The assertion for \(\ell\ne p\) is A.41. Put
\(\ell=p\), and make the finite global field restriction
used there, so the marked model is split and semistable,
its graph has the stated rational \(P\)--\(Q\) path,
the proper homology image has connected Zariski closure,
and \(\Gamma=\overline{[G_L,G_L]}\) lies in the
specified open derived subgroup. Write \(F=L_v\) for
the local field and \(F_0\) for its maximal unramified
subfield. Put
\[
 H=\bigl(\pi_1(X_{\bar L})^{(p)}\bigr)^{s(\Gamma)}.
\]
Let \(\mathfrak h\) be the Lie algebra of its actual
unipotent closure, formed before taking finite weight
quotients, exactly as in A.38.

\paragraph{Crystalline leading subobjects killed by local monodromy.}
The general comparison proved in the first two paragraphs
of A.31 applies to this marked semistable curve and every
packet at its selected cusp. It makes every finite rational
unipotent Lie quotient semistable over \(F\), and identifies
its \(D_{\rm st}\) realization with the log-crystalline
one, whose monodromy is the actual edge-twist residue.
Those arguments were stated for an arbitrary residue prime
and marked semistable curve; their comparison and packet
proofs do not use the later special \(S_3\) models.
In particular arbitrary rationalized cusp Kummer twists
remain semistable and contribute a cusp-inner residue.

Use the weight filtration of A.41: proper homology has
weight one and cusp generators have weight two. Its
intrinsic second step is the kernel of the map to proper
abelian homology; subsequent steps are generated by the
bracket and by the weight-two cusp images. The comparison
is compatible with this filtration by its functoriality
for proper compactification, brackets, and the cusp maps.
Thus no crystalline splitting of ordinary open homology
is required.

Every finite quotient of \(\mathfrak h\) is a globally
stable subrepresentation of a finite ambient Lie quotient.
Its \(G_L\)-action factors through \(G_L^{\rm ab}\),
because \(\Gamma\) fixes the actual subgroup pointwise.
It is semistable locally by closure of semistability under
subquotients, and crystalline by the lemma in A.31 that a semistable
representation with abelian image is crystalline. Hence
Fontaine's \(N\) vanishes on its whole \(D_{\rm st}\)
realization.

The proper Jacobian has good reduction, since the special
curve is of compact type. Thus
\(V=H_1(\overline X_{\bar L},\mathbf Q_p)\) is
crystalline at \(v\). The independent boundary space
\(T\) is a sum of \(\mathbf Q_p(1)\)'s and is also
crystalline. Every finite weighted graded piece of the
ambient free Lie algebra, and every subobject
\[
                          W_d=\operatorname{gr}_d\mathfrak h,
\]
is consequently crystalline. Here and below the free Lie
algebra uses \(V\) in weight one and \(T\) in weight two.

In the log-crystalline realization write
\[
 V_{\rm c}=D_{\rm cris}(V)=\bigoplus_i V_{i,\rm c},
 \quad T_{\rm c}=D_{\rm cris}(T),\qquad
 \Omega_{\rm c}=\sum_i\Omega_{i,\rm c}.              \tag{A.42.1}
\]
The blocks come from the positive-genus components of the
special compact-type curve and are orthogonal symplectic.
Choose a basis of the one-dimensional crystalline Tate
fibre to view their polarization classes as bivectors.
The edge formula of A.31, with its node-thickness
coefficients, gives exactly (A.41.4)--(A.41.6) over
\(F_0\), including the same positive distances \(r_i\).
The actual tangent has \(N(t_P)=0\). A possible inner
weight-one discrepancy centralizing \(t_P\) is zero,
and a weight-two discrepancy belongs to
\(F_0t_P\). Thus the first possible weighted residue
is the genuine derivation
\[
                  N_2=D_{\mathrm{graph}}+
                                      \lambda\operatorname{ad}(t_P),
                         \qquad\lambda\in F_0,       \tag{A.42.2}
\]
up to a common nonzero normalization. There is no
weight-zero residue because the weighted generators
are crystalline, and no weight-one outer residue for
the separating graph. The displayed inner comparison
also excludes a weight-one based residue.

Exactness of \(D_{\rm st}\), applied to a sufficiently
large finite quotient of \(\mathfrak h\) and its
induced weight filtration, now gives
\[
                         D_{\rm cris}(W_d)\subset\ker N_2.
                                                               \tag{A.42.3}
\]
More explicitly, lift a vector of the left side to the
weight-\(d\) part of that finite \(D_{\rm st}\) quotient.
Its monodromy is zero; its component in weight \(d+2\)
is \(N_2\) of the leading vector because all smaller
weight shifts vanish. Surjectivity of the finite
unipotent-closure transitions, as proved in A.38, ensures
that the full \(W_d\) is represented in this procedure.

\paragraph{Transport the global group through a genuine fibre functor.}
Let \(\mathcal C\) be the full tensor subcategory of
global \(\mathbf Q_p\)-representations generated by
\(V\) and \(\mathbf Q_p(1)\), closed under subquotients
and duals. Its usual fibre functor \(\omega\) has group
\(G\), the connected global algebraic monodromy group.
The proper polarization makes \(G\) a subgroup of
\(\operatorname{GSp}(V)\), with multiplier the Tate
character. Every object of \(\mathcal C\) is locally
crystalline. Therefore
\[
                    \eta=D_{\rm cris}|_{\mathcal C}
                               :\mathcal C\longrightarrow\operatorname{Vec}_{F_0}
                                                               \tag{A.42.4}
\]
is a faithful exact tensor fibre functor, by
\cite[Theorem 5.2.1 and Proposition 9.1.6]{fghBrinonConrad}.
Its tensor automorphism group \(G_{\rm c}\) is an
inner form of \(G_{F_0}\): the fibre-comparison torsor
\(\operatorname{Isom}^{\otimes}(\omega_{F_0},\eta)\)
is faithfully flat by
\cite[Theorem 3.2, p.31]{fghDeligneMilneTC}.
In particular \(G_{\rm c}\) is connected. It acts on
\(V_{\rm c}\) by symplectic similitudes and on
\(T_{\rm c}\) by its Tate character \(\chi_{\rm c}\).
Each \(W_d\) is an object of \(\mathcal C\), so its
crystalline fibre in (A.42.3) is \(G_{\rm c}\)-stable.

This statement can be checked directly with the period
comparison, avoiding any ambiguity about individual
matrices. Let \(E=\operatorname{Frac}(B_{\rm cris})\).
The functorial crystalline comparison gives tensor
isomorphisms
\[
 c:E\otimes_{F_0}D_{\rm cris}(V)
                       \xrightarrow{\sim}E\otimes_{\mathbf Q_p}V,
                                                               \tag{A.42.5}
\]
and likewise for \(T\), all their tensor/Lie
constructions, and every subobject \(W_d\). Thus
\(c\) identifies the crystalline leading subspace
with \(E\otimes W_d\). Conjugating \(G_E\) by
\(c\) gives the connected comparison group acting
on that same crystalline ambient space. This is not
an assertion that a global matrix is defined on the
\(F_0\)-fibre. The map \(N_2\) in (A.42.3) is
extended \(E\)-linearly for this linear-algebra
calculation; one does not differentiate period
coefficients or conjugate Fontaine's differential
operator onto etale vectors.

\paragraph{The diagonal argument and descent of scalar characters.}
Apply the all-weight calculation of A.41 over \(F_0\)
to (A.42.2), with
\(U_{\rm c}=\langle\Omega_{i,\rm c}\rangle\).
It gives
\[
 D_{\rm cris}(W_d)\subset
                  \mathbb L(T_{\rm c}\oplus U_{\rm c})_d.
                                                               \tag{A.42.6}
\]
Because its left side is \(G_{\rm c}\)-stable, the
exact tensor-intersection argument puts it in
\[
 \mathbb L(T_{\rm c}\oplus U_{\rm c,*})_d,
           \qquad U_{\rm c,*}=\bigcap_{g\in G_{\rm c}}gU_{\rm c}.
                                                               \tag{A.42.7}
\]
Compute these intersections and connected orbits after extending
\(F_0\) to an algebraic closure; all ensuing scalar identities
descend faithfully. The connected diagonal-projector lemma of A.41 applies
to \(G_{\rm c}\) and the local orthogonal component
blocks. It proves that \(G_{\rm c}\) acts on
\(U_{\rm c,*}\) by \(\chi_{\rm c}\). All generators
on the right of (A.42.7) therefore have weight two and
that character. It follows that odd \(W_d\)'s vanish
and \(G_{\rm c}\) acts on \(D_{\rm cris}(W_{2n})\)
by the scalar \(\chi_{\rm c}^n\).

A scalar-character identity is preserved by tensor
change of fibre and faithfully flat descent. Concretely,
trivialize the comparison torsor, or use (A.42.5) over
\(E\). Conjugation of a scalar matrix leaves it
unchanged. The preceding identity therefore becomes
that the global \(G\)-representation on
\(W_{2n}\) is \(\chi_p^n\) after a faithful scalar
extension, and hence already over \(\mathbf Q_p\).
This is precisely the Tate-grading hypothesis of A.38
for the closure of the actual subgroup \(H\). That
criterion proves \(H=I_{P,p}\). Enlarging back to the
specified open derived subgroup, which fixes the same
inertia, proves the assertion.

The algebraic kernel and intersection certificates of
\texttt{code/a\_compact\_type\_fixed.py} apply unchanged:
all their identities are over characteristic-zero fields.
The additional step here is the functorial fibre change
and its scalar-character descent, proved above.
\end{proof}

% A.43: full independent parent proof audit and code rerun passed; registered.
\subsection{An unbounded family of dihedral covers}

\begin{statement}{A.43}{Proved}{Every odd dihedral two-cusp cover has cusp fixed groups at all primes}
For every odd integer \(n\ge3\), let
\[
 C_n:\ v^2=w^{2n+1}-w,
 \qquad X_n=C_n\setminus\{P_0,P_\infty\},
                                                               \tag{A.43.1}
\]
where \(P_0\) and \(P_\infty\) are the unique points
above \(w=0\) and \(w=\infty\). This is a genus-\(n\)
curve and admits a finite etale map \(X_n\to U\) of
degree \(2n\). Geometrically that map is Galois with
dihedral group of order \(2n\). At two its pointed
stable reduction, after finite extension, has a rational
centre carrying both cusps and \(n\) supersingular
elliptic tails; its proper Jacobian has good supersingular
reduction and \([P_0-P_\infty]\) has exact order two.

For every coefficient prime \(\ell\), every cusp packet
at either cusp, and every indicated open derived subgroup,
\[
             \left(\pi_1(X_{n,\bar K})^{(\ell)}\right)^{s(\Gamma_0)}
                                    =I_{P,\ell}.     \tag{A.43.2}
\]
This applies in particular to every induced original packet
and permitted open \(\Delta_0\) of Part A. It proves
Task A.5 for this family of unbounded genus, without a
primality hypothesis on \(n\). It asserts no normalizer
classification for the new higher-genus curves.

The proof includes a general construction. Suppose a smooth
proper positive-genus curve over a number field has good reduction at residue
characteristic \(p\), distinct rational points \(P,Q\)
with the same specialization, and a rational function
\(f\) with \(\operatorname{div}(f)=d(P-Q)\), where
\(d>0\) (exchange \(P,Q\) if necessary).
For any integer \(n\ge2\) with \((n,pd)=1\), the normalized cover \(z^n=f\)
has, after finite extension, a rational component carrying
its two points above \(P,Q\), with \(n\) copies of the
original positive-genus component attached to it. For the
complement of those two points, A.42 gives the same fixed
inertia conclusion at every coefficient prime.
\end{statement}

\begin{proof}
\paragraph{The actual two-branch normalization.}
We first prove the general construction. Work over a finite
local extension at which the proper curve \(C\) has a
smooth model and the distinct marked sections \(P,Q\)
have the same closed point. Separate them by blowing up
until they are disjoint. There is one positive-genus
component \(C_0\), and a chain of rational components
ending in a rational tail carrying both marks. Stabilizing
the pointed model contracts the unmarked valence-two
components and leaves one rational tail \(R_0\) joined
to \(C_0\) by one edge of positive length. The two
marks and that node are three distinct points of \(R_0\).
Use the regular model with the intervening chain retained
for the following divisor calculation.

The horizontal divisor \(dP-dQ\) has degree zero on
every component, because both marked sections meet the
same last rational component. Write \(a_v\) for the
vertical order of \(f\) along component \(v\).
Intersect its principal divisor with each component.
The intersection equations are the graph Laplacian
\[
                       \sum_{u\sim v}(a_u-a_v)=0.
                                                               \tag{A.43.3}
\]
The kernel of the Laplacian on a connected finite graph
consists of constants: multiply each equation by \(a_v\),
sum, and obtain the negative of
\(\sum_{\{u,v\}}(a_u-a_v)^2\). Thus all vertical
orders equal one integer \(a\). If \(\varpi\) is a
uniformizer, \(f_0=\varpi^{-a}f\) has precisely the
horizontal divisor \(dP-dQ\), with no vertical component.
It is a unit at every node. On \(C_0\) its reduction
has divisor zero and is a nonzero constant \(c\); on
all intervening rational components the same is true.
Choose a coordinate \(u\) on the final tail for which
\(P=0\), \(Q=\infty\), and its attaching node is
\(u=1\). There
\[
                         \overline f_0=c u^d.        \tag{A.43.4}
\]
This calculation identifies the reduction of the given
function; it is not a choice of another cover with the
same graph.

Make a finite base extension containing an \(n\)-th
root of \(\varpi^a\), the needed \(n\)-th root of
\(c\), and all \(n\)-th roots of unity. Since
\(n\) is prime to \(p\), the derivative in the
Kummer equation is a unit wherever \(f_0\) is a
unit. On the formal neighbourhood of \(C_0\), the
reduction of \(f_0\) is constant. Hensel's lemma
therefore gives \(n\) separate roots, compatibly
on overlapping formal charts: the root with a fixed
reduction is unique. Thus the actual normalized cover
splits into \(n\) copies of \(C_0\) there. At each
node the complete local ring has equation
\(xy=\varpi^e\), and \(f_0\) is a unit with
constant nonzero reduction. The same Hensel argument
splits the cover into \(n\) copies of this nodal
chart. In particular every connecting annulus is
unramified and keeps its positive length.

On \(R_0\), equation (A.43.4) gives the connected
normalization of
\(z^n=u^d\), after scaling the constant. Because
\((n,d)=1\), choose integers \(a',b'\) with
\(da'+nb'=1\). The parameter
\(s=z^{a'}u^{b'}\) satisfies
\[
                           u=s^n,\qquad z=s^d.
                                                               \tag{A.43.5}
\]
Its smooth compactification is a rational curve. It is
fully ramified at the two marks and unramified at the
node \(u=1\), which has \(n\) distinct preimages.
Near a marked section the same normalized local
parameter gives a smooth chart over the valuation
ring; the ramification is tame. The remaining charts
are etale. Consequently normalization of the actual
model, with the harmless chain components retained,
is semistable. Contracting unmarked valence-two
rational components gives precisely one rational
centre and \(n\) copies of \(C_0\), each attached
at one of those distinct node preimages. The cover
is connected, since valuation \(d\) at a branch
point is coprime to \(n\).

The special dual graph is a tree. Riemann--Hurwitz,
with two totally ramified points, gives
\[
       2g(C')-2=n(2g(C)-2)+2(n-1),
                          \qquad g(C')=n g(C).       \tag{A.43.6}
\]
This agrees with the sum of the genera of its \(n\)
positive-genus components. Its two removed points
are both on the rational centre. Thus A.41--A.42
apply to either selected cusp, proving the general
construction and its all-prime fixed-group assertion.
Additional finite global extensions realizing the
required local extension are harmless, by the field
restriction argument of A.41.

\paragraph{The explicit cover of the once-punctured elliptic curve.}
For the family (A.43.1), put
\[
 E:\ y^2=x^3-x,\qquad T=(0,0),
 \qquad q_n:C_n\longrightarrow E,
          \quad x=w^n,\quad y=w^{(n-1)/2}v.           \tag{A.43.7}
\]
The equation follows on squaring the last expression.
It is the normalization of the cyclic cover \(w^n=x\).
On \(E\),
\(\operatorname{div}(x)=2T-2O\). Since \(n\)
is odd, \(q_n\) has degree \(n\), with total
ramification precisely at \(T,O\). It has no other
ramification. The polynomial defining \(C_n\) has
\(2n+1\) simple roots, and its hyperelliptic model
has genus \(n\), as also follows from (A.43.6).

Over \(K=\mathbf Q(i)\), multiplication by \(1+i\)
on \(E\) is
\[
 [1+i](x,y)=
 \left(\frac{x-x^{-1}}{2i},
             -\frac{1+i}{4}\,y(1+x^{-2})\right).
                                                               \tag{A.43.8}
\]
Its kernel is \(\{O,T\}\). Hence \([1+i]\circ q_n\)
has degree \(2n\) and is etale after precisely the
two points \(P_0,P_\infty\) are deleted and the
point \(O\) is deleted downstairs. This is the
claimed finite etale map \(X_n\to U\).

For the dihedral assertion, after adjoining
\(\zeta_n\) the cyclic generator is
\[
 r(w,v)=(\zeta_n w,\zeta_n^{(n+1)/2}v).
\]
Translation by \(T\) on \(E\), which is
\((x,y)\mapsto(-x^{-1},y x^{-2})\), lifts to
\[
 h(w,v)=\left(-w^{-1},
           (-1)^{(n-1)/2}\frac{v}{w^{n+1}}\right).
                                                               \tag{A.43.9}
\]
Direct substitution gives \(h^2=1\) and
\(hrh=r^{-1}\). These transformations form a
faithful group of order \(2n\) and preserve
\([1+i]\circ q_n\). Since that map already has
that degree, it is the geometric Galois cover with
this dihedral group. In particular this assertion
holds for composite odd \(n\) as well.

\paragraph{Its actual stable reduction at two.}
The elliptic curve \(E\) acquires good supersingular
reduction at two, as in A.40. Its point \(T\) is
nonzero two-torsion, so specializes to \(O\): the
geometric special fibre has no nonzero two-primary
torsion. Apply the general normalization argument
with \(f=x\), \(d=2\), and this odd \(n\).
It gives \(n\) copies of the supersingular elliptic
special fibre attached to one rational centre, with
both marked points on that centre. The Jacobian is
therefore an abelian scheme with special fibre the
product of those \(n\) elliptic curves, by the
compact-type product recalled in
\cite[Section 1]{fghDupuyRabinoff}.
In particular its special \(2\)-rank is zero.
The proper Jacobian has dimension \(n\) and
supersingular reduction.

Finally, on \(C_n\) the function \(w\) has
\[
                  \operatorname{div}(w)=2P_0-2P_\infty.
                                                               \tag{A.43.10}
\]
The divisor class has exact order two: if the two
points were linearly equivalent, a degree-one function
would make the positive-genus curve rational.
Thus both the direct rational-path condition and
the torsion-divisor condition of A.41 are verified
on this actual family. Applying A.42 yields (A.43.2)
for all coefficient primes, including two and every
prime dividing \(n\).

The program \texttt{code/a\_odd\_dihedral\_family.py}
checks the isogeny and cover identities exactly,
the involution and dihedral exponents for sample
odd composite and prime values, the normalized
rational-tail parameter, and the connected graph
Laplacian. The model and all-prime arguments are
proved uniformly above.
\end{proof}

% A.44: full independent parent proof/source audit and exact code rerun passed.
\subsection{Dihedral covers of two-power degree}

\begin{statement}{A.44}{Proved}{The four-cusp two-power dihedral family at every coefficient prime}
Let \(m=2^s\), \(s\ge1\), and put
\[
 B_m:\ v^2=w^{2m}-1,
 \qquad Y_m=B_m\setminus\{A,B,C,D\},               \tag{A.44.1}
\]
where \(A,B\) have \(w=0\), \(v=i,-i\), and
\(C,D\) are the two infinity points with
\(v/w^m\to1,-1\), respectively. Then \(B_m\)
has genus \(m-1\), and \(Y_m\to U\) is a finite
etale geometric dihedral cover of degree \(2m\).
Its proper Jacobian has potentially good reduction at
two, with special \(2\)-rank zero. Every difference
of two of the four cusps is killed by \(4m\).
Consequently the four cusp vertices of its pointed
stable reduction are joined entirely through rational
components.

For every coefficient prime \(\ell\), every cusp
packet at any one of these four cusps, and every
indicated open derived subgroup,
\[
                   \left(\pi_1(Y_{m,\bar K})^{(\ell)}\right)^{s(\Gamma_0)}
                                      =I_{P,\ell}.   \tag{A.44.2}
\]
This includes all induced original packets and permitted
open \(\Delta_0\) of Part A. It proves the fixed-group
part of Task A.5 for this second family of unbounded
genus. No normalizer classification for these new
higher-genus curves is asserted.
\end{statement}

\begin{proof}
\paragraph{The actual cover and its dihedral group.}
There is a map to \(E:y^2=x^3-x\) given by
\[
                 x=w^m,\qquad y=w^{m/2}v.            \tag{A.44.3}
\]
It has degree \(m\): the extension over \(K(x)\)
through \(K(w,v)\) has degree \(2m\), and the
intermediate elliptic function field has degree two.
It is the normalization of \(w^m=x\). The zeros
and poles of \(x\) on \(E\) have order two, so
there are two points above each of \(T=(0,0)\) and
\(O\), each with ramification index \(m/2\), and
no other ramification. Composing with the degree-two
isogeny \([1+i]\) in (A.43.8) gives the asserted
finite etale map to \(U\), of degree \(2m\),
with precisely the four displayed cusps removed.
The hyperelliptic polynomial has degree \(2m\) and
simple roots; hence the proper genus is \(m-1\).

Choose a primitive \(2m\)-th root \(\eta\) with
\(\eta^{m/2}=i\), after a finite extension, and set
\(\zeta_m=\eta^2\). The cyclic generator and a
lift of translation by \(T\) are
\[
 r(w,v)=(\zeta_m w,-v),
 \qquad h(w,v)=(\eta/w,-i v/w^m).                   \tag{A.44.4}
\]
The relation \(\eta^m=-1\) verifies the curve
equation, \(h^2=1\), and \(hrh=r^{-1}\).
Their induced maps on \(E\) are respectively the
identity and \((x,y)\mapsto(-x^{-1},yx^{-2})\).
They form a faithful group of order \(2m\) preserving
the composite map. Therefore the geometric Galois group
is dihedral of that order. For \(m=2\), this means
the group of order four. In every case it is a
\(2\)-group.

\paragraph{Explicit two-primary cusp divisors.}
The following divisors are computed on the smooth
compactification:
\[
 \begin{split}
 \operatorname{div}(w)&=A+B-C-D,\\
 \operatorname{div}\left(\frac{v-i}{v+i}\right)&=2m(A-B),\\
 \operatorname{div}(v-w^m)&=m(C-D).
 \end{split}                                       \tag{A.44.5}
\]
For the second identity use
\((v-i)(v+i)=w^{2m}\): at \(A\) one factor
is a unit and the other has order \(2m\), with
the opposite assertion at \(B\). Their poles at
the two infinities cancel in the quotient. For the
third use \((v-w^m)(v+w^m)=-1\); the respective
orders at \(C,D\) are \(m,-m\).
The first identity is the ordinary degree-two
hyperelliptic map at zero and infinity.

Twice \(m\) times the first divisor, plus the
second and minus twice the third, is \(4m(A-C)\).
With the second subtracted instead it is \(4m(B-C)\).
The third already kills \(C-D\) after multiplication
by \(m\). Thus every cusp difference is killed
by \(4m\), which is a power of two. We only need
this annihilator, and make no claim that all these
orders equal \(4m\).

\paragraph{Complex multiplication and good reduction.}
Let
\(\tau(w,v)=(\zeta_{2m}w,v)\). A holomorphic
basis is
\(w^j\,dw/v\), \(0\le j\le m-2\), so the
characters on this basis have exponents
\(1,\ldots,m-1\). On the full complex homology
we obtain each exponent
\[
              1,\ldots,m-1,m+1,\ldots,2m-1
\]
once. Thus the faithful endomorphism algebra generated
by \(\tau\) is
\[
 \mathbf Q[\tau]
   =\mathbf Q[t]\bigg/\frac{t^{2m}-1}{t^2-1}
   =\prod_{a=2}^{s+1} F_a,
                  \qquad F_a=\mathbf Q(\zeta_{2^a}).  \tag{A.44.6}
\]
Its dimension is \(2m-2=2g(B_m)\). The rational
idempotents split the Jacobian, up to isogeny after
a finite field extension, into abelian varieties
\(J_a\) with
\[
       \dim J_a=2^{a-2},\qquad
       F_a\hookrightarrow\operatorname{End}^0(J_a),
       \qquad [F_a:\mathbf Q]=2\dim J_a.            \tag{A.44.7}
\]
Indeed the \(F_a\)-summand in rational homology
has precisely that field's degree. No simplicity
of \(J_a\) is needed. These are abelian varieties
with complex multiplication in the usual degree-
\(2\dim J_a\) sense.

By \cite[Theorem 6(a), printed p.503]{fghSerreTate1968},
every \(J_a\) has potentially good reduction at two.
After a common finite extension so does their product
and its isogenous Jacobian. Good reduction is invariant
under isogeny by \cite[Corollary 2 to Theorem 1, printed p.493]{fghSerreTate1968};
equivalently apply the Neron--Ogg--Shafarevich criterion
to the rational Tate module at an auxiliary prime different from two. Make all the \(F_a\)-actions
defined before taking these good models. The same primary
proof of Theorem 6 specializes these field actions to
their good reductions using extension of endomorphisms
to the abelian schemes.

\paragraph{An elementary proof of special two-rank zero.}
For every \(a\ge2\), the polynomial
\[
                  \Phi_{2^a}(1+z)=(1+z)^{2^{a-1}}+1
                                                               \tag{A.44.8}
\]
is Eisenstein at two: the constant term is two, all
other nonleading coefficients are even, and the leading
coefficient is one. Hence
\(F_a\otimes\mathbf Q_2\) is a field of degree
\(2\dim J_a\).

Let \(\widetilde J_a\) be its good special fibre,
of dimension \(g_a\), and let
\[
 T_2^{\rm et}(\widetilde J_a)
       =\varprojlim_n\widetilde J_a[2^n](\bar\kappa),
 \qquad V_2^{\rm et}=T_2^{\rm et}\otimes\mathbf Q_2.
\]
Its dimension is the special \(2\)-rank \(f_a\),
which satisfies \(0\le f_a\le g_a\). One way
to see the upper bound is to factor \([2]\) through
relative Frobenius of purely inseparable degree
\(2^{g_a}\); its remaining Verschiebung factor
has degree \(2^{g_a}\), so the geometric kernel
of \([2]\) has at most \(2^{g_a}\) points.

The specialized \(F_a\)-action makes
\(V_2^{\rm et}\) a unital module over
\(F_a\otimes\mathbf Q_2\). If it were nonzero,
its \(\mathbf Q_2\)-dimension would be at least
\([F_a:\mathbf Q]=2g_a\), contradicting
\(f_a\le g_a\). Therefore \(f_a=0\).
Two-rank is additive in products and invariant under
isogeny, so the good special fibre of the full Jacobian
has two-rank zero. This argument requires neither a
Shimura--Taniyama slope computation nor a claim about
its complete Newton polygon.

\paragraph{The marked paths and the actual fixed groups.}
Take a semistable model of \(B_m\) after this finite
extension. Its dual graph is a tree, because the torus
rank of its semistable Jacobian is the graph's first
Betti number, whereas its Jacobian has good reduction.
This is the same compact-type criterion used in A.25
and A.40: the generalized special Jacobian is an
extension of the product of the component Jacobians
by \(\mathbf G_m^{b_1}\), as follows by gluing
line bundles along the nodes.

All cusp differences in (A.44.5) are two-primary
torsion. The special Jacobian has no nonzero geometric
two-primary torsion, so all four cusp Abel images
specialize to the same point. The extended Abel map
and its componentwise description in A.41 now show
that the entire path between every pair of marked
cusps consists of rational components. In particular
for each selected cusp there is a distinct cusp
satisfying the hypotheses of A.41--A.42. They prove
(A.44.2) at every coefficient prime.

At the coefficient prime two there is also an
independent direct verification: the geometric Galois
closure of this cover is the displayed dihedral
\(2\)-group. Thus A.22 embeds its full geometric
pro-two group as an open subgroup of the original
one and gives exactly the same inertia fixed group.
For odd coefficient primes the prime-to-residue-
characteristic statement A.41 already suffices.
The field and packet restrictions are precisely
those handled in A.41.

The exact program \texttt{code/a\_two\_power\_dihedral.py}
checks the curve and deck identities, the cusp-divisor
integer combinations, the multiplicity-one character
polynomial, and the Eisenstein certificates for several
powers of two. Their proofs above are uniform in \(s\).
\end{proof}

\subsection{The universal covering question remains separate}

\begin{statement}{A.45}{Open}{The full scope of Task A.5}
For arbitrary open \(H\subset\pi\), arbitrary coefficient prime
\(\ell\), and every original cusp packet preserving the indicated
arithmetic cover, we have not determined both the full fixed subgroup
of \(H^{(\ell)}\) under \(s_b(\Delta_H)\) and all its closed
procyclic subgroups with open arithmetic normalizer.
In particular we neither prove nor refute the universal assertion
that both kinds of subgroup are cuspidal. The proved families below
do not form a cofinal family of all open geometric subgroups.
\end{statement}

\begin{proof}[Cases proved and the remaining comparisons]
The actual full profinite fixed subgroup is known for every such
\(H\): it is \(H\cap I\), by A.9. This gives the exact
criterion for a fixed pro-\(\ell\) element to have a fixed lift
to \(H\). It does not compute fixed elements of the quotient
\(H^{(\ell)}\), since the quotient map need not be surjective
on fixed subgroups.

For the quotient problem, A.10 completes all first covers in
Task A.4, including every prime and packet. A.40 completes every
cover with genus-one compactification, with arbitrary isogeny
degree. A.22 also supplies the coefficient prime whenever the
geometric Galois closure is a group of that prime's power order.
The all-prime theorem A.42 supplies the fixed subgroup for any
compact-type reduction with a rational joining path from the
selected cusp to a second cusp. A.43 and A.44 give two explicit
families of unbounded genus satisfying that criterion. A.47 also
completes the normalizer classification at every prime for the
entire odd dihedral family of A.43. A.48 completes every
normalizer in the two-power dihedral family of A.44, including
the previously displayed genus-seven case at coefficient seventeen.
A.49 supplies complete fixed-group and normalizer calculations
for explicit covers in every proper genus, including all even genera;
its extra cusps lie on a rational core in an actual supersingular
Kummer model.
A.50 and A.51 complete both all-prime assertions in the
degree-twenty-seven Heisenberg cover of genus ten. Its residual
surface has actual quaternion multiplication, which supplies
the required reduction and cusp paths at three.
A.52 computes that three-adic reduction directly and completes both
all-prime assertions for every enlargement
\(C\setminus q^{-1}(E[3m])\), for all positive integers \(m\).
Odd multiples use its explicit wild model at three, and even
multiples use the smooth tame map at two.
A.53 proves that every proper compactification occurring in Task A.5
admits infinitely many torsion enlargements with the full all-prime
fixed-inertia theorem. The comparison to the original affine curve is
an open immersion and the covering map is changed by an elliptic
isogeny; this supplies neither a finite refinement over the old map
nor a lift of its fixed elements. Normalizer recognition, when already
known on the original curve, does pass to these enlargements by A.11.
A.54 gives the actual degree-one-hundred-twenty-five Heisenberg cover
of genus fifty-one. Its eight-dimensional multiplicity factor has no CM
subvariety, proved by two persistent global Frobenius fields. All
normalizers are cuspidal for coefficient primes not congruent to one
modulo five; the full fixed group at five follows from A.22.
A.55 completes every remaining split-prime normalizer, including
eleven, using actual cyclic prime-power covers of the genus-two
three-cusp quotient \(w^5=t(t-1)\). Their Pryms have strictly
interior Newton slopes, and an integral transfer calculation detects
even divisible proper-homology classes. The deck translates of
the actual quotient maps detect all remaining quartic-CM lines.
A.58 now completes the full fixed-subgroup comparison at every prime
for this same Heisenberg cover. Its exact five-adic charts give one
elliptic main component, five genus-six and ten genus-two
Artin--Schreier tails. Their genera exhaust fifty-one; all twenty-five
cusps occur in rational clusters away from those tails. Thus the
rational-path criterion A.42 applies to every cusp and packet.

A.56 also gives a normalizer criterion for every proper genus at least two
with potentially good curve reduction at the coefficient prime and
special p-rank at least two. For covers of the original elliptic curve,
at Gaussian split primes it removes the special-rank condition and applies
uniformly to every degree-\(d\) cover when \(p>d\) and
\(p\equiv1\pmod4\). These are normalizer results only.

A.57 completes the special-rank-one case for arbitrary proper curves,
so potential good curve reduction at the coefficient prime is now
sufficient without any rank condition. Consequently every degree-\(d\)
cover has cuspidal normalizers at every \(p>d\), with no congruence
restriction. The outstanding normalizer cases therefore require the
coefficient prime to be at most the degree and the proper covering curve
to have no potentially good reduction at the places above that prime,
unless one of the other proved criteria applies.


A.59 further completes every coefficient-prime normalizer for
every cover whose geometric Galois closure is finite nilpotent,
including all non-Galois covers with such a closure. Its actual
intermediate cover has prime-to-coefficient Galois closure, while
the upper cover has coefficient-prime-group closure and induces
an open pro-prime embedding. No such factorization is asserted
for arbitrary finite groups, and this normalizer theorem does
not determine the full fixed groups of all nilpotent covers.


A.60 completes both assertions at every coefficient prime for
the next actual Heisenberg cover, of degree343, proper genus148
and49cusps. Its explicit seven-adic model has one supersingular
elliptic main component and seven genus21 Artin--Schreier tails;
all marks occur on a rational tree over the origin, disjoint
from those tails. A.61 extends both conclusions uniformly to
every inert-prime Heisenberg cover
\(z^p=\psi_p(x)\), \(p\equiv3\pmod4\), of degree
\(p^3\), genus \(1+p^2(p-1)/2\), and \(p^2\) cusps.
Its divided differential operator proves the local coefficient
identity for all such primes, and the exact charts give one
supersingular elliptic main component and \(p\) tails of genus
\(p(p-1)/2\). This is an unbounded family, with no cofinality
claim and no fixed-group assertion for every nilpotent cover.

A.62 completes the ordinary-prime half by an actual unramified
translation comparison modulo \(p^2\) and separable-isogeny descent.
It gives \(p\) quartic and \(p(p-3)\) quadratic Artin--Schreier
tails, with all cusps on rational clusters of size \(p\).
Thus both assertions are now proved uniformly for the actual
exponent-\(p\) Heisenberg cover at every odd prime \(p\).
This remains a particular unbounded family of geometric subgroups.

A.63 further proves both assertions for every isogeny enlargement
\(C_p\setminus q_p^{-1}(E[pm])\) when \(p\equiv1\pmod4\),
for every \(m\ge1\). An actual noncentral deck subgroup fixes
each relevant special component and supplies a rational companion
for every cusp, including those whose images meet a positive-genus
tail. The comparison to the smaller marked curve is an open immersion.

A.64 removes the congruence restriction in this enlargement conclusion:
every odd Heisenberg prime and every positive enlargement integer now
satisfy both assertions. More generally, if an enlargement integer has
a prime outside the original geometric closure order, the actual tame
good model at that prime gives rational cusp companions and full
fixed inertia at all coefficient primes, including that residue prime.
This includes residue two. It is still an enlargement of the removed
set on the same compactification, with no fixed-lift or cofinality claim.

A.65 also settles every intermediate of each odd Heisenberg cover,
equivalently every cover whose geometric closure has odd-prime exponent
and nilpotency class at most two. The noncentral order-prime quotients
are treated through their actual stable-model maps, rather than by
transferring fixedness through a surjective pro-prime quotient.

A.66 allows an additional abelian geometric factor of order prime to
the Heisenberg prime, including all nonnormal intermediate subgroups.
The actual elliptic-isogeny base change preserves one connected
radicial main and copies each rational marked joining path.
Products of two nonabelian Sylow factors are still outside this
argument; no general nilpotent fixed-group assertion follows.

A.67 supplies a non-nilpotent family: every tetrahedral
\(A_4\) closure and every intermediate has both assertions at
all coefficient primes. The actual three-adic model has a
rational six-cusp centre and four supersingular elliptic
tails. An explicit residue-field argument treats the wild
order-three quotient and its rational two-cusp centre.
This does not extend automatically to coprime elliptic-isogeny
pullbacks: their three-adic dual graph can have cycles.

A.68 completes the two nonstable genus-seven dihedral covers
\(v^2=w^{16}\pm6w^8+1\), and every actual intermediate
of their degree-sixteen closures, at all coefficient primes.
An actual genus-two quotient, exact Igusa invariants and Liu's
residue-two criterion give potential good Jacobians of
\(2\)-rank zero. Explicit principal divisors kill every cusp
difference by eight; norm maps and the Frattini inertia calculation
carry the rational-path criterion to each intermediate.
A.69 now extends this conclusion uniformly to
\(v^2=w^{2m}\pm6w^m+1\) for every two-power \(m\),
including every actual intermediate. The primitive cyclotomic
new factor has special \(2\)-rank zero or is ordinary. The ordinary
alternative contradicts the actual cyclic component action,
finite-flat lifting of smooth fixed points, and the genus of the
rooted marked tree.
A.70 now classifies every actual two-power dihedral closure of
the fixed elliptic base and proves both assertions for all its
intermediates. The three characteristic rotation quotients are
realized by explicit elliptic two-isogenies with distinct dual
kernels; the order-four abelian case is treated separately.
A.71 permits an additional odd-order abelian factor, with every
actual intermediate included. The connected elliptic-isogeny
pullback has a separately proved compact-type model, and its
rational cusp paths lift as literal copies of the original
marked subtree. No special two-rank assertion for the enlarged
Jacobian is needed.
A.72 settles every actual generalized-quaternion two-group
closure and every intermediate. Its central quotient is one
of the classified dihedral covers. A primitive cyclotomic
anti-invariant factor could only have rank zero or be ordinary;
the ordinary case would force a component orbit twice as large
as its total genus, using the quaternion group's unique
involution and the ordinary fourth-root obstruction. The
actual elliptic abelianization quotient gives the required
two-primary cusp divisors.
A.73 supplies the same complete result for every actual
semidihedral two-group closure and its intermediates. Here
the ordinary alternative would force one orbit of elliptic
components; rooting the reduction tree at the whole group's
invariant centre makes them equidistant. Actual fixed-point
lifting and disjoint unmarked child subtrees then contradict
the total genus.
A.74 extends the odd-abelian preservation theorem to all three
two-group families, including every quotient closure and
every intermediate. The same actual connected isogeny
pullback, separately proved tree condition and literal
rational-path lift apply to quaternion and semidihedral
factors as well.
A.75 completes every actual modular two-group closure and
every intermediate. Its unmarked proper stable-tree argument
rules out an ordinary anti-factor by a generically fixed-point-free
outside involution; no marked-fibre slope assertion is assumed.
A.76 combines these four families into the exact class of
two-groups with a cyclic subgroup of index at most two, including
the small orders and every odd abelian extension through the
actual connected isogeny pullback.
A.77 further completes every geometric closure whose commutator
subgroup has at most two elements, with all intermediates and
the automatic odd abelian factors. Actual faithful quotient
curves supply a surjective sum of Jacobian pullbacks, giving
special rank zero for the two-group source. The order-thirty-two
example has genus nine and sixteen cusps and lies outside the
cyclic-index-two class. Products of two nonabelian Sylow factors,
arbitrary two-groups and the universal question remain outside
these conclusions.

A.78 now proves potential good Jacobian reduction of special
rank zero for every prime-power Galois cover of a good
positive-genus rank-zero curve branched at most at one
horizontal section in mixed characteristic. A.79 combines
this with the actual cusp-divisor relation for cyclic-derived
closures: for \(p=2\) or odd \(p\equiv3\pmod4\),
every quotient of a cyclic-derived \(p\)-group times a
prime-to-\(p\) abelian group satisfies both assertions,
including all intermediates. Its actual connected isogeny
pullback preserves each rational companion path.
The rank theorem alone does not give cusp torsion or a
rational companion path for an arbitrary prime-power group;
that fixed-group frontier and the universal question remain
open.


A.80 settles the actual natural genus-two degree-eight intermediate
of the wreath closure \(C_2\wr C_4\), whose derived subgroup
is noncyclic. The remaining marked configuration has component
degrees \((2,0)\), and a cusp can lie on the positive-genus
Prym tail. An exact partial-distance kernel removes the original
elliptic handles while retaining the Prym handles; its non-CM
\(\mathrm{SL}_2\) action then makes every fixed grade Tate.
This settles that actual intermediate at every coefficient prime,
without claiming fixedness for all subgroups of its genus-seventeen
Galois closure.

A.81 proves the same conclusions for any actual genus-two double
cover of an elliptic curve two-primary isogenous to the base,
branched at two isogeny-kernel points and with non-CM elliptic
Prym. In particular all wild two-power torsion enlargements of
the A.80 curve are included, with degree \(2\cdot4^{r+1}\)
and \(2\cdot4^{r+1}-2\) cusps. Each is treated through its
actual degree-two model and two-group closure; no fixed-element
transfer through the affine open immersion is asserted.








A.46 supplies a general
residue-prime normalizer criterion when the potentially good
Jacobian has special p-rank zero. The remaining
normalizer assertions have their own hypotheses: A.11 reduces
a noncuspidal normalized procyclic subgroup to a nonzero proper
homology line, whereas a fixed-group calculation alone makes
no such classification.

The unresolved geometric cases include reductions with a cycle
in their dual graph, and compact-type reductions in which every
path from the selected cusp to a different cusp meets a
positive-genus component. The first-core proof of A.41 uses
positive distances from that joining path to every positive-genus
component. On the path that distance vanishes; for a graph with
cycles the additional graph homology also changes the argument.
Thus one cannot extend A.42 just by allowing zero coefficients
or by applying it to an unspecified semistable alteration.
A.80 supplies a partial-distance replacement for its actual non-CM
elliptic Prym configuration. No general replacement covering all
these remaining cases is proved here.

For normalizers, one must exclude the remaining actual rational
lines in proper homology, not only list their CM characters.
An isogeny factor of the Jacobian does not supply a morphism to
a punctured elliptic target: the inverse image of its missing
point must consist entirely of cusps. A.11 uses that actual
morphism hypothesis. Neither the Jacobian decomposition nor
the absence of new fixed lifts supplies it automatically.

The former example \(Y_8:v^2=w^{16}-1\) at coefficient
seventeen is now settled by A.48 and is not an open case.
The universal problem still includes covers outside all the
proved families and geometric hypotheses above. No cofinality
argument for those families, or least-index determination of
an undecided cover among all open \(H\), has been supplied.
The next required step is an actual normalizer obstruction for
the remaining proper-homology lines, or a constructed noncuspidal
subgroup with its open arithmetic normalizer. Neither an
abstract character nor a vector in a graded completion is such
a counterexample.
\end{proof}

% Independent full-draft audit passed; see notes/de_A46_audit.md.
\subsection{A residue-prime normalizer criterion from zero p-rank}

\begin{statement}{A.46}{Proved}{Zero p-rank excludes noncuspidal pro-p normalizers}
Let \(Y\) be an affine hyperbolic curve over a number field \(L\),
let \(C\) be its smooth proper compactification, and put
\(J=\operatorname{Jac}(C)\). Suppose that at some place of \(L\)
above a rational prime \(p\), the abelian variety \(J\) has
potentially good reduction with geometric special \(p\)-rank zero.
In the geometrically pro-\(p\) arithmetic fundamental group of \(Y\),
every nontrivial closed procyclic subgroup of the geometric kernel
whose arithmetic normalizer has open image in \(G_L\) is contained
in a conjugate of a cusp inertia group.

The representation-theoretic assertion used here is stronger:
\(V_pJ=T_pJ\otimes_{\mathbf Z_p}\mathbf Q_p\) has no
one-dimensional \(\mathbf Q_p\)-subquotient after restriction to
\(G_{L'}\), for any finite extension \(L'/L\).
Both assertions include \(p=2\). No splitting of the arithmetic
normalizer is assumed.
\end{statement}

\begin{proof}
We first prove the representation assertion locally. Enlarge the
completion to a finite extension \(k/\mathbf Q_p\) over which \(J\)
has good reduction, and let \(\widetilde J\) be its special fibre.
Write \(k_0=W(\kappa)[1/p]\) for the maximal unramified subfield,
and \(f=[k_0:\mathbf Q_p]\). The covariant crystalline functor
on cohomology identifies, as Frobenius modules,
\[
 D_{\mathrm{cris},k}\bigl(H^1_{\acute{\mathrm e}\mathrm t}
              (J_{\bar k},\mathbf Q_p)\bigr)
       \simeq H^1_{\mathrm{cris}}(\widetilde J/W(\kappa))[1/p].
                                                        \tag{A.46.1}
\]
Here the good-reduction comparison is applied with constant
coefficients to the abelian scheme: in
\cite[Section 5, Corollary 4 and the subsequent Remarks 1--5,
printed pp.256--260]{RootFaltings2002}, invert the crystalline period
in the comparison and its inverse, and take Galois invariants.
There is no boundary divisor or logarithmic monodromy in this
smooth proper case. Equivalently this is the usual crystalline
comparison for the first cohomology of an abelian scheme.

The Newton slopes on the right of (A.46.1) lie in \([0,1]\).
The multiplicity of slope zero is the special \(p\)-rank, and
polarization pairs every slope \(a\) with \(1-a\).
Thus all slopes lie strictly between zero and one. These facts,
including invariance under finite residue extension, are recalled
in \cite[Sections 2.2.1 and 2.2.3]{RootAchterPries2013}.

Crystalline representations are closed under subquotients, and
\(D_{\mathrm{cris},k}\) is exact and preserves dimension;
see \cite[Theorem 5.2.1(1)--(3) and Proposition 9.1.6]{fghBrinonConrad}.
Suppose a one-dimensional subquotient existed. Duality converts
one of \(V_pJ\) into one of its dual
\(H^1_{\acute{\mathrm e}\mathrm t}(J_{\bar k},\mathbf Q_p)\).
Its crystalline module \(D\) is a rank-one \(k_0\)-isocrystal.
Choose a basis \(e\), and write
\[
       \varphi(e)=a e,\quad a\in k_0^\times,
       \qquad
       \varphi^f(e)=
         \left(\prod_{j=0}^{f-1}\sigma^j(a)\right)e .
                                                        \tag{A.46.2}
\]
For the normalization \(v_p(p)=1\), its unique Newton slope is
\[
 \frac1f v_p\left(\prod_{j=0}^{f-1}\sigma^j(a)\right)
                         =v_p(a)\in\mathbf Z,          \tag{A.46.3}
\]
because \(k_0/\mathbf Q_p\) is unramified. This remains true
however ramified \(k/\mathbf Q_p\) is. Exactness makes \(D\)
a subquotient of the isocrystal in (A.46.1). Its slope must
therefore be one of that isocrystal's slopes: after taking
\(\varphi^f\), characteristic polynomials multiply in an exact
sequence. This contradicts the strict interval \((0,1)\).
No extension of the coefficient field \(\mathbf Q_p\) has been
made; a line only after such an extension would not be excluded
by (A.46.3).

For a finite global extension \(L'/L\), choose a completion above
the indicated place and then a further finite local extension
giving good reduction. A \(\mathbf Q_p\)-line or subquotient for
\(G_{L'}\) remains one on this local group. Good reduction and
special \(p\)-rank zero persist under these extensions. The local
contradiction therefore proves the asserted global statement
for every \(L'\).

Finally let \(A\subset\pi_1(Y_{\bar L})^{(p)}\) have the
normalizer in the statement. If it were noncuspidal, A.11 would
give a nonzero image of \(A\) in \(V_pJ\), whose rational span
is a one-dimensional \(\mathbf Q_p\)-subrepresentation of an open
arithmetic subgroup. More explicitly, geometric inner actions
are trivial on proper homology, so the normalizer's action on
this line factors through its open image in \(G_L\).
This is precisely a line over a finite extension excluded above.
The argument uses the actual proper-homology map and requires
neither a section of the normalizer nor any lifting of fixed
vectors from a rational completion.
\end{proof}

% Registered A.47: full independent parent proof audit and code rerun PASS.
\subsection{Normalizers in every odd dihedral cover}

\begin{statement}{A.47}{Proved}{All-prime normalizer recognition for the odd dihedral family}
For every odd \(n\ge3\), every prime \(\ell\), and
\(X_n\) of A.43, the nontrivial closed procyclic
subgroups of \(\pi_1(X_{n,\bar K})^{(\ell)}\)
whose arithmetic normalizer has open Galois image are
exactly the nontrivial closed subgroups of conjugates
of the two cusp inertia groups. Together with A.43,
this proves both the full fixed-group and normalizer
assertions of Task A.5 for the entire odd dihedral
family, including composite \(n\).

The tame part uses the following general obstruction:
in compact-type reduction, a nonzero proper homology
class of an arithmetically normalized procyclic subgroup
cannot have nonzero projections to two different
positive-genus components, when the coefficient prime
differs from the residue characteristic.
\end{statement}

\begin{proof}
\paragraph{The compact-type component-support obstruction.}
We prove the stated general obstruction first. Work at
residue characteristic \(p\) with coefficient
\(\ell\ne p\), after finite extension so the
semistable model and its components are split. Write
\[
               V=\bigoplus_{i=1}^r V_i,
               \qquad\Omega=\sum_{i=1}^r\Omega_i,
                                                               \tag{A.47.1}
\]
for the proper homology and its local orthogonal
symplectic decomposition. The \(V_i\)'s belong to
the positive-genus components of the compact-type
special fibre. Choose an ordinary base point at a
smooth point on any root component. Root and orient
the graph away from it; let \(m_e>0\) be the edge
lengths and \(i_e\) the indicator that component
\(i\) is downstream from edge \(e\).

The actual based edge formula of
\cite[Definition 3.1.5 and Theorem 3.1.6]{fghBettsDogra}
gives, on the proper rational unipotent group, zero
first Johnson degree and first possible degree two
\[
 D(v_i)=[W_i,v_i],\qquad
 W_i=\sum_e m_ei_e\delta_e,\qquad
 \delta_e=\sum_j j_e\Omega_j.                        \tag{A.47.2}
\]
A nonzero tame inertia power multiplies this entire
derivation by a common nonzero scalar, which does not
affect the argument. The rational graded Lie algebra
through the required degrees is the usual surface
presentation
\[
                           \mathfrak l=\mathbb L(V)/(\Omega),
                                                               \tag{A.47.3}
\]
as used in A.25. The choice of representative of
outer monodromy changes it only by an inner term.

Suppose an arithmetically normalized procyclic
subgroup has nonzero proper homology class \(v\).
Its image in the proper pro-\(\ell\) group is still
nontrivial, with the same open arithmetic normalizer
image. Restrict the number field so that its absolute
Galois group lies in that image. The local inertia
action on \(V\) is trivial, because the proper
Jacobian has good reduction. Therefore any lift of
inertia in the normalizer acts on the procyclic
subgroup by exponent one: its exponent is also the
scalar on the nonzero homology vector \(v\).
An actual nonzero power of tame inertia must thus
carry a generator \(u\) to a conjugate of itself
under any representative of its proper outer action.

We recall why this imposes a precise degree-three
condition without any integral lifting assumption.
Use the representative (A.47.2), with zero Johnson
degree one. Comparing degree two in the rational
conjugacy equation gives
\[
                 [q_1,v]=0\quad\text{in }\Lambda^2V/\mathbf Q_\ell\Omega,
                                                               \tag{A.47.4}
\]
where \(q_1\) is the first term of a conjugator.
If two positive-genus components occur, the total
genus is at least two. The decomposable bivector
\(q_1\wedge v\), of rank at most two, cannot be
a nonzero multiple of the nondegenerate rank-
\(2g\) bivector \(\Omega\). Thus \(q_1\) is
proportional to \(v\). Multiplying the rational
conjugator on the right by the corresponding inverse
rational power of \(u\) removes \(q_1\) without
changing its conjugation of \(u\). In degree three,
all terms involving higher parts of \(\log u\)
now disappear, and one obtains the necessary condition
\[
                              D(v)\in[v,\mathfrak l_2].
                                                               \tag{A.47.5}
\]
This uses a rational quotient only for a necessary
equation. It claims neither that a rational power
lies in the integral group nor that an actual
normalizer extension splits. Adding an inner
representative to \(D\) does not change (A.47.5).

Suppose \(v_i\ne0\) and \(v_j\ne0\) for distinct
components \(i,j\). Extend each of these vectors
to a symplectic basis of its block. There is a Lie
map from (A.47.3) to the free Lie algebra on \(x,y\)
which sends
\[
 v_i\longmapsto x,\quad v_j\longmapsto-x,\qquad
 \Omega_i\longmapsto[x,y],\quad
 \Omega_j\longmapsto-[x,y],                         \tag{A.47.6}
\]
and all other blocks to zero. Within either selected
block choose the companion of its indicated vector
so that the first symplectic pair maps to the shown
commutator, and kill all its other pairs. This
constructs the map even if a component has genus
greater than one. It kills both \(v\) and the
surface relation \(\Omega\).

Put \(Z=[x,y]\). Formula (A.47.2) now gives
\[
 \begin{split}
 D(v)&\longmapsto
          \left(\sum_e m_e(i_e-j_e)^2\right)[Z,x]\\
      &=\operatorname{dist}(C_i,C_j)\,[[x,y],x]\ne0.
 \end{split}                                                   \tag{A.47.7}
\]
An edge has \((i_e-j_e)^2=1\) exactly when it
separates the two components, so the coefficient is
the positive distance between distinct vertices.
The Lie word is nonzero: in the free associative
algebra it is \(2xyx-yxx-xxy\). But (A.47.6)
kills the right side of (A.47.5), since it kills
\(v\). This contradiction proves that a candidate
nonzero homology class is supported on a single
positive-genus component. The proof works for every
compact-type tree, not only for a star.

\paragraph{Apply this obstruction to the odd dihedral covers.}
Suppose that a noncuspidal nontrivial procyclic subgroup
of the geometric pro-\(\ell\) group of \(X_n\)
has open arithmetic normalizer. By A.11 a generator
has nonzero proper homology class
\[
                         v\in V_\ell\operatorname{Jac}(C_n).
\]
The actual punctured morphism
\(f=[1+i]\circ q_n:X_n\to U\) of A.43 implies
\(f_*v=0\). Indeed, if that proper image were
nonzero, the image of the procyclic subgroup in the
target pro-\(\ell\) group would be nontrivial and
would still have open arithmetic normalizer. A.12
puts it in target cusp inertia, whose proper homology
is zero, a contradiction. This step uses the actual
map of punctured curves and the induced arithmetic
homomorphism, not just a homomorphism of Jacobians.

Let \(\ell\) be odd and use the actual model at
two from A.43. Its \(n\) positive-genus components
are copies of the original supersingular elliptic
special fibre. The normalization proof identifies
the map of each such component to the main component
of the base model with the identity. Hence, under
compact-type specialization,
\[
             V_\ell\operatorname{Jac}(C_n)
                         =\bigoplus_{i=1}^n V_i,
\]
and every restriction
\(f_*|_{V_i}:V_i\to V_\ell E\) is an isomorphism:
it is the component identity followed by the
isogeny \([1+i]\). This compatibility also follows
by functoriality of the good Jacobian models. The
possible inseparability of that special isogeny at
two does not affect its prime-to-two rational Tate
module isomorphism.

A nonzero vector in \(\ker f_*\) cannot be
supported on one \(V_i\), since that restriction
is injective. Thus \(v\) has nonzero projections
onto at least two components. The compact-type
component-support obstruction just proved, with
\(p=2\), contradicts this. It excludes every
noncuspidal normalized procyclic subgroup at every
odd coefficient prime, including the primes dividing
\(n\); only nonzero rational coefficients were used.

At \(\ell=2\), A.43 proves that the proper Jacobian
has potentially good reduction at two with special
two-rank zero. The general normalizer criterion
A.46 therefore excludes a noncuspidal pro-two
subgroup directly. It applies over every finite
field extension needed to contain the arithmetic
normalizer image.

Conversely, after making a cusp rational its inertia
has a decomposition group with open arithmetic
projection. Every nonzero closed subgroup of its
pro-\(\ell\) inertia is characteristic there and
has the same open normalizer property. This proves
the exact classification in the statement.

The program \texttt{code/a\_odd\_dihedral\_normalizers.py}
checks the distance coefficient on every labelled
tree of small order, the exact two-block Lie witness,
and the component-map kernel implication for the
explicit stars. The obstruction above is uniform
in the number of components and in the tree.
\end{proof}

% Registered A.48 after checkpoint07 release; full parent proof/source audit and independent certificate PASS.
\subsection{Normalizers in every two-power dihedral cover}

\begin{statement}{A.48}{Proved}{All-prime normalizers for the four-cusp two-power family}
For every \(m=2^s\), \(s\ge1\), every coefficient
prime \(\ell\), and the four-cusp curve \(Y_m\) of
A.44, the nontrivial closed procyclic subgroups of
\(\pi_1(Y_{m,\bar K})^{(\ell)}\) with open
arithmetic normalizer are exactly the nontrivial closed
subgroups of conjugates of cusp inertia. Together with
A.44 this proves both assertions of Task A.5 for the
entire two-power dihedral family. In particular it
settles the previously displayed case
\(v^2=w^{16}-1\) at \(\ell=17\).

The proof combines the actual compact-type
component-support obstruction of A.47 with a bound
on the multiplicity of a Galois character in the
primitive cyclotomic factor. It does not require an
explicit stable graph for these higher-genus curves.
\end{statement}

\begin{proof}
\paragraph{The primitive factor in the actual double-cover tower.}
Write \(B_m:v^2=w^{2m}-1\) and
\(\tau(w,v)=(\zeta_{2m}w,v)\). All constructions
may be defined over a finite extension containing the
indicated root of unity. There is an actual finite
etale map of punctured curves
\[
 q_m:Y_m\longrightarrow Y_{m/2},\qquad
                         (w,v)\longmapsto(w^2,v),
                    \qquad m\ge4.                  \tag{A.48.1}
\]
The corresponding proper map has degree two and
ramifies only at the four removed points over zero
and infinity. Its deck involution is \(\tau^m\).
The double-cover pullback and trace on rational
homology satisfy the usual identities
\(q_{m,*}q_m^*=2\) and
\(q_m^*q_{m,*}=1+\tau^m\). Consequently
\[
 W_m:=\ker(q_{m,*}:V_\ell J(B_m)\to V_\ell J(B_{m/2}))
                 =\{v:\tau^m v=-v\}.              \tag{A.48.2}
\]
This is the primitive \(\mathbf Q(\zeta_{2m})\)
factor: after adjoining the eigenvalues of \(\tau\),
it has exactly one eigenline for each odd exponent
modulo \(2m\). Every such eigenvalue has exact
order \(2m\). The eigenvalue assertion follows also
from the differential basis
\(w^{j-1}dw/v\), \(1\le j\le m-1\), and its
complex conjugate, as in A.44.

We emphasize that (A.48.1) is used covariantly on
fundamental groups and on homology. There is no
claim that an unramified double cover gives an
inclusion of pro-\(\ell\) groups when \(\ell\)
is odd.

\paragraph{A Galois character occurs at most twice in \(W_m\).}
Assume \(m\ge4\). We prove this assertion over every
finite number field over which \(\tau\) is defined,
and after every further finite extension. Enlarge
that field, if needed, to a finite Galois field \(L\)
containing \(\mathbf Q(\zeta_{2m})\). Enlarge the
coefficient field to a finite extension of
\(\mathbf Q_\ell\) containing all \(2m\)-th roots
of unity. Since Galois commutes with \(\tau\),
it acts on each of the one-dimensional \(\tau\)
eigenspaces by a character. No simplicity assertion
for the CM factor is required.

Here is the exact bound on equality of these
characters. It is convenient first to use the dual
etale cohomology, with its functorial action by
\(\tau^*\); passing to homology inverts Galois
characters and preserves their multiplicities.
Set
\[
 \Phi_m=\{1,3,5,\ldots,m-1\}
                \subset(\mathbf Z/2m)^\times.       \tag{A.48.3}
\]
The holomorphic eigenspaces of the primitive factor
are exactly those indexed by \(\Phi_m\).
For an embedding of \(L\) into
\(\overline{\mathbf Q}_\ell\) sending the chosen
\(\zeta_{2m}\) to \(\zeta_{2m}^a\), the
Hodge--Tate weight-one eigenspaces are indexed by
\(a\Phi_m\), and their complement has weight zero.
This follows directly from the displayed differential
basis and the canonical, functorial Hodge--Tate
comparison
\cite[Theorem 2.2.3, printed p.13]{fghBrinonConrad}.
All units \(a\bmod 2m\) occur, because \(L\)
contains the cyclotomic field. The comparison applies
at every place over \(\ell\), with its labelled
field embeddings. The weights remain unchanged on
finite extensions by
\cite[Theorem 2.4.6, printed p.20]{fghBrinonConrad}.

Thus, if the two global characters indexed by
\(j,k\in(\mathbf Z/2m)^\times\) are equal on
\(G_L\), then
\[
       j\in a\Phi_m\quad\Longleftrightarrow\quad
       k\in a\Phi_m\qquad\text{for every }a.
                                                               \tag{A.48.4}
\]
Equivalently \(kj^{-1}\) stabilizes \(\Phi_m\)
under multiplication. This is a necessary condition
for equality, which suffices for the bound; no
converse or assertion about finite-order twists is
needed.

The stabilizer is exactly
\[
       \operatorname{Stab}(\Phi_m)=\{1,m-1\}.       \tag{A.48.5}
\]
To prove this without a CM-type classification, put
\(z=\zeta_{2m}\). The finite geometric series gives
\[
               \sum_{j\in\Phi_m}z^j
                         =\frac{2z}{1-z^2}.
\]
If multiplication by \(a\) preserves \(\Phi_m\),
the same sum with \(z\) replaced by \(z^a\)
is equal to it. Hence
\(z^a-z^{-a}=z-z^{-1}\). For \(t=z^a\),
this factors as
\[
                         (t-z)(t+z^{-1})=0.
\]
Therefore \(a=1\) or \(a=m-1\pmod{2m}\).
Both do preserve \(\Phi_m\): for odd
\(1\le j<m\), multiplication by \(m-1\)
gives \(m-j\bmod 2m\). This proves (A.48.5).
Together with (A.48.4), it bounds the multiplicity
of every global Galois character in \(W_m\) by two,
even after any finite extension. In particular, if
\(0\ne v\in W_m\) spans a Galois-stable rational
line, then
\[
          \dim_{\mathbf Q_\ell}
              \operatorname{span}\{\tau^k v:k\in\mathbf Z\}
                              \le2.                \tag{A.48.6}
\]
Indeed all nonzero eigencomponents of \(v\) must
carry its single Galois character. Extending the
coefficient field for this argument does not change
the dimension of the rational span.

\paragraph{A wild cyclic action forces a large genus.}
Let \(Z\) be a smooth proper positive-genus curve
over an algebraically closed field of characteristic
two, with two-rank zero, and let a cyclic group of
order \(N=2^r\), \(r\ge1\), act faithfully on it.
Then
\[
                         g(Z)\ge\frac{(N-1)(N-2)}6.
                                                               \tag{A.48.7}
\]
For completeness, the Deuring--Shafarevich formula
\cite[formula (6), printed p.5; Lemma 3.2 and
Remark 3.1, printed pp.7--8]{fghGiuliettiKorchmaros}
shows that this action has one short orbit, consisting
of one fixed point. Directly, if \(f'\) is the
quotient two-rank and the short orbit sizes are
\(d_j\), then
\[
                    -1=N(f'-1)+\sum_j(N-d_j).
\]
It forces \(f'=0\). Every short orbit contributes
at least \(N/2\), so there is exactly one, of
size one. There is thus exactly one ramification
point, with full inertia group.

Let \(u_1,\ldots,u_r\) be its upper jumps.
By \cite[Theorem 1.1(c), manuscript p.2]{fghObusPries},
\(u_1\ge1\) and \(u_i\ge2u_{i-1}\), so
\(u_i\ge2^{i-1}\). Hilbert's different formula
and the upper-to-lower conversion in
\cite[Section 2.4, manuscript p.4]{fghObusPries}
give, after telescoping,
\[
 \delta=(N-1)+\sum_{i=1}^r2^{i-1}u_i
       \ge (N-1)+\sum_{i=1}^r4^{i-1}
        =\frac{N^2+3N-4}{3}.                       \tag{A.48.8}
\]
Riemann--Hurwitz, with nonnegative quotient genus,
gives \(2g(Z)-2\ge-2N+\delta\), proving
(A.48.7). The same argument includes a supersingular
elliptic component; it does not require genus at
least two in the Deuring--Shafarevich step.

\paragraph{The induction and its component-orbit contradiction.}
The assertion for \(m=2\) is A.40, and that for
\(m=4\) is the stable \(C_4\) case of A.14.
At coefficient prime two it holds for every \(m\)
by A.44 and A.46. Fix an odd \(\ell\), assume
\(m\ge8\), and suppose the assertion known for
\(m/2\).

Suppose a noncuspidal procyclic subgroup of the
geometric pro-\(\ell\) group of \(Y_m\) has open
arithmetic normalizer. A.11 gives a nonzero proper
homology vector \(v\). Over a finite field
extension contained in the normalizer's arithmetic
image, its rational line is Galois-stable. This
uses only the action on homology: geometric inner
automorphisms act trivially there, so no splitting
of the normalizer extension is assumed.

The actual map (A.48.1) and the induction hypothesis
force \(q_{m,*}v=0\). Otherwise the image procyclic
subgroup is nontrivial and has the same open
arithmetic-normalizer property; in the target it is
cuspidal and hence has zero proper homology, a
contradiction. Thus \(v\in W_m\), so (A.48.6)
applies.

At two choose a split stable model over a common
finite extension. A.44 makes its graph a tree, and
every positive-genus component has two-rank zero,
since the good special Jacobian is the product of
the component Jacobians and has two-rank zero. The
compact-type obstruction of A.47, with coefficient
\(\ell\ne2\), says that \(v\) is supported
on a single positive-genus component \(Z\).
The automorphism \(\tau\) extends to the stable
model and permutes the corresponding orthogonal
homology blocks. This is compatible with the actual
good-Jacobian specialization.

Let \(d\) be the length of the orbit of \(Z\)
under \(\tau\). The first \(d\) translates of
\(v\) are nonzero vectors in different homology
blocks, hence linearly independent. Formula
(A.48.6) gives \(d\le2\), and so \(d=1\)
or \(d=2\).

The restriction of \(\tau^d\) to \(Z\) has
exact order
\[
                              N=2m/d.              \tag{A.48.9}
\]
It has order dividing this number because \(\tau\)
has order \(2m\). Conversely every nonzero
\(\tau\)-eigencomponent of \(v\in W_m\)
has primitive \(2m\)-th eigenvalue, so
\((\tau^d)^k v=v\) is possible only when
\(N\mid k\). An identity automorphism of \(Z\)
would fix its whole homology block and hence \(v\).
This proves the exact order assertion without
assuming that \(\tau\) acts faithfully on every
component of the model.

All \(d\) components in this orbit have the same
genus. By (A.48.7), their contribution to the total
genus \(m-1\) is at least
\[
 \begin{cases}
  \displaystyle\frac{(2m-1)(2m-2)}6,&d=1,\\[4pt]
  \displaystyle\frac{(m-1)(m-2)}3,&d=2.
 \end{cases}                                                   \tag{A.48.10}
\]
Both exceed \(m-1\) when \(m\ge8\): their
differences are respectively
\(2(m-1)(m-2)/3\) and \((m-1)(m-5)/3\).
This contradiction proves the induction step.
It excludes all noncuspidal candidates at every odd
coefficient prime, including primes at which the
cyclotomic factor does not split over
\(\mathbf Q_\ell\).

Every nonzero closed subgroup of a cusp inertia
group has open arithmetic normalizer after making
the cusp rational, as in A.47. This proves the
classification. For the concrete previously open
case \(m=8\), an orbit of one would require a
component of genus at least \(35\), and an orbit
of two would require two components of genus at
least \(7\), whereas the total genus is \(7\).

The exact certificate \texttt{code/a\_two\_power\_normalizers.py} checks
the differential characters and primitive kernel,
the interval stabilizer and labelled weight-profile
multiplicities, the Hilbert--Herbrand telescoping
identity, and both genus contradictions. The proof
itself is uniform in \(m\).
\end{proof}

% Registered A.49 after full independent root proof/source audit PASS.
\subsection{Complete cover calculations in every proper genus}

\begin{statement}{A.49}{Proved}{A supersingular Kummer family in every genus}
Let \(n\ge2\), choose a prime \(p\equiv3\pmod4\)
with \(p\nmid n\), and let \(a\ge1\). On
\(E:y^2=x^3-x\), choose a point \(Q\) of exact
order \(p\) and a rational function \(f\) with
\[
                         \operatorname{div}(f)=p(O-Q).
                                                               \tag{A.49.1}
\]
All objects are defined over a suitable finite number
field. Let \(C\) be the smooth normalization of
\(z^n=f\), let \(q:C\to E\) be its projection,
and remove the full inverse image of \(E[p^a]\):
\[
                  X=C\setminus q^{-1}(E[p^a]).       \tag{A.49.2}
\]
Then \([p^a]\circ q:X\to U\) is an actual
connected finite etale cover of degree \(np^{2a}\),
with proper genus \(n\) and exactly
\(2+n(p^{2a}-2)\) cusps.

For every coefficient prime, every selected cusp
packet, and every indicated open derived subgroup,
its full geometric pro-prime fixed subgroup is the
selected cusp inertia. For every coefficient prime,
its nontrivial closed procyclic subgroups with open
arithmetic normalizer are exactly the nontrivial
closed subgroups of cusp inertia conjugates.
Thus both assertions of Task A.5 hold in this family
in every proper genus \(n\ge2\), including even
and composite \(n\). The resulting higher-genus
Jacobian need not be assumed to have complex
multiplication.
\end{statement}

\begin{proof}
\paragraph{The actual number-field cover.}
Take a finite extension over which \(Q\), the
\(p^a\)-torsion, and the necessary roots of unity
are defined. Since \([p](O-Q)=0\) in the
elliptic Picard group, the degree-zero divisor
\(p(O-Q)\) is principal. A defining function is
available over a further finite extension; this
justifies the actual rational function in (A.49.1).
The Kummer cover is connected and has degree \(n\),
because the valuation of \(f\) at \(O\) is
\(p\), coprime to \(n\). It is totally ramified
of index \(n\) at \(O,Q\), and unramified
elsewhere. Riemann--Hurwitz gives
\[
                       2g(C)-2=2(n-1),
                      \qquad g(C)=n.                \tag{A.49.3}
\]
In characteristic zero \([p^a]\) is an etale
isogeny of degree \(p^{2a}\). Its origin fibre
is precisely \(E[p^a]\), so removing (A.49.2)
removes all branch points and precisely the origin
fibre of the composite. There is one point over
each of \(O,Q\), and \(n\) points over every
other point of \(E[p^a]\), proving the degree
and cusp count. The Euler characteristic checks
\[
       2-2n-\bigl(2+n(p^{2a}-2)\bigr)=-np^{2a}.
                                                               \tag{A.49.4}
\]
There are infinitely many primes congruent to three
modulo four, by the usual Euclidean proof using
\(4\prod p_i-1\). Hence a suitable \(p\)
exists for every \(n\), and the construction is
not restricted to a congruence class of genera.

\paragraph{The actual model and all its extra cusp sections.}
The displayed elliptic equation has good reduction
at every odd prime. At our chosen \(p\), its
reduction is supersingular: the coefficient criterion
in \cite[Theorem V.4.1(a) and Example V.4.5]{fghSilvermanAEC}
applied to
\(x^{(p-1)/2}(x^2-1)^{(p-1)/2}\) gives zero
coefficient at \(x^{p-1}\), because
\((p-1)/2\) is odd. Thus the special elliptic
curve has \(p\)-rank zero. Every point of
\(E[p^a]\) specializes to its origin: its
specialization is geometric \(p\)-power torsion,
and the special elliptic curve has no nonzero such
points. In particular \(O\) and \(Q\) collide.

Apply the actual two-branch normalization proved in
A.43, with \(d=p\) and \(\gcd(n,p)=1\).
We recall its map, not just its abstract graph.
Starting from the good elliptic model, separate
\(O,Q\) by a chain of blowups. The result consists
of the elliptic main component and a rational chain,
whose last component carries the two distinct marks.
The graph-Laplacian calculation makes the vertical
orders of \(f\) constant. After removing this
common vertical factor and making a finite extension,
its reduction on the main component and the
intermediate chain is a nonzero constant with an
\(n\)-th root. Hensel's lemma, since \(n\) is
invertible in the residue field, splits the actual
normalization there into \(n\) copies.

On the final rational component choose \(u\) with
the branch marks at zero and infinity and the
attaching point at one. Up to a constant the
normalization is
\[
                  z^n=u^p,\qquad u=s^n,
                  \qquad z=s^p.                   \tag{A.49.5}
\]
Bezout's identity for \(n,p\) proves this is the
normalization, and it is a rational component.
There are \(n\) distinct points over \(u=1\).
Thus the actual normalization has \(n\)
supersingular elliptic components, each mapping
isomorphically to the elliptic main component,
joined through rational chains to one rational
centre carrying the two branch marks. The map
\(q\) extends on this normalization to the
blown-up base model, and composing with its
blowdown gives an actual morphism to the original
good elliptic model. Contracting unnecessary chains
is not required in this argument.

It remains to check every extra cusp in (A.49.2),
not only the two branch points. Its image is a
point of \(E[p^a]\), so its special image in
the original good elliptic model is the origin.
On each positive-genus component the extended
map to that model is the isomorphism just described.
The unique point of this component mapping to the
origin is exactly its attaching node with the
rational part of the model. Hence an extra cusp
can specialize only on the rational part or at
one of these attaching nodes. It cannot specialize
at any other smooth point of a positive-genus
component.

After a further finite extension make all cusp
sections rational, and separate them in this
semistable model. These are modifications supported
on the rational part and its attaching nodes.
Separating sections at a smooth rational point
adds rational trees; separating sections at a node
after semistable base change inserts rational chains
and their rational branches. The original elliptic
components persist, and no modification at any other
point of them is needed. Equivalently, those \(n\)
elliptic components already account for all genus
in (A.49.3), so semistable separation cannot introduce
a further positive-genus component or a cycle.
The rational part remains connected, with each
elliptic component attached as a tail. Every cusp
is a smooth mark on that rational part. Thus the
path between any two cusp vertices contains only
rational components, including both endpoints.
This establishes the precise pointed-model hypothesis
for every selected cusp, not merely for the branch
pair.

The proper special Jacobian is the product of these
\(n\) supersingular elliptic Jacobians. It has
good reduction and \(p\)-rank zero. A.41 and
A.42 now give the full fixed-subgroup assertion at
every coefficient prime, for every cusp packet and
every indicated open derived subgroup.

\paragraph{All normalizers, using the same actual component maps.}
Suppose a noncuspidal nontrivial procyclic subgroup
has open arithmetic normalizer. By A.11 its proper
homology has a nonzero generator \(v\), after
restricting the arithmetic field if necessary.
The actual punctured map
\([p^a]\circ q:X\to U\) and A.12 imply that
its proper homology image is zero: a nonzero image
would yield a normalized nontrivial procyclic
subgroup of the target with nonzero proper homology,
contrary to target cuspidality.

If \(\ell\ne p\), the compact-type specialization
gives
\[
                     V_\ell J(C)=\bigoplus_{i=1}^nV_i.
\]
On every component the proper map to the target
elliptic curve is its component isomorphism followed
by \([p^a]\). Thus its map on each \(V_i\)
is a rational Tate-module isomorphism. A nonzero
vector in the total kernel cannot be supported
on a single component. This contradicts the
component-support obstruction of A.47, applied at
residue characteristic \(p\). Its proof is
independent of the residue characteristic and only
requires \(\ell\ne p\).

If \(\ell=p\), the good special Jacobian's
\(p\)-rank zero lets A.46 exclude a noncuspidal
normalized subgroup directly. Conversely every
nonzero closed subgroup of a cusp inertia group
has open arithmetic normalizer after making that
cusp rational. This proves the exact classification.

\paragraph{A concrete even-genus instance.}
Take \(p=3\), \(a=1\), and \(n=4\).
Choose \(Q=(q,r)\) with
\[
          3q^4-6q^2-1=0,
          \qquad r^2=q^3-q,
          \qquad \lambda=\frac{3q^2-1}{2r}.
\]
Then \(Q\) has order three and its tangent
line has divisor \(3Q-3O\). In fact
\[
       x^3-x-\bigl(r+\lambda(x-q)\bigr)^2=(x-q)^3,
                                                               \tag{A.49.6}
\]
using the two displayed relations. One may take
\(f=(y-r-\lambda(x-q))^{-1}\). The smooth
normalization of \(z^4=f\), with the inverse
image of \(E[3]\) removed, has genus four,
\(30\) cusps, and degree \(36\) over \(U\).
Both fixed groups and normalizers have the asserted
classification at every coefficient prime. The
accompanying exact certificate verifies this tangent
identity, the ramification and Euler counts, and
the tame normalization exponents for sample parameters;
the preceding proof applies to every \(n,p,a\)
in the statement.
\end{proof}

% Registered A.50: full independent root and PBC proof/source/code audits PASS.
\subsection{The degree-twenty-seven Heisenberg cover}

\begin{statement}{A.50}{Proved}{All normalizers in an explicit genus-ten cover}
Let \(C\) be the smooth proper normalization of
\[
 y^2=x^3-x,\qquad z^3=\psi_3(x):=3x^4-6x^2-1,
 \qquad q:C\longrightarrow E.                         \tag{A.50.1}
\]
Put \(S=q^{-1}(E[3])\) and \(X=C\setminus S\).
Then \([3]q:X\to U\) is a connected geometric Galois
cover with Heisenberg deck group of order \(27\),
proper genus \(10\), and nine cusps.
For every coefficient prime, the nontrivial closed
procyclic subgroups of its geometric pro-prime group
with open arithmetic normalizer are precisely the
nontrivial closed subgroups of cusp inertia conjugates.
Its full pro-three fixed subgroup under every indicated
open derived group and cusp packet is the selected
cusp inertia, by A.22.

The proper Jacobian has an isogeny decomposition
\[
              J_C\sim E\times E_0^3\times B^3,
       \qquad E_0:3s^2=z^3+4,\quad \dim B=2,           \tag{A.50.2}
\]
after finite number-field extension. The surface
\(B\) has an action by \(\mathbf Q(\sqrt{-3})\)
and has no geometric CM abelian subvariety.
The \(E_0^3\) part is detected by actual punctured
maps, not merely by correspondences of Jacobians.
The curve also has potentially good reduction at two
with special two-rank zero.
No fixed-subgroup assertion at coefficient primes
other than three is inferred here from the normalizer
calculation.
\end{statement}

\begin{proof}
\paragraph{The actual cover and its deck group.}
The division polynomial has divisor
\[
 \operatorname{div}_E(\psi_3)
       =\sum_{T\in E[3]\setminus\{O\}}(T)-8(O).
                                                               \tag{A.50.3}
\]
Its eight finite zeroes and its pole have orders prime
to three. Thus \(q\) is connected of degree three,
totally ramified at precisely the nine points of
\(E[3]\). Each has a unique preimage. Riemann--Hurwitz
therefore gives \(2g(C)-2=9(3-1)=18\).
The composite with the etale degree-nine isogeny
\([3]\) ramifies only above \(O\). Removing its
full fibre gives the asserted degree and nine cusps.

Here is an explicit construction of all the deck lifts.
Let \(T=(a,b)\in E[3]\setminus\{O\}\), so
\(3a^4-6a^2-1=0\). The quotient isogeny with kernel
\(\langle T\rangle\) has coordinates
\[
 \begin{split}
 X_a&=\frac{2(x^3-x+a^3-a)}{(x-a)^2}-x-4a,\\
 Y_a&=y\frac{dX_a}{dx},\\
 E_a&:Y_a^2=X_a^3+(9-30a^2)X_a+42a-70a^3.
 \end{split}                                                   \tag{A.50.4}
\]
The first expression is the sum-of-translates formula
\(x(P)+x(P+T)+x(P-T)-2x(T)\); hence it is invariant
under translation by \(T\). Direct substitution gives
\[
               \psi_3(x)=3(x-a)^3(X_a+3a).          \tag{A.50.5}
\]
Consequently translation by \(T\) lifts to
\[
       \widetilde T(P,z)=
       \left(P+T,\;z\frac{x(P+T)-a}{x(P)-a}\right).
                                                               \tag{A.50.6}
\]
The product of the three displayed ratios telescopes,
so \(\widetilde T^3=1\). This formula is a rational
function-field automorphism and therefore extends to
\(C\). It commutes with the central automorphism
\(c:z\mapsto\zeta_3z\).

Choose \(T\ne O\) and put \(T'=iT\). The two
points form a basis of \(E[3]\), since \(i\) has
irreducible polynomial \(t^2+1\) over \(\mathbf F_3\).
Their lifts and \(c\) give all \(27\) deck
automorphisms of \([3]q\): their images contain all
nine translations, and their central kernel contains
\(\mu_3\). This group cannot be abelian. Otherwise
the inertia about the sole branch point on \(E\)
would be trivial, since that inertia is a commutator
in the once-punctured elliptic fundamental group.
It instead has order three by (A.50.3).
The two lifts have order three, their commutator is
central of order three, and the usual class-two
power formula shows every element has order dividing
three. This is precisely the Heisenberg group
\(H_3\) of order \(27\).
All these maps are defined over
\(L=\mathbf Q(i,E[3])\), which contains \(\zeta_3\)
by the Weil pairing; enlarging \(L\) later is harmless.
A.22 now gives the stated full pro-three fixed subgroup.

\paragraph{An actual once-punctured CM target.}
The order-four automorphism
\[
               \tau(x,y,z)=(-x,iy,z)
                                                               \tag{A.50.7}
\]
has quotient map
\[
          q_0:C\longrightarrow E_0,
          \qquad (x,y,z)\longmapsto(z,s=x^2-1),
          \qquad 3s^2=z^3+4.                     \tag{A.50.8}
\]
The identity is immediate from (A.50.1). Generically
there are four choices of \((x,y)\), permuted
transitively by \(\tau\), so this is the exact
quotient of degree four. With \(O\) its point at
infinity, \(E_0\) has full CM by
\(\mathbf Z[\zeta_3]\); equivalently it is
\(Y^2=X^3+108\) under \(X=3z,Y=9s\).
Its automorphism \([\zeta_3]\) is
\((z,s)\mapsto(\zeta_3z,s)\).

The three points
\[
          O,\quad (0,2/\sqrt3),\quad(0,-2/\sqrt3)
\]
are exactly \(\ker[1-\zeta_3]\). The two finite
points have four preimages each, the eight zeroes in
(A.50.3); infinity has the unique remaining cusp,
with ramification index four. Hence
\[
      [1-\zeta_3]\circ q_0:X\longrightarrow E_0\setminus\{O\}
                                                               \tag{A.50.9}
\]
is an actual punctured morphism of degree twelve:
its entire origin fibre is precisely \(S\).
Every composite with a deck automorphism of \(X\)
has the same property.

We also record why the target normalizer theorem
holds at every prime in this CM case. The proof of
A.12 applies with \(\mathbf Q(\sqrt{-3})\) in
place of its two listed fields. At a nonsplit prime,
including two and the ramified prime three, an
ordinary good Frobenius and all its positive powers
generate the quadratic CM algebra over
\(\mathbf Q_\ell\). Its eigenvalue quotient has
infinite order, as their ideals have different
support. Thus no rational proper homology line is
stable under an open Galois subgroup, and A.11 applies.
Every split prime is at least seven. The hypotheses
of \cite[Theorem 4.5 and Corollary 4.10]{Ishii2025}
are satisfied: class number one, the full integer
ring, and potentially ordinary reduction at the
split coefficient prime. The index \((2,2)\) obeys
the unit congruence modulo six. Those results make
\(\kappa_{(2,2)}\) nonzero. Nakamura's actual
metabelian formula
\cite[Theorem A and Corollary 4.12]{Nakamura1995}
and the three distinct CM weights isolate the same
middle special derivation \(T\) of A.12, for which
\(T(x)\notin[x,L_4]\) and \(T(y)\notin[y,L_4]\).
The actual lower-inner-term removal, torus projection,
and necessary normalizer equation are exactly those
proved there. They exclude both split eigenlines.
Nonvanishing persists after every finite field
extension. Thus (A.50.9) has the required all-prime
normalizer target, including coefficients two and three.

\paragraph{The rational isotypic factors and the exact span of the maps.}
We give the representation calculation because merely
counting an elliptic quotient would not detect the
correct multiplicity. Let \(F=\mathbf Q(\zeta_3)\).
The nontrivial-central part of \(\mathbf Q[H_3]\)
is \(M_3(F)\). Explicitly its three-dimensional
representation over \(F\) is generated by
\[
 U=\operatorname{diag}(1,\zeta_3,\zeta_3^2),
       \qquad V e_j=e_{j+1}\quad(j\bmod3),
       \qquad UV=\zeta_3VU.                       \tag{A.50.10}
\]
The remaining rational summands are the characters
of \(H_3/\langle c\rangle\). Their representation
on \(H_1(C,\mathbf Q)\) factors through
\(H_1(C/\langle c\rangle,\mathbf Q)=H_1(E,\mathbf Q)\),
on which the quotient group acts by translations and
therefore trivially. Since \(\dim H_1(C)=20\), the
remaining part has rational dimension eighteen,
or three copies of the six-dimensional rational
simple module belonging to \(M_3(F)\).
Rational matrix idempotents therefore give
\(J_C\sim E\times B_3^3\), where
\(\dim B_3=3\) and \(F\) acts on \(B_3\).
These are actual isogeny factors: clear denominators
in the idempotents and take their images in \(J_C\).

The lift (A.50.6) and its \(\tau\)-conjugate can
be chosen so that \(\tau\) sends them to one
another and sends the second to the inverse of the
first. Indeed \(\tau^2\) lies over \([-1]\),
fixes \(x,z\), and conjugates (A.50.6) to its
inverse. In (A.50.10) the induced automorphism is
therefore \(U\mapsto V,V\mapsto U^{-1}\), after
choosing the central generator consistently.
A convenient intertwiner is
\[
 A=\frac1{-\sqrt{-3}}
       (\zeta_3^{-jk})_{0\le j,k\le2}.
                                                               \tag{A.50.11}
\]
Direct multiplication gives
\[
 AUA^{-1}=V,\quad AVA^{-1}=U^{-1},\quad A^4=1,
 \quad\det(t-A)=(t-1)(t^2+1).                   \tag{A.50.12}
\]
On the multiplicity space of dimension three over
\(F\), write \(\tau=A\otimes B_\tau\).
This is possible by the matrix-algebra commutant;
\(B_\tau^4=1\), and its characteristic polynomial
has coefficients in \(F\). Equivalently
\(A^{-1}\tau\) is an actual rational endomorphism
commuting with \(M_3(F)\), hence an endomorphism
of the isogeny factor \(B_3\).

The \(\tau\)-invariant space is precisely
\(q_0^*H_1(E_0,\mathbf Q)\). It has dimension
one over \(F\), because the central automorphism
on this elliptic quotient is multiplication by
\(\zeta_3\). Let \(n_1,n_{-1},k,k\) be the
multiplicities of \(1,-1,i,-i\) in \(B_\tau\).
The last two are equal because \(i\notin F\).
From (A.50.12), the invariant dimension is
\(n_1+2k=1\), while \(n_1+n_{-1}+2k=3\).
Thus
\[
                         n_1=1,\quad n_{-1}=2,\quad k=0.
                                                               \tag{A.50.13}
\]
In particular the invariant vector is a pure tensor:
the one-eigenline of \(A\) tensored with the
one-eigenline of \(B_\tau\). Its \(H_3\)-span
is exactly one copy of the three-dimensional
\(F\)-module, rather than an unspecified larger
isotypic span. The corresponding abelian subvariety
is isogenous to \(E_0^3\), since it is generated
by the images of the actual maps
\(q_0g\), \(g\in H_3\). The complementary
minus-one multiplicity factor has dimension two;
call it \(B\). It inherits the \(F\)-action.
This proves (A.50.2) and the exact dimension of the
homology detected by the punctured maps (A.50.9).
All projectors and maps are defined over \(L\).

\paragraph{A genus-four quotient gives exact Frobenius certificates.}
Fix the lift \(\widetilde T\) in (A.50.6), with
\(x(T)=a\), and set
\(u=z/(x-a)\), \(\eta=Y_a\).
Equations (A.50.4)--(A.50.5) give the actual quotient
\(D=C/\langle\widetilde T\rangle\) as
\[
 \eta^2=\frac{u^9}{27}-a u^6+(3-a^2)u^3+15a-7a^3.
                                                               \tag{A.50.14}
\]
It has genus four. For instance the noncentral
order-three subgroup acts freely on \(C\), since
all inertia in the full cover is central, and
Riemann--Hurwitz gives that genus. Its invariant
space in each three-dimensional representation
(A.50.10) has dimension one. Hence the actual
isogeny decomposition over \(L\) is
\[
                              J_D\sim E\times E_0\times B.
                                                               \tag{A.50.15}
\]
This identification uses the same matrix idempotents
and the same surface \(B\), not an unrelated
surface reconstructed from point counts.

The primes \(61\) and \(313\) split in \(L\):
\(i\) is rational modulo either, and the exact
certificate verifies that all eight nonzero points
of \(E[3]\) are rational over each residue field.
Choose places sending \(a\) to \(4\) and
\(26\), respectively. All curves and maps just
used have good reduction there. Direct exhaustive
counts of (A.50.14) over both the prime field and
its quadratic extension give
\[
\begin{array}{c|r|r|r|r|r}
p&a&\#D(\mathbf F_p)&\#D(\mathbf F_{p^2})
  &\operatorname{tr}(E)&\operatorname{tr}(E_0)\\\hline
61&4&96&3696&-10&-13\\
313&26&312&97872&26&-13
\end{array}                                                     \tag{A.50.16}
\]
The certificate constructs each quadratic field as
\(\mathbf F_p[t]/(t^2-d)\) with a specified
nonsquare \(d\), evaluates the quadratic character
at every argument, and includes the single point
at infinity of the odd-degree hyperelliptic model.
It separately verifies separability of that model
and all curve-map identities.

Subtract the two elliptic traces and their square
traces from (A.50.16). Reciprocity then determines
the two degree-four Frobenius polynomials of the
same surface \(B\):
\[
\begin{aligned}
 P_{61}(t)&=t^4+11t^3+60t^2+671t+3721,\\
 P_{313}(t)&=t^4+11t^3-192t^2+3443t+97969.
\end{aligned}                                                  \tag{A.50.17}
\]
Both have the form
\[
                  t^4+b t^3+(b^2-p)t^2+pb t+p^2,
                         \qquad b=11.              \tag{A.50.18}
\]
If \(\alpha,\bar\alpha\) solve
\(t^2-bt+p=0\), its four roots are
\(\zeta_3\alpha,\zeta_3^2\alpha,
  \zeta_3\bar\alpha,\zeta_3^2\bar\alpha\).
Thus the third-power Frobenius polynomial is
\[
             \bigl(t^2-(b^3-3pb)t+p^3\bigr)^2.
                                                               \tag{A.50.19}
\]
The two ordinary imaginary quadratic fields are
\[
             \mathbf Q(\sqrt{-123}),\qquad
             \mathbf Q(\sqrt{-1131}),               \tag{A.50.20}
\]
which differ. These fields survive every further
positive power: \(p\nmid b\), so the ideals of
\(\alpha\) and \(\bar\alpha\) are supported
at different primes above \(p\). Their ratio is
not a root of unity. Consequently \(\alpha^r\)
is nonreal and generates its same imaginary
quadratic field for every \(r>0\).

\paragraph{Why these counts exclude every CM factor of the surface.}
We spell out the required implication. A surface
with an action by an imaginary quadratic field
\(F\) which has any geometric CM subvariety
must itself be isogenous to the square of a CM
elliptic curve, unless it is geometrically simple
of CM type. In the nonsimple case, a CM elliptic
factor and a non-CM elliptic factor would make the
\(F\)-action project to an impossible unital
embedding \(F\hookrightarrow\mathbf Q\).
Two nonisogenous CM elliptic factors would force
both endomorphism fields to equal \(F\), making
them isogenous after all. The only possibility
is therefore an elliptic square with CM by some
imaginary quadratic field \(F'\).

The geometrically simple CM alternative is impossible
in dimension two with this \(F\)-action. Its
endomorphism field would be a quartic CM field
containing \(F\), hence the compositum of \(F\)
with its real quadratic subfield, a biquadratic
field. Every CM type of a biquadratic CM field is
induced from one of its imaginary quadratic
subfields. The rational Hodge structure is then
two copies of that of a CM elliptic curve, giving
an elliptic-square isogeny after clearing lattice
denominators, as in the Hodge argument of A.3.
This contradicts geometric simplicity.

Suppose, then, that \(B\) had any CM factor.
Over a finite number field it would be isogenous
to \(F_1^2\), with \(F_1\) an elliptic curve
with CM by a fixed field \(F'\). At each of
(61) and (313), pass to a residue extension which
defines this isogeny and all endomorphisms, and
then take a further multiple of three. Formula
(A.50.19) gives ordinary elliptic Frobenius roots
\(\alpha^r,\bar\alpha^r\). The endomorphism
field of an ordinary CM elliptic reduction is
its imaginary quadratic Frobenius field, by
\cite[Theorem V.3.1]{fghSilvermanAEC}; since
it contains \(F'\), it equals \(F'\).
By (A.50.20) and its persistence under powers,
this would force \(F'\) to equal two different
fields. This contradiction proves that \(B\)
has no geometric CM subvariety.

It follows at every coefficient prime that
\(V_\ell(B^3)\) has no \(\mathbf Q_\ell\)-line
stable under an open absolute Galois subgroup.
Indeed such a line
is fixed by that subgroup's closed commutator.
Scale a nonzero vector into its Tate lattice and
reduce modulo \(\ell^n\); this gives torsion
of unbounded order over the maximal abelian
extension of the corresponding finite number
field. \cite[\S1, Corollary 1]{Ruppert1998}
would then give a CM subvariety, which has just
been excluded. No assertion about a conjectural
monodromy group of \(B\) is needed.

\paragraph{All procyclic normalizers.}
Suppose a noncuspidal nontrivial procyclic subgroup
has open arithmetic normalizer. A.11 supplies a
nonzero proper homology vector \(v\) spanning
a rational line stable under an open Galois
subgroup. The actual map \([3]q:X\to U\)
and A.12 make its \(E\)-projection zero.
Every actual map (A.50.9), precomposed with a deck
automorphism, has the CM target recognition proved
above. Each therefore kills the corresponding
proper homology image of \(v\): otherwise the
image procyclic subgroup would be nontrivial,
would retain open arithmetic normalizer, and
would have nonzero proper homology, contrary to
cuspidality on the target.

The exact span calculation (A.50.13) says that these
maps detect all of \(E_0^3\). Thus \(v\) lies
in \(V_\ell(B^3)\), which has no such rational
line. This is a contradiction for every prime,
including three. Notice that A.22 was used only
for the separate full fixed-group assertion at
three, not as a normalizer theorem. Conversely
all nonzero closed cusp inertia subgroups have
open arithmetic normalizer after finite extension
making the cusp rational. This proves the exact
all-prime normalizer classification.

\paragraph{An independent residue-two check.}
After a finite extension the original \(E\) has
good supersingular reduction at two. Its three-torsion
sections remain distinct. Formula (A.50.3), after
removing a constant vertical factor and a finite
extension extracting its cube root, gives a tame
cyclic-three normalization with smooth reduction:
locally at each branch section its order is prime
to three, and elsewhere it is etale. Thus \(C\)
has good reduction. Its order-four automorphism
extends to this model. The special \([-1]\) on
\(E\) has only the origin as a fixed geometric
point, because the elliptic special fibre has
two-rank zero. There is just one point of \(C\)
over that origin. Hence \(\tau^2\), and also
\(\tau\), have exactly this one fixed point;
\(\tau^2\) is nontrivial since it lies over
the nontrivial negation of the elliptic curve.
The special \(\tau\)-action therefore has
order four and just one short orbit, of size one.

Take also a good model of \(E_0\), whose reduction
at two is supersingular. The actual map \(q_0\)
extends to it by the Neron mapping property.
The pullback of a line bundle of degree one has
degree four on every fibre, so the special map is
nonconstant of degree four. It is invariant under
the faithful special \(C_4\)-action. Factoring
through its degree-four function-field quotient
leaves a map of degree one, so that quotient is
exactly the supersingular \(E_0\), without an
inseparable remaining factor. The Deuring--Shafarevich
formula \cite[formula (6), p.~5]{fghGiuliettiKorchmaros}
gives
\[
                  f_C-1=4(0-1)+(4-1)=-1.
                                                               \tag{A.50.21}
\]
Thus \(f_C=0\), and A.46 independently proves
the pro-two normalizer assertion. The nine cusps
are distinct on this good special curve, so this
calculation does not meet the rational joining-path
hypothesis of A.42 and is not promoted to a
fixed-group conclusion at other primes.
\end{proof}

% Registered A.51: full independent root and DE proof/source/code audits PASS.
\subsection{Quaternion multiplication completes the Heisenberg fixed groups}

\begin{statement}{A.51}{Proved}{Full fixed inertia in the genus-ten Heisenberg cover}
For the curve \(X=C\setminus S\) of A.50,
every coefficient prime, every selected cusp packet,
and every indicated open derived subgroup, the full
geometric pro-prime fixed subgroup is the selected
cusp inertia. Together with A.50, both assertions
of Task A.5 hold for this degree-twenty-seven cover.

More precisely, the same surface \(B\) in
(A.50.2) has rational quaternion multiplication by
\[
                         (-3,2)_{\mathbf Q}.
                                                               \tag{A.51.1}
\]
The Jacobian \(J_C\) has potentially good
reduction at three with special three-rank zero,
and every difference of two of the nine cusp
points is annihilated by nine in \(J_C\).
\end{statement}

\begin{proof}
\paragraph{Identify the polarized surface by an exact cellular calculation.}
We compute the rational cup form, not a numerical
period approximation. Write
\[
 H_3=\{(a,b,c):a,b,c\in\mathbf F_3\},\qquad
 (a,b,c)(a',b',c')=(a+a',b+b',c+c'+ab'),
\]
and put \(x=(1,0,0)\), \(y=(0,1,0)\),
\(c=(0,0,1)=[x,y]\). The Heisenberg quotient
of the rank-two free group is its class-two,
exponent-three quotient. In particular every
epimorphism onto \(H_3\) from two handle
generators has the same kernel up to the evident
deck labelling. The topological cover underlying
(A.50.1), whose degree and deck group were proved
there, is therefore the following square cover
of the square torus.

Label the \(27\) oriented squares by \(g\in H_3\).
The right neighbour is \(gx\), and the upper
neighbour is \(gy\). There are nine original
vertices, labelled by \(g\bmod\langle c\rangle\).
Write \(v(g)\) for such a vertex. The oriented
horizontal and vertical edges are
\[
 h_g:v(g)\longrightarrow v(gx),\qquad
 k_g:v(g)\longrightarrow v(gy).
\]
The oriented boundary of square \(g\) is
\[
                       h_g+k_{gx}-h_{gy}-k_g.       \tag{A.51.2}
\]
Subdivide each square by adding its centre \(m_g\)
and the four radial edges \(r_{g,j}\), directed
from the centre to corner \(j\), for corners
\[
        v(g),\quad v(gx),\quad v(gxy),\quad v(gy)
                     \qquad(j=0,1,2,3).
\]
There are \(36\) vertices, \(162\) edges, and
\(108\) positively oriented triangles. The
boundary of triangle \((g,j)\) is
\[
                   r_{g,j}+e_{g,j}-r_{g,j+1},
      \qquad(e_{g,0},e_{g,1},e_{g,2},e_{g,3})
                   =(h_g,k_{gx},-h_{gy},-k_g),       \tag{A.51.3}
\]
with radial subscripts read modulo four.
These formulas specify every entry of the two
integral cellular boundary matrices.

For clarity the cup signs are also explicit.
Order the nine original vertices lexicographically
before all the centres. For two rational cellular
one-cocycles \(a,b\), a triangle with oriented
outside edge \(e=\epsilon e_+\) contributes
\[
                -\epsilon\,a(e_+)\,b(r_{\rm high})
                                                               \tag{A.51.4}
\]
to \(\langle a\smile b,[C]\rangle\), where
\(r_{\rm high}\) is the radial edge ending at
the larger of the two original vertices. This
is the ordered-triangle cup formula: if the
outside vertices are increasingly ordered, the
ordered triangle is their pair followed by its
centre; reversing their order changes both signs
as in (A.51.4). It also follows on barycentrically
subdividing this ordered triangular complex.
Thus (A.51.4) computes the actual topological
polarization on cohomology.

Deck transformations act by left multiplication
on all labels. A lift of quarter rotation has
\[
 \alpha(a,b,c)=(-b,a,c-ab),\qquad
 h_g\longmapsto k_{\alpha(g)x},\quad
 k_g\longmapsto-h_{\alpha(g)},\quad
 r_{g,j}\longmapsto r_{\alpha(g),j+1}.              \tag{A.51.5}
\]
The formulas respect (A.51.2)--(A.51.3), have
order four, and conjugate the deck generators
by \(x\mapsto y,y\mapsto x^{-1}\).
They give a holomorphic lift of quarter rotation
of the square elliptic curve. It and the lift
(A.50.7) are conjugate by a deck transformation:
both are complements of the odd group \(H_3\)
in the extension by \(C_4\), and coprime
complement conjugacy applies. This comparison is
also why the following projector selects the
same factor \(B\) as A.50. Conjugating the lift
merely replaces \(q_0\) by one of its deck
translates, whose entire span is unchanged.

Let \(U,V,Z,R\) denote the resulting rational
operators on \(H^1(C,\mathbf Q)\), using inverse
pullback so that they form a covariant group
representation. They satisfy
\[
 UV=ZVU,\quad RUR^{-1}=V,\quad RVR^{-1}=U^{-1}.
\]
Put
\[
\begin{split}
 P&=1-\frac{1+Z+Z^2}{3},\qquad
 P_0=\frac{1+U+U^2}{3},\\
 A&=\frac{1+2Z}{3}
       \sum_{j,k=0}^2 Z^{-jk}V^jP_0V^{-k}\,P,
       \qquad T=A^3RP,\\
 P_-&=\frac{P-T}{2},\qquad e=P_0P_- .
\end{split}                                                   \tag{A.51.6}
\]
The matrix units \(V^jP_0V^{-k}\) show that
\(A\) is exactly the Fourier intertwiner
(A.50.11) on the nontrivial-central summand.
In particular
\[
 A^4=P,\quad T^2=P,\quad
 \operatorname{rank}P=18,\quad
 \operatorname{rank}\tfrac{P+T}{2}=6,\quad
 \operatorname{rank}P_-=12,\quad\operatorname{rank}e=4.
                                                               \tag{A.51.7}
\]
The plus part is the span of the actual elliptic
quotient maps from A.50; the minus part is
\(H^1(B^3,\mathbf Q)\). Taking \(P_0\), the
invariants of the noncentral order-three subgroup,
selects one copy of \(H^1(B,\mathbf Q)\).
These identifications use actual rational
endomorphisms of the Jacobian. They do not assign
a new geometric meaning to an arbitrary matrix
subspace.

Here is the complete small output needed from the
exact certificate. On \(eH^1(C,\mathbf Q)\),
there are vectors \(v_1,v_2\) forming a basis
over \(F=\mathbf Q(\zeta_3)\), with \(Z\)
acting as \(\zeta_3\), such that the cup matrix
in the rational basis \((v_1,Zv_1,v_2,Zv_2)\) is
\[
 \Psi_B=
 \begin{pmatrix}
 0&-1/6&0&0\\1/6&0&0&0\\
 0&0&0&1/3\\0&0&-1/3&0
 \end{pmatrix}.                                  \tag{A.51.8}
\]
The file \texttt{a\_heisenberg\_polarization.json}
retains all four cellular cocycles as exact rational
values on the labelled edges. The accompanying
program constructs (A.51.2)--(A.51.4), fixes a
spanning-tree gauge for coboundaries, and computes
a basis of the resulting twenty-dimensional
cocycle space. Its cup determinant is one.
It checks the complete deck and rotation relations,
symplectic preservation, every projector identity
and rank in (A.51.7), and (A.51.8) on the four
retained cocycles. This is a finite exact rational
certificate for the displayed pairing, with all
entries specified by the cellular formulas; it
uses no bounded-degree extrapolation.

Write \(\kappa=Z-Z^2=\sqrt{-3}\). The
corresponding Hermitian form is
\[
 h(v,w)=\frac{\Psi_B(v,\kappa w)}2
                  -\frac{\kappa\Psi_B(v,w)}2.
                                                               \tag{A.51.9}
\]
It is linear in its first variable and conjugate
linear in its second. Formula (A.51.8) gives
\[
                  h=\operatorname{diag}(-1/6,1/3),
                    \qquad-\det h=1/18.             \tag{A.51.10}
\]
Scaling the displayed polarization by a rational
number, or changing the \(F\)-basis, changes this
discriminant only by
an \(F\)-norm, so this particular exact choice
is sufficient for the quaternion calculation.

\paragraph{Construct the rational Hodge endomorphism.}
We also verify the Hodge input explicitly.
For the central automorphism \(z\mapsto\zeta_3z\)
in (A.50.1), the holomorphic eigendifferentials
of nontrivial character have dimensions three
and six. One may use the following direct bases:
\[
\begin{array}{c|l}
\zeta_3^2&dx/z,\quad dx/(yz),\quad x\,dx/(yz)\\
\zeta_3&dx/z^2,\quad x\,dx/z^2,
      \quad x^jdx/(yz^2)\ (0\le j\le3).
\end{array}                                                     \tag{A.51.11}
\]
At a finite branch point, \(x-a\) has order
three and \(z\) order one. At the single point
above infinity, \(x,y,z,dx\) have orders
\(-6,-9,-8,-7\), respectively. These facts
check holomorphicity; the eigenvalues and distinct
functions check independence. Together with the
one invariant elliptic differential they give
all ten differentials. Thus the threefold factor
\(B_3\) in A.50 has \(F\)-signature \((2,1)\),
up to exchanging the two embeddings. The elliptic
quotient differential is
\(q_0^*(dz/s)=4x\,dx/z^2\), of the larger
character. Removing its three deck translates
leaves the surface \(B\) with signature \((1,1)\).
Inverse pullback exchanges the two embeddings
and leaves this balanced signature unchanged.

The Hodge circle on \(H^1(B,\mathbf Q)\)
preserves \(h\), and its determinant over
\(F\) is one: at either embedding its two
circle characters are \(z\) and \(z^{-1}\).
It therefore lies in \(\operatorname{SU}(h)\).
This is the dimension-two Weil-type condition;
no equality with the full special unitary group
is required.

Define the following \(F\)-antilinear rational
operator on the basis of (A.51.10):
\[
                  j(v_1)=-v_2,\qquad j(v_2)=-2v_1,
                      \qquad j(av)=\bar a\,j(v).
                                                               \tag{A.51.12}
\]
A direct two-by-two calculation shows that it
commutes with \(\operatorname{SU}(h)\).
Indeed for a diagonal form \(\operatorname{diag}(a,b)\),
the antilinear operator with matrix
\(\bigl(\begin{smallmatrix}0&-b\\a&0\end{smallmatrix}\bigr)\)
commutes with every determinant-one isometry;
(A.51.12) is six times this operator for
\(a=-1/6,b=1/3\).
It follows that \(j\) is a rational Hodge
endomorphism. Weight-one Hodge realization is
fully faithful for complex abelian varieties,
by \cite[Theorems 4.4 and 7.12, pp.~489 and 508]{fghMilneSVModuli}.
After clearing lattice denominators this gives
an actual algebraic endomorphism of the complex
abelian variety. Endomorphisms of a number-field
abelian variety descend to a finite algebraic
extension; enlarge the field to define these
finitely many endomorphisms. If one works
contravariantly on cohomology, the resulting
endomorphism algebra is the opposite algebra,
which is isomorphic to the same quaternion algebra.

The exact relations are
\[
             \kappa^2=-3,\qquad j^2=2,
                         j\kappa=-\kappa j.
                                                               \tag{A.51.13}
\]
The four operators \(1,\kappa,j,\kappa j\)
are rationally independent: the first two are
\(F\)-linear and the last two \(F\)-antilinear,
and neither pair has a nonzero relation.
This proves the actual multiplication (A.51.1).

\paragraph{The division algebra forces the required reduction at three.}
The quaternion algebra \((-3,2)\) is division
over \(\mathbf Q_3\). To check this directly,
\(\mathbf Q_3(\sqrt{-3})/\mathbf Q_3\) is
ramified. A unit norm reduces to a square in
\(\mathbf F_3^\times\), hence to one; two is
not a norm. In the quaternion presentation the
reduced norm of \(a+bj\) is
\(N(a)-2N(b)\). A nonzero zero of this norm
would make two a norm, so the reduced norm is
anisotropic and the algebra is division.
Every nonzero module over this division algebra
has \(\mathbf Q_3\)-dimension at least four.

Pass to a finite local extension over which
\(B\) is semistable and all its endomorphisms
are defined. Endomorphisms extend to the Neron
model and preserve its torus. If that torus
had positive rank \(r\le2\), its character
space tensored with \(\mathbf Q_3\) (or the
contravariant dual) would be a nonzero unital module over the
four-dimensional division algebra. This is
impossible. The toric rank is therefore zero,
so semistability is good reduction. This is
the same extension-and-specialization argument
used in A.44, now with a division algebra in
place of a full CM field; the extension facts
are those in \cite[proof of Theorem 6(a), p.~503]{fghSerreTate1968}.

On that good special fibre the rational etale
Tate space of its three-divisible group has
dimension equal to the three-rank, at most two.
It too is preserved by every endomorphism.
If nonzero it would be a unital module over
\((-3,2)_{\mathbf Q_3}\), again impossible.
Thus the special three-rank of \(B\) is zero.
The word ``unital'' matters: the identity
endomorphism restricts to the identity on any
nonzero toric or etale part, so this cannot be
an action which kills the division algebra.

The original \(E\) has good supersingular
reduction at three. The CM elliptic curve
\(E_0\) has potentially good reduction by
\cite[Theorem 6(a)]{fghSerreTate1968}; its ramified
quadratic CM field over \(\mathbf Q_3\)
cannot act on a nonzero etale three-Tate space
of dimension at most one, so its special
three-rank is zero as well. By (A.50.2) the
same potential good reduction and three-rank
zero hold for \(J_C\).

\paragraph{All cusp differences and the full integral fixed subgroup.}
For each \(T\in E[3]\), let \(P_T\) be
its unique preimage on \(C\). If \(T,U\)
are two such points, \(3(T-U)\) is principal
on \(E\). Total ramification of \(q\)
gives
\[
             q^*\bigl(3(T-U)\bigr)
                           =9(P_T-P_U).             \tag{A.51.14}
\]
Thus every cusp difference is killed by nine,
not merely the sum of the eight finite cusps.
After good reduction of \(J_C\) at three,
all these differences specialize to zero,
since its special three-rank is zero.

The stable reduction of \(C\) is of compact
type, because its Jacobian is good. Separate
all cusp sections in a semistable pointed model.
The Abel-map criterion used in A.42 shows that
any two of these marks are joined entirely
through rational components, including their
marked endpoint components: the only connected
fibres of the special Abel map are the rational
subtrees. Equivalently one may apply directly
the p-primary cusp-difference sufficient
criterion of A.42. There are nine distinct
cusps, so every selected cusp has another one
to which this criterion applies.
A.42 then proves the full fixed inertia assertion
at every coefficient prime, for every indicated
packet and open derived subgroup. This final
step is the proved group theorem, including its
integral cusp-power subtraction; no rational
Malcev fixed element is silently lifted.
\end{proof}

% Registered A.52: independent full root and DE proof/source/code audits PASS.
\subsection{An explicit wild model and more Heisenberg cusps}

\begin{statement}{A.52}{Proved}{The three-adic Heisenberg model and all torsion enlargements}
Let \(C\), \(q:C\to E\), and \(E:y^2=x^3-x\) be as in A.50.
After a finite extension at three, the positive-genus components of
its stable reduction are one supersingular elliptic curve and three
copies of the genus-three Artin--Schreier curve
\[
                         V^3-V=X^4.                    \tag{A.52.1}
\]
The elliptic component maps purely inseparably with degree three to
\(E\); each of the other three components maps to a different point
of \(E[2]\setminus\{O\}\). The reduction is of compact type and its
Jacobian has three-rank zero. This gives a geometric proof of the
reduction assertion in A.51 independent of its quaternion calculation.

For every positive integer \(m\), put
\[
 X_m=C\setminus q^{-1}(E[3m]),\qquad
                 X_m\xrightarrow{[3m]q}E\setminus\{O\}.
                                                               \tag{A.52.2}
\]
This is a connected finite etale cover of degree \(27m^2\), proper
genus ten, and \(27m^2-18\) cusps. For every coefficient prime its
nontrivial closed procyclic subgroups with open arithmetic normalizer
are precisely the nontrivial subgroups of conjugates of cusp inertia.
For every positive \(m\), the full geometric pro-prime fixed subgroup,
for every selected cusp packet and every indicated open derived group,
is also the selected cusp inertia.
\end{statement}

\begin{proof}
All fields of definition may be enlarged finitely. First choose a
local extension at three over which the proper curve has semistable
reduction, and then enlarge it to contain the finitely many constants
below. The residue field can be treated geometrically. Write
\(f(x)=3x^4-6x^2-1\), so that \(C\) has function field
\[
                       y^2=x^3-x,\qquad z^3=f(x).
\]

\paragraph{Three actual algebraic charts.}
For each \(a\in\{0,1,-1\}\), set
\(\delta=3a^2-1\). Choose constants
\[
         \eta^2=-3,\qquad \rho^4=\eta,
                       \qquad z_0^3=f(a),
\]
and use the invertible generic coordinate changes
\[
       x=a+\rho^2U,\qquad y=\rho X,
                              \qquad z=z_0(1+\eta V).
\]
The exact integral equations are
\[
\begin{aligned}
 X^2&=\delta U+3a\rho^2U^2+\rho^4U^3,\\
 z_0^3(V+\eta V^2-V^3)
       &=2\delta U^2+4a\rho^2U^3+\eta U^4.
\end{aligned}                                                   \tag{A.52.3}
\]
For example
\(f(a+d)-f(a)=6\delta d^2+12ad^3+3d^4\), and division by
\(3\eta\), with \(d=\rho^2U\), gives the second equation.
This verifies (A.52.3) as an exact equation in the actual function
field, not a truncated expansion.

The algebra in (A.52.3) is free over the valuation-ring polynomial
algebra in \(U\), using the monic equations in \(X\) and \(V\).
It is therefore flat, and its generic fibre is the corresponding
open subset of \(C\). On the special fibre \(\delta=-1\) and
\(z_0^3=-1\). The first equation eliminates \(U=-X^2\), and
the second becomes exactly (A.52.1). The partial derivatives with
respect to \(U\) in the first equation and \(V\) in the second
are units on the special fibre, so this special affine curve is
smooth. It is geometrically integral: a polynomial of degree four
cannot be of the form \(g^3-g\) in \(\overline{\mathbf F}_3(X)\),
since a pole of \(g^3-g\) has order divisible by three. Its
smooth projective completion has a unique point at infinity and
genus three. One direct ramification calculation is as follows.
The cover of \(\mathbf P^1_X\) has degree three, is etale at
finite points, and has pole order four at infinity. Its single
lower jump is four: with \(t=X^{-1}\), the valuation on the
extension has \(v(t)=3\), \(v(V)=-4\), and \(t^3V^2\) is a
uniformizer whose difference under \(V\mapsto V+1\) has order
five. The different has order \(5(3-1)=10\), and Hurwitz gives
\(2g-2=-6+10=4\).

We spell out why these are components of the stable model. Compactify
and normalize each integral affine chart, and take a regular resolution
of the resulting proper model. The smooth generic point of its special
curve retains the same residue function field. Birational morphisms
between regular arithmetic surfaces factor into blowups of closed
points, and only rational curves occur in their exceptional loci;
see \cite[Theorem 9.2.2 and Remark 9.2.6, pp.~395--396,
Theorem 9.3.21, p.~422]{fghLiu}. Passing to the minimal regular
model and then contracting the rational components to the stable model
therefore retains this positive-genus component. Equivalently one can
dominate these models and a regular resolution of the stable model by
a common regular model; the same factorization argument gives the
claim. The three valuations are distinct, since their restrictions to
\(E\) have centers \((a,0)\) at three different special points.
Each corresponding component maps to that point.

\paragraph{The remaining elliptic component.}
Choose \(\pi^3=3\) and substitute \(z=-1+\pi u\). The exact main
chart is
\[
       u^3-\pi^2u^2+\pi u=x^4-2x^2,
                         \qquad y^2=x^3-x.           \tag{A.52.4}
\]
Its special function field is
\(\overline{\mathbf F}_3(E)((x^4+x^2)^{1/3})\). This is a field:
\[
                 d(x^4+x^2)=(x^3-x)\,dx=y^2\,dx\ne0.
\]
Over a perfect residue field, a one-variable function field has
\([k(E)^{1/3}:k(E)]=3\). Thus this extension is precisely
\(k(E)^{1/3}\); its smooth projective curve is the Frobenius twist
of \(E\), of genus one, and its map to \(E\) is purely inseparable
of degree three. The main chart is smooth at the generic point of
this special component. The same compactification and factorization
argument retains it in the stable model. Its valuation lies over the
generic point of the original elliptic special fibre, so it is
distinct from all three preceding valuations.

The sum of the four component genera is already
\(1+3+3+3=10=g(C)\). The genus formula for the connected stable
fibre consequently forces the dual graph to be a tree and every
other component to be rational. In particular the Jacobian has good
reduction after this extension. The main elliptic curve is
supersingular at three. Each Artin--Schreier component has three-rank
zero: Deuring--Shafarevich for its cyclic-three map to a rational
curve with one totally ramified point gives
\(f-1=3(-1)+(3-1)=-1\); see
\cite[formula (6), p.~5]{fghGiuliettiKorchmaros}.
The Jacobian of the compact-type fibre is the product of the component
Jacobians, hence has three-rank zero. This part of the proof did not
use A.51.

For use with markings, the actual map \(q\) extends to the good
elliptic model. One can see this directly on a regular semistable
model of \(C\): resolve the indeterminacy of the rational map to
that proper elliptic scheme by \cite[Theorem 9.2.7, pp.~396--397]{fghLiu}. Each exceptional component is rational,
so its map to the smooth elliptic special fibre is constant; the map
descends through the blowups. It also descends through stable
contraction, for the same reason. On the special fibre its degree on
the main component is three, as (A.52.4) proves. A degree-one line
bundle on the target has pullback of degree three on every fibre.
Consequently all other components have degree zero and map constantly
to the elliptic curve. Every connected subtree off the main component
attaches at one point, and maps to the image of that point. Since the
main map is radicial, a given reduced point of \(E\) has just one
possible attachment point on the main component.

\paragraph{The enlarged markings and their rational paths.}
The proper degree-three map \(q\) ramifies exactly above \(E[3]\),
with one point of ramification index three in each of its nine fibres.
The isogeny \([3m]\) is etale in characteristic zero. Removing its
whole origin fibre therefore gives (A.52.2), with degree \(27m^2\)
and
\[
        \#q^{-1}(E[3m])=9+3(9m^2-9)=27m^2-18.       \tag{A.52.5}
\]
Its Euler characteristic is \(2-20-(27m^2-18)=-27m^2\), as required.

Suppose now that \(m\) is odd, and write
\(3m=3^h n\), where \(h\ge1\) and \(n\) is odd and prime to
three. Because \(E\) is supersingular, its geometric three-primary
torsion specializes to the origin. Prime-to-three torsion specializes
injectively. Thus the reduction support of \(E[3m]\) is exactly
\(E[n]\), and each such reduction class contains \(3^{2h}\)
distinct characteristic-zero torsion points. None of these reduced
points is in \(E[2]\setminus\{O\}\).

All marks in a given class can therefore specialize only on the main
component or on rational subtrees attached at its unique point above
that class. They cannot specialize on any of the three genus-three
components, or on any connecting subtree leading to such a component:
those components map to a nonzero two-torsion point. For a nonzero
reduced class every relevant \(q\)-fibre is unramified and contains
three points, giving \(3^{2h+1}\) marks. In the origin class the
nine branch fibres contribute nine marks and the others contribute
three each, giving
\[
 \#\text{marks in a class}=
 \begin{cases}
    3^{2h+1}-18,&\text{the origin class},\\
    3^{2h+1},&\text{a nonzero class}.
 \end{cases}                                                   \tag{A.52.6}
\]
In particular every class contains at least nine marks.

After a further finite extension make all marks into sections and
separate them on a regular semistable model. The inverse image of
any one such main-component point is a connected rational tree.
All marks in that class lie on this tree: if any remained on the
main component at the attachment point, it would meet a node rather
than the required smooth marked locus. Repeated separation and stable
marked contraction only change rational components. Hence for every
selected cusp there is a second cusp in the same class such that
the entire path joining their marked components is rational. This
is the pointed hypothesis of A.42; it does not require the paths
between different reduction classes to be rational. Applying A.42
with the compact-type, three-rank-zero reduction just proved gives
the full integral fixed-inertia assertion for every coefficient prime,
including three, every packet, and every indicated open derived group.

\paragraph{Even multiples use the smooth model at two.}
It remains to treat even \(m\). Write \(m=2^a b\), with \(a\ge1\)
and \(b\) odd. Use the actual model constructed in the residue-two
paragraph of A.50: after finite extension \(E\) has good supersingular
reduction, all its three-torsion sections are distinct, and normalizing
its good model in \(z^3=f\), after removing the constant Kummer twist,
gives a smooth proper model of \(C\). The map \(q\) remains a finite
cyclic-three morphism of these models. Its local form at each branch
section is a unit times a cube of a source parameter; away from those
sections it is etale. Thus its special map has degree three, is tamely
and totally ramified at the nine distinct three-torsion points, and
is etale elsewhere. These properties concern the actual normalized
map used in A.50, not just the abstract good-reduction curve.

The reduction support of \(E[3m]\) at two is \(E[3b]\). Each reduced
target class contains \(4^a\) characteristic-zero torsion points.
Above a class outside \(E[3]\), the special map \(q\) has three
distinct preimages. In the strict henselian completed neighborhood
of each one, its etale local map is an isomorphism. Consequently
each of those three source points receives exactly \(4^a\) marked
sections. Above a reduced point of \(E[3]\), the special source
point is unique. One of the \(4^a\) target sections is the branch
section; it contributes one source mark. The other \(4^a-1\) target
sections have three distinct generic preimages, all reducing to
that same source point. Its marking count is therefore
\[
                         1+3(4^a-1)=3\cdot4^a-2.    \tag{A.52.7}
\]
These two cases account for every mark. Each source reduction
cluster has at least four marks. Since the proper source model is
smooth, separating each such cluster produces a rational marked
tree attached at its common smooth point. For every selected mark,
choose a second one in the same cluster; their whole joining path
is rational. The pointed model is of compact type, and A.42 at
residue two proves the full fixed-inertia conclusion for every
coefficient prime and packet. Together with the odd case, this
settles every positive \(m\).

\paragraph{The normalizers for arbitrary \(m\).}
This part does not need \(m\) odd. Every actual punctured elliptic
map used in A.50 remains a morphism from the smaller curve \(X_m\)
to the same once-punctured target. Its proper-homology map is unchanged.
The Jacobian decomposition of A.50, the deck translates detecting
its CM elliptic factors, and the absence of a potentially stable
rational line in its remaining factor are also unchanged. Thus
A.11 and the exact A.50 map-detection argument exclude every
noncuspidal procyclic subgroup with open arithmetic normalizer.
Conversely every nontrivial subgroup of cusp inertia has such a
normalizer after a finite extension making the cusp rational.

The accompanying exact certificate checks the two polynomial chart
identities for all three values of \(a\), the reduced equation,
the main chart, and the degree and marking-count identities at both
residue primes.
\end{proof}

% Registered A.53: independent root and DE full proof/source/code audits PASS.
\subsection{A fixed-subgroup enlargement on every cover compactification}

\begin{statement}{A.53}{Proved}{Every cover compactification admits a torsion enlargement with full fixed inertia}
Let \(Y\to U\) be any connected finite etale cover from Task A.5,
with proper compactification \(C\), degree \(d\), proper genus
\(g\), and \(r\) cusps. Let \(q:C\to E\) be its extended
proper map. There are infinitely many rational primes
\(p\equiv3\pmod4\) such that
\[
       Z_p=C\setminus q^{-1}(E[p]),\qquad
                    Z_p\xrightarrow{[p]q}U             \tag{A.53.1}
\]
is a connected finite etale cover of degree \(p^2d\), proper genus
\(g\), and \(p^2d-2g+2\) cusps, and its full geometric
pro-\(\ell\) fixed group is the selected cusp inertia for every
coefficient prime \(\ell\), every cusp packet, and every indicated
open derived subgroup.

If \(Y\) has cuspidal procyclic normalizer recognition at a given
coefficient prime, then every further puncturing of \(Y\), in
particular \(Z_p\), has that recognition at the same prime.
Thus any previously proved all-prime normalizer theorem for \(Y\)
also holds for (A.53.1).

The comparison \(Z_p\hookrightarrow Y\) is an open immersion;
the new map to \(U\) is the composite with \([p]\). This construction
asserts neither finite refinement over the original map to \(U\)
nor lifting of fixed elements from \(Y\). In particular it does
not prove the universal assertion of Task A.5.
\end{statement}

\begin{proof}
Enlarge the number field so that all original cusps are rational.
Their proper genus satisfies \(g\ge1\), since \(q\) is a finite
map to an elliptic curve. The map is ramified only above the origin,
and its full origin fibre is precisely the cusp set. Let \(e_c\)
be its ramification index at cusp \(c\). Thus
\[
                   \sum_c e_c=d,\qquad
                         2g-2=\sum_c(e_c-1)=d-r.    \tag{A.53.2}
\]

\paragraph{Choose an actual good tame model.}
Spread out the proper curve, the finite morphism, and the finitely
many cusp sections over an open subset of the integer ring. After
removing finitely many primes, both proper curves are smooth, the
cusp sections are pairwise disjoint, \(q\) is finite flat of degree
\(d\), and it is etale off the origin section. Moreover the equality
of Cartier divisors
\(q^*(O)=\sum_c e_c(c)\) holds on these models.
These assertions can be obtained by clearing denominators in finite
presentations for the morphism and the displayed divisor equality,
and then removing the closed nonsmooth, nonflat, and unwanted
ramification loci, which have empty generic fibre. In particular
only finitely many rational primes have been excluded. The
finite-bad-reduction assertion is also
\cite[Proposition 10.1.21(a), pp.~462--463]{fghLiu}.

Choose a rational prime \(p\equiv3\pmod4\) outside that finite set
and with \(p>d\). There are infinitely many choices. This inequality
makes every index \(e_c\) prime to \(p\); it avoids inferring
individual tame indices merely from a degree coprimality. At a place
above \(p\), \(E:y^2=x^3-x\) has supersingular good reduction;
see \cite[Theorem V.4.1(a) and Example V.4.5, pp.~148 and 152]{fghSilvermanAEC}.
After a further finite extension make all points of \(E[p]\) and
their inverse images rational. Good reduction and the stated finite
map persist. Every point of \(E[p]\) specializes to the origin,
since a supersingular elliptic curve has no nonzero geometric
\(p\)-torsion in its special fibre.

\paragraph{Every cusp has a companion in its own rational cluster.}
Fix an original cusp \(c\). In the completed strict henselian
neighborhoods of the source and target sections choose parameters
\(u,t\). The Cartier-divisor equality gives
\[
                              t=u^{e_c}h(u),
                                                               \tag{A.53.3}
\]
where \(h\) is a unit. Since \(e_c\) is invertible in the residue
field, extract an \(e_c\)-th root of this unit and change \(u\).
The map is then exactly \(t=u^{e_c}\). Each of the \(p^2-1\)
nonzero characteristic-zero points of \(E[p]\) has nonzero parameter
of positive valuation. Its \(e_c\) distinct roots all reduce to
\(c\). The origin itself has one source point in this neighborhood.
Consequently the number of marked points specializing to \(c\) is
\[
                              1+e_c(p^2-1)\ge p^2.  \tag{A.53.4}
\]
There are no other reduction classes: all target points specialize
to the origin, and the special origin fibre is the disjoint set of
original cusp points with its indicated multiplicities. These local
counts sum to all points of \(q^{-1}(E[p])\).

The proper source model is smooth. Separating each of these clusters
by blowups, after finite extension if necessary, creates rational
marked trees attached at the corresponding original smooth source
point. For any chosen new cusp, another mark lies in its cluster by
(A.53.4), and their entire joining path, including both marked
components, is rational. The proper model has compact type because
it is smooth; inserting these rational trees preserves that property.
This is exactly the pointed geometric hypothesis of A.41--A.42.
Those statements give the full integral fixed-inertia assertion at
every coefficient prime, including \(\ell=p\), for every packet
and every indicated open derived group. No restriction on the
\(p\)-rank or endomorphism algebra of \(J_C\) is required here.

The characteristic-zero map \([p]q\) has degree \(p^2d\) and is
unramified after removal of its full origin fibre. The new affine
curve is geometrically connected because it is an open subset of
the same smooth connected proper curve. Its cusp count is
\[
 \sum_c\bigl(1+e_c(p^2-1)\bigr)
       =r+(p^2-1)d=p^2d-2g+2.                       \tag{A.53.5}
\]
This also verifies its Euler characteristic \(-p^2d\).

\paragraph{Normalizer recognition survives further puncturing.}
More generally let \(Z\hookrightarrow Y\) be an open immersion of
smooth affine curves with the same proper compactification. Suppose
that \(Y\) satisfies the pro-\(\ell\) normalizer criterion. If a
noncuspidal normalized procyclic subgroup existed in the geometric
pro-\(\ell\) group of \(Z\), A.11 would give it nonzero image in
the proper Tate homology of that common compactification. Its image
under the actual morphism \(Z\to Y\) is therefore nontrivial and
has the same nonzero proper-homology image. The arithmetic normalizer
maps to an arithmetic normalizer with open projection. The criterion
on \(Y\) puts this image in cusp inertia, whose proper-homology
image is zero. This contradiction proves the assertion.
The converse direction, that subgroups of cusp inertia have open
arithmetic normalizer after finite field restriction, is the usual
cusp decomposition-group inclusion.

Finally, the actual comparison to \(Y\) is the open immersion, which
induces a surjection on geometric pro-prime fundamental groups.
The fixed-group functor is not right exact, so the fixed theorem
for \(Z_p\) is not being pushed down to \(Y\). The map (A.53.1)
is a new covering map to \(U\), rather than a factorization through
\(Y\to U\) as a finite cover of that same base.
\end{proof}

% Registered A.54 after full independent root, DE and PBC audits; exact reruns PASS.
\subsection{The degree-one-hundred-twenty-five Heisenberg cover}

\begin{statement}{A.54}{Proved}{An actual non-CM factor and the nonsplit-prime normalizers}
Let \(C\) be the smooth proper curve
\[
 \begin{aligned}
 y^2&=x^3-x,\qquad z^5=\psi _5(x),\\
 \psi _5(x)&=5x^{12}-62x^{10}-105x^8+300x^6-125x^4+50x^2+1,
 \end{aligned}
                                                               \tag{A.54.1}
\]
and let \(q:C\to E\), \(S=q^{-1}(E[5])\), and \(X=C\setminus S\).
The map \([5]q:X\to U\) is a geometric Galois cover
with Heisenberg group of order \(125\), genus \(51\), and \(25\) cusps.
After a finite number-field extension,
\[
 J_C\sim E\times A_2^5\times B_8^5,\qquad
 A_2\sim J(v^2=w^5-1),\quad \dim B_8=8.             \tag{A.54.2}
\]
Here \(A_2\) has primitive quartic CM by \(\mathbf Q(\zeta _5)\),
and \(B_8\) has an action by \(K=\mathbf Q(\zeta _{20})\)
but has no geometric CM abelian subvariety.
The last assertion follows from an actual genus-eleven quotient of \(C\)
and two persistent, unequal quadratic Frobenius extensions of \(K\).

For every coefficient prime \(\ell\not\equiv1\pmod5\),
the nontrivial procyclic subgroups of the geometric pro-\(\ell\) group
with open arithmetic normalizer are exactly the nontrivial subgroups
of cusp inertia conjugates. The full pro-five fixed subgroup for every
packet and open derived subgroup in Part A is the selected cusp inertia.
The normalizer question for \(\ell\equiv1\pmod5\) is completed
by A.55 below. The full fixed-subgroup comparison at every other
prime is completed separately by the explicit model in A.58.
\end{statement}

\begin{proof}
\paragraph{The actual cover.}
The division-polynomial divisor is
\[
 \operatorname{div}(\psi _5)=D_{E[5]}-25(O),
 \qquad D_T:=\sum_{P\in T}(P).                     \tag{A.54.3}
\]
Its orders are \(1\) at the twenty-four nonzero points and \(-24\)
at \(O\). The Kummer cover is connected, of degree five and totally
ramified at exactly these twenty-five points. Consequently
\(2g(C)-2=25(5-1)=100\); removing their unique preimages gives
the asserted degree and number of cusps.

Every translation by \(T\in E[5]\) lifts after a finite constant
extension. Indeed, the divisor of \(\psi _5(P+T)/\psi _5(P)\) is
\(25((O)-(-T))\). Since \(5((O)-(-T))\) is principal, this
ratio is a constant times a fifth power of a rational function.
Taking its fifth root gives the lift. The lifts and
\(c:z\mapsto\zeta _5z\) form a central extension of \(E[5]\)
by \(\mu _5\), of order \(125\). It is nonabelian: in an abelian
cover of a once-punctured elliptic curve, the sole peripheral
commutator has trivial image, whereas here its inertia has order five.

This group has exponent five. In a class-two group with central
kernel of order five, the fifth-power map on lifts defines an additive
map \(E[5]\to\mu _5\). The actual involution
\(\sigma(x,y,z)=(x,-y,z)\) acts by \(-1\) on \(E[5]\) and
trivially on the central kernel. Equivariance therefore makes that
additive map equal to its negative, hence zero. The nonabelian
exponent-five central extension is the Heisenberg group \(H_5\).
All lifts, constants, and later projectors are defined over a finite
number field. A.22 applies to this actual five-group cover, giving
the asserted full pro-five fixed subgroup.

\paragraph{The multiplicity factors are algebraic factors of this curve.}
Put \(F=\mathbf Q(\zeta _5)\). The nontrivial-central component
of \(\mathbf Q[H_5]\) is \(M_5(F)\): its standard representation
on \(F^5\) is generated by
\[
 Ue_j=e_{j+1},\qquad Ve_j=\zeta _5^j e_j
                         \quad(j\bmod5).
\]
The central-invariant homology is \(H_1(E,\mathbf Q)\); all
translations act trivially there. The remaining rational dimension
is \(102-2=100\), or five copies of the simple rational module
of dimension twenty. Matrix idempotents, cleared of denominators,
thus give actual isogeny factors \(J_C\sim E\times B_{10}^5\),
where \(\dim B_{10}=10\) and \(\dim_F H_1(B_{10})=5\).

The actual order-four automorphism is
\(\tau(x,y,z)=(-x,iy,z)\), with \(\tau^2=\sigma\).
Choose a basis \(T,iT\) of \(E[5]\), avoiding its two
\(i\)-eigenlines. A central adjustment of a lift of \(T\)
ensures that \(\sigma\) takes that lift to its inverse:
the adjustment solves one linear equation with coefficient two
in \(\mathbf F_5\). Taking its \(\tau\)-conjugate for the
second lift gives \(U\mapsto V,V\mapsto U^{-1}\), with a
consistent choice of central generator. The Fourier intertwiner
is
\[
 A=\frac1{\sqrt5}(\zeta _5^{jk})_{0\le j,k\le4},
 \qquad \sqrt5=1+2\zeta _5+2\zeta _5^4.
                                                               \tag{A.54.4}
\]
Direct multiplication gives \(AUA^{-1}=V\),
\(AVA^{-1}=U^{-1}\), and \(A^2e_j=e_{-j}\).
Thus \(A^4=1\), and its multiplicities at \(1,-1,i,-i\)
are \((2,1,1,1)\). Write \(\tau=A\otimes T\) on the
noncentral homology. Equivalently, \(A^{-1}\tau\) is an actual
rational endomorphism commuting with \(M_5(F)\) and acting
as \(T\) on \(B_{10}\).

The quotients by \(\tau\) and \(\sigma\) have equations
\(z^5=P_6(s)\), \(s=x^2\), and \(z^5=\psi _5(x)\),
respectively. Their genera are ten and twenty-two: the polynomials
have six and twelve simple roots, with degrees coprime to five.
The original elliptic summand contributes no invariants under either
automorphism. If the multiplicities of \(T\) at \(1,-1,i,-i\)
are \(a,b,d,d\), equality of the last two follows from \(i\notin F\),
and the actual quotient genera give
\[
 a+b+2d=5,\qquad 2a+b+2d=5,
 \qquad3(a+b)+4d=11.
\]
Hence \((a,b,d,d)=(0,1,2,2)\). The \(T=-1\) factor has
dimension two and rank one over \(F\), so is CM; the \(T^2=-1\)
factor has dimension eight, with \(K=F(i)\) acting at rank two.
These are the factors \(A_2,B_8\) in (A.54.2).

\paragraph{An explicit quotient identifies the same factors over a number field.}
Take
\[
 \alpha=2-i,\qquad \lambda_H=\frac{1-2i}{5},\qquad
 \lambda=\frac{1+2i}{5},\qquad h_H=x^2-\lambda_H.
\]
The nonzero points of \(H=\ker\alpha\) have \(x^2=\lambda_H\).
This follows either from the duplication law, with
\([2]P=[i]P\), or from the following exact coordinate identity:
\[
 \begin{aligned}
 x(\alpha P)&=
 \frac{x(3+4i)(x^2-1-2i)^2}{25(x^2-\lambda_H)^2},
 \\
 \psi _5&=(-35-120i)\bigl(x(\alpha P)^2-\lambda\bigr)h_H^5.
 \end{aligned}
                                                               \tag{A.54.5}
\]
Its divisor explanation is equally exact:
\(\operatorname{div}(h_H)=D_H-5(O)\), while
\(x^2-\lambda\) has divisor \(D_{\ker\bar\alpha}-5(O)\).
Pulling the latter back by \(\alpha\) gives
\(D_{E[5]}-5D_H\). The remaining ratio is constant, and
the expansion at \(O\) identifies it as \(5\alpha^4=-35-120i\).

After adjoining \(\gamma\) with \(\gamma^5=-35-120i\),
the actual map
\[
 \begin{gathered}
 C\longrightarrow D_\lambda:\qquad
 (P,z)\longmapsto\left(\alpha P,\frac{z}{\gamma h_H(P)}\right),
 \\
 D_\lambda:\quad y^2=x^3-x,\quad z^5=x^2-\lambda.
 \end{gathered}
                                                               \tag{A.54.6}
\]
has degree five. The identity makes its function-field extension
the pullback along \(\alpha\), so it is the quotient by a
noncentral order-five subgroup, denoted again by \(H\).
That subgroup acts freely since all inertia of the full cover is
central. It is \(\tau\)-stable, and (A.54.6) is equivariant
for \(\tau\) and \(c\). In particular \(D_\lambda\) has
genus eleven. This explicit identity, rather than branch data alone,
identifies the quotient and its constant twist.

In the standard five-dimensional module, the invariant line of
every \(\sigma\)-stable noncentral subgroup has \(A^2=+1\).
For completeness, such a subgroup is generated by
\(\zeta _5^{ab/2}U^aV^b\); if \(a\ne0\), its invariant
vector has coordinates \(\zeta _5^{bj^2/(2a)}\), which are
even in \(j\), and for \(a=0\) it is \(e_0\).
Thus on the noncentral homology of this actual quotient,
\(\sigma=T^2\). It follows that
\[
 J_{D_\lambda}\sim E\times A_2\times B_8,\qquad
 D_\lambda/\sigma:\quad x^2=z^5+\lambda.             \tag{A.54.7}
\]
The last curve becomes \(v^2=w^5-1\) after scaling over
a finite extension, proving the asserted identification of \(A_2\).
Its holomorphic differentials \(dz/x,z\,dz/x\) give the
quartic CM type \(\{1,2\}\subset(\mathbf Z/5)^\times\).

\paragraph{The actual primitive-twenty factor and its Hodge signature.}
On \(D_\lambda\), put \(s=x^2\) and \(W=yz\). Then
\[
 W^{20}=s^5(s-1)^{10}(s-\lambda)^4.                 \tag{A.54.8}
\]
This is a birational cyclic-twenty model. Indeed
\(y^4=s(s-1)^2\), \(z^5=s-\lambda\), and
\(y=W^5/(s(s-1)^2(s-\lambda))\), recovering the two
coprime subextensions. Also \(x=y^2/(s-1)\).
Choose once and for all
\[
 K=\mathbf Q(\zeta),\quad \zeta=\zeta_{20},
 \quad i=\zeta^5,\quad \zeta_5=\zeta^4.
\]
The automorphism \(\delta=\tau c^4\) sends \(W\) to
\(\zeta W\); it satisfies \(\delta^4=c\),
\(\delta^5=\tau\), and \(\delta^{10}=\sigma\).
Consequently its primitive-twenty rational projector selects
exactly the same \(B_8\) as (A.54.7). It defines that factor
over \(K\) itself, independently of the finite extension used
to compare it with \(C\).

The four local exponents in (A.54.8) are \((5,10,4,1)\)
modulo twenty, the last at infinity. The space of holomorphic
differentials with \(\delta\)-eigenvalue \(\zeta^{-k}\)
has dimension
\[
 h_k=-1+\sum_{m\in\{5,10,4,1\}}\left\{\frac{km}{20}\right\}.
                                                               \tag{A.54.9}
\]
One obtains this directly by writing a differential as
\(r(s)\,ds/W^k\) and imposing the four valuation inequalities;
the integral parts give a line bundle on \(\mathbf P^1\)
of degree \(-2+\sum\{km/20\}\), whose space of sections
has the stated dimension. For primitive \(k\), the values are
\[
\begin{array}{c|rrrrrrrr}
k&1&3&7&9&11&13&17&19\\\hline
h_k&0&1&1&1&1&1&1&2.
\end{array}                                                     \tag{A.54.10}
\]
In particular there are balanced embeddings with dimensions
\((1,1)\), and one unbalanced conjugate pair \((0,2)\).
No balanced-determinant or quaternion argument is being substituted
for the following arithmetic calculation.

\paragraph{Exact Frobenius fields at two places of the same field \(K\).}
Both \(41\) and \(61\) split completely in \(K\).
For \(q\) equal to either, choose a primitive root \(g\) modulo
\(q\), and the place with
\(\zeta\mapsto g^{(q-1)/20}\). This fixes
\(i\mapsto g^{(q-1)/4}\) and \(\lambda\mapsto(1+2i)/5\)
in the reduction. The four branch sections in (A.54.8) remain
distinct and all ramification is tame, so the smooth normalization
and its Jacobian have good reduction. Define
\(\chi(g)=\zeta\) and extend \(\chi(0)=0\). On the
primitive eigenspace with \(\delta^*\)-eigenvalue \(\zeta\),
the trace of Frobenius is
\[
 a_q=-\sum_{u\in\mathbf F_q}
       \chi\bigl(u^5(u-1)^{10}(u-\lambda)^4\bigr).  \tag{A.54.11}
\]
This follows from the curve trace formula applied to the
Kummer eigensheaf; every local inertia character is nontrivial,
so compactly supported and proper primitive cohomology agree.
Equivalently, average the ordinary fixed-point formula over
the twenty deck twists. The convention \(\chi(g)=\zeta\)
agrees with the displayed \(\delta^*\) eigenvalue, since
the Frobenius permutation of a fibre is multiplication by
\(f(u)^{(q-1)/20}\).
These are direct applications of
\cite[Theorem 25.1, p.147, and Theorem 29.4, p.165]{PBC_MilneLEC}.

For the square Frobenius use \(\chi\circ N_{\mathbf F_{q^2}/\mathbf F_q}\)
in the same sum. Writing \(u=a+b\sqrt\nu\), each norm is
an elementary quadratic expression, so the retained certificate
uses only integer arithmetic and a specified nonsquare \(\nu\).
Set \(d_q=(a_q^2-a_{q^2})/2\),
\(\Delta_q=a_q^2-4d_q\); the relative Frobenius polynomial is
\(t^2-a_qt+d_q\). The exact sums give
\[
\begin{aligned}
a_{41}&=\zeta^5-\zeta^4+\zeta^3-2\zeta^2+2\zeta-1,\\
a_{61}&=\zeta^7+\zeta^6-2\zeta^5-4\zeta^4
                  +\zeta^3+4\zeta^2-4\zeta-6,\\
\Delta_{41}&=80\zeta^7+60\zeta^6-40\zeta^5-85\zeta^4\\
           &\hspace{1em}+160\zeta^3+75\zeta^2-70\zeta-55,\\
\Delta_{61}&=50\zeta^7+85\zeta^6-170\zeta^5-105\zeta^4\\
           &\hspace{1em}+60\zeta^3+115\zeta^2-170\zeta-30.
\end{aligned}                                                  \tag{A.54.12}
\]
All expressions are reduced modulo
\(\Phi_{20}(\zeta)=\zeta^8-\zeta^6+\zeta^4-\zeta^2+1\).
The complete twenty-bin counts, determinants, and extension-field
parameters are retained with the executed certificate.

The following tiny residue certificates show that the two
quadratic extensions are fields and are unequal as extensions
of the fixed field \(K\):
\[
\begin{array}{c|r|r|r|r}
\text{element}&r&\zeta\bmod r&\text{residue}&
                 \text{quadratic character}\\\hline
\Delta_{41}&61&33&17&-1\\
\Delta_{61}&41&39&15&-1\\
\Delta_{41}&41&39&5&1\\
\Delta_{41}\Delta_{61}&41&39&34&-1
\end{array}                                                     \tag{A.54.13}
\]
Every indicated residue is a unit and every indicated root
satisfies \(\Phi_{20}\). A square in \(K\) which is a unit
at such a place has square residue. Therefore
\[
 L_{41}=K(\sqrt{\Delta_{41}}),\qquad
 L_{61}=K(\sqrt{\Delta_{61}}),\qquad L_{41}\not\simeq_K L_{61}.
                                                               \tag{A.54.14}
\]

These fields survive every positive Frobenius power. For
\(q=41\), at the \(K\)-place \(\zeta\mapsto39\), the
trace is \(13\) modulo \(41\) and the determinant is zero.
For \(q=61\), the corresponding witness is
\(\zeta\mapsto24\), trace \(5\), determinant zero.
The determinant is a nonzero algebraic integer. The Newton
polygon of \(t^2-a_qt+d_q\) at this witness therefore has
one root of valuation zero and one of positive valuation.
Their quotient cannot be a root of unity. Thus, for every
\(n\ge1\), \(K[\pi_q^n]=K[\pi_q]=L_q\).
This is precisely what is needed when the endomorphisms of
\(B_8\) are defined only after finite extension.

\paragraph{Why the field certificates exclude every CM factor.}
First, if \(B_8\) has a nonzero CM isogeny factor, its maximal
CM sub-Hodge structure is preserved by every endomorphism,
hence by \(K\). Its rational dimension is therefore a positive
multiple of eight. A proper such part has abelian dimension four;
its quotient also has dimension four and inherits a faithful
unital \(K\)-action. A rank-one Hodge structure over the CM
field \(K\) is of CM type: its Mumford--Tate group is contained
in \(\operatorname{Res}_{K/\mathbf Q}\mathbf G_m\).
Semisimplicity of polarizable weight-one Hodge structures then
makes the whole \(B_8\) CM. This argument uses the quotient,
and does not assume the existence of a chosen \(K\)-stable
polarization complement.

Suppose, therefore, that \(B_8\) is CM. Put
\[
 \mathcal D=\operatorname{End}_{\mathrm{Hdg},K}
                  (H_1(B_8,\mathbf Q)).
\]
At a balanced embedding in (A.54.10), its complexification
preserves two distinct one-dimensional Hodge summands; hence
\(\dim_K\mathcal D\le2\). The Mumford--Tate group of a CM
Hodge structure is a torus, whose commutant on a two-dimensional
space has dimension at least two. Thus \(\mathcal D\) is a
quadratic etale \(K\)-algebra. These are actual rational
endomorphisms of the abelian variety, by the weight-one full
faithfulness in \cite[Theorems 4.4 and 7.12, pp.489 and 508;
CM discussion on pp.509 and 519]{fghMilneSVModuli}.
Choose a finite extension over which a basis of \(\mathcal D\)
is defined.

At either chosen good place, these endomorphisms specialize
injectively and commute with some positive power \(\pi_q^n\).
Indeed, they extend to the abelian scheme, and prime-to-\(q\)
Tate modules identify generic and special homomorphisms
faithfully. The required integral faithfulness is exactly
\cite[Theorem 10.15, pp.49--50]{fghMilneAV2008}; no surjectivity
part of a Tate isogeny theorem is used here.
In the rational endomorphism algebra of the special fibre consider
\[
 \mathcal C_q=\operatorname{Cent}_{\operatorname{End}^0(\overline B_{8,q})}
                          (K,\pi_q^n).
\]
For an auxiliary prime different from \(q\), its faithful
Tate realization lies in the centralizer of a nonscalar
separable quadratic operator on a rank-two
\(K\otimes\mathbf Q_\ell\)-module. That centralizer has
\(K\)-dimension two. Thus \(\dim_K\mathcal C_q\le2\);
it already contains \(K[\pi_q^n]=L_q\), so
\(\mathcal C_q=L_q\) as a global \(K\)-algebra.
The specialized \(\mathcal D\) has the same dimension and
therefore gives \(\mathcal D\simeq_K L_q\) at both places.
This contradicts (A.54.14).
The residue test at \(r=41\) in (A.54.13) is a number-field
square-class test, not an assertion using a prime-to-characteristic
Tate module at that same prime. This proves that \(B_8\)
has no geometric CM factor.

\paragraph{The normalizers away from split quartic primes.}
After any finite extension defining the factors, the Tate module
of \(B_8^5\) has no Galois-stable rational line:
\cite[\S1, Corollary 1]{Ruppert1998} would otherwise give a
nonzero CM abelian subvariety. For \(A_2\), the CM type
\(\{1,2\}\) has trivial stabilizer in
\((\mathbf Z/5)^\times\). Hence the four characters obtained
by splitting \(F\otimes\overline{\mathbf Q}_\ell\)
remain pairwise distinct on every open global Galois subgroup.
Here is the needed labelled argument. Equality of two characters
on such a subgroup implies equality of their Hodge--Tate weights
at every embedding of a finite field of definition into
\(\overline{\mathbf Q}_\ell\). The labelled weights are the
translates of the indicator of \(\{1,2\}\); its trivial
stabilizer forces the two embeddings to be identical.
The comparison and invariance under finite extension are
\cite[Theorems 2.2.3 and 2.4.6, pp.13 and 20]{fghBrinonConrad}.

If \(\ell\not\equiv1\pmod5\), every local field factor of
\(F\otimes\mathbf Q_\ell\) has degree greater than one
(including the totally ramified factor at five). A stable
\(\mathbf Q_\ell\)-line in \(V_\ell(A_2^5)\) would, after
projection to a nonzero factor and extension of scalars, have
the same character in every embedding in that local orbit.
The orbit has size greater than one, contradicting the
distinctness just proved. Thus no such line exists.

Now a noncuspidal normalized procyclic subgroup in \(X\)
has, by A.11, a nonzero proper-homology vector spanning a
line under an open arithmetic subgroup. Its projections to
\(A_2^5\) and \(B_8^5\) vanish by the preceding paragraphs.
It therefore projects nontrivially to the original elliptic
factor. The actual map \([5]q:X\to U\) is nonzero on precisely
that factor. Its image procyclic subgroup is nontrivial and
still has open arithmetic normalizer. The original normalizer
theorem forces that image to be cuspidal, with zero proper
homology, a contradiction. This proves the claimed normalizer
classification. Conversely every nontrivial cusp-inertia subgroup
has the required open arithmetic normalizer after making its
cusp rational.

At split primes the remaining possible proper line lies in
\(A_2^5\). The actual genus-two quotient in (A.54.7), with
its three cusps \(z=0\) and infinity, has a degree-five etale
map to the tripod via
\(t=(x+\sqrt\lambda)/(2\sqrt\lambda)\): its three missing
values are \(0,1,\infty\), and after scaling its equation is
\(w^5=t(t-1)\). The deck translates of the actual map
\(C\to D_\lambda/\sigma\) jointly detect the whole
\(A_2^5\) homology: their image is the \(H_5\)-span of the
\(H\)-invariant line in the simple module, tensored with the
one-dimensional \(T=-1\) multiplicity space. Thus a normalizer
theorem for this specific three-punctured genus-two target
settles the remaining split-prime source question once the target
criterion of A.55 is applied.
For \(\ell\ne5\) its degree-five map to the tripod does not
make its pro-\(\ell\) group an open subgroup of the tripod
group, so no split-prime conclusion follows from that map alone.
The executed exact certificate is retained with this proof.
\end{proof}

% A.55: independent root and DE full proof/source audits and exact reruns PASS.
\subsection{Cyclic auxiliary covers at the coefficient prime}

\begin{statement}{A.55}{Proved}{All-prime normalizers for the quartic-CM genus-two target and the Heisenberg-five cover}
Let
\[
 G:\ v^2=w^5-1,
 \qquad G^\circ=G\setminus\{(0,i),(0,-i),\infty\}.
                                                        \tag{A.55.1}
\]
Over a number field defining this curve, for every prime \(\ell\),
every nontrivial closed procyclic subgroup of
\(\pi_1(G^\circ_{\overline K})^{(\ell)}\) whose arithmetic
normalizer has open image is contained in a conjugate of cusp inertia.
The assertion is unchanged by a finite field extension or a twist
which becomes isomorphic to this punctured curve after such an extension.

Consequently the normalizer assertion of A.54 holds at \emph{every}
coefficient prime, including all \(\ell\equiv1\pmod5\). In particular
it holds at \(\ell=11\) for the actual genus-51, degree-125
Heisenberg cover in A.54. This statement adds no fixed-subgroup
claim to the pro-five fixed-subgroup result already proved there;
the all-prime full fixed-group comparison is supplied separately by A.58.
\end{statement}
\begin{proof}
\paragraph{The integral cyclic-cover calculation.}
We first record precisely the integral detector that will be used.
Let \(q:Y\to X\) be the connected unramified cyclic cover of
degree \(m=\ell^n\) of a proper genus-two curve defined by a
primitive character
\(\eta:H_1(X,\Z_\ell)\to\Z/m\).
Choose a symplectic basis \(a,b,c,d\) so that
\(\eta(b)=1\) and \(\eta(a)=\eta(c)=\eta(d)=0\).
The line generated by \(a\) modulo \(m\) is the Poincare-dual
line of \(\eta\), with an immaterial sign convention.
Then
\[
 q_*\bigl(H_1(Y,\Z_\ell)^{\operatorname{Deck}(q)}\bigr)
   =\Z_\ell a\oplus m\Z_\ell b
       \oplus m\Z_\ell c\oplus m\Z_\ell d.             \tag{A.55.2}
\]
Here and below these homology groups can equivalently be taken as
the integral Tate homology of the respective Jacobians.

For completeness this is an integral surface-group calculation.
Use the presentation
\(\langle a,b,c,d\mid[a,b][c,d]=1\rangle\) and
coset representatives \(1,b,\ldots,b^{m-1}\).
The Schreier generators are
\[
 a_j=b^jab^{-j},\quad c_j=b^jcb^{-j},\quad
 d_j=b^jdb^{-j}\quad(0\le j<m),\qquad B=b^m.
\]
After abelianization the \(m\) lifted surface relations are
\(a_j-a_{j+1}=0\), with indices taken cyclically. Their relation
lattice is saturated: eliminate \(a_1,\ldots,a_{m-1}\), leaving
the single generator \(h=[a_0]\). Thus an integral basis of
\(H_1(Y,\Z_\ell)\) is
\(h,B,[c_j],[d_j]\), of rank \(2m+2\).
The deck generator fixes \(h,B\) and cyclically permutes each
of the two \(m\)-tuples. Its invariant lattice has basis
\(h,B,\sum_j[c_j],\sum_j[d_j]\). Their pushforwards are
\(a,mb,mc,md\), proving (A.55.2).
In particular, a class whose rational image is in the old subspace
\(q^*H_1(X,\Q_\ell)\) satisfies
\[
      q_*[u]\bmod m\ \in\ (\Z/m)\operatorname{PD}(\eta).
                                                        \tag{A.55.3}
\]
Indeed the old rational subspace is exactly the deck-invariant
subspace, and a deck-invariant rational integral vector is
integrally invariant because the integral homology is torsion-free.
This implication is integral; it does not assume that \([u]\) or
\(q_*[u]\) is primitive.

\paragraph{The local auxiliary covers and their Pryms.}
We next prove the local geometric input, for arbitrary \(n\).
Let \(\mathcal X/R\) be a smooth proper curve over a complete
mixed-characteristic discrete valuation ring of residue
characteristic \(p\), and suppose a degree-zero line bundle
\(\mathcal L\) has order \(m=p^n\), with special line bundle
\(\overline{\mathcal L}\) still of exact order \(m\).
After a finite extension choose an isomorphism
\(\mathcal L^{\otimes m}\simeq\mathcal O\). It defines a finite
flat \(\mu_m\)-torsor
\[
 \mathcal Y=\operatorname{Spec}_{\mathcal X}
        \left(\bigoplus_{j=0}^{m-1}\mathcal L^{\otimes j}\right)
                         \longrightarrow\mathcal X,       \tag{A.55.4}
\]
where the multiplication uses this trivialization; reversing
\(\mathcal L\) changes only the convention. Its generic fibre
is a connected etale cyclic cover of degree \(m\).

The special fibre of (A.55.4) is integral and its normalized
function field is the purely inseparable extension
\(k(X)^{1/p^n}\) of \(k(X)\). To see both assertions, represent
\(\overline{\mathcal L}\) by a divisor \(D\) and its trivialization
by \(f\) with \(\operatorname{div}(f)=mD\). If \(f=g^p\),
then \(p^{n-1}D=\operatorname{div}(g)\), contrary to the exact
order of \(\overline{\mathcal L}\). Hence \(f\notin k(X)^p\),
and adjoining an \(m\)-th root gives degree \(m\).
A one-variable function field over a perfect field has
\([k(X)^{1/p^n}:k(X)]=p^n\), so the resulting field is exactly
the displayed Frobenius extension. Finite flatness makes each
special affine ring torsion-free over the base affine ring;
its injection into this generic field proves that it is a domain.

The total space \(\mathcal Y\) is normal. Locally it is the
hypersurface \(z^m=u\), with \(u\) a unit, over a regular
two-dimensional base, so it satisfies \(S_2\). Horizontal
codimension-one points lie in the smooth generic fibre. At the
unique vertical generic point the special fibre is generically
reduced, so the local maximal ideal is generated by a base
uniformizer and the local ring is regular. This proves \(R_1\),
hence normality. The same reasoning applies after every finite
extension of the valuation ring, since the special order of
\(\overline{\mathcal L}\) is unchanged.

Apply \citet[Theorem 1, printed p.179]{fghRaynaud1990Pgroups}
to this actual unramified \(p\)-group cover of the generic fibre
of the smooth proper model. Its Jacobian has potentially good
reduction. In a semistable model the graph is a tree, and the
components which map to points of the original special fibre
have crystalline slopes strictly between zero and one.
There is exactly one nonexceptional component: it is the
normalization \(X^{(p^{-n})}\) just identified, with its radicial
degree-\(m\) map to \(X\). The reduced special fibre is
geometrically unibranch because this normalization is radicial;
thus there is no additional nonexceptional component hidden at
a singular point.

Let \(P_q=(\ker(q_*:J_Y\to J_X))^0\) be the Prym.
After taking good models, the norm extends as a homomorphism of
abelian schemes. The special Jacobian of the compact-type fibre
is the product of its component Jacobians. On the unique main
factor the special norm is the isogeny induced by the radicial
Frobenius map; on every exceptional factor it is zero, because
that component maps to a point. It follows, up to isogeny, that
\[
        \overline{P_q}\ \sim\
           \prod_{Z\ \mathrm{exceptional}}\operatorname{Jac}(Z).
                                                        \tag{A.55.5}
\]
This can also be read from the rational idempotent
\(1-m^{-1}q^*q_*\), which extends after multiplication by \(m\).
In particular, \(P_q\) has potentially good reduction and
special \(p\)-rank zero. The main factor need not have
\(p\)-rank zero; it is precisely the old factor removed in
(A.55.5).

The rank-one crystalline argument of A.46 now gives
\[
 V_p(P_q)\text{ has no nonzero }\Q_p\text{-line stable under
 any open local or global Galois subgroup}.              \tag{A.55.6}
\]
Indeed, dualizing a stable homological line gives a rank-one
crystalline quotient of \(H^1\). In the cohomological convention
its Hodge--Tate weights are zero or one; the rank-one
weak-admissibility argument of A.46 would give Newton slope
zero or one. Neither slope occurs in (A.55.5). This conclusion is unchanged
by further finite extensions. Notice that Raynaud's theorem is
used for the full cyclic group of order \(p^n\); a calculation
only for degree \(p\) would not suffice below.

\paragraph{The split coefficient primes on the actual genus-two curve.}
The Jacobian of \(G\) has full CM by \(\Z[\zeta_5]\), with
type \(\{1,2\}\). A.54 proves that its four labelled
one-dimensional CM characters remain distinct on every open
Galois subgroup. At a nonsplit coefficient prime there is no
\(\Q_\ell\)-line stable under an open subgroup, by that same
argument, so A.11 already proves the assertion.

It remains to consider \(p=\ell\equiv1\pmod5\).
After a finite field extension defining all endomorphisms, the
integral Tate lattice has four rank-one eigensummands for
\(\Z_p[\zeta_5]\simeq\Z_p^4\). They form the two perfect
symplectic pairs with labels \(\{1,4\}\) and \(\{2,3\}\).
The integral decomposition is valid because \(p\ne5\).
The curve has good ordinary reduction at a place over \(p\).
For a direct check, write \(p=10k+1\). On the basis
\(dw/v,w\,dw/v\), the Cartier matrix of \(v^2=w^5-1\)
over \(\F_p\) is diagonal, with nonzero diagonal entries
\[
        (-1)^{3k}\binom{5k}{2k},\qquad
        (-1)^k\binom{5k}{4k}.                         \tag{A.55.7}
\]
Indeed its entries are the coefficients of \(w^{pi-j}\)
in \((w^5-1)^{(p-1)/2}\), with indices \(1\le i,j\le2\);
the off-diagonal exponents are not divisible by five. The
binomial coefficients are units since \(5k<p\).
The coefficient formula follows directly by applying Cartier
to \(w^{j-1}v^{-p}(w^5-1)^{(p-1)/2}dw\), as in
\cite[Section 3, pp.5--6]{fghAchterHowe2017}.
Invertibility is the ordinary criterion.

The four CM idempotents extend to the good abelian scheme and
split its \(p\)-divisible group into height-one factors.
Polarization makes the two factors in each symplectic pair
Cartier dual; consequently exactly one in each pair is etale
and the other multiplicative. In particular, in either chosen
pair there are compatible Picard torsion points \(L_n\) of
order \(p^n\) whose specializations have exactly that order.
They are defined over finite number fields at each finite
level. Their line bundles extend over the good curve model:
use the relative degree-zero Picard scheme and rigidify at a
rational point, after another finite extension. The trivialization
\(L_n^{\otimes p^n}\simeq\mathcal O\) also extends after
rigidification; any scalar discrepancy is absorbed after a
finite extension. Thus these are the actual line bundles
required in (A.55.4), not merely classes in a local vector space.

Suppose a noncuspidal procyclic subgroup \(A\) of the
geometric pro-\(p\) group of \(G^\circ\) has open arithmetic
normalizer. Choose a generator \(u\). By A.11 its proper
homology class \(v\) is nonzero, and its rational span is
stable under an open Galois subgroup. Distinctness of the CM
characters puts \(v\) in a single one of the four integral
eigenlines. Write \(v=p^r v_0\), where \(v_0\) is primitive.
Choose the etale Picard eigenline in the \emph{other} symplectic
pair and take \(n>r\).

Kummer duality and the principal polarization identify
\(L_n\in J_G[p^n]\) with the character
\(\eta_n:H_1(G,\Z_p)\to\mu_{p^n}\) obtained by Weil pairing.
After fixing a primitive root, identify its target with
\(\Z/p^n\). Its Poincare-dual line is the chosen Picard
line, up to a unit and the harmless choice of signs. Hence
\[
       \eta_n(v)=0,
       \qquad
       v\bmod p^n\notin
          (\Z/p^n)\operatorname{PD}(\eta_n).          \tag{A.55.8}
\]
The second assertion uses \(n>r\) and the direct integral
decomposition into the two different symplectic pairs. This
is why the other pair was chosen; choosing the same dual line
could give no contradiction to (A.55.3).

Let \(q_n:Y_n\to G\) be the actual cyclic unramified proper
cover supplied by this Kummer torsor, and set
\(Y_n^\circ=q_n^{-1}(G^\circ)\). The kernel of \(\eta_n\)
contains \(A\); the geometric pro-\(p\) group of
\(Y_n^\circ\) is that open index-\(p^n\) subgroup, because
the pro-\(p\) topology on an open subgroup is induced.
Choose the corresponding lift of the base point. Descend the
cover and this point to a finite field extension and intersect
the original normalizer with its arithmetic fundamental group.
This intersection still has open arithmetic image and still
normalizes \(A\). Thus the proper class
\(\widetilde v\in H_1(Y_n,\Z_p)\) of the \emph{same}
generator \(u\) spans a stable line, with
\(q_{n,*}\widetilde v=v\ne0\). These are actual integral
classes in the covering group, not rational Malcev lifts.

The projection of this line to \(V_p(P_{q_n})\) must vanish
by (A.55.6). Therefore \(\widetilde v\) is in the old rational
homology, and (A.55.3) gives the opposite of (A.55.8).
This contradiction proves the normalizer criterion at every
split prime. The choice of \(n\) also handles arbitrary
divisibility of the original proper class.

\paragraph{Return to the Heisenberg cover.}
For the actual cover of A.54, its proper homology is
\(V_\ell E_i\oplus V_\ell A_2^{\oplus5}
\oplus V_\ell B_8^{\oplus5}\).
The \(B_8\) summand has no potentially stable rational line,
as proved there. A line projecting nontrivially to \(E_i\)
is excluded by the original punctured elliptic covering map.
The remaining lines are detected by the deck translates of
the actual maps to \(D_\lambda/\sigma\) constructed in A.54;
their proper targets are twists of \(G\), and their removed
fibres are exactly the three points in (A.55.1).
Indeed, after a finite scaling, the target is
\(u^5=t(t-1)\), related to (A.55.1) by
\(t=(v+i)/(2i)\) and \(w=\beta u\), \(\beta^5=-4\).
The maps are therefore maps of the actual punctured curves.
Their nonzero images would have open arithmetic normalizer
in the target pro-\(\ell\) group and nonzero proper homology,
contradicting the criterion just proved. Statement A.11
finishes the all-prime assertion for the Heisenberg cover.

The exact certificate \texttt{code/a\_quartic\_cm\_normalizers.py} checks the
Schreier relations, integral deck invariant lattice and its
pushforward for prime-power degrees, and (A.55.7) at all
split primes below a fixed bound. The uniform statements are
proved above and are not inferred from the finite checks.
\end{proof}

% A.56 independently audited by root and DE; see notes/fgh_A56_audit.md.
\subsection{A reduction criterion for arbitrary finite covers}

\begin{statement}{A.56}{Proved}{Cuspidal normalizers from special rank and at all sufficiently large Gaussian split primes}
Let \(Y\) be an affine hyperbolic curve over a number field,
with smooth proper compactification \(C\) of genus \(g\ge2\),
and let \(p\) be a prime. Suppose that \(C\) has potentially
good \emph{curve} reduction at some place above \(p\), and that
the good special Jacobian has \(p\)-rank at least two.
Then every nontrivial closed procyclic subgroup of
\(\pi_1(Y_{\overline K})^{(p)}\) with open arithmetic normalizer
is contained in cusp inertia.

More specifically, let \(Y\to U=E\setminus\{O\}\) be any
geometrically connected finite etale cover, where
\(E:y^2=x^3-x\), and let \(C\to E\) be its proper
compactification. If \(p\equiv1\pmod4\) and \(C\) has
potentially good curve reduction at a place above \(p\),
the same normalizer conclusion holds, with no special-rank
hypothesis on \(C\).

In particular, for every such cover of degree \(d\), its
pro-\(p\) normalizer criterion holds for every
\[
                         p>d,\qquad p\equiv1\pmod4.
                                                        \tag{A.56.1}
\]
These assertions concern actual normalized procyclic subgroups.
They do not assert full fixed-subgroup comparisons or a
cofinal family of covers.
\end{statement}
\begin{proof}
\paragraph{Choosing the Picard torsion without a CM splitting.}
After a finite number-field extension take a good proper curve
model at the selected place, and let \(J\) be its Jacobian.
Put
\[
 T=T_pJ,\qquad
 s:T\longrightarrow T_p^{\rm et}\overline J,
 \qquad K=\ker s.                                    \tag{A.56.2}
\]
The specialization map is surjective and its kernel is saturated.
Indeed the connected--etale sequence of the \(p\)-divisible group
over the complete strictly henselian valuation ring gives this
sequence on geometric generic Tate modules; see
\cite[Sections 1.4 and 2.2, printed pp.160 and164,
and Section 2.4, p.169]{fghTate1967Pdivisible}.
Surjectivity at each geometric generic finite level and the
surjective transition maps on the connected kernels give
surjectivity on their inverse limits. The etale quotient
specializes isomorphically to the geometric special Tate module.
If the special \(p\)-rank is \(f\), then
\(\operatorname{rank}T=2g\) and
\(\operatorname{rank}K=2g-f\). We do not assert that this
sequence splits as integral arithmetic representations.

Suppose \(A\) is a noncuspidal procyclic subgroup with open
arithmetic normalizer, choose a generator \(u\), and let
\(v=p^r v_0\ne0\) be its proper integral homology class,
with \(v_0\) primitive. Its existence and the stability of
its rational span follow from A.11. Use the principal
polarization to regard \(T\) as the symplectic integral
homology lattice. In \(T/pT\), the hyperplane
\(\overline v_0^\perp\) has dimension \(2g-1\).
Its intersections with \(K/pK\) and with the line
\(\F_p\overline v_0\) are both proper subspaces: for the
first use \(f\ge2\), and for the second use \(g\ge2\).
The union of two proper linear subspaces cannot be the whole
space, even over \(\F_2\). Choose
\[
 \overline w\in\overline v_0^\perp,
 \qquad \overline w\notin K/pK,
 \qquad \overline w\notin\F_p\overline v_0.             \tag{A.56.3}
\]
Since \(v_0\) is primitive and the symplectic pairing perfect,
the integral map \(T\to\Z_p(1)\) given by pairing with
\(v_0\) is surjective. Its kernel reduces onto the displayed
hyperplane. Thus \(\overline w\) lifts to
\(w\in v_0^\perp\), and \(s(w)\) is primitive.

For \(n>r\), let \(L_n\in J[p^n]\) be the point represented
by \(w\bmod p^n\). Its specialization has exact order
\(p^n\). Each finite-level point is defined over a finite
number field. After rigidification at a rational point it gives
an actual degree-zero line bundle on the good curve model,
with its \(p^n\)-th power trivialized, exactly as in A.55.
The resulting character \(\eta_n\), defined by Kummer
duality and polarization, satisfies
\[
       \eta_n(v)=0,
       \qquad v\bmod p^n\notin
               (\Z/p^n)\operatorname{PD}(\eta_n).     \tag{A.56.4}
\]
For the second assertion, an equality
\(p^r v_0=a w\pmod {p^n}\) forces \(p^r\mid a\),
because \(w\) is primitive. Divide by \(p^r\) and reduce
modulo \(p\), contradicting (A.56.3).

\paragraph{The actual cover and its new homology.}
The \(\mu_{p^n}\)-torsor defined by \(L_n\) gives an
unramified proper cyclic cover \(q_n:C_n\to C\). The
normality and radicial-main-component proof in A.55 applies
without change for every genus. The main component is
\(\overline C^{(p^{-n})}\); the norm on its Jacobian is a
Frobenius isogeny. Raynaud's Theorem 1 therefore identifies
the special Prym, up to isogeny, with the product of
exceptional component Jacobians, all with slopes strictly
between zero and one. By A.46 its Tate representation has
no rational line stable under an open Galois subgroup.

The subgroup \(A\) lies in the index-\(p^n\) covering
subgroup because \(\eta_n(v)=0\). Pass to the actual
punctured curve \(q_n^{-1}Y\), and intersect the normalizer
with its arithmetic fundamental group after a finite
descent extension. The arithmetic image remains open.
The class \(\widetilde v\) of the same generator \(u\)
in \(H_1(C_n,\Z_p)\) spans a stable line and pushes
forward to \(v\). Its Prym projection vanishes, so it is
in the old rational homology and is integrally invariant
under the cyclic deck group.

The integral calculation (A.55.2) works in every genus.
For a symplectic basis \(a_1,b_1,\ldots,a_g,b_g\)
adapted to \(\eta_n(b_1)=1\), the Schreier homology has
the two fixed generators \(h,B\), and two regularly
permuted \(p^n\)-tuples for each \(j\ge2\). Thus
\[
 q_{n,*}H_1(C_n,\Z_p)^{\operatorname{Deck}(q_n)}
 =\Z_p a_1\oplus p^n\Z_p b_1
       \oplus\bigoplus_{j=2}^g
             (p^n\Z_p a_j\oplus p^n\Z_p b_j).         \tag{A.56.5}
\]
This contradicts (A.56.4), proving the first assertion.

\paragraph{The once-punctured Gaussian elliptic target.}
Now let \(q:C\to E\) compactify the given cover. Its
nonconstant map makes \(E\) an isogeny factor of \(J_C\):
\(q_*q^*=[\deg q]\), so
\[
                     J_C\sim E\times P_q.              \tag{A.56.6}
\]
At \(p\equiv1\pmod4\), the good elliptic reduction of
\(E\) is ordinary. One can check this directly from the
nonzero central coefficient of \((x^3-x)^{(p-1)/2}\);
equivalently it is the usual split-prime ordinary
criterion for CM by \(\Z[i]\).
After a finite extension giving the good curve model of
\(C\), all factors in (A.56.6) have good reduction.
Special \(p\)-rank is additive in products and invariant
under isogeny, hence
\(f(J_C)=1+f(P_q)\ge1\).

If \(f(J_C)\ge2\), the first assertion applies.
If \(f(J_C)=1\), the Prym has special \(p\)-rank zero.
Statement A.46 forces any stable rational line in
\(V_pJ_C\) to project trivially to \(V_pP_q\).
It therefore projects nontrivially under \(q_*\) to
\(V_pE\). The actual punctured morphism
\(Y\to E\setminus\{O\}\) carries a proposed normalized
procyclic subgroup with this line to a nontrivial
arithmetically normalized subgroup with nonzero proper
homology. The already proved original elliptic criterion
puts its image in cusp inertia, whose proper homology
vanishes, a contradiction. This also covers proper
genus one, when \(q_*\) is an isogeny. Thus no genus
restriction is needed for the second assertion.

\paragraph{A degree bound which ensures good curve reduction.}
Suppose the original connected cover has degree \(d\),
and \(p>d\). The geometric Galois closure has group a
subgroup of \(S_d\), so its order is prime to \(p\).
The marked elliptic curve \((E,O)\) has good reduction
at this odd prime. The prime-to-\(p\) specialization
theorem for a smooth proper curve with a relative
normal-crossing boundary identifies its geometric
generic prime-to-\(p\) covers with the corresponding
special covers, preserving inertia; see
\cite[Expos\'e XIII, Corollary 2.9, Section 2.10,
and the proof of Corollary 2.12, new pp.288--292]{fghSGA1}.
Consequently, after a finite extension, the actual
Galois closure extends tamely over the good pair
\((\mathcal E,\mathcal O)\).

This extension has smooth proper source. Away from the
marked section it is etale. At a point above the section,
after an etale base change its completed local map is
\(R[[t]]\to R[[z]],\ t=z^e\), with \(e\) prime to
\(p\); the source is smooth over \(R\). The same
local description, or the subgroup quotient of this
tame cover, gives a smooth model for the original
non-Galois curve \(C\). Its good reduction is thus
potential, with no additional unspecified exceptional
prime set. Applying the second assertion gives
(A.56.1).
\end{proof}

% A.57 full independent root and DE audits PASS; see notes/fgh_A57_audit.md.
\subsection{The rank-one case and all sufficiently large coefficient primes}

\begin{statement}{A.57}{Proved}{Potential good curve reduction excludes every noncuspidal pro-p normalizer}
Let \(Y\) be an affine hyperbolic curve over a number field, with
smooth proper compactification \(C\) of genus at least two.
If \(C\) has potentially good curve reduction at a place above
the coefficient prime \(p\), every nontrivial closed procyclic
subgroup of \(\pi_1(Y_{\overline K})^{(p)}\) whose arithmetic
normalizer has open image is contained in cusp inertia.
There is no condition on the special \(p\)-rank, and \(p=2\)
is included.

Consequently, for every geometrically connected degree-\(d\)
finite etale cover \(Y\to E\setminus\{O\}\), where
\(E:y^2=x^3-x\), the normalizer conclusion holds at every
coefficient prime
\[
                              p>d.                     \tag{A.57.1}
\]
More generally it holds whenever its proper compactification
has potentially good curve reduction at a place above \(p\).
These are assertions about actual normalized procyclic subgroups;
no fixed-subgroup calculation or fixed-element lifting is asserted.
\end{statement}

\begin{proof}
The special-rank-zero case is A.46, and special rank at least
two is A.56. We prove the remaining rank-one case, including
the precise reduction fact required for arbitrary cyclic covers.

\paragraph{An unramified p-cover above a rank-one special curve.}
Let \(R\) be a complete discrete valuation ring of mixed
characteristic \((0,p)\), let \(\mathcal C/R\) be a smooth
proper curve with geometrically connected fibres, and suppose
\(f(\mathcal C_{\overline k})=1\). Let
\(q:C'_{K_R}\to C_{K_R}\) be a geometrically connected finite
etale Galois cover of \(p\)-power degree; the cyclic case will
suffice below. Raynaud's Theorem 1 gives potentially good
reduction of \(J_{C'}\) and a semistable source with a tree
as dual graph. Every component contracted to a point of
\(\mathcal C_k\) has crystalline slopes strictly between
zero and one \cite[Theorem 1, printed p.179]{fghRaynaud1990Pgroups}.
We claim that the good special Jacobian of \(C'\) has
\(p\)-rank exactly one.

First make a finite extension on which the source is semistable,
and then work over the completion of the maximal unramified
extension. Its residue field is algebraically closed. This
last extension only computes geometric components and ranks;
all good-reduction assertions already descend to finite
extensions. Re-normalize the smooth base in the extended
generic covering field, obtaining a finite normal map
\[
                   \mathcal T\longrightarrow\mathcal C.
                                                        \tag{A.57.2}
\]
Finiteness follows from excellence of a complete discrete
valuation ring. Its generic source remains geometrically
connected. Properness gives a finite \(R\)-algebra
\(H^0(\mathcal T,\mathcal O_{\mathcal T})\); restriction to
the dense generic fibre embeds this algebra in \(K_R\),
because the latter fibre is a smooth geometrically connected
proper curve. Every element is integral over \(R\), so normality
of \(R\) gives
\[
                  H^0(\mathcal T,\mathcal O_{\mathcal T})=R.
                                                        \tag{A.57.3}
\]
Stein factorization therefore makes the special fibre connected.
This argument applies after every finite extension and after
the indicated geometric residue extension: it proves the
geometric connectedness needed here, rather than only
connectedness over the original finite residue field.

If (A.57.2) is etale, its special fibre is already a connected
etale Galois \(p\)-cover of \(\mathcal C_k\). Otherwise
purity of the branch locus on the regular surface
\(\mathcal C\) says that its reduced branch support is
the entire special fibre. Indeed the generic map is etale,
and the only possible vertical prime divisor is the smooth
irreducible fibre \(\mathcal C_k\); there are no isolated
codimension-two branch points. This branch support is
geometrically unibranch.

Apply Raynaud's Proposition 1 to (A.57.2). Its part (i)
makes the reduced inverse image of this branch support
geometrically unibranch. In particular it is locally
irreducible. A locally irreducible noetherian scheme has
pairwise disjoint irreducible components, which are open
and closed. Connectedness from (A.57.3) therefore gives a
\emph{single} irreducible special component \(T_0\).
Its decomposition group is the entire Galois group \(G\).
Part (ii) makes its generic inertia \(I\) constant at all
points of \(T_0\); it is normal in \(G\). Put
\(\mathcal T'=\mathcal T/I\), and let \(T'_0\) be the
image of \(T_0\). The same proposition states that
\(\mathcal T'\to\mathcal C\) is etale on a neighborhood
of \(T'_0\), and that \(T'_0\to\mathcal C_k\) is an
etale Galois cover of group \(G/I\)
\cite[Proposition 1(i)--(ii), printed pp.179--180]{fghRaynaud1990Pgroups}.

Here the neighborhood assertion is sufficient globally:
\(T'_0\) is the \emph{whole} special fibre as a topological
space, because \(\mathcal T\) has only one special component
and the finite quotient is surjective. The quotient map is
also etale on the entire generic fibre. Every point is in
one of these fibres, so
\[
                  \mathcal T'\longrightarrow\mathcal C
                  \quad\hbox{is finite etale everywhere.}
                                                        \tag{A.57.4}
\]
Thus \(T'_0\) is smooth, proper and connected. The map
\(T_0\to T'_0\) is the homeomorphism in Proposition 1;
on function fields its extension is purely inseparable.
Equivalently the inertia quotient has removed the whole
separable residue extension. The normalization \(Z_0\)
of \(T_0\) is consequently a finite radicial cover of
\(T'_0\). Over the perfect residue field a finite purely
inseparable map of smooth curves is a power of relative
Frobenius followed by an isomorphism: their function fields
have \(p\)-basis of cardinality one. Hence
\[
                            f(Z_0)=f(T'_0).             \tag{A.57.5}
\]
Alternatively their etale sites agree under this universal
homeomorphism, giving the same \(H^1(-,\F_p)\).
No equality between the radicial degree and \(|I|\) is
needed; possible vertical multiplicities play no role.

The unramified Deuring--Shafarevich formula gives
\[
                  f(T'_0)-1
                    =|G/I|\bigl(f(\mathcal C_k)-1\bigr)=0.
                                                        \tag{A.57.6}
\]
For the precise unramified formula see
\cite[Corollary 1.8 and equation (1.8.2), printed pp.35--36]
{fghCrew1984}. In the etale case of (A.57.2), apply this
formula directly with \(G\), obtaining the same conclusion.

The semistable source has just this one noncontracted
component: its function field is the residue function field
of the unique vertical valuation of \(\mathcal T\).
Every other component is contracted to a point of the
smooth base. The tree assertion and strict slopes in
Raynaud's Theorem 1 therefore give
\[
           f(\overline J_{C'})=f(Z_0)+
                  \sum_{Z\ {
m contracted}}f(J_Z)=1.  \tag{A.57.7}
\]
Here the special Jacobian in a compact-type model is the
product of the component Jacobians, and strict slopes
imply zero \(p\)-rank. This proves the claim.

\paragraph{The new Prym has no rational arithmetic line.}
For any such cyclic cover of degree \(m=p^n\), put
\(P_q=\ker(q_*:J_{C'}\to J_C)^0\).
The identity \(q_*q^*=[m]\) gives the actual isogeny
\[
                           J_{C'}\sim J_C\times P_q.   \tag{A.57.8}
\]
All its factors have good reduction after a finite extension,
by Raynaud and the good reduction of \(J_C\). Special
\(p\)-rank is additive in products and invariant under
isogeny. Equations (A.57.7)--(A.57.8) show
\(f(\overline P_q)=1-1=0\).
Statement A.46 then excludes every one-dimensional
\(\Q_p\)-subquotient of \(V_pP_q\), even after arbitrary
finite global or local extension. In particular no new
homology line can be preserved by an open arithmetic group.
This assertion holds for \emph{every} geometrically connected
cyclic \(p^n\)-cover of the good rank-one curve; we have not
imposed a specialization-order condition on its Picard point.

\paragraph{Choosing an arbitrary detecting cyclic cover.}
Now suppose the proposed noncuspidal normalized subgroup
has generator \(u\). By A.11 its proper homology class is
\(v=p^r v_0\ne0\), with \(v_0\) primitive in the
symplectic lattice \(T=H_1(C_{\overline K},\Z_p)\).
Since \(g\ge2\), its mod-\(p\) perpendicular space has
dimension \(2g-1\ge3\). Choose
\[
       \overline w\in\overline v_0^\perp
                           \setminus\F_p\overline v_0.
                                                        \tag{A.57.9}
\]
Perfect integral duality lifts \(\overline w\) to a
primitive \(w\in v_0^\perp\), exactly as in A.56.
Here no condition on the special image of \(w\) is needed.
For \(n>r\), the primitive character
\(\eta_n:T\to\Z/p^n\), given by pairing with
\(w\bmod p^n\) after choosing a root-of-unity generator,
defines a geometrically connected degree-\(p^n\) etale
proper cyclic cover. It descends with its deck action
after a finite number-field extension. Its Picard point
and an appropriate line-bundle trivialization also give
this cover directly. By construction
\[
       \eta_n(v)=0,
       \qquad v\bmod p^n\notin
                    (\Z/p^n)\operatorname{PD}(\eta_n).
                                                        \tag{A.57.10}
\]
The latter assertion follows by dividing a hypothetical
equality \(p^rv_0=a w\pmod {p^n}\) by \(p^r\)
and reducing modulo \(p\); primitivity of \(w\)
first forces \(p^r\mid a\).

The same generator \(u\) lies in the actual covering
subgroup of the pro-\(p\) fundamental group because
\(\eta_n(v)=0\). This covering subgroup is the pro-\(p\)
group of the actual punctured cover \(q^{-1}Y\), since
the cover has \(p\)-power Galois degree. Intersect the
arithmetic normalizer with its arithmetic covering group,
after the finite descent extension. The intersection has
finite index in the original normalizer, hence still has
open arithmetic image. The proper integral class
\(\widetilde v\) of \(u\) on \(C'\) pushes to \(v\)
and spans a stable rational line. Its projection to the
Prym in (A.57.8) must vanish by the preceding paragraph.
Thus \(\widetilde v\) lies in the rational old homology,
which equals the deck-invariant rational subspace.
The integral homology is torsion-free, so it is also
integrally invariant under the deck group.

The actual integral Schreier calculation of A.55--A.56,
in a basis adapted to \(\eta_n(b_1)=1\), gives
\[
 \begin{split}
 q_*H_1(C',\Z_p)^{\operatorname{Deck}(q)}
   ={}&\Z_p a_1\oplus p^n\Z_p b_1\\
     &{}\oplus\bigoplus_{j=2}^g
                  (p^n\Z_p a_j\oplus p^n\Z_p b_j).
 \end{split}                                            \tag{A.57.11}
\]
Modulo \(p^n\) this is precisely the
\(\operatorname{PD}(\eta_n)\)-line. Its membership for
\(v=q_*\widetilde v\) contradicts (A.57.10).
The choice \(n>r\) is why this argument handles a
generator with divisible proper homology; no extraction
of a nonexistent root of \(u\) is used. This completes
the rank-one case and the first assertion.

\paragraph{Arbitrary covers of the punctured elliptic curve.}
A nonconstant proper map to \(E\) has source genus at
least one. Genus one is already A.40 at every coefficient
prime. For genus at least two the first assertion applies
whenever the stated potential good curve reduction holds.
If the cover has degree \(d\ge2\) and \(p>d\), then
\(p\) is odd and \((E,O)\) has good marked reduction.
Its geometric Galois closure embeds in \(S_d\), so has
order prime to \(p\). The prime-to-\(p\) specialization
and tame local model proved explicitly in A.56 give
potentially good reduction of the proper covering curve.
That argument imposed no congruence on \(p\); the
congruence in A.56 was only used for its earlier rank-one
case. Thus the present criterion proves (A.57.1) for
every \(p>d\). For \(d=1\), including \(p=2\), use
the original elliptic criterion or A.40 directly.
\end{proof}

% A.58 full independent root and DE proof/source/certificate audits PASS.
\subsection{The actual five-adic Heisenberg model}

\begin{statement}{A.58}{Proved}{All-prime fixed groups in the genus-fifty-one Heisenberg cover}
Let \(C\), \(q:C\to E\), and \(Y=C\setminus q^{-1}(E[5])\)
be the actual degree-one-hundred-twenty-five Heisenberg cover of A.54:
\[
              E:y^2=x^3-x,\qquad z^5=\psi _5(x).
\]
After a finite extension at five its proper stable reduction has
one ordinary elliptic component, five genus-six components of the
form \(V^5-V=U^4\), and ten genus-two components of the form
\(V^5-V=U^2\), over the algebraic residue field. Every other
component is rational, and the dual graph is a tree. The elliptic
component maps radicially with degree five to \(E\); each other
positive-genus component maps to a distinct point outside the
reduced support of \(E[5]\).

Every selected cusp of \(Y\) has a second cusp joined to it by
a rational path in the pointed stable fibre. Consequently, for
every coefficient prime, every selected cusp packet, and every
indicated open derived subgroup, the full geometric pro-prime
fixed subgroup is exactly the selected cusp inertia. Together
with A.54--A.55, both assertions of Task A.5 are therefore proved
at every prime for this actual Heisenberg cover.
\end{statement}

\begin{proof}
Put
\[
\begin{aligned}
 g(x)&=3x^2+1,\\
 h(x)&=\frac{\psi _5(x)-g(x)^5}{5}\\
     &=x^{12}-61x^{10}-102x^8+6x^6-43x^4+7x^2.
\end{aligned}                                                   \tag{A.58.1}
\]
These are exact polynomial identities in the actual covering
field. At five the indicated elliptic model is smooth and ordinary.
For example its Cartier coefficient is the coefficient of \(x^4\)
in \((x^3-x)^2\), namely \(-2\ne0\pmod5\).
We work over a finite local extension containing the constants
used below and admitting semistable reduction; such constants
may also be included in a finite global field of definition.

\paragraph{The actual elliptic main component.}
Choose \(\pi^5=5\), and set \(z=g(x)+\pi u\). The exact equation is
\[
 u^5+\pi g^4u+2\pi^2g^3u^2+2\pi^3g^2u^3+\pi^4gu^4=h(x),
                   \qquad y^2=x^3-x.                 \tag{A.58.2}
\]
Its special function field is \(k(E)(h^{1/5})\). In
characteristic five the exact factorization is
\[
          dh=2x(x^2-2)^3(x^2-x+1)(x^2+x+1)\,dx\ne0.
                                                        \tag{A.58.3}
\]
The one-variable function field \(k(E)\) has \(p\)-basis
of cardinality one. Hence this extension is \(k(E)^{1/5}\),
and its smooth projective normalization is the Frobenius twist
of \(E\), of genus one. The chart is smooth at its special
generic point: the nonzero differential in (A.58.3) gives the
required Jacobian rank together with the elliptic equation.
Thus this is an actual positive-genus component of a proper
model of \(C\), with degree-five radicial map to \(E\).

\paragraph{The fifteen residue centres.}
On the open subset where \(g\) is a unit put
\[
              F=\frac{\psi _5}{g^5}=1+5H,
                         \qquad H=\frac h{g^5}.        \tag{A.58.4}
\]
In characteristic five, \(dH=dh/g^5\). Its zeros, counted
on the elliptic curve with a genuine local parameter, are as follows:
\[
\begin{array}{c|c|c|c}
 \text{centres}&\text{number of points}&\operatorname{ord}(dH)&k\\
 \hline
 (0,0)&1&3&4\\
 x^2=2&4&3&4\\
 (1,0),\ (-1,0)&2&1&2\\
 x^4+x^2+1=0&8&1&2
\end{array}                                                   \tag{A.58.5}
\]
At a two-torsion point \((a,0)\), the parameter is \(y\),
and \(x-a\) has order two while \(dx\) has order one.
At all other displayed points \(x-a\) is a parameter.
The polynomial factors in (A.58.3) are square-free and
pairwise disjoint except for the displayed multiplicity three;
they meet neither \(g=0\) nor any additional two-torsion point.
This proves (A.58.5) directly. Its zero orders total twenty-five.

Choose lifts of these centres in a finite \emph{unramified}
extension \(R_0/\Z_5\). For the non-two-torsion centres take
\(a^2=2\) or \(a^4+a^2+1=0\), respectively, and a root
\(y_0^2=a^3-a\). All required residue roots are simple,
and \(y_0\) is a unit. The other lifts are the actual
points \((0,0),(1,0),(-1,0)\). The unramified choice
ensures that a coefficient reducing to zero is divisible
by \(5\), not merely by a smaller uniformizer after a
ramified extension.

\paragraph{Exact algebraic tail charts.}
For any selected lift \((a,y_0)\), let
\[
 F_0=F(a),\qquad H_0=H(a),\qquad
 N_a(x)=h(x)-H_0g(x)^5,
 \qquad z_0^5=F_0.                                    \tag{A.58.6}
\]
Here \(F_0\in R_0^\times\), \(F_0\equiv1\pmod5\),
and \(N_a(a)=0\) exactly. Enlarge the field to choose
\(\eta^4=-5\) and \(\rho^k=\eta\), where \(k\) is
the last column of (A.58.5). The exact covering-coordinate change is
\[
                    z=g(x)z_0(1+\eta V).
\]
After subtraction of the constant term and division by \(5\eta\),
the second equation is
\[
 g(x)^5F_0
   \bigl(V+2\eta V^2+2\eta^2V^3+\eta^3V^4-V^5\bigr)
                         =\frac{N_a(x)}{\eta}.        \tag{A.58.7}
\]
The equality follows by the binomial theorem and \(\eta^4=-5\);
no formal truncation has been made.

For \(y_0\ne0\), substitute
\(x=a+\rho U,\ y=y_0+\rho W\). The first exact equation is
\[
 2y_0W+\rho W^2
       =(3a^2-1)U+3a\rho U^2+\rho^2U^3.               \tag{A.58.8}
\]
Write \(N_a(a+t)=\sum_{j\ge1}n_jt^j\). The order
in (A.58.5), and \(k<5\), imply
\[
 n_j\in5R_0\ (1\le j<k),\qquad n_k\in R_0^\times.
                                                        \tag{A.58.9}
\]
Indeed \(N_a=g^5(H-H_0)\), and in characteristic five
the first nonconstant term of \(H\) has degree \(k\).
Thus \(N_a(a+\rho U)/\eta\) is integral. For \(j<k\)
its coefficient has valuation at least
\(1+(j-k)/(4k)>0\); for \(j>k\) it has valuation
at least \((j-k)/(4k)>0\). Its reduction is exactly
\(\overline n_k U^k\). Equation (A.58.8) eliminates
\(W\) linearly on the special fibre.

For \(y_0=0\), put \(\delta=3a^2-1\), a unit, and use
\(y=\rho U,\ x=a+\rho^2W\). The first equation now is
\[
           U^2=\delta W+3a\rho^2W^2+\rho^4W^3.        \tag{A.58.10}
\]
The first nonzero term of \(\overline N_a(a+t)\) has
degree \(k/2\), with nonzero coefficient. All lower
coefficients are divisible by \(5\). Substitution in
(A.58.7), with \(\rho^k=\eta\), is again integral;
its reduction is a nonzero multiple of \(W^{k/2}\).
Equation (A.58.10) eliminates
\(W=\overline\delta^{-1}U^2\).

In both cases the resulting special equation is therefore
\[
                         V-V^5=c_aU^k,
                         \qquad c_a\in k^\times.      \tag{A.58.11}
\]
The two equation partials in \(W,V\) are units on the
special fibre, after inverting \(g(x)\), whose reduction
there is a nonzero constant. These are smooth algebraic
charts over the enlarged valuation ring. Their generic
coordinate changes are invertible on the actual curve.
The special affine curve is geometrically integral: a
polynomial of degree \(k=2\) or \(4\) is not an
Artin--Schreier coboundary, as its pole order is not
divisible by five. Its smooth completion has one point
over infinity. To compute the ramification directly, put
\(t=U^{-1}\), so the extension valuation gives
\(v(t)=5\) and \(v(V)=-k\). The element
\(tV^{4/k}\) is a uniformizer. Under a nontrivial
translation \(V\mapsto V+b\) its difference has
valuation \(k+1\), since \(4/k\) is a unit modulo
five. The lower jump is therefore \(k\), the different
has degree \(4(k+1)\), and Hurwitz gives genus
\((5-1)(k-1)/2=2(k-1)\). Over the algebraic residue
field rescaling \(U\) gives the models in the statement.

\paragraph{Persistence and exhaustion of the proper genus.}
Compactify these algebraic charts and take regular
resolutions. Each smooth special generic point defines
a component with the displayed positive genus. The
factorization of birational maps of regular arithmetic
surfaces into point blowups preserves every positive-genus
component on passage to the stable model; see
\cite[Theorem 9.2.2 and Remark 9.2.6, pp.395--396,
Theorem 9.3.21, p.422]{fghLiu}, as used explicitly in A.52.
The fifteen tail valuations are distinct because their
restrictions to the good elliptic model have the fifteen
different centres (A.58.5). The main valuation lies over
its special generic point, so is distinct from all of them.

Their genera already sum to
\[
                         1+5\cdot6+10\cdot2=51=g(C).
                                                        \tag{A.58.12}
\]
The genus formula for a connected semistable fibre forces
every further component to be rational and the dual graph
to have first Betti number zero. Thus the reduction is of
compact type. Each Artin--Schreier tail has \(p\)-rank
zero by the one-branch-point Deuring--Shafarevich formula
\cite[Corollary 1.8, pp.35--36]{fghCrew1984};
the main elliptic component is ordinary. In particular the
good special Jacobian has five-rank one, although this
rank calculation is not needed for the fixed-group criterion.
No unramified-cover version of Raynaud's theorem is
applied to this ramified degree-five morphism.

\paragraph{The actual map and the cusp clusters.}
The proper map \(q\) extends from a regular semistable
model of \(C\) to the good elliptic model. Resolve its
indeterminacy, and then descend through the resulting
point blowups: their exceptional curves are rational,
so every map from one of them to the elliptic special
fibre is constant. The same applies to stable contraction;
compare \cite[Theorem 9.2.7, pp.396--397]{fghLiu} and
the detailed argument in A.52. The main component has
degree five by (A.58.2). Pulling back a degree-one line
bundle on the target gives total degree five on every
fibre, so all other components map constantly. Each
connected subtree off the main component attaches at one
point, and maps to the image of that point. The main
map is radicial, so a given reduced point of \(E\)
has a unique possible attachment point on that component.

On the ordinary elliptic curve the twenty-five geometric
generic points of \(E[5]\) specialize five at a time
to its five geometric special points, by the connected--etale
sequence for its ordinary five-divisible group
\cite[Sections 1.4 and 2.2]{fghTate1967Pdivisible}.
In these coordinates
their reduced support is precisely
\[
                    \{O\}\ \cup\ \{x^2=3\}.          \tag{A.58.13}
\]
The finite support follows also from
\(\overline\psi _5=(3x^2+1)^5\); its four finite
points together with \(O\) exhaust the ordinary
special five-torsion. This support is disjoint from
every centre in (A.58.5).

The map \(q\) is totally ramified over each generic
point of \(E[5]\), with one cusp in each fibre, by
the division-polynomial divisor of A.54. Its twenty-five
cusps consequently form five clusters of five marks over
the five points (A.58.13). None can lie on a positive-genus
tail or on its connecting subtree, since those subtrees
map to one of the different points (A.58.5).
Every cluster therefore lies over a single main-component
point with only rational subtrees above it.

Separate the marked sections on a semistable model. Since
each such cluster has at least two marks, its main-component
point must be replaced by a rational marked tree. No mark
can remain at its attachment point on the main component,
because that point is a node in the separated marked model.
All joining paths within that cluster consequently lie
on rational components, including their endpoints.
Stable marked contractions preserve this property.

For each chosen cusp select another in the same cluster.
The precise pointed hypothesis of A.41--A.42 now holds.
Those statements give the full integral fixed-inertia
theorem for every coefficient prime, every packet, and
every permitted open derived subgroup. The all-prime
normalizer theorem was already proved through the actual
quotient maps in A.54--A.55. This completes both assertions
for the specified Heisenberg cover, without a cofinality
or fixed-lifting assertion for other covers.
\end{proof}

% A.59: registered after full root and DE proof/source/certificate audits PASS.
\subsection{Normalizers for every cover with nilpotent Galois closure}

\begin{statement}{A.59}{Proved}{All-prime normalizer recognition for nilpotent geometric covers}
Let \(U=E\setminus\{O\}\), with \(E:y^2=x^3-x\), and let
\(Y\to U\) be a geometrically connected finite etale cover over
a number field. Suppose that its geometric Galois closure has
finite nilpotent group. Then, for every coefficient prime
\(\ell\), every nontrivial closed procyclic subgroup of
\(\pi_1(Y_{\overline K})^{(\ell)}\) whose arithmetic normalizer
has open image is contained in a conjugate of cusp inertia.
In particular this holds for every cover whose geometric Galois
closure is a group of prime-power order, including non-Galois covers.

More generally, for a fixed coefficient prime \(\ell\), it is
enough to have an actual factorization
\[
                           Y\longrightarrow V\longrightarrow U
                                                        \tag{A.59.1}
\]
in which the geometric Galois closure of the first map is an
\(\ell\)-group and that of the second has order prime to
\(\ell\). All maps and normalizers here are actual arithmetic
fundamental-group maps after finite descent. No full fixed-group
assertion or cofinality among arbitrary covers is claimed.
\end{statement}

\begin{proof}
We first prove the assertion for (A.59.1), then construct this
factorization for every coefficient prime in the nilpotent case.
Finite extensions of the number field are harmless: intersecting
an arithmetic normalizer with the inverse image of an open
Galois subgroup preserves openness of its arithmetic image,
while the geometric inertia subgroups do not change.
We may therefore arrange that the covers, their Galois closures
and all intermediate covers used below are defined over one
finite field extension.

\paragraph{The prime-to-coefficient intermediate curve has good reduction.}
The elliptic curve \(E\) has potentially good reduction at
every finite place, including the places above two: its
\(j\)-invariant is the integer \(1728\), and one may use
\cite[Proposition VII.5.5]{fghSilvermanAEC}. Once its good
model is chosen, the origin is a smooth section, so
\((E,O)\) is a good marked pair.

At a place of residue characteristic \(\ell\), apply the
prime-to-\(\ell\) specialization theorem to this good pair
and the geometric Galois closure of \(V\to U\). After a
finite extension the actual cover extends tamely over the pair;
see \cite[Expos\'e XIII, Corollary 2.9, Section 2.10 and
the proof of Corollary 2.12]{fghSGA1}, with the explicit
comparison already proved in A.56. Its proper source is smooth:
away from the origin the map is etale, and above the origin
its completed local equation after an etale base change is
\(R[[t]]\to R[[z]],\ t=z^e\), with \(e\) prime to
\(\ell\). The same statement holds for subgroup quotients.
Consequently the proper compactification \(\overline V\)
has potentially good \emph{curve} reduction at this place.

A nonconstant map \(\overline V\to E\) forces its genus
to be at least one. If that genus is at least two, A.57
proves its pro-\(\ell\) normalizer criterion. If it is
one, A.40 gives the criterion for the actual isogeny
cover with its removed origin fibre. Thus the normalizer
criterion holds for \(V\) in every case, including
\(\ell=2\). This step needs no restriction on the
degree of \(V\); only its prime-to-\(\ell\) Galois
closure is used.

\paragraph{The first map gives an actual open pro-ell embedding.}
Put \(\Pi=\pi_1(V_{\overline K})\), and let
\(H\subset\Pi\) be the open subgroup corresponding to
\(Y\). Its normal core \(C\) has \(\Pi/C\) an
\(\ell\)-group. The pro-\(\ell\) topology of \(H\)
is induced from \(\Pi\). For completeness, if
\(N\triangleleft H\) is open with \(H/N\) an
\(\ell\)-group, then \(N\cap C\) is normal in
\(C\). Intersect its finitely many conjugates under
\(\Pi/C\), obtaining \(M\triangleleft\Pi\) with
\(M\subset N\). The quotient \(C/M\) embeds in the
product of the \(\ell\)-groups
\(C/g(N\cap C)g^{-1}\). Hence \(\Pi/M\) is an
\(\ell\)-group. This proves the required induced topology
and the continuous open injection
\[
               H^{(\ell)}\lhook\joinrel\longrightarrow
                                  \Pi^{(\ell)}.       \tag{A.59.2}
\]
This is the general completion argument in A.22; it uses
an \(\ell\)-group Galois closure, not merely an
\(\ell\)-power degree.

The actual arithmetic map induced by \(Y\to V\)
has (A.59.2) on its geometric kernels and the identity
on the common field's absolute Galois quotient. It is
therefore itself injective, with open image. In particular
the arithmetic normalizer of a subgroup of \(H^{(\ell)}\)
maps into its arithmetic normalizer in the geometrically
pro-\(\ell\) group of \(V\), with the same open
arithmetic image.

Let \(A\subset H^{(\ell)}\) be a nontrivial closed
procyclic subgroup with that property. By the criterion
for \(V\), its image is contained in some geometric
inertia conjugate \(J\). The intersection
\[
                         H^{(\ell)}\cap J             \tag{A.59.3}
\]
is exactly the full pro-\(\ell\) inertia of a cusp
of the actual curve \(Y\). Indeed the usual geometric
covering correspondence describes the branches above a
cusp by double cosets, and the inertia of the branch
is the corresponding subgroup intersection. Equivalently,
the local map \(t=z^e\) identifies it with
\(e\Z_\ell(1)\subset\Z_\ell(1)\).
Here \(e\) is an \(\ell\)-power, since the Galois
closure of \(Y\to V\) is an \(\ell\)-group. The
induced-topology comparison makes this exactly the
inertia in (A.59.2), with no additional quotient kernel.
Thus \(A\) is cuspidal on \(Y\), proving the general
factorization assertion.

\paragraph{Every nilpotent Galois closure has the required factorization.}
Write \(\pi=\pi_1(U_{\overline K})\), let \(H\subset\pi\)
correspond to the original cover, let \(N\) be its normal
core, and put \(G=\pi/N\), \(A=H/N\). For the chosen
coefficient prime write the finite nilpotent group as
\[
                   G=S\times T,
 \qquad S\text{ an }\ell\text{-group},\quad
              |T|\text{ prime to }\ell.               \tag{A.59.4}
\]
Every subgroup splits accordingly:
\[
             A=A_S\times A_T,
             \quad A_S=A\cap S,\quad A_T=A\cap T.
                                                        \tag{A.59.5}
\]
To see this without a normality assumption on \(A\),
take \((s,t)\in A\) and choose an integer congruent
to one modulo \(|S|\) and zero modulo \(|T|\).
That power of \((s,t)\) is \((s,1)\in A\); the
complementary power gives \((1,t)\in A\). This
proves (A.59.5) for arbitrary subgroups.

Let \(B=S\times A_T\), and let \(V\to U\) be
the actual intermediate cover corresponding to the
inverse image of \(B\). Its geometric Galois closure
has group
\[
 \begin{aligned}
 G/\operatorname{core}_G(B)
    &\simeq T/\operatorname{core}_T(A_T),\\
 B/\operatorname{core}_B(A)
    &\simeq S/\operatorname{core}_S(A_S).
 \end{aligned}                                          \tag{A.59.6}
\]
The first group has order prime to \(\ell\); the
second, which is the Galois closure group of
\(Y\to V\), is an \(\ell\)-group.
The formulas follow directly by taking conjugates
in the direct products; in the second formula
\(A_T\) is normal in the \(T\)-factor \(A_T\)
of \(B\). Thus (A.59.1) has exactly the required
two Galois-closure hypotheses, and the preceding
argument applies.

This proves the all-prime assertion simultaneously
for every cover with finite nilpotent geometric
Galois closure. In particular it gives an independent
normalizer proof for all prime-power Heisenberg
covers previously considered. Their full fixed
subgroups still require separate arguments such
as A.22 and A.41--A.42. The finite nilpotent
quotients are not asserted to be cofinal among
all finite geometric quotients.
The exact certificate \texttt{code/a\_nilpotent\_cover\_normalizers.py}
checks every subgroup and both core formulas in three mixed-prime
examples, including nonnormal subgroups. The general proof above
is independent of that finite sample.
\end{proof}

% A.60: registered after full independent root and DE proof/source/certificate audits PASS.
\subsection{The next Heisenberg cover: an exact seven-adic model}

\begin{statement}{A.60}{Proved}{Both all-prime assertions for the genus-one-hundred-forty-eight Heisenberg cover}
Let \(E:y^2=x^3-x\), and let \(\psi _7\) be its seventh
division polynomial. Let \(C_7\) denote the smooth proper
normalization of the displayed function field. The actual
curve and punctured map
\[
 C_7:\ y^2=x^3-x,\ z^7=\psi _7(x),\qquad
 Y_7=C_7\setminus q^{-1}(E[7])
              \xrightarrow{[7]q}E\setminus\{O\}       \tag{A.60.1}
\]
give a geometrically connected Galois cover with Heisenberg
group of order \(343\), proper genus \(148\), and \(49\)
cusps. For every coefficient prime, its nontrivial closed
procyclic subgroups with open arithmetic normalizer are
cuspidal. For every selected cusp packet and every indicated
open derived subgroup, its full geometric pro-prime fixed
subgroup is exactly the selected cusp inertia.

The fixed-group assertion comes from an explicit seven-adic
compact-type model. Its positive-genus components are one
supersingular elliptic main component and seven copies of
\[
                              V^7-V=U^8,               \tag{A.60.2}
\]
each of genus twenty-one. All forty-nine cusps lie on a
rational marked tree attached to the main component, away
from the seven positive-genus tails.
\end{statement}

\begin{proof}
We first specify the actual cover and its group, so the
model calculation concerns an identified instance of Task A.5.
Use the exact polynomials
\[
\begin{aligned}
 \psi _3&=3x^4-6x^2-1,\\
 B&=x^6-5x^4-5x^2+1,\\
 \psi _7&=\psi _5\psi _3^3-128(x^3-x)^2B^3,
 \\ h&=\frac{\psi _7+1}{7},
\end{aligned} \tag{A.60.3}
\]
where \(\psi _5\) is the polynomial of A.54. These are the
standard division polynomials of \cite[Exercise III.3.7,
pp.105--106]{fghSilvermanAEC}. The middle
identity is the division-polynomial recurrence
\(\psi _7=\psi _5\psi _3^3-\psi _2\psi _4^3\), with
\(\psi _2=2y\) and \(\psi _4=4yB\).
It gives \(h\in\Z[x]\), and
\[
 \begin{aligned}
 \overline h={}&x^{24}-2x^{22}-2x^{20}+x^{18}-x^{14}\\
               &-3x^{10}-x^8-x^6-3x^4
                   \quad\text{in }\F_7[x].
 \end{aligned}
                                                        \tag{A.60.4}
\]

\paragraph{The actual geometric Galois group.}
The divisor is
\(\operatorname{div}(\psi _7)=D_{E[7]}-49(O)\),
where \(D_{E[7]}\) includes \(O\); see
\cite[Exercise III.3.7(f), p.106]{fghSilvermanAEC}. Its order at each
nonzero seven-torsion point is one, and its order at
\(O\) is \(-48\), also prime to seven. Thus the
Kummer cover \(q\) is connected of degree seven and
totally ramified precisely at the forty-nine points of
\(E[7]\). Hurwitz gives
\[
                2g(C_7)-2=49(7-1)=294,
                           \qquad g(C_7)=148.          \tag{A.60.5}
\]
The characteristic-zero isogeny \([7]\) is etale, so
(A.60.1) has degree \(343\), has forty-nine cusps and
is etale away from them.

Every translation by \(T\in E[7]\) lifts to \(C_7\)
after finite extension. Indeed the divisor of
\(\psi _7(P+T)/\psi _7(P)\) is
\(49(O)-49(-T)\), seven times a principal divisor
\(7(O)-7(-T)\). Choose a function \(r_T\) for that
divisor and a seventh root of the remaining constant.
Then \((P,z)\mapsto(P+T,a_T r_T(P)z)\) is an actual
lift. These give all \(343\) automorphisms of the
degree-\(343\) composite, with exact sequence
\[
                 1\longrightarrow\mu _7
                 \longrightarrow G\longrightarrow E[7]
                 \longrightarrow1.                   \tag{A.60.6}
\]
The extension is central. It is nonabelian: the original
once-punctured elliptic inertia is a commutator, whose
image in this cover has order seven, whereas it would
vanish in an abelian quotient.

For precision, the group has exponent seven, not
forty-nine. The evenness of \(\psi _7\) gives the actual
automorphism \(\rho(x,y,z)=(-x,iy,z)\), which acts by
\([i]\) on \(E[7]\) and trivially on \(\mu _7\).
In a central extension (A.60.6), seventh powers define
an additive map \(E[7]\to\mu _7\): the commutator
correction has exponent \(\binom72\), divisible by
seven. This map is invariant under \([i]\). Since
\([i]-1\) is invertible on \(E[7]\), it vanishes.
Choose a basis of \(E[7]\) with nonzero commutator
pairing and lift it; together with the central generator
its relations are precisely those of the Heisenberg
group of order \(343\). This also puts the normalizer
assertion within A.59. All these maps and groups descend
after a finite extension, which is sufficient for the
arithmetic normalizer comparison.

\paragraph{The radicial main component and its seven critical centres.}
Choose \(\pi^7=7\) and substitute \(z=-1+\pi u\).
The exact equation is
\[
 \begin{gathered}
 u^7-\pi^6u^6+3\pi^5u^5-5\pi^4u^4\\
             {}+5\pi^3u^3-3\pi^2u^2+\pi u=h(x),
                \qquad y^2=x^3-x.
 \end{gathered}                     \tag{A.60.7}
\]
Modulo seven it has function field \(k(E)(h^{1/7})\).
The exact differential factorization in characteristic seven is
\[
                   dh=3x^3(x^2-1)^3(x^2-3)^7\,dx\ne0.
                                                        \tag{A.60.8}
\]
Thus this field is \(k(E)^{1/7}\). Its normalized curve
is an elliptic Frobenius twist, and (A.60.7) is smooth
at its special generic point, giving an actual
genus-one main component mapping radicially with degree
seven to \(E\). The elliptic special fibre is
supersingular: the coefficient of \(x^6\) in
\((x^3-x)^3\) is zero.

The zeros of (A.60.8) on \(E\) are exactly
\[
                (0,0),(1,0),(-1,0),
                            \qquad x^2=3,              \tag{A.60.9}
\]
the latter condition giving four points. Every zero
has order seven. At a two-torsion point use \(y\)
as parameter: \(x-a\) has order two and \(dx\)
order one. At each of the four other points, \(x-a\)
is a parameter. These are seven distinct finite points;
none is the origin.

\paragraph{The degree-seven coefficient which must also vanish.}
The differential alone would not exclude a seventh-power
term below the desired degree-eight term. We verify that
term explicitly. Choose the centres (A.60.9) in a finite
\emph{unramified} extension \(R_0/\Z_7\), taking
\(a^2=3\) and a unit root \(y_0^2=a^3-a\) for the
four non-two-torsion points. Formula (A.60.4) gives
\[
 \begin{aligned}
  \overline h(a+t)-\overline h(a)&\equiv2t^8\pmod {t^9},
                                      &&a^2=3,\\
  \overline h(a+t)-\overline h(a)&\equiv b_at^4\pmod {t^5},
                                      &&a=0,1,-1,\\
                         b_0&=4,\qquad b_1=b_{-1}=2.
 \end{aligned}                                          \tag{A.60.10}
\]
In particular every coefficient of degree one through
seven in the first line vanishes, including degree
seven. These are identities in
\(\F_7[a,t]/(a^2-3)\), not only checks at one residue
embedding. The certificate expands the entire
polynomial and verifies these identities exactly.
All corresponding lower coefficients over \(R_0\)
are therefore divisible by \(7\).

\paragraph{Exact smooth algebraic tail charts.}
Choose \(\eta^6=-7\), \(\rho^8=\eta\), and
\(z_0^7=\psi _7(a)\). The latter is a unit congruent
to \(-1\). Use \(z=z_0(1+\eta V)\). The exact
covering equation, after division by \(7\eta\), is
\[
 \begin{aligned}
 \mathcal B_\eta(V)&=V+3\eta V^2+5\eta^2V^3+5\eta^3V^4\\
                  &\qquad+3\eta^4V^5+\eta^5V^6-V^7,\\
 z_0^7\mathcal B_\eta(V)&=\frac{h(x)-h(a)}{\eta}.
 \end{aligned}                                            \tag{A.60.11}
\]
For \(y_0\ne0\), set \(x=a+\rho U\),
\(y=y_0+\rho W\). The elliptic equation is
\[
 2y_0W+\rho W^2
              =(3a^2-1)U+3a\rho U^2+\rho^2U^3.       \tag{A.60.12}
\]
By (A.60.10), the right side of (A.60.11) is
integral and reduces to \(2U^8\). A lower coefficient
of degree \(j<8\) has valuation at least
\(1+(j-8)/48>0\); a higher coefficient has valuation
at least \((j-8)/48>0\). The elliptic equation
eliminates \(W\) by its unit linear coefficient.

At \(a=0,1,-1\), instead set \(y=\rho U\),
\(x=a+\rho^2W\). Put \(\delta=3a^2-1\), a unit.
The first equation becomes
\[
                U^2=\delta W+3a\rho^2W^2+\rho^4W^3.
                                                        \tag{A.60.13}
\]
The second line of (A.60.10) makes the right side
of (A.60.11) integral, with reduction
\(b_aW^4=b_a\overline\delta^{-4}U^8\).
Again all other coefficients have strictly positive
valuation at the indicated scale.

The special equations in both cases are therefore
\(V^7-V=c_aU^8\), with \(c_a\ne0\). The two
equation partials in \(W,V\) are units on the
special fibre, so these are actual smooth algebraic
charts with integral special curve. Their generic
coordinate changes are invertible on \(C_7\).
The Artin--Schreier polynomial is irreducible because
its pole order eight is prime to seven. Over the
algebraic residue field rescaling \(U\) gives
(A.60.2).

For its ramification calculation put \(t=U^{-1}\).
At the unique point above infinity, \(v(t)=7\),
\(v(V)=-8\), and \(t^7V^6\) is a uniformizer.
Under a nontrivial translation \(V\mapsto V+b\),
its difference has valuation \(49-40=9\). The
lower jump is eight and the different is \(6\cdot9\).
Hurwitz gives genus twenty-one. The one-branch-point
Deuring--Shafarevich formula gives seven-rank zero,
as in \cite[Corollary 1.8]{fghCrew1984}.

\paragraph{Stable components and all forty-nine marked sections.}
The actual positive-genus valuations furnished by
these seven smooth charts persist in the stable
model, by the regular-surface factorization and
positive-genus persistence argument of A.52 and
A.58; the precise sources are
\cite[Theorem 9.2.2, Remark 9.2.6 and Theorem 9.3.21]
{fghLiu}. Their restrictions to the good elliptic
model have the distinct centres (A.60.9), while
the main valuation lies above its generic point.
Their genera already exhaust the generic genus:
\[
                              1+7\cdot21=148.          \tag{A.60.14}
\]
The genus formula forces every further component
to be rational and the dual graph to be a tree.
All component Jacobians have seven-rank zero,
including the supersingular elliptic main, so
the good special Jacobian has seven-rank zero.

The actual degree-seven map \(q\) extends to the
good elliptic model by the same resolution-and-descent
argument as A.58. It has degree seven on the
radicial main component. Constancy of the degree
of a pulled-back line bundle makes every other
component map constantly. Every off-main subtree
has one attachment point and maps to its image.

All forty-nine points of the generic \(E[7]\)
specialize to \(O\), since the elliptic special
fibre is supersingular. Their ramified \(q\)-fibres
are the forty-nine distinct cusps. None can
specialize on one of the seven positive-genus
tails or its connecting subtree, since all
those subtrees map to the finite points (A.60.9).
The radicial main has a unique point over \(O\).
After separating the cusp sections, all forty-nine
marks therefore lie on the rational marked tree
attached at that point. No mark can remain at
the main attachment node. Thus every selected
cusp has a rational joining path to another,
with rational endpoints.

The all-prime full fixed-inertia conclusion is
exactly A.41--A.42. The actual cover has a
seven-group geometric Galois closure, so its
all-prime normalizer conclusion follows from
A.59. Alternatively away from seven its proper
curve has potential good reduction by the tame
argument and A.57 applies, while at seven the
explicit zero-rank Jacobian gives A.46. No
characteristic-zero Jacobian decomposition or
unverified arithmetic character is required. The exact
certificate \texttt{code/a\_heisenberg\_seven.py} verifies
(A.60.3)--(A.60.4), (A.60.7)--(A.60.13),
the missing degree-seven coefficient, all scaling valuations,
and the genus and cusp counts; the geometric conclusions
use the full model argument above.
\end{proof}

% A.61: root, DE and PBC full independent proof/source/certificate audits PASS.
\subsection{An unbounded Heisenberg family from a divided differential operator}

\begin{statement}{A.61}{Proved}{Both assertions for every inert-prime Heisenberg cover}
For every prime \(p\equiv3\pmod4\), let \(C_p\) be the
smooth proper normalization of the function field
\[
 E:y^2=x^3-x,\qquad z^p=\psi_p(x),
 \qquad q:C_p\longrightarrow E,
                                                        \tag{A.61.1}
\]
where \(\psi_p\) is the usual division polynomial. Then
\[
 Y_p=C_p\setminus q^{-1}(E[p])
           \xrightarrow{[p]q}U
                                                        \tag{A.61.2}
\]
is an actual degree-\(p^3\) Heisenberg cover, with
\(p^2\) cusps and proper genus
\(1+p^2(p-1)/2\). It satisfies both assertions of
Task A.5 at every coefficient prime: all nontrivial
arithmetically normalized pro-\(\ell\) cyclic subgroups are
cuspidal, and every permitted cusp packet and indicated
open derived subgroup has full fixed group exactly
the selected cusp inertia.

At a place above \(p\), an explicit stable model has
one supersingular elliptic main component and \(p\)
positive-genus tails, each geometrically isomorphic to
\[
                         V^p-V=U^{p+1},
 \qquad g=\frac{p(p-1)}2.                              \tag{A.61.3}
\]
Every other component is rational, the dual graph is
a tree, and all \(p^2\) cusps lie on a rational marked
tree over the main point above \(O\). The proof is
uniform in \(p\); no finite polynomial pattern is
being substituted for it.
\end{statement}

\begin{proof}
We establish a uniform local coefficient identity before
constructing the model. Throughout its proof all reduced
functions and constants are defined over \(\F_p\),
even when their zeros are subsequently separated over
a finite extension.

\paragraph{Frobenius, and the integral divided polynomial.}
Write \(f=x^3-x\). Pairing \(x\) with \(-x\) in the
quadratic-character point count gives
\(\#E(\F_p)=p+1\), since \(-1\) is a nonsquare.
The Frobenius endomorphism \(F\) therefore satisfies
\[
                           F^2=[-p].                  \tag{A.61.4}
\]
This is the actual endomorphism identity of
\cite[Theorem V.2.3.1(b), p.142]{fghSilvermanAEC}, not
only a characteristic-polynomial statement on one
Tate module. It also proves supersingularity: the
multiplication-by-\(p\) morphism is purely inseparable
of degree \(p^2\), so its geometric kernel is \(\{O\}\).

The exact division formulas are
\[
 x\circ[p]=\frac{\phi_p}{\psi_p^2},\qquad
 \phi_p=x\psi_p^2-\psi_{p+1}\psi_{p-1}.
                                                        \tag{A.61.5}
\]
The polynomial \(\phi_p\) is monic of degree \(p^2\);
the numerator and denominator are relatively prime on
a nonsingular Weierstrass curve. These assertions and
the division recurrence are
\cite[Exercise III.3.7(a)--(d), pp.105--106]{fghSilvermanAEC}.
Reducing (A.61.5) and using (A.61.4) gives
\(\overline\phi_p=x^{p^2}\overline\psi_p^{\,2}\).
The denominator is nonzero: otherwise relative primality
would fail for the nonconstant monic numerator. Comparing
degrees and leading coefficients shows
\(\overline\psi_p=c\), a nonzero constant with \(c^2=1\).

In fact \(\psi_n(0)=(-1)^{(n-1)/2}\) for odd \(n\).
This follows directly by induction from the odd division
recurrence at the two-torsion point \((0,0)\): the even
division functions vanish there; for an odd middle index
only the first product remains, and for an even middle
index only the negative second product remains. The
initial values are \(\psi_1(0)=1\), \(\psi_3(0)=-1\).
Thus in the present case \(c=-1\), and
\[
             h=\frac{\psi_p+1}{p}\in\Z[x],
       \qquad \deg h=\frac{p^2-1}{2},
       \qquad \text{leading coefficient }1.           \tag{A.61.6}
\]
We retain the notation \(c=-1\) in the identities below
to make their origins visible.

\paragraph{The invariant derivation and its divided power.}
Let
\[
 \mathcal A=\Z_{(p)}[x,y]/(y^2-x^3+x),\qquad
 \omega=\frac{dx}{2y},\qquad
 D(x)=2y,\quad D(y)=3x^2-1.                            \tag{A.61.7}
\]
On the smooth elliptic curve the derivation \(D\) is
the nowhere-vanishing invariant vector field dual to
\(\omega\). Its reduction satisfies \(D^p=0\).
Here is a direct verification for this curve. In
characteristic \(p\), the \(p\)-th power of a derivation
is a derivation; it is still translation-invariant,
so it is a scalar multiple of \(D\). Over \(\F_{p^2}\),
the automorphism \(\rho(x,y)=(-x,iy)\) conjugates
\(D\) to \(-iD\), and therefore \(D^p\) to
\((-i)^pD^p=iD^p\). A scalar multiple of \(D\)
instead transforms by \(-i\). Since \(p\) is odd,
these are different, proving \(D^p=0\).

Consequently \(D^p\mathcal A\subset p\mathcal A\),
and the operation
\[
                \mathcal E(\overline b)
                  =\overline{D^p b/p}
                  \quad(b\in\mathcal A)               \tag{A.61.8}
\]
is a well-defined \(\F_p\)-linear operator on
\(\overline{\mathcal A}\). Changing the lift by
\(p b_1\) changes the right side by
\(\overline{D^p b_1}=0\). It commutes with \(D\).
It is not being treated as a derivation. Its essential
identity is
\[
                    \mathcal E(b^p)=-(Db)^p.           \tag{A.61.9}
\]
To prove it, expand \(D^p\) on a product of \(p\)
identical factors. The distribution with one derivative
on every factor contributes \(p!(Db)^p\), which after
division by \(p\) reduces to \(-(Db)^p\). The
distributions with all \(p\) derivatives on one factor
sum to \(p b^{p-1}D^p b\), and vanish after division
and reduction. Every other distribution has all entries
less than \(p\), so its multinomial coefficient is
divisible by \(p\). Its orbit under cyclic permutation
of the \(p\) identical factors has length \(p\);
the additional orbit factor makes it divisible by
\(p^2\). These are all cases, proving (A.61.9).

If a reduced function \(R\) satisfies \(DR=a\in\F_p\),
the same expansion on \(p+1\) factors gives
\[
 \mathcal E(R^{p+1})
       =R^p\mathcal E(R)-R a^p.                       \tag{A.61.10}
\]
Indeed the single \(p\)-derivative terms give the first
term. The distributions with \(p\) first derivatives
give the second by Wilson's congruence. Every other
term contains some \(D^jR\), \(2\le j<p\), which
vanishes modulo \(p\), in addition to the factor
\(p\) already present in the multinomial coefficient.
This proves (A.61.10) with arbitrary integral lifts.

\paragraph{The division-polynomial differential identity.}
In characteristic zero one has the exact rational-function
identity
\[
 \psi_p D^2\psi_p-(D\psi_p)^2
       =p^2\psi_{p+1}\psi_{p-1}.                      \tag{A.61.11}
\]
For completeness, equivalently
\(D^2\log\psi_p=p^2(x-x\circ[p])\).
In a formal coordinate \(s\) with \(ds=\omega\),
the principal part on either side at every nonzero
\(p\)-torsion point is \(-s^{-2}\); at \(O\) it is
\((p^2-1)s^{-2}\). There are no other poles. Their
difference is therefore a constant. At \(O\),
\(x=s^{-2}+O(s^2)\). Moreover \(\psi_p\) is an
even polynomial in \(x\): pullback by \([i]\)
preserves its divisor and its leading coefficient,
since \((p^2-1)/2\) is even. Thus its next-to-leading
coefficient vanishes, and
\(\psi_p=p s^{1-p^2}(1+O(s^4))\).
Both constant terms in the displayed logarithmic
identity are zero. This proves it, and (A.61.5)
gives (A.61.11). The formal Laurent assertion for
\(x\) also follows immediately from
\((Dx)^2=4(x^3-x)\); no transcendental comparison
is required.

Substitute \(\psi_p=-1+ph\) in (A.61.11).
Reduction first proves \(D^2\overline h=0\).
Hence
\[
                 D\overline h=R^p
               \quad\text{for a unique }R\in\F_p(E).
                                                        \tag{A.61.12}
\]
The kernel of a nonzero derivation on the one-variable
function field \(\F_p(E)\) is its subfield of
\(p\)-th powers. The function on the left is regular
away from \(O\), so normality makes \(R\) regular
there too. Its pole has order exactly \(p\): the
leading coefficient and degree in (A.61.6) give
\(D\overline h\) a pole of order \(p^2\).
The function \(R\) is odd under \([-1]\), since
\(\overline h\) is a polynomial in \(x\).

Put \(T=D^2h/p\in\mathcal A\). Its integrality
also follows from (A.61.11), since the reduction of
\(\psi_p\) is a unit. Reducing the full identity and
using (A.61.4)--(A.61.5) yields
\[
 c\overline T-R^{2p}=c^2(x-x^{p^2}),
                         \qquad c=-1.                 \tag{A.61.13}
\]
Apply \(D^{p-2}\), which is a positive power even
when \(p=3\). The \(p\)-th-power terms disappear,
and (A.61.8) gives
\[
                   \mathcal E(\overline h)
                           =cD^{p-2}x.                \tag{A.61.14}
\]
Apply \(D\) and use (A.61.9) and commutation:
\[
                      -(DR)^p=cD^{p-1}x.               \tag{A.61.15}
\]
The function \(D^{p-1}x\) is killed by \(D\) and
has pole order at most \(p+1<2p\). Its \(p\)-th
root can have pole order at most one, so it is
constant on the genus-one curve. Thus
\[
                         DR=a\in\F_p^\times.          \tag{A.61.16}
\]
The constant is nonzero: otherwise \(R\) would be
a \(p\)-th power whose root has a single simple
pole, which is impossible on an elliptic curve.
Because \(a\in\F_p\), equation (A.61.15) says
\(cD^{p-1}x=-a\). Therefore
\(cD^{p-2}x+R\) is killed by \(D\) and has pole
order at most \(p\). Its \(p\)-th root is again
constant. Both summands are odd under \([-1]\),
so that constant is zero. We have proved
\[
               R=-cD^{p-2}x,
                \qquad \mathcal E(\overline h)=-R.    \tag{A.61.17}
\]
This identifies the actual divided polynomial's
derivative; \(D^2\overline h=0\) alone would not
have sufficed.

\paragraph{The hidden degree-p Taylor coefficient vanishes.}
By (A.61.12) and (A.61.16), there is a unique
\(S\in\F_p(E)\) such that
\[
            \overline h=\frac{R^{p+1}}a+S^p.           \tag{A.61.18}
\]
Its only possible pole is at \(O\), of order at
most \(p+1\). Since \(\mathcal E\) commutes with
\(D\) and \(DR=a\), the function
\(\mathcal E(R)\) is killed by \(D\). Its pole
order is at most \(2p\): take an integral lift of
\(R\) of pole order at most \(p\), and apply
\(D^p/p\). Consequently
\[
           \mathcal E(R)=A^p,
                      \qquad A\in L(2O)(\F_p).        \tag{A.61.19}
\]
Apply \(\mathcal E\) to (A.61.18), using
(A.61.9)--(A.61.10), (A.61.17), and
\(a^{p-1}=1\). It gives
\[
 -R=\left(\frac{RA}{a}\right)^p-R-(DS)^p,
                  \qquad DS=\frac{RA}{a}.             \tag{A.61.20}
\]
The last equality follows from injectivity of
Frobenius. Every constant and every function in
(A.61.12)--(A.61.20) is defined over \(\F_p\);
in particular the use of \(a^{p-1}=1\) is justified.

The differential \(dR=a\omega\) never vanishes.
Since \(R\) has one pole of order \(p\), it has
exactly \(p\) distinct simple zeros, all away from
\(O\). At any such point \(P\), equation
(A.61.20) says \(dS(P)=0\). For a local parameter
\(t\), equation (A.61.18) now gives
\[
 \begin{aligned}
 \overline h(P+t)-\overline h(P)
       &=\frac{R(P+t)^{p+1}}a+(S(P+t)-S(P))^p,\\
 \operatorname{ord}_P(\overline h-\overline h(P))&=p+1.
 \end{aligned}
                                                        \tag{A.61.21}
\]
Indeed the first term has exactly that order and
the second has order at least \(2p>p+1\).
Thus every Taylor coefficient of degrees one
through \(p\), including the degree invisible
to ordinary differentiation, is zero.

More explicitly \(R=yQ(x)\), with
\(\deg Q=(p-3)/2\). Equation (A.61.16) is
\[
       (3x^2-1)Q+2(x^3-x)Q'=a.                         \tag{A.61.22}
\]
It makes \(Q\) squarefree and disjoint from
\(x(x^2-1)\). The zeros of \(R\) consist of
the three finite two-torsion points and
\(p-3\) other points, two over each root of
\(Q\). There are no omitted critical centres.

\paragraph{Actual algebraic main and tail charts.}
Choose \(\pi^p=p\) and set \(z=-1+\pi u\).
The exact expansion of (A.61.1), after division
by \(p\), has reduction \(u^p=\overline h\).
Its differential is nonzero by (A.61.12), so its
normalized special function field is
\(\F_p(E)^{1/p}\). It gives a genus-one main
component mapping radicially with degree \(p\)
to the good elliptic special fibre. Its special
generic point is smooth, using the nonzero
differential and the elliptic equation.

Choose the \(p\) residue centres just found over
a finite \emph{unramified} extension \(R_0/\Z_p\).
For a non-two-torsion centre, lift its \(x\)-coordinate
as a simple root of an integral lift of \(Q\), then
lift \(y\) by the unit-derivative elliptic equation.
At a two-torsion point take \((a,0)\) with
\(a=0,1,-1\). Equation (A.61.21) therefore implies
divisibility of every lower integral Taylor
coefficient by \(p\), not merely by a ramified
uniformizer.

Take \(\eta^{p-1}=-p\), \(\rho^{p+1}=\eta\),
\(z_0^p=\psi_p(a)\), and put
\(z=z_0(1+\eta V)\). The exact equation is
\[
 \begin{aligned}
 \mathcal B_\eta(V)
    &=\sum_{j=1}^{p-1}\frac{\binom pj}{p}
                       \eta^{j-1}V^j-V^p,\\
 z_0^p\mathcal B_\eta(V)
    &=\frac{h(x)-h(a)}{\eta}.
 \end{aligned} \tag{A.61.23}
\]
At a centre \((a,y_0)\) with \(y_0\ne0\), set
\(x=a+\rho U\), \(y=y_0+\rho W\). The elliptic
equation becomes
\[
 2y_0W+\rho W^2
             =(3a^2-1)U+3a\rho U^2+\rho^2U^3.         \tag{A.61.24}
\]
The leading term on the right of (A.61.23) is
a nonzero multiple of \(U^{p+1}\).
For a lower degree \(1\le j<p+1\), its valuation
after division by \(\eta\) is at least
\[
                   1+\frac{j-(p+1)}{p^2-1}>0;
                                                        \tag{A.61.25}
\]
for a higher degree it is at least
\((j-(p+1))/(p^2-1)>0\).

At a two-torsion centre use instead
\(y=\rho U\), \(x=a+\rho^2W\), with
\(\delta=3a^2-1\) a unit. Then
\[
                 U^2=\delta W+3a\rho^2W^2+\rho^4W^3.
                                                        \tag{A.61.26}
\]
The first nonzero degree of
\(\overline h(a+X)-\overline h(a)\) is
\((p+1)/2\), since \(y\) is the local parameter.
The same valuation calculation, replacing \(j\)
by \(2j\), leaves a nonzero multiple of
\(U^{p+1}\). All other terms have strictly
positive valuation.

The special equations are therefore
\(V^p-V=c_PU^{p+1}\), \(c_P\ne0\).
The partials in \(W,V\) are units. These are
actual smooth algebraic charts with integral
special curves, and their generic coordinate
changes are invertible on (A.61.1). The
Artin--Schreier polynomial is irreducible because
its pole order \(p+1\) is prime to \(p\).
Over the algebraic residue field rescaling
\(U\) gives (A.61.3).

At its unique point over infinity put \(t=U^{-1}\).
One has \(v(t)=p\), \(v(V)=-(p+1)\), and
\(t^pV^{p-1}\) is a uniformizer. Its difference
under \(V\mapsto V+b\), \(b\ne0\), has
valuation \(p+2\). The lower jump is \(p+1\),
the different is \((p-1)(p+2)\), and Hurwitz
gives genus \(p(p-1)/2\). Its \(p\)-rank is
zero by the one-branch-point Deuring--Shafarevich
formula, \cite[Corollary 1.8]{fghCrew1984}.

\paragraph{The actual cover, all components, and every cusp.}
The division divisor is
\(D_{E[p]}-p^2(O)\). Thus \(q\) is connected
and totally ramified at every point of \(E[p]\),
with orders one and \(1-p^2\), both prime to
\(p\). It has genus \(1+p^2(p-1)/2\), and
(A.61.2) has degree \(p^3\) and \(p^2\) cusps.
Every translation by \(T\in E[p]\) lifts:
the divisor of its division-function ratio is
\(p^2(O)-p^2(-T)\), which is \(p\) times
the principal divisor \(p(O)-p(-T)\).
Adjoining roots of the finitely many remaining
constants gives the actual lifts over a number
field. The deck group is a central extension
of \(E[p]\) by \(\mu_p\). It is nonabelian
because peripheral commutator inertia has
order \(p\). The automorphism
\((x,y,z)\mapsto(-x,iy,z)\) fixes \(\mu_p\)
and acts by \([i]\) on \(E[p]\).
The \(p\)-th-power map \(E[p]\to\mu_p\)
is additive, since \(p\mid\binom p2\), and
is \([i]\)-invariant. As \([i]-1\) is
invertible on \(E[p]\), this map is zero.
The group is therefore the exponent-\(p\)
Heisenberg group of order \(p^3\), exactly as
in the individual verification of A.60.

The positive-genus valuations furnished by
the main chart and the \(p\) tail charts
persist in the stable model. They are distinct:
the tail restrictions have the \(p\) different
finite elliptic centres, and the main restriction
is the generic point. The precise common-regular-
resolution and positive-genus persistence argument
is that of A.52, A.58 and A.60, using
\cite[Theorem 9.2.2, Remark 9.2.6 and Theorem 9.3.21]
{fghLiu}. Their genera already sum to
\[
                     1+p\frac{p(p-1)}2
                        =1+\frac{p^2(p-1)}2.          \tag{A.61.27}
\]
The semistable genus formula forces compact
type and makes every remaining component
rational. No unramified-cover theorem is
being applied to the ramified map \(q\).

The actual map to the good elliptic model
extends by resolving its indeterminacy and
contracting rational exceptional components,
as in A.58. Its main component already has
degree \(p\). Constancy of the degree of a
pulled-back line bundle forces every other
component to map to a point. Every off-main
subtree has one attachment point and maps
constantly to that image.

All generic \(p\)-torsion points specialize
to \(O\), by (A.61.4). Each has a unique
ramified cusp above it. The positive-genus
tails lie above zeros of \(R\), all different
from \(O\), and hence no cusp lies on those
tails or their connecting subtrees. The
radicial main map has one point over \(O\).
Separating the \(p^2\) sections therefore
puts every mark on a rational marked tree
attached there. No mark remains at its
attachment node. Every chosen cusp has a
rational joining path to another, including
both endpoints. Statements A.41--A.42 now
give the full fixed-inertia conclusion at
every coefficient prime and every permitted
packet. Statement A.59 gives every normalizer
because the actual geometric Galois group
is a \(p\)-group.

The exact program \texttt{code/a\_heisenberg\_inert.py}
checks the
division identities, divided operators,
Frobenius roots \(R,A,S\), the identity
\(DS=RA/a\), and all chart coefficients at
several inert primes. It is a certificate
for the formulas; the uniform argument
(A.61.8)--(A.61.21) is what proves the
assertion for every inert prime.
\end{proof}

% A.62: full independent root and PBC proof/source/certificate audits PASS.
\subsection{The ordinary-prime Heisenberg family by actual isogeny descent}

\begin{statement}{A.62}{Proved}{Both assertions for every ordinary-prime Heisenberg cover}
For every prime \(p\equiv1\pmod4\), let \(C_p\) be the smooth
proper normalization of \(z^p=\psi_p(x)\) over
\(E:y^2=x^3-x\), and let
\[
 Y_p=C_p\setminus q^{-1}(E[p])\xrightarrow{[p]q}U.
                                                        \tag{A.62.1}
\]
This actual degree-\(p^3\) Heisenberg cover has proper genus
\(1+p^2(p-1)/2\) and \(p^2\) cusps. Both assertions of
Task A.5 hold at every coefficient prime and for every
permitted cusp packet and open derived subgroup.

At a place above \(p\), its proper stable reduction has one
ordinary elliptic main component, \(p\) tails of the form
\(V^p-V=U^4\), and \(p(p-3)\) tails of the form
\(V^p-V=U^2\), over the algebraic residue field. Every
remaining component is rational and the dual graph is a tree.
The \(p^2\) marks form \(p\) rational clusters of \(p\)
marks each, disjoint from every positive-genus tail.

Together with A.61, this proves both all-prime assertions for
the actual exponent-\(p\) Heisenberg cover at every odd prime
\(p\). It makes no cofinality assertion about all finite covers.
\end{statement}

\begin{proof}
All special-fibre functions below are initially defined over
\(\F_p\); finite unramified extensions will separate the
points used in the actual charts. Since \(p\equiv1\pmod4\),
one has \(p\ge5\).

\paragraph{The separable isogeny and the actual reduced division polynomial.}
The coefficient of \(x^{p-1}\) in
\((x^3-x)^{(p-1)/2}\) is, up to sign,
\(\binom{(p-1)/2}{(p-1)/4}\), and is nonzero modulo \(p\).
Thus the good elliptic reduction is ordinary by
\cite[Theorem V.4.1(a), p.148]{fghSilvermanAEC}.
Its reduced \(p\)-torsion subgroup \(H\) has order \(p\),
and the inseparable degree of \([p]\) is \(p\).

There is an actual degree-\(p\) CM endomorphism
\(\alpha:E\to E\), over \(\Z_p\), with unit differential.
Indeed choose a Gaussian factor of \(p\) that is a unit
under the chosen embedding \(\Z[i]\to\Z_p\), and apply
that endomorphism to the smooth elliptic group scheme.
Duality sends the degree-one automorphism \([i]\)
to its inverse \([-i]\) and is additive. Hence the
dual of this Gaussian endomorphism is its conjugate;
their composite is \([p]\). The degree is therefore
\(p\), by \cite[Theorem III.6.2(a),(c),(d), p.83]{fghSilvermanAEC}.
Its kernel is finite etale over \(\Z_p\), since its
differential is a unit, and its special kernel is exactly
\(H\). Write \(\overline\alpha^*\omega=\mu\omega\),
where \(\omega=dx/(2y)\) and \(\mu\in\F_p^\times\).
The endomorphism commutes with the quarter-rotation
\(\rho(x,y)=(-x,iy)\).

The exact numerator and denominator in
\[
 x\circ[p]=\phi_p/\psi_p^2
                                                        \tag{A.62.2}
\]
are relatively prime after reduction; \(\phi_p\) is monic
of degree \(p^2\), by
\cite[Exercise III.3.7(a)--(d), pp.105--106]{fghSilvermanAEC}.
At every point of \(H\), the special map \([p]\) has
ramification index \(p\). At each nonzero such point,
the denominator therefore has a zero of order \(2p\).
At \(O\), the pole order of the quotient is \(2p\).
Comparing degrees in (A.62.2) gives
\[
 \overline\psi_p=\overline g^p,\qquad
 \deg\overline g=(p-1)/2,\qquad
 \operatorname{div}(\overline g)=D_H-p(O).
                                                        \tag{A.62.3}
\]
Here the coefficients have unique Frobenius roots in
\(\F_p\). The zeros of \(\overline g\) are the nonzero
points of \(H\), each simple. As in A.61, \(\psi_p\)
is even in \(x\); therefore so is \(\overline g\).
Choose an even integral polynomial lift \(g\) of the
same degree and set
\[
 \begin{gathered}
 h=\frac{\psi_p-g^p}{p}\in\Z_p[x],\qquad H_0=\frac{h}{g^p},\\
 \frac{\psi_p}{g^p}=1+pH_0.
 \end{gathered}
                                                        \tag{A.62.4}
\]
The subscript distinguishes this rational function from
the reduced torsion group. A different integral lift of
\(\overline g\) gives the same \(\overline h\), since
\((g+pb)^p\equiv g^p\pmod{p^2}\).
The degree of \(h\) is \((p^2-1)/2\), with leading
coefficient one. In particular \(\overline H_0\) has
an exact pole of order \(p-1\) at \(O\).

\paragraph{Actual translation comparison modulo p squared.}
We prove that \(d\overline H_0\) is invariant under
every translation in \(H\). This comparison uses the
actual division function, not only its branch divisor
after reduction.

After a finite unramified extension \(R/\Z_p\), every
point of \(\ker\alpha\) is an \(R\)-section. Fix one,
\(T\), and write \(K=\operatorname{Frac}(R)\).
There is a function \(k_T\in K(E)^\times\) with
\(\operatorname{div}(k_T)=p(O)-p(-T)\), since \(pT=0\).
The characteristic-zero division divisor consequently gives
\[
 \frac{\psi_p(P+T)}{\psi_p(P)}=c_T k_T(P)^p,
                         \qquad c_T\in K^\times.
                                                        \tag{A.62.5}
\]
Let \(v_G\) be the Gauss valuation at the generic point
of the smooth elliptic special fibre. Its value group
on \(K(E)\) is \(\Z\), with uniformizer \(p\).
Scale \(k_T\) by a power of \(p\) so that \(v_G(k_T)=0\).
The two division functions in (A.62.5) are Gauss units,
by (A.62.3), so \(c_T\in R^\times\).

Choose \(\gamma_T\in R^\times\) with
\(\overline\gamma_T^p=\overline c_T\), possible
because the residue field is perfect. Reduction of
(A.62.5) and injectivity of Frobenius in its function
field show that
\[
 q_T(P)=\frac{g(P+T)}{g(P)k_T(P)\gamma_T}
                  \equiv1\pmod p
                       \quad\text{in the Gauss DVR}.
                                                        \tag{A.62.6}
\]
Thus \(q_T=1+pb_T\) with \(b_T\) in that DVR, and
\(q_T^p\equiv1\pmod{p^2}\). Also
\(c_T/\gamma_T^p=1+p\epsilon_T\), with
\(\epsilon_T\in R\). The exact equations now give
\[
 \begin{aligned}
 1+pH_0(P+T)
   &=\frac{c_T}{\gamma_T^p q_T(P)^p}(1+pH_0(P)),\\
 \overline H_0(P+\overline T)
   &=\overline H_0(P)+\overline\epsilon_T.
 \end{aligned}
                                                        \tag{A.62.7}
\]
All divisibilities here take place over an unramified
field, so reduction modulo \(p^2\) has its literal
meaning. The possibly nonzero constant in the last
line is harmless; its differential is zero.
The residues of these sections run through all of
\(H\), proving the desired invariance.

\paragraph{Descent, pole order, and every critical zero.}
Let \(D\) be the invariant derivation
\(Dx=2y\), \(Dy=3x^2-1\), and put
\[
        F=D\overline H_0,
        \qquad U=D\overline g/\overline g.
                                                        \tag{A.62.8}
\]
The exact characteristic-zero identity (A.61.11), with
\(\psi_p=g^p+ph\), gives on division by \(p\) and
reduction
\[
 \begin{gathered}
 D^2\overline h
   =\overline g^{p-2}(D\overline g)^2
      -\overline g^{p-1}D^2\overline g,\\
 DF=-DU.
 \end{gathered}                       \tag{A.62.9}
\]
Translation invariance from (A.62.7) proves that the
actual function \(F\) descends through
\(\overline\alpha\): write \(F=f\circ\overline\alpha\).

We also determine \(DU\) exactly. For \(T\in H\),
the divisor of \(\overline g(P+T)/\overline g(P)\)
is \(p(O)-p(-T)\). Its logarithmic differential has
no pole, because all these orders are divisible by
\(p\). It is therefore a constant multiple of
\(\omega\). Thus \(U(P+T)-U(P)\) is constant and
\(DU\) is \(H\)-invariant. It has at most double
poles, precisely over \(H\), so its descent belongs
to \(L(2O)=\langle1,x\rangle\) on the quotient elliptic
curve, identified with \(E\) by \(\overline\alpha\).
The quarter-rotation sends \(DU\) to \(-DU\),
because \(\overline g\) is invariant and \(\rho\)
scales \(D\) by \(-i\). It sends the quotient
\(x\)-coordinate to \(-x\). The constant term in
that two-dimensional space must consequently vanish.

At \(O\), the pole order of \(\overline g\) is
\(p-1\), so in an invariant local parameter to
the required first order \(U=s^{-1}+O(1)\) and
\(DU=-s^{-2}+O(s^{-1})\). Since
\(x\circ\overline\alpha=\mu^{-2}s^{-2}+O(s^{-1})\),
the leading coefficient determines
\[
       DU=-\mu^2x\circ\overline\alpha,
          \qquad Df=\mu x.                           \tag{A.62.10}
\]
Only first-order local expansions are used here;
no characteristic-\(p\) formal logarithm is assumed.

The only possible poles of \(F\) lie in \(H\).
Its pole at \(O\) has exact order \(p\), by the
pole order \(p-1\) of \(\overline H_0\), which is
not divisible by \(p\). The isogeny is etale;
therefore \(f\) has exactly one pole, at \(O\),
of exact order \(p\). The function \(F\), and
hence \(f\), is odd under \([-1]\).
In particular \(f(0,0)=0\). Equation (A.62.10)
has a zero of order two there, so \(f\) has a
zero of order three. Indeed \(3<p\), and the
first nonzero differentiated coefficient is of
degree two. At every other zero of \(f\), the
right side is nonzero, so the zero is simple.
Its zero divisor has total degree \(p\). Thus
there is one triple zero and \(p-3\) simple zeros.
Pullback through the etale degree-\(p\) isogeny gives
exactly \(p\) triple and \(p(p-3)\) simple zeros
of \(F\), all outside \(H\).

The differential \(d\overline H_0=F\omega\)
has these same zero orders. At a zero of order
\(k-1\), with \(k=4\) or \(2\), its primitive
\(\overline H_0-\overline H_0(P)\) has exact
order \(k\): every smaller degree is less than
\(p\), so differentiation detects it. In particular
there is no hidden lower degree-\(p\) coefficient
in this case. The two-torsion centres use \(y\)
as parameter; at all others \(x-x(P)\) is a parameter.

\paragraph{Exact main and all tail charts.}
Choose \(\pi^p=p\) and set \(z=g(x)+\pi u\).
The exact equation, divided by \(p\), has special
function field
\[
                         u^p=\overline h(x).
                                                        \tag{A.62.11}
\]
Its differential is nonzero because \(F\ne0\).
This is the degree-\(p\) radicial field
\(k(E)^{1/p}\), with ordinary elliptic normalization.
The chart is smooth at its special generic point,
by the differential and the elliptic equation.
It therefore supplies an actual elliptic component,
not just a formal equation.

Separate all critical centres over a finite
unramified extension of \(\Z_p\). At a non-two-torsion
centre lift its \(x\)-coordinate arbitrarily in that
unramified ring and lift \(y\) by Hensel's lemma;
its \(y\)-coordinate is a unit. At a two-torsion
centre use \((a,0)\), \(a=0,1,-1\), itself.
Every such \(g(a)\) is a unit, since the centres
avoid \(H\). Set
\[
 \begin{aligned}
 H_a&=h(a)/g(a)^p,&
 N_a(x)&=h(x)-H_a g(x)^p,\\
 F_a&=1+pH_a,& z_a^p&=F_a,\\
 \eta^{p-1}&=-p,& \rho^k&=\eta,
                \qquad z=g(x)z_a(1+\eta V).
 \end{aligned}                                         \tag{A.62.12}
\]
The exact binomial equation is
\[
 \begin{aligned}
 \mathcal B_\eta(V)
   &=\sum_{j=1}^{p-1}\frac{\binom pj}{p}\eta^{j-1}V^j-V^p,\\
 g(x)^pF_a\mathcal B_\eta(V)&=N_a(x)/\eta.
 \end{aligned}                                         \tag{A.62.13}
\]
At a non-two-torsion lift \((a,y_a)\), put
\(x=a+\rho U\), \(y=y_a+\rho W\). The exact
elliptic equation is
\[
 2y_aW+\rho W^2
          =(3a^2-1)U+3a\rho U^2+\rho^2U^3.
                                                        \tag{A.62.14}
\]
The first nonconstant Taylor degree of
\(\overline N_a\) is \(k\), with a nonzero
coefficient. Every smaller integral coefficient
is divisible by \(p\), because the lifts are
unramified. After division by \(\eta\), a lower
degree \(j<k\) has valuation at least
\[
                 1+\frac{j-k}{(p-1)k}>0;
                                                        \tag{A.62.15}
\]
a higher degree has valuation at least
\((j-k)/((p-1)k)>0\). Thus only degree \(k\)
survives. At a two-torsion centre put instead
\(y=\rho U\), \(x=a+\rho^2W\), obtaining
\[
 U^2=(3a^2-1)W+3a\rho^2W^2+\rho^4W^3.
                                                        \tag{A.62.16}
\]
The first nonzero \(x-a\) degree is \(k/2\).
The same valuation estimates with \(j\) replaced
by \(2j\) leave a unit multiple of \(U^k\).

In both cases the special equations have unit
partials in \(W,V\), and the resulting curve is
\(V^p-V=cU^k\), \(c\ne0\). These are integral
smooth algebraic charts whose generic coordinate
changes are invertible on the actual covering field.
The special curve is geometrically integral since
\(k\) is prime to \(p\). Over the algebraic residue
field its nonzero coefficient can be absorbed in \(U\).
At infinity, with \(t=U^{-1}\), one has
\(v(t)=p\) and \(v(V)=-k\). As \(k\mid p-1\),
the element \(tV^{(p-1)/k}\) is a uniformizer;
its difference under a nonzero translation of \(V\)
has valuation \(k+1\). The different is
\((p-1)(k+1)\), and Hurwitz gives genus
\((p-1)(k-1)/2\). The one-branch-point
Deuring--Shafarevich formula gives special \(p\)-rank
zero for every tail \cite[Corollary 1.8]{fghCrew1984}.

\paragraph{Genus exhaustion, the actual map, and all marked points.}
The characteristic-zero division divisor, translation
lifts and quarter-rotation argument of A.61 identify
the actual deck group as the exponent-\(p\) Heisenberg
group and give the asserted degree, cusps and genus.
No feature of that argument requires inertness.

The positive-genus main and all tail valuations just
constructed persist in the stable model, by the same
common-regular-resolution argument as A.58 and A.61,
using \cite[Theorem 9.2.2, Remark 9.2.6 and
Theorem 9.3.21]{fghLiu}. They are distinct since their
restrictions to the elliptic model have distinct
finite centres, or its generic point for the main.
Their genera sum to
\[
 1+p\frac{3(p-1)}2+p(p-3)\frac{p-1}2
                  =1+\frac{p^2(p-1)}2.                \tag{A.62.17}
\]
The semistable genus formula therefore forces compact
type and makes every remaining component rational.
The actual map \(q\) extends to the good elliptic
model after the same regular resolution and descent
of rational exceptional curves. The main component
already maps with degree \(p\), so the degree of a
pulled-back ample line bundle makes every other
component map constantly. Each off-main subtree
has one attaching point and the same constant image.

The ordinary CM factorization splits the generic
\(p\)-torsion into the etale factor of order \(p\)
and the connected factor of order \(p\). The first
specializes bijectively to \(H\), and the second
specializes to \(O\). Thus exactly \(p\) distinct
generic points of \(E[p]\) specialize to each point
of \(H\). Each has one ramified cusp on \(C_p\).
The main radicial map has one point above each
point of \(H\), while all positive-genus tails
and their connecting subtrees lie over centres
outside \(H\). Separating the \(p\) marks at
each main point therefore puts them on a rational
marked tree. No mark stays at its attachment node.
Every selected cusp is joined to another by a path
consisting entirely of rational components, including
both endpoints.

Statement A.42 now gives full fixed inertia at every
coefficient prime and permitted packet, and A.59
gives all normalizers. The exact certificate
\path{code/a_heisenberg_ordinary.py} checks the division differential,
the actual separable-isogeny descent and every
critical multiplicity at several ordinary primes.
The unramified translation comparison
(A.62.5)--(A.62.10) proves the assertion uniformly.
\end{proof}

% A.63: full independent root and DE proof/source/certificate audits PASS.
\subsection{Every isogeny enlargement of an ordinary Heisenberg cover}

\begin{statement}{A.63}{Proved}{All ordinary Heisenberg enlargements have cusp fixed groups and normalizers}
Let \(p\equiv1\pmod4\) be prime. Let \(C_p\) be the smooth
proper normalization of \(z^p=\psi_p(x)\) over
\(E:y^2=x^3-x\), with its actual degree-\(p\) map \(q\).
For every integer \(m\ge1\), put
\[
 S_{p,m}=q^{-1}(E[pm]),\qquad
 Y_{p,m}=C_p\setminus S_{p,m}
       \xrightarrow{[pm]q} E\setminus\{O\}.
                                                        \tag{A.63.1}
\]
This is a finite etale cover of degree \(p^3m^2\), with
\[
 g(C_p)=1+\frac{p^2(p-1)}2,\qquad
 |S_{p,m}|=p^3m^2-p^3+p^2.                              \tag{A.63.2}
\]
Both assertions of Task A.5 hold for every
coefficient prime, every selected cusp packet, and every
permitted open derived subgroup for these actual covers.

More precisely, at a place above \(p\), every cusp has a
distinct cusp with the same specialization on the unmarked
proper stable model. In the stable marked model their joining
path contains only rational components. No assertion that
odd-order torsion avoids the positive-genus tails is used.
There is no assertion that this family is cofinal among all
finite covers.
\end{statement}

\begin{proof}
We work after a finite number-field extension making the
actual deck group and the finite marking set rational, and
after a further extension at a place above \(p\) admitting
the proper stable model. All valuations below are normalized
by \(v(p)=1\). Such finite extensions are permitted in the
open-derived comparison of A.41--A.42.

\paragraph{The actual cover and its finite deck action.}
By A.62, the degree-\(p^3\) map \([p]q\) has the actual
exponent-\(p\) Heisenberg group \(G\). Its centre \(Z\)
has order \(p\), is generated by
\(\sigma:z\mapsto\zeta_pz\), and fits in
\[
 1\longrightarrow Z\longrightarrow G
       \longrightarrow E[p]\longrightarrow1.
                                                        \tag{A.63.3}
\]
The last map records the actual translation on the target of
\(q\). The only nontrivial stabilizers of geometric points
of \(C_p\) under \(G\) are exactly \(Z\), at the original
cusps \(q^{-1}(E[p])\). Indeed a fixed point under a deck
element must map under \(q\) to a point fixed by its target
translation; a nonzero translation has no fixed point.
The residual possibility is the central cyclic cover \(q\),
whose branch locus is exactly \(E[p]\), with index \(p\)
at every point. This follows also from
\(\operatorname{div}(\psi_p)=D_{E[p]}-p^2(O)\).

The marking set \(S_{p,m}\) is \(G\)-stable, since
translations by \(E[p]\) preserve \(E[pm]\).
The composite in (A.63.1) is etale away from this set:
\([pm]\) is etale in characteristic zero, and the entire
branch set of \(q\) has been removed. Its degree is
\(p(pm)^2\). There are \(p^2\) branch fibres with one
point each, and \(p^2m^2-p^2\) other fibres with \(p\)
points each. This proves (A.63.2), including the exact cusp
count and the Euler identity \(2g-2+|S_{p,m}|=p^3m^2\).

\paragraph{The small connected translations.}
Let \(E[p]^0\subset E[p](\overline K_v)\) be the subgroup
specializing to the origin of the ordinary good reduction.
It has order \(p\). For a formal parameter \(t\) at \(O\),
every nonzero \(T\in E[p]^0\) satisfies
\[
                         v(t(T))=\frac1{p-1}.           \tag{A.63.4}
\]
Here is the exact valuation argument. The formal group is
defined over \(\Z_p\), and ordinarity makes its reduction
have height one. Thus
\[
 [p](t)=pt+\sum_{2\le j<p}p b_jt^j+u_pt^p
                       +\sum_{j>p}u_jt^j,
 \quad b_j,u_j\in\Z_p,\quad u_p\in\Z_p^\times.
                                                        \tag{A.63.5}
\]
The divisibility of the coefficients with \(2\le j<p\)
is \cite[Corollary IV.4.4, p.126]{fghSilvermanAEC};
the height and inseparable-degree comparison is
\cite[Theorem IV.7.4, p.134]{fghSilvermanAEC}.
For a nonzero root with \(v(t)>0\), the terms with
\(2\le j<p\) have valuation strictly larger than
\(1+v(t)\), and those with \(j>p\) have valuation
strictly larger than \(p v(t)\). The two remaining terms
must have equal valuation. This proves (A.63.4).

Define the actual order-\(p^2\) subgroup
\[
             G_c=\operatorname{pr}^{-1}(E[p]^0)\subset G.
                                                        \tag{A.63.6}
\]
It contains \(Z\) and is isomorphic to
\((\Z/p)^2\): every group of order \(p^2\) is abelian,
and the ambient group has exponent \(p\).

\paragraph{Its action on every positive-genus component.}
Use the actual proper stable fibre constructed in A.62.
It has one elliptic main component, mapping radicially of
degree \(p\) to the good elliptic special fibre, and the
following other positive-genus components:
\[
 V^p-V=U^k,\qquad k=2\text{ or }4.                       \tag{A.63.7}
\]
There is one such tail above each critical centre specified
in A.62. The genus exhaustion there proves that these are
all positive-genus components and that the dual graph is a
tree. Any remaining proper-stable component is rational.
The group action extends to this stable model by
\cite[Proposition 10.1.16 and Theorem 10.3.34(b),
pp.460--461 and 518--519]{fghLiu}.

Every element of \(G_c\) reduces to the identity translation
on the base elliptic curve. Hence it fixes the main component
and acts identically on its function field: the inclusion
from the base is purely inseparable, which admits no
nontrivial automorphism over the base.

We check the tail action in its actual scaled chart. Let
\(s\) be the good integral parameter at its unramified
centre used in A.62: it is \(x-a\) away from two-torsion,
and \(y\) at the relevant two-torsion centre. The tail
valuation is represented by
\[
 s=\rho U,\qquad
 v(\rho)=\frac1{k(p-1)},\qquad
 z=g z_a(1+\eta V),\quad v(\eta)=\frac1{p-1}.
                                                        \tag{A.63.8}
\]
In the two-torsion chart \(x-a\) has order \(s^2\), as
in A.62; the parameter in (A.63.8) is still \(s=y\).
Any constant rescaling making (A.63.7) monic is a unit and
does not change these valuations.

The smooth integral group law gives, on the completed
neighbourhood of this centre,
\[
                 s(P+T)-s(P)\in t(T)R[[s,t(T)]].        \tag{A.63.9}
\]
To justify this inclusion, write translation by a formal
point with parameter \(t\) as an integral power series
in \(s,t\). At \(t=0\) it is the identity, so the
difference is divisible by \(t\). All coefficients are
integral because the centre is a smooth point of the good
elliptic group scheme. Evaluating the tail Gauss valuation
and using (A.63.4) gives
\[
 v\bigl(s(P+T)-s(P)\bigr)\ge\frac1{p-1}
           >\frac1{k(p-1)}=v(\rho).                    \tag{A.63.10}
\]
Thus translation by \(T\in E[p]^0\) preserves the actual
tail disc and induces the identity on its reduced variable
\(U\). There is exactly one positive-genus component
above this centre, so every lift in \(G_c\) preserves the
tail. Its action on (A.63.7) fixes \(U\); consequently it
is \(V\mapsto V+c\) for some \(c\in\F_p\), because
these are exactly the automorphisms of the degree-\(p\)
Artin--Schreier extension over \(\overline k(U)\).

The central generator acts nontrivially on this tail. In
the actual chart (A.63.8) it sends
\[
       V\longmapsto\zeta_pV+\frac{\zeta_p-1}{\eta}.
                                                        \tag{A.63.11}
\]
The second summand is a unit, since
\(v(\zeta_p-1)=v(\eta)=1/(p-1)\); its reduction is a
nonzero element of \(\F_p\), as also follows by
preserving (A.63.7). The tail-action homomorphism is
therefore a surjection
\[
 G_c\longrightarrow\F_p,\qquad
 K_B:=\ker(G_c\longrightarrow\operatorname{Aut}(B)),
 \quad |K_B|=p,\quad K_B\cap Z=1.                       \tag{A.63.12}
\]
Here \(B\) denotes that individual tail. In particular
\(K_B\) is a noncentral subgroup acting identically on
the whole smooth proper tail, including its point at
infinity. This is an actual group kernel, not a kernel of
an approximate arithmetic representation.

\paragraph{Rational components and the exact marked specialization.}
Every positive-genus vertex of the proper stable tree is
fixed by \(G_c\), by the preceding analysis. Every leaf
of that tree is positive-genus: an unmarked rational leaf
would be unstable. An automorphism of a finite tree fixing
all its leaves fixes the whole tree. For example, every
non-leaf vertex is determined by the paths between the
leaves in its incident branches; vertices of valence two
also occur on such fixed paths. Consequently \(G_c\)
fixes all vertices and edges. Each rational proper-stable
component has at least three distinct nodes, and these
nodes are fixed individually. An automorphism of
\(\mathbf P^1\) fixing three points is the identity.
Thus \(G_c\) acts pointwise on every rational component
and fixes every node of the proper stable fibre.

Let \(P\in S_{p,m}\), and let \(x\) be its specialization
on this unmarked stable model. If \(x\) is a smooth point
of an Artin--Schreier tail \(B\), choose a nonidentity
element of \(K_B\). If \(x\) is on the main component,
on a rational component, or is a node, choose any element
of \(G_c\setminus Z\). In either case the element fixes
\(x\) and does not fix \(P\), by the exact generic
stabilizer statement following (A.63.3). We obtain
\[
                    Q\ne P,\qquad Q\in S_{p,m},
                    \qquad \operatorname{sp}(Q)=x.
                                                        \tag{A.63.13}
\]
In fact the chosen cyclic subgroup gives \(p\) distinct
marked points with that specialization. The choice may
depend on the component; no single noncentral subgroup is
asserted to act identically on every tail at once.

We spell out why (A.63.13) is the pointed hypothesis
needed in A.41--A.42. Forget the marks in a semistable
marked model and stabilize. This contracts only rational
components. Above a smooth point of the unmarked model
the exceptional fibre is a rational tree attached at that
point, and above a node it is a rational chain with
possibly attached rational trees. The fibres are
connected, and the ambient dual graph is a tree. Since
the marked sections are disjoint and meet only smooth
points, two distinct sections specializing to the same
unmarked point are separated on this rational exceptional
fibre. Neither can stay at an attaching node of a
positive-genus component. Their joining path, including
both endpoints, therefore contains only rational
components. The same conclusion is immediate if their
unmarked component is already rational. This describes
the contraction fibre, not just the image point on the
elliptic base.

\paragraph{The actual fixed group and normalizer conclusions.}
The rational companion just constructed works for every
selected cusp. A.41--A.42 apply directly to \(Y_{p,m}\),
at residue characteristic \(p\), and give
\[
 \bigl(\pi_1(Y_{p,m,\overline L})^{(\ell)}\bigr)^{s(\Delta_0)}
                        =I_{P,\ell}                    \tag{A.63.14}
\]
for every coefficient prime \(\ell\), every permitted
packet, and every permitted open derived subgroup.
As there, choose a finite field of definition whose
closed derived subgroup lies in the given \(\Delta_0\);
the latter already fixes the selected inertia, so the
conclusion passes to it. This is a direct geometric
application, not a fixed-group lift through a quotient.

For normalizers use the actual open immersion
\(Y_{p,m}\hookrightarrow Y_{p,1}\). It induces the
identity on the proper Jacobian. If a noncuspidal closed
procyclic subgroup of the geometric pro-\(\ell\) group
had open arithmetic normalizer, A.11 would give it
nonzero image in proper rational homology. Its image in
the target pro-\(\ell\) group is then infinite
procyclic, still has open arithmetic normalizer, and has
the same nonzero proper-homology image. A.62 puts this
image in target cusp inertia, which has zero proper
homology, a contradiction. No assertion about the
nilpotent Galois closure of the enlarged composite is
needed.

The exact certificate \path{code/a_heisenberg_ordinary_enlargements.py}
checks the counting identities, both strict radius
inequalities, all tail-kernel characters in several
actual finite Heisenberg models, and an ordinary
critical point of exact odd order. The uniform argument
above proves the statement for all \(p,m\).
\end{proof}

% A.64: root and PBC full proof/source/certificate audits PASS.
\subsection{A prime-to-closure enlargement criterion and all Heisenberg enlargements}

\begin{statement}{A.64}{Proved}{Torsion enlargement outside the geometric closure order}
Let \(Y\to U=E\setminus\{O\}\) be any geometrically
connected finite etale cover from Task A.5, with degree
\(d\), proper compactification \(C\), genus \(g\), and
extended map \(q:C\to E\). Let \(N\) be the order of
its geometric Galois closure. If an integer \(m\ge1\)
has a prime divisor \(r\nmid N\), then
\[
 Y_m=C\setminus q^{-1}(E[m])
                \xrightarrow{[m]q}U                   \tag{A.64.1}
\]
has full pro-\(\ell\) fixed group equal to the selected
cusp inertia, for every coefficient prime \(\ell\),
every cusp packet, and every permitted open derived
subgroup. It has degree \(m^2d\), proper genus \(g\),
and \(m^2d-2g+2\) cusps. The residue prime \(r=2\)
is allowed. If \(Y\) has pro-\(\ell\) cuspidal
normalizer recognition, then so does \(Y_m\).

As a consequence, for every odd prime \(p\) and every
integer \(m\ge1\), both assertions of Task A.5 hold
for the actual cover
\[
 C_p\setminus q_p^{-1}(E[pm])
                  \xrightarrow{[pm]q_p}U,             \tag{A.64.2}
\]
where \(C_p\) is the smooth proper normalization of
\(z^p=\psi_p(x)\) and \(q_p:C_p\to E\).
Its degree is \(p^3m^2\), genus
\(1+p^2(p-1)/2\), and cusp count
\(p^3m^2-p^3+p^2\).

The map \(Y_m\hookrightarrow Y\) is an open immersion.
These assertions do not give a finite refinement over the
original map \(Y\to U\), and they do not lift fixed
elements from a quotient fundamental group. The universal
Task A.5 remains open.
\end{statement}

\begin{proof}
\paragraph{An actual good tame model at every allowed residue prime.}
Choose a place above \(r\). The elliptic curve
\(E:y^2=x^3-x\) has potentially good reduction there
because \(j(E)=1728\) is integral; this includes
\(r=2\), by \cite[Proposition VII.5.5, p.197,
and Appendix A.1.4(b)]{fghSilvermanAEC}.
After finite extension, take its smooth proper group
scheme with the smooth origin section.

Since the actual geometric Galois closure has order
prime to \(r\), prime-to-\(r\) specialization for
the good pair \((E,O)\) extends that cover tamely
over this model after finite extension. The relevant
comparison, including inertia, is
\cite[Expos\'e XIII, Corollary 2.9, Section 2.10 and
the proof of Corollary 2.12]{fghSGA1}. This is the
same actual finite-cover comparison proved in A.56
and used in A.59: a geometric prime-to-\(r\)
cover of the generic fibre corresponds to a tame
cover of the special fibre, and finite descent
defines the required cover and its generic
identification over a finite extension. A possible
arithmetic twist is removed by that extension.

Normalize the good elliptic model in the actual
function field of \(C\). The resulting proper
source is smooth. Away from the origin it is
etale. At a point above the origin, after strict
henselian localization, the tame local map is
\[
                    R[[t]]\longrightarrow R[[u]],
                              t=u^e,\quad r\nmid e.
                                                        \tag{A.64.3}
\]
The indices divide the prime-to-\(r\) closure
order, so this argument applies to subgroup
quotients as well as to the Galois closure.
Thus the conclusion is potentially good
\emph{curve} reduction of \(C\), with the actual
map finite and etale outside the origin, rather
than good reduction of its Jacobian alone.
Nothing in this argument excludes residue two.

Enlarge the field further to split every point
in the finite set \(q^{-1}(E[m])\). The good
source and target and the tame local descriptions
persist under this base change. A nonconstant map
to an elliptic curve forces \(g\ge1\).

\paragraph{Every target torsion point has a nontrivial specialization cluster.}
Write \(m=r^a n\), where \(a\ge1\) and \(r\nmid n\).
Let \(f\in\{0,1\}\) be the \(r\)-rank of the good
elliptic special fibre. The geometric special
\(r^a\)-torsion has \(r^{af}\) points, while its
generic fibre has \(r^{2a}\) points. These facts
hold also in characteristic two; see
\cite[Theorem V.3.1, p.145]{fghSilvermanAEC}.

For precision, every special point lifts to a
generic torsion point after finite extension,
and every specialization fibre has the same size.
Indeed the finite flat group scheme \(E[r^a]\)
over a strict henselian discrete valuation ring
is the disjoint union of the finite flat pieces
whose special fibres are its individual geometric
points: lift the idempotents of its special
coordinate algebra. Each piece has a nonempty
generic fibre. Translation by a generic lift of
a special point permutes these pieces, so all
pieces have the same rank. Their total rank is
\(r^{2a}\). The generic fibre is etale, and hence
each piece contains exactly
\[
                             c=r^{a(2-f)}\ge2           \tag{A.64.4}
\]
generic geometric points. The \(n\)-torsion is
finite etale and specializes bijectively after
splitting. Decomposing \(E[m]\) into its
\(r\)-primary and prime-to-\(r\) parts therefore
proves that \emph{every} geometric specialization
cluster of \(E[m]\) has exactly \(c\) points.
This includes clusters away from the origin when
the reduction is ordinary; they are not omitted.

\paragraph{The exact source clusters.}
Let \(b\) be a special point of the target
underlying one of these torsion clusters.
If \(b\ne O\), the finite source map is etale
above its whole formal disc. Each special source
point above \(b\) receives precisely one lift
of each of the \(c\) generic target points,
so its source cluster has size \(c\ge2\).

If \(b=O\), let \(P\) be an original cusp
with index \(e_P\). The target cluster contains
the origin itself and \(c-1\) other points.
Equation (A.64.3) gives one point above the
origin and \(e_P\) distinct characteristic-zero
points above each other target point. All of
them specialize to \(P\). The source cluster
size is exactly
\[
                              1+e_P(c-1)\ge2.          \tag{A.64.5}
\]
The unit extraction used in (A.64.3) is legitimate
because \(e_P\) is invertible in the residue
field. These two cases account for every mark
of \(q^{-1}(E[m])\).

The proper source model is smooth. Distinct
marked sections specializing to one of its
points separate, after semistable marked
resolution, on a rational tree attached at
that point. Since every cluster has at least
two marks, every selected cusp has a distinct
companion whose entire joining path, including
the endpoints, is rational. This is the exact
pointed hypothesis of A.41--A.42. It yields the
full integral fixed-inertia conclusion at
every coefficient prime, including \(\ell=r\),
and for every packet. The finite global field
restriction can be chosen with its closed
derived subgroup inside the prescribed open
derived subgroup, exactly as in A.41.

The original indices satisfy
\(\sum_P e_P=d\) and
\(2g-2=\sum_P(e_P-1)=d-|q^{-1}(O)|\).
In characteristic zero, all other target
points are unramified. Hence
\[
 \begin{split}
 |q^{-1}(E[m])|
   &=|q^{-1}(O)|+d(m^2-1)\\
   &=m^2d-2g+2.
 \end{split}                                            \tag{A.64.6}
\]
The map (A.64.1) is finite etale of degree
\(m^2d\), and its source remains geometrically
connected. The normalizer assertion is the
actual-open-immersion argument of A.53 using
A.11 and unchanged proper homology; it does
not require a fixed-element lifting assertion.

\paragraph{Every odd Heisenberg prime and every enlargement.}
Let \(Y_p=C_p\setminus q_p^{-1}(E[p])\).
Its degree-\(p^3\) map to \(U\) is
\(f_p=[p]q_p\), whose geometric Galois closure
has the actual order \(p^3\), by A.61--A.62.
Moreover
\[
                   f_p^{-1}(E[m])=q_p^{-1}(E[pm]).       \tag{A.64.7}
\]
If \(m\) has a prime factor different from
\(p\), apply the general conclusion above to
\(f_p\). This treats every such \(m\), with
no congruence restriction on that other prime.
If \(p\equiv1\pmod4\), A.63 treats all
remaining \(m\) as well.

It remains to take \(p\equiv3\pmod4\) and
\(m=p^a\), \(a\ge0\). Use the actual
compact-type model of A.61 at residue \(p\).
The target elliptic curve is supersingular,
so all points of \(E[p^{a+1}]\) specialize to
\(O\). Its proper source has one elliptic
main component mapping radicially to the
good target, and its other positive-genus
components map constantly to the finite
zeros of the function \(R\) in A.61.
None of these centres is \(O\), where
\(R\) has its pole.

Consequently no enlarged mark can specialize
on any of those positive-genus tails. On the
main component there is exactly one point
mapping to \(O\), by radiciality. Every
rational component carrying an enlarged
mark lies in the fibre over \(O\): its map
to a good elliptic curve is constant. The
path from such a component to the main
cannot pass through a positive-genus tail,
because the map at each connecting node
would then have both values \(O\) and a
different critical centre. All rational
parts of this fibre attach at the unique
main point over \(O\). Nodalness allows
only one branch at that point on the main,
so these parts form one connected rational
subtree, if they are present.

If there is no such rational part in the
unmarked model, all marks specialize to
that single smooth main point and marked
separation creates a rational tree there.
If there is already a rational part, the
main attachment is a node, so the smooth
marked sections lie on that rational part
or on rational components inserted above
its nodes. Thus every joining path between
the enlarged marks is rational. The set
contains the original \(p^2\ge9\) cusps,
so every mark has a companion. A.41--A.42
again prove all-prime fixedness.

A.61--A.62 give all-prime normalizer
recognition on \(Y_p\); A.53's actual open
immersion argument gives it on every
enlargement (A.64.2). This completes the
asserted consequence while retaining the
universal open frontier stated above.

The executed certificate
\path{code/a_torsion_enlargement_closure.py}
checks both specialization-cluster counts
and the Euler identity for actual
Heisenberg degrees at several primes and
enlargements, including residue two.
\end{proof}

% A.65: root and DE full proof/source/certificate audits PASS.
\subsection{Every intermediate cover of an odd Heisenberg cover}

\begin{statement}{A.65}{Proved}{All intermediate exponent-prime class-two covers}
Let \(p\) be any odd prime, and let
\(Y_p\to U\) be the actual degree-\(p^3\)
Heisenberg cover of A.61--A.62. Every geometrically
connected intermediate cover of \(Y_p\to U\)
satisfies both assertions of Task A.5 at every
coefficient prime, for every cusp packet and every
permitted open derived subgroup.

More intrinsically, this holds for every finite cover
of \(U\) whose geometric Galois closure has exponent
\(p\) and nilpotency class at most two. The new
nonabelian intermediate type has degree \(p^2\),
proper genus \(1+p(p-1)/2\), and \(p\) cusps.
Its closure is the full Heisenberg group. This does
not settle all nilpotent covers or the universal
Task A.5.
\end{statement}

\begin{proof}
\paragraph{The actual subgroup alternatives.}
Write \(C=C_p\), \(q:C\to E\),
\(G=H_p\), and \(Z=Z(G)\). The actual map
to \(U\) is \([p]q\), and
\[
  1\longrightarrow Z\longrightarrow G
       \xrightarrow{\;\tau\;} E[p]\longrightarrow1,
  \qquad C/Z=E.                                      \tag{A.65.1}
\]
The commutator on \(E[p]\simeq\mathbf F_p^2\)
is a nondegenerate alternating form with values
in \(Z\simeq\mathbf F_p\). If a subgroup
\(S\le G\) contains \(Z\), its proper quotient
is \(E/(S/Z)\). The corresponding affine cover
has genus one and is covered by A.40.

If \(S\cap Z=1\), then \([S,S]=1\) and
the image \(\tau(S)\) is isotropic. In dimension
two an isotropic subspace has dimension at most
one. Thus \(S=1\), the case of A.61--A.62, or
\(S\) is a noncentral subgroup of order \(p\).
We treat the latter without replacing its quotient
by a pro-prime inclusion.

Every geometric point stabilizer for the actual
action of \(G\) on \(C\) is either trivial or
\(Z\). Indeed an element with nonzero translation
\(\tau(g)\) cannot fix a point over the elliptic
curve, and the cyclic degree-\(p\) map \(q\)
is totally ramified exactly above \(E[p]\).
Consequently \(S\) acts freely on the entire
generic proper curve. For
\[
       \pi:C\longrightarrow D=C/S
\]
we have an unramified degree-\(p\) map and
\[
 g(D)=1+\frac{p(p-1)}2,\qquad
 \#\{\text{cusps of }D\}=p.                         \tag{A.65.2}
\]
The genus follows from Riemann--Hurwitz and
\(g(C)=1+p^2(p-1)/2\); the cusp count follows
from the free \(S\)-action on the \(p^2\)
cusps. In particular there are at least three
distinct cusps.

\paragraph{A genuine elliptic quotient square.}
Put \(T=\tau(S)\), let
\(\phi:E\to A=E/T\) be the degree-\(p\)
isogeny, and let \(\widehat\phi:A\to E\)
be its dual. Actual equivariance of \(q\)
gives a degree-\(p\) map \(q_D:D\to A\)
with
\[
       q_D\pi=\phi q,\qquad
       \widehat\phi q_D\pi=[p]q.                    \tag{A.65.3}
\]
Thus the intermediate affine map to \(U\) is
\(\widehat\phi q_D\), of degree \(p^2\).
Its removed set is exactly
\(q_D^{-1}(\ker\widehat\phi)\). There is one
point above each of the \(p\) points of
\(\ker\widehat\phi\), with ramification
index \(p\). All these curves, groups, maps,
and marked points are defined over a suitable
finite number field. We may make subsequent
finite extensions without changing the geometric
subgroup or the selected packet.

\paragraph{Ordinary primes: choose a place with nonzero reduced translation.}
Suppose \(p\equiv1\pmod4\), and factor
\((p)=(\varpi)(\bar\varpi)\) in
\(\mathbf Z[i]\). The two lines
\(E[\varpi]\) and \(E[\bar\varpi]\)
are distinct in \(E[p]\). At a place above
\(\varpi\), the connected kernel of reduction
on \(E[p]\) is \(E[\varpi]\), while at a
place above \(\bar\varpi\) it is
\(E[\bar\varpi]\). To recall the exact
reason, the Gaussian endomorphisms extend to
the good elliptic model. The derivative of
\([\varpi]\) on an invariant differential is
\(\varpi\). Its degree is \(p\), so its
special map is purely inseparable at the first
place and etale at the second. The other factor
has the reverse behavior. These are the formal
CM identifications used in A.40, with
\cite[Theorem III.6.2 and Theorem V.3.1]{fghSilvermanAEC}.

Since \(T\) is one line, it cannot equal both
connected lines. Choose a place where it does
not equal the connected line. Then every
nonzero element of \(T\) has nonzero reduction.
This remains true after further finite extension.

Use the actual A.62 semistable model, with its
map to the good elliptic model and with all
original cusps separated. The \(G\)-action
extends to it; one may take the stable marked
model and then an equivariant resolution. For
each \(1\ne s\in S\), equivariance gives
\[
       \bar q(sP)=\bar q(P)+\overline{\tau(s)}
                         \ne\bar q(P)                \tag{A.65.4}
\]
at every geometric special point. Thus \(S\)
has no special stabilizer, including at nodes.
There is also no generic stabilizer. The finite
constant-group quotient is therefore an etale
\(S\)-torsor, even though \(|S|=p\): the
constant group scheme \(\mathbf Z/p\) is etale.
The quotient is a semistable proper model of
\(D\), by etale descent.

A.62 has \(p\) rational marked clusters above
the \(p\) distinct reduced points of \(E[p]\),
and each cluster contains \(p\) cusps. The
reduced translations of \(T\) permute these
clusters freely and transitively. Choose a lift
of any selected cusp of \(D\), and choose a
different cusp in its same rational cluster.
The two lifts are not \(S\)-equivalent, because
every nonidentity element changes the cluster.
Their rational joining path maps to a rational
joining path between distinct quotient cusps:
the quotient identifies the \(p\) copies of
this cluster and is an isomorphism on each one.
Thus A.42 gives the full fixed-inertia conclusion
at every coefficient prime.

\paragraph{Inert primes: compact type descends and no new positive component is hidden.}
Now suppose \(p\equiv3\pmod4\). The elliptic
curve \(E\) has good supersingular reduction.
After finite extension so does its isogenous
curve \(A\). The maps \(\phi\) and
\(\widehat\phi\) extend to the good group
models and retain degree \(p\): the degree
of the pullback of a relative degree-one
line bundle is constant. Their special maps
are purely inseparable, since a supersingular
elliptic curve has no nonzero geometric
\(p\)-torsion. Every point of
\(\ker\widehat\phi\) specializes to \(O\).

The norm on Jacobians is surjective, and
\(\pi_*\pi^*=[p]\). The potentially good
reduction of \(J_C\), supplied by A.61,
therefore implies potential good reduction
of \(J_D\), by
\cite[Theorem 1 and Corollary 2, p.493]{fghSerreTate1968}.
After choosing semistable reduction of \(D\),
its dual graph is consequently a tree. More
explicitly, the special identity component
of its Neron model is the special \(\operatorname{Pic}^0\),
and its torus rank is the graph cycle rank;
good Jacobian reduction forces this rank to
vanish. The exact comparison and rank formula
are recalled in \cite[Remark 10.4.24 and
Lemma 7.5.18(b)]{fghLiu}.

We justify the required actual model maps.
Start with regular semistable models and
resolve any rational map to a projective
target by point blowups. The relevant precise
result is \cite[Theorem 9.2.7, pp.396--397;
Proposition 9.2.5]{fghLiu}; every added
component is a projective line. In particular
the rational map from a regular model of
\(D\) to the good elliptic model \(A\)
extends without adding a positive component,
and in fact descends to the original model.
Here is the descent argument. At the last
blowup its exceptional projective line maps
constantly to the elliptic special fibre.
The image of the graph of the resolved map
is proper and quasi-finite over the original
normal surface: above the centre its fibre
is one point, and elsewhere it is one point.
It is therefore finite birational and hence
an isomorphism. Descending successively proves
the claimed extension of \(q_D\). This
argument avoids applying the Neron mapping
property to a source that is not smooth
over the valuation ring.

Next resolve the actual \(\pi:C\to D\)
by point blowups of the regular source model.
The resulting proper map is surjective: its
closed image contains the entire generic
fibre of the integral target model. Every
positive-genus component of the target
special fibre is therefore dominated by
one of the original positive-genus source
components. Indeed the added components are
rational, and a projective line cannot
dominate a positive-genus normalization,
by Luroth's theorem. This proves the
comparison for every positive target
component, including ones whose appearance
might otherwise be obscured by a wild quotient.
We have made no assertion that wild quotient
commutes with reduction.

A.61 has one elliptic main component whose
map to \(E\) is radicial. All other
positive-genus source components map
constantly to finite zeros of its function
\(R\), none of which is \(O\). Since the
special \(\phi\) is radicial, their images
in \(A\) are still different from \(O\).
The composite from the source main to
\(A\) is radicial and nonconstant. In
the resolved diagram (A.65.3) it factors
through precisely one target component,
which we call \(D_0\). Both factors are
radicial: their function fields are
intermediate fields of a purely inseparable
extension. In particular \(D_0\) has
exactly one point above \(O\).

No other positive target component maps
nonconstantly to \(A\), since the unique
source main cannot dominate two components
and all source tails have constant images.
Every other positive target component
maps constantly to a point different from
\(O\), by the domination argument. Also
\(D_0\) has genus one: it is dominated
by an elliptic curve and maps nonconstantly
to an elliptic curve.

\paragraph{All quotient cusps have rational joining paths.}
Every cusp of \(D\) specializes above
\(O\in A\). It cannot lie on any of the
other positive-genus components. A rational
component carrying such a mark has constant
map \(O\). Its path in the special-fibre
tree to \(D_0\) cannot pass through a
positive component with a different
constant image: equality of images at
each intervening node would be violated.
All such rational parts attach to the
unique point of \(D_0\) above \(O\).
At a nodal point only one other branch can
meet this component, so these parts form
one connected rational subtree.

If this subtree is absent before marking,
the \(p\ge3\) cusps all specialize to
the same smooth point of \(D_0\), and
their separation inserts a rational tree.
If it is present, its attachment on
\(D_0\) is a node; disjoint smooth marked
sections lie on the rational subtree or
on rational components inserted above
its nodes. In either case every selected
cusp has a distinct companion with a
fully rational joining path. A.41--A.42
give the desired full integral fixed
group for every coefficient prime and
every packet. Finite field restriction
inside a prescribed open derived subgroup
is exactly the one in A.41.

Every intermediate cover has nilpotent
geometric Galois closure, so A.59 gives
the normalizer assertion at all coefficient
primes. We have used no transfer of fixed
subgroups through a surjective pro-prime
quotient.

\paragraph{Intrinsic closure formulation.}
The geometric fundamental group of \(U\)
is free profinite on two generators. Its
maximal class-two exponent-\(p\) quotient
has the presentation
\[
 \langle x,y,z\mid z=[x,y],\ z\text{ central},
                      \ x^p=y^p=z^p=1\rangle.         \tag{A.65.5}
\]
Every word has the form \(x^a y^b z^c\),
with exponents modulo \(p\); the upper
unitriangular matrix model realizes
\(p^3\) distinct words. For odd \(p\),
that matrix group has exponent \(p\),
so (A.65.5) is exactly \(H_p\).
Every nonabelian quotient still has a
two-dimensional abelianization and
nontrivial one-dimensional commutator,
hence has order \(p^3\) and is \(H_p\).
The kernel of any surjection to \(H_p\)
is consequently the same characteristic
kernel of this universal quotient.

The actual A.61--A.62 cover realizes that
quotient. Thus every cover with geometric
closure as in the statement is geometrically
an intermediate of this one, or has abelian
closure and genus one. Any isomorphism and
finite descent data are defined after a
finite number-field extension. The preceding
all-packet, open-derived conclusions pass
to the original arithmetic setting exactly
as in A.41. This proves the stated scope
without a cofinality assertion.

The executed certificate
\path{code/a_heisenberg_intermediates.py}
checks the exact finite Heisenberg subgroup
classification, cusp orbits, and the two
ordinary connected-line choices. It is
support for the uniform proof, not a
substitute for the model comparison.
\end{proof}

% A.66: full root and DE independent proof/source/certificate audits PASS.
\subsection{An abelian factor of order prime to the Heisenberg prime}

\begin{statement}{A.66}{Proved}{Heisenberg times prime-to-prime abelian closure}
Let \(p\) be an odd prime and let \(B\) be a finite abelian group
of order prime to \(p\). Every geometrically connected finite
cover of \(U=E\setminus\{O\}\), \(E:y^2=x^3-x\), whose geometric
Galois closure is a quotient of
\[
                         H_p\times B
\]
satisfies both assertions of Task A.5, at every coefficient
prime, every original cusp packet, and every permitted open
derived subgroup. Here \(H_p\) is the exponent-\(p\), order-\(p^3\)
Heisenberg group.

More concretely, take any intermediate \(D^\circ\to U\) of
A.65 and any elliptic isogeny \(\psi:E'\to E\) of degree \(h\)
prime to \(p\). The actual fibre product
\[
 Y=D^\circ\times_U
       \bigl(E'\setminus\psi^{-1}(O)\bigr)
                                                               \tag{A.66.1}
\]
is geometrically connected and has both assertions. Its proper
compactification \(D'\to D\) is finite étale of degree \(h\), and
\[
 \deg(Y/U)=h\deg(D^\circ/U),\quad
 g(D')=1+h(g(D)-1),\quad
 \#(D'\setminus Y)=h\#(D\setminus D^\circ).
                                                               \tag{A.66.2}
\]
In particular the noncentral order-\(p\) quotient gives degree
\(hp^2\), genus \(1+hp(p-1)/2\), and \(hp\) cusps; the full
Heisenberg case gives degree \(hp^3\), genus
\(1+hp^2(p-1)/2\), and \(hp^2\) cusps.

This is a fixed-subgroup theorem for actual composite covers,
not a transfer through a surjective pro-prime quotient. It does
not treat products of two nonabelian Sylow factors or all finite
nilpotent closures. The universal Task A.5 remains open.
\end{statement}

\begin{proof}
All geometric covers, groups, isogenies and marked points can
first be defined over a finite number-field extension. The
comparison with every original packet and every permitted open
derived subgroup is the finite-extension comparison of
A.41--A.42. We establish the actual pointed stable-model
hypotheses before using it.

\paragraph{Coprime factors and the actual function fields.}
If \(S\le H_p\times B\), then
\[
                         S=S_p\times S_B.
                                                               \tag{A.66.3}
\]
Indeed, for \((a,b)\in S\), choose an integer that is one
modulo the order of \(a\) and zero modulo the order of \(b\);
raising \((a,b)\) to that integer isolates \((a,1)\in S\).
The reverse choice isolates \((1,b)\). This proves (A.66.3)
without a normality assumption on \(S\).

An abelian geometric quotient of the rank-two fundamental
group kills its cusp commutator, so its covering extends
étale over \(E\). After making a point above \(O\) rational
and using it as origin, this proper cover is an elliptic
isogeny. Its deck group is its kernel. The intrinsic
Heisenberg comparison in A.65 identifies the \(H_p\)-quotient
with the actual cover of A.61--A.62.

The two Galois function fields have coprime degrees. Their
intersection over the base function field is trivial and,
since they are Galois, they are linearly disjoint. Their
compositum therefore has Galois group \(H_p\times B\).
Taking fixed fields under (A.66.3) identifies each intermediate
with the actual fibre product (A.66.1), with the abelian
isogeny replaced by its quotient under \(S_B\).
This argument proves geometric connectedness as well as the
identification of the covering map. It is not an identification
made from group orders alone.

The same reasoning applies if the closure is only a quotient
of \(H_p\times B\): its normal kernel splits into the two
coprime factors. A proper nonabelian quotient of \(H_p\) does
not exist, by A.65; the remaining case is abelian and hence
has genus-one compactification. That case is already A.40.

The projection of (A.66.1) to \(D^\circ\) extends to
\[
                         D'=D\times_E E'\longrightarrow D.
                                                               \tag{A.66.4}
\]
It is finite étale because \(\psi\) is. Thus \(D'\) is smooth
proper and geometrically connected, and there are exactly
\(h\) points above every old cusp. Riemann--Hurwitz proves
all formulas (A.66.2). When \(D\) has genus one this again
falls under A.40. It remains to handle \(S_p=1\) and a
noncentral order-\(p\) subgroup.

\paragraph{The one nonconstant component in the actual models.}
In these two cases A.61--A.65 supply a factorization
\[
       D\xrightarrow{q_D} A\xrightarrow{\lambda}E,
       \qquad f_D=\lambda q_D ,
                                                               \tag{A.66.5}
\]
where \(A\) is an elliptic curve isogenous to \(E\) and
\(\lambda\) has \(p\)-power degree. For the full Heisenberg
cover, \(A=E,\ q_D=q_p,\ \lambda=[p]\). For a noncentral
subgroup, \(T\subset E[p]\) is its translation image,
\(A=E/T\), and \(\lambda=\widehat\phi\), with
\(\phi:E\to E/T\) as in (A.65.3).

Choose the place over \(p\) used in A.61 or A.62. In the
ordinary noncentral case choose the Gaussian place where
\(T\) has nonzero reduction, exactly as in A.65. After a
finite extension, \(A\) has good reduction and the pointed
semistable model \(\mathcal D\) has the following properties:

\begin{enumerate}
\item its dual graph is a tree;
\item the actual map \(q_D\) extends to the good model
\(\mathcal A\);
\item there is exactly one component \(M\) on which the
special map is nonconstant, and
\(M\to\overline A\) is finite radicial;
\item every other component maps constantly to
\(\overline A\);
\item every cusp has a distinct cusp joined to it by a path
consisting entirely of rational components, including the
components containing the two marks.
\end{enumerate}

For the full cover these are the explicit main, tail and
marking conclusions of A.61--A.62. For an inert-prime
noncentral quotient the actual resolution and domination
argument in A.65 proves (ii)--(iv), and its rational fibre
over the origin proves (v). Compact type gives (i).

At the ordinary noncentral place, property (i) follows
as well from the actual surjective norm \(J_C\to J_D\):
the source has good reduction by A.62, so the target
does too by the same Serre--Tate Corollary 2. The
semistable target graph therefore has cycle rank zero.

Here is the additional verification of (iii) for the
ordinary noncentral quotient. In A.65 the order-\(p\)
group acts freely on the entire semistable model because
its translation reduction is nonzero. It preserves the
unique nonconstant component of the source, which is an
elliptic curve \(M_C\) with radicial degree-\(p\) map
\(M_C\to\overline E\). Its quotient
\(M_C\to M_C/S_p\) is finite étale of degree \(p\).
The special isogeny
\(\overline\phi:\overline E\to\overline A\) is also
étale of degree \(p\). The commuting square
\[
 \begin{array}{ccc}
 M_C&\longrightarrow&\overline E\\
 \downarrow&&\downarrow\overline\phi\\
 M_C/S_p&\longrightarrow&\overline A
 \end{array}
                                                               \tag{A.66.6}
\]
therefore makes the bottom map degree \(p\) with separable
degree one: the composite has total degree \(p^2\) and
separable degree \(p\), and the left map is étale of degree
\(p\). Hence the bottom map is radicial. Every other source
component is constant over \(\overline A\), so no other
nonconstant target component appears. The quotient model
and its marking comparison are exactly those of A.65.
This proves all five properties in this case too.

If point blowups are used to separate marks or resolve a
map, every new component is rational and maps constantly
to the elliptic special fibre. Thus the same properties
hold on one common pointed semistable model.

\paragraph{An actual connected elliptic base change.}
Form the proper fibre product
\[
                         A'=A\times_E E',
                                                               \tag{A.66.7}
\]
using \(\lambda\) and \(\psi\). The two function fields
have coprime degrees, so \(A'\) is geometrically connected.
The projection \(A'\to A\) is finite étale of degree \(h\).
The point \((O,O)\) is an origin, making \(A'\) an elliptic
curve and this projection an isogeny.
Associativity of the actual fibre products gives
\[
                         D'=D\times_A A'.
                                                               \tag{A.66.8}
\]

At the chosen place, \(A'\) has good reduction because it is
isogenous to the good elliptic curve \(A\);
\cite[Theorem 1 and Corollary 2, p.~493]{fghSerreTate1968}.
The isogeny and its dual extend between the good abelian
schemes. Their composite is multiplication by \(h\), which
is étale since \(p\nmid h\). Invariant differentials show
that both extended maps have invertible differential.
Their kernels are finite, being contained in the
finite étale \(h\)-torsion. Properness makes the maps
finite, and they are étale. Their degree is the generic
degree \(h\), by flatness. Thus the actual projection
extends as a finite étale isogeny
\[
                    \mathcal A'\longrightarrow\mathcal A .
                                                               \tag{A.66.9}
\]
This is a comparison of good elliptic models, not a
prime-to-prime assertion about an arbitrary singular source.

\paragraph{The pulled-back graph and the marked paths.}
Take the actual model
\[
              \mathcal D'=\mathcal D
                              \times_{\mathcal A}\mathcal A'.
                                                               \tag{A.66.10}
\]
It is proper and semistable because it is finite étale
over \(\mathcal D\). Its generic fibre is the connected
curve (A.66.8). We work over an algebraically closed
residue field to describe its components.

The component over \(M\) is connected. Indeed
\[
                   M'=M\times_{\overline A}\overline A'
\]
maps radicially and surjectively to the connected elliptic
curve \(\overline A'\). A finite radicial surjection is a
universal homeomorphism, and this statement is preserved
by base change. The map \(M'\to M\) is étale; hence
\(M'\) is smooth. There is exactly one component over \(M\).

Every other component of \(\mathcal D\) maps to a single
point of \(\overline A\), so its pullback is a disjoint
union of \(h\) identical components. Removing the vertex
of \(M\) from the tree graph of \(\mathcal D\) leaves
subtrees, each attached to \(M\) in one node. The map to
\(\overline A\) is constant on each whole subtree: adjacent
constant components agree at their common node.
Consequently its pullback consists of \(h\) disjoint
copies, attached at the \(h\) points above that original
node on \(M'\). It follows that the whole special graph
of \(\mathcal D'\) is a tree. In particular the new
Jacobian has good reduction; no cycle is created by the
prime-to-\(p\) base change.

Now fix any marked cusp of \(\mathcal D'\). The chosen
old cusp lies in a rational joining path with a distinct
old cusp, by property (v). The elliptic map is constant
on this entire path. Its pullback is therefore \(h\)
disjoint identical copies. The selected new cusp belongs
to one copy, and the other endpoint of that copy is a
distinct marked cusp. Both endpoints and the entire
joining path lie on rational components. This argument
works for every lift of every cusp, not just for one
chosen component.

Equivalently, any subsequent stabilization contracts only
rational components, and marked sections can be separated
on the same rational trees. The actual pointed criterion
of A.41--A.42 is therefore satisfied by \(\mathcal D'\).
It proves the full fixed-subgroup equality at every
coefficient prime, including the residue prime \(p\),
and for every permitted packet and open derived subgroup.

\paragraph{Normalizers and scope.}
The geometric closure of (A.66.1) is a subgroup quotient
of \(H_p\times B\), hence is nilpotent. A.59 supplies the
normalizer assertion at every coefficient prime. The
fixed-subgroup argument above used the actual model,
not this nilpotence alone. In particular it supplies no
claim about a product of two nonabelian prime factors,
where the additional cover is not a proper elliptic
isogeny and the constant-component pullback argument
does not apply.
The exact certificate
\path{code/a_heisenberg_abelian_products.py}
exhausts the subgroups in three actual coprime product groups
and checks the degree, genus, cusp and copied-tree identities.
The actual model comparison is the proof above.
\end{proof}

% A.67: complete root and DE independent proof/source/certificate audits PASS.
\subsection{The tetrahedral geometric closures}

\begin{statement}{A.67}{Proved}{All intermediates of every tetrahedral closure}
Every geometrically connected finite cover of
\(U=E\setminus\{O\}\), \(E:y^2=x^3-x\), whose geometric Galois
closure is a quotient of \(A_4\), satisfies both assertions of
Task A.5 at every coefficient prime, every original cusp packet,
and every permitted open derived subgroup.

For an actual \(A_4\)-Galois cover \(X^\circ\to U\), the proper
curve \(X\) has genus four. Its intermediate covers have the
following invariants; subgroups in the same row are conjugate.
\[
\begin{array}{c|r|r|r}
 S\le A_4&[A_4:S]&g(X/S)&\#((X/S)\setminus(X^\circ/S))\\ \hline
 1&12&4&6\\
 C_2&6&2&4\\
 C_3&4&2&2\\
 V_4&3&1&3\\
 A_4&1&1&1
\end{array}
\]
The proof identifies every such closure with an actual cover
constructed from a degree-three isogeny. At residue characteristic
three its pointed model has a rational component containing all
its cusps. The same property is proved for its genus-two
intermediates, including the wild \(C_3\) quotient. This is a
non-nilpotent family; the universal Task A.5 remains open.
\end{statement}

\begin{proof}
We may make a finite number-field extension to define all covers,
deck transformations, isogenies and marked points. At the end we
use the packet and open-derived-subgroup comparison in A.41--A.42,
so this preliminary extension does not restrict the asserted
original packets. We first identify the actual curves and maps.

\paragraph{The isogeny and the functions defining the cover.}
Let \(\phi:E'\to E\) be a degree-three isogeny, with kernel
\(T=\{O,P,-P\}\). A plane embedding by \(|3O|\), followed by
translation of coordinates putting \(P=(0,0)\) and its tangent
at \(y=0\), gives
\[
 E':\quad y^2+a\,xy+b\,y=x^3,\qquad P=(0,0),
 \qquad b(a^3-27b)\ne0.                                      \tag{A.67.1}
\]
Indeed, a point of order three is an inflection point. In the
general Weierstrass equation, passing through \((0,0)\) first
kills \(a_6\), the tangent condition kills \(a_4\), and triple
intersection with \(y=0\) kills \(a_2\). Smoothness at \(P\)
forces \(a_3=b\ne0\). Thus (A.67.1) is a coordinate construction,
not an assumed classification formula.

For a variable point \(Q=(x,y)\), the chord-and-tangent law gives
\[
\begin{aligned}
 f_0(Q)&=x,&
 f_1(Q)=x(Q+P)&=-\frac{by}{x^2},\\
 f_2(Q)=x(Q-P)&=\frac{b(y+b+ax)}{x^2},&
 f_0f_1f_2&=-b^2.
\end{aligned}                                                 \tag{A.67.2}
\]
Their divisors are respectively
\[
 P+(-P)-2O,\qquad O+P-2(-P),\qquad O+(-P)-2P.
                                                               \tag{A.67.3}
\]
The two independent odd-divisor vectors in (A.67.3) prove that
adjoining two of the three square roots gives a connected
\(V_4\)-cover. Define \(X\) as the smooth proper normalization
of this function field, choosing roots with
\[
 u_j^2=f_j,\qquad u_0u_1u_2=ib.
                                                               \tag{A.67.4}
\]
Translation by \(P\) cyclically permutes the \(f_j\), so cyclically
permuting the \(u_j\) lifts it to an automorphism \(\tau\) of order
three. The even sign changes of the roots form \(V_4\), and
\[
 V_4\rtimes\langle\tau\rangle=A_4
                                                               \tag{A.67.5}
\]
is the deck group of the actual composite \(X\to E'\to E\).
The only branch points of \(X\to E'\) are \(T\); at each the
inertia has order two, and the three inertia subgroups are
distinct. Thus \(X\to E\) is unramified away from \(O\).
Riemann--Hurwitz gives \(2g(X)-2=3\cdot2=6\), and there are six
points over \(O\). Every nontrivial generic stabilizer is a
subgroup of \(V_4\) of order two.

\paragraph{Why this includes every \(A_4\) closure.}
The cusp of the rank-two geometric fundamental group is a
commutator. In a surjection to \(A_4\) its image is a nonzero
element of \(V_4\): otherwise the two generating images would
commute. The \(C_3=A_4/V_4\) quotient is abelian, kills the cusp,
and therefore extends to a degree-three elliptic isogeny
\(E'\to E\), after choosing an origin. Over \(E'\) its three
cusp inertia images are the three distinct nonzero elements
of \(V_4\), cycled by translation by \(P\).

Here is uniqueness of the equivariant quadratic cover, including
possible unramified twists. Over the geometric field the exact
residue sequence is
\[
 0\longrightarrow H^1(E',\mathbf F_2)
 \longrightarrow H^1(E'\setminus T,\mathbf F_2)
 \longrightarrow \mathbf F_2[T]^0\longrightarrow0.
                                                               \tag{A.67.6}
\]
Translation acts trivially on the first term, of dimension two.
The last term, the sum-zero subspace on the three points, is the
irreducible two-dimensional \(C_3\)-module over \(\mathbf F_2\).
Since \(3\) is invertible in \(\mathbf F_2\), the sequence splits
as \(C_3\)-modules. Consequently its unique nontrivial
two-dimensional isotypic summand is the image of the dual of
any equivariant \(V_4\)-quotient with the specified inertia.
Changing its basis only changes the labels of the deck group.
In particular the trivial unramified summand cannot supply an
additional equivariant quadratic twist. The lift of \(C_3\)
is unique up to \(V_4\)-conjugacy, since
\(H^1(C_3,V_4)=0\), by averaging. Equations (A.67.2)--(A.67.5)
therefore give every geometric \(A_4\) closure, as an actual
cover of \(E\), after a finite field extension.

\paragraph{An exact quotient equation, rather than an assumed \(j\)-formula.}
Summing the translated coordinates and subtracting the kernel
coordinates gives the invariant functions
\[
\begin{split}
 X_0&=x+\frac{ab}{x}+\frac{b^2}{x^2},\\
 Y_0&=\frac{(x^3-abx-2b^2)y-b(ax+b)^2}{x^3}.
\end{split}                                                   \tag{A.67.7}
\]
For example \(X_0=x(Q)+x(Q+P)+x(Q-P)\), since the two finite
kernel \(x\)-coordinates are zero. The addition law gives
\[
 y(Q+P)=\frac{by(y+ax)}{x^3}-b,\qquad
 y(Q-P)=-\frac{b(y+b+ax)^2}{x^3},
\]
and \(Y_0=y(Q)+y(Q+P)+y(Q-P)+b\), proving its invariance too.
Substitution in (A.67.1), without division by a numerical
specialization, verifies the equation
\[
 Y_0^2+aX_0Y_0+bY_0
       =X_0^3-5abX_0-(a^3b+7b^2).
                                                               \tag{A.67.8}
\]
The function \(X_0\) has a double pole at each of the three
kernel points and no others. Its degree is six. The smooth
elliptic equation (A.67.8) has \(X_0\)-degree two, so the map
has degree three and is exactly the quotient by \(T\).

The usual Weierstrass invariant polynomials applied to
(A.67.8) give
\[
 c_4=a(a^3+216b),\qquad
 \Delta=b(a^3-27b)^3.
                                                               \tag{A.67.9}
\]
Set \(r=a^3/b\). The quotient is isomorphic over a finite
extension to the prescribed \(E\), if and only if
\[
 \frac{r(r+216)^3}{(r-27)^3}=1728.
\]
The exact polynomial identity
\[
 r(r+216)^3-1728(r-27)^3=(r^2-540r-5832)^2
                                                               \tag{A.67.10}
\]
therefore yields
\[
                         r=270\pm162\sqrt3.
                                                               \tag{A.67.11}
\]
In particular \(a\ne0\). These computations use the actual
isogeny, not merely a remembered modular correspondence.

\paragraph{The complementary elliptic factor and residue two.}
One of the three double quotients of \(X\to E'\) is obtained by
putting \(z^2=x\) and \(Y=2y+az^2+b\):
\[
 D_0:\quad Y^2=4z^6+a^2z^4+2abz^2+b^2.
                                                               \tag{A.67.12}
\]
Its involution \((z,Y)\mapsto(-z,Y)\) has quotient \(E'\).
The involution \((z,Y)\mapsto(-z,-Y)\) has the actual quotient
\[
 t=z^{-2},\quad w=Y/z^3,\qquad
 F:\ w^2=4+a^2t+2ab t^2+b^2t^3.
                                                               \tag{A.67.13}
\]
After \(X_F=b^2t,\ Y_F=b^2w\), this is
\[
 Y_F^2=X_F^3+2abX_F^2+a^2b^2X_F+4b^4.
                                                               \tag{A.67.14}
\]
The two elliptic differentials pull back to independent
multiples of \(z\,dz/Y\) and \(dz/Y\). Hence
\(J_{D_0}\sim E'\times F\). The three nontrivial \(V_4\)-character
spaces in \(H^1(X)\) are cyclically permuted by the actual
automorphism \(\tau\); each is the complement to \(H^1(E')\)
in one of these double quotients. Consequently
\[
                         J_X\sim E'\times F^3.
                                                               \tag{A.67.15}
\]
This identifies the factor \(F\) using actual maps and
differentials.

For (A.67.14), direct invariant calculation gives
\[
 c_4(F)=16a^2b^2,\quad
 \Delta(F)=256b^7(a^3-27b),\quad
 j(F)=\frac{16r^2}{r-27}=4608\pm2688\sqrt3.
                                                               \tag{A.67.16}
\]
At every place above two, \(j(F)=128(36\pm21\sqrt3)\) has
positive valuation, since the factor in parentheses is a unit.
Thus \(F\) has potentially good reduction and its good special
fibre has \(j=0\), by Proposition VII.5.5, including its
characteristic-two clause in Appendix A.1.4(b)
\cite[pp.~197, 199]{fghSilvermanAEC}. An elliptic curve with
\(j=0\) in characteristic two is supersingular: over the
algebraic closure it is \(y^2+y=x^3\), whose inversion
\((x,y)\mapsto(x,y+1)\) has no affine fixed point. Hence it has
no nonzero geometric point of order two.

The curve \(E'\), being isogenous to \(E\), also has potentially
good supersingular reduction at two. Good reduction is
preserved under isogeny and quotient
\cite[Theorem 1 and Corollary 2, p.~493]{fghSerreTate1968}.
The isogeny (A.67.15) therefore proves that \(J_X\) has potentially
good reduction of two-rank zero. Every intermediate Jacobian
is an actual quotient of \(J_X\), by the norm map, and has the
same two properties. Statement A.46 supplies the pro-\(2\)
normalizer assertion for each of these curves.

\paragraph{A three-adic marked base model.}
At a place above three, make a finite extension so that \(E'\)
has good supersingular reduction. This is possible because it
is isogenous to \(E\); supersingularity is preserved by the
specialized isogeny. The three points \(O,P,-P\) all reduce to
\(O\). In a good formal parameter \(t\), put \(c=t(P)\). Then
\[
 v(c)>0,\qquad t(-P)=-c+O(c^2).
                                                               \tag{A.67.17}
\]
In the cluster chart \(t=c\,u\), the three marks have distinct
reductions \(0,1,-1\). The integral formal addition law shows
that translation by \(P\) reduces to \(u\mapsto u+1\):
all terms of total degree at least two become divisible by
\(c\) after division by \(c\). The marked base thus consists of
the good elliptic component and a rational component with
marks \(0,1,-1\), attached at \(u=\infty\); regular resolution
may insert rational chains. No assertion about a common
radius of all three-torsion points is needed.

Each function \(f_j\) has degree-zero horizontal divisor on
the rational component, no horizontal divisor on the main
component, and no other horizontal zeros or poles. The
intersection equations for its principal divisor imply that
its vertical orders are constant along this tree. Explicitly,
the graph Laplacian of the vertical-order vector is zero,
because the horizontal degree on every component is zero;
the kernel of the Laplacian of a connected finite tree consists
of constants. After multiplying by a constant and extending
the field to extract its square root, the vertical orders
are therefore zero. On the main elliptic component each
reduced \(f_j\) is a nonzero constant. On the rational
component their reduced divisors, and hence their functions
up to nonzero constants, are
\[
 \frac{u^2-1}{u^2},\qquad
 \frac{u(u-1)}{(u+1)^2},\qquad
 \frac{u(u+1)}{(u-1)^2}.
                                                               \tag{A.67.18}
\]
All constants may be removed after a finite extension.

Since two is invertible, Hensel's lemma splits the normalization
over the elliptic component into four copies of the good
elliptic curve. At the attaching node all three functions are
units and have zero slopes, so the same quadratic covers are
étale and split there. Over the rational component, (A.67.18)
defines a connected tame \(V_4\)-cover branched at \(0,1,-1\),
with three distinct inertia groups of order two. Its genus is
zero by Riemann--Hurwitz, and there are four unramified points
over infinity. Thus normalization of the actual global
function field, followed by regular resolution, gives
\[
 \begin{gathered}
 \text{one rational centre, four supersingular elliptic tails,}\\
 \text{and six distinct cusp sections on the centre.}
 \end{gathered}                                               \tag{A.67.19}
\]
The only possible inserted components are rational chains.
These local normalizations agree on the node collars because
they are restrictions of the same Kummer field (A.67.4).
The four tail genera already sum to \(g(X)=4\); the graph is
a tree. This also proves potentially good reduction of \(J_X\)
with three-rank zero.

\paragraph{The exact residual centre and its automorphisms.}
The preceding rational cover can be written over the algebraic
closure of \(\mathbf F_3\) as
\[
             u=\frac{(v^2+1)^2}{v^3-v}.
                                                               \tag{A.67.20}
\]
Direct substitution gives
\[
 u(v+1)=u(v)+1,\qquad u(-1/v)=u(v).
                                                               \tag{A.67.21}
\]
The conjugates of \(v\mapsto-1/v\) generate \(V_4\);
\(v\mapsto v+1\) gives the order-three lift. Any other such
lift is conjugate to it by \(V_4\), as already proved by
\(H^1(C_3,V_4)=0\). The four nodes are
\(\mathbf P^1(\mathbf F_3)\). The derivative numerator is
\[
 v^6+v^4+v^2+1=\frac{v^9-v}{v^3-v},
                                                               \tag{A.67.22}
\]
so the six cusps are exactly \(\mathbf F_9\setminus\mathbf F_3\).
They have stabilizer of order two in
\(\operatorname{PSL}_2(\mathbf F_3)=A_4\).

In particular a \(C_3\) fixes one of the four nodes and cycles
the other three. It has no fixed cusp. Its quotient coordinate
\[
                              w=v^3-v
                                                               \tag{A.67.23}
\]
sends the two cusp orbits to \(i,-i\), where \(i^2=-1\), and
the two node orbits to \(0,\infty\). A subgroup \(C_2\le V_4\)
fixes two cusps and no node. Its quotient centre has four
cusp images and two node images. All of these images are
distinct as indicated.

\paragraph{The tame genus-two intermediates.}
For \(S=C_2\), quotienting the actual three-adic model is tame.
The four tails form two free \(C_2\)-orbits, giving two elliptic
tails. The rational centre has the four distinct marks and
two nodes just computed. Local quotient singularities at
nodes are resolved by rational chains. Thus every cusp lies
on this rational marked centre, and joining two cusps does
not pass through a positive-genus component. Its genus is
two, also directly from Riemann--Hurwitz on \(X\to X/C_2\).

\paragraph{The wild genus-two intermediate: residue fields and resolution.}
Put \(D=X/C_3\). The action is free on the generic proper curve,
since the only nontrivial generic stabilizers have order two.
Hence \(g(D)=2\), and its two cusps are the two \(C_3\)-orbits
already computed. The norm \(J_X\to J_D\) is surjective. Since
\(J_X\) has good reduction after the chosen extension, so does
\(J_D\), by the good-reduction criterion cited above. Therefore
any semistable model of \(D\) has a tree as its dual graph.

Take the finite normal quotient of the actual model of \(X\)
by \(C_3\), and then resolve and semistabilize, after a further
finite extension if necessary. We do not assume that forming
this quotient commutes with reduction. Instead, consider the
component valuations and their residue fields. Distinct
orbits of components give distinct valuations of the quotient
field. On the rational centre the residual action is
\(v\mapsto v+1\). The quotient residue field is contained in
\(k(v)^{C_3}=k(v^3-v)\). If \(A\) is an integral lift of \(v\)
at this valuation, the coefficients of
\[
                      \prod_{g\in C_3}(T-gA)
\]
belong to the quotient field and are integral. Their
reductions give \(T^3-T-(v^3-v)\), so \(v^3-v\) belongs to the
quotient residue field. The centre image is consequently
the rational component with coordinate (A.67.23), not an
unspecified reduction of a quotient.

The orbit of three elliptic tails has trivial decomposition
group and gives one elliptic component isomorphic to a tail.
For the fixed tail, the map \(X\to E'\) is an isomorphism on
that component. Translation by \(P\) on the good model of
\(E'\) reduces to the identity because \(P\) reduces to \(O\).
Thus the residual action on the fixed elliptic tail is
trivial. For any residue function \(a\) on that tail, lift it
to an integral element \(A\). The invariant polynomial above
now reduces to \((T-a)^3=T^3-a^3\), proving that the quotient
residue field contains all cubes. It lies between \(k(C)^3\)
and \(k(C)\), where \(C\) is the elliptic tail. Since \(k\)
is perfect and \(k(C)\) has transcendence degree one, this
extension is purely inseparable of degree at most three.
The normalization of the image is therefore an elliptic
curve, either \(C\) or its Frobenius twist.

The two tail orbits have thus produced two distinct
positive-genus components, each of genus one. Their genera
exhaust \(g(D)=2\). The graph is a tree by good reduction of
\(J_D\); hence all further components in resolution and
semistabilization are rational. This genus argument is what
excludes extra positive-genus exceptional components of the
wild quotient.

Finally the two cusp orbits are free already on the residual
centre. Near each such orbit the quotient is the quotient of
three disjoint smooth neighbourhoods permuted freely, so the
marked images remain smooth and distinct. They are at
\(w=i,-i\), whereas the two node orbits are at \(0,\infty\).
The two nodes lead to the two elliptic image components,
possibly through rational chains. Thus the rational centre
has two marks and two distinct outgoing branches and survives
marked stabilization. Both cusps lie on it, and their joining
path is rational. This proves the exact pointed hypothesis
needed for \(D\), including the wild order-three comparison.
The resolution and stabilization operations used here are
the usual point blow-ups and contraction of unstable rational
components; see \cite[9.2.5, 9.2.7, 10.4.24]{fghLiu}.

\paragraph{Fixed subgroups and the remaining normalizer primes.}
For \(X\) and both types of genus-two intermediate, the
three-adic constructions give a compact-type pointed model
in which every cusp has another cusp joined to it entirely
through rational components. Apply A.41--A.42 to conclude
that every permitted fixed subgroup, at every coefficient
prime, is exactly the chosen cusp inertia. The two
genus-one rows are A.40.

For normalizers at coefficient prime three, all intermediate
Jacobians have potentially good reduction of three-rank zero:
this follows from (A.67.19) and the actual norm quotients.
Apply A.46. The coefficient-two assertion was proved using
(A.67.15)--(A.67.16) and A.46. At any coefficient prime
\(\ell\ge5\), the geometric Galois closure has order dividing
twelve, prime to \(\ell\). The good one-marked elliptic model
of \(E\), after a finite extension, admits the actual
prime-to-\(\ell\) tame extension of this cover. The
normalization is a smooth proper curve with the prescribed
marked branch fibres; equivalently this is the specialization
equivalence of SGA~1, Expos\'e XIII, 2.9--2.12
\cite{fghSGA1}, as in A.59. Every intermediate proper curve
therefore has potentially good curve reduction at \(\ell\).
Apply A.57 for genus at least two and A.40 in genus one.
This proves the normalizer assertion for all primes.

The subgroup table follows directly from the subgroup types
of \(A_4\), the six-point cusp action, and Riemann--Hurwitz;
the exact identities \(2g-2+r=\deg(Y/U)\) provide an
additional check. Proper quotients of \(A_4\) are trivial or
\(C_3\), whose covers are elliptic isogenies and are A.40.
The finite-extension comparison of A.41--A.42 now returns
both assertions to every original packet and permitted
open derived subgroup, completing the statement.

The executed exact certificate is
\path{code/a_tetrahedral_covers.py}. It verifies the translation
and three-isogeny identities modulo the actual cubic equation,
the complementary elliptic invariants, both values of \(r\),
the residual rational map and all \(A_4\) orbits on its four
nodes and six cusps. The geometric arguments above, including
the wild quotient and the passage to all covers, are separate
from these finite algebraic checks.
\end{proof}

% A.68: full root and DE proof/source/exact-certificate audits PASS, including every intermediate.
\subsection{The two nonstable dihedral genus-seven covers}

\begin{statement}{A.68}{Proved}{The nonstable genus-seven dihedral covers and their intermediates}
Let
\[
 C_{\pm}:v^2=w^{16}\pm6w^8+1,
 \qquad S=\{P_{0,+},P_{0,-},P_{\infty,+},P_{\infty,-}\},
\]
where the signs distinguish the values of \(v\) at zero and of
\(v/w^8\) at infinity. Equip these punctured curves with the actual
degree-sixteen covers of \(E\setminus\{O\}\) described below.
For every coefficient prime, every induced cusp packet, and every
open derived Galois subgroup preserving the selected cover and cusp,
the fixed geometric pro-\(\ell\) subgroup is exactly the selected
cusp inertia. The same conclusion holds for every actual intermediate
cover of either degree-sixteen dihedral closure. Their procyclic normalizer assertions also hold, by
the nilpotent-closure theorem A.59.

The new geometric input is that their proper Jacobians have potential
good reduction of \(2\)-rank zero at two and that every cusp difference
is annihilated by eight. This is a result for these actual covers;
it makes no claim about arbitrary higher dihedral covers.
\end{statement}

\begin{proof}
We first work with the minus-six member \(C\). The substitution
\(w\mapsto\zeta_{16}w\), after a finite field extension, identifies
the two members and their four cusp sets.

\paragraph{The actual cover and its quotient maps.}
The degree-sixteen map to the original elliptic curve is
\[
 X=\frac{(w^8-1)^2}{4w^8},\qquad
 Y=\frac{v(w^{16}-1)}{8w^{12}},\qquad Y^2=X^3-X.
                                                               \tag{A.68.1}
\]
It also factors as elliptic duplication after the actual degree-four map
\[
 x=\frac{w^8-1+v}{2},\qquad y=w^4(x-1),
 \qquad y^2=x^3-x,\qquad w^4=\frac{y}{x-1}.           \tag{A.68.2}
\]
Substitution and the elliptic duplication formulas verify
\([2](x,y)=(X,Y)\). The divisor of \(y/(x-1)\) on the elliptic
curve is
\[
                 T_0+T_{-1}-T_1-O.
\]
Consequently (A.68.2) is totally ramified of index four at precisely
the four cusps. Duplication is unramified, so (A.68.1) is finite
\'{e}tale away from these cusps. Its deck group is the actual dihedral
group of order sixteen generated by
\[
 r(w,v)=(\zeta_8w,-v),\qquad s(w,v)=(1/w,-v/w^8).
\]
Indeed both preserve (A.68.1), satisfy \(r^8=s^2=1\) and
\(srs=r^{-1}\), and give sixteen distinct automorphisms.
The hyperelliptic equation has sixteen distinct roots, so \(g(C)=7\).

Besides the original deck action, put
\(\mu(w,v)=(\zeta_8w,v)\). The quotient by \(\mu^4\) is the
actual genus-three curve
\[
 C_4:\ v^2=u^8-6u^4+1,\qquad u=w^2.                 \tag{A.68.3}
\]
There is also an actual map to the genus-two curve
\[
 \begin{split}
 B:&\quad W^2=(T+2)(T^4-4T^2-4),\\
 &T=w^2+w^{-2},\qquad W=(w+w^{-1})v/w^4.
 \end{split}                                                   \tag{A.68.4}
\]
The polynomial on the right has discriminant \(-2^{20}\),
so its smooth proper model has genus two. To see the quotient
construction explicitly, first quotient by
\(h(w,v)=(1/w,v/w^8)\), obtaining
\[
 t=w+w^{-1},\quad z=v/w^4,\qquad
 z^2=t^8-8t^6+20t^4-16t^2-4.
\]
Then take invariants \(u=t^2,W=tz\) under \((t,z)\mapsto(-t,-z)\)
and set \(T=u-2\). These substitutions prove (A.68.4), including
the function-field identification rather than merely its branch data.

\paragraph{The full Jacobian decomposition comes from those maps.}
Use the basis
\(\omega_j=w^{j-1}dw/v\), \(1\le j\le7\), of
\(H^0(C,\Omega^1)\). The pullbacks from \(C_4\) span
\(\omega_2,\omega_4,\omega_6\). For (A.68.4), direct differentiation gives
\[
 \frac{dT}{W}=2(\omega_5-\omega_3),\qquad
 \frac{T\,dT}{W}=2(\omega_7-\omega_5+\omega_3-\omega_1).
                                                               \tag{A.68.5}
\]
Thus its pullback plane is spanned by
\(\omega_1-\omega_7\) and \(\omega_3-\omega_5\).
Since \(\mu^*\omega_j=\zeta_8^j\omega_j\), this plane and its
\(\mu^*\)-translate are complementary in the four-dimensional
odd-index space: for each pair \((j,8-j)=(1,7),(3,5)\),
the two eigenvalues are distinct. The three actual maps therefore
induce an isomorphism on rational cotangent spaces of their product
Jacobians, hence an isogeny
\[
                       J_C\sim J_{C_4}\times J_B^2.            \tag{A.68.6}
\]
All maps and the isogeny are defined after a finite number-field
extension; this is sufficient for potential reduction and the
derived-group formulation in A.41--A.42.

\paragraph{Exact Igusa invariants of the new genus-two factor.}
We use the standard Igusa invariants \(J_{2i}\), with
\(J_{10}=2^{-12}\operatorname{disc}(f)\) for our monic quintic.
For clarity the intermediate Igusa--Clebsch invariants are denoted
\(I^{\rm IC}_{2i}\), distinct from Liu's auxiliary invariants below.
Homogenize the quintic in (A.68.4) to a binary sextic \(F\).
With the normalized transvectant
\[
 (f,g)_k=\frac{(m-k)!(n-k)!}{m!n!}
 \sum_{j=0}^k(-1)^j\binom{k}{j}
 \frac{\partial^kf}{\partial x^{k-j}\partial z^j}
 \frac{\partial^kg}{\partial x^j\partial z^{k-j}},
                                                               \tag{A.68.7}
\]
where \(m=\deg f\) and \(n=\deg g\), put \(i=(F,F)_4\), \(d=(i,i)_2\), and
\(A=(F,F)_6,B_0=(i,i)_4,C_0=(i,d)_4\). Exact polynomial
differentiation gives
\[
 A=-8/5,\qquad B_0=1792/1875,\qquad C_0=-149504/46875.
\]
The standard conversion in \cite{fghSageIgusa}
\[
 \begin{split}
 I^{\rm IC}_2&=-120A,\\
 I^{\rm IC}_4&=-720A^2+6750B_0,\\
 I^{\rm IC}_6&=8640A^3-108000AB_0+202500C_0
 \end{split}
\]
gives \((192,4608,-516096)\), while
\(I^{\rm IC}_{10}=\operatorname{disc}(f)=-1048576\).
Using the conversions in \cite[equation (3), p.3]{fghDabrowskiSadek},
\[
 \begin{aligned}
 J_2&=I^{\rm IC}_2/8,\\
 J_4&=(4J_2^2-I^{\rm IC}_4)/96,\\
 J_6&=(8J_2^3-160J_2J_4-I^{\rm IC}_6)/576,\\
 J_8&=(J_2J_6-J_4^2)/4,
 \end{aligned}
\]
we obtain
\[
                 (J_2,J_4,J_6,J_8,J_{10})
                    =(24,-24,1248,7344,-256).                 \tag{A.68.8}
\]
The exact certificate also reproduces the official PARI example
\cite{fghPARIHyperelliptic}
\(x^5+3x^2-4\), after the explicitly checked weighted rescaling from \(P\) to \(4P\).
Every criterion below is a quotient of equal weights, so this
independent normalization check rules out a scaling discrepancy.

\paragraph{Liu's residue-two criterion determines the stable fibre.}
Liu's criterion \cite[Theorem 1, printed p.204]{fghLiu1993Stable} applies in residue characteristic two
with \(\varepsilon=4\) and \(I_{2\varepsilon}=J_8\). His auxiliary
invariants, defined on printed page 203, are
\[
 \begin{aligned}
 \mathcal I_4&=J_2^2-24J_4=1152,\\
 \mathcal I_{12}&=-8J_4^3+9J_2J_4J_6-27J_6^2-J_2^2J_8\\
                 &=-52641792.
 \end{aligned}
\]
Their valuations, with \(v_2(2)=1\), are
\(v_2(\mathcal I_4)=7\), \(v_2(\mathcal I_{12})=14\),
\(v_2(J_8)=4\), and \(v_2(J_{10})=8\). Therefore the three
expressions in Theorem 1\((V_*)\) have valuations
\[
 \begin{array}{c|ccc}
 \text{expression}&\mathcal I_4^4/J_8^2&J_{10}^4/J_8^5&
                              \mathcal I_{12}^4/J_8^6\\\hline
 v_2&20&12&32.
 \end{array}                                                   \tag{A.68.9}
\]
They lie in the maximal ideal, so the stable fibre has two
irreducible components meeting in one point. The two additional
expressions in part \((V)\) have valuations
\[
 v_2\!\left(\frac{\mathcal I_4^{12}}{J_{10}^4J_8}\right)=48,
 \qquad
 v_2\!\left(\frac{\mathcal I_{12}^{4}}{J_{10}^4J_8}\right)=20.
                                                               \tag{A.68.10}
\]
Both are integral, so the two components are smooth elliptic curves.
Moreover the theorem's formulas are
\[
 \begin{aligned}
 (j_1j_2)^4&=
  \overline{\mathcal I_4^{12}/(J_{10}^4J_8)},\\
 (j_1+j_2)^4&=2^6\cdot3^3+
  \overline{\mathcal I_{12}^{4}/(J_{10}^4J_8)}.
 \end{aligned}
\]
Both right-hand sides vanish in characteristic two. Hence
\(j_1=j_2=0\), and both elliptic components are supersingular.
Thus \(J_B\) has potential good reduction with \(2\)-rank zero.

The factor \(J_{C_4}\) has this same property by the proper-factor
calculation in A.37: its elliptic factors have integral \(j\)-invariants
reducing to zero at two. Potential good reduction is preserved by
isogeny (equivalently use an auxiliary prime in the
N\'eron--Ogg--Shafarevich criterion), and the \(2\)-rank of good
abelian reduction is isogeny invariant. Equation (A.68.6) therefore
proves potential good reduction of \(J_C\) with \(2\)-rank zero.
In particular the stable proper fibre of \(C\) is of compact type.

\paragraph{All four cusp differences are actual eight-torsion.}
The following identities hold in the function field of \(C\):
\[
 \begin{split}
 (v-w^8+1)(v+w^8-1)&=-4w^8,\\
 (v-w^8-1)(v+w^8+1)&=-8w^8.
 \end{split}                                                   \tag{A.68.11}
\]
There are no zeros or poles outside the four cusps. Expanding
\(\sqrt{1-6u^8+u^{16}}=1-3u^8-4u^{16}+O(u^{24})\) at zero,
and using \(u=w^{-1}\) at infinity, determines the exact divisors:
\[
 \begin{split}
 \operatorname{div}(v-w^8+1)&=8P_{0,-}-8P_{\infty,-},\\
 \operatorname{div}(v+w^8-1)&=8P_{0,+}-8P_{\infty,+},\\
 \operatorname{div}(v-w^8-1)&=8P_{0,+}-8P_{\infty,-},\\
 \operatorname{div}(v+w^8+1)&=8P_{0,-}-8P_{\infty,+}.
 \end{split}                                                   \tag{A.68.12}
\]
These four relations imply \(8[P-Q]=0\) for any two cusps.
After a finite extension giving the good Jacobian and separated
marked stable curve, each such divisor specializes to zero because
the special Jacobian has no nonzero geometric \(2\)-power torsion.
The torsion-divisor criterion proved in A.41 then says that the
joining path between any two cusp components contains only rational
components. Every cusp has a distinct companion cusp. A.42 applies
to every coefficient prime, every cusp packet and every open derived
subgroup, proving the stated full fixed-subgroup equality. A.59
independently supplies the normalizer assertion for the actual
dihedral closure (A.68.1).

\paragraph{Every actual intermediate cover of either dihedral closure.}
The same fixed-group conclusion holds for every intermediate
\(D^\circ\to E\setminus\{O\}\) of either degree-sixteen cover.
To prove this, extend the number field so that its actual geometric
subgroup \(H\subset G=D_{16}\), all cusps, and all quotient maps are
defined. Let \(D\) be the smooth proper normalization of the quotient
function field \(\overline K(C)^H\). There are actual finite maps
\[
                   C\xrightarrow{f}D\longrightarrow E,
\]
and the missing points of \(D^\circ\) are precisely the images of
the four cusps of \(C\), equivalently the inverse image of \(O\).
The compact genus of \(D\) is at least one, by Riemann--Hurwitz
for its nonconstant finite map to \(E\).

At \(P_{0,+}\), the elliptic parameter \(-X/Y\) from (A.68.1)
has order four: its leading term is \(2w^4\). The actual inertia
is therefore
\[
 I=\langle r^2\rangle\simeq C_4.
\]
This also follows directly from the displayed deck action:
\(r\) interchanges the two cusps over zero, \(r^2\) fixes
each, and every reflection sends zero to infinity. In particular
\(I\) is normal in \(G\) and is the inertia at every cusp.
Moreover
\[
 \Phi(G)=G^2[G,G]=\langle r^2\rangle=I.
\]
One can check the required Frattini property without invoking
a classification theorem. A subgroup whose image is all of
\(G/I\simeq C_2^2\) contains an odd power of \(r\) and a
reflection: if two independent images are both reflections,
their product gives that odd rotation. These elements generate
\(G\). Consequently \(HI=G\) forces \(H=G\).

The cusps of \(D\) are the actual orbit set
\(H\backslash G/I\). Normality of \(I\) gives its cardinality
\([G:HI]\), which is at least two for every proper subgroup
\(H<G\). For \(H=G\), the intermediate is the original
one-cusp elliptic curve, whose assertion is already proved.

For a proper \(H\), the norm
\(f_*:J_C\to J_D\) is surjective: its composite with pullback
is multiplication by \(\deg f\). Thus \(J_D\) is an isogeny
quotient of \(J_C\). After a finite extension at two it has
good reduction, by the auxiliary-prime
N\'eron--Ogg--Shafarevich criterion, and its special \(2\)-rank
is zero. Indeed an abelian quotient of a good abelian variety
is good, and its rational \(2\)-divisible group is an isogeny
factor; it cannot acquire a slope-zero part absent in \(J_C\).
The proper stable reduction of \(D\) is therefore of compact type.

Let \(P',Q'\) be any two distinct cusps of \(D\). Choose
cusps \(P,Q\) of \(C\) mapping respectively to them. The
relations (A.68.12) give a rational function \(g\) on \(C\)
with \(\operatorname{div}(g)=8(P-Q)\). The function-field
norm satisfies
\[
 \operatorname{div}(\operatorname{Nm}_{K(C)/K(D)}g)
       =f_*\operatorname{div}(g)=8(P'-Q').
\]
There is no ramification multiplier in this pushforward: over
an algebraically closed constant field, the residue degree
of each point is one. This is a norm of an actual function,
not an inference from pullback torsion alone.

Hence every distinct pair of target cusps has difference
annihilated by eight. Its specialization in the good
\(2\)-rank-zero target Jacobian is zero. Separate all marks
in a stable marked model of \(D\); adding rational components
for this separation preserves compact type. The torsion-divisor
criterion of A.41 proves that the unique joining path between
the components carrying \(P'\) and \(Q'\), including both
endpoints, is rational. Since every selected cusp has a
distinct companion, A.42 applies at every coefficient prime.
Its finite-field-restriction argument supplies the indicated
original open-derived subgroups and arbitrary induced packets.
Finally A.59 supplies all normalizers, because the actual
geometric Galois closure of an intermediate is a quotient
of this finite two-group. No pro-odd injectivity of the map
from an intermediate into \(C\) is used.

All finite identities used here are checked in the retained programs
\path{code/a_dihedral_eight_fixed.py},
\path{code/a_dihedral_eight_maps.py}, and
\path{code/a_dihedral_eight_intermediates.py}; the general
fixed-group implication is the proved A.41--A.42 theorem, rather than
an inference from a finite computation.
\end{proof}

% A.69: full root, DE and PBC proof/source/exact-certificate audits PASS.
\subsection{The nonstable dihedral family at every two-power level}

\begin{statement}{A.69}{Proved}{Every two-power member of the nonstable dihedral family}
Let \(m\ge2\) be a power of two, and let \(C_m\) be the smooth proper
normalization of
\[
                    v^2=w^{2m}-6w^m+1.
\]
Give \(C_m\setminus S_m\), where \(S_m\) consists of the four points
above \(w=0,\infty\), the actual degree-\(2m\) map to
\(E\setminus\{O\}\) below. Its proper Jacobian has potential good
reduction of \(2\)-rank zero at two. Every cusp difference is killed
by \(m\). Consequently both all-prime assertions of Task A.5 hold
for this cover and for every actual intermediate of its dihedral
Galois closure: every induced cusp packet has precisely its selected
inertia as the fixed geometric pro-\(\ell\) subgroup under every
open derived Galois subgroup, and every arithmetically normalized
procyclic pro-\(\ell\) subgroup is cuspidal. The same assertions hold
for the plus-six family, identified with this one after a finite
extension by \(w\mapsto\zeta_{2m}w\).

This proves a uniform family and all its intermediates; the universal
cover question A.45 remains open.
\end{statement}

\begin{proof}
All curves, automorphisms, quotient maps, and chosen cusp sections can
be defined over one finite number-field extension. Further finite
extensions used for reduction are harmless in the formulation of
A.41--A.42. We always distinguish the dihedral deck action from the
additional cyclic automorphism used to control the Jacobian.
At the chosen place over two we enlarge this field so that the
elliptic base has good reduction, the Jacobians below have good
reduction, and all displayed endomorphisms and automorphisms are
defined. The local actions throughout fix the resulting coefficient
valuation ring.

\paragraph{The actual cover.}
The equation is smooth of genus \(m-1\), and the actual map is
\[
 X=\frac{(w^m-1)^2}{4w^m},\qquad
 Y=\frac{v(w^{2m}-1)}{8w^{3m/2}},\qquad Y^2=X^3-X.
                                                        \tag{A.69.1}
\]
It factors through elliptic duplication after
\[
 x=\frac{w^m-1+v}{2},\qquad y=w^{m/2}(x-1),
 \qquad y^2=x^3-x,\qquad w^{m/2}=\frac{y}{x-1}.
                                                        \tag{A.69.2}
\]
The divisor of the last elliptic function is
\(T_0+T_{-1}-T_1-O\). Hence the first map is cyclic of degree
\(m/2\), totally ramified at those four points, and unramified
elsewhere; duplication is unramified. Thus (A.69.1) has degree
\(2m\), its four cusp indices are \(m/2\), and it is finite
\'{e}tale off \(S_m\). Its actual deck transformations
\[
 r(w,v)=(\zeta_m w,-v),\qquad s(w,v)=(1/w,-v/w^m)
                                                        \tag{A.69.3}
\]
satisfy \(r^m=s^2=1\), \(srs=r^{-1}\), and give \(2m\)
distinct transformations. The deck group is therefore
\(D_{2m}\). For \(m=2\) this is the unramified degree-four
elliptic isogeny case, so \(J_{C_2}\) is isogenous to \(E\).

\paragraph{Potential good reduction: the precise ramified input.}
After a finite extension at two the original elliptic curve has a
smooth proper model. The branch divisor consists of the single
section \(O\), so its closure is \'{e}tale over the valuation ring.
The finite dihedral group is nilpotent. The ramified
nilpotent-cover theorem, as stated precisely in
Lehr--Matignon, Theorem 6.13, applies: the branch closure is \'{e}tale
over the valuation ring, and generic tameness means tameness in
characteristic zero, which holds here. It gives potential good
reduction of \(J_{C_m}\), equivalently a tree as the dual graph of
the proper stable fibre. Sa\"idi's statement of the same ramified
extension, printed page 662, confirms this scope.
\cite{fghLehrMatignon2005,fghSaidi2007}
We use only this tree/potential-good conclusion. Raynaud's
unramified theorem about strictly interior slopes of contracted
components is not being extended to a marked branch fibre.

\paragraph{The new factor and its integral characteristic-two restriction.}
Put
\[
 \mu(w,v)=(\zeta_mw,v),\qquad
 \sigma=\mu^{m/2},\qquad C_m/\langle\sigma\rangle=C_{m/2}
                                                        \tag{A.69.4}
\]
for \(m\ge4\); the quotient coordinate is \(w^2\).
Let \(A_m\) be the anti-invariant isogeny factor of \(J_{C_m}\)
for \(\sigma\). Then
\[
 J_{C_m}\sim J_{C_{m/2}}\times A_m,\qquad
 \dim A_m=m/2,
 \qquad \mathbf Q(\zeta_m)\hookrightarrow\operatorname{End}^0(A_m).
                                                        \tag{A.69.5}
\]
Indeed the regular differentials \(w^{j-1}dw/v\),
\(1\le j\le m-1\), have \(\mu\)-eigenvalues \(\zeta_m^j\).
The anti-invariant indices are exactly the odd indices; since \(m\)
is a power of two, every such eigenvalue is primitive of order
\(m\). This proves the field action, including its unit, and shows
that \(\mu^k-1\) is invertible on every rational realization of
\(A_m\) whenever \(m\nmid k\).

After the indicated finite extension giving good reduction, all
these endomorphisms extend to the abelian scheme and the same
unital field acts on the rational \'{e}tale Tate space of the special
\(2\)-divisible group. Since
\(\mathbf Q(\zeta_m)\otimes\mathbf Q_2\) is a field of degree
\(\varphi(m)=m/2\), the special \(2\)-rank of \(A_m\) is either
zero or \(m/2\). In the latter case \(A_m\) is ordinary.
This is a dimension argument in the actual special \'{e}tale Tate
space, not a conclusion from a Hodge character pattern.

We prove by induction that the ordinary alternative cannot occur.
The base \(m=2\) was identified with an isogeny of the original
supersingular elliptic curve. Assume the result for \(m/2\), and
suppose for a contradiction that \(A_m\) is ordinary.

\paragraph{Componentwise consequences of an ordinary anti-invariant factor.}
Choose the stable marked model of \((C_m,S_m)\) over a complete
valuation ring \(R\), after the finite extension needed to define
all automorphisms and sections. Its residue field may be replaced
by an algebraic closure. Its dual graph \(\Gamma\) is a tree;
inserting marked rational components does not change this fact or
the product decomposition of its identity-component Picard scheme.
Thus the good special Jacobian is the product of the Jacobians of
the positive-genus components. The involution \(\sigma\) acts on
this actual product.

If \(\sigma\) exchanges two positive-genus components, their
diagonal and anti-diagonal factors are both isogenous to the same
component Jacobian. The diagonal is a factor of the specialization
of \(J_{C_{m/2}}\), whose \(2\)-rank is zero, while the
anti-diagonal is a factor of the ordinary \(A_m\). This is
impossible. A component on which \(\sigma\) acts identically
contributes only to the invariant factor; no restriction on its
geometric genus is inferred from that trivial action.

Consider a component \(B\) fixed by \(\sigma\), on which it acts
nontrivially. Its quotient \(B'=B/\langle\sigma\rangle\) is a
smooth proper curve, and its invariant Jacobian is isogenous to
\(J_{B'}\), by the actual pullback and norm identities. Hence
\(f(B')=0\), where \(f\) denotes \(2\)-rank. The anti-invariant
Prym is a factor of \(A_m\) and is ordinary whenever it has
positive dimension.

The Deuring--Shafarevich formula and Riemann--Hurwitz give, if
\(r_B\) is the number of fixed points and \(d_P\) the local
different exponent,
\[
 f(B)=r_B-1,\qquad
 g(B)-g(B')=g(B')-1+\frac12\sum_{P\ {\rm fixed}}d_P,
 \qquad d_P\ge2.                                  \tag{A.69.6}
\]
There must be a fixed point, since an unramified connected
double cover of a \(2\)-rank-zero curve would give negative
\(2\)-rank. If the Prym had dimension zero, (A.69.6) would force
\(g(B')=0\), \(r_B=1\), and \(g(B)=0\); it contributes no
positive genus. Otherwise ordinarity says
\(f(B)=g(B)-g(B')\). Consequently
\[
      r_B=g(B')+\frac12\sum d_P
      \quad\Longrightarrow\quad
      g(B')=0,\qquad d_P=2\ \hbox{for every fixed }P.
                                                        \tag{A.69.7}
\]
Thus every positive-genus component on which \(\sigma\) is
nontrivial is ordinary, and \(\sigma\) is its hyperelliptic
involution. The Deuring--Shafarevich input is the ordinary
curve formula, not a formula applied to the marked vertical fibre.
\cite{fghCrew1984,fghYang2024}

\paragraph{An ordinary hyperelliptic curve has no cyclic fourth root of its involution.}
Here is the required elementary ramification argument. If a
\(2\)-group acts faithfully on an ordinary smooth curve, apply
Deuring--Shafarevich and Riemann--Hurwitz and subtract twice the
first formula from the second. The result is
\[
 0=2|G|\bigl(g(B/G)-f(B/G)\bigr)
       +\sum_{P\in B/G}\frac{|G|}{e_P}
                          \bigl(d_P-2(e_P-1)\bigr).
                                                        \tag{A.69.8}
\]
Every term is nonnegative: for wild \(2\)-group inertia the
zeroth and first ramification groups coincide, so
\(d_P=\sum_{i\ge0}(|G_i|-1)\ge2(e_P-1)\).
Indeed the local minimal polynomial of a uniformizer has derivative
\(\prod_{g\ne1}(t-g(t))\); its valuation is
\(\sum_{g\ne1}\operatorname{ord}_t(t-g(t))\), which is this
ramification-group sum. We use the classical lower numbering,
as in \cite[\S2.2, equation (3) and Exercise 2.12(c)]{fghBouwWewers}.
Equality forces the second ramification
group to be trivial. The first ramification quotient embeds in
the additive residue field via the coefficient in
\(t\mapsto t+a t^2+\cdots\); it has exponent two. Thus no
cyclic inertia group of order four can fix a point of an ordinary
curve.

Suppose an automorphism \(a\) of an ordinary curve has order four
and \(a^2=\sigma\), with \(B/\sigma\simeq\mathbf P^1\).
Its induced involution on \(\mathbf P^1\) has a unique fixed
point \(Q\) in characteristic two. If \(Q\) were unramified in
the double cover, \(a\) would permute its two points, so \(a^2\)
would fix them individually; but the deck involution \(\sigma\)
exchanges them. Thus \(Q\) is ramified, its unique point above it
is fixed by \(a\), and the cyclic inertia of order four
contradicts (A.69.8).

\paragraph{The primitive cyclic action forces one full orbit of elliptic components.}
Take a positive-genus component \(B\) contributing to the
anti-invariant factor, and let \(d\) be its orbit length under
\(\mu\). The restriction of \(\mu^d\) to \(B\) has exact order
\(m/d\). Otherwise a smaller power of \(\mu\) would be the
identity on the nonzero Jacobian contribution of this entire
orbit, contrary to the primitive field action in (A.69.5).
Moreover \(\sigma\) acts nontrivially on \(B\), so \(m/d\ge2\).
The induced automorphism on \(B/\sigma=\mathbf P^1\) has order
\(m/(2d)\). A finite \(2\)-power-order element of
\(\operatorname{PGL}_2\) in characteristic two has order at most
two. Hence \(m/d\) is two or four; the fourth-root argument
excludes four. Therefore \(d=m/2\).

The sum of the genera of these components is precisely
\(\dim A_m=m/2\): exchanged positive-genus pairs were excluded,
and each invariant quotient here has genus zero. There is
therefore exactly one orbit of \(m/2\) components, each an
ordinary elliptic curve. On each of them \(\sigma\) has exactly
two fixed points by (A.69.6).

\paragraph{A smooth fixed point really lifts in this wild situation.}
Suppose \(P\) is a smooth special point fixed by \(\sigma\) on
a component on which \(\sigma\) is nontrivial. The action fixes
the coefficient ring \(R\). In a smooth formal coordinate its
completed local ring is \(R[[t]]\), and
\(F(t)=\sigma(t)-t\) has nonzero reduction of finite order.
Weierstrass preparation writes \(F\) as a unit times a
distinguished monic polynomial. Thus the fixed locus in this
formal disc is finite flat of positive degree over \(R\), and
has generic fixed points specializing to \(P\). The smooth
formal-disc identification realizes its roots as points on the
generic curve, and properness extends them to sections after a
finite extension. This is the usual
one-variable Weierstrass preparation theorem, with no uniformizer
factor because the reduction is nonzero.
\cite[Corollary 2.16]{fghVonk2021}
For the ordinary
elliptic components the degree is two, since their different
exponent is two. There is no vertical factor: nonzero reduction
of \(F\) is exactly what rules it out.

The generic fixed points of \(\sigma(w,v)=(-w,v)\) are exactly
the four cusps. Each of those cusps is fixed individually by
\(\mu\). Therefore a smooth \(\sigma\)-fixed point cannot lie
on one of the nontrivially permuted elliptic components: its
lifted cusp section would force that component to be fixed by
\(\mu\). This argument concerns smooth special points only;
no lifting assertion is made at a node or on a component on
which \(\sigma\) is the identity.

\paragraph{The rooted-tree genus contradiction.}
Root \(\Gamma\) at the component containing one selected marked
cusp. This vertex is \(\mu\)-fixed. All the \(m/2\) ordinary
elliptic vertices have the same positive distance from this root.
On each such vertex \(B\), the parent edge and its attaching
point are fixed by \(\sigma\), because \(\sigma\) fixes both
the root and \(B\). This accounts for one of its two fixed
points. Its other fixed point cannot be smooth by the preceding
paragraph, so it is a different node, lying on a child edge.

The descendant subtree beyond that child edge contains no marked
cusp: if it did, the path from the root to that \(\mu\)-fixed
mark would be fixed pointwise as a graph by \(\mu\), forcing
\(B\) itself to be fixed. It contains no member of the ordinary
elliptic orbit, since every member of that orbit has the same
distance from the root. A nonempty finite unmarked subtree
attached by one edge must contain a positive-genus component;
otherwise it has a rational leaf with fewer than three special
points, contrary to stability of the marked model. These
descendant subtrees are disjoint for the distinct ordinary
elliptic vertices, again because those vertices have equal
distance from the root. Their contributions, together with
the elliptic vertices themselves, give
\[
             g(C_m)\ \ge\ \frac m2+\frac m2=m,
                                                        \tag{A.69.9}
\]
contrary to \(g(C_m)=m-1\). This excludes the ordinary
alternative and completes the induction proving special
\(2\)-rank zero.

\paragraph{Cusps and every actual intermediate.}
The identities
\[
 \begin{aligned}
 (v-w^m+1)(v+w^m-1)&=-4w^m,\\
 (v-w^m-1)(v+w^m+1)&=-8w^m
 \end{aligned}                                           \tag{A.69.10}
\]
have exactly the cusp divisors used in A.68 with eight replaced
by \(m\). More explicitly, write \(P_{0,\pm}\) for the points
with \(v=\pm1\), and \(P_{\infty,\pm}\) for those with
\(v/w^m\to\pm1\). The expansion
\(\sqrt{1-6u^m+u^{2m}}=1-3u^m-4u^{2m}+O(u^{3m})\)
at zero, and at infinity with \(u=w^{-1}\), gives
\[
 \begin{aligned}
 \operatorname{div}(v-w^m+1)&=mP_{0,-}-mP_{\infty,-},\\
 \operatorname{div}(v+w^m-1)&=mP_{0,+}-mP_{\infty,+},\\
 \operatorname{div}(v-w^m-1)&=mP_{0,+}-mP_{\infty,-},\\
 \operatorname{div}(v+w^m+1)&=mP_{0,-}-mP_{\infty,+}.
 \end{aligned}                                             \tag{A.69.12}
\]
These connect all four cusps and show that every cusp difference
is annihilated by \(m\).
The cusp inertia in (A.69.3) is
\(I=\langle r^2\rangle=\Phi(D_{2m})\), including the trivial
Frattini subgroup when \(m=2\). It is normal. For every proper
subgroup \(H<D_{2m}\), the actual intermediate
\(D=C_m/H\) has
\[
        \#S_D=\#(H\backslash D_{2m}/I)
                 =[D_{2m}:HI]\ge2.                 \tag{A.69.11}
\]
For completeness, if \(HI=D_{2m}\), surjectivity onto
\(D_{2m}/I\simeq C_2^2\) supplies in \(H\) an odd rotation
and a reflection. The odd rotation generates \(\langle r\rangle\),
so these two elements generate \(D_{2m}\) and \(H=D_{2m}\).
This proves the displayed strict inequality without an assumption
about chosen lifts. The actual compact curve \(D\) maps finitely to
\(E\), so has genus at least one. Its Jacobian is an isogeny
quotient of \(J_{C_m}\) by the actual norm, hence has potential
good reduction and special \(2\)-rank zero. Over an algebraic
closure the pushforward of a point is its image with coefficient
one; ramification indices do not enter this pushforward.
Taking norms of the functions in (A.69.10) therefore shows that
every difference of cusps on \(D\) is killed by \(m\).

Their reductions in the good special Jacobian are zero, because
a \(2\)-rank-zero abelian variety has no nonzero geometric
\(2\)-power torsion points. The compact-type Abel-path criterion
of A.41 gives a rational joining path to a distinct companion
for every cusp. A.42 therefore proves the fixed-subgroup
assertion at every coefficient prime, including two. The
normalizer assertion follows from A.59, since every such
closure is a \(2\)-group. The case \(H=D_{2m}\) is the already
proved original punctured elliptic curve. This completes the
proof for every intermediate and both signs.
The retained exact support is \path{code/a_dihedral_two_power_fixed.py};
its finite identities do not replace the unbounded geometric induction.
\end{proof}

% A.70: full root and DE proof/source/exact-certificate audits PASS.
\subsection{Every two-power dihedral geometric closure}

\begin{statement}{A.70}{Proved}{Both cusp assertions for every two-power dihedral closure}
Let \(Y\to U=E\setminus\{O\}\) be an actual connected finite
\'{e}tale cover whose geometric Galois closure has dihedral group
\(D_{2m}\), where \(m\) is a power of two. Then both assertions
of Task A.5 hold for \(Y\), at every coefficient prime, every
induced cusp packet, and every permitted open derived subgroup.
In particular its full fixed geometric pro-\(\ell\) subgroup is
its selected cusp inertia, and every nontrivial arithmetically
normalized procyclic pro-\(\ell\) subgroup is cuspidal.

For \(m\ge4\), there are exactly three geometric Galois covers
with this group. As actual covers of the fixed elliptic base they
are the three models
\[
 C_{m,a,b}:\ v^2=w^{2m}+aw^m+b,\qquad
 (a,b)\in\{(0,-1),(6,1),(-6,1)\},                 \tag{A.70.1}
\]
with the maps below. Every intermediate of each is included.
The order-four case has one geometric \(V_4\) cover and is
handled separately. This conclusion does not cover arbitrary
nilpotent closures or products of two nonabelian Sylow groups;
the universal statement A.45 remains open.
\end{statement}

\begin{proof}
\paragraph{Three actual rotation quotients.}
For each model in (A.70.1), put
\[
 t=w^m,\qquad u=w^{m/2}v,\qquad
 E_{a,b}:u^2=t^3+at^2+bt.                         \tag{A.70.2}
\]
This is a degree-\(m\) map. Indeed the function-field extension
over \(\overline K(t)\) through \(\overline K(w,v)\) has
degree \(2m\), while \(\overline K(E_{a,b})\) has degree two.
The polynomial in (A.70.1) has simple roots, and is not a square
in \(\overline K(w)\). Thus the compact source is connected,
smooth after normalization, and has genus \(m-1\).
The cyclic rotation group is generated by
\[
                 r(w,v)=(\zeta_mw,-v).            \tag{A.70.3}
\]
It fixes \((t,u)\) and has order \(m\), so (A.70.2) is its
actual quotient, not merely an isogenous elliptic factor.

The degree-two isogeny with kernel \(\{O,(0,0)\}\) on
\(E_{a,b}\) is given by
\[
 t'=t+a+b/t,\qquad u'=u(1-b/t^2),\qquad
 u'^2=t'^3-2at'^2+(a^2-4b)t'.                    \tag{A.70.4}
\]
For \(a=\pm6,b=1\), identify its target with the fixed
\(E:Y^2=X^3-X\) by
\[
                  X=(t'-2a/3)/4,\qquad Y=u'/8.
                                                               \tag{A.70.5}
\]
For \((a,b)=(0,-1)\), use
\[
                  X=-it'/2,\qquad Y=-(1+i)u'/4.
                                                               \tag{A.70.6}
\]
These are the explicit isogenies already checked in A.4. Their
composites with (A.70.2) have degree \(2m\) and have exactly
four points above \(O\), namely the two points above each of
\(w=0,\infty\). The cyclic map is ramified only there, with
index \(m/2\), because \(\operatorname{div}(t)=2(0,0)-2O\)
on \(E_{a,b}\). The isogeny is unramified.

Choose \(\eta\) with \(\eta^m=b\), after a finite extension.
The actual reflection
\[
 s(w,v)=\left(\eta/w,\ -\frac{b}{\eta^{m/2}}\frac v{w^m}\right)
                                                               \tag{A.70.7}
\]
satisfies \(s^2=1\) and \(srs=r^{-1}\). It induces
\[
                   (t,u)\longmapsto(b/t,-bu/t^2),
                                                               \tag{A.70.8}
\]
which is translation by \((0,0)\) on \(E_{a,b}\).
Thus the composite to the fixed base is invariant under a faithful
\(D_{2m}\) action of order equal to its degree, and is the actual
Galois cover claimed in the statement.

\paragraph{The three degree-two covers of the fixed base are distinct.}
The nonzero kernel point of the dual of (A.70.4) is \((t',u')=(0,0)\).
One can see this directly: the two nonkernel nonzero two-torsion
points on \(E_{a,b}\) satisfy \(t^2+at+b=0\), so both map
to \((0,0)\); the pullback--dual identity says that this image
is precisely the dual kernel. This is the usual actual dual-isogeny
identity \cite[III.6.1--6.2]{fghSilvermanAEC}.
Under (A.70.5)--(A.70.6) these
three dual kernels on the \emph{fixed} elliptic curve are
\[
 \begin{array}{c|c}
 (a,b)&\text{nonzero dual-kernel point on }E\\ \hline
 (0,-1)&(0,0)\\
 (6,1)&(-1,0)\\
 (-6,1)&(1,0).
 \end{array}                                                   \tag{A.70.9}
\]
An isomorphism of degree-two covers over the identity of \(E\)
preserves this dual kernel. Hence the three rotation quotients
are different covers of the specified base. The two nonstable
source curves may be isomorphic over an algebraic extension;
that isomorphism covers \([i]\) on \(E\), not its identity,
and does not identify their marked covering kernels.

\paragraph{Exhaustion of the geometric dihedral closures.}
The geometric fundamental group of \(U\) is free profinite on
two generators. Connected geometric Galois covers with group
\(D_{2m}\) correspond to continuous epimorphisms onto this
finite group, modulo its automorphism group. Equivalently they
are generating ordered pairs modulo automorphisms.
For \(m\ge4\), the rotation subgroup \(R=\langle r\rangle\)
is characteristic: it is generated by the elements of order
strictly greater than two. Every automorphism has the form
\[
       r\longmapsto r^e,\qquad s\longmapsto r^f s,
       \qquad e\in(\mathbf Z/m)^\times,\ f\in\mathbf Z/m.
                                                               \tag{A.70.10}
\]
A generating pair must have at least one reflection. If only
one entry is a reflection, the other is an odd rotation. If
both entries are reflections, their exponents differ by an
odd integer. Formula (A.70.10) is transitive on each of the
three resulting types, with representatives
\[
                      (r,s),\qquad(s,r),\qquad(s,rs).
                                                               \tag{A.70.11}
\]
These types cannot be equivalent because their images in the
characteristic quotient \(D_{2m}/R\simeq C_2\) are the three
different nonzero ordered binary pairs. Thus there are exactly
three geometric covering kernels, distinguished by the actual
rotation quotient. The three models above have distinct such
quotients, so they exhaust these kernels.

For \(m=2\), the group is \(V_4\), its rotation subgroup is
not characteristic, and \(\operatorname{Aut}(V_4)\) is transitive
on generating ordered pairs. There is one cover, the maximal
exponent-two abelian cover, which extends to the elliptic
isogeny \([2]\). Its intermediates have genus one and are
settled by A.40. A cyclic group of order two, if included under
the degenerate dihedral notation, is likewise an elliptic-isogeny
case.

\paragraph{Every intermediate, including the stable family.}
For the two nonstable models, A.69 already proves both assertions
for every actual intermediate. Its potential-good and cusp-torsion
statements are intrinsic to the marked compact curve, so they
apply to the explicit maps (A.70.5), including the plus-six map.
For the stable model, A.44 proves potential good reduction of its
Jacobian at two, special \(2\)-rank zero, and annihilation of
every cusp difference by \(4m\). We spell out why the same
hypotheses reach every intermediate, rather than infer fixedness
through a pro-prime quotient.

At every cusp of (A.70.1), the actual inertia is
\(I=\langle r^2\rangle\), normal in \(D_{2m}\).
For every proper subgroup \(H<D_{2m}\), the quotient curve
\(D=C_{m,0,-1}/H\) has cusp count
\[
                   [D_{2m}:HI]\ge2.              \tag{A.70.12}
\]
Indeed \(I=\Phi(D_{2m})\); more elementarily, \(HI=D_{2m}\)
would supply an odd rotation and a reflection in \(H\), forcing
\(H=D_{2m}\). The actual norm map makes \(J_D\) an isogeny
quotient of the source Jacobian, so it has potential good
reduction with special \(2\)-rank zero. Over an algebraically
closed field the pushforward of a point is its image with
coefficient one. Norms of the actual functions annihilating
source cusp differences therefore kill every difference of
target cusps by \(4m\). The compact genus of \(D\) is at least
one, because it has a nonconstant finite map to \(E\).

Consequently the torsion-divisor criterion in A.41 gives a
rational joining path to a distinct companion for every cusp
of \(D\). A.42 proves full fixed inertia at every coefficient
prime; A.59 proves every normalizer assertion, since the actual
closure remains a two-group. The quotient \(H=D_{2m}\) is the
original \(U\), already settled.

Finally any cover in the statement becomes, after a finite
number-field extension, an actual intermediate of one of these
three geometric closures. The covering isomorphisms, subgroup
and marked points can all be defined over a finite extension.
A geometric path used to identify base points need not be
defined over that field: it gives the arithmetic comparison up
to geometric inner conjugacy, under which the packet, fixed-group
and normalizer assertions are invariant.
The finite-field-restriction argument of A.41--A.42
then gives the original permitted open derived subgroups and
packets; normalizer recognition also descends by intersecting
its open arithmetic image with this finite-index Galois subgroup.
No assertion that fixed points commute with a quotient is used.
The exact isogeny identities and finite generating-pair support
are retained in \path{code/a_all_dihedral_two_power.py}.
\end{proof}

% A.71: full root, DE and PBC proof/source audits PASS; exact certificate rerun PASS.
\subsection{Odd abelian factors over two-power dihedral covers}

\begin{statement}{A.71}{Proved}{Dihedral two-power closures with an odd abelian factor}
Both all-prime
assertions of Task A.5 hold for every actual cover whose geometric
Galois closure is a quotient of
\(D_{2m}\times B\), where \(m\) is a power of two and \(B\)
is an abelian group of odd order. No assertion is made for a
product with a second nonabelian Sylow factor. The universal
question A.45 remains open.
\end{statement}

\begin{proof}
\paragraph{The actual group and function-field factorization.}
A quotient of a direct product of groups of coprime orders is
the product of the two corresponding quotients. Indeed every
subgroup of such a product splits into its two intersections:
for an element \((x,y)\), suitable powers, using Bezout for the
two group orders, give \((x,1)\) and \((1,y)\) in the subgroup.
Apply this also to the normal kernel of the quotient. Every
quotient of a dihedral two-group is dihedral or abelian; the
abelian cases will be genus-one covers. Thus the nonabelian
case has actual closure group \(G_2\times B_0\), with \(G_2\)
dihedral of two-power order and \(B_0\) abelian of odd order.
Every actual covering subgroup has the form \(H_2\times H_B\).

The abelian cover of \(U\) corresponding to \(B_0/H_B\)
extends to an actual elliptic isogeny
\[
                         \eta:E'\longrightarrow E.
\]
The peripheral element is a commutator and is killed by every
abelian quotient, so this extension is unramified on the proper
elliptic curve. After a finite number-field extension choose an
origin above \(O\); the finite morphism is then an isogeny.
Let \(Z\to E\) be the compactification of the actual
\(H_2\)-intermediate of the dihedral closure. The given cover
has function field equal to the compositum of those of \(Z\)
and \(E'\) over \(E\). This is a field: the normal closure
of the first extension has two-power degree, and the second is
Galois of odd degree, so their intersection is the base field.
Consequently the given compact curve is exactly the connected
normalization of
\[
                         Z\times_E E'.
\]
The removed set is the full inverse image of \(O\), namely the
pairs of a cusp of \(Z\) and a point of \(\eta^{-1}(O)\).
This identifies the actual covering map, not merely its degree
or geometric representation.

\paragraph{Good models of the odd isogeny and of the actual map.}
If \(Z\) has genus one, so does its unramified isogeny pullback,
and A.40 applies. Otherwise A.70, through A.44 or A.69 and the
actual norm maps, gives a potential-good Jacobian at two and a
stable marked model \(\mathcal Z/R\) of compact type in which
every cusp has a distinct companion joined by rational
components. Enlarge the local field at two so that this model,
all marks and the good elliptic model \(\mathcal E/R\) are
defined. The isogenous elliptic curve \(E'\) also has good
reduction, and \(\eta\) extends to an isogeny
\(\mathcal E'\to\mathcal E\) of odd degree. Its finite flat
kernel has order invertible in the residue field, so this
extended isogeny is finite \'{e}tale.

The actual generic map \(q:Z\to E\) extends over
\(\mathcal Z\). Here is a construction avoiding an assertion
about an arbitrary rational map on a singular surface. Use a
marked section \(P\) with \(q(P)=O\). The compact-type
Abel map based at \(P\) extends to
\(\mathcal Z\to\operatorname{Pic}^0(\mathcal Z/R)\), as in
the precise Abel-map input of A.41
\cite[Theorem 1.1 and Proposition 3.2]{fghDupuyRabinoff}.
The actual homomorphism
\(J_Z\to E\) extends between the good abelian schemes.
Their composition restricts generically to \(q\), and is the
required extension. Inserted rational components are mapped
constantly by this same construction.

It follows that
\[
              \mathcal Y=\mathcal Z\times_{\mathcal E}\mathcal E'
\]
is finite \'{e}tale over \(\mathcal Z\). It is a proper flat
semistable model of the actual connected generic pullback.
The generic-fibre field comparison above is what proves its
geometric connectedness; no disconnected formal fibre product
is selected without justification.

\paragraph{Why the pulled-back dual graph is a tree.}
One must not infer compact type just from the finite \'{e}tale
map \(\mathcal Y\to\mathcal Z\): covers of component curves
can join several copies of the original graph and create cycles.
Instead take the actual geometric Galois closure over \(E\).
Its group \(G_2\times B_0\) is nilpotent. Its sole possible
branch point is \(O\), since the odd isogeny is unramified.
Over the good elliptic model the branch closure is an \'{e}tale
section. Generic characteristic zero supplies tameness.
Lehr--Matignon Theorem 6.13, with exactly the scope checked in
A.69, gives a potential-good Jacobian of this actual closure
\cite{fghLehrMatignon2005}.
The actual norm to the given intermediate then gives potential
good reduction of \(J_Y\). After a further finite extension,
the semistable special fibre of \(\mathcal Y\) therefore has a
tree dual graph. No claim about its special two-rank is needed.

\paragraph{Every selected cusp has a lifted rational companion.}
Let a chosen cusp of the pullback lie above a cusp \(P\) of
\(Z\). Choose a distinct companion \(Q\) on \(Z\), and let
\(\mathcal T\) be the connected rational subtree in the
geometric special fibre joining their marked vertices. It
includes both endpoint components. Each component is a
projective line, so its map to the good elliptic curve is
constant; the constants agree across the nodes. Since the
cusp sections map to \(O\), the whole subtree maps to the
single special point \(O\).

The finite \'{e}tale isogeny has, over this algebraically closed
residue field, a discrete fibre of \(\deg\eta\) distinct
points above \(O\). Therefore its pullback over \(\mathcal T\)
is literally the disjoint union of \(\deg\eta\) copies of
\(\mathcal T\). The selected marked section lies on exactly
one such copy. That same copy contains a marked section above
\(Q\): sections above a cusp are the pairs with the kernel
sections of \(\eta\), and prime-to-two kernel sections have
distinct reductions. These two generic marked points are
distinct because their images \(P,Q\) are distinct.

Thus the selected cusp is joined to a genuine companion by an
entirely rational path in \(\mathcal Y\). Stabilization may
contract rational components but cannot insert a positive-genus
component into this path. Since the whole graph is a tree, this
is the unique joining path required by A.42. The argument works
for every selected cusp and does not require the new Jacobian
to have two-rank zero.

A.42 now proves the full fixed-inertia assertion at every
coefficient prime. A.59 proves every normalizer assertion for
the nilpotent geometric closure. As in A.70, isomorphisms,
subgroups and marked points need only a finite field of
definition; geometric base-point paths are used up to inner
conjugacy, and the proved finite-field-restriction argument
supplies every original permitted packet and open derived
subgroup. The exact subgroup and curve-degree support is retained in
\path{code/a_dihedral_abelian_products.py}.
\end{proof}

% A.72: full root and DE proof/source/exact-certificate audits PASS.
\subsection{Every generalized-quaternion geometric closure}

\begin{statement}{A.72}{Proved}{Both cusp assertions for every generalized-quaternion closure}
Let \(Y\to U=E\setminus\{O\}\) be an actual connected finite
\'{e}tale cover, where \(E:y^2=x^3-x\), whose geometric Galois
closure has group
\[
 Q_{4n}=\langle r,s\mid r^{2n}=1,\ s^2=r^n,
                         \ srs^{-1}=r^{-1}\rangle,
 \qquad n\ge2\text{ a power of two}.                  \tag{A.72.1}
\]
Then both assertions of Task A.5 hold for \(Y\), at every
coefficient prime, every induced cusp packet, and every permitted
open derived subgroup. In particular the full fixed geometric
pro-\(\ell\) subgroup is the selected cusp inertia, and every
nontrivial arithmetically normalized procyclic pro-\(\ell\)
subgroup is cuspidal.

More precisely, the smooth proper Galois closure \(C\) has genus
\(2n-1\), four cusps, potential good Jacobian reduction of
\(2\)-rank zero at two, and every cusp difference is killed by
\(2n\). Every actual intermediate inherits the good-reduction
and torsion-divisor properties needed for the fixed-group
conclusion. This settles this family of geometric closures;
the universal question A.45 remains open.
\end{statement}

\begin{proof}
We start with the actual geometric Galois closure of the given
cover. After a finite number-field extension it and its deck
transformations, quotient maps, and cusps are defined over that
field. No identification with an unproved explicit curve model
is needed. All further reduction arguments are at a place over
two, after a finite extension making the elliptic base good,
the necessary stable models available, and all endomorphisms
defined. Local automorphisms fix the coefficient valuation ring.

\paragraph{The actual inertia subgroup and the four cusps.}
Write elements of \(G=Q_{4n}\) uniquely as \(r^a s^b\), with
\(a\in\mathbf Z/(2n)\) and \(b\in\{0,1\}\). Their product is
\[
 (a,b)(c,d)=\bigl(a+(-1)^b c+n b d,\ b+d\bmod2\bigr),
                                                        \tag{A.72.2}
\]
where the first coordinate is reduced modulo \(2n\).
The geometric fundamental group of the once-punctured elliptic
curve is free profinite on two generators; its peripheral
generator is their commutator, up to inverse and conjugacy.
Thus the image of a generating pair in this actual quotient
determines its actual cusp inertia.

A generating pair cannot consist of two rotations. If one entry
is \(r^a\) and the other \(r^b s\), it generates precisely
when \(a\) is odd. If both entries are outside
\(\langle r\rangle\), say \(r^a s,r^b s\), it generates
precisely when \(a-b\) is odd: their quotient is a rotation
with that exponent, and their squares are \(r^n\).
These assertions follow because \(2n\) is a power of two.
For the commutator convention \([x,y]=xyx^{-1}y^{-1}\),
\[
 [r^a,r^b s]=r^{2a},\qquad
 [r^a s,r^b s]=r^{2(a-b)}.                         \tag{A.72.3}
\]
Reversing the entries in the first formula takes the inverse.
Consequently every generating pair has commutator generating
the same normal subgroup
\[
                 I=\langle r^2\rangle,
 \qquad |I|=n,\qquad G/I\simeq C_2^2.             \tag{A.72.4}
\]
This is therefore the actual inertia at every cusp. There are
four cusps and their ramification indices are all \(n\).
Riemann--Hurwitz for the actual degree-\(4n\) map \(C\to E\)
gives
\[
                    2g(C)-2=4(n-1),
                    \qquad g(C)=2n-1.             \tag{A.72.5}
\]

The same elementary description proves the needed Frattini
property. The squares and commutators generate \(I\), since
all outside elements square to \(r^n\), while the square of
\(r\) is \(r^2\). Equivalently, if a subgroup \(H\) satisfies
\(HI=G\), it contains an odd rotation and an outside element,
which generate \(G\). Hence \(I=\Phi(G)\), and no proper
subgroup surjects onto \(G/I\). This includes \(n=2\), where
the group is \(Q_8\).

\paragraph{Compact type and the actual central quotient.}
The group \(G\) is nilpotent. The elliptic base has potential
good reduction at two, the sole branch section \(O\) is
\'{e}tale over the valuation ring, and generic ramification is
tame in characteristic zero. Lehr--Matignon, Theorem 6.13,
therefore gives potential good reduction of \(J_C\), or
equivalently a tree as the dual graph of the proper stable
fibre. We use exactly this ramified nilpotent-cover conclusion,
also recorded by Sa\"idi on printed page 662.
\cite{fghLehrMatignon2005,fghSaidi2007}
No slope statement for an exceptional component over a marked
branch point is inferred from it.

The unique involution of \(G\) is
\(\sigma=r^n\). It is central and lies in \(I\).
The actual quotient \(C_0=C/\langle\sigma\rangle\) is a
cover of the same fixed elliptic base, with group
\(G/\langle\sigma\rangle\simeq D_{2n}\), inertia of order
\(n/2\), four cusps, and genus \(n-1\). A.70 identifies
every such actual dihedral cover with one of its three models;
the reduction assertions of A.44 and A.69 show that
\(J_{C_0}\) has potential good reduction with special
\(2\)-rank zero. When \(n=2\), the quotient group is
\(V_4\), and \(C_0\to E\) is an elliptic isogeny, giving
the same conclusion from the supersingularity of the good
reduction of \(E\) at two.

\paragraph{The anti-invariant factor has only two possible ranks.}
Let \(A\) be the anti-invariant isogeny factor for \(\sigma\).
The actual pullback and norm maps give
\[
             J_C\sim J_{C_0}\times A,
             \qquad \dim A=n.                    \tag{A.72.6}
\]
On \(A\) we have \(r^n=-1\). Since \(n\) is a power of
two, \(x^n+1=\Phi_{2n}(x)\) is irreducible over \(\mathbf Q\).
The action is unital and \(A\ne0\), so it gives an injection
\[
             \mathbf Q(\zeta_{2n})
                       \hookrightarrow\operatorname{End}^0(A).
                                                        \tag{A.72.7}
\]
This argument uses the actual deck endomorphism, without any
assumption about a Hodge character decomposition.

After the indicated good-reduction extension, endomorphisms
extend to the abelian scheme. Thus the field
\(\mathbf Q_2(\zeta_{2n})\), of degree \(n\) over
\(\mathbf Q_2\), acts with its unit on the rational \'{e}tale
Tate space of the special \(2\)-divisible group of \(A\).
If this space is nonzero its dimension is divisible by \(n\).
Its dimension is the special \(2\)-rank \(f(A)\), at most
\(\dim A=n\). Hence
\[
                    f(A)=0\quad\hbox{or}\quad f(A)=n.
                                                        \tag{A.72.8}
\]
The second alternative means that \(A\) is ordinary. We now
exclude it by the actual action on stable components.

\paragraph{The component calculation under the ordinary alternative.}
Suppose \(A\) is ordinary, and choose the stable marked model
of \(C\) after the finite extension already described. Over
an algebraically closed residue field its graph is a tree, and
its good special Jacobian is the product of the Jacobians of
its positive-genus components. The specialization of the
\(\sigma\)-invariant factor has \(2\)-rank zero, whereas
the anti-invariant factor is ordinary.

If \(\sigma\) exchanges two positive-genus components, their
invariant diagonal and anti-invariant diagonal are both
isogenous to the same component Jacobian. A positive-dimensional
abelian variety cannot have both rank zero and be ordinary,
so no such pair occurs. A component on which \(\sigma\)
acts identically contributes only to the invariant factor;
we do not infer that its genus is zero.

Let \(B\) be a component fixed by \(\sigma\), with a
nontrivial induced involution, and put \(B'=B/\sigma\).
The actual pullback--norm identities identify the invariant
Jacobian factor with \(J_{B'}\) up to isogeny, hence
\(f(B')=0\). Its anti-invariant Prym is ordinary when
positive-dimensional. Write \(r_B\) for the number of fixed
points of \(\sigma\) on \(B\), and \(d_P\) for their
different exponents. The smooth-curve Deuring--Shafarevich
and Riemann--Hurwitz formulas give
\[
 \begin{aligned}
 f(B)&=r_B-1,\\
 g(B)-g(B')&=g(B')-1+\frac12\sum_{P\ {\rm fixed}}d_P,
                    \qquad d_P\ge2.
 \end{aligned}                                         \tag{A.72.9}
\]
Here \(r_B\ge1\), since otherwise the first equation gives
negative rank. If the Prym has dimension zero, these equations
force \(g(B')=g(B)=0\) and \(r_B=1\). If it has positive
dimension, its ordinarity implies
\(f(B)=g(B)-g(B')\). Therefore
\[
 r_B=g(B')+\frac12\sum d_P
       \quad\Longrightarrow\quad
 g(B')=0,\qquad d_P=2\text{ for every fixed point}.
                                                        \tag{A.72.10}
\]
Every positive-genus component contributing to the anti-invariant
factor is thus ordinary; its whole Jacobian is anti-invariant,
and \(\sigma\) has genus-zero quotient. The sum of the genera
of these components is exactly \(\dim A=n\).
The cited Deuring--Shafarevich formula applies to the smooth
component curves, not to an arbitrary marked vertical fibre.
\cite[Corollary 1.8]{fghCrew1984}
\cite[Proposition 1.4]{fghYang2024}

\paragraph{Why a fourth root cannot act on such a component.}
We repeat the elementary ramification argument at its required
scope. For a faithful finite \(2\)-group \(T\) acting on an
ordinary smooth proper curve \(B\), subtract twice
Deuring--Shafarevich from Riemann--Hurwitz. With branch indices
\(e_Q\) and different exponents \(d_Q\), it gives
\[
 0=2|T|\bigl(g(B/T)-f(B/T)\bigr)
      +\sum_Q\frac{|T|}{e_Q}
                         \bigl(d_Q-2(e_Q-1)\bigr).
                                                        \tag{A.72.11}
\]
All summands are nonnegative. Indeed, for wild \(2\)-group
inertia the classical lower groups have \(T_0=T_1\), and
\(d_Q=\sum_{j\ge0}(|T_j|-1)\ge2(e_Q-1)\).
The equality forces \(T_2=1\) at each branch point. The
coefficient of \(t^2\) in the local action
\(t\mapsto t+c t^2+\cdots\) embeds \(T_1/T_2\) in the
additive residue field. Thus a cyclic subgroup of order four
cannot fix a point on an ordinary curve. This uses the
classical equal-characteristic lower numbering.
\cite[\S2.2, equation (3) and Exercise 2.12(c)]{fghBouwWewers}

If an order-four automorphism \(a\) has square \(\sigma\)
and \(B/\sigma\simeq\mathbf P^1\), its induced involution
on \(\mathbf P^1\) has a unique fixed point \(Q\) in
characteristic two. This point cannot be unramified in
\(B\to B/\sigma\): on its two-point fibre the square of
any permutation is trivial, whereas \(a^2=\sigma\) would
exchange the points. Hence the fibre is a single ramified
point fixed by \(a\), contradicting (A.72.11). There is
therefore no such fourth root on an ordinary component.

\paragraph{The unique involution gives the contradiction.}
Take a positive-genus anti-invariant component \(B\), which
exists because \(\dim A=n>0\). Its stabilizer \(G_B\)
contains \(\sigma\). Its action on \(B\) is faithful:
every nontrivial subgroup of the quaternion \(2\)-group
contains its unique involution \(\sigma\), whereas
\(\sigma\) acts nontrivially on \(B\).
For completeness, a nonidentity rotation has a power equal
to \(r^n\), and every outside element has square \(r^n\);
this proves the subgroup assertion directly.

If \(G_B\) were larger than \(\langle\sigma\rangle\),
it would contain an element of order at least four. A power
of that element has order four and square \(\sigma\).
Faithfulness would give precisely the fourth root just
excluded. Thus \(G_B=\langle\sigma\rangle\), and the
orbit of \(B\) has \(2n\) components. Since \(\sigma\)
is central, all these components are anti-invariant and
have positive genus. Their genera already sum to at least
\(2n\), contrary to the total anti-invariant dimension \(n\).
This excludes the ordinary alternative in (A.72.8).
Consequently \(f(A)=0\), and \(J_C\) has special
\(2\)-rank zero. No smooth fixed-point lifting assertion
or unproved marked-fibre slope assertion enters this argument.

\paragraph{The actual elliptic quotient gives cusp torsion.}
Let \(E'=C/I\), with its actual map to \(E\). All cusp
inertia is killed in this quotient, so it extends to a
connected unramified degree-four cover of the proper
elliptic curve. Choose a point over \(O\) as its origin.
It is then an elliptic isogeny with kernel \(C_2^2\), hence
with kernel exactly \(E'[2]\). Write its four kernel points
as \(Q_1,\ldots,Q_4\), and the cusps of \(C\) above them
as \(P_1,\ldots,P_4\).

The actual map \(h:C\to E'\) has degree \(n\). At each
\(Q_j\) its inertia is the whole group \(I\), so there is
one point above it and
\[
                     h^*(Q_j)=nP_j.
                                                        \tag{A.72.12}
\]
Every \(Q_i-Q_j\) is a two-torsion divisor class on \(E'\).
Choose an actual rational function \(f_{ij}\) with
\(\operatorname{div}(f_{ij})=2(Q_i-Q_j)\), after the same
permitted finite extension. Pulling it back gives
\[
       \operatorname{div}(f_{ij}\circ h)=2n(P_i-P_j).
                                                        \tag{A.72.13}
\]
Thus the required torsion comes from the actual intermediate
elliptic map, with its total ramification, not just from an
isogeny factor of the Jacobian.

\paragraph{Every actual intermediate and the arithmetic conclusion.}
Let \(H<G\) be proper and \(D=C/H\). Because \(I\) is
normal, its cusps are the \(H\)-orbits on \(G/I\), so
their number is \([G:HI]\ge2\). Its compact genus is
positive, since it has a nonconstant finite map to \(E\).
The actual norm map makes \(J_D\) an isogeny quotient
of \(J_C\), preserving potential good reduction and special
\(2\)-rank zero. For any two cusps of \(D\), choose cusps
above them in \(C\) and take the norm of the function in
(A.72.13). Over an algebraically closed field the pushforward
of a point is its image with coefficient one, regardless
of ramification. Hence every target cusp difference is
also killed by \(2n\).

In the compact-type marked model of \(D\), these cusp
differences specialize to zero, because an abelian variety
of \(2\)-rank zero has no nonzero geometric \(2\)-power
torsion points. The extended Abel map and A.41 therefore
give a rational joining path from every selected cusp to
any distinct companion cusp. A.42 proves the complete
fixed-group assertion for every coefficient prime, including
two. A.59 gives all the normalizer assertions, since the
actual geometric Galois closure is a two-group. The quotient
by \(H=G\) is the already settled original curve \(U\).

Finally finite field extensions are used only to define
curves, deck transformations, subgroup quotients, and marked
points and to obtain reduction. A geometric path identifying
base points need not be defined over a finite extension;
the resulting comparison is up to geometric inner conjugacy,
which preserves the packet and inertia formulations. The
restriction argument in A.41--A.42 handles the original
permitted open derived subgroup, and normalizers descend by
intersecting their open arithmetic image with the field
extension. At no point is fixedness transferred through a
noninjective geometric quotient.

The exact multiplication, commutator, subgroup, cusp-count,
cyclotomic-dimension and component-orbit checks are retained
in \path{code/a_quaternion_two_power.py}. The written arguments
establish the assertion for every power of two, independently
of the finite verification range.
\end{proof}

% A.73: full root and DE proof/source/exact-certificate audits PASS.
\subsection{Every semidihedral geometric closure}

\begin{statement}{A.73}{Proved}{Both cusp assertions for every semidihedral two-group closure}
Let \(Y\to E\setminus\{O\}\) be an actual connected finite
\'{e}tale cover whose geometric Galois closure has group
\[
 G=SD_{4n}=\langle r,s\mid r^{2n}=s^2=1,
                              \ srs^{-1}=r^{n-1}\rangle,
 \qquad n\ge4\text{ a power of two}.                 \tag{A.73.1}
\]
Then both assertions of Task A.5 hold for \(Y\) at every
coefficient prime, every induced cusp packet, and every
permitted open derived subgroup. Thus the full fixed geometric
pro-\(\ell\) subgroup is its selected cusp inertia and every
nontrivial arithmetically normalized procyclic pro-\(\ell\)
subgroup is cuspidal.

The proper Galois closure has genus \(2n-1\), four cusps,
potential good Jacobian reduction of special \(2\)-rank zero
at two, and cusp differences killed by \(2n\). The proof
includes every actual intermediate. The universal question
A.45 remains open.
\end{statement}

\begin{proof}
Let \(C\to E\) be the smooth proper compactification of
the actual geometric Galois closure. Enlarge the number field
finitely to define it, its deck transformations, quotient maps
and marked points. At a place over two enlarge further so that
the original elliptic curve has good reduction and the stable
models and endomorphisms used below are defined. In particular
every local automorphism fixes the coefficient valuation ring.

\paragraph{Exact group structure and actual cusp inertia.}
Put \(k=n-1\). Every group element is uniquely \(r^a s^b\),
with \(a\bmod2n\) and \(b\in\{0,1\}\), and
\[
                (a,b)(c,d)=(a+k^b c,\ b+d\bmod2).
                                                        \tag{A.73.2}
\]
Here \(k^2\equiv1\pmod{2n}\). A pair with one rotation
and one outside element generates \(G\) precisely when
the rotation exponent is odd. A pair \(r^a s,r^b s\)
generates precisely when \(a-b\) is odd. Indeed their
quotient is \(r^{a-b}\); in the even-difference case
all available rotations have even exponent, since
\((r^a s)^2=r^{na}\). Two rotations cannot generate.

Direct multiplication gives
\[
 [r^a,r^b s]=r^{(2-n)a},\qquad
 [r^a s,r^b s]=r^{(2-n)(a-b)}.                    \tag{A.73.3}
\]
For a generating pair these exponents are twice odd integers,
because \(4\mid n\). The peripheral commutator in the
geometric fundamental group of the punctured elliptic curve
therefore gives the actual inertia subgroup
\[
             I=\langle r^2\rangle,
             \qquad |I|=n,\qquad G/I\simeq C_2^2.
                                                        \tag{A.73.4}
\]
It is normal, so the actual Galois cover has four cusps,
each of ramification index \(n\), and
\[
                       g(C)=2n-1.                 \tag{A.73.5}
\]
This follows from \(2g(C)-2=4(n-1)\). Also \(I=\Phi(G)\):
the quotient by \(I\) is elementary abelian, squares and
commutators lie in \(I\), and the square of \(r\) generates
\(I\). Explicitly, \(HI=G\) forces a subgroup \(H\) to
contain an odd rotation and an outside element, hence \(H=G\).

\paragraph{The actual old factor and the special rank dichotomy.}
The nilpotent ramified-cover theorem of Lehr--Matignon applies
to the actual \(C\to E\): the base has a good smooth model,
the sole branch section \(O\) is \'{e}tale over its valuation
ring, and ramification in generic characteristic zero is tame.
Its Theorem 6.13 gives potential good reduction of \(J_C\)
and a tree as the proper stable dual graph.
\cite{fghLehrMatignon2005,fghSaidi2007}
Only this conclusion is used; no marked-fibre slope theorem
is assumed.

The central involution is \(\sigma=r^n\). The actual
quotient \(C_0=C/\langle\sigma\rangle\) has dihedral
group \(D_{2n}\), since \(n-1\equiv-1\pmod n\).
Its four cusp indices are \(n/2\), and \(g(C_0)=n-1\).
The actual classification in A.70 and the reduction statements
of A.44/A.69 show that its Jacobian has potential good
reduction of special \(2\)-rank zero.

Let \(A\) be the anti-invariant isogeny factor for
\(\sigma\). Actual pullback and norm give
\[
                J_C\sim J_{C_0}\times A,
                \qquad \dim A=n.
                                                        \tag{A.73.6}
\]
On this nonzero factor \(r^n=-1\). Since
\(x^n+1=\Phi_{2n}(x)\), the actual deck endomorphism
therefore gives a unital field injection
\(\mathbf Q(\zeta_{2n})\hookrightarrow\operatorname{End}^0(A)\).
After good reduction this acts on the rational \'{e}tale
Tate space of the special \(2\)-divisible group. The local
field \(\mathbf Q_2(\zeta_{2n})\) has degree \(n\),
whereas that Tate space has dimension at most \(\dim A=n\).
Consequently \(A\) has special \(2\)-rank zero or is ordinary.
We exclude the ordinary alternative.

\paragraph{What ordinarity says about the actual stable components.}
Take the stable marked model of \((C,S)\), with the four
cusps marked, over the resulting complete valuation ring.
Its graph is a tree, and its good special Jacobian is the
product of the Jacobians of its positive-genus components.
Suppose \(A\) is ordinary. The invariant factor for
\(\sigma\) has rank zero. A pair of positive-genus components
exchanged by \(\sigma\) is impossible: its invariant and
anti-invariant factors are both isogenous to the same
component Jacobian. Components on which \(\sigma\) is
the identity are retained in the invariant factor.

For a component \(B\) fixed by \(\sigma\) with a
nontrivial induced involution, let \(B'=B/\sigma\).
Actual pullback--norm identities give \(f(B')=0\), and
its anti-invariant Prym is ordinary if nonzero. Write
\(r_B\) for the number of fixed points and \(d_P\) for
the different exponents. The smooth-curve formulas are
\[
 \begin{aligned}
 f(B)&=r_B-1,\\
 g(B)-g(B')&=g(B')-1+\frac12\sum_{P\ {\rm fixed}}d_P,
                   \qquad d_P\ge2.
 \end{aligned}                                         \tag{A.73.7}
\]
They are Deuring--Shafarevich and Riemann--Hurwitz on the
smooth component curves.
\cite[Corollary 1.8]{fghCrew1984}
\cite[Proposition 1.4]{fghYang2024}
There is at least one fixed point. A zero-dimensional Prym
would force \(g(B')=g(B)=0\). Otherwise ordinarity gives
\(f(B)=g(B)-g(B')\), and hence
\[
                    r_B=g(B')+\frac12\sum d_P.
\]
It follows that \(g(B')=0\) and every \(d_P=2\).
Thus every positive-genus anti-invariant component is
ordinary, with \(\sigma\) its genus-zero-quotient involution;
its entire Jacobian is anti-invariant. The genera of all such
components sum to \(n\).

An ordinary curve with this involution admits no order-four
automorphism whose square is \(\sigma\). Here is the precise
input, also proved in A.69. For an ordinary curve with a
faithful \(2\)-group action, subtracting twice
Deuring--Shafarevich from Riemann--Hurwitz makes each
nonnegative term \(d_P-2(e_P-1)\) vanish. The classical
different sum and \(G_0=G_1\) then give \(G_2=1\).
The coefficient of \(t^2\) injects \(G_1/G_2\) into the
additive residue field, so an order-four subgroup cannot
fix a point. A fourth root of \(\sigma\) induces an
involution of \(\mathbf P^1=B/\sigma\), which has one
fixed point in characteristic two. Its fibre must be
ramified: otherwise the square of its permutation on a
two-point fibre would be trivial while \(\sigma\) exchanges
those two points. Thus the fourth root fixes the unique
point above it, a contradiction.
\cite[\S2.2, equation (3) and Exercise 2.12(c)]{fghBouwWewers}

\paragraph{The only possible component orbit.}
Let \(B\) be a positive-genus anti-invariant component, and
let \(G_B\) be its stabilizer. It contains \(\sigma\).
Every rotation outside \(\{1,\sigma\}\) has a power
of order four with square \(\sigma\), so such a rotation
cannot belong to \(G_B\). Every outside element \(r^a s\)
with \(a\) odd has square \(\sigma\), so it too is
excluded. These exclusions do not presume faithful action
of the whole stabilizer: an element whose square acts as
the nontrivial involution necessarily acts with order four.
Therefore
\[
             G_B\cap\langle r\rangle=\langle\sigma\rangle,
             \qquad |G_B|\le4.                       \tag{A.73.8}
\]
The second assertion also follows from the index-two rotation
subgroup. The orbit of \(B\) has at least \(n\) components,
all with positive anti-invariant genus. Their total genus is
only \(n\). Thus there is exactly one full \(G\)-orbit of
\(n\) such components, all ordinary elliptic curves, each
with stabilizer of order four. There are no other
positive-genus anti-invariant components.

Put \(\mu=r^2\). It belongs to the normal inertia subgroup
\(I\), so it fixes all four generic cusp points individually.
For each ordinary elliptic component \(B\), (A.73.8) gives
\[
            |\langle\mu\rangle B|=n/2\ge2.
                                                        \tag{A.73.9}
\]
In particular \(\mu\) moves every component in this orbit.
On each such ordinary elliptic curve, (A.73.7) says that
\(\sigma\) has exactly two fixed points.

\paragraph{A smooth fixed point would lift to an individually fixed cusp.}
Neither of those two points can be a smooth point of the
whole special fibre. Suppose it were, and complete the
smooth local chart there as \(R[[t]]\). The local
automorphism \(\sigma\) fixes \(R\). Its restriction
to \(B\) is nontrivial, so \(\sigma(t)-t\) has nonzero
reduction in the residue power-series ring. This reduction
vanishes to a positive finite order at the fixed point.
Weierstrass preparation therefore writes the difference
as a unit times a positive-degree distinguished polynomial,
with no uniformizer factor. It gives a nonempty finite flat
horizontal fixed locus in this formal disc.
\cite[Corollary 2.16]{fghVonk2021}
After a finite extension one of its generic points is an
actual point of the smooth proper generic curve, specializing
to the chosen point; this is the usual smooth formal-disc
identification and properness.

Any generic point fixed by \(\sigma\) is a cusp: the
actual cover \(C\to E\) is \'{e}tale off those points,
so its deck action is free there. All such cusps are fixed
by \(\mu\). Their specializations must consequently be
\(\mu\)-fixed. A smooth point on \(B\) cannot be fixed
by an automorphism moving its unique component, contrary
to (A.73.9). This proves that both \(\sigma\)-fixed
points on each ordinary elliptic component are nodes.
It is an explicit finite-flat lifting argument, not an
assumption that every wild special fixed point lifts.

\paragraph{The root must be fixed by the whole group.}
Let \(\Gamma\) be the finite underlying dual tree of
the stable marked model. Its centre is invariant under
every automorphism: repeatedly delete all terminal vertices
until a single vertex or a single edge remains. Choose
this vertex as root, or the midpoint of this edge. In the
second case an abstract barycentric subdivision makes it
a vertex; this does not alter the geometric components,
their genera, their nodes or their stability. Thus the
root is fixed by \emph{all} of \(G\), not just by
\(\mu\).

All \(n\) ordinary elliptic vertices belong to one
\(G\)-orbit, so they have the same positive distance
from this root. In particular none lies below another.
The root and each such vertex are fixed by \(\sigma\),
so their unique joining path is fixed pointwise. The node
on that elliptic component toward the root is one of its
two \(\sigma\)-fixed nodes. The other is a child node.
Its descendant subtree is nonempty and contains no other
ordinary elliptic vertex, because all those vertices have
the same depth.

This child subtree contains no marked cusp. Indeed \(\mu\)
fixes every cusp vertex and the chosen \(G\)-fixed root.
An automorphism of a finite tree fixing the endpoints of a
path fixes that path pointwise. A root-to-cusp path through
the moving elliptic vertex would contradict (A.73.9).
This is why a root fixed only by \(\mu\) would be
insufficient: its two elliptic orbits would not necessarily
have the same depth.

An unmarked finite child subtree with just one attachment
to its parent cannot consist entirely of rational components.
If it had \(v\) such vertices, its total valence within
the full graph would be \(2v-1\), whereas marked stability
requires at least \(3v\). Thus it contains a positive-genus
component. The child subtrees from the \(n\) equidistant
elliptic vertices are pairwise disjoint. None contains another
one of those elliptic vertices. Summing genera therefore gives
\[
                  g(C)\ge n+n=2n,
\]
contrary to (A.73.5). This excludes the ordinary anti-invariant
factor and proves that the actual good special Jacobian has
\(2\)-rank zero.

\paragraph{Actual cusp divisors and every intermediate.}
The quotient \(E'=C/I\) is an actual connected unramified
degree-four cover of \(E\), since the whole cusp inertia
has been killed. With an origin chosen above \(O\), it is
an elliptic isogeny with kernel \(C_2^2=E'[2]\).
For its four kernel points \(Q_i\), the actual map
\(h:C\to E'\) has one point \(P_i\) above each and
\(h^*Q_i=nP_i\). Pulling back a rational function with
divisor \(2(Q_i-Q_j)\) gives
\[
                    2n(P_i-P_j)\text{ principal}.
                                                        \tag{A.73.10}
\]
This uses the actual total ramification of the intermediate
map, not a substitute Jacobian factor.

For every proper subgroup \(H<G\), the cusp number of
\(D=C/H\) is \([G:HI]\ge2\), by the normal Frattini
property. Its genus is positive because it maps finitely
and nonconstantly to \(E\). The actual Jacobian norm
makes \(J_D\) an isogeny quotient of \(J_C\), hence it
has potential good reduction and special \(2\)-rank zero.
Norms of the functions in (A.73.10) kill every cusp
difference of \(D\) by \(2n\): pushforward of a point
over an algebraically closed field has coefficient one,
independently of ramification.

The torsion-divisor criterion in A.41 gives a rational
joining path to a distinct companion for every cusp of
\(D\). A.42 gives its full fixed inertia at every
coefficient prime, and A.59 gives every normalizer assertion
for its two-group closure. The quotient \(H=G\) is the
already settled original punctured elliptic curve.
Finite extensions are used for the algebraic data and
reduction only. Base-point paths are compared up to inner
conjugacy, and A.41--A.42's restriction argument handles
the original packets and permitted open derived groups.
No quotient-to-fixed-subgroup surjectivity is asserted.

The exact group and dimension checks are retained in
\path{code/a_semidihedral_two_power.py}. The tree and local
fixed-locus arguments above establish the unbounded
conclusion independently of that finite verification range.
\end{proof}

% A.74: full root and DE proof/source audits PASS; reuses the exact A71--A73 certificates.
\subsection{Odd abelian factors preserve the three two-group families}

\begin{statement}{A.74}{Proved}{Odd abelian extensions of dihedral, quaternion and semidihedral covers}
Let \(F\) be a two-power dihedral group, a generalized
quaternion two-group as in A.72, or a semidihedral two-group
as in A.73. Let \(B\) be a finite abelian group of odd order.
Both all-prime assertions of Task A.5 hold for every actual
connected finite \'{e}tale cover of \(E\setminus\{O\}\)
whose geometric Galois closure is a quotient of \(F\times B\).
Every actual intermediate, every induced cusp packet and
every permitted open derived subgroup are included.

This is a preservation theorem for an additional abelian
factor. It makes no assertion for a product of two nonabelian
Sylow factors or for arbitrary two-groups. The universal
question A.45 remains open.
\end{statement}

\begin{proof}
We apply the actual fibre-product construction of A.71,
checking its hypotheses for all three families and retaining
its separate tree and marked-path comparisons.

\paragraph{Closure groups and actual intermediate maps.}
Every subgroup of a direct product of finite groups of coprime
orders splits into its two factor intersections. For an element
\((x,y)\), Bezout powers for the two group orders produce
\((x,1)\) and \((1,y)\) in that subgroup. Applying this to
a normal quotient kernel shows that the actual closure group
has the form \(F_0\times B_0\), with \(F_0\) a quotient
of \(F\) and \(B_0\) abelian of odd order.

Every quotient \(F_0\) is again one of the three listed
types or is abelian. For a dihedral group this follows
directly by passing its rotation-and-reflection presentation
to the quotient. For a quaternion group, every nontrivial
subgroup contains its unique involution, so every proper
quotient factors through its dihedral central quotient.
For a semidihedral group its centre is exactly
\(\{1,\sigma\}\). Every nontrivial normal subgroup of a
finite \(2\)-group meets the centre nontrivially: conjugation
of the whole group on that subgroup has nonfixed orbits of
even cardinality, so its fixed subgroup has even order and
contains more than the identity. Thus every proper
semidihedral quotient also factors through its dihedral
central quotient. These are group-theoretic statements about
the actual closure, independent of a choice of source model.

Every actual covering subgroup splits as \(H_2\times H_B\).
The quotient \(B_0/H_B\) corresponds to an abelian cover of
the punctured elliptic curve. Its boundary commutator is
trivial, so it extends to a connected unramified cover of
the proper elliptic curve. After choosing an origin over
\(O\), it is an actual odd-degree elliptic isogeny
\(\eta:E'\to E\).

Let \(Z\to E\) be the actual \(H_2\)-intermediate of
the \(F_0\) closure. The function field of the given
compact curve \(Y\) is the compositum of the function
fields of \(Z\) and \(E'\) over that of \(E\). The
normal closure of the first has two-power degree, while
the second is Galois of odd degree; their intersection is
the base field. Thus the actual pullback
\[
                         Y=Z\times_E E'
\]
after smooth proper normalization is geometrically connected.
The removed points are precisely the pairs of a cusp of
\(Z\) and a kernel point of \(\eta\). We have identified
the actual cover, rather than selected an unspecified
component of a formal fibre product.

\paragraph{The precise geometric hypothesis inherited from the three families.}
If \(F_0\) is abelian, the whole given cover extends to
an elliptic isogeny and A.40 applies. More generally if
\(Z\) has genus one, so does its unramified pullback
along \(\eta\), and A.40 again applies.

For positive genus greater than one, A.70, A.72 and A.73
give the same geometric property for every actual
\(H_2\)-intermediate: after finite extension at two it
has compact-type reduction, and each selected cusp has
a distinct companion joined by rational components,
including the marked endpoint components. In those proofs
this follows from potential-good rank-zero Jacobian
reduction, actual two-primary cusp divisors and the
extended Abel-map criterion in A.41. The existence of
this rational path, rather than rank zero by itself,
is the input to the present preservation argument.

\paragraph{The actual model extends and its pullback is semistable.}
Use the actual compact-type model \(\mathcal Z/R\) just
obtained and good elliptic models \(\mathcal E,\mathcal E'\).
As in A.71, the actual map \(Z\to E\) extends to
\(\mathcal Z\to\mathcal E\): base its extended Abel map
at a cusp mapping to \(O\), and compose with the actual
extended Jacobian homomorphism. The required relative Abel
map, including the nodes, is supplied by
\cite[Theorem 1.1 and Proposition 3.2]{fghDupuyRabinoff}.
Homomorphisms between the good abelian varieties extend
over their abelian schemes. This constructs the morphism;
no general rational-map extension on a singular surface
is being assumed.

The odd-degree isogeny extends to a finite \'{e}tale
isogeny of the good elliptic models. Its finite flat
kernel has order invertible in residue characteristic two.
Therefore
\[
             \mathcal Y=\mathcal Z\times_{\mathcal E}\mathcal E'
\]
is a proper flat semistable model of the actual connected
generic curve, finite \'{e}tale over \(\mathcal Z\).
This is precisely the local model comparison in A.71.

\paragraph{Compact type is a separate consequence.}
An \'{e}tale pullback of a compact-type curve need not
have a tree dual graph. Here the actual total Galois
closure has nilpotent group \(F_0\times B_0\), and its
only possible branch point on the good elliptic base is
\(O\). The branch closure is an \'{e}tale section and
generic ramification is tame in characteristic zero.
Lehr--Matignon Theorem 6.13 therefore gives potential
good reduction of the Jacobian of this actual closure.
\cite{fghLehrMatignon2005}
The actual norm to \(Y\) makes \(J_Y\) an isogeny
quotient, hence also potentially good. After further
finite extension the above semistable model has a tree
as its dual graph. No special \(2\)-rank claim for
\(J_Y\) is needed or made.

\paragraph{The selected marked path lifts literally.}
Choose a cusp of \(Y\) above a cusp \(P\) of \(Z\).
Choose its distinct companion \(Q\) on \(Z\), and
let \(T\) be their rational joining subtree in the
geometric special fibre. Every rational component maps
constantly to the smooth elliptic curve, and the
constants agree across its nodes. The cusp sections
map to \(O\), so all of \(T\) maps to that point.
Its pullback along the finite \'{e}tale odd isogeny is
therefore literally a disjoint union of \(\deg\eta\)
copies of \(T\), indexed by the distinct special
kernel points.

The selected cusp lies on one such copy. That same copy
contains a section over \(Q\), because the odd kernel
sections have distinct reductions after making them
rational. The two generic cusps are distinct since
\(P\ne Q\). Stabilization can contract rational
components but cannot introduce a positive-genus
component into their joining path. The separately
proved tree condition makes this the unique path
required by A.42.

Thus A.42 gives the full fixed-inertia assertion at
every coefficient prime for every cusp. A.59 supplies
all normalizer assertions for the actual nilpotent
closure. Finite extensions define only the algebraic
data and models; base-point paths are used up to inner
conjugacy, and the field-restriction argument of
A.41--A.42 preserves the original packets and open
derived subgroups. This proves the statement.

No new numerical assertion is needed for this deduction.
The exact coprime-power and curve-degree support of
A.71 is retained in
\path{code/a_dihedral_abelian_products.py}; the actual
quaternion and semidihedral group checks are retained
in their respective A.72 and A.73 certificates. The
proof above verifies the quotient and model comparisons
for all group orders without extrapolating from them.
\end{proof}

% A.75: full root and DE proof/source/exact-certificate audits PASS.
\subsection{Every modular two-group geometric closure}

\begin{statement}{A.75}{Proved}{Both cusp assertions for modular two-group closures}
Let \(Y\to E\setminus\{O\}\) be an actual connected finite
\'{e}tale cover whose geometric Galois closure has group
\[
 M_{4n}=\langle r,s\mid r^{2n}=s^2=1,\quad
                        srs^{-1}=r^{n+1}\rangle,
 \qquad n\ge4\text{ a power of two}.                 \tag{A.75.1}
\]
Then both assertions of Task A.5 hold for \(Y\) at every
coefficient prime, every induced cusp packet and every permitted
open derived subgroup. In particular its full fixed geometric
pro-\(\ell\) subgroup is the selected cusp inertia, and every
nontrivial arithmetically normalized procyclic pro-\(\ell\)
subgroup is cuspidal.

The proper Galois closure has genus \(n+1\), \(2n\) cusps,
potential good Jacobian reduction of special \(2\)-rank zero
at two, and cusp differences killed by \(2n\). These properties
pass to every actual intermediate in the precise form used below.
The universal question A.45 remains open.
\end{statement}

\begin{proof}
Let \(C\to E\) be the actual smooth proper Galois closure.
Enlarge its number field finitely to define the cover, its deck
group, quotient maps, marked points and relevant endomorphisms.
At a place over two enlarge further to obtain good reduction
of the original elliptic curve and the stable models below.
Every deck automorphism then fixes the coefficient valuation
ring. We use an \emph{unmarked proper stable model} in the
ordinary-component contradiction. Marked models enter only
at the final cusp-path comparison.

\paragraph{Actual inertia, genus and the abelian quotient.}
Write \(k=n+1\). The group consists of the pairs
\((a,b)\), with \(a\bmod2n\), \(b\bmod2\), and multiplication
\[
 (a,b)(c,d)=(a+k^b c,b+d),\qquad k^2\equiv1\pmod{2n}.
                                                        \tag{A.75.2}
\]
A rotation and an outside element generate if and only if
the rotation exponent is odd. Two outside elements
\(r^a s,r^b s\) generate if and only if \(a-b\) is odd:
their quotient is \(r^{a-b}\), and all available rotations
otherwise have even exponent. Two rotations cannot generate.
Direct multiplication gives
\[
 [r^a,r^b s]=r^{-na},\qquad
 [r^a s,r^b s]=r^{-n(a-b)}.                            \tag{A.75.3}
\]
Thus every generating pair has commutator
\(\sigma=r^n\), the central involution. The actual cusp
inertia of the punctured elliptic cover is consequently
\[
 I=\langle\sigma\rangle,\quad |I|=2,\quad
 G/I\simeq C_n\times C_2,
 \quad |S_C|=2n,
 \quad g(C)=n+1.                                      \tag{A.75.4}
\]
The genus follows from \(2g(C)-2=2n(2-1)\).
The actual quotient \(E'=C/I\to E\) is connected and
unramified, hence, after choosing an origin above \(O\),
an elliptic isogeny with kernel \(C_n\times C_2\).
In particular it has potential good supersingular reduction
at two, since it is isogenous to the original \(E\).

Here \(\Phi(G)=\langle r^2\rangle\), and only the inclusion
\(I\subset\Phi(G)\) is needed. Squares and commutators lie
in \(\langle r^2\rangle\), the quotient is elementary
abelian of rank two, and the square of \(r\) generates
that subgroup. Alternatively, \(H\Phi(G)=G\) gives two
independent parity vectors in \(H\); the generating-pair
criterion above then gives an odd rotation and an outside
element in \(H\), so \(H=G\). Thus a proper subgroup
satisfies \(HI\ne G\). We do not identify \(I\) with
\(\Phi(G)\).

\paragraph{Compact type and the two possible special ranks.}
Apply Lehr--Matignon's nilpotent ramified-cover theorem to
the actual map \(C\to E\). The base has a smooth good model,
the sole branch section \(O\) is \'{e}tale over the valuation
ring, the deck group is nilpotent, and generic characteristic
zero ramification is tame. Its Theorem 6.13 gives potential
good reduction of \(J_C\) and a tree as the dual graph
of its proper stable reduction.
\cite[Theorem 6.13]{fghLehrMatignon2005}
\cite[p.662]{fghSaidi2007}
No assertion about slopes of marked exceptional components
is inferred from that theorem.

Actual pullback and norm for the quotient by \(\sigma\)
give
\[
 J_C\sim E'\times A,\qquad \dim A=n,\qquad
                     r^n|_A=-1.                       \tag{A.75.5}
\]
Since \(X^n+1=\Phi_{2n}(X)\), the nonzero anti-invariant
factor \(A\) has a unital action of the field
\(\mathbf Q(\zeta_{2n})\). After good reduction its special
rational \'{e}tale \(2\)-divisible Tate space is a module
over \(\mathbf Q_2(\zeta_{2n})\), a field of degree \(n\).
Its dimension is the special \(2\)-rank of \(A\), and
is at most \(n\). It is therefore zero or \(n\).
The latter alternative says that \(A\) is ordinary.
We now exclude it on the actual proper stable model.

\paragraph{The positive-genus components in the ordinary alternative.}
Suppose \(A\) ordinary. The good special Jacobian is the
product of the Jacobians of the positive-genus components
of the proper stable fibre. Its \(\sigma\)-invariant
factor is the supersingular elliptic curve \(E'\), and
its anti-invariant factor is ordinary. A pair of positive-genus
components exchanged by \(\sigma\) is impossible: its
invariant and anti-invariant factors are isogenous to the
same component Jacobian. Components fixed pointwise by
\(\sigma\) contribute to the invariant factor.

Let \(B\) be a component fixed by \(\sigma\) on which
the induced involution is nontrivial, and put
\(h=g(B/\sigma)\), \(d=g(B)-h\). Pullback and norm
identify its invariant Jacobian with that of \(B/\sigma\)
up to isogeny, so the latter has \(2\)-rank zero. Its
anti-invariant factor is ordinary, giving \(f(B)=d\).
For \(r_B\) fixed points with different exponents \(d_P\),
Deuring--Shafarevich and Riemann--Hurwitz give
\[
 d=r_B-1,\qquad
 d=h-1+\frac12\sum_P d_P,
                         \qquad d_P\ge2.             \tag{A.75.6}
\]
\cite[Corollary 1.8]{fghCrew1984}
\cite[Proposition 1.4]{fghYang2024}
These are formulas on smooth component curves. They imply
\(h=0\) and every \(d_P=2\). If \(B\) has positive
genus then \(d>0\), so \(B\) is ordinary, its entire
Jacobian is anti-invariant, and \(\sigma\) has genus-zero
quotient. The genera of all such components sum to \(n\).
The remaining invariant genus is exactly one.

We recall the elementary fourth-root obstruction at its exact
scope. On an ordinary smooth curve with a faithful finite
\(2\)-group action, subtract twice Deuring--Shafarevich
from Riemann--Hurwitz. The resulting terms
\(2|H|(g(B/H)-f(B/H))\) and
\(d_P-2(e_P-1)\) are nonnegative and sum to zero.
Thus each second lower ramification group is trivial.
The coefficient of the quadratic term of a uniformizer
action embeds \(H_1/H_2\) in the additive residue field;
in particular every wild point stabilizer is elementary
abelian. If an automorphism of order four squared to our
\(\sigma\), its induced involution on \(B/\sigma=\mathbf P^1\)
would have one fixed point. That fibre cannot be unramified:
an order-four permutation of a two-point fibre has square
identity, whereas \(\sigma\) exchanges those points.
The fibre is therefore one ramified point fixed by the
order-four automorphism, contradicting the stabilizer condition.
\cite[\S2.2, equation (3) and Exercise 2.12(c)]{fghBouwWewers}

\paragraph{The only possible orbit and its fixed parent nodes.}
Every rotation other than \(1,\sigma\) has a power of
order four with square \(\sigma\). Also
\[
                    (r^a s)^2=r^{(n+2)a}.             \tag{A.75.7}
\]
As \(\gcd(n+2,2n)=2\), the only outside involutions
are \(s\) and \(\sigma s\); every other outside element
has a power of order four with square \(\sigma\).
For any positive-genus anti-invariant component \(B\),
its stabilizer therefore lies in
\[
              H_0=\{1,\sigma,s,\sigma s\}\simeq C_2^2.
                                                        \tag{A.75.8}
\]
This conclusion requires no faithful action of the whole
stabilizer: an element whose square acts as the nontrivial
\(\sigma\) necessarily acts with order four.

The full orbit has at least \(n\) positive-genus components,
while their total genus is only \(n\). Hence there is
exactly one orbit of \(n\) ordinary elliptic components,
each stabilized by \(H_0\), and no other anti-invariant
positive-genus components. The invariant genus one is thus
carried by a unique elliptic component \(B_0\), on which
\(\sigma\) acts identically. It is supersingular.
The whole group fixes its vertex, because it is the unique
component of this kind.

Root the proper stable tree at \(B_0\). All the ordinary
elliptic vertices are in a single \(G\)-orbit and have
the same positive distance from this root. The stabilizer
\(H_0\) fixes the unique path from the root to each such
vertex. In particular it fixes the parent node on the
corresponding elliptic component \(B\).

\paragraph{The stabilizer action has a kernel outside generic inertia.}
A faithful \(C_2^2\) action on an ordinary elliptic curve
cannot have a common fixed point. Indeed a fixed point
rules out a genus-one quotient by Riemann--Hurwitz, so
the quotient is rational. Deuring--Shafarevich would then
require the total contribution \(\sum_P(e_P-1)\) to
be four. A common fixed point contributes three; each
other ramified orbit contributes either two (stabilizer
of order two) or three (stabilizer of order four). The
necessary equation
\[
                         4=3+2a+3b,\qquad a,b\ge0,
                                                        \tag{A.75.9}
\]
is impossible. Applied at the fixed parent node, this
shows that the action of \(H_0\) on \(B\) has a kernel.
The kernel has order two, since \(\sigma\) acts
nontrivially. Write it as \(\langle k_B\rangle\).
It is outside \(I=\langle\sigma\rangle\), and
\(t_B=\sigma k_B\) is also outside \(I\), while its
induced action on \(B\) is precisely \(\sigma\).

Every generic point stabilizer for the actual cover
\(C\to E\) is trivial or the central inertia group \(I\).
Consequently \(t_B\) has no fixed point on the smooth
proper generic curve, including at the cusps.

\paragraph{No smooth special fixed point of this outside involution.}
Suppose a \(\sigma\)-fixed point on \(B\) were smooth
on the whole proper stable special fibre. It would also
be fixed by \(t_B\). Complete the smooth chart as
\(R[[z]]\). The automorphism fixes \(R\), and
\(t_B(z)-z\) has nonzero reduction because its action
on \(B\) is nontrivial. The reduced series vanishes to
a positive finite order at the fixed point. Weierstrass
preparation writes it as a unit times a positive-degree
distinguished polynomial, with no uniformizer factor.
It therefore defines a nonempty finite flat horizontal
fixed locus in this disc.
\cite[Corollary 2.16]{fghVonk2021}
After a finite extension a generic root is an actual
point of the smooth proper curve, by the smooth formal-disc
identification and properness. This contradicts the absence
of generic fixed points of \(t_B\).

The argument is on the \emph{unmarked proper model}.
It is not enough merely to separate a pair of marked
cusps on an inserted rational component: contracting such
a component would return the forbidden smooth fixed point.
On the ordinary elliptic \(B\), formula (A.75.6) gives
exactly two \(\sigma\)-fixed points. Both must therefore
be nodes of the proper stable model.

\paragraph{The proper stable tree now gives a contradiction.}
One of these two nodes is the parent node toward \(B_0\).
The other is a child node. Its nonempty descendant subtree
contains neither the root \(B_0\) nor any ordinary
elliptic vertex, since all the latter have equal depth.
Every positive-genus component has already been accounted
for. Thus this descendant subtree consists entirely of
rational components and has one attachment to its parent.
If it has \(v\ge1\) vertices, its total valence in the
whole tree is \(2v-1\). Proper unmarked stability requires
at least \(3v\), an impossibility. This excludes the
ordinary alternative and proves that the actual good
special Jacobian of \(C\) has \(2\)-rank zero.

\paragraph{Actual cusp functions and all intermediates.}
The kernel of \(E'\to E\) has exponent \(n\). For
any two of its points \(Q_i,Q_j\), the divisor
\(n(Q_i-Q_j)\) is principal. The actual degree-two map
\(h:C\to E'\) has one point \(P_i\) above each
\(Q_i\), with \(h^*Q_i=2P_i\). Pulling back the
principal divisor gives
\[
                     2n(P_i-P_j)\text{ principal}.
                                                        \tag{A.75.10}
\]
For every proper subgroup \(H<G\), the actual quotient
\(D=C/H\) has \([G:HI]\ge2\) cusps, since
\(I\subset\Phi(G)\). It has positive genus, as it
maps finitely and nonconstantly to \(E\). The actual
Jacobian norm makes \(J_D\) an isogeny quotient of
\(J_C\), hence gives potential good reduction and
special \(2\)-rank zero. Norms of the functions in
(A.75.10) kill every cusp difference of \(D\) by
\(2n\). Over an algebraically closed field pushforward
of a point has coefficient one, independently of ramification.

The torsion-divisor criterion of A.41 therefore supplies
a rational joining path to a distinct companion for every
cusp on the appropriate stable marked model of \(D\).
A.42 proves the full fixed-inertia assertion at every
coefficient prime, and A.59 proves all normalizer assertions
for the two-group closure. The original quotient \(H=G\)
is already covered by A.40. Finite extensions define the
algebraic data and reduction; geometric base-point paths
are compared only up to inner conjugacy. A.41--A.42's
restriction argument retains all original packets and
permitted open derived groups. No lifting of fixed elements
through a quotient is used.

The exact finite group and dimension checks are retained
in \path{code/a_modular_two_power.py}. The component,
local fixed-locus and tree arguments establish the result
for every power of two, independently of that finite range.
\end{proof}

% A.76: full root and DE proof/source/exact-certificate audits PASS.
\subsection{Cyclic-index-two closures with an odd abelian factor}

\begin{statement}{A.76}{Proved}{The whole cyclic-index-two class and its odd abelian extensions}
Let \(F\) be a finite two-group with a cyclic subgroup of
index at most two, and let \(B\) be a finite abelian group
of odd order. Every actual connected finite \'{e}tale cover
of \(E\setminus\{O\}\) whose geometric Galois closure is
a quotient of \(F\times B\) satisfies both assertions
of Task A.5 at every coefficient prime. Every actual
intermediate, every original induced cusp packet and every
permitted open derived subgroup is included.

The nonabelian possibilities for \(F\) are exactly the
dihedral, generalized-quaternion, semidihedral and modular
families of A.70, A.72, A.73 and A.75. The statement
includes the small orders explicitly below. It is a
preservation theorem for an odd abelian factor; it asserts
neither cofinality nor either universal assertion of A.45,
and does not cover a product of two nonabelian Sylow factors.
\end{statement}

\begin{proof}
\paragraph{An elementary classification, including the small orders.}
A cyclic \(F\) is abelian. Otherwise choose a cyclic
subgroup \(\langle r\rangle\) of index two and order
\(2^a\), and an element \(s\) outside it. This subgroup
is normal, and the actual extension has a presentation
\[
 r^{2^a}=1,\qquad srs^{-1}=r^k,\qquad s^2=r^h,
 \quad k^2\equiv1\pmod{2^a},\quad (k-1)h\equiv0\pmod{2^a}.
                                                        \tag{A.76.1}
\]
Indeed the listed relations reduce every word to one of
\(2^{a+1}\) words \(r^i s^j\), and the actual group
has that order. Changing \(s\) to \(r^b s\) replaces
\(h\) by \(h+(1+k)b\) modulo \(2^a\).

For \(a\ge3\), write \(n=2^{a-1}\). The only roots
of \(k^2=1\) modulo \(2^a\) are
\[
                         1,\ -1,\ 1+n,\ -1+n.
                                                        \tag{A.76.2}
\]
To see this, \(k\) is odd, the two even numbers \(k-1,k+1\)
have greatest common divisor two, and one of them has
exact two-valuation one. Their product being divisible
by \(2^a\), the other is divisible by \(2^{a-1}\).
This gives exactly the four residues in (A.76.2).
The remaining normalizations are exact:
\begin{itemize}
\item If \(k=1\), the group is abelian. If \(h\) is odd,
\(s\) generates a cyclic group of order \(2^{a+1}\).
If \(h\) is even, a change of \(s\) sets \(h=0\),
and the group is \(C_{2^a}\times C_2\).
\item If \(k=-1\), the consistency condition gives
\(h=0\) or \(h=n\). These are the dihedral and
quaternion presentations, respectively.
\item If \(k=n-1\), consistency gives \(h=0\) or
\(h=n\), because \(\gcd(n-2,2n)=2\). Changing
\(s\) by \(rs\) interchanges these values, since
\(1+k=n\). This gives the semidihedral presentation.
\item If \(k=n+1\), consistency says that \(h\)
is even. Since \(\gcd(n+2,2n)=2\), changing \(s\)
sets \(h=0\). This is the modular presentation.
\end{itemize}

For \(a=2\), the only roots are \(k=1,-1\); the
same first two cases give the abelian groups, \(D_8\)
and \(Q_8\). For \(a=1\), \(k=1\) and the groups
are \(C_4\) and \(C_2^2\). For \(a=0\), the group
has order at most two and is cyclic. Thus no small
nonabelian case is omitted, and the semidihedral and
modular families first enter when \(n\ge4\).
Every quotient again has a cyclic subgroup of index at
most two: take the image of the chosen cyclic subgroup.
These are group-theoretic statements about the actual
closure, not identifications of curves from their genera.

\paragraph{The actual connected coprime pullback.}
A subgroup of a direct product of finite groups of coprime
orders is the product of its intersections with those
factors. Indeed powers of \((x,y)\), with exponents
chosen by Bezout, give \((x,1)\) and \((1,y)\) in the
same subgroup. Apply this also to a normal quotient kernel.
The actual closure group therefore has the form
\(F_0\times B_0\), where \(F_0\) is a quotient of
\(F\) and \(B_0\) is odd abelian. The actual covering
subgroup is \(H_2\times H_B\).

The abelian cover corresponding to \(B_0/H_B\) is
unramified over the compact elliptic curve: the peripheral
commutator is killed by an abelian quotient. After choosing
an origin it is an actual elliptic isogeny
\(\eta:E'\to E\) of odd degree. Let \(Z\) be the
proper compactification of the actual \(H_2\)-intermediate
of the \(F_0\)-closure. Its field has two-power degree
over the function field of \(E\), while the isogeny
field is Galois of odd degree. They are linearly disjoint,
so their actual compositum is the field of the geometrically
connected fibre product
\[
                              Y=Z\times_E E'.          \tag{A.76.3}
\]
Proper smooth normalization is understood. The removed
points are exactly pairs of a cusp of \(Z\) and a
point of \(\ker\eta\). This identifies the actual
covering map as well as the field; it does not select
an unspecified component of a formal fibre product.

\paragraph{The geometric input retained from the four families.}
If \(Z\) has genus one, its pullback is an unramified
isogeny cover of the proper elliptic curve, so A.40 applies.
If \(F_0\) is abelian the same case holds. Otherwise
A.70, A.72, A.73 or A.75 applies to the actual cover
\(Z\), including every intermediate. Those proofs give
potential good Jacobian reduction of special two-rank zero,
actual two-primary torsion cusp differences, and hence a
stable marked model at two in which every cusp has a
distinct companion joined by rational components. The
marked endpoint components are included in this rational
path. This geometric property, not merely an already
known fixed-group equality, is the input needed here.

\paragraph{The map extends, and compact type is checked separately.}
Make one finite local extension defining the good models
and all marked points. Base the Abel map of \(Z\) at
a cusp sent to \(O\), and compose its relative extension
with the actual homomorphism \(J_Z\to E\). The
compact-type Abel-map theorem extends this morphism
across nodes and any inserted rational components.
\cite[Theorem 1.1 and Proposition 3.2]{fghDupuyRabinoff}
Homomorphisms between the good abelian varieties extend
over their abelian schemes. This constructs the extension
\(\mathcal Z\to\mathcal E\) of the actual generic map.

The odd isogeny extends to a finite \'{e}tale isogeny
of good elliptic models: its finite flat kernel has order
invertible in the residue field. Thus
\(\mathcal Y=\mathcal Z\times_{\mathcal E}\mathcal E'\)
is a proper flat semistable model of the actual connected
generic curve, finite \'{e}tale over \(\mathcal Z\).
Compact type does not follow from this pullback assertion
alone. Instead apply the nilpotent ramified-cover theorem
to the actual geometric Galois closure over \(E\),
whose group is \(F_0\times B_0\). Its only branch
section is \(O\), with \'{e}tale closure on the good
elliptic model. Lehr--Matignon Theorem 6.13 supplies
potential good reduction of the closure Jacobian.
\cite[Theorem 6.13]{fghLehrMatignon2005}
Its actual norm onto \(J_Y\) gives potential good
reduction there as well, so the semistable dual graph
of \(\mathcal Y\) is a tree. No special-rank-zero
assertion for \(J_Y\) is needed.

\paragraph{Every selected rational path lifts literally.}
Take a selected cusp above \(P\) on \(Z\) and a
distinct companion \(Q\). Let \(T\) be the connected
rational subtree joining their marked vertices in the
geometric special fibre. Each component of \(T\) maps
constantly to the good elliptic curve, and those constants
agree across the nodes. Since its endpoint sections map
to \(O\), the whole subtree maps to the special \(O\).
The fibre of the good \'{e}tale odd isogeny above \(O\)
is a discrete set of \(\deg\eta\) reduced points.
The pullback of \(T\) is therefore literally that many
disjoint copies of \(T\).

The selected marked section chooses one copy. The same
copy contains a section above \(Q\), with the same
kernel point of \(\eta\). Distinct prime-to-two kernel
sections have distinct reductions; the two generic marked
points are distinct because \(P\ne Q\). They are
joined entirely by rational components. Stabilization
can contract rational components but cannot introduce a
positive-genus component on this path. The separately
proved tree property makes this the unique joining path
required by A.42.

A.42 now proves all full fixed-inertia assertions, and
A.59 proves the normalizer assertions for the nilpotent
closure. The finite extensions define the algebraic data
and reduction only. Geometric paths are used up to inner
conjugacy; the established finite-field restriction argument
retains every original cusp packet and permitted open
derived subgroup. This is the actual preservation argument
of A.71/A.74 applied to the classified fourth family as well.
No fixed element is lifted through a quotient, and no
cofinality statement follows.

The elementary congruence and small-order checks are retained
in \path{code/a_cyclic_index_two_abelian.py}. The complete
classification and geometric preservation are proved above,
independently of the finite range of that certificate.
\end{proof}

% A.77: full root and DE proof/source/exact-certificate audits PASS.
\subsection{Closures with at most two commutator elements}

\begin{statement}{A.77}{Proved}{Both cusp assertions when the geometric commutator has order at most two}
Let \(Y\to E\setminus\{O\}\) be an actual connected finite
\'{e}tale cover whose finite geometric Galois closure group
\(G\) satisfies \(|G'|\le2\). Then both assertions of
Task A.5 hold for \(Y\), at every coefficient prime,
for every original induced cusp packet and every permitted
open derived subgroup. Every actual intermediate is included.

For a nonabelian two-group closure \(G\), its proper
Galois source \(C\) has potential good Jacobian reduction
of special two-rank zero at two. Its cusp inertia is
\(G'\), its genus is \(1+|G|/4\), and its cusp
number is \(|G|/2\). If \(e\) is the exponent of
\(G/G'\), every cusp difference is killed by \(2e\).
The rank-zero and torsion assertions here concern the
two-group source; the odd abelian extension below needs
only the actual rational-path preservation.

This includes, for example, the actual order-thirty-two
closure
\[
 \langle x,y,z\mid x^4=y^4=z^2=1,\quad
             [x,y]=z,\quad [x,z]=[y,z]=1\rangle.       \tag{A.77.1}
\]
Its full proper source has genus nine and sixteen cusps.
Its largest cyclic subgroup has order four, so it is
outside the cyclic-index-at-most-two class of A.76.
The universal question A.45 remains open.
\end{statement}

\begin{proof}
An actual geometric Galois closure is a quotient of the
rank-two free profinite group of the punctured elliptic
curve. Thus \(G\) is generated by the actual images
of two peripheral-handle generators. Pass to a finite
number-field extension defining its cover, deck group,
quotients and marked points. All later comparison maps
are actual morphisms of these curves or Jacobians.

If \(G\) is abelian, the cover compactifies to an
elliptic isogeny and A.40 applies. Suppose \(G'\simeq C_2\).
Conjugation on this group is trivial, since
\(\operatorname{Aut}(C_2)=1\). Hence \(G'\) is central,
\(G\) has nilpotency class two, and the finite group
is the direct product of its Sylow subgroups. Its odd
Sylow subgroups are abelian, since their commutators
lie in both an odd group and \(G'\). Write
\[
                            G=G_2\times B,
                                                        \tag{A.77.2}
\]
where \(B\) is odd abelian. Projection of the actual
two generators generates \(G_2\). We first prove the
stronger geometric statements for this two-group and
all its intermediates, then retain the odd factor by
an actual elliptic-isogeny pullback.

\paragraph{Every faithful irreducible quotient has a cyclic subgroup of index two.}
Let \(T\) be a two-generated finite two-group with
\(T'\simeq C_2\), and let \(\rho\) be an irreducible representation over
an algebraic closure of an auxiliary \(\mathbf Q_\ell\),
with \(\ell\ne2\). Put \(N_\rho=\ker\rho\)
and \(Q=T/N_\rho\). This is an actual finite group
quotient, with a faithful irreducible representation.
If it is abelian, it is cyclic, because a faithful
one-dimensional finite scalar image in characteristic zero is cyclic.
Suppose \(Q\) nonabelian. Then \(Q'\) is still
of order two and central. Schur's lemma embeds its
centre \(Z\) in that scalar group, so \(Z\)
is cyclic of order \(2^c\), with \(c\ge1\).

Every square in \(Q\) is central, since
\([a^2,b]=[a,b]^2=1\). The two actual generators
\(u,v\) therefore give an elementary abelian quotient
\(Q/Z\) generated by at most two elements. It cannot
be cyclic, because a group cyclic modulo its centre
is abelian. Consequently
\[
                Q/Z\simeq C_2^2,\qquad |Q|=4|Z|.
                                                        \tag{A.77.3}
\]
Let \(\sigma=[u,v]\) be its nontrivial commutator.
The class-two word collection is exact: every word
is \(u^a v^b\sigma^d\). It commutes with \(u\)
and \(v\) exactly when \(a,b\) are even. Thus
\[
                         Z=\langle u^2,v^2,\sigma\rangle.
                                                        \tag{A.77.4}
\]
If \(c\ge2\), \(\sigma\) belongs to the square
subgroup \(Z^2\). If both \(u^2,v^2\) also belonged
to \(Z^2\), (A.77.4) would give \(Z\subset Z^2\),
which is impossible. One of those squares therefore
generates the cyclic \(Z\), and its square root in
\(Q\) has order \(2|Z|\). It generates a cyclic
subgroup of index two. If \(c=1\), the same argument
works unless \(u^2=v^2=1\); in that case
\((uv)^2=\sigma\), and \(uv\) has order four and
index two in the order-eight group \(Q\).

This proves the claimed faithful-quotient property
without a classification of representations or an
assumption that the original centre is cyclic.
A.76's elementary classification identifies every
nonabelian \(Q\) with one of the already proved
cyclic-index-two families. In this class-two situation
only the relevant small dihedral/quaternion or modular
cases can occur, but that further restriction is unnecessary.

\paragraph{The quotient Jacobians are actual and jointly surject.}
Let \(C\to E\) be the proper Galois source for \(T=G_2\).
For each irreducible representation over that splitting field choose its
kernel \(N\), with repeated kernels used only once.
There are finitely many such normal subgroups. Form the
actual quotient curve and quotient map
\[
                    q_N:C\longrightarrow C_N=C/N.
                                                        \tag{A.77.5}
\]
The map \(C_N\to E\) is the actual Galois cover
with group \(T/N\), not merely a curve with an
abstractly matching representation. If that group is
abelian its proper source is an elliptic isogeny cover.
Otherwise its cyclic-index-two group falls under A.70,
A.72, A.73 or A.75. Those proofs give potential good
Jacobian reduction of special two-rank zero at two.
Thus every \(J_{C_N}\) has this property.

Consider the sum of the actual Jacobian pullbacks,
\[
       q^*:\prod_N J_{C_N}\longrightarrow J_C,
                  \qquad (a_N)_N\longmapsto\sum_Nq_N^*a_N.
                                                        \tag{A.77.6}
\]
It is surjective. To verify this, use rational Tate
homology at any auxiliary prime different from two.
Actual pullback and norm satisfy
\[
                 q_N^*q_{N*}=\sum_{h\in N}h,
          \qquad q_{N*}q_N^*=|N|.                    \tag{A.77.7}
\]
These identities already hold on divisors: the pullback of
 the image of a point counts each point in its orbit with
 its ramification index, exactly the stabilizer multiplicity
 in the sum over \(N\); the degree of every fibre is \(\lvert N\rvert\).
 They therefore descend to the actual Picard pullback and norm.
 Consequently the image of \(q_N^*\) is precisely
the \(N\)-invariant subspace of \(V_\ell J_C\).
After a splitting scalar extension, this finite-group
representation is semisimple. Every irreducible summand
is fixed by its own kernel, and that kernel is one of
the \(N\)'s in (A.77.5). The kernels were chosen in exactly this auxiliary-adic
splitting field. Their invariant images therefore span
all of \(V_\ell J_C\), by descent of a spanning
identity after scalar extension.

The abelian image of (A.77.6) has full rational Tate
homology and hence full dimension, proving surjectivity.
Its kernel may have positive dimension: we do \emph{not}
claim that the whole product is isogenous to \(J_C\).
Poincare complete reducibility instead makes \(J_C\)
an isogeny quotient or factor of that product. Potential
good reduction descends to such a quotient, and after
common good reduction the special two-rank is isogeny
invariant and additive. Thus \(J_C\) has potential
good reduction of special two-rank zero.
The quotient-good-reduction input is precisely the
one used in A.40 and A.68--A.75.
\cite[Corollary 2 to Theorem 1]{fghSerreTate1968}

\paragraph{Actual cusp inertia and functions on the two-group source.}
Since \(T\) is class two and generated by the actual
handle images \(a,b\), its commutator subgroup is
generated by \([a,b]\). It is nontrivial of order
two. The actual peripheral relation therefore gives
\[
 I=T'=\langle\sigma\rangle,
 \qquad |S_C|=|T|/2,
 \qquad g(C)=1+|T|/4.                                 \tag{A.77.8}
\]
Here \(I\) is central. The actual quotient
\(E'=C/I\to E\) is an unramified elliptic isogeny
whose kernel is \(T/I\). Let \(e\) be the exponent
of this two-group. For kernel points \(Q_i,Q_j\),
the divisor \(e(Q_i-Q_j)\) is principal on \(E'\).
The degree-two map \(h:C\to E'\) has one cusp
\(P_i\) above each \(Q_i\), with \(h^*Q_i=2P_i\).
Thus
\[
                        2e(P_i-P_j)\text{ is principal}.
                                                        \tag{A.77.9}
\]

For any proper subgroup \(H<T\), the cusp count of
the actual quotient \(D=C/H\) is \([T:HI]\).
It is at least two, because the commutator subgroup
\(I\) lies in the Frattini subgroup of the finite
two-group. Explicitly every maximal subgroup has index
two and an abelian quotient, so contains \(I\);
if \(HI=T\), a maximal subgroup containing the
proper \(H\) would give a contradiction. The genus
of \(D\) is positive because it maps finitely to
\(E\). Its Jacobian is an actual norm quotient of
\(J_C\), and hence has the same potential-good,
special-rank-zero properties. Norms of (A.77.9) kill
every cusp difference on \(D\) by \(2e\), since
pushforward of a point over an algebraically closed
field has coefficient one, independently of ramification.

A.41's torsion-divisor criterion now gives a rational
path to a distinct companion for every cusp of every
such \(D\), including the marked endpoint components.
A.42 gives all full fixed-inertia conclusions. The
quotient \(H=T\) is the original curve and is already
covered by A.40. A.59 gives all normalizer assertions
for these nilpotent closures.

\paragraph{Retaining the odd abelian factor by the actual path construction.}
Return to (A.77.2). Every covering subgroup splits as
\(H_2\times H_B\), by the coprime-power argument of
A.71. Its odd abelian part is an actual elliptic
isogeny \(\eta:E''\to E\). Let \(Z\) be the
proper \(H_2\)-intermediate of the two-group source.
The two actual function fields have coprime degrees,
with the isogeny field Galois. Their compositum is
therefore the geometrically connected actual pullback
\(Y=Z\times_EE''\), with removed points exactly
the pairs over \(O\).

If \(Z\) has genus one, A.40 applies. Otherwise the
stronger geometric input just proved supplies its
compact-type marked model at two and rational cusp
companion paths. Base the Abel map at a cusp sent
to \(O\), extend it over that model, and compose
with the actual extended homomorphism \(J_Z\to E\).
This gives the morphism of models used to pull back
\(\eta\), including all nodes and inserted rational
components.
\cite[Theorem 1.1 and Proposition 3.2]{fghDupuyRabinoff}
The good odd-degree isogeny is finite \'{e}tale, so
its pullback is a semistable model of the actual
connected source. Separately, the geometric closure
has nilpotent group \(G_2\times B\), good elliptic
base and sole branch section \(O\). Lehr--Matignon
Theorem 6.13, followed by the actual Jacobian norm,
proves that this pullback has tree dual graph.
\cite[Theorem 6.13]{fghLehrMatignon2005}

A connected rational cusp-joining subtree on \(Z\)
maps constantly to \(O\). Over it the \'{e}tale
isogeny pullback is a disjoint union of literal copies
of that subtree. The copy selected by a given marked
point also contains a distinct companion, indexed by
the same kernel point of \(\eta\). The generic
points are distinct because their \(Z\)-coordinates
are distinct. Stabilization cannot insert positive
genus into their rational path. A.42 therefore proves
all fixed-inertia assertions for \(Y\); A.59 proves
all normalizer assertions. This is the actual A.71/A.76
preservation proof with its geometric hypotheses checked,
not a transfer of fixed elements through a quotient.

All finite field extensions concern algebraic data and
reduction. Geometric paths are compared up to inner
conjugacy, and the established A.41--A.42 restriction
argument retains the original packets and open derived
subgroups.

\paragraph{The explicit order-thirty-two illustration.}
The group (A.77.1) has normal forms
\(x^i y^j z^k\), with \(i,j\bmod4\), \(k\bmod2\),
and multiplication
\[
 (i,j,k)(i',j',k')=(i+i',j+j',k+k'+ji')
                                                        \tag{A.77.10}
\]
with the stated moduli. It has centre
\(\langle x^2,y^2,z\rangle\simeq C_2^3\), commutator
\(\langle z\rangle\), and exponent four. Thus its
largest cyclic subgroup has index eight. The four
actual quotients obtained by imposing
\(x^2=z^\epsilon\), \(y^2=z^\delta\), with
\(\epsilon,\delta\in\{0,1\}\), are three dihedral
order-eight groups and one quaternion order-eight group.
Their four degree-two irreducible characters, together
with the sixteen linear characters of \(C_4^2\),
exhaust the representation degrees. This illustrates
the actual quotient-Jacobian spanning mechanism without
replacing the general argument by a finite example.

The exact group, quotient and character checks are
retained in \path{code/a_commutator_order_two.py}.
\end{proof}

% A.78: full root, DE and PBC proof/source audits PASS.
\subsection{A single horizontal branch section in mixed characteristic}

\begin{statement}{A.78}{Proved}{A single horizontal branch section and special rank zero}
Let \(R\) be a complete discrete valuation ring of mixed
characteristic \((0,p)\), with algebraically closed residue field,
and let \(X/R\) be a smooth proper geometrically connected curve
of genus at least one. Let \(x\in X_s\) be a closed point.
Suppose an actual finite geometrically connected Galois covering of \(X_K\), with
\(p\)-group \(G\), has at most one horizontal branch section
specializing to \(x\), and that this section is a smooth
\(R\)-section. After finite extension all the other branch
sections are assumed to have reductions different from \(x\).
Then every positive-genus component of a semistable source
contracted to \(x\) has \(p\)-rank zero.

If moreover \(f(X_s)=0\) and the generic cover is branched
at most at one \(R\)-section on all of \(X\), its
Jacobian has potential good reduction with special
\(p\)-rank zero. The local assertion allows other branch
sections away from \(x\); the global conclusion does not.
\end{statement}

This statement concerns contracted components. It does not
identify the genus of a marked vertical fibre or a joining
path between its marks. No slope assertion for an arbitrary
marked fibre is imported from Raynaud.

\begin{proof}[Proof of the local assertion]
If there is no branch section through \(x\), the assertion is
Raynaud's nonmarked smooth-point theorem, in the precise form
of Yang, Definition~1.5, Proposition~1.6 and Theorem~1.9
\cite{fghYang2024}. Suppose there is a branch section \(P\).
Make a finite extension to define its geometric inertia and
the roots of unity used below. The inertia at \(P_K\) is
cyclic, because this is a discrete valuation of residue
characteristic zero. Write its order as \(e=p^a\).

\paragraph{An actual global Kummer base change with a smooth local model.}
Choose a horizontal relatively ample divisor \(D\) disjoint
from \(P\) and from \(x\). For a sufficiently large multiple
of \(D\), relative Serre vanishing gives the surjection onto
the first two jets along \(P\). Choose a rational function
\(t\) which is regular near \(P\), vanishes exactly to
order one there and has unit first derivative at \(x\).
Equivalently, in a neighbourhood of \(x\), the zero divisor
of \(t\) is exactly \(P\). Such a choice may be seen by
prescribing \(t\) modulo the ideal of \(2P\), using
\(H^1(X,\mathcal O_X(nD-2P))=0\). Its other zeros and
poles have reductions away from \(x\).

Normalize the actual proper curve \(X\) in
\[
                    K(X)(z),\qquad z^e=t.                 \tag{A.78.1}
\]
The valuation at \(P_K\) makes this extension connected
of degree \(e\). After adjoining \(\mu_e\), it is cyclic.
At the unique closed point \(x'\) over \(x\), the
normalization is already regular and smooth over \(R\):
\[
 \widehat{\mathcal O}_{X,x}=R[[t]],\qquad
 R[[t]][z]/(z^e-t)=R[[z]].                               \tag{A.78.2}
\]
Here the second expression means its completed local ring.
The finite algebra before completion is regular at that
point as well. The special inverse image of \(x\) has
the single underlying point \(z=0\).

After a further finite extension choose a semistable model
\(X'\) of the new generic curve which retains this smooth
neighbourhood \emph{and an actual morphism to \(X\)}.
Here is the model comparison in detail. Under every finite
extension of the coefficient DVR the local ring in (A.78.2)
remains smooth. Resolve the proper normalization away from
this open set to obtain a regular proper model. Its main
residue extension is
\(z^{p^a}=\overline t\). Since \(d\overline t\ne0\),
this is the full purely inseparable extension of degree
\(p^a\); its smooth normalization has the same positive
genus as \(X_s\). In particular the component containing
\(x'\) cannot be contracted as a rational exceptional curve.

After a finite extension giving semistable reduction,
contract exceptional curves of the first kind to the
minimal regular model. The minimal regular model is
semistable: see \cite[Lemmas 55.8.5 and 55.14.5, Tags 0CD9 and 0CDG]{fghStacksKummerModels}.
Every contracted curve is rational. Its morphism to the
smooth positive-genus curve \(X_s\) is constant, even
in the inseparable case, by L\"uroth. Thus the actual map
to \(X\) descends through each contraction. For clarity,
this descent can be checked on its graph: the proper graph
image over the contracted regular model has one-point
fibres over the contraction centre, hence is finite
birational and is an isomorphism by normality. Away from
the centre the descent is already an isomorphism.

No such exceptional curve passes through \(x'\), which
is a smooth point of the retained positive-genus component.
Consequently this contraction procedure is an isomorphism
near \(x'\), and its inverse image of \(x\) remains
exactly \(\{x'\}\). Separating the other branch sections
and any auxiliary sections takes place away from this
neighbourhood. Alternatively, in a common regular model,
all additional components over \(x'\) are rational point-
blowup trees with one attachment and no retained marks;
the pointed stabilization contracts that entire tree.
The contraction map still descends to \(X\), by the same
constant-image argument. We obtain a pointed semistable
base with the actual properties
\[
 X'\longrightarrow X,\qquad (X'\longrightarrow X)^{-1}(x)
       =\{x'\},\qquad
 \widehat{\mathcal O}_{X',x'}=R[[z]].                    \tag{A.78.3}
\]
The resolution and common-model statements used here are
\cite[Lemma 54.4.3 and \S54.17, Tags 0C5H and 0C5Q]{fghStacksKummerModels}, or
\cite[\S9.2 and \S10.3]{fghLiu}. No assertion that
an arbitrary finite generic map extends finitely between
stable models is being made; the displayed map is proper.

\paragraph{The horizontal branch is killed only on the generic fibre.}
Take any connected component of the normalization of the
generic pullback to \(X'_K\), and denote it by \(Y'_K\).
Its cover of \(X'_K\) is Galois with a subgroup of \(G\),
hence with a \(p\)-group. At the horizontal valuation
\(P_K\), (A.78.1) has ramification index \(e\).
The characteristic-zero discrete-valuation Abhyankar lemma
therefore kills the original inertia after this base change;
see \cite[Lemma 15.116.4, Tag 0BRM]{fghStacksKummerModels}. Thus \(Y'_K\to X'_K\)
is finite \'etale throughout the generic neighbourhood
above \(x'\). All its remaining branch sections have
reductions away from \(x'\), since their images in \(X_s\)
are different from \(x\).

This is an application at the horizontal valuation, whose
residue field has characteristic zero. It does \emph{not}
claim that the mixed-characteristic surface normalization
is \'etale or regular after base change. Vertical ramification
and singularities in that normalization may remain.
They are exactly what Raynaud's nonmarked smooth-point
theorem permits.

Exclude \(x'\) from the new branch marking and extend the
actual generic cover to a pointed semistable covering of
\(X'\), using Yang Proposition~1.6. Yang Theorem~1.9
then gives
\[
                    f(Y'_{x'})=0.                        \tag{A.78.4}
\]
In particular every positive-genus component in this
vertical fibre has \(p\)-rank zero. Adding or removing
auxiliary marks changes only rational exceptional components.

\paragraph{Actual domination of each old vertical component.}
Let \(B\) be a positive-genus component of the original
semistable source which maps to \(x\). The chosen connected
generic pullback has a finite surjective map to the original
generic source. Resolve the rational map between their proper
models after resolving the nodal surfaces. By \cite[Lemma 54.4.3, Tag 0C5H]{fghStacksKummerModels} this requires
only point blowups of a regular source surface; all newly
exceptional curves are rational. The common resolving model
need not itself have reduced or semistable special fibre;
only its new exceptional curves are used here. The resulting proper
generically finite morphism is surjective, so some component
\(B'\) dominates \(B\). A rational exceptional curve
cannot dominate a positive-genus smooth curve, even
inseparably, by L\"uroth's theorem. Hence \(B'\) is the
strict transform of a positive-genus component already
present in the semistable model of \(Y'\).

The diagram of actual function fields commutes with (A.78.1).
Since \(B\) maps to \(x\), the image of \(B'\) in
\(X'\) is the unique point \(x'\) above \(x\).
Thus (A.78.4) applies to \(B'\). The induced finite morphism
between the smooth normalizations of \(B'\) and \(B\)
forces \(f(B)\le f(B')=0\). For its inseparable part,
\(p\)-rank is invariant; for its separable part, the
Jacobian norm is surjective and \(p\)-rank cannot increase
on an abelian quotient. This proves the local assertion.
\end{proof}

\paragraph{The global conclusion at a rank-zero base.}
Under the same mixed-characteristic hypotheses, suppose
\(f(X_s)=0\), and the actual \(p\)-group cover is
branched at most at one \(R\)-section \(P\).
Then its Jacobian has potential good reduction of
special \(p\)-rank zero.

\begin{proof}[Proof of the global conclusion]
Potential good Jacobian reduction, equivalently a tree
for the semistable dual graph, follows from the ramified
nilpotent-cover theorem of Lehr--Matignon,
Theorem~6.13 \cite{fghLehrMatignon2005}. Its hypotheses
are the good base and the single \'etale horizontal branch
section. The local assertion just proved treats every
component contracted to \(P_s\); Raynaud's nonmarked
theorem treats components contracted to other points.

It remains to treat a component dominating \(X_s\).
Normalize \(X\) in the actual covering function field
and restrict to the regular open surface \(X\setminus P\).
Purity places its ramification support, if nonempty, on
the single smooth irreducible vertical divisor
\(X_s\setminus\{P_s\}\). Apply Raynaud,
Proposition~1(ii), \emph{on this open surface only}
\cite{fghRaynaud1990Pgroups}. The horizontal and vertical branch
divisors meet at \(P_s\), so the proposition is not
being applied across that crossing.

For each main component, its inertia is constant on this
open part. Quotient by it gives a connected finite \'etale
cover of \(X_s\setminus\{P_s\}\) with group \(D/I\),
where \(D\) is its decomposition group. The original
residue component maps radicially to that quotient.
After smooth proper normalization, the latter is an actual
connected \(p\)-group cover \(B_0\to X_s\) branched
at most at \(P_s\). This also holds if the vertical
ramification support is empty, when the cover is already
\'{e}tale there. The finite radicial map preserves
\(p\)-rank; it need not have degree exactly \(|I|\).

Write \(d=|D/I|\), and let \(r\ge1\) be the number
of points above \(P_s\). Deuring--Shafarevich \cite[Corollary 1.8]{fghCrew1984} gives
\[
 f(B_0)-1=d(0-1)+\sum_{Q\mid P_s}(e_Q-1)
              =-d+(d-r)=-r.                            \tag{A.78.5}
\]
The same formula includes the unramified case, with
\(r=d\). Since \(f(B_0)\ge0\), necessarily \(r=1\)
and \(f(B_0)=0\). Every main component therefore has
rank zero. All component ranks and the graph Betti number
are zero, proving the stated Jacobian conclusion.
\end{proof}


% A.79: full root, DE and PBC proof/source audits PASS.
\subsection{Cyclic-derived closures at a supersingular place}

\begin{statement}{A.79}{Proved}{Cyclic-derived prime-power closures and coprime abelian factors}
Let \(p=2\) or let \(p\) be an odd prime with
\(p\equiv3\pmod4\). Let \(P\) be a finite
\(p\)-group with cyclic derived subgroup, and let
\(B\) be a finite abelian group of order prime to \(p\).
Every actual connected finite \'etale cover of
\(E\setminus\{O\}\) whose geometric Galois closure
is a quotient of \(P\times B\) satisfies both assertions
of Task A.5 at every coefficient prime. Every actual
intermediate, original induced cusp packet and permitted
open derived subgroup is included.

The reduction place in this argument is a supersingular
place of the fixed CM elliptic curve. The conclusion
neither treats arbitrary nilpotent closures nor asserts
cofinality. In particular, the cyclic-derived hypothesis
supplies an actual cusp-divisor relation which the rank-zero
theorem A.78 alone does not supply for an arbitrary group.
\end{statement}

\begin{proof}
\paragraph{The actual prime-power closure and its inertia.}
First suppose the actual Galois closure has group \(P_0\),
a quotient of \(P\). Its derived subgroup remains cyclic.
The two actual handle generators \(a,b\) generate \(P_0\).
Their commutator \(c=[a,b]\) normally generates \(P_0'\),
because the quotient killing it is abelian. Since every
subgroup of cyclic \(P_0'\) is characteristic in
\(P_0'\), the subgroup \(\langle c\rangle\) is
normal in \(P_0\). Hence
\[
                   I=\langle c\rangle=P_0'.             \tag{A.79.1}
\]
This is the actual geometric cusp inertia. In particular
it is normal. No inertia is inferred from an abstract
curve isomorphism or a dimension count.

Let \(C\to E\) be the proper closure. The actual
quotient \(E_1=C/I\to E\) is unramified on the
proper curves, and becomes an elliptic isogeny after a
choice of origin. Its kernel is \(P_0/I\). Write
\(d\) for the exponent of this kernel. Each point of
the kernel has a unique cusp of \(C\) above it,
with ramification index \(|I|\). If \(Q,Q'\) are
kernel points and \(P_Q,P_{Q'}\) their cusps, then
pullback of the actual principal elliptic divisor gives
\[
 d(Q-Q')\sim0
 \quad\Longrightarrow\quad
 d|I|(P_Q-P_{Q'})\sim0.                                \tag{A.79.2}
\]
Thus every cusp difference is killed by a power of \(p\).

At the chosen place above \(p\), the fixed elliptic
curve has potential good supersingular reduction. After
a finite local extension A.78 therefore applies to the
actual \(p\)-group cover, whose only branch section
is \(O\). It gives potential good reduction of \(J_C\)
of special \(p\)-rank zero.

\paragraph{Every actual intermediate retains companions.}
For an actual subgroup \(H<P_0\), let \(Z=C/H\).
The norm \(J_C\to J_Z\) is surjective, so potential
good reduction and special rank zero descend to \(J_Z\).
Over an algebraic closure the pushforward of a point
under \(C\to Z\) has residue degree one. Pushing
forward (A.79.2) therefore proves the same torsion
relation for any two cusps of \(Z\); no ramification
factor is inserted into this pushforward.

Since \(I\) is normal, the number of cusps of \(Z\)
is \([P_0:HI]\). Every proper subgroup of a finite
\(p\)-group is contained in a maximal subgroup, and
\(P_0'\subset\Phi(P_0)\) lies in each such maximal
subgroup. Thus \(HI\ne P_0\), and the number of
cusps is at least \(p\). The compact genus is at
least one, because there is a finite nonconstant map to
\(E\). If it is one, A.40 already applies.
Otherwise the rank-zero abelian special fibre has no
nonzero geometric \(p\)-power torsion. All cusp differences
therefore specialize to zero in the good Jacobian.
The compact-type Abel-map and rational-path argument of
A.41 identifies their marked joining paths as entirely
rational, including the marked endpoint components.
Every cusp has a distinct companion. A.42 gives all-prime
fixed inertia, and A.59 gives all normalizers. The base
quotient \(Z=E\), corresponding to \(H=P_0\), and every abelian \(P_0\)
case are the genus-one assertion A.40.

\paragraph{An actual connected coprime abelian pullback.}
Now allow the factor \(B\). Every subgroup of a direct
product of finite groups of coprime orders splits as the
product of its intersections with the factors: Bezout
powers extract \((x,1)\) and \((1,y)\) from \((x,y)\).
This applies to the normal quotient kernel and the actual
covering subgroup. Consequently the actual closure group
is \(P_0\times B_0\), with \(P_0'\) cyclic and
\(B_0\) abelian of order prime to \(p\), and its
covering subgroup is \(H_p\times H_B\).

The proper \(B_0/H_B\)-cover is an actual elliptic
isogeny \(\eta:E'\to E\) of degree prime to \(p\):
the peripheral commutator vanishes in the abelian quotient.
Let \(Z\to E\) be the actual \(H_p\)-intermediate
of the \(P_0\)-closure. Its degree is a power of \(p\).
The isogeny extension of function fields is Galois of
coprime degree, so these two fields are linearly disjoint.
Their actual compositum is therefore the field of the
geometrically connected fibre product
\[
                              Y=Z\times_E E'.           \tag{A.79.3}
\]
Proper smooth normalization is understood. The removed
points are exactly the pairs of a cusp of \(Z\) and
a point of \(\ker\eta\). If \(Z\) has genus one,
this is an unramified proper elliptic cover and A.40
applies. Otherwise the preceding argument supplies the
actual rational companion paths on \(Z\).

\paragraph{Extension of the map and a separate compact-type argument.}
Choose a finite extension defining all algebraic maps,
marked points and good reduction data. On the compact-type
marked model \(\mathcal Z\), extend the Abel map based
at a cusp using \cite[Theorem 1.1 and Proposition 3.2]{fghDupuyRabinoff}.
Compose with the actual homomorphism \(J_Z\to E\),
which extends between the good abelian schemes. This gives
an actual morphism \(\mathcal Z\to\mathcal E\)
across the nodes and the inserted rational components.
The prime-to-\(p\) isogeny extends to a finite \'etale
isogeny of the good elliptic models. Thus
\[
 \mathcal Y=\mathcal Z\times_{\mathcal E}\mathcal E'
                                                               \tag{A.79.4}
\]
is a proper flat semistable model of the actual connected
curve, finite \'etale over \(\mathcal Z\).

Its graph is not declared a tree merely from finite
\'{e}taleness. Instead apply Lehr--Matignon Theorem 6.13
\cite[Theorem 6.13]{fghLehrMatignon2005} to the actual nilpotent Galois closure over the good
elliptic model, with its sole \'etale branch section
\(O\). The closure Jacobian has potential good reduction,
and its norm quotient \(J_Y\) does too. Hence the
semistable dual graph of \(\mathcal Y\) is a tree.
No special-rank-zero assertion for \(J_Y\) is needed.

\paragraph{The selected companion path lifts as a literal copy.}
Let \(T\) be the rational subtree joining a cusp
\(P\) of \(Z\) to a distinct companion \(Q\).
Every component of \(T\) maps constantly to the good
elliptic special fibre, and those constants agree at the
nodes. Its endpoint sections map to \(O\), so its
whole image is \(O\). The fibre of the good \'etale
isogeny above \(O\) is a reduced discrete set of
\(\deg\eta\) points. Its pullback over \(T\) is
literally that many disjoint copies of \(T\).

A selected cusp of \(Y\) chooses one copy. In that
copy the lift of \(Q\) is a distinct marked point,
with the same kernel coordinate. Distinct prime-to-\(p\)
kernel sections have distinct reductions. The joining
path remains rational after marked stabilization, and
the separately verified tree property makes it the
unique path. This verifies the exact hypothesis of A.42
for each selected cusp of \(Y\). A.42 proves all
fixed-inertia conclusions and A.59 supplies normalizers
for the nilpotent closure.

This is the actual preservation argument of A.66/A.76
with residue prime \(p\) in place of two. Its changed
hypotheses have all been verified above: good reduction
of the elliptic target, prime-to-\(p\) good isogeny,
connected actual compositum, compact type independently,
and literal companion-path copies. Finite extensions
are used only to define the algebraic and reduction data.
The arithmetic comparisons are up to inner conjugacy,
with the established restriction to the permitted open
derived subgroups; no geometric path is asserted to be
defined over a finite field extension. Thus all original
cusp packets are retained. No quotient fixed element is
lifted, and no cofinality conclusion is asserted.
\end{proof}

% A.80: full root and independent proof/source/code audits PASS.
\subsection{A genus-two intermediate with a noncyclic derived closure}

\begin{statement}{A.80}{Proved}{The natural genus-two wreath cover has the required fixed groups}
Choose \(s\) with \(s^4+2s^2+2=0\), and put
\[
 i=-1-s^2,\qquad q=3+2s+2s^2,\qquad
 E_4:y^2=x^3-6x^2+x,\qquad Q=(q,-2sq).
                                                        \tag{A.80.1}
\]
Then \(2Q=R=(1,2i)\), and \(R\) has order four.
Let \(Y\) be the smooth proper normalization of
\[
 v^2=(u^2+q)^3-6(u^2+q)^2+(u^2+q).
                                                        \tag{A.80.2}
\]
Its actual map to \(E_4\) is
\(q_Y(u,v)=(u^2+q,v)+Q\). Define
\[
 \phi(x,y)=\left(\frac{(x-1)^2}{4x},
                  \frac{y(1-x^2)}{8x^2}\right),\qquad
 \eta=[1+i]\circ\phi:E_4\longrightarrow E,
 \quad E:y^2=x^3-x.
                                                        \tag{A.80.3}
\]
The map \(f=\eta q_Y\) has degree eight and is ramified only
above \(O\). Its geometric Galois closure has group
\(C_2\wr C_4=(C_2)^4\rtimes C_4\), of order sixty-four.
The affine curve \(X=Y\setminus f^{-1}(O)\) has genus two
and six cusps, with ramification indices \(2,2,1,1,1,1\).

For every original packet of Part A, every coefficient prime
\(\ell\), and every permitted open \(\Delta_0\) preserving
this actual cover and the selected cusp lift \(P\),
\[
       \pi_1(X_{\overline K})^{(\ell),s_b(\Delta_0)}
                              =I_{P,\ell}.              \tag{A.80.4}
\]
All its pro-\(\ell\) normalizer assertions follow from A.59.
At odd \(\ell\), the fixed-group proof also holds for every
packet at any cusp over a finite number field of definition,
with that field's closed derived subgroup or any open subgroup
of it. This statement concerns this actual genus-two intermediate;
it does not assert fixedness for every intermediate of the
genus-seventeen Galois closure.
\end{statement}

\begin{proof}
\paragraph{The actual elliptic isogeny and wreath closure.}
The displayed identities give \(i^2=-1\), \(2Q=R\),
and \(2R=(0,0)=T\). The rational functions in (A.80.3)
satisfy the equation of \(E\), define a degree-two isogeny
with kernel \(\{O,T\}\), and send \(R\) to \((0,0)\).
Consequently \(\eta\) is a cyclic degree-four isogeny with
\[
                     \ker\eta=K=\{O,R,T,-R\}.
                                                        \tag{A.80.5}
\]
These are identities of actual maps, not only identities of
elliptic \(j\)-invariants.

Let \(\ell_Q\) be the tangent line to \(E_4\) at \(Q\).
The function \(a=(x-1)/\ell_Q\) has divisor
\(R+O-2Q\). The connected double cover \(z^2=a\) is therefore
branched exactly at \(O,R\) and has genus two. On translating
the elliptic variable by \(-Q\), the function \(x-q\)
has the same divisor as \(a\). Their ratio is a nonzero
constant, which can be made a square after a finite number-field
extension. This identifies its smooth normalization with
(A.80.2) and its actual elliptic map with \(q_Y\).
The two remaining points of \(K\) each have two unramified
preimages, proving the stated six-cusp description.

For the full Galois closure, take the four translates
\(a_j(P)=a(P+jR)\), \(0\le j<4\). Their square classes
are linearly independent. Indeed the kernel of the cyclic
map \(\mathbb F_2[C_4]\to\langle[a_j]\rangle\) is an
ideal of \(\mathbb F_2[t]/((t+1)^4)\). Every nonzero ideal
contains \((t+1)^3=1+t+t^2+t^3\). A nonzero kernel would
make \(\prod_j a_j\) a square. But half its divisor is
\(D_K-D_{Q+K}\), whose degree-zero class on \(E_4\) is
\(-4Q=-2R=T\ne0\). Thus that product is not a square.
Adjoining the four roots gives the full \((C_2)^4\)-extension;
translation by \(R\) lifts by cyclically permuting the four
roots, literally with order four. The group over \(E\) is
therefore the stated wreath product. Its natural degree-eight
action is on four pairs; the cusp commutator exchanges one
pair in each of two adjacent blocks. This gives the displayed
cycle lengths. Riemann--Hurwitz gives genus seventeen and
thirty-two cusps for the full closure.

\paragraph{The actual complementary elliptic factor is non-CM.}
Put
\[
 A=3q-6,\qquad B=3q^2-12q+1,\qquad C=q^3-6q^2+q.
                                                        \tag{A.80.6}
\]
The two involutions \(\sigma(u,v)=(-u,v)\) and
\(h(u,v)=(u,-v)\) give the actual elliptic quotients
\(E_4\) and
\[
 F:\quad V^2=t(t^3+At^2+Bt+C),\qquad t=u^2,\quad V=uv.
                                                        \tag{A.80.7}
\]
The second quotient is by \(\sigma h\). The change
\(X=C/t,\ Y_F=CV/t^2\) gives
\[
             Y_F^2=X^3+BX^2+ACX+C^2.                  \tag{A.80.8}
\]
The two quotient pullbacks are the \(\sigma\)-eigenspaces
of proper homology. They have dimension two each, and give
an actual isogeny \(J_Y\sim E_4\times F\). In particular
the complementary factor and its projector are defined by
the displayed maps.

Here is an exact ordinary-reduction certificate for non-CM.
At the degree-one primes of \(\mathbb Q(s)\) specified below,
(A.80.8) has the following reductions:
\[
\begin{array}{c|c|c|c|c|c}
 p&s\bmod p&(B,AC,C^2)\bmod p&\#F(\mathbb F_p)&a_p&a_p^2-4p\\\hline
29&4&(15,25,24)&34&-4&-100\\
37&17&(17,21,33)&42&-4&-132
\end{array}                                                   \tag{A.80.9}
\]
Both discriminants of the cubic models are nonzero, and both
reductions are ordinary since \(p\nmid a_p\). The exact
Legendre-symbol sums for every field element are retained in
the certificate. By \cite[Theorem V.3.1, printed p.145]
{fghSilvermanAEC}, their geometric rational endomorphism
algebras are respectively \(\mathbb Q(i)\) and
\(\mathbb Q(\sqrt{-33})\). The Frobenius eigenvalue ratio
has nonzero valuation at a prime above \(p\), because the
ordinary slopes are zero and one; it is not a root of unity.
Thus every positive Frobenius power still generates the same
quadratic field. If \(F\) had CM, specialization after a
finite extension defining its endomorphisms would embed the
same quadratic CM field into both of these distinct quadratic
fields, which is impossible. Hence \(F\) is non-CM.
This proof does not infer non-CM from a numerical approximation
to \(j\), nor from a conjectural classification of CM points.

\paragraph{The proper compact-type model and its actual map.}
Work at residue characteristic two. The wreath cover is an
actual one-point Galois two-group cover of the potential-good
supersingular elliptic curve \(E\). A.78 gives potential
good Jacobian reduction with special two-rank zero for its
full closure. The norm onto \(J_Y\) is surjective; hence
the same properties hold for \(J_Y\). After finite extension,
the proper stable model of \(Y\) has a tree, and every
positive-genus component has two-rank zero. Its only possible
proper stable shapes are one smooth genus-two component or
two supersingular elliptic components joined at a node.
Indeed a stable rational vertex has valence at least three,
whereas a tree with total positive genus two has at most two
positive-genus leaves.

The actual map \(q_Y\) extends to the good elliptic model
of \(E_4\), by the extended Abel--Jacobi map and the
homomorphism of good Jacobians; see
\cite[Theorem 1.1 and Proposition 3.2]{fghDupuyRabinoff}.
The degree of the pullback of a degree-one line bundle is
two on each fibre. Rational components map constantly, so
the sum of the degrees on positive-genus components is two.
All points of \(K\) specialize to \(O\), since the good
elliptic special fibre is supersingular and \(K\) is
two-primary. The automorphism \(\sigma\) extends to the
proper stable model. It cannot act identically on every
positive-genus component: the good-Jacobian Tate comparison
would then make it the identity on proper homology, whereas
the actual quotient decomposition above has a nonzero
\(-1\)-eigenspace \(V_F\). This supplies the precise
specialization faithfulness needed below.

If the special fibre is smooth of genus two, its degree-two
map to the elliptic target is separable. A purely inseparable
degree-two map between smooth curves over the algebraically
closed residue field would preserve genus. The two generic
ramification cusps specialize to fixed points of \(\sigma\).
Deuring--Shafarevich, in
\cite[Corollary 1.8, printed pp.35--36]{fghCrew1984},
with source and quotient of two-rank
zero, gives exactly one such point. Its image is \(O\),
and it is the unique point of the special fibre above \(O\).
Thus all six cusps specialize there. After separating the
marks, their connecting paths contain only rational components.

If the two component degrees are \((1,1)\), \(\sigma\)
exchanges the components. Otherwise its restriction to each
would be the identity, since each degree-one map is an
isomorphism, contradicting faithful specialization. The
two ramification cusps specialize to the node, so its image
is \(O\). Each elliptic component has just that point
above \(O\). Again all marks lie in the rational tree
inserted over that node. In both cases A.41 applies at every
odd coefficient prime and to every selected cusp.

The remaining case has degrees \((2,0)\). Denote the main
component by \(B_E\) and the contracted tail by \(B_F\).
The map \(B_E\to(E_4)_s\) is radicial: a degree-two
separable isogeny from a supersingular elliptic curve does
not exist. Therefore \(\sigma\) is the identity on
\(B_E\), and acts nontrivially on \(B_F\). It fixes
their node, and this is its unique fixed point on the
supersingular elliptic tail, again by Deuring--Shafarevich.

If the image of the node is different from \(O\), all six
cusps specialize to the single point of \(B_E\) over
\(O\), and A.41 applies. Suppose instead that the node
maps to \(O\). The two ramification cusps must specialize
to that node, since they are fixed by \(\sigma\). Any
selected cusp with the same proper specialization has a
rational companion path. A remaining selected cusp can
specialize smoothly on \(B_F\). Choose one of the two
ramification cusps as \(Q_0\). In the pointed stable
tree the path from the selected cusp to \(Q_0\) avoids
the positive-genus vertex \(B_E\). The vertex carrying
\(Q_0\) lies in the rational tree inserted over the old
node, not on either elliptic component. Thus the distance
from \(B_E\) to that path is strictly positive.

For clarity, the local handle blocks in this last configuration
are the actual global elliptic factors. The specialization of
\((q_Y)_*\) is a nonzero isogeny on \(B_E\) and zero on
\(B_F\). Its adjoint pullback therefore identifies the
global \(V_E=V_\ell E_4\) with the first local block and
the global Prym \(V_F=V_\ell F\) with the second. This
uses the actual map and its polarization adjoint. It does
not apply arbitrary global endomorphisms to a local fibre.

\paragraph{A partial-distance kernel in every weight.}
We prove the precise extension of the calculation in A.41
needed here. Let a compact-type pointed graph have two
handle blocks \(V_E,V_F\). Root it at the selected cusp
\(P\), eliminate the generator of a distinct cusp \(Q_0\),
and retain the other cusp generators \(T\), including
\(t_P\). Give handles weight one and cusps weight two.
Assume only that \(B_E\) is off the \(P\)--\(Q_0\)
path. No condition is imposed on the location of \(B_F\).
The actual first weighted graph derivation, including the
arbitrary packet term, has the form
\[
 Dv=[W_E,v]\ (v\in V_E),\quad
 Dv=[W_F,v]\ (v\in V_F),\quad
 Dt_j=[W_j,t_j],
                                                        \tag{A.80.10}
\]
where each \(W\) is a linear combination of
\(\Omega_E,\Omega_F,T\). Its coefficient of \(\Omega_E\)
in \(W_E\) is
\[
 r_E=\sum_e m_e E_e(1-Q_{0,e})>0.                       \tag{A.80.11}
\]
This is exactly the positive distance just established.
Adding \(\lambda t_P\) to every \(W\), for arbitrary
\(\lambda\in\mathbb Q_\ell\), does not change it.
The based comparison and the absence of a weight-one inner
discrepancy are precisely A.41.4--A.41.6, using
\cite[Definition 3.1.5 and Theorem 3.1.6]{fghBettsDogra}.

Put \(k=\mathbb Q_\ell\), and inside the free associative
enveloping algebra set
\[
             B=k\langle T,V_F,\Omega_E\rangle.
                                                        \tag{A.80.12}
\]
These are free primitive generators. The component algebra
\(k\langle V_E\rangle\) has the triangular bimodule
basis A.41.8 over \(k[\Omega_E]\); take its free product
with the free algebra on \(T,V_F\). This gives a normal
basis in which every word outside \(B\) has a first
nonboundary \(E\)-handle core preceded by a \(B\)-word.
Unlike A.41, the individual \(F\)-handle letters belong
to the coefficient algebra; this change is essential.

The derivation preserves \(B\): \(\Omega_F\) already
belongs to the algebra generated by \(V_F\), and
\(D\Omega_E=[W_E,\Omega_E]\). For a fixed first core
and its following tail, the operator which raises the
weight of its initial coefficient slot by two is
\[
                         K=D_B+R_{W_E}.                \tag{A.80.13}
\]
All terms acting after that core preserve the initial-slot
weight. They cannot remove the core; the triangular
bimodule basis retains end powers of \(\Omega_E\),
and any other coefficient letter separates components.
Thus the top initial-slot operator acts independently on
each fixed core tail. We use weighted slot degree here,
since \(F\)-handle letters have weight one.

Count occurrences of the single free letter \(\Omega_E\)
in \(B\). The part of \(K\) increasing that count is
\[
       K^+=D_c+r_E R_{\Omega_E},\qquad
       D_c\Omega_E=0,\qquad
       D_c z=c_z[\Omega_E,z]\quad(z\in T\cup V_F).
                                                        \tag{A.80.14}
\]
All omitted terms preserve this count. Fix the word
\(z_1\cdots z_n\) in the other free letters and insert
powers of \(\Omega_E\) in its \(n+1\) slots. The
resulting polynomial module \(k[Z_0,\ldots,Z_n]\) has
\(K^+\) equal to multiplication by
\[
 c_{z_1}Z_0+
 \sum_{h=1}^{n-1}(c_{z_{h+1}}-c_{z_h})Z_h
                  +(r_E-c_{z_n})Z_n.                  \tag{A.80.15}
\]
For \(n=0\) this is \(r_EZ_0\). In every case the
coefficients sum to \(r_E\ne0\), so the polynomial is
nonzero and multiplication is injective. The maximum
\(\Omega_E\)-count of a finite input proves that \(K\)
is injective. The maximum initial-slot weight then proves,
in every weighted degree,
\[
       \ker D\subset B\cap\mathbb L(V_E\oplus V_F\oplus T)
                       =\mathbb L(T,V_F,\Omega_E).
                                                        \tag{A.80.16}
\]
The last equality takes primitive elements in the free
tensor Hopf subalgebra (A.80.12). This proves an unbounded
kernel inclusion; no finite-degree rank extrapolation is used.

\paragraph{The non-CM block makes every surviving fixed grade Tate.}
Choose a finite number field \(L\) defining all maps,
cusps and the model, with connected algebraic monodromy
and \(\Gamma=\overline{[G_L,G_L]}\) contained in the
prescribed open derived subgroup. Put
\(H=\pi_1(X_{\bar L})^{(\ell),s(\Gamma)}\), and take
the Lie algebra \(\mathfrak h\) of the unipotent closures
of this actual subgroup, with its induced weight filtration.
It is stable under \(G_L\), whose action on it is abelian.
The actual local Frobenius--tame-inertia commutator in
\(\Gamma\) has first weighted derivation a nonzero common
multiple of (A.80.10). Hence (A.80.16) bounds the graded
pieces of \(\mathfrak h\), exactly as in A.41.

The global proper representation is the actual direct sum
\(V_E\oplus V_F\). The first factor has CM over a finite
extension; the second is non-CM by (A.80.9). Serre's open
image theorem, \cite[Theorem 3, printed p.299]
{A_Serre1972}, shows that the connected derived algebraic
monodromy group is
\[
                         \{1\}\times\mathrm{SL}(V_F).
                                                        \tag{A.80.17}
\]
Indeed its projection to the CM torus is trivial and its
projection to \(\mathrm{GL}(V_F)\) has derived group
\(\mathrm{SL}(V_F)\). The actual \(\Gamma\)-fixed
graded pieces are therefore invariant under that full
\(\mathrm{SL}_2\).

In (A.80.16), every generator of \(T\) and \(\Omega_E\)
has weight two and character \(\chi_\ell\). A tensor
summand with \(r\) letters in \(V_F\) and \(a\) of
these Tate generators has weight \(r+2a\). Its
\(\mathrm{SL}(V_F)\)-invariants vanish for odd \(r\);
for even \(r\), \(\mathrm{GL}(V_F)\) acts on them
by the scalar \(\det^{r/2}\). One may see this directly:
the action factors through \(\mathrm{GL}_2/\mathrm{SL}_2\),
and a scalar matrix acts on the tensor power by its
\(r\)-th power. Since \(\det V_F=\chi_\ell\), the
entire invariant summand has scalar character
\(\chi_\ell^{a+r/2}\). The same holds for its primitive
Lie subspace. Consequently
\[
 \operatorname{gr}_{2m+1}\mathfrak h=0,\qquad
 G_L\text{ acts on }\operatorname{gr}_{2m}\mathfrak h
                           \text{ by }\chi_\ell^m.
                                                        \tag{A.80.18}
\]
Apply the actual-subgroup cyclotomic criterion of A.38.
It gives \(H=I_{P,\ell}\), with no lifting of arbitrary
rational Lie elements. This handles the only configuration
not already covered by rational companion paths and proves
the odd-prime assertions for every cusp packet.

For the original Part A packets at \(\ell=2\), the actual
wreath closure is a two-group, so A.22 gives the required
open embedding and exact fixed inertia. Normalizers at all
primes follow from A.59. Arithmetic comparisons use actual
cover maps after finite field extension and geometric paths
only up to inner conjugacy; no arbitrary profinite path is
asserted to be defined over a finite extension.

\paragraph{Exact certificates and scope.}
The retained programs
\path{code/a_wreath64_frontier.py},
\path{code/a_wreath64_map.py}, and
\path{code/a_wreath64_prym.py}
check the finite group/action, rational map identities,
quotient equations and both complete finite-field counts.
Their JSON and stdout are retained; the read-only aggregate replay is
\path{code/a_wreath64_genus_two.py}. The additional exact
\(j\)-polynomial and Sturm sequence are corroboration only;
the non-CM proof uses the two ordinary fields above.
The graph and all-weight kernel arguments are proved here,
independently of finite algebraic checks. The full closure
and its other intermediate fixed groups remain outside
this statement.
\end{proof}

% A.81: full root and DE proof audits PASS; exact statement precision audited.
\subsection{Double covers with a non-CM elliptic Prym}

\begin{statement}{A.81}{Proved}{A general genus-two family and its wild torsion enlargements}
Let \(E:y^2=x^3-x\), and let
\(\psi:E'\to E\) be an actual isogeny of degree
\(n=2^a\ge2\), defined over a finite number field \(L\supset\mathbb Q(i)\).
Let \(q:Y\to E'\) be an actual connected double cover
of smooth proper curves, branched at exactly two distinct
points of \(\ker\psi\). Then \(Y\) has genus two.
Let \(\sigma\) be its deck involution and \(h\) its
hyperelliptic involution. The other elliptic quotient
\[
                  q_F:Y\longrightarrow F=Y/\langle\sigma h\rangle
                                                        \tag{A.81.1}
\]
is made an elliptic curve after a finite extension and a
choice of origin. Suppose that this actual elliptic curve
\(F\) is non-CM. Equivalently the elliptic Prym of \(q\)
is non-CM. Put
\[
       f=\psi q,\qquad S=q^{-1}(\ker\psi),\qquad X=Y\setminus S.
                                                        \tag{A.81.2}
\]
The actual degree-\(2n\) cover \(X\to E\setminus\{O\}\)
has \(2n-2\) cusps and two-group geometric Galois closure.
For every coefficient prime \(\ell\), original Part A packet
and permitted open \(\Delta_0\) preserving this cover and a
selected cusp lift \(P\),
\[
 \left(\pi_1(X_{\overline L})^{(\ell)}\right)^{s_b(\Delta_0)}
                              =I_{P,\ell}.             \tag{A.81.2a}
\]
Every nontrivial closed procyclic subgroup of
\(\pi_1(X_{\overline L})^{(\ell)}\) whose arithmetic
normalizer has open image in \(G_L\) is contained in a
conjugate of cusp inertia.
At odd coefficient primes the fixed assertion also holds
for every cusp packet over a finite number field of definition
and every open subgroup of its closed derived Galois group.

In particular, keep the actual curve \(Y\), map \(q_Y\)
and cyclic four-isogeny \(\eta\) of A.80. For every
\(r\ge0\), define
\[
 \psi_r=[2^r]\eta,\qquad
 X_r=Y\setminus q_Y^{-1}(\ker\psi_r).
                                                        \tag{A.81.3}
\]
Both assertions hold at every coefficient prime for these
actual covers. Their degrees and cusp counts are
\[
              \deg(X_r/E\setminus\{O\})=2\cdot4^{r+1},
              \qquad \#(Y\setminus X_r)=2\cdot4^{r+1}-2.
                                                        \tag{A.81.4}
\]
This is an unbounded family on one genus-two compactification.
It is not asserted cofinal among finite covers. The comparison
\(X_r\subset X_0\) is an open immersion with a changed
elliptic covering map, not a finite refinement or a fixed-lift
principle for arbitrary covers.
\end{statement}

\begin{proof}
\paragraph{Actual geometric quotients and closure.}
Make one finite extension defining the maps, branch points,
involutions and all cusps. Riemann--Hurwitz for the double
cover gives \(2g(Y)-2=2\), hence \(g(Y)=2\).
The hyperelliptic involution is central and acts as \(-1\)
on proper homology. The \(\sigma\)-invariant homology is
the dimension-two pullback from \(E'\); therefore the
\(\sigma h\)-invariant part is the complementary
dimension-two eigenspace. The actual quotient in (A.81.1)
has genus one. Its pullback and that of \(q\) give an
isogeny
\[
                       J_Y\sim E'\times F.             \tag{A.81.5}
\]
The image of \(q_F^*\) is the connected kernel of the
norm to \(E'\), up to isogeny. Thus the assumed Prym is
identified by an actual quotient map, not by an unspecified
isogeny factor. Translation of origins has no effect on
the proper homology or the non-CM hypothesis.

There is also an actual two-group closure. Over an algebraic
closure write \(k(Y)=k(E')(\sqrt a)\). For every
\(T\in\ker\psi\), adjoin the square root of the
translate \(a(P+T)\) in a common algebraic closure of
\(k(E')\). Their compositum is a connected field and
is Galois over \(k(E)\): the translations permute the
quadratic subextensions, and each quadratic extension is
already normal. Its group fits into an exact sequence
with elementary abelian two-group kernel and quotient
\(\ker\psi\), a two-group. Consequently the actual
normal closure of \(Y\to E\), a subextension of this
field, has two-group Galois group. No splitting of this
extension and no independence of the translated square
classes is needed.

The isogeny is unramified in characteristic zero. Both
branch points of \(q\) map to \(O\), so \(X\to E\setminus
\{O\}\) is finite etale. Above the \(n\) points of
\(\ker\psi\), the two branch points each have one
preimage and the remaining \(n-2\) points each have two.
This proves the exact cusp count \(2+2(n-2)=2n-2\).
In particular every selected cusp has at least one distinct
possible companion.

\paragraph{Every reduction hypothesis in the A.80 argument persists.}
At a place above two, \(E'\) has potential good
supersingular reduction, because it is isogenous to \(E\).
After finite extension all points of \(\ker\psi\)
specialize to the origin of its good special fibre.
A.78 applies to the actual one-point two-group closure
just constructed over \(E\). Norm maps make \(J_Y\)
a quotient of its Jacobian, so \(J_Y\) has potential
good reduction with special two-rank zero. Its proper
stable fibre therefore has exactly the alternatives used
in A.80: one smooth genus-two component, or two
supersingular elliptic components joined at a node.

The actual degree-two morphism \(q\) extends through
the relative Abel map. The sum of its positive component
degrees is two, by the constant degree of the pullback
of a degree-one line bundle on the good elliptic target.
These are precisely the actual degree configurations
\((2)\), \((1,1)\) and \((2,0)\) in A.80.
The involution has a nonzero anti-invariant proper
homology block, so cannot become the identity on every
positive-genus component. The two ramification cusps
are actual \(\sigma\)-fixed sections and both map to
the origin after reduction. None of these assertions
depends on the number of other cusps.

For completeness, the resulting path comparison is the
following. In the smooth genus-two case, the special
degree-two map is separable and its rank-zero
Deuring--Shafarevich formula gives a unique fixed point
over the origin. In degree \((1,1)\), the involution
exchanges the components and the ramification cusps
specialize to their node, which is the only point over
the origin on either component. All cusps thus lie in
a single rational marked insertion in these cases.
In degree \((2,0)\), the main elliptic component maps
radicially to \(E'_s\), while the involution on the
contracted supersingular elliptic tail has exactly one
fixed point, its node. If that node maps away from the
origin, all cusps lie over the single main point above
the origin. If it maps to the origin, both ramification
cusps specialize to the node. A selected cusp on the
tail may have no rational companion path, but a path
to either ramification cusp avoids the main elliptic
component with strictly positive distance.

In this last case the actual \(q_*\), together with
its polarization adjoint, identifies the global
\(V_\ell E'\) block with the main component and the
global \(V_\ell F\) block with the tail. The new
partial-distance proof A.80.10--A.80.16 therefore applies
with exactly these blocks and with all \(2n-3\)
independent cusp generators retained. It imposes no
bound on that number. It gives
\[
              \ker D\subset\mathbb L(T,V_\ell F,\Omega_{E'}).
                                                        \tag{A.81.6}
\]
Every packet term is allowed in that proof. The
coefficient of \(\Omega_{E'}\) in its own block
conjugator is the same strictly positive distance, so
the all-weight slot injection is unchanged.

The first global elliptic factor is CM, being isogenous
to \(E\), while \(F\) is non-CM by hypothesis.
The actual connected derived monodromy is therefore
\(\{1\}\times\mathrm{SL}(V_\ell F)\), as in
A.80.17. Its invariant tensors in (A.81.6) have even
weight and the corresponding scalar power of the
cyclotomic character. A.38, applied to the closure of
the actual fixed subgroup, proves fixed inertia at
every odd coefficient prime. The rational-path cases
are already covered by A.41. At coefficient two,
A.22 applies to the original Part A packets through
the actual two-group closure. A.59 supplies all the
normalizer assertions. Field restriction and geometric
paths are used exactly in their inner-invariant senses
in A.80; no rational lifting of an arbitrary Lie
element is asserted.

Finally, (A.81.3) preserves the two original branch
points and the actual non-CM quotient \(F\) of A.80.
The degree of \(\psi_r\) is \(4^{r+1}\), so the
general degree and cusp formulas give (A.81.4).
This is a fresh application to each actual enlarged
cover, rather than a transfer of fixedness through
the open immersion. No new numerical premise is
used: the exact non-CM and map certificates are those
already retained and independently rerun for A.80.
\end{proof}


\begin{statement}{A.6}{Proved}{Sections agreeing on the abelian-field kernel}
Fix $s_b$ as above, and require agreement on $\Delta$ as actual homomorphisms.
Put $\Gamma=G_K/\Delta=\Gal(K^{\rm ab}/K)$. The set of such sections
is canonically a torsor under $Z^1_{\cts}(\Gamma,I)$, by
$a\longmapsto(g\mapsto a(\overline g)s_b(g))$. Their classes under
geometric conjugation are a torsor under
\[
                     H^1_{\cts}(\Gamma,I)\simeq\Z/4.     \tag{A.2}
\]
Every such section belongs to the same cuspidal packet. The four classes
are the twists of $s_b$ by the Kummer classes of $1,i,-1,-i$.
After choosing compatible generators of the Tate modules, the actual
cocycle group is parametrized by
\[
 \left(\tfrac14\Z_2\right)\times\prod_{\ell\text{ odd}}\Z_\ell,
 \qquad
 (t_\ell)_\ell\longmapsto
 \bigl[g\longmapsto((\chi_\ell(g)-1)t_\ell)_\ell\bigr]. \tag{A.3}
\]
\end{statement}
\begin{proof}
Let $s$ be another section agreeing with $s_b$ on $\Delta$, and put
$a(g)=s(g)s_b(g)^{-1}\in\pi$. For $d\in\Delta$, the section identities
and the normality of $\Delta$ give
\[
 s(g)s_b(d)s(g)^{-1}=s_b(gdg^{-1})
                 =s_b(g)s_b(d)s_b(g)^{-1}.
\]
Thus $a(g)$ centralizes $s_b(\Delta)$, and A.1 gives $a(g)\in I$.
It is a continuous cyclotomic cocycle and is zero on $\Delta$.
Since $\Delta$ acts trivially on $I$, it descends to a continuous cocycle
on $\Gamma$; continuity of the descended map follows from the quotient
topology of the compact quotient. Conversely every such cocycle gives an
actual section with the required agreement. A geometric conjugator between
two such sections centralizes their common restriction to $\Delta$, so it
lies in $I$ by A.1. Therefore precisely the usual coboundaries identify
their classes.

We compute the continuous cocycles, without replacing continuous cohomology
by finite-coefficient limits. For an abelian profinite group $\Gamma$ and a
continuous cocycle $b_\ell:\Gamma\to\Z_\ell(1)$, commutativity gives
\[
 (\chi_\ell(g)-1)b_\ell(h)
       =(\chi_\ell(h)-1)b_\ell(g).                       \tag{A.4}
\]
The cyclotomic image of $G_{\Q(i)}$ is
\[
 (1+4\Z_2)\times\prod_{\ell\text{ odd}}\Z_\ell^\times.
\]
This follows by restricting the cyclotomic extension
$\Q(\mu_\infty)/\Q$ to its subfield $\Q(i)$.
For odd $\ell$ choose $g$ with $\chi_\ell(g)-1$ a unit; for $\ell=2$
choose $g$ with $\chi_2(g)=5$. Equation (A.4), in the torsion-free module
$\Z_\ell(1)$, gives exactly
$b_\ell(h)=(\chi_\ell(h)-1)t_\ell$, with
$t_\ell\in\Z_\ell$ for odd $\ell$ and $t_2\in\tfrac14\Z_2$.
Every such formula is continuous and integral-valued. Continuity into the
product is equivalent to continuity of every coordinate, so there is no
additional convergence condition in (A.3). Coboundaries correspond to
$(t_\ell)_\ell\in\prod_\ell\Z_\ell$. Their quotient is
$(\tfrac14\Z_2)/\Z_2\simeq\Z/4$.

For identification with roots of unity, choose compatible $2^n$-th roots of
$i$. Their Kummer cocycle is
$(\chi_2-1)/4$, which represents a generator of the displayed quotient;
the odd components are coboundaries. Thus the four classes are exactly
those stated. This calculation also shows why passage to an arbitrary
completed tangent parameter introduces no extra same-$\Delta$ classes.
\end{proof}

\begin{statement}{A.8}{Proved}{A finite cover detects every nonidentity element}
Let $G$ be any profinite group and $u\ne1$. There exist an open subgroup
$H\subset G$ containing $u$ and a prime $\ell$ such that the image of
$u$ in the maximal pro-$\ell$ quotient $H^{(\ell)}$ is nontrivial.
If $G$ is free profinite of finite rank, its image in
$H^{\rm ab}\otimes_{\wh\Z}\Z_\ell$ is already nontrivial.
This proves Lead 4.
\end{statement}
\begin{proof}
Choose a finite continuous quotient $q:G\twoheadrightarrow D$ detecting $u$.
Let $H=q^{-1}(\langle q(u)\rangle)$ and choose a prime divisor $\ell$
of the order of $q(u)$. Projection of the cyclic group
$\langle q(u)\rangle$ onto its nontrivial $\ell$-primary factor defines
a finite abelian $\ell$-quotient of $H$ detecting $u$. It factors through
both the pro-$\ell$ quotient and the indicated abelianization. No
cofinality assertion about the original group's pro-prime quotients is used.
\end{proof}

\begin{statement}{A.9}{Proved}{Lead 1 and the full fixed groups in covers}
In the notation of Lead 1, $\operatorname{Fix}=I$ and $\Theta=1$.
Consequently the asserted semidirect product is the trivial-kernel one,
and its freeness assertion holds in rank zero. If an open subgroup
$H\subset\pi$ is normalized by $s_b(G_{K_H})$, then
\[
                  H^{s_b(\Delta_H)}=H\cap I.             \tag{A.5}
\]
In particular, a $\Delta_H$-fixed element of $H^{(\ell)}$ has a fixed
lift in $H$ if and only if it is in the image of $H\cap I$.
\end{statement}
\begin{proof}
A.1 gives the fixed group and hence $\Theta=I\cap\Omega=1$:
each prime component of inertia injects into its pro-prime image, since
the boundary commutator has nonzero primitive degree-two class.
For (A.5), the subgroup $\Delta_H=G_{K_H}\cap\Delta$ contains the closed
commutator subgroup of $G_{K_H}$. The proof of A.1 over $K_H$ therefore
gives $\pi^{s_b(\Delta_H)}=I$. Intersect with $H$.
The assertion about a fixed lift follows by applying the quotient map to
(A.5); the converse uses the actual fixed elements of $H\cap I$.
It does not identify all fixed elements of $H^{(\ell)}$ and makes no
claim that fixed points commute with taking a quotient.
\end{proof}

\subsection{The residual full-group problem and the local kernels}

\begin{statement}{A.2}{Open}{The full profinite procyclic normalizer criterion}
We do not prove or refute that every nontrivial closed procyclic
$A\subset\pi$ with open arithmetic normalizer is cuspidal.
The full fixed-group theorem A.1 does not supply a section of that
normalizer. The residual possibility is contained in
$R_{\rm nil}(\pi')$, and the prime-complement configuration of
LM6 A.62 remains an embedding problem inside the actual arithmetic
curve group. The first-cover results A.11--A.14 rule out the
specified additional pro-prime detections but do not form a cofinal
family of all finite covers.
\end{statement}
\begin{proof}[Attempt and remaining obstruction]
The strengthened finite-cover argument is A.15 below. The structural
reduction has $N_0=P\rtimes B$, geometric kernel $P\times C$,
$P\simeq\Z_\ell$, and $C$ pro-prime-to-$\ell$. It supplies a
section of $N_0/C$, not one of $N_0$. The explicit nonsplit abstract
extensions in LM6 A.63 show why prime support alone does not close
this step. The action on $P$ is constrained by the CM-character
argument of LM6 A.53, but that is a statement about the quotient
action and does not annul the extension class of $B$.

We also tried a direct local valuation argument. At a good place
of residue characteristic different from $\ell$, the preimage in
$B$ of a pro-$\ell$ Sylow subgroup of the local Galois group has a
pro-$\ell$ Sylow complement mapping isomorphically to that subgroup:
its kernel is pro-prime-to-$\ell$. This gives an actual local
section on that Sylow subgroup. Tame inertia acts trivially on the
unramified CM character of $P$. The uniqueness theorem used in A.1
then puts $P$ in the stabilizer of the resulting primitive valuation.
Here, however, no point or cusp supplied a horizontal valuation in
advance. A valuation on a reduction component can have a nontrivial
geometric stabilizer. Thus this does not identify a cusp.
At the primes above $\ell$, splitting over the cyclotomic kernel
is also insufficient: a section over that infinite-extension group
need not be a point or cusp section over the local number field.

The next concrete attack is to use the embedding of the same
$P\times C$ in all finite cover groups, together with these local
valuation stabilizers, to force the compatible extension classes
of $B$ to vanish. Neither a compatible family of such vanishings
nor a configuration embedded in $\Pi$ has been constructed here.
\end{proof}

\begin{statement}{A.7}{Open}{The remaining assertion in R27}
The claim that every configuration $(P,C,B)$ in the actual curve
group virtually splits remains open in this return. The residual
and local-kernel statements below, and the audited CM restrictions
and prime-complement construction, do not prove that claim.
The full-group local section theorem proved here in D.56
does imply the cuspidality of a virtually split normalizer; this
is a different direction from the required virtual splitting.
The alleged density gap in that preprint's Section 8 has the
source mismatch described in A.20.
\end{statement}
\begin{proof}[Scope of the audit]
The proofs of LM6 A.36, A.50--A.53, A.61--A.65 were read in the
attached file \texttt{LM6\_POINTS\_AS\_SECTIONS\_paper.tex}.
The arithmetic commutant arguments A.54--A.60 were independently
checked, including fresh finite-field computations; their results
belong to Part B of this return. We do not treat a statement-only
intake report as evidence for a proof. The independent arguments
and the exact source hypotheses used in this audit are recorded
in the following statements and in the work report.
\end{proof}

\begin{statement}{A.15}{Proved}{The residual group in LM6 A.36}
Let $P_n=\pi'\pi^n$ be the geometric group of the multiplication
cover $E\setminus E[n]\to E\setminus\{O\}$. Inside $\pi'$,
\[
 R_{\rm nil}(\pi')
  =\bigcap_{n\ge1,\ \ell}\ker(P_n\longrightarrow P_n^{(\ell)}).
                                                               \tag{A.15.1}
\]
Any noncuspidal closed procyclic subgroup with open arithmetic
normalizer is contained in this group, which lies in $\pi''$
and is strictly smaller than $R_{\rm nil}(\pi)$.
\end{statement}
\begin{proof}
If the subgroup has nonzero image in some $\pi^{(\ell)}$, the
criterion for the CM curve puts that image in inertia. Its
nontrivial $\ell$-component maps injectively to this procyclic
image, so its full normalizer acts cyclotomically on that
component. Nakamura's full profinite recognition theorem
\cite[Theorem 2.1.1]{Nakamura1994} puts that component in actual
full inertia. The full centralizer of any nonidentity inertia
element is that inertia; hence the whole procyclic subgroup is
cuspidal. This argument uses the theorem for the full class of
finite groups, not a lifting of conjugacy from a pro-prime quotient.

In the remaining case $A\subset R_{\rm nil}(\pi)\subset\pi'$,
so $A\subset P_n$ for all $n$. Suppose it survives in
$P_n^{(\ell)}$. Its proper homology maps to zero under the map
to $\pi^{(\ell)}$. On the proper Tate modules that map is
multiplication by $n$, which is injective even if $\ell\mid n$.
Thus its proper homology was zero. By A.11 its nonzero pro-$\ell$
image must be cuspidal, and the same full recognition and
centralizer argument gives a contradiction. This proves the
containment in the right side of (A.15.1).

Every map $P_n\to P_n^{(\ell)}$ kills $R_{\rm nil}(\pi')$ on
restriction to $\pi'$. Conversely, if an element of $\pi'$ is
not in $R_{\rm nil}(\pi')$, it survives in a finite quotient
$G$ of $\pi/R_{\rm nil}(\pi')$. This quotient exists because
$R_{\rm nil}(\pi')$ is characteristic in $\pi'$ and normal in
$\pi$. The image $G'$ is nilpotent. Take $n$ divisible by the
exponent of $G/G'$. The image of $P_n$ is exactly $G'$; a Sylow
projection of $G'$ detects the selected element. This proves
(A.15.1). Its containment in $\pi''$ follows because finite
abelian quotients are among the nilpotent quotients of $\pi'$.

For strictness use $G=\mathbb F_{25}\rtimes C_3$, with $C_3$
acting by a primitive cube root in $\mathbb F_{25}^{\times}$.
This is a two-generated quotient of $\pi$.
The executed program \texttt{a\_normalizer\_certificates.py}
checks its order $75$, $|G'|=25$, $G''=1$, and
$[G',G]=G'$. Thus its nilpotent residual is $G'\ne1$ while the
nilpotent residual of $G'$ is trivial. The image of a profinite
nilpotent residual in a finite quotient is the finite nilpotent
residual: use the descending closed lower-central series and
compactness to commute its intersection with the finite image.
Consequently the two residuals of $\pi$ in question are distinct.
\end{proof}

\begin{statement}{A.16}{Proved}{The CM and prime-support restrictions in LM6 A.52--A.53}
Let $A$ be a noncuspidal procyclic subgroup of the full geometric
group of a number-field hyperbolic curve, with normalizer image
$G_L$. Then $L$ contains a CM subfield. If $A_\ell\ne1$ and
$M/\Q$ is finite Galois containing $L$, some CM subfield of $M$
embeds into $\Q_\ell$. In particular, for imaginary quadratic
$L$ such a prime must split in $L$.
If $\ell$ is unramified in $M$ and its Frobenius conjugacy class
is that of a complex conjugation, then $A_\ell=1$. These excluded
primes have positive Dirichlet density. Thus a subgroup with
cofinite prime support, in particular one isomorphic to
$\widehat\Z$, must be cuspidal. The same is true with arbitrary
prime support if $L$ is totally real.
\end{statement}
\begin{proof}
We include the character descent that is essential in the full
group. Put $D=N_\Pi(A)\cap\pi$. For a nonzero component
$A_\ell\simeq\Z_\ell$, let $\theta$ denote its normalizer
character. Every pro-$\ell$ subgroup of $D$ containing $A_\ell$
is free pro-$\ell$ in the affine case. Its conjugation character
on $A_\ell$ is trivial: the first nonzero Magnus term of a
generator is unchanged by an inner automorphism, and cancellation
in its torsion-free graded piece forces exponent one. It follows
that the pro-$\ell$ Sylow subgroup of $\theta(D)$ is trivial.
Since $\Z_\ell^\times$ has prime-to-$\ell$ torsion of exponent
$\ell-1$ for odd $\ell$, and is pro-$2$ at two,
\[
 \lambda=\theta^d:G_L\longrightarrow\Z_\ell^\times,
 \qquad d=\ell-1\ (\ell>2),\quad d=1\ (\ell=2)
                                                               \tag{A.16.1}
\]
is a well-defined continuous character. For proper curves the
same argument is applied to each finitely generated pro-$\ell$
subgroup, using A.23 below.

Detect a generator of $A_\ell$ in proper homology of a finite
cover, using Nakamura's Lemma 2.1.4. Such a detector exists
because otherwise $A_\ell$, and hence $A$, would be cuspidal.
Intersect its arithmetic group with the normalizer and replace
the number field by the resulting finite extension $L'/L$.
The detected rational line has a character $\psi$ with
\[
                    \lambda|_{G_{L'}}=\psi^d.
                                                               \tag{A.16.2}
\]
It has Frobenius weight one, and its Hodge--Tate exponent relative
to the cyclotomic character is either zero or one at each
embedding. Hodge--Tate representations descend through finite
extensions, so $\lambda$ itself is Hodge--Tate; see
\cite[Chapter III, Appendix, Corollary 2]{A_Serre}.
The one-dimensional local algebraicity theorem and the global
unit relations in \cite[Chapter III, \S1.2; Chapter II, \S3.1]{A_Serre}
therefore apply. The character is unramified outside finitely
many primes, as is seen after the finite extension from the
Jacobian representation.

Write $G=\Gal(M/\Q)$, $H=\Gal(M/L)$. The algebraic type of
$\lambda|_{G_M}$ is an integer-valued function $n:G\to\Z$,
right $H$-invariant by norm pullback. The global unit relation
says that, with $a=|G|^{-1}\sum_g n(g)$ and $b=n-a$,
\[
                b(cg)=-b(g)
       \quad\text{for every complex conjugation }c.
                                                               \tag{A.16.3}
\]
The function $b$ cannot be zero. If it were, $a$ would be an
integer and the character would be finite times $\chi_\ell^a$:
the quotient has trivial algebraic type, is trivial on open
local units at $\ell$, and factors through a finite ray class
group. Kill the finite character and compare (A.16.2).
Frobenius absolute values give $2a=d$, whereas Hodge--Tate
exponents give $a\in\{0,d\}$, impossible for $d>0$.

Let $J\lhd G$ be generated by all products of two complex
conjugations. Equation (A.16.3) makes $b$ left $J$-invariant,
hence also right $J$-invariant by normality. It is therefore
right $JH$-invariant. Modulo $J$, all complex conjugations have
the same central image. None lies in $JH$, since it would then
both fix and negate the nonzero function $b$. Consequently
$M^{JH}\subset L$ is totally imaginary with a common
complex-conjugation involution and totally real fixed field;
it is a CM field.

For the local assertion take the left stabilizer
$P=\operatorname{Stab}_G(b)$. It contains $J$ and contains no
complex conjugation, so the same argument makes $M^P$ CM.
The algebraic type is defined over $\Q_\ell$; hence the local
decomposition subgroup $D_v\subset G$ fixes $b$ and lies in
$P$. In the selected completion, $M^P$ therefore embeds in
$M_v^{D_v}=\Q_\ell$. This includes ramified primes. If $L$ is
imaginary quadratic, take $M=L$; it embeds in $\Q_\ell$ exactly
at its split primes.

If Frobenius at an unramified $\ell$ is complex conjugation,
choose its prolongation so $D_v$ contains that conjugation.
Then (A.16.3) and $D_v$-invariance force $b=0$, contradiction.
Chebotarev gives density $|\operatorname{Cl}(c)|/|G|>0$.
Finally a totally real field contains no CM subfield. These
observations prove all the remaining assertions.
\end{proof}

\begin{statement}{A.17}{Proved}{The prime-complement configuration and the direction supplied by the local theorem}
For a noncuspidal normalizer in the affine setting there are,
after replacing it by an open subgroup and the field by a
finite extension, closed subgroups satisfying
\[
 \begin{gathered}
 N_0\cap\pi=P\times C,\qquad P\simeq\Z_\ell,
 \qquad 1\ne P\subset A_\ell,\qquad C\text{ pro-prime-to-}\ell,\\
 N_0=P\rtimes B,\qquad 1\to C\to B\to G_L\to1.
 \end{gathered}                                         \tag{A.17.1}
\]
The construction of (A.17.1) uses no section conjecture.
The full-group case of D.56 implies that a procyclic
normalizer is virtually split if and only if it is cuspidal.
Thus in (A.17.1) $C\ne1$ and $B\to G_L$ has no section
after any finite extension. This verifies precisely the two
structural and local assertions of LM6 A.61--A.62.
\end{statement}
\begin{proof}
First shrink the normalizer so that its geometric kernel $D$
centralizes $A_\ell$, using the finite character image in
(A.16.1). A pro-$\ell$ Sylow subgroup $P^*$ of $D$ contains
the normal subgroup $A_\ell$ and is free pro-$\ell$. Its
nontrivial center makes it procyclic; thus $A_\ell$ is open
in $P^*$. The compact set $P^*\setminus A_\ell$ avoids one.
Choose an open normal subgroup $N_1$ of the normalizer avoiding
this set. Then $P_1=P^*\cap N_1$ is a nonzero central Sylow
subgroup of $D_1=D\cap N_1$, so it is the unique Sylow subgroup.
At every finite quotient a central Sylow subgroup has a unique
prime-to-$\ell$ complement: coprime cohomology gives existence,
and the absence of homomorphisms into the Sylow subgroup gives
uniqueness. These complements are compatible. Hence
$D_1=P_1\times C_1$, and uniqueness makes $C_1$ normal in $N_1$.

Choose a finite cover detecting $P_1$ in homology and intersect
its arithmetic group with $N_1$. Closed subgroups of a product
of a pro-$\ell$ and a pro-prime-to-$\ell$ group split into
their primary factors, by taking limits of integer powers of
each element. We retain the form $D=P\times C$ with $P\ne1$.
Its homology image must have nonzero proper projection:
otherwise it would be cyclotomic and full inertia recognition
would make $A$ cuspidal. Thus $P$ injects equivariantly into
$T_\ell J$ for the compactified cover's Jacobian.
Tate semisimplicity gives a rational equivariant retraction.
Clearing denominators gives $r:T_\ell J\to P$ with
$r|_P=\ell^k$ for some $k\ge0$.

Make a point of the compactified cover rational and use its
Abel--Jacobi map to the actual extension
$(\prod_qT_qJ)\rtimes G_L$. Restriction to $N$ and composition
with $r$ kills $C$ and embeds
\[
                   N/C\hookrightarrow P\rtimes G_L.
\]
Its image meets $P$ in $\ell^kP$ and surjects to $G_L$.
Modulo $\ell^kP$ it is the graph of a finite continuous
cocycle. Its zero set is an open subgroup. After restricting
there, subtracting the geometric coordinate of each lift gives
all $(0,g)$ in the image. Compactness makes the inverse onto
this closed image continuous, so $N/C$ has an actual section.
Let $B$ be its inverse image in $N$. Then $N=P\rtimes B$,
$B\cap\pi=C$, and (A.17.1) follows.

We give a shorter independent proof of the final assertion
than the infinite-index Hom argument in LM6 A.61. Suppose an
actual section $s:G_L\to N_\Pi(A)$ exists. Localize at one
finite place. D.56, with the formation of all finite groups
and its affine neighborhood consequence, makes this section
a point or cusp section.
Here is that consequence: use a cofinal system of characteristic
open geometric subgroups $N$ and the actual open arithmetic
subgroups $N\rtimes s(G_{L_v})$. These are section neighborhoods
over the same local field. One can first choose a member whose
compactification has genus at least two: take a characteristic
cover in which every cusp inertia has sufficiently large order,
and apply Riemann--Hurwitz. Apply the proper theorem to every
further member of this cofinal system. The
resulting compactification points are compatible, by injectivity
in the proper theorem. Their finite fibres above the first
point have nonempty compatible inverse limit. The section
fixes this lifted horizontal point; its stabilizer is the
ordinary decomposition group or the cuspidal decomposition
group according to whether the point is in the affine curve.
This uses full etale covers and does not replace them by
pro-prime covers.

Conjugation on the procyclic group $A$ has abelian automorphism
group. Thus the actual section's restriction to the local
commutator centralizes $A$. The local assertion proved in A.1
now gives $A=1$ for an ordinary point section, or $A\subset I$
for a cusp section. The former is excluded. Conversely a
rational tangent at a rational cusp gives a section normalizing
every closed subgroup of its inertia, after finite extension.
This proves the equivalence. Applying it to (A.17.1) proves
the stated failure of virtual splitting and $C\ne1$.
Only the last part of the proof uses the local section
theorem, supplied here by D.56. The structural quotient
splitting does not use it, and neither part uses the
global density inference discussed in A.20. Thus this
assertion no longer assumes the main theorem of the
supplied local preprint.
\end{proof}

\begin{statement}{A.18}{Proved}{Why the abstract kernel and local splitting data do not suffice}
For every prime $\ell$ and an odd prime $q\ne\ell$, there is
an extension of $G_{\Q(i)}$ by $\Z_q(1)$ with no section over
any finite extension. It can be chosen split at every place
above $\ell$. For a split $\ell$, adjoining the semidirect
factor $\Z_\ell(\psi)$ for a CM character gives an abstract
configuration with geometric kernel $\Z_\ell\times\Z_q(1)$,
which itself embeds as a closed subgroup of the free profinite
geometric group. No embedding of the whole arithmetic
configuration in $\Pi$ is asserted. This proves the obstruction
claimed in LM6 A.63 and isolates the additional embedding issue.
\end{statement}
\begin{proof}
Choose two finite places $v,w$ away from $\ell$ and $q$.
The Brauer invariant-sum sequence supplies compatible classes
$\beta_n\in H^2(G_K,\mu_{q^n})$ with local invariants
$q^{-n}$ and $-q^{-n}$ at $v,w$, and zero elsewhere.
The coefficient map $\mu_{q^{n+1}}\to\mu_{q^n}$ corresponds
to multiplication by $q$ on Brauer classes, so these classes
are compatible. Starting with a continuous normalized cocycle
representing $\beta_n$, choose one for $\beta_{n+1}$; their
difference after reduction is a coboundary of a finite-valued
continuous one-cochain. Lift that cochain value by value and
subtract its coboundary. Induction gives compatible actual
cocycles, whose inverse limit is a continuous
$z:G_K^2\to\Z_q(1)$.

On the product space define
\[
              (a,g)(b,h)=(a+g b+z(g,h),gh).
                                                               \tag{A.18.1}
\]
The cocycle identity gives a profinite group $B$. At a place
$v'$ above $v$ of a finite extension $L/K$, restriction of
$\beta_n$ has invariant $[L_{v'}:K_v]/q^n$, nonzero for
sufficiently large $n$. A section of $B\to G_K$ over $G_L$
would make all these finite reductions coboundaries, a
contradiction. This verifies continuous nonsplitting directly;
it does not infer a continuous obstruction from an unspecified
inverse limit of vanishing finite obstructions.

At a place above $\ell$ all the finite local classes are zero.
The local finite groups $H^1(G_{K_v},\mu_{q^n})$ satisfy
Mittag--Leffler, so the continuous comparison has no
$\varprojlim{}^1$ kernel. Thus the continuous local class is
zero and the extension splits there.
If $\ell$ splits in $K$, choose one CM Tate character $\psi$
of $E$ and put $N=\Z_\ell(\psi)\rtimes B$, with $B$ acting
through its quotient $G_K$. Its geometric kernel is the direct
product of the two different prime factors. Abstractly this is
a closed procyclic subgroup of $\widehat\Z$, hence embeds in
a free profinite generator subgroup and is projective. The
action, centralization and local splitting properties therefore
do not alone rule out the nonsplitting. The missing condition
is an embedding of this entire extension, with its arithmetic
projection and action, in the particular curve group $\Pi$.
\end{proof}

\begin{statement}{A.23}{Proved}{The proper-curve version of the audit}
Every finitely generated pro-$\ell$ subgroup of the full
geometric group of a proper hyperbolic curve is free pro-$\ell$.
The CM and prime-support restrictions of A.16 and the
prime-complement configuration of A.17 consequently hold
also for proper hyperbolic curves. There is no boundary
alternative. By the proper full-group case of D.56, their
normalizers cannot virtually split. These are the assertions
of LM6 A.64--A.65 with the dependencies separated.
\end{statement}
\begin{proof}
Let $Q$ be such a finitely generated pro-$\ell$ subgroup of
a proper surface group $\Pi_g$, $g\ge2$.
Continuous finite-coefficient cohomology of $Q$ is the direct
limit of that of open overgroups $V$ in $\Pi_g$.
The subgroup $Q$ cannot be open: every open surface group has
nontrivial abelian quotients at primes other than $\ell$.
If the supernatural index $[\Pi_g:Q]$ had finite $\ell$-adic
valuation, choose an open overgroup $V$ for which the
$\ell$-part of $[\Pi_g:V]$ is maximal, equivalently the
remaining $\ell$-part of $[V:Q]$ is minimal. All further
restriction maps toward $Q$ then have prime-to-$\ell$ degree
on a cofinal subsystem and are injective on $H^1(-,\F_\ell)$
by corestriction. The genera of these covers are unbounded,
so $H^1(Q,\F_\ell)$ would have unbounded dimension, contrary
to finite generation. Thus $[\Pi_g:Q]$, and hence every
$[V:Q]$, has infinite $\ell$-adic valuation.
On $H^2(V,\F_\ell)$ restriction is
multiplication by the covering degree, hence is eventually
zero. Thus $H^2(Q,\F_\ell)=0$, which is the free pro-$\ell$
criterion.

For the Sylow step of A.17 no finite generation of the whole
Sylow group is needed. Each subgroup generated by $A_\ell$
and one additional Sylow element is finitely generated,
free pro-$\ell$ by the preceding argument, and has central
nonzero $A_\ell$ after the character shrinking. It is cyclic.
Every finite image of the whole Sylow group is therefore
cyclic, so the Sylow group is procyclic. The rest of the
prime-complement and semisimple-retraction construction is
unchanged, with proper homology as the detector from the
outset. Finally a local section of a proper curve is an
ordinary point section by D.56. Its local
derived centralizer is trivial by A.1, contradicting any
nonzero normalized procyclic subgroup. This proves the last
assertion with the stated internal theorem dependency.
\end{proof}

\begin{statement}{A.19}{Proved}{Exactly what the two local kernels differ by}
Let $v\mid p$ be a finite place of $K=\Q(i)$ and put $k=K_v$.
Under a compatible decomposition-group embedding,
\[
 D_v\cap\Delta=\overline{[G_k,G_k]},\qquad
 \frac{G_{k(\mu_\infty)}}{\overline{[G_k,G_k]}}
 \simeq\ker\bigl(N_{k/\Q_p}:\mathcal O_k^{\times}
                                  \to\Z_p^{\times}\bigr).
                                                               \tag{A.19.1}
\]
The second group is trivial at split primes, is noncanonically
$\mu_{p+1}\times\Z_p$ at odd inert primes, and is
$\mu_4\times\Z_2$ at the ramified prime two.
For a finite extension $L/K$ and a place $w\mid v$, the kernel
using the original $\Delta$ is
$G_{L_w}\cap\overline{[G_k,G_k]}$; the quotient lost by passing
from $G_{L_w(\mu_\infty)}$ to this kernel is an open subgroup
of the norm kernel in (A.19.1).
\end{statement}
\begin{proof}
The easy containment of the local commutator in $\Delta$ follows
from abelianization. For equality, every finite abelian extension
of $k$ occurs as the completion of a finite abelian extension of
$K$: the Grunwald--Wang theorem has no exceptional condition here
because $i\in K$. See \cite[Theorems 1--2]{A_LorenzRoquette}.
Thus the map $G_k^{\rm ab}\to G_K^{\rm ab}$ is injective,
which gives the first equality. The same local decomposition
identification is stated in \cite[Chapter III, \S2.1]{A_Serre}.

Local reciprocity identifies $G_k^{\rm ab}$ with the profinite
completion of $k^{\times}$. Its action on all roots of unity
is the restriction map to $G_{\Q_p}^{\rm ab}$, identified by
local Kronecker--Weber with all the local cyclotomic characters.
Reciprocity is compatible with norms. Hence the quotient in
(A.19.1) is the kernel of the completed norm
$\widehat{k^{\times}}\to\widehat{\Q_p^{\times}}$.
On valuation factors the norm is multiplication by the residue
degree, injective on $\widehat\Z$. Its kernel is therefore
exactly the displayed unit norm kernel, with no valuation factor.

For $k=\Q_p$ it is zero. For the unramified quadratic extension
at odd $p$, the norm on residue roots of unity has kernel
$\mu_{p+1}$, and on principal units the logarithm identifies
the norm with trace on a rank-two $\Z_p$-lattice, giving a
rank-one kernel. For $k=\Q_2(i)$ the kernel has torsion
exactly $\mu_4$, and its principal-unit part has rank one,
again by the logarithm on a sufficiently small open subgroup.
The classification of finitely generated abelian pro-$2$ groups
gives the stated noncanonical product.

Finally the chosen local embedding for $L_w$ is an open subgroup
of $G_k$. Intersecting the first equality in (A.19.1) gives the
claim about the original $\Delta$. The image of an open subgroup
under a surjective homomorphism of profinite groups is open,
which proves the last assertion. If one instead replaces
$\Delta$ by the new global commutator of $G_L$, its local
intersection is $G_{L_w}'$ and the principal norm-kernel rank
is $[L_w:\Q_p]-1$; these are different comparisons.

The loss is consequently explicit. In A.1 it loses no uniqueness
hypothesis: the smaller local derived subgroup already has a
nonzero tame $\Z_r$ quotient for every $r\ne p$. For an arbitrary
normalizer, replacing the larger local kernel by the smaller one
still does not supply a point or cusp section to which the
centralizer theorem could be applied.
\end{proof}

\begin{statement}{A.20}{Refuted}{The density intersection inference, and its correct attribution}
The inference that a set of strictly positive Dirichlet density
must meet a set of potentially density one is false. The
counterexample below refutes that inference only. It does not
refute a section theorem or a curve finiteness theorem.
The attached local section-conjecture preprint's Section 8
does not contain the alleged inference or a Corollary 7.10.
LM6 D.6 attributes that step to the other 2026 source, Porowski's
unnumbered corollary following Corollary 7.10.
\end{statement}
\begin{proof}
Let $L/K$ be a nontrivial quadratic extension. The unramified
inert primes of $K$ have density $1/2$. The completely split
primes also have density $1/2$, and are disjoint from them.
But the pullback of the completely split set to $L$ has density
one: primes over inert primes have relative residue degree two
and contribute a bounded series at the Dirichlet-density pole.
Thus the split set has potentially density one and does not meet
the positive-density inert set.

Replacing strictly positive by strongly positive density repairs
this particular inference: by definition its pullback has
positive upper density after every finite extension, so it
meets any set that becomes density one in some such extension.
This is exactly the hypothesis retained by the valid uniform
reduction argument in LM6 D.6. No counterexample to the stronger
arithmetic conclusion is inferred from the failed weaker proof.

We checked the whole supplied file
\texttt{sections/08-global-sections.tex}. It proceeds by
Harari--Stix finite descent, Stix's finite-support result and
Stoll's finite-descent criteria; it explicitly retains the
finite-cover equality as a hypothesis. It has no use of the
displayed false intersection step. The local Theorem 1.1 is
proved before that section. Consequently this source mismatch
does not invalidate that local input. The arguments A.1, A.3,
A.6, A.11--A.15, and A.19 do not use either the local section
theorem or the disputed global density inference.
\end{proof}

